% GENERATED FILE -- DO NOT EDIT.
%
% Produced by harness/assemble_contract.py from module.yaml and claims/*.tex.
% Edit the claim files and regenerate; `--check` fails if this file is stale.
\documentclass[11pt,a4paper]{article}
\pdftrailerid{a65a4f3477e6c083ec8b03ebfa3d06e7}
%% --- 1. Mathematics (order matters: physics after braket wins \ket, \tr) ---
\usepackage{amsmath}
\usepackage{amssymb}
\usepackage{amsfonts}
\usepackage{amsthm}
\usepackage{mathtools}
\usepackage{dsfont}
\usepackage{braket}
\usepackage{physics}
%% --- 2. Page furniture ------------------------------------------------------
\usepackage{xcolor}
\usepackage{array}    % >{...} column prefixes in the wide tables of some claims
\usepackage{enumitem}
\usepackage{graphicx}
\usepackage{booktabs}
\usepackage{longtable}
\usepackage{xltabular}              % front- and back-matter tables across pages
\usepackage{needspace}
\usepackage{xurl}
\usepackage[noautomatic]{imakeidx}  % defined-terms index; latexmk runs makeindex
\makeindex[title=Index of defined terms,intoc]
\usepackage{tcolorbox}
\tcbuselibrary{skins,breakable}
%% --- 3. Links and navigation (hyperref late; bookmark, refcount after it) ---
\usepackage{hyperref}
\usepackage{bookmark}
\usepackage{refcount}
%% --- 4. Colours: every distinction also has a shape (rule, caps, size, weight)
\definecolor{nnGrey}{gray}{0.40}      % non-normative text
\definecolor{metaGrey}{gray}{0.30}    % the generated meta line
\definecolor{techGrey}{gray}{0.45}    % heading of a technical claim
\definecolor{mainRule}{RGB}{0,110,140}
\definecolor{reqAccent}{RGB}{165,70,0}
\definecolor{linkBlue}{RGB}{0,62,150}
\definecolor{extLink}{RGB}{120,40,140}  % links into another PDF
\hypersetup{colorlinks=true, linkcolor=linkBlue, citecolor=linkBlue,
  urlcolor=linkBlue, filecolor=extLink, bookmarksnumbered=true,
  bookmarksopen=true, bookmarksopenlevel=1, pdfdisplaydoctitle=true}
% A link into another PDF (\href to a .pdf, \hyperref{file}..., xr-hyper) is a
% relative URI, <file>#nameddest=<destination>, in the file-link colour:
% browser viewers (Edge, Chrome) ignore hyperref's default GoToR action.
\makeatletter
\renewcommand*{\hyper@linkfile}[3]{% anchor text, file name, destination
  \begingroup
    \ifx\\#3\\\def\CQ@target{#2}\else\def\CQ@target{#2\hyper@hash nameddest=#3}\fi
    \let\@urlcolor\@filecolor
    \hyper@linkurl{#1}{\CQ@target}%
  \endgroup}
\makeatother
%% --- 5. The mapped Unicode set (harness/contract_navigation.py, UNICODE_TEX)
\DeclareUnicodeCharacter{03B1}{\ensuremath{\alpha}}
\DeclareUnicodeCharacter{03B2}{\ensuremath{\beta}}
\DeclareUnicodeCharacter{03B3}{\ensuremath{\gamma}}
\DeclareUnicodeCharacter{03B4}{\ensuremath{\delta}}
\DeclareUnicodeCharacter{03B5}{\ensuremath{\varepsilon}}
\DeclareUnicodeCharacter{03F5}{\ensuremath{\epsilon}}
\DeclareUnicodeCharacter{03B6}{\ensuremath{\zeta}}
\DeclareUnicodeCharacter{03B7}{\ensuremath{\eta}}
\DeclareUnicodeCharacter{03B8}{\ensuremath{\theta}}
\DeclareUnicodeCharacter{03D1}{\ensuremath{\vartheta}}
\DeclareUnicodeCharacter{03B9}{\ensuremath{\iota}}
\DeclareUnicodeCharacter{03BA}{\ensuremath{\kappa}}
\DeclareUnicodeCharacter{03BB}{\ensuremath{\lambda}}
\DeclareUnicodeCharacter{03BC}{\ensuremath{\mu}}
\DeclareUnicodeCharacter{03BD}{\ensuremath{\nu}}
\DeclareUnicodeCharacter{03BE}{\ensuremath{\xi}}
\DeclareUnicodeCharacter{03C0}{\ensuremath{\pi}}
\DeclareUnicodeCharacter{03C1}{\ensuremath{\rho}}
\DeclareUnicodeCharacter{03F1}{\ensuremath{\varrho}}
\DeclareUnicodeCharacter{03C3}{\ensuremath{\sigma}}
\DeclareUnicodeCharacter{03C2}{\ensuremath{\varsigma}}
\DeclareUnicodeCharacter{03C4}{\ensuremath{\tau}}
\DeclareUnicodeCharacter{03C5}{\ensuremath{\upsilon}}
\DeclareUnicodeCharacter{03C6}{\ensuremath{\varphi}}
\DeclareUnicodeCharacter{03D5}{\ensuremath{\phi}}
\DeclareUnicodeCharacter{03C7}{\ensuremath{\chi}}
\DeclareUnicodeCharacter{03C8}{\ensuremath{\psi}}
\DeclareUnicodeCharacter{03C9}{\ensuremath{\omega}}
\DeclareUnicodeCharacter{0393}{\ensuremath{\Gamma}}
\DeclareUnicodeCharacter{0394}{\ensuremath{\Delta}}
\DeclareUnicodeCharacter{0398}{\ensuremath{\Theta}}
\DeclareUnicodeCharacter{039B}{\ensuremath{\Lambda}}
\DeclareUnicodeCharacter{039E}{\ensuremath{\Xi}}
\DeclareUnicodeCharacter{03A0}{\ensuremath{\Pi}}
\DeclareUnicodeCharacter{03A3}{\ensuremath{\Sigma}}
\DeclareUnicodeCharacter{03A5}{\ensuremath{\Upsilon}}
\DeclareUnicodeCharacter{03A6}{\ensuremath{\Phi}}
\DeclareUnicodeCharacter{03A8}{\ensuremath{\Psi}}
\DeclareUnicodeCharacter{03A9}{\ensuremath{\Omega}}
\DeclareUnicodeCharacter{00E1}{\'a}
\DeclareUnicodeCharacter{00E0}{\`a}
\DeclareUnicodeCharacter{00E2}{\^a}
\DeclareUnicodeCharacter{00E4}{\"a}
\DeclareUnicodeCharacter{00E5}{\aa{}}
\DeclareUnicodeCharacter{00E9}{\'e}
\DeclareUnicodeCharacter{00E8}{\`e}
\DeclareUnicodeCharacter{00EA}{\^e}
\DeclareUnicodeCharacter{00EB}{\"e}
\DeclareUnicodeCharacter{00ED}{\'{\i}}
\DeclareUnicodeCharacter{00EF}{\"{\i}}
\DeclareUnicodeCharacter{00F3}{\'o}
\DeclareUnicodeCharacter{00F2}{\`o}
\DeclareUnicodeCharacter{00F4}{\^o}
\DeclareUnicodeCharacter{00F6}{\"o}
\DeclareUnicodeCharacter{00F8}{\o{}}
\DeclareUnicodeCharacter{00FA}{\'u}
\DeclareUnicodeCharacter{00F9}{\`u}
\DeclareUnicodeCharacter{00FB}{\^u}
\DeclareUnicodeCharacter{00FC}{\"u}
\DeclareUnicodeCharacter{00F1}{\~n}
\DeclareUnicodeCharacter{00E7}{\c{c}}
\DeclareUnicodeCharacter{00DF}{\ss{}}
\DeclareUnicodeCharacter{00C4}{\"A}
\DeclareUnicodeCharacter{00C5}{\AA{}}
\DeclareUnicodeCharacter{00C9}{\'E}
\DeclareUnicodeCharacter{00D6}{\"O}
\DeclareUnicodeCharacter{00DC}{\"U}
\DeclareUnicodeCharacter{2013}{\textendash{}}
\DeclareUnicodeCharacter{2014}{\textemdash{}}
\DeclareUnicodeCharacter{2018}{`}
\DeclareUnicodeCharacter{2019}{'}
\DeclareUnicodeCharacter{201C}{``}
\DeclareUnicodeCharacter{201D}{''}
\DeclareUnicodeCharacter{2026}{\ldots{}}
\DeclareUnicodeCharacter{00A0}{~}
\DeclareUnicodeCharacter{00A7}{\S{}}
\DeclareUnicodeCharacter{2264}{\ensuremath{\le}}
\DeclareUnicodeCharacter{2265}{\ensuremath{\ge}}
\DeclareUnicodeCharacter{2260}{\ensuremath{\neq}}
\DeclareUnicodeCharacter{2248}{\ensuremath{\approx}}
\DeclareUnicodeCharacter{2261}{\ensuremath{\equiv}}
\DeclareUnicodeCharacter{223C}{\ensuremath{\sim}}
\DeclareUnicodeCharacter{2243}{\ensuremath{\simeq}}
\DeclareUnicodeCharacter{2245}{\ensuremath{\cong}}
\DeclareUnicodeCharacter{221D}{\ensuremath{\propto}}
\DeclareUnicodeCharacter{226A}{\ensuremath{\ll}}
\DeclareUnicodeCharacter{226B}{\ensuremath{\gg}}
\DeclareUnicodeCharacter{2192}{\ensuremath{\to}}
\DeclareUnicodeCharacter{2190}{\ensuremath{\leftarrow}}
\DeclareUnicodeCharacter{2194}{\ensuremath{\leftrightarrow}}
\DeclareUnicodeCharacter{21D2}{\ensuremath{\Rightarrow}}
\DeclareUnicodeCharacter{21D0}{\ensuremath{\Leftarrow}}
\DeclareUnicodeCharacter{21D4}{\ensuremath{\Leftrightarrow}}
\DeclareUnicodeCharacter{21A6}{\ensuremath{\mapsto}}
\DeclareUnicodeCharacter{2208}{\ensuremath{\in}}
\DeclareUnicodeCharacter{2209}{\ensuremath{\notin}}
\DeclareUnicodeCharacter{2286}{\ensuremath{\subseteq}}
\DeclareUnicodeCharacter{2282}{\ensuremath{\subset}}
\DeclareUnicodeCharacter{2287}{\ensuremath{\supseteq}}
\DeclareUnicodeCharacter{2283}{\ensuremath{\supset}}
\DeclareUnicodeCharacter{2229}{\ensuremath{\cap}}
\DeclareUnicodeCharacter{222A}{\ensuremath{\cup}}
\DeclareUnicodeCharacter{2216}{\ensuremath{\setminus}}
\DeclareUnicodeCharacter{2205}{\ensuremath{\emptyset}}
\DeclareUnicodeCharacter{2200}{\ensuremath{\forall}}
\DeclareUnicodeCharacter{2203}{\ensuremath{\exists}}
\DeclareUnicodeCharacter{2115}{\ensuremath{\mathbb{N}}}
\DeclareUnicodeCharacter{2124}{\ensuremath{\mathbb{Z}}}
\DeclareUnicodeCharacter{211A}{\ensuremath{\mathbb{Q}}}
\DeclareUnicodeCharacter{211D}{\ensuremath{\mathbb{R}}}
\DeclareUnicodeCharacter{2102}{\ensuremath{\mathbb{C}}}
\DeclareUnicodeCharacter{00B7}{\ensuremath{\cdot}}
\DeclareUnicodeCharacter{00D7}{\ensuremath{\times}}
\DeclareUnicodeCharacter{2212}{\ensuremath{-}}
\DeclareUnicodeCharacter{00B1}{\ensuremath{\pm}}
\DeclareUnicodeCharacter{2213}{\ensuremath{\mp}}
\DeclareUnicodeCharacter{221A}{\ensuremath{\surd}}
\DeclareUnicodeCharacter{2211}{\ensuremath{\sum}}
\DeclareUnicodeCharacter{220F}{\ensuremath{\prod}}
\DeclareUnicodeCharacter{222B}{\ensuremath{\int}}
\DeclareUnicodeCharacter{221E}{\ensuremath{\infty}}
\DeclareUnicodeCharacter{2202}{\ensuremath{\partial}}
\DeclareUnicodeCharacter{2207}{\ensuremath{\nabla}}
\DeclareUnicodeCharacter{2297}{\ensuremath{\otimes}}
\DeclareUnicodeCharacter{2295}{\ensuremath{\oplus}}
\DeclareUnicodeCharacter{2218}{\ensuremath{\circ}}
\DeclareUnicodeCharacter{2308}{\ensuremath{\lceil}}
\DeclareUnicodeCharacter{2309}{\ensuremath{\rceil}}
\DeclareUnicodeCharacter{230A}{\ensuremath{\lfloor}}
\DeclareUnicodeCharacter{230B}{\ensuremath{\rfloor}}
\DeclareUnicodeCharacter{27E8}{\ensuremath{\langle}}
\DeclareUnicodeCharacter{27E9}{\ensuremath{\rangle}}
\DeclareUnicodeCharacter{2016}{\ensuremath{\|}}
\DeclareUnicodeCharacter{2113}{\ensuremath{\ell}}
\DeclareUnicodeCharacter{00B0}{\ensuremath{^{\circ}}}
\DeclareUnicodeCharacter{00B2}{\ensuremath{^{2}}}
\DeclareUnicodeCharacter{00B3}{\ensuremath{^{3}}}
\DeclareUnicodeCharacter{00B9}{\ensuremath{^{1}}}
\DeclareUnicodeCharacter{2080}{\ensuremath{_{0}}}
\DeclareUnicodeCharacter{2081}{\ensuremath{_{1}}}
\DeclareUnicodeCharacter{2082}{\ensuremath{_{2}}}
\DeclareUnicodeCharacter{207F}{\ensuremath{^{n}}}
\DeclareUnicodeCharacter{2096}{\ensuremath{_{k}}}
%% --- 6. Numbering: one counter per section -----------------------
\numberwithin{equation}{section}
\setcounter{tocdepth}{1}
%% --- 7. The project macro set, then the claim environments -----------------
% spec/macros.tex, read relative to this module's directory (latexmk -cd).
% =============================================================================
%  spec/macros.tex -- the one macro set every generated contract inputs.
%
%  Two parts: part A, the macros spec/notation.md catalogues, and part B,
%  the macros local to one module, grouped by module.
%  harness/assemble_contract.py emits % =============================================================================
%  spec/macros.tex -- the one macro set every generated contract inputs.
%
%  Two parts: part A, the macros spec/notation.md catalogues, and part B,
%  the macros local to one module, grouped by module.
%  harness/assemble_contract.py emits % =============================================================================
%  spec/macros.tex -- the one macro set every generated contract inputs.
%
%  Two parts: part A, the macros spec/notation.md catalogues, and part B,
%  the macros local to one module, grouped by module.
%  harness/assemble_contract.py emits \input{../../macros} for it, so the
%  file is read from spec/modules/<id>/.
%
%  RULES
%    * A macro that spec/notation.md catalogues is a verbatim copy of it
%      (part A). Changing one is a notation change: bump spec/notation.md.
%    * Part B holds the module-local macros. A macro defined by two modules
%      with the same expansion appears once, under the first module (in the
%      order below) that defines it. Two macros with the same name and
%      DIFFERENT expansions are both kept under a module prefix, and the
%      claims write the prefixed names:
%        \Chain  -> \redChain (reductions, \mathfrak{C}: the reduction chain)
%                   \pcChain  (protocol-compiler, \mathfrak{R}: the two-segment
%                             chain; \Chsrc and \Chrec are built on it)
%        \Nmsm   -> \redNmsm  (reductions, \mathcal{N}: a POVM measurement
%                             channel)
%                   \pcNmsm   (\mathcal{N}_{\mathsf A}: Alice's measurement
%                             channel, written by
%                             transmitter/thm:certificate-sufficiency,
%                             transmitter/rem:xreg and
%                             reductions/lem:source-transfer)
%    * Deliberately omitted: a \provenance{tag}{details} macro. Provenance
%      is a header field of every claim file (spec/claim-format.md section
%      3); a claim body that calls such a macro fails to compile, which is
%      the intended signal.
%
%  REQUIRED PACKAGES, loaded by the assembler's preamble before this file:
%    amsmath, amssymb, amsfonts, mathtools, dsfont, braket, physics, xcolor
%    (braket before physics, so that physics wins for \ket, \bra, \tr, ...).
%
%  This file contains only definitions: no \usepackage, no document body.
% =============================================================================


% =============================================================================
%  PART A -- notation macros, the ones spec/notation.md catalogues, each a
%  verbatim copy of its catalogue row; grouped under numbered headings.
% =============================================================================

% -----------------------------------------------------------------------------
% 1. State spaces and operator sets (a "system label" -> spaces of operators)
% -----------------------------------------------------------------------------
\newcommand{\Hs}[1]{\mathcal{H}(#1)}        % Hilbert space of a system label
\newcommand{\Lin}[1]{\mathcal{L}(#1)}       % bounded linear operators on H(#1)
\newcommand{\Linn}[2]{\mathcal{L}(#1,#2)}   % linear maps H(#1) -> H(#2)
\newcommand{\Pos}[1]{\mathrm{Pos}(#1)}      % positive semidefinite operators on H(#1)
\newcommand{\Density}{\mathrm{D}}           % density operators: \Density(A) = D(A)
\newcommand{\Seq}[1]{S_{=}(#1)}             % normalized states (trace = 1)
\newcommand{\Sleq}[1]{S_{\leq}(#1)}         % subnormalized states (0 <= trace <= 1)
\newcommand{\id}{\mathds{1}}                % identity operator
\newcommand{\idon}[1]{\mathds{1}_{#1}}      % identity on a named system, e.g. \idon{C}

% DEPRECATED: applies \mathcal to both arguments itself, so it takes bare
% letters, not system labels. Prefer \Linn; do not use in new contracts.
\providecommand{\Lop}[2]{\mathcal{L}(\mathcal{#1}, \mathcal{#2})}


% -----------------------------------------------------------------------------
% 2. Composite systems
% -----------------------------------------------------------------------------
\newcommand{\dsum}{\oplus}                  % direct sum of system labels
% Tensor product: use \otimes directly.
% n-fold direct sum of an operator space: \Lin{A}^{\oplus n}


% -----------------------------------------------------------------------------
% 3. Events, conditioning and classical registers (for STATES)
%    Convention:  \wedge  -> subnormalized "and-event" component, e.g. \rho_{A\wedge x}
%                 |        -> normalized conditional state,        e.g. \rho_{A|x}
% (Write these directly in subscripts; no extra macros needed.)
% -----------------------------------------------------------------------------


% -----------------------------------------------------------------------------
% 4. Photon-number subspaces / direct-sum blocks
%    Convention:  operators carry the restriction "on top", in brackets;
%                 states carry it as a \wedge subscript (see Sec. 3).
% -----------------------------------------------------------------------------
\newcommand{\Pnum}[1]{P_{#1}}               % photon-number projector (system from context)
\newcommand{\Pnumon}[2]{P_{#1}^{#2}}        % with explicit system superscript, e.g. P_n^C
\newcommand{\blk}[2]{#1^{(#2)}}             % operator block, e.g. \blk{M}{\leqN} = M^{(<=N)}
\newcommand{\leqN}{\leq N}                  % use as \blk{M}{\leqN}
\newcommand{\gtN}{> N}                      % use as \blk{M}{\gtN}
\newcommand{\eqn}{= n}                      % use as \blk{M}{\eqn}


% -----------------------------------------------------------------------------
% 5. System labels used in the protocol
% -----------------------------------------------------------------------------
\newcommand{\Chan}{C}                       % quantum-channel (Eve) input to Bob
\newcommand{\Prb}{P}                        % probe input
\newcommand{\Flag}{{\mathrm{Flg}}}          % flag system (EXISTS ONLY AFTER squashing)
\newcommand{\Loss}{{\mathrm{Lss}}}          % loss system
\newcommand{\Vc}{{\mathrm{Vcm}}}            % vacuum system
\newcommand{\Nout}{N_{\mathrm{out}}}
\newcommand{\Nin}{N_{\mathrm{in}}}


% -----------------------------------------------------------------------------
% 6. Alphabets
% -----------------------------------------------------------------------------
\newcommand{\Yset}{\mathcal{Y}}             % Bob's (coarse-grained) detection alphabet
\newcommand{\Yraw}{\mathcal{Y}_{\mathrm{raw}}}% raw click patterns {0,1}^4
\newcommand{\Yscset}{\mathcal{Y}_{\singlec}}% single-click outcomes
\newcommand{\Ymcset}{\mathcal{Y}_{\mc}}     % multi-click outcomes
\newcommand{\Tset}{\mathcal{T}_{\mathrm{set}}}% private setting set (the probe settings in the optical instance)
\newcommand{\Xset}{\mathcal{X}}             % Alice's alphabet
\newcommand{\Dset}{\{H,V,+,-\}}             % polarization / detector labels
\newcommand{\pass}{\top}                    % pass sifting


% -----------------------------------------------------------------------------
% 7. Mode / outcome abbreviations
% -----------------------------------------------------------------------------
\newcommand{\on}{{\mathrm{on}}}
\newcommand{\off}{{\mathrm{off}}}
\newcommand{\dc}{{\mathrm{dc}}}             % dark count
\newcommand{\sq}{\mathrm{sq}}               % squashed
\newcommand{\mc}{\mathrm{mc}}               % multi-click
\newcommand{\nc}{\mathrm{nc}}               % no-click
\newcommand{\singlec}{\mathrm{sc}}          % single-click
\newcommand{\vac}{\mathrm{vac}}             % vacuum
\newcommand{\kout}{{k_\mathrm{out}}}
\newcommand{\kin}{{k_\mathrm{in}}}


% -----------------------------------------------------------------------------
% 8. Maps, channels and miscellaneous operators
% -----------------------------------------------------------------------------
\newcommand{\Prob}{\mathrm{Prob}}           % probability (use \Prob[\cdot] or \Prob(\cdot))
\newcommand{\lmin}{\lambda_{\min}}          % minimum eigenvalue
\newcommand{\Wflag}{W_{\Flag}}              % flag-space weight
\newcommand{\transpose}{\mathsf{T}}
\newcommand{\innerp}[2]{\left\langle #1 \vert #2 \right\rangle}
\renewcommand{\vec}[1]{\boldsymbol{#1}}     % NB: overrides the LaTeX arrow accent


% -----------------------------------------------------------------------------
% 9. Key-rate numerics
% -----------------------------------------------------------------------------
\newcommand{\cG}{\mathcal{G}}               % protocol map
\newcommand{\cZ}{\mathcal{Z}}               % pinching map
\newcommand{\Amap}{\mathcal{A}_{\mathrm{ann}}} % announcement + measurement map
\newcommand{\Amapnb}{\mathcal{A}_{\mathrm{ann}|\mathrm{nb}}}
\newcommand{\simpl}{}                       % empty suffix for simplified block maps
\newcommand{\acSet}{\mathcal{S}_{\mathrm{acc}}} % acceptance set
\newcommand{\testE}{\mathcal{E}_{\mathrm{test}}} % test-statistics map
\newcommand{\ecLeak}{\mathrm{leak}^{\mathrm{EC}}_{\mathrm{obs}}}
\newcommand{\pgen}{p_{\mathrm{gen}}}
\newcommand{\askeyr}{R_{\infty}}             % asymptotic key rate
\newcommand{\ppass}{p_{\mathrm{pass}}}       % sifting pass probability


% -----------------------------------------------------------------------------
% 10. Finite-size analysis / Renyi entropies
% -----------------------------------------------------------------------------

% --- Round types: test vs. key generation (register T, alphabet \Fset) --------
\newcommand{\tst}{\mathrm{test}}            % test round
\newcommand{\gen}{\mathrm{gen}}             % key-generation round
\newcommand{\Fset}{\mathcal{T}}             % round-type alphabet {test, gen} (not \Tset)
\newcommand{\Chat}{\widehat{C}}             % acceptance-test register
\newcommand{\Chatset}{\widehat{\mathcal{C}}}% its alphabet

% --- Basis / key extraction functions ----------------------------------------
\newcommand{\bA}{b_{A}}                     % Alice's basis-extraction function
\newcommand{\bB}{b_{B}}                     % Bob's basis-extraction function
\newcommand{\kB}{k_{B}}                     % Bob's key-extraction function

% --- Maps --------------------------------------------------------------------
\newcommand{\Mprot}{\mathcal{M}}            % single-round protocol map
\newcommand{\Phimap}{\boldsymbol{\Phi}}     % test-statistics map
\newcommand{\Pdeph}{\mathcal{P}}            % dephasing (pinching) map on \Blreg

% --- Isometries --------------------------------------------------------------
\newcommand{\Vmsm}{V_{\mathrm{msm}}}        % measurement isometry
\newcommand{\Vann}{V_{\mathrm{ann}}}        % announcement isometry
\newcommand{\VI}{V_{I}}                     % Eve's copy of the announcements
\newcommand{\VChat}{V_{\Chat}}              % acceptance-data isometry
\newcommand{\VR}{V_{R}}                     % sifting + key map -> raw key R
\newcommand{\VZ}{V_{Z}}                     % raw key -> key register Z
\newcommand{\Wprot}{W}                      % full protocol isometry

% --- Renyi entropies ---------------------------------------------------------
% SECURITY-CRITICAL: see spec/notation.md §4. Never change an arrow to
% make an inequality work out; never convert silently between the sandwiched
% (tilde) and Petz (overline) families.
\newcommand{\Hsandup}{\widetilde{H}^{\uparrow}_{\alpha}}   % sandwiched, up
\newcommand{\Hsanddown}{\widetilde{H}^{\downarrow}_{\alpha}}% sandwiched, down
\newcommand{\Hpetzup}{\overline{H}^{\uparrow}_{\alpha}}    % Petz, up
\newcommand{\Hpetzdown}{\overline{H}^{\downarrow}_{\alpha}}% Petz, down
\newcommand{\hup}{h^{\uparrow}_{\alpha}}                   % single-round rate
\newcommand{\hqkd}{h^{\uparrow,\mathrm{QKD}}_{\alpha}}     % single-round QKD rate
\newcommand{\hlow}{\underline{h}_{\alpha}}                 % lower bound on \hqkd
% --- spec/notation.md 3.5 and 4.5 -------------------------------------------
% \hlb  is the first-order linearization lower bound on \hlow: the headline
%       key rate (the `first-order` tier of foundations/def:tiers). NOT \hlow
%       itself.
% \tnum is the instance's one declared numerical tolerance. Its
%       near-collision with the declared marginal \tau_A is disclosed in
%       spec/notation.md 3.5; bare \tau is reserved and denotes neither.
\newcommand{\hlb}{\underline{h}^{\mathrm{FO}}_{\alpha}}    % first-order lower bound on \hlow
\newcommand{\tnum}{\tau_{\mathrm{num}}}                    % declared numerical tolerance


% =============================================================================
%  PART B -- module-local macros, grouped by module. Deprecated aliases,
%  defined so that a text written with them compiles, used in no new contract: \Lop,
%  \testE, \Pdeph, \pass, \ecLeak (spec/notation.md section 3.9).
% =============================================================================

% --- meat-core -----------------------------------------------------------------
\newcommand{\SecSet}{\mathbb{S}}            % sector index set of Q
\newcommand{\Qsys}{Q}                       % declared sector-carrying measured system
\newcommand{\Vset}{\mathcal{V}}             % declared joint outcome alphabet of the Q-POVM
\newcommand{\CAset}{\mathcal{C}_{A}}        % Alice announcement alphabet
\newcommand{\CBset}{\mathcal{C}_{B}}        % Bob announcement alphabet
\newcommand{\Party}{\mathfrak{P}}           % declared party partition
\newcommand{\Oset}[1]{\mathcal{O}_{#1}}     % readout alphabet of party #1
\newcommand{\rd}[1]{\pi_{#1}}               % readout map of party #1
\newcommand{\Sadm}{\mathcal{S}_{A}}         % admissible Alice subspace
\newcommand{\nomarg}{\mathsf{none}}         % "no trusted marginal declared"
\newcommand{\Zset}{\mathcal{Z}_{\perp}}     % raw-key alphabet incl. \perp
\newcommand{\Zval}{\mathcal{Z}}             % key alphabet without \perp
\newcommand{\Etil}{\widetilde{E}}           % single-round purifying register
\newcommand{\Ehat}{\widehat{E}}             % Eve's n-round register
\newcommand{\OmAT}{\Omega_{\mathrm{AT}}}    % acceptance event
\newcommand{\epa}{\varepsilon_{\mathrm{PA}}}
\newcommand{\eev}{\varepsilon_{\mathrm{EV}}}
\newcommand{\lEC}{\lambda_{\mathrm{EC}}}
\newcommand{\Decl}{\mathfrak{D}}            % protocol declaration
\newcommand{\CList}{\mathfrak{L}}           % declared constraint list
\newcommand{\PinS}{\mathcal{P}_{\SecSet}}   % sector pinching on Q
\newcommand{\Feas}{\mathcal{W}}             % feasible set of block states
\newcommand{\Chkey}{\Chatset_{\mathrm{key}}}% key-live acceptance values
\newcommand{\Kann}{\mathcal{K}}             % announcement-class alphabet
\newcommand{\hold}{\mathsf{h}}              % declared key holder (A or B)
\newcommand{\sfD}{D}                        % classical KL divergence (base 2)
\newcommand{\Dtwo}{D}                       % quantum relative entropy (base 2)
\newcommand{\Dsand}{\widetilde{D}_{\alpha}} % sandwiched Renyi divergence
\newcommand{\idmap}{\mathrm{id}}            % identity channel
\newcommand{\Aev}{\mathcal{A}}              % Eve's per-round attack channel
\newcommand{\AttSet}{\mathfrak{A}}          % licensed attack-channel set
\newcommand{\Chacc}{\Chatset_{\mathrm{acc}}}% acceptance support of \acSet
\newcommand{\epsstar}{\varepsilon^{\star}}      % distance to an exact POVM
\newcommand{\dfH}{\Delta_{\mathrm{herm}}}       % Hermiticity defect
\newcommand{\dfP}{\Delta_{\mathrm{pos}}}        % negative part
\newcommand{\dfC}{\Delta_{\mathrm{cpl}}}        % completeness defect
\newcommand{\lmoq}{g}                       % acceptance-side gradient block
\newcommand{\lmoJ}{G}                       % feasible-side gradient block
\newcommand{\bacc}{\beta_{\mathrm{acc}}}    % certified acceptance-side minimum
\newcommand{\Itset}{\mathfrak{I}}           % finite index set of iterates (C22)

% --- reductions ----------------------------------------------------------------
\newcommand{\Bsys}{B}                        % physical (pre-squashing) system
\newcommand{\Tcl}[1]{\mathcal{T}(#1)}        % trace-class operators on H(#1)
\newcommand{\Ins}{\mathcal{I}}               % generalized-instrument branch map
\newcommand{\redNmsm}{\mathcal{N}}              % POVM measurement channel
\newcommand{\Eprot}{\mathcal{E}}             % n-round protocol channel
\newcommand{\Grnd}{\mathsf{G}}               % one round map of a protocol
\newcommand{\Kfin}{\mathsf{K}}               % end-of-protocol processing
\newcommand{\Crit}{\mathsf{sec}}             % value of a security criterion
\newcommand{\Pflag}{\Pi_{\Flag}}             % flag-sector projector on Q
\newcommand{\Pleq}{\Pi_{\leqN}}              % preserved-block projector
\newcommand{\Pgt}{\Pi_{\gtN}}                % above-cutoff projector
\newcommand{\Sph}{S^{\mathrm{ph}}}           % family of physical inputs
\newcommand{\Srch}{S^{\mathrm{rch}}}         % reachable squashed states
\newcommand{\Ssq}{S^{\sq}}                   % squashed admissible set
\newcommand{\Ffss}{F^{\mathrm{FSS}}}         % flag-state-squasher target
\newcommand{\Fwp}{F^{\mathrm{WP}}}           % multi-bound WPFSS target
\newcommand{\Mwp}{M^{\mathrm{WP}}}           % renormalized above-cutoff POVM
\newcommand{\Mdec}{{M'}^{\mathrm{ph}}}       % declared physical family on B
\newcommand{\Ppvm}{P^{\Flag}}                % flag PVM
\newcommand{\Lmp}{{\Lambda_{\mathrm{mp}}}}   % measure-and-prepare channel
\newcommand{\Lwp}{{\Lambda^{\mathrm{WP}}}}   % multi-bound WPFSS channel
\newcommand{\Lfss}{{\Lambda^{\mathrm{FSS}}}} % flag-state-squasher channel
\newcommand{\lmax}{\lambda_{\max}}           % maximum eigenvalue
\newcommand{\Lnoise}{\widetilde{\Lambda}}    % flag-noise channel
\newcommand{\cg}{\varpi}                     % declared coarse-graining map
\newcommand{\MCK}{\mathsf{MC}}               % multiclick raw-pattern set
\newcommand{\redChain}{\mathfrak{C}}            % reduction chain
\newcommand{\Lnk}{\ell}                      % a single typed link
\newcommand{\dsq}{\preceq}                      % the squashing relation
\newcommand{\opn}[1]{\left\lVert #1\right\rVert}           % operator (spectral) norm
\newcommand{\mxn}[1]{\left\lVert #1\right\rVert_{\max}}    % entrywise max modulus
\newcommand{\Fld}{\mathbb{F}}                   % exact ordered substrate field
\newcommand{\Fldi}{\mathbb{F}_{i}}              % its Gaussian extension F + iF
\newcommand{\lat}[1]{\big\lfloor #1\big\rfloor_{\kappa}} % round DOWN onto the lattice

% --- protocol-compiler ---------------------------------------------------------
\newcommand{\Spec}{\mathfrak{S}}             % protocol specification (compiler input)
\newcommand{\pcChain}{\mathfrak{R}}            % reduction chain (two typed segments)
\newcommand{\Chsrc}{\pcChain_{\mathrm{src}}}   % its Alice-side source-map segment
\newcommand{\Chrec}{\pcChain_{\mathrm{rec}}}   % its receiver-side squashing segment
\newcommand{\rvirt}{\rho^{\mathrm{v}}}       % virtual (PRE-source-map) signal state
\newcommand{\Fph}{F^{\mathrm{ph}}}           % physical (pre-reduction) POVM element
\newcommand{\Vph}{\mathcal{V}^{\mathrm{ph}}} % physical outcome alphabet
\newcommand{\Tpub}{\mathcal{T}_{\mathrm{pub}}}% public-choice alphabet
\newcommand{\Uset}{\mathcal{U}}              % residual outcome alphabet
\newcommand{\idxm}{\mathsf{idx}}             % label-to-index map
\newcommand{\Asig}{A'}                       % transmitted (signal) system -- the REAL one
\newcommand{\Avirt}{A''}                     % virtual transmitted system (source descriptor of Xi_1)
\newcommand{\Abar}{\bar{A}}                  % Alice's measured register
\newcommand{\Ahat}{\hat{A}}                  % Alice's shield register
\newcommand{\pcNmsm}{\mathcal{N}_{\mathsf A}}  % Alice's measurement channel (see the RULES block)
\newcommand{\Phip}{\Phi^{+}}                 % maximally entangled vector label
\newcommand{\Smap}{\mathfrak{s}}             % the emitted sector map (fraktur: a compile-time object)

% --- receiver ------------------------------------------------------------------
\newcommand{\Modes}{\mathbb{M}}              % declared ordered mode set
\newcommand{\Inp}{\mathbb{I}}                % declared input modes
\newcommand{\InpC}{\mathbb{I}_{\Chan}}       % channel input modes
\newcommand{\InpS}{\mathbb{I}_{\Prb}}        % probe (source) input modes
\newcommand{\Env}{\mathbb{E}}                % dilation (environment) modes
\newcommand{\Port}[1]{\mathbb{P}_{#1}}       % mode set of detector port #1
\newcommand{\Emd}[1]{E_{#1}}                 % coordinate projector on modes #1
\newcommand{\Gmat}{G}                        % instrument (transfer) matrix
\newcommand{\Fock}{\mathfrak{F}}             % Fock / second-quantization functor
\newcommand{\Symn}[1]{\mathrm{Sym}^{#1}}     % symmetric tensor power
\newcommand{\normord}[1]{{:}\mspace{2mu}#1\mspace{2mu}{:}}   % normal ordering
\newcommand{\rk}{\operatorname{rank}}        % rank (physics.sty owns \rank)
\newcommand{\SCK}{\mathsf{SC}}               % single-click pattern set
\newcommand{\Recv}{\mathfrak{R}}             % receiver instrument tuple
\newcommand{\Cat}{\mathfrak{C}}              % component catalogue
\newcommand{\Grph}{\mathfrak{E}}             % declared element graph
\newcommand{\defic}{\mathbb{L}}              % loss deficiency operator
\newcommand{\Pvac}[1]{\Pi_{\vac}(#1)}        % joint no-click element on ports #1
\newcommand{\pmc}{p_{\mc}}                   % uniform high-photon mc weight
\newcommand{\Bnd}{\beta}                     % the scalar bound function beta(n)
\newcommand{\Rot}{R}                         % planar rotation matrix
\newcommand{\Wplate}{W}                      % retarder (wave-plate) matrix
\newcommand{\Bsp}{B}                         % beamsplitter core matrix

% --- statistics ----------------------------------------------------------------
\newcommand{\honest}{p^{\mathrm{hon}}}          % honest per-round distribution
\newcommand{\ecomp}{\varepsilon_{\mathrm{comp}}}% completeness budget
\newcommand{\simplex}[1]{\Delta(#1)}            % probability simplex
\newcommand{\Chatcon}{\Chatset_{\mathrm{con}}}  % constrained outcomes
\newcommand{\Ireg}[3]{\mathcal{I}_{#1}(#2,#3)}  % regularized incomplete beta
\newcommand{\Ncount}[1]{\mathsf{N}_{#1}}        % per-outcome count
\newcommand{\qlo}{\mathbf{q}^{-}}               % box lower bounds
\newcommand{\qhi}{\mathbf{q}^{+}}               % box upper bounds
 for it, so the
%  file is read from spec/modules/<id>/.
%
%  RULES
%    * A macro that spec/notation.md catalogues is a verbatim copy of it
%      (part A). Changing one is a notation change: bump spec/notation.md.
%    * Part B holds the module-local macros. A macro defined by two modules
%      with the same expansion appears once, under the first module (in the
%      order below) that defines it. Two macros with the same name and
%      DIFFERENT expansions are both kept under a module prefix, and the
%      claims write the prefixed names:
%        \Chain  -> \redChain (reductions, \mathfrak{C}: the reduction chain)
%                   \pcChain  (protocol-compiler, \mathfrak{R}: the two-segment
%                             chain; \Chsrc and \Chrec are built on it)
%        \Nmsm   -> \redNmsm  (reductions, \mathcal{N}: a POVM measurement
%                             channel)
%                   \pcNmsm   (\mathcal{N}_{\mathsf A}: Alice's measurement
%                             channel, written by
%                             transmitter/thm:certificate-sufficiency,
%                             transmitter/rem:xreg and
%                             reductions/lem:source-transfer)
%    * Deliberately omitted: a \provenance{tag}{details} macro. Provenance
%      is a header field of every claim file (spec/claim-format.md section
%      3); a claim body that calls such a macro fails to compile, which is
%      the intended signal.
%
%  REQUIRED PACKAGES, loaded by the assembler's preamble before this file:
%    amsmath, amssymb, amsfonts, mathtools, dsfont, braket, physics, xcolor
%    (braket before physics, so that physics wins for \ket, \bra, \tr, ...).
%
%  This file contains only definitions: no \usepackage, no document body.
% =============================================================================


% =============================================================================
%  PART A -- notation macros, the ones spec/notation.md catalogues, each a
%  verbatim copy of its catalogue row; grouped under numbered headings.
% =============================================================================

% -----------------------------------------------------------------------------
% 1. State spaces and operator sets (a "system label" -> spaces of operators)
% -----------------------------------------------------------------------------
\newcommand{\Hs}[1]{\mathcal{H}(#1)}        % Hilbert space of a system label
\newcommand{\Lin}[1]{\mathcal{L}(#1)}       % bounded linear operators on H(#1)
\newcommand{\Linn}[2]{\mathcal{L}(#1,#2)}   % linear maps H(#1) -> H(#2)
\newcommand{\Pos}[1]{\mathrm{Pos}(#1)}      % positive semidefinite operators on H(#1)
\newcommand{\Density}{\mathrm{D}}           % density operators: \Density(A) = D(A)
\newcommand{\Seq}[1]{S_{=}(#1)}             % normalized states (trace = 1)
\newcommand{\Sleq}[1]{S_{\leq}(#1)}         % subnormalized states (0 <= trace <= 1)
\newcommand{\id}{\mathds{1}}                % identity operator
\newcommand{\idon}[1]{\mathds{1}_{#1}}      % identity on a named system, e.g. \idon{C}

% DEPRECATED: applies \mathcal to both arguments itself, so it takes bare
% letters, not system labels. Prefer \Linn; do not use in new contracts.
\providecommand{\Lop}[2]{\mathcal{L}(\mathcal{#1}, \mathcal{#2})}


% -----------------------------------------------------------------------------
% 2. Composite systems
% -----------------------------------------------------------------------------
\newcommand{\dsum}{\oplus}                  % direct sum of system labels
% Tensor product: use \otimes directly.
% n-fold direct sum of an operator space: \Lin{A}^{\oplus n}


% -----------------------------------------------------------------------------
% 3. Events, conditioning and classical registers (for STATES)
%    Convention:  \wedge  -> subnormalized "and-event" component, e.g. \rho_{A\wedge x}
%                 |        -> normalized conditional state,        e.g. \rho_{A|x}
% (Write these directly in subscripts; no extra macros needed.)
% -----------------------------------------------------------------------------


% -----------------------------------------------------------------------------
% 4. Photon-number subspaces / direct-sum blocks
%    Convention:  operators carry the restriction "on top", in brackets;
%                 states carry it as a \wedge subscript (see Sec. 3).
% -----------------------------------------------------------------------------
\newcommand{\Pnum}[1]{P_{#1}}               % photon-number projector (system from context)
\newcommand{\Pnumon}[2]{P_{#1}^{#2}}        % with explicit system superscript, e.g. P_n^C
\newcommand{\blk}[2]{#1^{(#2)}}             % operator block, e.g. \blk{M}{\leqN} = M^{(<=N)}
\newcommand{\leqN}{\leq N}                  % use as \blk{M}{\leqN}
\newcommand{\gtN}{> N}                      % use as \blk{M}{\gtN}
\newcommand{\eqn}{= n}                      % use as \blk{M}{\eqn}


% -----------------------------------------------------------------------------
% 5. System labels used in the protocol
% -----------------------------------------------------------------------------
\newcommand{\Chan}{C}                       % quantum-channel (Eve) input to Bob
\newcommand{\Prb}{P}                        % probe input
\newcommand{\Flag}{{\mathrm{Flg}}}          % flag system (EXISTS ONLY AFTER squashing)
\newcommand{\Loss}{{\mathrm{Lss}}}          % loss system
\newcommand{\Vc}{{\mathrm{Vcm}}}            % vacuum system
\newcommand{\Nout}{N_{\mathrm{out}}}
\newcommand{\Nin}{N_{\mathrm{in}}}


% -----------------------------------------------------------------------------
% 6. Alphabets
% -----------------------------------------------------------------------------
\newcommand{\Yset}{\mathcal{Y}}             % Bob's (coarse-grained) detection alphabet
\newcommand{\Yraw}{\mathcal{Y}_{\mathrm{raw}}}% raw click patterns {0,1}^4
\newcommand{\Yscset}{\mathcal{Y}_{\singlec}}% single-click outcomes
\newcommand{\Ymcset}{\mathcal{Y}_{\mc}}     % multi-click outcomes
\newcommand{\Tset}{\mathcal{T}_{\mathrm{set}}}% private setting set (the probe settings in the optical instance)
\newcommand{\Xset}{\mathcal{X}}             % Alice's alphabet
\newcommand{\Dset}{\{H,V,+,-\}}             % polarization / detector labels
\newcommand{\pass}{\top}                    % pass sifting


% -----------------------------------------------------------------------------
% 7. Mode / outcome abbreviations
% -----------------------------------------------------------------------------
\newcommand{\on}{{\mathrm{on}}}
\newcommand{\off}{{\mathrm{off}}}
\newcommand{\dc}{{\mathrm{dc}}}             % dark count
\newcommand{\sq}{\mathrm{sq}}               % squashed
\newcommand{\mc}{\mathrm{mc}}               % multi-click
\newcommand{\nc}{\mathrm{nc}}               % no-click
\newcommand{\singlec}{\mathrm{sc}}          % single-click
\newcommand{\vac}{\mathrm{vac}}             % vacuum
\newcommand{\kout}{{k_\mathrm{out}}}
\newcommand{\kin}{{k_\mathrm{in}}}


% -----------------------------------------------------------------------------
% 8. Maps, channels and miscellaneous operators
% -----------------------------------------------------------------------------
\newcommand{\Prob}{\mathrm{Prob}}           % probability (use \Prob[\cdot] or \Prob(\cdot))
\newcommand{\lmin}{\lambda_{\min}}          % minimum eigenvalue
\newcommand{\Wflag}{W_{\Flag}}              % flag-space weight
\newcommand{\transpose}{\mathsf{T}}
\newcommand{\innerp}[2]{\left\langle #1 \vert #2 \right\rangle}
\renewcommand{\vec}[1]{\boldsymbol{#1}}     % NB: overrides the LaTeX arrow accent


% -----------------------------------------------------------------------------
% 9. Key-rate numerics
% -----------------------------------------------------------------------------
\newcommand{\cG}{\mathcal{G}}               % protocol map
\newcommand{\cZ}{\mathcal{Z}}               % pinching map
\newcommand{\Amap}{\mathcal{A}_{\mathrm{ann}}} % announcement + measurement map
\newcommand{\Amapnb}{\mathcal{A}_{\mathrm{ann}|\mathrm{nb}}}
\newcommand{\simpl}{}                       % empty suffix for simplified block maps
\newcommand{\acSet}{\mathcal{S}_{\mathrm{acc}}} % acceptance set
\newcommand{\testE}{\mathcal{E}_{\mathrm{test}}} % test-statistics map
\newcommand{\ecLeak}{\mathrm{leak}^{\mathrm{EC}}_{\mathrm{obs}}}
\newcommand{\pgen}{p_{\mathrm{gen}}}
\newcommand{\askeyr}{R_{\infty}}             % asymptotic key rate
\newcommand{\ppass}{p_{\mathrm{pass}}}       % sifting pass probability


% -----------------------------------------------------------------------------
% 10. Finite-size analysis / Renyi entropies
% -----------------------------------------------------------------------------

% --- Round types: test vs. key generation (register T, alphabet \Fset) --------
\newcommand{\tst}{\mathrm{test}}            % test round
\newcommand{\gen}{\mathrm{gen}}             % key-generation round
\newcommand{\Fset}{\mathcal{T}}             % round-type alphabet {test, gen} (not \Tset)
\newcommand{\Chat}{\widehat{C}}             % acceptance-test register
\newcommand{\Chatset}{\widehat{\mathcal{C}}}% its alphabet

% --- Basis / key extraction functions ----------------------------------------
\newcommand{\bA}{b_{A}}                     % Alice's basis-extraction function
\newcommand{\bB}{b_{B}}                     % Bob's basis-extraction function
\newcommand{\kB}{k_{B}}                     % Bob's key-extraction function

% --- Maps --------------------------------------------------------------------
\newcommand{\Mprot}{\mathcal{M}}            % single-round protocol map
\newcommand{\Phimap}{\boldsymbol{\Phi}}     % test-statistics map
\newcommand{\Pdeph}{\mathcal{P}}            % dephasing (pinching) map on \Blreg

% --- Isometries --------------------------------------------------------------
\newcommand{\Vmsm}{V_{\mathrm{msm}}}        % measurement isometry
\newcommand{\Vann}{V_{\mathrm{ann}}}        % announcement isometry
\newcommand{\VI}{V_{I}}                     % Eve's copy of the announcements
\newcommand{\VChat}{V_{\Chat}}              % acceptance-data isometry
\newcommand{\VR}{V_{R}}                     % sifting + key map -> raw key R
\newcommand{\VZ}{V_{Z}}                     % raw key -> key register Z
\newcommand{\Wprot}{W}                      % full protocol isometry

% --- Renyi entropies ---------------------------------------------------------
% SECURITY-CRITICAL: see spec/notation.md §4. Never change an arrow to
% make an inequality work out; never convert silently between the sandwiched
% (tilde) and Petz (overline) families.
\newcommand{\Hsandup}{\widetilde{H}^{\uparrow}_{\alpha}}   % sandwiched, up
\newcommand{\Hsanddown}{\widetilde{H}^{\downarrow}_{\alpha}}% sandwiched, down
\newcommand{\Hpetzup}{\overline{H}^{\uparrow}_{\alpha}}    % Petz, up
\newcommand{\Hpetzdown}{\overline{H}^{\downarrow}_{\alpha}}% Petz, down
\newcommand{\hup}{h^{\uparrow}_{\alpha}}                   % single-round rate
\newcommand{\hqkd}{h^{\uparrow,\mathrm{QKD}}_{\alpha}}     % single-round QKD rate
\newcommand{\hlow}{\underline{h}_{\alpha}}                 % lower bound on \hqkd
% --- spec/notation.md 3.5 and 4.5 -------------------------------------------
% \hlb  is the first-order linearization lower bound on \hlow: the headline
%       key rate (the `first-order` tier of foundations/def:tiers). NOT \hlow
%       itself.
% \tnum is the instance's one declared numerical tolerance. Its
%       near-collision with the declared marginal \tau_A is disclosed in
%       spec/notation.md 3.5; bare \tau is reserved and denotes neither.
\newcommand{\hlb}{\underline{h}^{\mathrm{FO}}_{\alpha}}    % first-order lower bound on \hlow
\newcommand{\tnum}{\tau_{\mathrm{num}}}                    % declared numerical tolerance


% =============================================================================
%  PART B -- module-local macros, grouped by module. Deprecated aliases,
%  defined so that a text written with them compiles, used in no new contract: \Lop,
%  \testE, \Pdeph, \pass, \ecLeak (spec/notation.md section 3.9).
% =============================================================================

% --- meat-core -----------------------------------------------------------------
\newcommand{\SecSet}{\mathbb{S}}            % sector index set of Q
\newcommand{\Qsys}{Q}                       % declared sector-carrying measured system
\newcommand{\Vset}{\mathcal{V}}             % declared joint outcome alphabet of the Q-POVM
\newcommand{\CAset}{\mathcal{C}_{A}}        % Alice announcement alphabet
\newcommand{\CBset}{\mathcal{C}_{B}}        % Bob announcement alphabet
\newcommand{\Party}{\mathfrak{P}}           % declared party partition
\newcommand{\Oset}[1]{\mathcal{O}_{#1}}     % readout alphabet of party #1
\newcommand{\rd}[1]{\pi_{#1}}               % readout map of party #1
\newcommand{\Sadm}{\mathcal{S}_{A}}         % admissible Alice subspace
\newcommand{\nomarg}{\mathsf{none}}         % "no trusted marginal declared"
\newcommand{\Zset}{\mathcal{Z}_{\perp}}     % raw-key alphabet incl. \perp
\newcommand{\Zval}{\mathcal{Z}}             % key alphabet without \perp
\newcommand{\Etil}{\widetilde{E}}           % single-round purifying register
\newcommand{\Ehat}{\widehat{E}}             % Eve's n-round register
\newcommand{\OmAT}{\Omega_{\mathrm{AT}}}    % acceptance event
\newcommand{\epa}{\varepsilon_{\mathrm{PA}}}
\newcommand{\eev}{\varepsilon_{\mathrm{EV}}}
\newcommand{\lEC}{\lambda_{\mathrm{EC}}}
\newcommand{\Decl}{\mathfrak{D}}            % protocol declaration
\newcommand{\CList}{\mathfrak{L}}           % declared constraint list
\newcommand{\PinS}{\mathcal{P}_{\SecSet}}   % sector pinching on Q
\newcommand{\Feas}{\mathcal{W}}             % feasible set of block states
\newcommand{\Chkey}{\Chatset_{\mathrm{key}}}% key-live acceptance values
\newcommand{\Kann}{\mathcal{K}}             % announcement-class alphabet
\newcommand{\hold}{\mathsf{h}}              % declared key holder (A or B)
\newcommand{\sfD}{D}                        % classical KL divergence (base 2)
\newcommand{\Dtwo}{D}                       % quantum relative entropy (base 2)
\newcommand{\Dsand}{\widetilde{D}_{\alpha}} % sandwiched Renyi divergence
\newcommand{\idmap}{\mathrm{id}}            % identity channel
\newcommand{\Aev}{\mathcal{A}}              % Eve's per-round attack channel
\newcommand{\AttSet}{\mathfrak{A}}          % licensed attack-channel set
\newcommand{\Chacc}{\Chatset_{\mathrm{acc}}}% acceptance support of \acSet
\newcommand{\epsstar}{\varepsilon^{\star}}      % distance to an exact POVM
\newcommand{\dfH}{\Delta_{\mathrm{herm}}}       % Hermiticity defect
\newcommand{\dfP}{\Delta_{\mathrm{pos}}}        % negative part
\newcommand{\dfC}{\Delta_{\mathrm{cpl}}}        % completeness defect
\newcommand{\lmoq}{g}                       % acceptance-side gradient block
\newcommand{\lmoJ}{G}                       % feasible-side gradient block
\newcommand{\bacc}{\beta_{\mathrm{acc}}}    % certified acceptance-side minimum
\newcommand{\Itset}{\mathfrak{I}}           % finite index set of iterates (C22)

% --- reductions ----------------------------------------------------------------
\newcommand{\Bsys}{B}                        % physical (pre-squashing) system
\newcommand{\Tcl}[1]{\mathcal{T}(#1)}        % trace-class operators on H(#1)
\newcommand{\Ins}{\mathcal{I}}               % generalized-instrument branch map
\newcommand{\redNmsm}{\mathcal{N}}              % POVM measurement channel
\newcommand{\Eprot}{\mathcal{E}}             % n-round protocol channel
\newcommand{\Grnd}{\mathsf{G}}               % one round map of a protocol
\newcommand{\Kfin}{\mathsf{K}}               % end-of-protocol processing
\newcommand{\Crit}{\mathsf{sec}}             % value of a security criterion
\newcommand{\Pflag}{\Pi_{\Flag}}             % flag-sector projector on Q
\newcommand{\Pleq}{\Pi_{\leqN}}              % preserved-block projector
\newcommand{\Pgt}{\Pi_{\gtN}}                % above-cutoff projector
\newcommand{\Sph}{S^{\mathrm{ph}}}           % family of physical inputs
\newcommand{\Srch}{S^{\mathrm{rch}}}         % reachable squashed states
\newcommand{\Ssq}{S^{\sq}}                   % squashed admissible set
\newcommand{\Ffss}{F^{\mathrm{FSS}}}         % flag-state-squasher target
\newcommand{\Fwp}{F^{\mathrm{WP}}}           % multi-bound WPFSS target
\newcommand{\Mwp}{M^{\mathrm{WP}}}           % renormalized above-cutoff POVM
\newcommand{\Mdec}{{M'}^{\mathrm{ph}}}       % declared physical family on B
\newcommand{\Ppvm}{P^{\Flag}}                % flag PVM
\newcommand{\Lmp}{{\Lambda_{\mathrm{mp}}}}   % measure-and-prepare channel
\newcommand{\Lwp}{{\Lambda^{\mathrm{WP}}}}   % multi-bound WPFSS channel
\newcommand{\Lfss}{{\Lambda^{\mathrm{FSS}}}} % flag-state-squasher channel
\newcommand{\lmax}{\lambda_{\max}}           % maximum eigenvalue
\newcommand{\Lnoise}{\widetilde{\Lambda}}    % flag-noise channel
\newcommand{\cg}{\varpi}                     % declared coarse-graining map
\newcommand{\MCK}{\mathsf{MC}}               % multiclick raw-pattern set
\newcommand{\redChain}{\mathfrak{C}}            % reduction chain
\newcommand{\Lnk}{\ell}                      % a single typed link
\newcommand{\dsq}{\preceq}                      % the squashing relation
\newcommand{\opn}[1]{\left\lVert #1\right\rVert}           % operator (spectral) norm
\newcommand{\mxn}[1]{\left\lVert #1\right\rVert_{\max}}    % entrywise max modulus
\newcommand{\Fld}{\mathbb{F}}                   % exact ordered substrate field
\newcommand{\Fldi}{\mathbb{F}_{i}}              % its Gaussian extension F + iF
\newcommand{\lat}[1]{\big\lfloor #1\big\rfloor_{\kappa}} % round DOWN onto the lattice

% --- protocol-compiler ---------------------------------------------------------
\newcommand{\Spec}{\mathfrak{S}}             % protocol specification (compiler input)
\newcommand{\pcChain}{\mathfrak{R}}            % reduction chain (two typed segments)
\newcommand{\Chsrc}{\pcChain_{\mathrm{src}}}   % its Alice-side source-map segment
\newcommand{\Chrec}{\pcChain_{\mathrm{rec}}}   % its receiver-side squashing segment
\newcommand{\rvirt}{\rho^{\mathrm{v}}}       % virtual (PRE-source-map) signal state
\newcommand{\Fph}{F^{\mathrm{ph}}}           % physical (pre-reduction) POVM element
\newcommand{\Vph}{\mathcal{V}^{\mathrm{ph}}} % physical outcome alphabet
\newcommand{\Tpub}{\mathcal{T}_{\mathrm{pub}}}% public-choice alphabet
\newcommand{\Uset}{\mathcal{U}}              % residual outcome alphabet
\newcommand{\idxm}{\mathsf{idx}}             % label-to-index map
\newcommand{\Asig}{A'}                       % transmitted (signal) system -- the REAL one
\newcommand{\Avirt}{A''}                     % virtual transmitted system (source descriptor of Xi_1)
\newcommand{\Abar}{\bar{A}}                  % Alice's measured register
\newcommand{\Ahat}{\hat{A}}                  % Alice's shield register
\newcommand{\pcNmsm}{\mathcal{N}_{\mathsf A}}  % Alice's measurement channel (see the RULES block)
\newcommand{\Phip}{\Phi^{+}}                 % maximally entangled vector label
\newcommand{\Smap}{\mathfrak{s}}             % the emitted sector map (fraktur: a compile-time object)

% --- receiver ------------------------------------------------------------------
\newcommand{\Modes}{\mathbb{M}}              % declared ordered mode set
\newcommand{\Inp}{\mathbb{I}}                % declared input modes
\newcommand{\InpC}{\mathbb{I}_{\Chan}}       % channel input modes
\newcommand{\InpS}{\mathbb{I}_{\Prb}}        % probe (source) input modes
\newcommand{\Env}{\mathbb{E}}                % dilation (environment) modes
\newcommand{\Port}[1]{\mathbb{P}_{#1}}       % mode set of detector port #1
\newcommand{\Emd}[1]{E_{#1}}                 % coordinate projector on modes #1
\newcommand{\Gmat}{G}                        % instrument (transfer) matrix
\newcommand{\Fock}{\mathfrak{F}}             % Fock / second-quantization functor
\newcommand{\Symn}[1]{\mathrm{Sym}^{#1}}     % symmetric tensor power
\newcommand{\normord}[1]{{:}\mspace{2mu}#1\mspace{2mu}{:}}   % normal ordering
\newcommand{\rk}{\operatorname{rank}}        % rank (physics.sty owns \rank)
\newcommand{\SCK}{\mathsf{SC}}               % single-click pattern set
\newcommand{\Recv}{\mathfrak{R}}             % receiver instrument tuple
\newcommand{\Cat}{\mathfrak{C}}              % component catalogue
\newcommand{\Grph}{\mathfrak{E}}             % declared element graph
\newcommand{\defic}{\mathbb{L}}              % loss deficiency operator
\newcommand{\Pvac}[1]{\Pi_{\vac}(#1)}        % joint no-click element on ports #1
\newcommand{\pmc}{p_{\mc}}                   % uniform high-photon mc weight
\newcommand{\Bnd}{\beta}                     % the scalar bound function beta(n)
\newcommand{\Rot}{R}                         % planar rotation matrix
\newcommand{\Wplate}{W}                      % retarder (wave-plate) matrix
\newcommand{\Bsp}{B}                         % beamsplitter core matrix

% --- statistics ----------------------------------------------------------------
\newcommand{\honest}{p^{\mathrm{hon}}}          % honest per-round distribution
\newcommand{\ecomp}{\varepsilon_{\mathrm{comp}}}% completeness budget
\newcommand{\simplex}[1]{\Delta(#1)}            % probability simplex
\newcommand{\Chatcon}{\Chatset_{\mathrm{con}}}  % constrained outcomes
\newcommand{\Ireg}[3]{\mathcal{I}_{#1}(#2,#3)}  % regularized incomplete beta
\newcommand{\Ncount}[1]{\mathsf{N}_{#1}}        % per-outcome count
\newcommand{\qlo}{\mathbf{q}^{-}}               % box lower bounds
\newcommand{\qhi}{\mathbf{q}^{+}}               % box upper bounds
 for it, so the
%  file is read from spec/modules/<id>/.
%
%  RULES
%    * A macro that spec/notation.md catalogues is a verbatim copy of it
%      (part A). Changing one is a notation change: bump spec/notation.md.
%    * Part B holds the module-local macros. A macro defined by two modules
%      with the same expansion appears once, under the first module (in the
%      order below) that defines it. Two macros with the same name and
%      DIFFERENT expansions are both kept under a module prefix, and the
%      claims write the prefixed names:
%        \Chain  -> \redChain (reductions, \mathfrak{C}: the reduction chain)
%                   \pcChain  (protocol-compiler, \mathfrak{R}: the two-segment
%                             chain; \Chsrc and \Chrec are built on it)
%        \Nmsm   -> \redNmsm  (reductions, \mathcal{N}: a POVM measurement
%                             channel)
%                   \pcNmsm   (\mathcal{N}_{\mathsf A}: Alice's measurement
%                             channel, written by
%                             transmitter/thm:certificate-sufficiency,
%                             transmitter/rem:xreg and
%                             reductions/lem:source-transfer)
%    * Deliberately omitted: a \provenance{tag}{details} macro. Provenance
%      is a header field of every claim file (spec/claim-format.md section
%      3); a claim body that calls such a macro fails to compile, which is
%      the intended signal.
%
%  REQUIRED PACKAGES, loaded by the assembler's preamble before this file:
%    amsmath, amssymb, amsfonts, mathtools, dsfont, braket, physics, xcolor
%    (braket before physics, so that physics wins for \ket, \bra, \tr, ...).
%
%  This file contains only definitions: no \usepackage, no document body.
% =============================================================================


% =============================================================================
%  PART A -- notation macros, the ones spec/notation.md catalogues, each a
%  verbatim copy of its catalogue row; grouped under numbered headings.
% =============================================================================

% -----------------------------------------------------------------------------
% 1. State spaces and operator sets (a "system label" -> spaces of operators)
% -----------------------------------------------------------------------------
\newcommand{\Hs}[1]{\mathcal{H}(#1)}        % Hilbert space of a system label
\newcommand{\Lin}[1]{\mathcal{L}(#1)}       % bounded linear operators on H(#1)
\newcommand{\Linn}[2]{\mathcal{L}(#1,#2)}   % linear maps H(#1) -> H(#2)
\newcommand{\Pos}[1]{\mathrm{Pos}(#1)}      % positive semidefinite operators on H(#1)
\newcommand{\Density}{\mathrm{D}}           % density operators: \Density(A) = D(A)
\newcommand{\Seq}[1]{S_{=}(#1)}             % normalized states (trace = 1)
\newcommand{\Sleq}[1]{S_{\leq}(#1)}         % subnormalized states (0 <= trace <= 1)
\newcommand{\id}{\mathds{1}}                % identity operator
\newcommand{\idon}[1]{\mathds{1}_{#1}}      % identity on a named system, e.g. \idon{C}

% DEPRECATED: applies \mathcal to both arguments itself, so it takes bare
% letters, not system labels. Prefer \Linn; do not use in new contracts.
\providecommand{\Lop}[2]{\mathcal{L}(\mathcal{#1}, \mathcal{#2})}


% -----------------------------------------------------------------------------
% 2. Composite systems
% -----------------------------------------------------------------------------
\newcommand{\dsum}{\oplus}                  % direct sum of system labels
% Tensor product: use \otimes directly.
% n-fold direct sum of an operator space: \Lin{A}^{\oplus n}


% -----------------------------------------------------------------------------
% 3. Events, conditioning and classical registers (for STATES)
%    Convention:  \wedge  -> subnormalized "and-event" component, e.g. \rho_{A\wedge x}
%                 |        -> normalized conditional state,        e.g. \rho_{A|x}
% (Write these directly in subscripts; no extra macros needed.)
% -----------------------------------------------------------------------------


% -----------------------------------------------------------------------------
% 4. Photon-number subspaces / direct-sum blocks
%    Convention:  operators carry the restriction "on top", in brackets;
%                 states carry it as a \wedge subscript (see Sec. 3).
% -----------------------------------------------------------------------------
\newcommand{\Pnum}[1]{P_{#1}}               % photon-number projector (system from context)
\newcommand{\Pnumon}[2]{P_{#1}^{#2}}        % with explicit system superscript, e.g. P_n^C
\newcommand{\blk}[2]{#1^{(#2)}}             % operator block, e.g. \blk{M}{\leqN} = M^{(<=N)}
\newcommand{\leqN}{\leq N}                  % use as \blk{M}{\leqN}
\newcommand{\gtN}{> N}                      % use as \blk{M}{\gtN}
\newcommand{\eqn}{= n}                      % use as \blk{M}{\eqn}


% -----------------------------------------------------------------------------
% 5. System labels used in the protocol
% -----------------------------------------------------------------------------
\newcommand{\Chan}{C}                       % quantum-channel (Eve) input to Bob
\newcommand{\Prb}{P}                        % probe input
\newcommand{\Flag}{{\mathrm{Flg}}}          % flag system (EXISTS ONLY AFTER squashing)
\newcommand{\Loss}{{\mathrm{Lss}}}          % loss system
\newcommand{\Vc}{{\mathrm{Vcm}}}            % vacuum system
\newcommand{\Nout}{N_{\mathrm{out}}}
\newcommand{\Nin}{N_{\mathrm{in}}}


% -----------------------------------------------------------------------------
% 6. Alphabets
% -----------------------------------------------------------------------------
\newcommand{\Yset}{\mathcal{Y}}             % Bob's (coarse-grained) detection alphabet
\newcommand{\Yraw}{\mathcal{Y}_{\mathrm{raw}}}% raw click patterns {0,1}^4
\newcommand{\Yscset}{\mathcal{Y}_{\singlec}}% single-click outcomes
\newcommand{\Ymcset}{\mathcal{Y}_{\mc}}     % multi-click outcomes
\newcommand{\Tset}{\mathcal{T}_{\mathrm{set}}}% private setting set (the probe settings in the optical instance)
\newcommand{\Xset}{\mathcal{X}}             % Alice's alphabet
\newcommand{\Dset}{\{H,V,+,-\}}             % polarization / detector labels
\newcommand{\pass}{\top}                    % pass sifting


% -----------------------------------------------------------------------------
% 7. Mode / outcome abbreviations
% -----------------------------------------------------------------------------
\newcommand{\on}{{\mathrm{on}}}
\newcommand{\off}{{\mathrm{off}}}
\newcommand{\dc}{{\mathrm{dc}}}             % dark count
\newcommand{\sq}{\mathrm{sq}}               % squashed
\newcommand{\mc}{\mathrm{mc}}               % multi-click
\newcommand{\nc}{\mathrm{nc}}               % no-click
\newcommand{\singlec}{\mathrm{sc}}          % single-click
\newcommand{\vac}{\mathrm{vac}}             % vacuum
\newcommand{\kout}{{k_\mathrm{out}}}
\newcommand{\kin}{{k_\mathrm{in}}}


% -----------------------------------------------------------------------------
% 8. Maps, channels and miscellaneous operators
% -----------------------------------------------------------------------------
\newcommand{\Prob}{\mathrm{Prob}}           % probability (use \Prob[\cdot] or \Prob(\cdot))
\newcommand{\lmin}{\lambda_{\min}}          % minimum eigenvalue
\newcommand{\Wflag}{W_{\Flag}}              % flag-space weight
\newcommand{\transpose}{\mathsf{T}}
\newcommand{\innerp}[2]{\left\langle #1 \vert #2 \right\rangle}
\renewcommand{\vec}[1]{\boldsymbol{#1}}     % NB: overrides the LaTeX arrow accent


% -----------------------------------------------------------------------------
% 9. Key-rate numerics
% -----------------------------------------------------------------------------
\newcommand{\cG}{\mathcal{G}}               % protocol map
\newcommand{\cZ}{\mathcal{Z}}               % pinching map
\newcommand{\Amap}{\mathcal{A}_{\mathrm{ann}}} % announcement + measurement map
\newcommand{\Amapnb}{\mathcal{A}_{\mathrm{ann}|\mathrm{nb}}}
\newcommand{\simpl}{}                       % empty suffix for simplified block maps
\newcommand{\acSet}{\mathcal{S}_{\mathrm{acc}}} % acceptance set
\newcommand{\testE}{\mathcal{E}_{\mathrm{test}}} % test-statistics map
\newcommand{\ecLeak}{\mathrm{leak}^{\mathrm{EC}}_{\mathrm{obs}}}
\newcommand{\pgen}{p_{\mathrm{gen}}}
\newcommand{\askeyr}{R_{\infty}}             % asymptotic key rate
\newcommand{\ppass}{p_{\mathrm{pass}}}       % sifting pass probability


% -----------------------------------------------------------------------------
% 10. Finite-size analysis / Renyi entropies
% -----------------------------------------------------------------------------

% --- Round types: test vs. key generation (register T, alphabet \Fset) --------
\newcommand{\tst}{\mathrm{test}}            % test round
\newcommand{\gen}{\mathrm{gen}}             % key-generation round
\newcommand{\Fset}{\mathcal{T}}             % round-type alphabet {test, gen} (not \Tset)
\newcommand{\Chat}{\widehat{C}}             % acceptance-test register
\newcommand{\Chatset}{\widehat{\mathcal{C}}}% its alphabet

% --- Basis / key extraction functions ----------------------------------------
\newcommand{\bA}{b_{A}}                     % Alice's basis-extraction function
\newcommand{\bB}{b_{B}}                     % Bob's basis-extraction function
\newcommand{\kB}{k_{B}}                     % Bob's key-extraction function

% --- Maps --------------------------------------------------------------------
\newcommand{\Mprot}{\mathcal{M}}            % single-round protocol map
\newcommand{\Phimap}{\boldsymbol{\Phi}}     % test-statistics map
\newcommand{\Pdeph}{\mathcal{P}}            % dephasing (pinching) map on \Blreg

% --- Isometries --------------------------------------------------------------
\newcommand{\Vmsm}{V_{\mathrm{msm}}}        % measurement isometry
\newcommand{\Vann}{V_{\mathrm{ann}}}        % announcement isometry
\newcommand{\VI}{V_{I}}                     % Eve's copy of the announcements
\newcommand{\VChat}{V_{\Chat}}              % acceptance-data isometry
\newcommand{\VR}{V_{R}}                     % sifting + key map -> raw key R
\newcommand{\VZ}{V_{Z}}                     % raw key -> key register Z
\newcommand{\Wprot}{W}                      % full protocol isometry

% --- Renyi entropies ---------------------------------------------------------
% SECURITY-CRITICAL: see spec/notation.md §4. Never change an arrow to
% make an inequality work out; never convert silently between the sandwiched
% (tilde) and Petz (overline) families.
\newcommand{\Hsandup}{\widetilde{H}^{\uparrow}_{\alpha}}   % sandwiched, up
\newcommand{\Hsanddown}{\widetilde{H}^{\downarrow}_{\alpha}}% sandwiched, down
\newcommand{\Hpetzup}{\overline{H}^{\uparrow}_{\alpha}}    % Petz, up
\newcommand{\Hpetzdown}{\overline{H}^{\downarrow}_{\alpha}}% Petz, down
\newcommand{\hup}{h^{\uparrow}_{\alpha}}                   % single-round rate
\newcommand{\hqkd}{h^{\uparrow,\mathrm{QKD}}_{\alpha}}     % single-round QKD rate
\newcommand{\hlow}{\underline{h}_{\alpha}}                 % lower bound on \hqkd
% --- spec/notation.md 3.5 and 4.5 -------------------------------------------
% \hlb  is the first-order linearization lower bound on \hlow: the headline
%       key rate (the `first-order` tier of foundations/def:tiers). NOT \hlow
%       itself.
% \tnum is the instance's one declared numerical tolerance. Its
%       near-collision with the declared marginal \tau_A is disclosed in
%       spec/notation.md 3.5; bare \tau is reserved and denotes neither.
\newcommand{\hlb}{\underline{h}^{\mathrm{FO}}_{\alpha}}    % first-order lower bound on \hlow
\newcommand{\tnum}{\tau_{\mathrm{num}}}                    % declared numerical tolerance


% =============================================================================
%  PART B -- module-local macros, grouped by module. Deprecated aliases,
%  defined so that a text written with them compiles, used in no new contract: \Lop,
%  \testE, \Pdeph, \pass, \ecLeak (spec/notation.md section 3.9).
% =============================================================================

% --- meat-core -----------------------------------------------------------------
\newcommand{\SecSet}{\mathbb{S}}            % sector index set of Q
\newcommand{\Qsys}{Q}                       % declared sector-carrying measured system
\newcommand{\Vset}{\mathcal{V}}             % declared joint outcome alphabet of the Q-POVM
\newcommand{\CAset}{\mathcal{C}_{A}}        % Alice announcement alphabet
\newcommand{\CBset}{\mathcal{C}_{B}}        % Bob announcement alphabet
\newcommand{\Party}{\mathfrak{P}}           % declared party partition
\newcommand{\Oset}[1]{\mathcal{O}_{#1}}     % readout alphabet of party #1
\newcommand{\rd}[1]{\pi_{#1}}               % readout map of party #1
\newcommand{\Sadm}{\mathcal{S}_{A}}         % admissible Alice subspace
\newcommand{\nomarg}{\mathsf{none}}         % "no trusted marginal declared"
\newcommand{\Zset}{\mathcal{Z}_{\perp}}     % raw-key alphabet incl. \perp
\newcommand{\Zval}{\mathcal{Z}}             % key alphabet without \perp
\newcommand{\Etil}{\widetilde{E}}           % single-round purifying register
\newcommand{\Ehat}{\widehat{E}}             % Eve's n-round register
\newcommand{\OmAT}{\Omega_{\mathrm{AT}}}    % acceptance event
\newcommand{\epa}{\varepsilon_{\mathrm{PA}}}
\newcommand{\eev}{\varepsilon_{\mathrm{EV}}}
\newcommand{\lEC}{\lambda_{\mathrm{EC}}}
\newcommand{\Decl}{\mathfrak{D}}            % protocol declaration
\newcommand{\CList}{\mathfrak{L}}           % declared constraint list
\newcommand{\PinS}{\mathcal{P}_{\SecSet}}   % sector pinching on Q
\newcommand{\Feas}{\mathcal{W}}             % feasible set of block states
\newcommand{\Chkey}{\Chatset_{\mathrm{key}}}% key-live acceptance values
\newcommand{\Kann}{\mathcal{K}}             % announcement-class alphabet
\newcommand{\hold}{\mathsf{h}}              % declared key holder (A or B)
\newcommand{\sfD}{D}                        % classical KL divergence (base 2)
\newcommand{\Dtwo}{D}                       % quantum relative entropy (base 2)
\newcommand{\Dsand}{\widetilde{D}_{\alpha}} % sandwiched Renyi divergence
\newcommand{\idmap}{\mathrm{id}}            % identity channel
\newcommand{\Aev}{\mathcal{A}}              % Eve's per-round attack channel
\newcommand{\AttSet}{\mathfrak{A}}          % licensed attack-channel set
\newcommand{\Chacc}{\Chatset_{\mathrm{acc}}}% acceptance support of \acSet
\newcommand{\epsstar}{\varepsilon^{\star}}      % distance to an exact POVM
\newcommand{\dfH}{\Delta_{\mathrm{herm}}}       % Hermiticity defect
\newcommand{\dfP}{\Delta_{\mathrm{pos}}}        % negative part
\newcommand{\dfC}{\Delta_{\mathrm{cpl}}}        % completeness defect
\newcommand{\lmoq}{g}                       % acceptance-side gradient block
\newcommand{\lmoJ}{G}                       % feasible-side gradient block
\newcommand{\bacc}{\beta_{\mathrm{acc}}}    % certified acceptance-side minimum
\newcommand{\Itset}{\mathfrak{I}}           % finite index set of iterates (C22)

% --- reductions ----------------------------------------------------------------
\newcommand{\Bsys}{B}                        % physical (pre-squashing) system
\newcommand{\Tcl}[1]{\mathcal{T}(#1)}        % trace-class operators on H(#1)
\newcommand{\Ins}{\mathcal{I}}               % generalized-instrument branch map
\newcommand{\redNmsm}{\mathcal{N}}              % POVM measurement channel
\newcommand{\Eprot}{\mathcal{E}}             % n-round protocol channel
\newcommand{\Grnd}{\mathsf{G}}               % one round map of a protocol
\newcommand{\Kfin}{\mathsf{K}}               % end-of-protocol processing
\newcommand{\Crit}{\mathsf{sec}}             % value of a security criterion
\newcommand{\Pflag}{\Pi_{\Flag}}             % flag-sector projector on Q
\newcommand{\Pleq}{\Pi_{\leqN}}              % preserved-block projector
\newcommand{\Pgt}{\Pi_{\gtN}}                % above-cutoff projector
\newcommand{\Sph}{S^{\mathrm{ph}}}           % family of physical inputs
\newcommand{\Srch}{S^{\mathrm{rch}}}         % reachable squashed states
\newcommand{\Ssq}{S^{\sq}}                   % squashed admissible set
\newcommand{\Ffss}{F^{\mathrm{FSS}}}         % flag-state-squasher target
\newcommand{\Fwp}{F^{\mathrm{WP}}}           % multi-bound WPFSS target
\newcommand{\Mwp}{M^{\mathrm{WP}}}           % renormalized above-cutoff POVM
\newcommand{\Mdec}{{M'}^{\mathrm{ph}}}       % declared physical family on B
\newcommand{\Ppvm}{P^{\Flag}}                % flag PVM
\newcommand{\Lmp}{{\Lambda_{\mathrm{mp}}}}   % measure-and-prepare channel
\newcommand{\Lwp}{{\Lambda^{\mathrm{WP}}}}   % multi-bound WPFSS channel
\newcommand{\Lfss}{{\Lambda^{\mathrm{FSS}}}} % flag-state-squasher channel
\newcommand{\lmax}{\lambda_{\max}}           % maximum eigenvalue
\newcommand{\Lnoise}{\widetilde{\Lambda}}    % flag-noise channel
\newcommand{\cg}{\varpi}                     % declared coarse-graining map
\newcommand{\MCK}{\mathsf{MC}}               % multiclick raw-pattern set
\newcommand{\redChain}{\mathfrak{C}}            % reduction chain
\newcommand{\Lnk}{\ell}                      % a single typed link
\newcommand{\dsq}{\preceq}                      % the squashing relation
\newcommand{\opn}[1]{\left\lVert #1\right\rVert}           % operator (spectral) norm
\newcommand{\mxn}[1]{\left\lVert #1\right\rVert_{\max}}    % entrywise max modulus
\newcommand{\Fld}{\mathbb{F}}                   % exact ordered substrate field
\newcommand{\Fldi}{\mathbb{F}_{i}}              % its Gaussian extension F + iF
\newcommand{\lat}[1]{\big\lfloor #1\big\rfloor_{\kappa}} % round DOWN onto the lattice

% --- protocol-compiler ---------------------------------------------------------
\newcommand{\Spec}{\mathfrak{S}}             % protocol specification (compiler input)
\newcommand{\pcChain}{\mathfrak{R}}            % reduction chain (two typed segments)
\newcommand{\Chsrc}{\pcChain_{\mathrm{src}}}   % its Alice-side source-map segment
\newcommand{\Chrec}{\pcChain_{\mathrm{rec}}}   % its receiver-side squashing segment
\newcommand{\rvirt}{\rho^{\mathrm{v}}}       % virtual (PRE-source-map) signal state
\newcommand{\Fph}{F^{\mathrm{ph}}}           % physical (pre-reduction) POVM element
\newcommand{\Vph}{\mathcal{V}^{\mathrm{ph}}} % physical outcome alphabet
\newcommand{\Tpub}{\mathcal{T}_{\mathrm{pub}}}% public-choice alphabet
\newcommand{\Uset}{\mathcal{U}}              % residual outcome alphabet
\newcommand{\idxm}{\mathsf{idx}}             % label-to-index map
\newcommand{\Asig}{A'}                       % transmitted (signal) system -- the REAL one
\newcommand{\Avirt}{A''}                     % virtual transmitted system (source descriptor of Xi_1)
\newcommand{\Abar}{\bar{A}}                  % Alice's measured register
\newcommand{\Ahat}{\hat{A}}                  % Alice's shield register
\newcommand{\pcNmsm}{\mathcal{N}_{\mathsf A}}  % Alice's measurement channel (see the RULES block)
\newcommand{\Phip}{\Phi^{+}}                 % maximally entangled vector label
\newcommand{\Smap}{\mathfrak{s}}             % the emitted sector map (fraktur: a compile-time object)

% --- receiver ------------------------------------------------------------------
\newcommand{\Modes}{\mathbb{M}}              % declared ordered mode set
\newcommand{\Inp}{\mathbb{I}}                % declared input modes
\newcommand{\InpC}{\mathbb{I}_{\Chan}}       % channel input modes
\newcommand{\InpS}{\mathbb{I}_{\Prb}}        % probe (source) input modes
\newcommand{\Env}{\mathbb{E}}                % dilation (environment) modes
\newcommand{\Port}[1]{\mathbb{P}_{#1}}       % mode set of detector port #1
\newcommand{\Emd}[1]{E_{#1}}                 % coordinate projector on modes #1
\newcommand{\Gmat}{G}                        % instrument (transfer) matrix
\newcommand{\Fock}{\mathfrak{F}}             % Fock / second-quantization functor
\newcommand{\Symn}[1]{\mathrm{Sym}^{#1}}     % symmetric tensor power
\newcommand{\normord}[1]{{:}\mspace{2mu}#1\mspace{2mu}{:}}   % normal ordering
\newcommand{\rk}{\operatorname{rank}}        % rank (physics.sty owns \rank)
\newcommand{\SCK}{\mathsf{SC}}               % single-click pattern set
\newcommand{\Recv}{\mathfrak{R}}             % receiver instrument tuple
\newcommand{\Cat}{\mathfrak{C}}              % component catalogue
\newcommand{\Grph}{\mathfrak{E}}             % declared element graph
\newcommand{\defic}{\mathbb{L}}              % loss deficiency operator
\newcommand{\Pvac}[1]{\Pi_{\vac}(#1)}        % joint no-click element on ports #1
\newcommand{\pmc}{p_{\mc}}                   % uniform high-photon mc weight
\newcommand{\Bnd}{\beta}                     % the scalar bound function beta(n)
\newcommand{\Rot}{R}                         % planar rotation matrix
\newcommand{\Wplate}{W}                      % retarder (wave-plate) matrix
\newcommand{\Bsp}{B}                         % beamsplitter core matrix

% --- statistics ----------------------------------------------------------------
\newcommand{\honest}{p^{\mathrm{hon}}}          % honest per-round distribution
\newcommand{\ecomp}{\varepsilon_{\mathrm{comp}}}% completeness budget
\newcommand{\simplex}[1]{\Delta(#1)}            % probability simplex
\newcommand{\Chatcon}{\Chatset_{\mathrm{con}}}  % constrained outcomes
\newcommand{\Ireg}[3]{\mathcal{I}_{#1}(#2,#3)}  % regularized incomplete beta
\newcommand{\Ncount}[1]{\mathsf{N}_{#1}}        % per-outcome count
\newcommand{\qlo}{\mathbf{q}^{-}}               % box lower bounds
\newcommand{\qhi}{\mathbf{q}^{+}}               % box upper bounds

\makeatletter
\newcommand*{\ClaimHeadFont}{\bfseries}   % switched by prominence
\newcommand*{\ClaimNoteFont}{\mdseries}
\newcommand*{\ClaimShortId}{}
\newcommand*{\SetShortId}[1]{\renewcommand*{\ClaimShortId}{#1}}
\newcommand*{\RestatedName}{}
% Upright bodies throughout: italic is reserved for non-normative text.
% \@endparpenalty\@M keeps the meta line on the statement's page.
\newtheoremstyle{claim}{\topsep}{\topsep}{\normalfont\@endparpenalty\@M}{}%
  {\ClaimHeadFont}{.}{.5em}{\thmname{#1}\thmnumber{ #2}\thmnote{ {\ClaimNoteFont(#3)}}}
\newtheoremstyle{claimass}{\topsep}{\topsep}{\normalfont\@endparpenalty\@M}{}%
  {\ClaimHeadFont}{.}{.5em}{\thmname{#1}\thmnumber{ #2}%
  \ifx\ClaimShortId\@empty\else\ [\textsf{\ClaimShortId}]\fi%
  \thmnote{ {\ClaimNoteFont(#3)}}}
\newtheoremstyle{claimremark}{.5\topsep}{.5\topsep}{\small\@endparpenalty\@M}{}%
  {\itshape}{.}{.5em}{\thmname{#1}\thmnumber{ #2}\thmnote{ {\upshape(#3)}}}
\newtheoremstyle{claimreq}{\topsep}{\topsep}{\normalfont\@endparpenalty\@M}{}%
  {\color{reqAccent}\scshape}{.}{.5em}%
  {\thmname{#1}\thmnumber{ #2}\thmnote{ {\upshape(#3)}}}
\makeatother
\theoremstyle{claim}
\newtheorem{theorem}{Theorem}[section]
\newtheorem{lemma}[theorem]{Lemma}
\newtheorem{proposition}[theorem]{Proposition}
\newtheorem{corollary}[theorem]{Corollary}
\newtheorem{fact}[theorem]{Fact}
\newtheorem{definition}[theorem]{Definition}
\newtheorem*{restatedclaim}{\RestatedName}
\theoremstyle{claimass}
\newtheorem{assumption}[theorem]{Assumption}
\newtheorem*{restatedassumption}{\RestatedName}
\theoremstyle{claimremark}
\newtheorem{remark}[theorem]{Remark}
\newtheorem{example}[theorem]{Example}
\theoremstyle{claimreq}
\newtheorem{requirement}[theorem]{Requirement}
\newtheorem*{restatedrequirement}{\RestatedName}
%% --- 8. Prominence: main results a margin rule and a bold heading ----------
\newtcolorbox{mainclaim}{enhanced jigsaw, breakable, blanker,
  borderline west={1.5pt}{-8pt}{mainRule}, left=0pt, right=0pt, top=0pt, bottom=0pt,
  boxsep=0pt, before skip=\topsep, after skip=\topsep, parbox=false,
  before upper={\renewcommand*{\ClaimNoteFont}{\bfseries}}}
\newenvironment{technicalclaim}{\renewcommand*{\ClaimHeadFont}{\mdseries\color{techGrey}}%
  \renewcommand*{\ClaimNoteFont}{}}{}
%% --- 9. Non-normative text, the meta line, navigation ------------------------
\makeatletter
% Tie what follows to this paragraph: no page break, no list penalty before
% the next theorem head (the \@nobreak route a section heading uses).
\newcommand{\keepwithnext}{\par\nobreak\vspace{3pt}\@nobreaktrue
  \everypar{\if@nobreak\@nobreakfalse\clubpenalty\@M\else\clubpenalty\@clubpenalty\everypar{}\fi}}
\newcommand{\claimsummary}[2]{% #1 claim id (named destination), #2 summary
  \par\addvspace{\bigskipamount}\Needspace{6\baselineskip}%
  {\small\itshape\color{nnGrey}\noindent\hypertarget{claim-#1}{}#2\par}\keepwithnext}
\newcommand{\claimanchor}[1]{% a claim without a summary: the destination alone
  \par\addvspace{\bigskipamount}\Needspace{6\baselineskip}%
  \phantomsection\hypertarget{claim-#1}{}\keepwithnext}
\newcommand{\claimleadin}[1]{{\color{nnGrey}\noindent#1\par}\keepwithnext}
\newcommand{\claimmeta}[1]{\par\nobreak\vspace{2pt}%
  {\footnotesize\color{metaGrey}\raggedright\interlinepenalty\@M\noindent#1\par}%
  \addvspace{\medskipamount}}
% After an assumption: the sentence generated from its verification field; after
% any claim with one: its plain-text reading (user_statement). Both
% non-normative, both grey.
\newcommand{\claimverification}[1]{\par\nobreak\vspace{2pt}%
  {\footnotesize\color{metaGrey}\noindent#1\par}\nobreak}
\newcommand{\claimreading}[1]{\par\nobreak\vspace{2pt}%
  {\footnotesize\color{nnGrey}\noindent\emph{Reading:}\enspace#1\par}\nobreak}
\makeatother
\newcommand{\claimbookmark}[4]{% #1 id, #2 kind word, #3 label, #4 plain title
  \bookmark[dest=claim-#1,level=2]{#2 \getrefnumber{#3} #4}}
\newcommand{\metasep}{\unskip\enspace\textperiodcentered\enspace\ignorespaces}
\newcommand{\metalabel}[1]{\emph{#1}\enspace}
\newcommand{\sectionintro}[1]{{\color{nnGrey}\noindent#1\par}\smallskip}
\newcommand{\sectionimports}[1]{{\footnotesize\color{metaGrey}\raggedright\noindent\metalabel{Imports}#1\par}\smallskip}
\newenvironment{overview}{\par\color{nnGrey}\itshape\noindent\ignorespaces}{\par}
% References: the whole "Def. 2.1" is the link; ids are the fallback text of
% every link into another PDF (a relative URI to a named destination).
\newcommand{\mref}[2]{\hyperref[#2]{#1~\ref*{#2}}}
\newcommand{\aref}[2]{\hyperref[#2]{\textsf{#1}}}
\DeclareUrlCommand{\claimid}{\urlstyle{sf}\def\UrlBreaks{\do\/\do\:\do\-\do\.\do\_}%
  \def\UrlBigBreaks{}\def\UrlNoBreaks{}}
\newcommand{\extref}[2]{\hyperref{#1}{}{claim-#2}{\claimid{#2}}}
\newcommand{\extaref}[3]{\hyperref{#1}{}{claim-#2}{\textsf{#3}}}
\newcommand{\coderef}[2]{\href{#1}{\claimid{#2}}}
% Claim-body commands (spec/claim-format.md section 4): \xref{<full id>} prints
% the kind, title and id of a claim of another contract, linked to its stable
% destination in that contract's PDF (one \DeclareXref per cited claim below);
% \term[gloss]{phrase} marks a defined phrase for the index (the gloss feeds
% the Terms table); \tref[phrase]{text} is a use of a defined term. The
% assembler rewrites both in the body into \termdef (the definition's named
% destination term-<slug>) and \termref / \exttermref (a link to it, in this
% or another contract's PDF); the fallbacks below serve text it copies
% elsewhere, such as the digest. \iface is an interface key.
\makeatletter
\newcommand{\DeclareXref}[4]{% id, kind word, title, contract PDF
  \@namedef{xref@#1}{\hyperref{#4}{}{claim-#1}{#2 ``#3'' ({\small\claimid{#1}})}}}
\newcommand{\xref}[1]{\@ifundefined{xref@#1}%
  {\textbf{[unresolved \nolinkurl{#1}]}\@latex@warning{Unresolved xref #1}}%
  {\@nameuse{xref@#1}}}
\makeatother
\DeclareUrlCommand{\iface}{\urlstyle{tt}}
\newcommand{\term}[2][]{\emph{#2}\index{#2}}
\newcommand{\tref}[2][]{#2}
\newcommand{\termdef}[2]{\hypertarget{#1}{}\emph{#2}\index{#2}}
\newcommand{\termref}[2]{\hyperlink{#1}{#2}}
\newcommand{\exttermref}[3]{\hyperref{#1}{}{#2}{#3}}
\hypersetup{pdftitle={MEAT key-rate engine (meat-core) -- contract 1.32.0}}

% Claims of other contracts cited here (\xref, [[...]]).
\DeclareXref{foundations/ass:snapped-measurement}{Assumption}{The device measures the snapped family where the run reads the snap}{../foundations/contract.pdf}
\DeclareXref{foundations/cor:least-eigenvalue-over-enclosure}{Corollary}{Least eigenvalue and rank over an enclosure}{../foundations/contract.pdf}
\DeclareXref{foundations/cor:rank-over-factor-enclosure}{Corollary}{Rank over a factor enclosure}{../foundations/contract.pdf}
\DeclareXref{foundations/def:binary64}{Definition}{Binary64 numbers, declared values and dyadic lattices}{../foundations/contract.pdf}
\DeclareXref{foundations/def:declared-order}{Definition}{Sums in a declared order and sums that round nowhere}{../foundations/contract.pdf}
\DeclareXref{foundations/def:evaluated-off-snap}{Definition}{A reference handed beside a snap; a block evaluated off its snap}{../foundations/contract.pdf}
\DeclareXref{foundations/def:exact-island}{Definition}{The exact island}{../foundations/contract.pdf}
\DeclareXref{foundations/def:exact-reference}{Definition}{Reference of a measurement family, exact or enclosed}{../foundations/contract.pdf}
\DeclareXref{foundations/def:factor-enclosure}{Definition}{Factor enclosure of a block}{../foundations/contract.pdf}
\DeclareXref{foundations/def:kl}{Definition}{Classical divergences and the quantum relative entropy}{../foundations/contract.pdf}
\DeclareXref{foundations/def:pinching}{Definition}{Operator conventions and pinching}{../foundations/contract.pdf}
\DeclareXref{foundations/def:renyi}{Definition}{Sandwiched Rényi divergence and conditional entropies}{../foundations/contract.pdf}
\DeclareXref{foundations/def:round-model}{Definition}{Round model, attack sequences and attack sets}{../foundations/contract.pdf}
\DeclareXref{foundations/def:sector-map}{Definition}{Sector maps and block-list operators}{../foundations/contract.pdf}
\DeclareXref{foundations/def:snap}{Definition}{Snapping declared data onto a lattice}{../foundations/contract.pdf}
\DeclareXref{foundations/def:tiers}{Definition}{The tiers of a reported key rate}{../foundations/contract.pdf}
\DeclareXref{foundations/def:tolerance}{Definition}{The declared tolerance and the defects it bounds}{../foundations/contract.pdf}
\DeclareXref{foundations/fact:concavity}{Fact}{Concavity of the pinched trace-power functional}{../foundations/contract.pdf}
\DeclareXref{foundations/fact:pinch-id}{Fact}{Pinched trace-power form of the conditional Rényi entropy}{../foundations/contract.pdf}
\DeclareXref{foundations/lem:floor-stricter}{Lemma}{The positivity floor against entrywise closeness to the cone}{../foundations/contract.pdf}
\DeclareXref{foundations/lem:kernel}{Lemma}{Extended divergence kernel}{../foundations/contract.pdf}
\DeclareXref{foundations/lem:least-eigenvalue-bound}{Lemma}{Least-eigenvalue lower bound from an approximate eigenbasis}{../foundations/contract.pdf}
\DeclareXref{foundations/lem:matrix-function-perturbation}{Lemma}{Perturbation bounds for matrix functions}{../foundations/contract.pdf}
\DeclareXref{foundations/lem:pinchmono}{Lemma}{Pinching lowers the trace power}{../foundations/contract.pdf}
\DeclareXref{foundations/lem:recentred-enclosure}{Lemma}{Zeroing the centres an enclosure cannot tell from zero}{../foundations/contract.pdf}
\DeclareXref{foundations/lem:register-compression}{Lemma}{Compressing a register onto a finite-dimensional subspace}{../foundations/contract.pdf}
\DeclareXref{foundations/rem:noise-floor}{Remark}{What domination controls, and why it needs a noise floor}{../foundations/contract.pdf}
\DeclareXref{foundations/req:consumer-decides-exactly}{Requirement}{Exact decision by every consumer}{../foundations/contract.pdf}
\DeclareXref{foundations/req:declared-tolerance}{Requirement}{One declared tolerance per instance}{../foundations/contract.pdf}
\DeclareXref{foundations/req:exact-emission}{Requirement}{Exact verification before emission}{../foundations/contract.pdf}
\DeclareXref{foundations/req:exact-island}{Requirement}{Exact arithmetic on the island, unbudgeted rounding elsewhere}{../foundations/contract.pdf}
\DeclareXref{foundations/req:norm-conversion}{Requirement}{Explicit conversion between norms}{../foundations/contract.pdf}
\DeclareXref{foundations/req:rounding-direction}{Requirement}{Rounding direction of emitted bounds}{../foundations/contract.pdf}
\DeclareXref{foundations/req:snap-reference}{Requirement}{Snaps of a measurement family: kept, or moved and disclosed}{../foundations/contract.pdf}
\DeclareXref{foundations/req:tiers}{Requirement}{Reporting of the tiers}{../foundations/contract.pdf}
\DeclareXref{foundations/req:tolerance-scope}{Requirement}{Scope of the declared tolerance}{../foundations/contract.pdf}
\DeclareXref{protocol-compiler/ass:alice-meas}{Assumption}{Alice's declared measurement where no source certificate is demanded}{../protocol-compiler/contract.pdf}
\DeclareXref{protocol-compiler/ass:atk}{Assumption}{Attack-set licence}{../protocol-compiler/contract.pdf}
\DeclareXref{protocol-compiler/ass:link}{Assumption}{Soundness of every reduction link}{../protocol-compiler/contract.pdf}
\DeclareXref{protocol-compiler/ass:rhom}{Assumption}{Round homogeneity of the declaration}{../protocol-compiler/contract.pdf}
\DeclareXref{protocol-compiler/ass:spec}{Assumption}{Specification faithful to the apparatus}{../protocol-compiler/contract.pdf}
\DeclareXref{protocol-compiler/def:compile}{Definition}{The compilation map}{../protocol-compiler/contract.pdf}
\DeclareXref{protocol-compiler/def:declaration}{Definition}{Protocol declaration}{../protocol-compiler/contract.pdf}
\DeclareXref{protocol-compiler/def:exact}{Definition}{Exact admissibility of declared operator data}{../protocol-compiler/contract.pdf}
\DeclareXref{protocol-compiler/def:feasible}{Definition}{Feasible set}{../protocol-compiler/contract.pdf}
\DeclareXref{protocol-compiler/def:homog}{Definition}{Licensed constraint list}{../protocol-compiler/contract.pdf}
\DeclareXref{protocol-compiler/def:target}{Definition}{Target conditions of a well-formed declaration}{../protocol-compiler/contract.pdf}
\DeclareXref{protocol-compiler/def:validator}{Definition}{The emitted validator}{../protocol-compiler/contract.pdf}
\DeclareXref{protocol-compiler/def:weighted}{Definition}{Choice-weighted families}{../protocol-compiler/contract.pdf}
\DeclareXref{protocol-compiler/lem:weighted}{Lemma}{Properties of the choice-weighted families}{../protocol-compiler/contract.pdf}
\DeclareXref{protocol-compiler/lem:weighted-factor}{Lemma}{A factor enclosure through a setting-dependent weighting}{../protocol-compiler/contract.pdf}
\DeclareXref{protocol-compiler/rem:block-restriction}{Remark}{Restriction of the feasible set to block-diagonal states}{../protocol-compiler/contract.pdf}
\DeclareXref{protocol-compiler/rem:engine-admission}{Remark}{Relation to the consuming engine's admission conditions}{../protocol-compiler/contract.pdf}
\DeclareXref{protocol-compiler/rem:obligations-discharged}{Remark}{Discharge of the chain obligations by the reductions module}{../protocol-compiler/contract.pdf}
\DeclareXref{protocol-compiler/rem:party-local-soundness}{Remark}{Soundness of a party-local entry}{../protocol-compiler/contract.pdf}
\DeclareXref{protocol-compiler/req:no-fixed-fraction}{Requirement}{No fixed fraction outside the measurement operators}{../protocol-compiler/contract.pdf}
\DeclareXref{protocol-compiler/req:reference-companion-compiled}{Requirement}{The reference companion compiled beside the declaration}{../protocol-compiler/contract.pdf}
\DeclareXref{protocol-compiler/thm:band-a}{Theorem}{The Band-A route}{../protocol-compiler/contract.pdf}
\DeclareXref{protocol-compiler/thm:feasibility}{Theorem}{Feasibility by an explicit or a verified candidate}{../protocol-compiler/contract.pdf}
\DeclareXref{receiver/ass:trust}{Assumption}{Trusted and exactly characterised receiver}{../receiver/contract.pdf}
\DeclareXref{receiver/req:eig-ledger-row}{Requirement}{Invoking settings on the eigenvalue-slack row of the ledger}{../receiver/contract.pdf}
\DeclareXref{receiver/req:evaluation-label}{Requirement}{The evaluation label on every artifact}{../receiver/contract.pdf}
\DeclareXref{receiver/thm:enclosed-weight}{Theorem}{The weight on enclosed bound data bounds the multiclick element}{../receiver/contract.pdf}
\DeclareXref{reductions/ass:recv}{Assumption}{Supplier artefact whose claims hold of the device}{../reductions/contract.pdf}
\DeclareXref{reductions/ass:tsm}{Assumption}{Output-state test statistics}{../reductions/contract.pdf}
\DeclareXref{reductions/cor:companion-factor-enclosure}{Corollary}{Factor enclosures carried to the reference companion}{../reductions/contract.pdf}
\DeclareXref{reductions/cor:reference-companion}{Corollary}{The reference companion of a reduction instance}{../reductions/contract.pdf}
\DeclareXref{reductions/cor:sequential}{Corollary}{Output-state equality against adaptive attack sequences}{../reductions/contract.pdf}
\DeclareXref{reductions/def:attset}{Definition}{Admissible attack sets in the consumer's format}{../reductions/contract.pdf}
\DeclareXref{reductions/def:protocol}{Definition}{Measurement-factorizing $n$-round protocol}{../reductions/contract.pdf}
\DeclareXref{reductions/def:reduction-instance}{Definition}{Reduction instance}{../reductions/contract.pdf}
\DeclareXref{reductions/def:roundmodel}{Definition}{Round model and attack sequences, typed by source mode}{../reductions/contract.pdf}
\DeclareXref{reductions/lem:row-relaxation}{Lemma}{A rounded constraint entry with its constant relaxed}{../reductions/contract.pdf}
\DeclareXref{reductions/lem:source-transfer}{Lemma}{Exactness of the source map in trusted-source mode}{../reductions/contract.pdf}
\DeclareXref{reductions/rem:why-exact}{Remark}{Exactness of the reductions and generality of the transfer}{../reductions/contract.pdf}
\DeclareXref{reductions/req:recv-row-provenance}{Requirement}{Omission of the supplier row for a receiver artefact the run built}{../reductions/contract.pdf}
\DeclareXref{reductions/req:reference-companion}{Requirement}{Emitting the reference companion beside the artefact}{../reductions/contract.pdf}
\DeclareXref{reductions/thm:licence}{Theorem}{The attack-set licence}{../reductions/contract.pdf}
\DeclareXref{reductions/thm:two-sided}{Theorem}{The two-sided reduction in trusted-source mode}{../reductions/contract.pdf}
\DeclareXref{reductions/wpfss/def:weight-route}{Definition}{The weight route}{../reductions/contract.pdf}
\DeclareXref{statistics/ass:reach}{Assumption}{An honest law produced by a feasible state}{../statistics/contract.pdf}
\DeclareXref{statistics/binomial-box/thm:completeness-bound}{Theorem}{The completeness bound}{../statistics/contract.pdf}
\DeclareXref{statistics/singleton/thm:singleton-completeness}{Theorem}{Completeness of a singleton acceptance set}{../statistics/contract.pdf}
\DeclareXref{transmitter/ass:ens}{Assumption}{Emission of the specified ensemble}{../transmitter/contract.pdf}
\DeclareXref{transmitter/ass:seq}{Assumption}{Sequential round structure}{../transmitter/contract.pdf}
\DeclareXref{transmitter/ass:src}{Assumption}{Retention of the system A in trusted-source mode}{../transmitter/contract.pdf}
\DeclareXref{transmitter/decoy/ass:wcp}{Assumption}{Emission of the declared weak-coherent-pulse states}{../transmitter/contract.pdf}
\DeclareXref{transmitter/def:source-segment}{Definition}{Source segment and its interface conditions}{../transmitter/contract.pdf}
\DeclareXref{transmitter/def:source-spec}{Definition}{Source specification}{../transmitter/contract.pdf}
\DeclareXref{transmitter/def:srcert}{Definition}{Source-replacement certificate}{../transmitter/contract.pdf}
\DeclareXref{transmitter/rem:source-mode}{Remark}{The two source modes}{../transmitter/contract.pdf}
\DeclareXref{transmitter/thm:certificate-sufficiency}{Theorem}{Checkability and sufficiency of the source-replacement certificate}{../transmitter/contract.pdf}

\begin{document}
\begin{center}
{\LARGE\bfseries MEAT key-rate engine (meat-core)}\\[4pt]
{\large Contract, version 1.32.0}
\end{center}
\begin{overview}
The MEAT core of the key-rate engine (the module meat-core) turns a finite protocol declaration, which the protocol compiler emits, and an acceptance set into a convex entropy program and a fixed-length key length; four of its main results form one chain from the single-round protocol channel the declaration generates to that key length, and the fifth is the security of the protocol that outputs it (\hyperlink{claim-meat-core/cor:composed-security}{Corollary~\ref*{cor:composed-security}}). Marginal-constrained entropy accumulation bounds the entropy of the accepted raw key from below by n times a single-round rate, less a penalty in the acceptance probability, against every adaptive attack with a finite-dimensional quantum memory (\hyperlink{claim-meat-core/thm:accumulated-entropy}{Theorem~\ref*{thm:accumulated-entropy}}). A licensed constraint list, one whose entries have a type the declared source mode admits, raises that rate only under the attack-set licence, the producer's certificate that every physical attack reduces to one that respects the list. The single-round rate is bounded below by the value h\_low (written as an underlined h\_\ensuremath{\alpha}) of a finite convex program over sector-block states (\hyperlink{claim-meat-core/thm:reduction}{Theorem~\ref*{thm:reduction}}), which a Choi parametrisation writes in solver-friendly variables without deciding the rank of Alice's marginal. The objective at any feasible point is an upper estimate of h\_low, never a lower bound on it (\hyperlink{claim-meat-core/lem:estimate}{Lemma~\ref*{lem:estimate}}), while any lower bound on h\_low, however obtained, yields a key length (\hyperlink{claim-meat-core/cor:keylength-bound}{Corollary~\ref*{cor:keylength-bound}}) through the \termref{term-fixed-length-post-processing-model}{fixed-length post-processing model}: the key is secret (\hyperlink{claim-meat-core/thm:keylength}{Theorem~\ref*{thm:keylength}}) and correct, and over an authentication scheme the protocol is secure with the authentication error added and its key consumption subtracted (\hyperlink{claim-meat-core/cor:composed-security}{Corollary~\ref*{cor:composed-security}}), against an adversary whose memory has any dimension. In trusted-source mode Alice's marginal is fixed, which needs the source licence of the transmitter module, the attestation that her system never reaches the adversary; untrusted-source mode leaves it free and never gives a larger value. The Frank-Wolfe submodule supplies the first-order lower bound on h\_low, with both linear minima bounded exactly, evaluates it, where the chain carries one, on the reference companion, the declaration of the reference model compiled beside the snapped one, and certifies it where every step from the declaration's data can be enclosed rigorously. The module admits declared data only when they are exactly admissible, the companion aside, which it admits beside its declaration and searches only in a short polish from the declaration's points; it characterises no device and constructs no honest statistics, and only the submodule's certified evaluation bounds the rounding of a computed number; where the honest statistics are supplied rather than computed from the declared devices, or the caller declares a constraint entry, it bounds the mismatch between the acceptance set and the feasible statistics from below, records the cap this proves on the probability that devices as declared pass the acceptance test, and refuses the run where that cap is at most a threshold, one half by default.
\end{overview}
\medskip
\noindent{\small\textcolor{nnGrey}{\itshape Grey text}, here and throughout, is non-normative: it explains, and nothing in it is assumed or proved.}

\smallskip
\noindent{\small The \href{https://github.com/fraunhofersg/certiqs-key/blob/v0.1.0/spec/modules/meat-core/design.md}{design notes} (non-normative) explain why the module is built as it is and draw its processing flowchart. The dependencies between the sections are drawn in \href{graphs/overview.svg}{the section dependency matrix}, and \href{https://github.com/fraunhofersg/certiqs-key/blob/v0.1.0/spec/modules/meat-core/graphs/README.md}{the graph index} opens every section's dependency graph. The \href{digest.pdf}{statements digest} holds every definition, assumption, result and requirement without the proofs. \href{https://github.com/fraunhofersg/certiqs-key/blob/v0.1.0/harness/contract-reading.md}{How to read a contract} explains this layout and its links.}

{\small\tableofcontents}

\section*{Main results}\addcontentsline{toc}{section}{Main results}
\sectionintro{The module's most important results, those its overview names as what it establishes, with every assumption each rests on, directly or through the claims it uses. Each links to its statement and proof; the statements digest holds them in full.}
\par\noindent\hyperref[thm:accumulated-entropy]{\textbf{Theorem~\ref*{thm:accumulated-entropy}}}\enspace Accumulated entropy bound under adaptive attacks\par\nobreak
{\small\itshape\color{nnGrey}\noindent On acceptance, the raw key's R\'enyi entropy given the adversary's information is at least n times the single-round rate less \ensuremath{\alpha}/(\ensuremath{\alpha}\ensuremath{-}1)\ensuremath{\cdot}log\ensuremath{_{2}}(1/Pr[accept]), for any adaptive attack with a finite-dimensional memory. Trusted-source mode reads the product marginal of \hyperlink{claim-meat-core/lem:product-marginal}{Lemma~\ref*{lem:product-marginal}}; untrusted mode needs none.\par}\nobreak
\claimmeta{\metalabel{Rests on}\extaref{../protocol-compiler/contract.pdf}{protocol-compiler/ass:atk}{A-ATK}, \extaref{../transmitter/contract.pdf}{transmitter/ass:ens}{A-ENS}, \aref{A-LAB}{ass:lab}, \extaref{../protocol-compiler/contract.pdf}{protocol-compiler/ass:rhom}{A-RHOM}, \aref{A-SEED}{ass:seed}, \extaref{../transmitter/contract.pdf}{transmitter/ass:seq}{A-SEQ}, \extaref{../transmitter/contract.pdf}{transmitter/ass:src}{A-SRC}}
\par\noindent\hyperref[thm:reduction]{\textbf{Theorem~\ref*{thm:reduction}}}\enspace Reduction of the rate to the convex program\par\nobreak
{\small\itshape\color{nnGrey}\noindent For a declaration that adds no constraints on the state (an empty constraint list), in either source mode, the value h\_low of the finite convex program bounds the single-round rate from below. Every step of the reduction is an equality or lowers the bound, so the computable program is conservative.\par}\nobreak
\claimmeta{\metalabel{Rests on}no assumption}
\par\noindent\hyperref[thm:keylength]{\textbf{Theorem~\ref*{thm:keylength}}}\enspace Secrecy of the key from the accumulated entropy\par\nobreak
{\small\itshape\color{nnGrey}\noindent For a R\'enyi order \ensuremath{\alpha} in (1, 2), a key of length \ensuremath{\ell} \ensuremath{\le} n\ensuremath{\cdot}h\_low \ensuremath{-} \ensuremath{\lambda}\_EC \ensuremath{-} \ensuremath{\lceil}log\ensuremath{_{2}}(1/\ensuremath{\varepsilon}\_EV)\ensuremath{\rceil} \ensuremath{-} \ensuremath{\alpha}/(\ensuremath{\alpha}\ensuremath{-}1)\ensuremath{\cdot}log\ensuremath{_{2}}(1/\ensuremath{\varepsilon}\_PA) + 1 is \ensuremath{\varepsilon}\_PA-secret on the accepted state at every all-at-once finite-register attack the reduction chain carries onto the round model (\hyperref{../protocol-compiler/contract.pdf}{}{claim-protocol-compiler/ass:link}{A-LINK}). \hyperlink{claim-meat-core/cor:composed-security}{Corollary~\ref*{cor:composed-security}} extends it to every attack.\par}\nobreak
\claimmeta{\metalabel{Rests on}\extaref{../protocol-compiler/contract.pdf}{protocol-compiler/ass:alice-meas}{A-ALICE-MEAS}, \extaref{../protocol-compiler/contract.pdf}{protocol-compiler/ass:atk}{A-ATK}, \aref{A-CONFORM}{ass:conform}, \extaref{../transmitter/contract.pdf}{transmitter/ass:ens}{A-ENS}, \aref{A-FIXED}{ass:fixed}, \aref{A-HASH}{ass:hash}, \aref{A-KEYREG}{ass:keyreg}, \aref{A-LAB}{ass:lab}, \extaref{../protocol-compiler/contract.pdf}{protocol-compiler/ass:link}{A-LINK}, \extaref{../protocol-compiler/contract.pdf}{protocol-compiler/ass:rhom}{A-RHOM}, \aref{A-SEED}{ass:seed}, \extaref{../transmitter/contract.pdf}{transmitter/ass:seq}{A-SEQ}, \extaref{../protocol-compiler/contract.pdf}{protocol-compiler/ass:spec}{A-SPEC}, \extaref{../transmitter/contract.pdf}{transmitter/ass:src}{A-SRC}}
\par\noindent\hyperref[cor:composed-security]{\textbf{Corollary~\ref*{cor:composed-security}}}\enspace Security of the protocol and its net key\par\nobreak
{\small\itshape\color{nnGrey}\noindent Secrecy and correctness compose to (\ensuremath{\varepsilon}\_PA + \ensuremath{\varepsilon}\_EV)-security with honest authentication; over an authentication scheme within \ensuremath{\varepsilon}\_auth that consumes k\_auth bits of pre-shared key, the protocol is (\ensuremath{\varepsilon}\_PA + \ensuremath{\varepsilon}\_EV + 2\ensuremath{\varepsilon}\_auth)-secure; when both parties output a key, its net key gain is \ensuremath{\ell} \ensuremath{-} k\_auth bits.\par}\nobreak
\claimmeta{\metalabel{Rests on}\extaref{../protocol-compiler/contract.pdf}{protocol-compiler/ass:alice-meas}{A-ALICE-MEAS}, \extaref{../protocol-compiler/contract.pdf}{protocol-compiler/ass:atk}{A-ATK}, \aref{A-AUTH}{ass:auth}, \aref{A-CONFORM}{ass:conform}, \extaref{../transmitter/contract.pdf}{transmitter/ass:ens}{A-ENS}, \aref{A-FIXED}{ass:fixed}, \aref{A-HASH}{ass:hash}, \aref{A-KEYREG}{ass:keyreg}, \aref{A-LAB}{ass:lab}, \extaref{../protocol-compiler/contract.pdf}{protocol-compiler/ass:link}{A-LINK}, \extaref{../protocol-compiler/contract.pdf}{protocol-compiler/ass:rhom}{A-RHOM}, \aref{A-SEED}{ass:seed}, \extaref{../transmitter/contract.pdf}{transmitter/ass:seq}{A-SEQ}, \extaref{../protocol-compiler/contract.pdf}{protocol-compiler/ass:spec}{A-SPEC}, \extaref{../transmitter/contract.pdf}{transmitter/ass:src}{A-SRC}}
\par\noindent\hyperref[cor:keylength-bound]{\textbf{Corollary~\ref*{cor:keylength-bound}}}\enspace Key length from any lower bound on the value\par\nobreak
{\small\itshape\color{nnGrey}\noindent Any h \ensuremath{\le} h\_low may replace h\_low in the key-length bound through max\{0, h\}, with the same security statement. This is where every solver obtains its key length; a bound h \ensuremath{\le} 0 yields no positive key when \ensuremath{\varepsilon}\_EV \textless{} 1.\par}\nobreak
\claimmeta{\metalabel{Rests on}\extaref{../protocol-compiler/contract.pdf}{protocol-compiler/ass:alice-meas}{A-ALICE-MEAS}, \extaref{../protocol-compiler/contract.pdf}{protocol-compiler/ass:atk}{A-ATK}, \aref{A-AUTH}{ass:auth}, \aref{A-CONFORM}{ass:conform}, \extaref{../transmitter/contract.pdf}{transmitter/ass:ens}{A-ENS}, \aref{A-FIXED}{ass:fixed}, \aref{A-HASH}{ass:hash}, \aref{A-KEYREG}{ass:keyreg}, \aref{A-LAB}{ass:lab}, \extaref{../protocol-compiler/contract.pdf}{protocol-compiler/ass:link}{A-LINK}, \extaref{../protocol-compiler/contract.pdf}{protocol-compiler/ass:rhom}{A-RHOM}, \aref{A-SEED}{ass:seed}, \extaref{../transmitter/contract.pdf}{transmitter/ass:seq}{A-SEQ}, \extaref{../protocol-compiler/contract.pdf}{protocol-compiler/ass:spec}{A-SPEC}, \extaref{../transmitter/contract.pdf}{transmitter/ass:src}{A-SRC}}
\par\noindent\hyperref[thm:firstorder]{\textbf{Theorem~\ref*{thm:firstorder}}}\enspace The estimate and the first-order lower bound\par\nobreak
{\small\itshape\color{nnGrey}\noindent In exact real arithmetic the best linearisation over the iterates, even ones feasible only up to windows, bounds h\_low from below given positive probabilities, maximal branch support, unperturbed gradients and bounds on both linear minima. An exactly feasible point's objective bounds it from above.\par}\nobreak
\claimmeta{\metalabel{Rests on}\extaref{../protocol-compiler/contract.pdf}{protocol-compiler/ass:atk}{A-ATK}, \extaref{../transmitter/contract.pdf}{transmitter/ass:ens}{A-ENS}, \aref{A-LAB}{ass:lab}, \extaref{../protocol-compiler/contract.pdf}{protocol-compiler/ass:rhom}{A-RHOM}, \aref{A-SEED}{ass:seed}, \extaref{../transmitter/contract.pdf}{transmitter/ass:seq}{A-SEQ}, \extaref{../transmitter/contract.pdf}{transmitter/ass:src}{A-SRC}}
\par\noindent\hyperref[cor:fo-dropped]{\textbf{Corollary~\ref*{cor:fo-dropped}}}\enspace First-order bound for the dropped program\par\nobreak
{\small\itshape\color{nnGrey}\noindent In either source mode the first-order bound holds for the dropped program (\hyperlink{claim-meat-core/def:dropped}{Definition~\ref*{def:dropped}}), with maximal support asked only of key-active pairs (acceptance value, joint sector) whose entropy is kept. It bounds h\_low from below, also at iterates feasible only up to windows.\par}\nobreak
\claimmeta{\metalabel{Rests on}\extaref{../protocol-compiler/contract.pdf}{protocol-compiler/ass:atk}{A-ATK}, \extaref{../transmitter/contract.pdf}{transmitter/ass:ens}{A-ENS}, \aref{A-LAB}{ass:lab}, \extaref{../protocol-compiler/contract.pdf}{protocol-compiler/ass:rhom}{A-RHOM}, \aref{A-SEED}{ass:seed}, \extaref{../transmitter/contract.pdf}{transmitter/ass:seq}{A-SEQ}, \extaref{../transmitter/contract.pdf}{transmitter/ass:src}{A-SRC}}
\par\noindent\hyperref[cor:fo-reference]{\textbf{Corollary~\ref*{cor:fo-reference}}}\enspace The first-order bound on the reference companion\par\nobreak
{\small\itshape\color{nnGrey}\noindent At any points, such as those of a search on the snapped declaration or of one continued on the companion, the bracket formed with the reference companion's objective, gradient, entries and linear minima bounds the companion's value from below; its key length is the reference model's.\par}\nobreak
\claimmeta{\metalabel{Rests on}\extaref{../protocol-compiler/contract.pdf}{protocol-compiler/ass:alice-meas}{A-ALICE-MEAS}, \extaref{../protocol-compiler/contract.pdf}{protocol-compiler/ass:atk}{A-ATK}, \aref{A-AUTH}{ass:auth}, \aref{A-CONFORM}{ass:conform}, \extaref{../reductions/contract.pdf}{reductions/ass:constdecl}{A-CONST}, \extaref{../transmitter/contract.pdf}{transmitter/ass:ens}{A-ENS}, \aref{A-FIXED}{ass:fixed}, \aref{A-HASH}{ass:hash}, \aref{A-KEYREG}{ass:keyreg}, \aref{A-LAB}{ass:lab}, \extaref{../protocol-compiler/contract.pdf}{protocol-compiler/ass:link}{A-LINK}, \extaref{../reductions/contract.pdf}{reductions/ass:meas}{A-MEAS}, \extaref{../protocol-compiler/contract.pdf}{protocol-compiler/ass:rhom}{A-RHOM}, \extaref{../reductions/contract.pdf}{reductions/ass:round}{A-ROUND-SQ}, \aref{A-SEED}{ass:seed}, \extaref{../transmitter/contract.pdf}{transmitter/ass:seq}{A-SEQ}, \extaref{../foundations/contract.pdf}{foundations/ass:snapped-measurement}{A-SNAP}, \extaref{../protocol-compiler/contract.pdf}{protocol-compiler/ass:spec}{A-SPEC}, \extaref{../transmitter/contract.pdf}{transmitter/ass:src}{A-SRC}}
\par\noindent\hyperref[cor:certified-keylength]{\textbf{Corollary~\ref*{cor:certified-keylength}}}\enspace Key length from the certified bracket\par\nobreak
{\small\itshape\color{nnGrey}\noindent The largest bracket formed by \hyperlink{claim-meat-core/frank-wolfe/lem:certified-bracket}{Lemma~\ref*{lem:certified-bracket}} over finitely many points is a first-order bound of \hyperlink{claim-meat-core/frank-wolfe/thm:firstorder}{Theorem~\ref*{thm:firstorder}}, so it is at most the value of the declaration it is formed on and gives that declaration's fixed-length key length exactly as h\_lb does.\par}\nobreak
\claimmeta{\metalabel{Rests on}\extaref{../protocol-compiler/contract.pdf}{protocol-compiler/ass:alice-meas}{A-ALICE-MEAS}, \extaref{../protocol-compiler/contract.pdf}{protocol-compiler/ass:atk}{A-ATK}, \aref{A-AUTH}{ass:auth}, \aref{A-CONFORM}{ass:conform}, \extaref{../transmitter/contract.pdf}{transmitter/ass:ens}{A-ENS}, \aref{A-FIXED}{ass:fixed}, \aref{A-HASH}{ass:hash}, \aref{A-KEYREG}{ass:keyreg}, \aref{A-LAB}{ass:lab}, \extaref{../protocol-compiler/contract.pdf}{protocol-compiler/ass:link}{A-LINK}, \extaref{../protocol-compiler/contract.pdf}{protocol-compiler/ass:rhom}{A-RHOM}, \aref{A-SEED}{ass:seed}, \extaref{../transmitter/contract.pdf}{transmitter/ass:seq}{A-SEQ}, \extaref{../protocol-compiler/contract.pdf}{protocol-compiler/ass:spec}{A-SPEC}, \extaref{../transmitter/contract.pdf}{transmitter/ass:src}{A-SRC}}

\section*{Assumptions at a glance}\addcontentsline{toc}{section}{Assumptions at a glance}
\sectionintro{Every assumption this contract states or its claims assume, with its category and the claims of this contract that assume it directly. An assumption of another contract links into that contract's PDF.}
{\footnotesize\begin{xltabular}{\linewidth}{@{}>{\raggedright\arraybackslash}p{1.6cm}>{\raggedright\arraybackslash}p{1.9cm}>{\raggedright\arraybackslash}X>{\raggedright\arraybackslash}p{2.3cm}>{\raggedright\arraybackslash}p{2.8cm}@{}}
\toprule
Assumption & Category & Title & Stated as & Assumed directly by\\
\midrule
\endhead
\extaref{../transmitter/contract.pdf}{transmitter/ass:ens}{A-ENS} & source & Emission of the specified ensemble & \extref{../transmitter/contract.pdf}{transmitter/ass:ens} & \mref{Lem.}{lem:product-marginal}, \mref{Thm.}{thm:accumulated-entropy}, \mref{Thm.}{thm:keylength}\\
\addlinespace
\extaref{../transmitter/contract.pdf}{transmitter/ass:src}{A-SRC} & source & Retention of the system A in trusted-source mode & \extref{../transmitter/contract.pdf}{transmitter/ass:src} & \mref{Lem.}{lem:range}, \mref{Lem.}{lem:product-marginal}, \mref{Thm.}{thm:accumulated-entropy}, \mref{Ex.}{ex:fss}, \mref{Rem.}{rem:not-licensed}, \mref{Rem.}{rem:modes-da1}, \mref{Rem.}{rem:cgate}\\
\addlinespace
\aref{A-LAB}{ass:lab} & laboratory & Each party's laboratory is sealed & \mref{Assumption}{ass:lab} & \mref{Lem.}{lem:product-marginal}, \mref{Thm.}{thm:accumulated-entropy}, \mref{Lem.}{lem:honest-attacks}, \mref{Thm.}{thm:keylength}\\
\addlinespace
\aref{A-AUTH}{ass:auth} & channel and adversary & Authentication of the classical channel by a keyed scheme & \mref{Assumption}{ass:auth} & \mref{Cor.}{cor:composed-security}\\
\addlinespace
\aref{A-EVE-DIM}{ass:eve-dim} & channel and adversary & The adversary's registers are finite-dimensional & \mref{Assumption}{ass:eve-dim} & ---\\
\addlinespace
\extaref{../protocol-compiler/contract.pdf}{protocol-compiler/ass:alice-meas}{A-ALICE-MEAS} & protocol structure & Alice's declared measurement where no source certificate is demanded & \extref{../protocol-compiler/contract.pdf}{protocol-compiler/ass:alice-meas} & \mref{Lem.}{lem:honest-attacks}, \mref{Thm.}{thm:keylength}\\
\addlinespace
\aref{A-FIXED}{ass:fixed} & protocol structure & Protocol parameters fixed before the run & \mref{Assumption}{ass:fixed} & \mref{Thm.}{thm:keylength}, \mref{Cor.}{cor:composed-security}\\
\addlinespace
\extaref{../protocol-compiler/contract.pdf}{protocol-compiler/ass:rhom}{A-RHOM} & protocol structure & Round homogeneity of the declaration & \extref{../protocol-compiler/contract.pdf}{protocol-compiler/ass:rhom} & \mref{Thm.}{thm:accumulated-entropy}\\
\addlinespace
\aref{A-SEPARATE}{ass:separate-executions} & protocol structure & Executions that share no key and no device state & \mref{Assumption}{ass:separate-executions} & \mref{Cor.}{cor:separate-executions}\\
\addlinespace
\extaref{../transmitter/contract.pdf}{transmitter/ass:seq}{A-SEQ} & protocol structure & Sequential round structure & \extref{../transmitter/contract.pdf}{transmitter/ass:seq} & \mref{Lem.}{lem:product-marginal}, \mref{Thm.}{thm:accumulated-entropy}, \mref{Lem.}{lem:honest-attacks}, \mref{Thm.}{thm:keylength}\\
\addlinespace
\extaref{../protocol-compiler/contract.pdf}{protocol-compiler/ass:spec}{A-SPEC} & protocol structure & Specification faithful to the apparatus & \extref{../protocol-compiler/contract.pdf}{protocol-compiler/ass:spec} & \mref{Thm.}{thm:keylength}\\
\addlinespace
\extaref{../reductions/contract.pdf}{reductions/ass:tsm}{A-TSM} & protocol structure & Output-state test statistics & \extref{../reductions/contract.pdf}{reductions/ass:tsm} & \mref{Rem.}{rem:not-licensed}\\
\addlinespace
\aref{A-SEED}{ass:seed} & local randomness & Private, independent and fresh local randomness & \mref{Assumption}{ass:seed} & \mref{Lem.}{lem:product-marginal}, \mref{Thm.}{thm:accumulated-entropy}, \mref{Lem.}{lem:ev-correctness}, \mref{Thm.}{thm:keylength}\\
\addlinespace
\aref{A-CONFORM}{ass:conform} & post processing & The implementation runs the analysed post-processing & \mref{Assumption}{ass:conform} & \mref{Lem.}{lem:ev-correctness}, \mref{Lem.}{lem:honest-attacks}, \mref{Thm.}{thm:keylength}, \mref{Cor.}{cor:composed-security}\\
\addlinespace
\aref{A-HASH}{ass:hash} & post processing & Ideal two-universal privacy amplification & \mref{Assumption}{ass:hash} & \mref{Lem.}{lem:ev-correctness}, \mref{Thm.}{thm:keylength}\\
\addlinespace
\aref{A-KEYREG}{ass:keyreg} & post processing & Dedicated key register and one-way reconciliation & \mref{Assumption}{ass:keyreg} & \mref{Thm.}{thm:keylength}\\
\addlinespace
\extaref{../protocol-compiler/contract.pdf}{protocol-compiler/ass:atk}{A-ATK} & supplied artefact & Attack-set licence & \extref{../protocol-compiler/contract.pdf}{protocol-compiler/ass:atk} & \mref{Lem.}{lem:range}, \mref{Rem.}{rem:no-constraints}, \mref{Thm.}{thm:constrained-accumulation}, \mref{Ex.}{ex:fss}, \mref{Rem.}{rem:not-licensed}, \mref{Thm.}{thm:keylength}, \mref{Cor.}{cor:constrained-key-length}, \mref{Rem.}{rem:cgate}\\
\addlinespace
\extaref{../protocol-compiler/contract.pdf}{protocol-compiler/ass:link}{A-LINK} & supplied artefact & Soundness of every reduction link & \extref{../protocol-compiler/contract.pdf}{protocol-compiler/ass:link} & \mref{Thm.}{thm:keylength}\\
\addlinespace
\aref{A-LMO}{ass:lmo} & numerical trust & Exact linear minimisation over the feasible set & \mref{Assumption}{ass:lmo} & ---\\
\bottomrule
\end{xltabular}}

\section*{Notation}\addcontentsline{toc}{section}{Notation}
\sectionintro{The symbols of the repository's notation catalogue (\href{../../notation.md}{spec/notation.md}) that this contract uses, in the catalogue's order, with their meaning.}
{\small\begin{xltabular}{\linewidth}{@{}>{\raggedright\arraybackslash}p{3.6cm}>{\raggedright\arraybackslash}X@{}}
\toprule
Symbol & Meaning\\
\midrule
\endhead
$\Hs{A}$ & Hilbert space of the system label $A$\\
$\Lin{A}$ & Bounded linear operators on $\mathcal{H}(A)$\\
$\Linn{A}{B}$ & Linear maps $\mathcal{H}(A)\to\mathcal{H}(B)$\\
$\Pos{A}$ & Positive semidefinite operators\\
$\Seq{A}$ & Normalized states\\
$\Sleq{A}$ & Subnormalized states\\
$\id$, $\idon{C}$ & Identity operator (contextual / named)\\
$\idmap$ & Identity channel\\
$\Decl$ & Protocol declaration $\big(A,\tau_A;\,Q,\mathbb{S};\,\{M_x\};\,\{F_v\};\,\mathfrak{P};\,a_A,a_B;\,\phi;\,r,\mathsf{h};\,\mathfrak{L}\big)$ \textemdash{} the tuple of \mref{Def.}{def:declaration}: the party partition $\mathfrak{P}$ is a slot, and the marginal slot $\tau_A$ takes a state or the no-marginal value (trusted- vs untrusted-source mode)\\
$\SecSet$ & Sector index set of $Q$\\
$\Qsys$ & The declared sector-carrying measured system\\
$\Party$ & The declared party partition: the readout maps $\pi_{\mathsf A}$, $\pi_{\mathsf B}$ onto the alphabets $\mathcal{O}_{\mathsf A}$, $\mathcal{O}_{\mathsf B}$, with a party label for $A$ and, for each sector $s$, a declared tensor factorisation of $\mathcal{H}(Q_s)$ with a party label per factor\\
$\Xset$ & Outcome alphabet of the measurement on $A$ \textemdash{} Alice's under the default partition\\
$\Vset$ & The measuring parties' declared joint outcome alphabet \textemdash{} Bob's alone under the default partition\\
$\CAset$, $\CBset$ & The alphabets of the first and second message of the transcript, which is the $L=2$ case\\
$\Kann$ & Transcript alphabet $\prod_{m=1}^{L}\mathcal{C}_m$, which at $L=2$ is $\mathcal{C}_A\times\mathcal{C}_B$\\
$\Chat$, $\Chatset$ & The acceptance-statistic register $\widehat{C}$ and its alphabet $\widehat{\mathcal{C}}$\\
$\Chkey$ & Key-live acceptance values\\
$\Zval$, $\Zset$ & Key alphabet $\{0,\dots,d_Z{-}1\}$ and $\mathcal{Z}\sqcup\{\perp\}$\\
$\hold$ & Declared key holder, $\mathsf{h}\in\{\mathsf A,\mathsf B\}$\\
$\CList$ & Declared constraint list $\{(\hat O_k,\bowtie_k,\Theta_k)\}$\\
$\acSet$ & Acceptance set\\
$\Vmsm$ & Measurement isometry ($\sqrt{M_x}\otimes\sqrt{F_v}$ form)\\
$\Vann$ & Announcement isometry\\
$\VI$ & Eve's copy of the announcements\\
$\VChat$ & Acceptance-statistic isometry\\
$\VR$, $\VZ$ & Raw-key isometry; raw key \ensuremath{\to} key register\\
$\Wprot$ & Full protocol isometry\\
$\Mprot$ & Single-round protocol channel $AQ\to Z\widehat{C}I$\\
$\Phimap$ & The test-statistics map $\mathcal{L}(AQ)\to\mathbb{R}^{\widehat{\mathcal{C}}}$\\
$\cZ$ & The key pinching, on the raw-key register $R$\\
$\PinS$ & Sector pinching on $Q$\\
$\cG$ & Branch (generation) maps $\mathcal{G}_{\hat c}$, sector components $\mathcal{G}_{\hat c,s}$\\
$\Feas$ & Feasible set of sector-block-diagonal marginals: $\operatorname{tr}\omega=1$ always; $\omega_A=\tau_A$ in trusted-source mode only; $\mathfrak{L}$ if declared (\mref{Def.}{def:feasible})\\
$\hup$ & Single-round MEAT rate (up variant)\\
$\hqkd$ & Single-round QKD rate \textemdash{} carries the $\uparrow$ arrow\\
$\hlow$ & Computable relaxation: lower bound on $h^{\uparrow,\mathrm{QKD}}_{\alpha}$; written \texttt{h\_low} in plain text\\
$\hlb$ & First-order lower bound on $\underline{h}_{\alpha}$: the maximum over iterates of the linearization of the objective, with lower bounds on its two linear minima\\
$\tnum$ & The instance's one declared numerical tolerance: the ceiling on the pre-snap defect an instance may repair in what it declares \textemdash{} a declared block whose measured defect (Hermiticity, positivity, completeness, column-stochasticity, a probability sum) exceeds it is a modelling error and is refused (a probability vector prepared from numerical inputs is admitted by the input preparation's rule instead, \extref{../foundations/contract.pdf}{foundations/req:declared-tolerance}); one within it is snapped onto an exactly admissible family and the snap is recorded with the defect it corrected\\
$\epa$, $\eev$ & Secrecy / correctness budgets\\
$\lEC$ & Error-correction transcript length, total bits, fixed-length register (shorter transcripts padded)\\
$\OmAT$ & Acceptance event $\operatorname{freq}_{\hat c_1^n}\in\mathcal{S}_{\mathrm{acc}}$\\
$\Etil$, $\Ehat$ & Purifying register; Eve's $n$-round register\\
$\sfD$ & Classical KL divergence, base 2; also the quantum relative entropy on operators\\
$\Dsand$ & Sandwiched R\'enyi divergence\\
$\Hsandup$ & Sandwiched R\'enyi conditional entropy, up\\
$\Hsanddown$ & Sandwiched R\'enyi conditional entropy, down\\
$\lmin$ & Minimum eigenvalue\\
$\transpose$ & Transpose, $X^{\mathsf T}$\\
\bottomrule
\end{xltabular}}

\section*{Terms}\addcontentsline{toc}{section}{Terms}
\sectionintro{The phrases this contract defines, and the terms of other contracts it uses, each with a one-line gloss and the definition that fixes it. In the text, a linked term opens its definition.}
{\small\begin{xltabular}{\linewidth}{@{}>{\raggedright\arraybackslash}p{3.6cm}>{\raggedright\arraybackslash}X>{\raggedright\arraybackslash}p{2.6cm}@{}}
\toprule
Term & Meaning & Defined in\\
\midrule
\endhead
\termref{term-acceptance-support}{acceptance support} & The outcomes of the acceptance alphabet to which some distribution of the acceptance set gives positive probability. & \mref{Definition}{def:first-order-setting}\\
\addlinespace[2pt]
\termref{term-acceptance-side-inequality}{acceptance-side inequality} & The inequality asking the acceptance-side linear form at the point used to be at most its minimum over the acceptance set. & \mref{Definition}{def:first-order-setting}\\
\addlinespace[2pt]
\termref{term-admissible-acceptance-set}{admissible acceptance set} & A closed convex subset of the probability simplex on the acceptance values. & \mref{Definition}{def:admissible-acceptance-set}\\
\addlinespace[2pt]
\exttermref{../foundations/contract.pdf}{term-admissible-attack-sequence}{admissible attack sequence} & an attack sequence whose every round lies in the attack set of a given constraint list & \extref{../foundations/contract.pdf}{foundations/def:round-model}\\
\addlinespace[2pt]
\termref{term-admissible-declaration}{admissible declaration} & A protocol declaration satisfying the eight conditions of this module: finite systems and total maps, a well-formed marginal slot, block-diagonal POVMs, announcement-determined sifting, constraints typed by the source mode, a non-empty feasible set and party computability. & \mref{Definition}{def:admissible-declaration}\\
\addlinespace[2pt]
\exttermref{../protocol-compiler/contract.pdf}{term-announcement-map}{announcement map} & A map from one party's readout alphabet to the alphabet of what that party announces. & \extref{../protocol-compiler/contract.pdf}{protocol-compiler/def:declaration}\\
\addlinespace[2pt]
\exttermref{../foundations/contract.pdf}{term-attack-sequence}{attack sequence} & one channel per round whose output is the per-round system, with Alice's system where the adversary supplies it, and Eve's new memory & \extref{../foundations/contract.pdf}{foundations/def:round-model}\\
\addlinespace[2pt]
\exttermref{../foundations/contract.pdf}{term-attack-set}{attack set} & the channels whose adjoint maps every listed operator into the stated relation with its constant & \extref{../foundations/contract.pdf}{foundations/def:round-model}\\
\addlinespace[2pt]
\termref{term-authenticated-classical-channel}{authenticated classical channel} & The classical channel that delivers each message either exactly as sent or as an authentication failure, never before it was sent, the adversary reading every message and choosing when it is delivered. & \mref{Definition}{def:authenticated-channel}\\
\addlinespace[2pt]
\termref{term-authentication-post-processing}{authentication post-processing} & The closing exchange of accept or abort decisions after which a party that saw an authentication failure, or whose partner did, holds no key. & \mref{Definition}{def:authenticated-channel}\\
\addlinespace[2pt]
\exttermref{../protocol-compiler/contract.pdf}{term-band-a}{Band A} & The direct route: a finite measurement declared on one sector, an empty source segment and one identity receiver link, given as data, with empty slots, that delivers the declared measurement. & \extref{../protocol-compiler/contract.pdf}{protocol-compiler/def:banda}\\
\addlinespace[2pt]
\exttermref{../protocol-compiler/contract.pdf}{term-basic-format}{basic format} & The format with one Alice sector, one measured system, Alice's measurement as a POVM, two messages neither reading the other, and no party-local entry. & \extref{../protocol-compiler/contract.pdf}{protocol-compiler/def:declaration}\\
\addlinespace[2pt]
\exttermref{../foundations/contract.pdf}{term-binary64-numbers}{binary64 numbers} & the finite values of the IEEE 754 binary64 format, each a dyadic rational & \extref{../foundations/contract.pdf}{foundations/def:binary64}\\
\addlinespace[2pt]
\exttermref{../foundations/contract.pdf}{term-block-list}{block list} & one square complex matrix per block of a sector map, of that block's dimension, identified with their direct sum & \extref{../foundations/contract.pdf}{foundations/def:sector-map}\\
\addlinespace[2pt]
\termref{term-branch-map}{branch map} & The pass map and the fail map of an acceptance value: the sums of the Kraus terms of its pass classes and of its fail classes. & \mref{Definition}{def:branch}\\
\addlinespace[2pt]
\exttermref{../foundations/contract.pdf}{term-certified}{certified} & the tier of a lower bound on the minimum of the objective of a declaration, formed with every arithmetic step from that declaration's data exact or rigorously enclosed, so that none of its hypotheses concerns arithmetic & \extref{../foundations/contract.pdf}{foundations/def:tiers}\\
\addlinespace[2pt]
\exttermref{../protocol-compiler/contract.pdf}{term-compilation-map}{compilation map} & The map from a protocol specification to the protocol declaration it fills, slot by slot. & \extref{../protocol-compiler/contract.pdf}{protocol-compiler/def:compile}\\
\addlinespace[2pt]
\exttermref{../foundations/contract.pdf}{term-computed-object}{computed object} & an object a module computes itself in the course of a computation, such as an iterate of an optimiser; it is not declared data & \extref{../foundations/contract.pdf}{foundations/def:declared-data}\\
\addlinespace[2pt]
\termref{term-constrained-single-round-rate}{constrained single-round rate} & The infimum, over acceptance distributions and states meeting the declared marginal and constraint list, of the KL penalty plus the key's conditional entropy. & \mref{Definition}{def:hqkd-list}\\
\addlinespace[2pt]
\termref{term-constraint-active}{constraint-active} & Said of a sector that is key-inert for every key-live acceptance value: it adds no deficit but remains a variable of the program. & \mref{Definition}{def:classes}\\
\addlinespace[2pt]
\exttermref{../foundations/contract.pdf}{term-consumer}{consumer} & a module other than a supplier that reads declared data & \extref{../foundations/contract.pdf}{foundations/def:declared-data}\\
\addlinespace[2pt]
\exttermref{../reductions/contract.pdf}{term-data-sector}{data sector} & A sector of a reduction instance's target whose blocks are the datum's matrices. & \extref{../reductions/contract.pdf}{reductions/def:reduction-instance}\\
\addlinespace[2pt]
\termref{term-declared-basis-of-a-0}{declared basis of $A_0$} & The images under the declared isometry of a fixed orthonormal basis of the Choi input register. & \mref{Definition}{def:support-frame}\\
\addlinespace[2pt]
\exttermref{../foundations/contract.pdf}{term-declared-data}{declared data} & the data an instance states and the data a producer emits as its own declaration, with the sets such data define & \extref{../foundations/contract.pdf}{foundations/def:declared-data}\\
\addlinespace[2pt]
\termref{term-declared-isometry}{declared isometry} & The isometry from the Choi input register into Alice's system whose range is the support frame. & \mref{Definition}{def:support-frame}\\
\addlinespace[2pt]
\exttermref{../reductions/contract.pdf}{term-declared-model}{declared model} & The handed model with its datum replaced by the declared datum, as the instance's model map defines it. & \extref{../reductions/contract.pdf}{reductions/def:emitnorm}\\
\addlinespace[2pt]
\exttermref{../foundations/contract.pdf}{term-declared-order}{declared order} & a total order on a finite alphabet, fixed together with the alphabet & \extref{../foundations/contract.pdf}{foundations/def:declared-order}\\
\addlinespace[2pt]
\exttermref{../foundations/contract.pdf}{term-declared-tolerance}{declared tolerance} & a binary64 number, at least zero and below one, that an instance states in its run configuration & \extref{../foundations/contract.pdf}{foundations/def:tolerance}\\
\addlinespace[2pt]
\exttermref{../foundations/contract.pdf}{term-declared-value}{declared value} & a binary64 number stated as a datum, which denotes exactly that dyadic rational and approximates nothing & \extref{../foundations/contract.pdf}{foundations/def:binary64}\\
\addlinespace[2pt]
\termref{term-diagonal-frame}{diagonal frame} & The support frame whose declared isometry is the coordinate injection onto the nonzero diagonal entries of a diagonal marginal: exact and tight. & \mref{Definition}{def:support-frame}\\
\addlinespace[2pt]
\exttermref{../reductions/contract.pdf}{term-domination-link}{domination link} & A receiver-stage typed link of kind family-dominating from a finite target family to a family it dominates at the domination weights, whose channel witness is the noise channel with added or shared flag sectors. & \extref{../reductions/contract.pdf}{reductions/domination/def:domination-link}\\
\addlinespace[2pt]
\termref{term-dropped-program}{dropped program} & The key-rate program with the deficits of a dropped set replaced by zero, over the unchanged feasible set. & \mref{Definition}{def:dropped}\\
\addlinespace[2pt]
\termref{term-dropped-set}{dropped set} & A set of pairs of a key-live acceptance value and a joint sector whose entropy deficits a run replaces by zero. & \mref{Definition}{def:dropped}\\
\addlinespace[2pt]
\termref{term-effective-outcome-set}{effective outcome set} & The outcome pairs of a pass class whose element of Alice's measurement is non-zero and whose element of the measured family is non-zero on the given sector. & \mref{Definition}{def:keyinert}\\
\addlinespace[2pt]
\termref{term-end-of-an-execution}{end of an execution} & The time at which both parties of an execution have sent and received every message the post-processing prescribes to them and fixed their outputs. & \mref{Definition}{def:execution-end}\\
\addlinespace[2pt]
\exttermref{../protocol-compiler/contract.pdf}{term-ensemble-form}{ensemble form} & The trusted-source form of Alice's measurement slot that holds positive blocks summing to the marginal, and denotes a POVM-form declaration. & \extref{../protocol-compiler/contract.pdf}{protocol-compiler/def:declaration}\\
\addlinespace[2pt]
\exttermref{../transmitter/contract.pdf}{term-ensemble-presentation}{ensemble presentation} & The certificate form that declares only the weighted virtual ensemble in sector blocks and defines the certificate from it. & \extref{../transmitter/contract.pdf}{transmitter/def:srcert}\\
\addlinespace[2pt]
\exttermref{../foundations/contract.pdf}{term-entrywise-enclosure}{entrywise enclosure} & a Gaussian-rational Hermitian centre together with symmetric non-negative rational radii that bound, entry by entry, the modulus of the difference between a matrix and that centre & \extref{../foundations/contract.pdf}{foundations/def:exact-reference}\\
\addlinespace[2pt]
\exttermref{../foundations/contract.pdf}{term-estimate}{estimate} & the tier of a key-rate number that is the value of the objective at the point an optimiser returned, claimed as no bound & \extref{../foundations/contract.pdf}{foundations/def:tiers}\\
\addlinespace[2pt]
\exttermref{../foundations/contract.pdf}{term-evaluated-off-its-snap}{evaluated off its snap} & said of a block of a snap in a run: no number from which the run computes its key length reads the snapped block, save the points at which those numbers are evaluated and the values other than the output length that the run fixes before its first round & \extref{../foundations/contract.pdf}{foundations/def:evaluated-off-snap}\\
\addlinespace[2pt]
\exttermref{../foundations/contract.pdf}{term-exact-island}{exact island} & the computations of the engine, its instances' declarations included, by which a consumer of data or the declarer about to declare them decides whether they satisfy the conditions their declaration states, a declarer computes and decides the reference of a measurement family, a consumer's replacement of an emitted scalar bound is decided to lie on the bound's weakening side, a producer constructs the operator data it emits, the acceptance side of the optimisation is constructed and certified, or a lower bound on a linear minimum over the feasible side of the optimisation is certified & \extref{../foundations/contract.pdf}{foundations/def:exact-island}\\
\addlinespace[2pt]
\exttermref{../foundations/contract.pdf}{term-exactly-admissible}{exactly admissible} & said of a declared operator family whose sector map agrees with the given one, whose blocks are exactly Hermitian and positive semidefinite, and whose sum in the declared order equals the constant times the identity on every block and rounds nowhere; the constant is one where none is named & \extref{../foundations/contract.pdf}{foundations/def:exact-admissibility}\\
\addlinespace[2pt]
\exttermref{../foundations/contract.pdf}{term-factor-enclosure}{factor enclosure} & a matrix box of columns together with an exact positive diagonal core, such that the matrix is the product of a matrix L the box contains, the core and the adjoint of L & \extref{../foundations/contract.pdf}{foundations/def:factor-enclosure}\\
\addlinespace[2pt]
\termref{term-fail-class}{fail class} & An announcement class of an acceptance value that the sifting predicate fails. & \mref{Definition}{def:branch}\\
\addlinespace[2pt]
\exttermref{../protocol-compiler/contract.pdf}{term-feasible-set}{feasible set} & The set of normalised block-list states on A and the measured system that meet the declared marginal and every constraint entry, each read exactly. & \extref{../protocol-compiler/contract.pdf}{protocol-compiler/def:feasible}\\
\addlinespace[2pt]
\termref{term-feasible-side-inequality}{feasible-side inequality} & The inequality asking the feasible-side linear form at a block list to be at most its minimum over the exact Choi feasible set. & \mref{Definition}{def:first-order-setting}\\
\addlinespace[2pt]
\exttermref{../foundations/contract.pdf}{term-first-order}{first-order} & the tier of a key-rate number that is a lower bound on the minimum of the objective, established by a claim under hypotheses that claim names & \extref{../foundations/contract.pdf}{foundations/def:tiers}\\
\addlinespace[2pt]
\termref{term-fixed-length-post-processing-model}{fixed-length post-processing model} & The classical processing after the rounds: acceptance test, one-way error correction with a padded fixed-length transcript, error verification with an announced flag, and privacy amplification, all messages public. & \mref{Definition}{def:postprocessing-model}\\
\addlinespace[2pt]
\termref{term-full-frame}{full frame} & The support frame whose declared isometry is the identity on Alice's system: always available, exact, and tight exactly when her marginal is nonsingular. & \mref{Definition}{def:support-frame}\\
\addlinespace[2pt]
\termref{term-fully-dropped}{fully dropped} & Said of a joint sector whose pairs with every key-live acceptance value all lie in the dropped set. & \mref{Definition}{def:dropped}\\
\addlinespace[2pt]
\exttermref{../protocol-compiler/contract.pdf}{term-general}{general} & Said of a constraint entry declaring a Hermitian matrix on Alice's system as its right-hand side. & \extref{../protocol-compiler/contract.pdf}{protocol-compiler/def:homog}\\
\addlinespace[2pt]
\exttermref{../foundations/contract.pdf}{term-handed-reference}{handed reference} & a reference of a snapped measurement family that its declarer passes, with its centres and radii and unaltered, to every consumer of the snap & \extref{../foundations/contract.pdf}{foundations/def:evaluated-off-snap}\\
\addlinespace[2pt]
\exttermref{../protocol-compiler/contract.pdf}{term-homogeneous}{homogeneous} & Said of a constraint entry, not declared party-local, whose right-hand side is a real scalar; in trusted-source mode it denotes that scalar times the marginal. & \extref{../protocol-compiler/contract.pdf}{protocol-compiler/def:homog}\\
\addlinespace[2pt]
\termref{term-honest-authentication-setting}{honest-authentication setting} & The setting in which every classical message reaches its recipient exactly as sent, in the order sent and after it was sent, the adversary receiving a copy, and no message fails authentication. & \mref{Definition}{def:postprocessing-model}\\
\addlinespace[2pt]
\exttermref{../protocol-compiler/contract.pdf}{term-index-map}{index map} & The emitted choice of a bijection from each alphabet and from the sector set onto an initial segment of the non-negative integers, through which every emitted list is indexed. & \extref{../protocol-compiler/contract.pdf}{protocol-compiler/def:compile}\\
\addlinespace[2pt]
\exttermref{../foundations/contract.pdf}{term-instance}{instance} & the party that fixes one protocol's device data, protocol data and run configuration & \extref{../foundations/contract.pdf}{foundations/def:declared-data}\\
\addlinespace[2pt]
\termref{term-joint-key-active-index-set}{joint key-active index set} & The joint sectors that are not key-inert with a given acceptance value. & \mref{Definition}{def:joint-sectors}\\
\addlinespace[2pt]
\exttermref{../protocol-compiler/contract.pdf}{term-joint-sector}{joint sector} & A pair of an Alice sector and a sector of the measured system; the joint sectors are ordered with the Alice sector slowest. & \extref{../protocol-compiler/contract.pdf}{protocol-compiler/def:declaration}\\
\addlinespace[2pt]
\termref{term-key-pinching}{key pinching} & The pinching of the raw-key register onto its computational basis. & \mref{Definition}{def:maps}\\
\addlinespace[2pt]
\termref{term-key-active}{key-active} & Said of a sector that fails to be key-inert for some key-live acceptance value, so that it carries entropy deficits in the objective. & \mref{Definition}{def:classes}\\
\addlinespace[2pt]
\termref{term-key-active-index-set}{key-active index set} & The set of sectors that are not key-inert with a given acceptance value. & \mref{Definition}{def:keyactive}\\
\addlinespace[2pt]
\termref{term-key-inert}{key-inert} & Said of a key-live acceptance value and a sector on which the raw-key map is constant on the effective outcome set of every pass class of that value. & \mref{Definition}{def:keyinert}\\
\addlinespace[2pt]
\termref{term-key-live-acceptance-value}{key-live acceptance value} & An acceptance value that some announcement class passing this module's sifting predicate maps to. & \mref{Definition}{def:declaration}\\
\addlinespace[2pt]
\termref{term-licensed-attack-set}{licensed attack set} & The attack set of a licensed constraint list, at the squashed system, in the declared source mode. & \mref{Definition}{def:roundmodel}\\
\addlinespace[2pt]
\exttermref{../protocol-compiler/contract.pdf}{term-licensed-for-the-declared-source-mode}{licensed for the declared source mode} & Said of a constraint list every entry of which is in the form the declared source mode admits, or is party-local. & \extref{../protocol-compiler/contract.pdf}{protocol-compiler/def:homog}\\
\addlinespace[2pt]
\exttermref{../protocol-compiler/contract.pdf}{term-lifted-operator}{lifted operator} & The operator on the whole measured system that a scoped entry's operator denotes: the identity on every other factor. & \extref{../protocol-compiler/contract.pdf}{protocol-compiler/def:declaration}\\
\addlinespace[2pt]
\exttermref{../protocol-compiler/contract.pdf}{term-marginal-slot}{marginal slot} & The slot of a declaration that holds Alice's marginal state, or the value marking untrusted-source mode. & \extref{../protocol-compiler/contract.pdf}{protocol-compiler/def:declaration}\\
\addlinespace[2pt]
\exttermref{../foundations/contract.pdf}{term-matrix-box}{matrix box} & a matrix of Gaussian-rational centres together with non-negative rational radii that bound, entry by entry, the modulus of the difference between a rectangular complex matrix and that centre & \extref{../foundations/contract.pdf}{foundations/def:factor-enclosure}\\
\addlinespace[2pt]
\termref{term-maximal-branch-support}{maximal branch support} & The span of the images, under the pass-class Kraus operators of an acceptance value and a sector, of the admissible Alice subspace tensored with that sector. & \mref{Definition}{def:branch-support}\\
\addlinespace[2pt]
\termref{term-maximal-key-length}{maximal key length} & The right-hand side of the key-length bound rounded down to an integer and clamped at zero. & \mref{Definition}{def:max-key-length}\\
\addlinespace[2pt]
\exttermref{../protocol-compiler/contract.pdf}{term-measured-factor}{measured factor} & A factor tagged Q: a system one party measures, with its own sector map, alphabet and family. & \extref{../protocol-compiler/contract.pdf}{protocol-compiler/def:declaration}\\
\addlinespace[2pt]
\exttermref{../reductions/contract.pdf}{term-measured-system}{measured system} & The quantum system the measuring parties receive in one round, and no other system; Bob's alone where he alone measures. & \extref{../reductions/contract.pdf}{reductions/def:protocol}\\
\addlinespace[2pt]
\exttermref{../reductions/contract.pdf}{term-measurement-channel}{measurement channel} & The channel that measures a POVM on a quantum input and writes the outcome into a classical register. & \extref{../reductions/contract.pdf}{reductions/def:msm}\\
\addlinespace[2pt]
\exttermref{../foundations/contract.pdf}{term-moved}{moved} & a block of a snap of an operator family read as a measurement that is not kept: the snap is not decided exactly to equal on it the data it is made from & \extref{../foundations/contract.pdf}{foundations/def:snap}\\
\addlinespace[2pt]
\termref{term-objective}{objective} & The function of an acceptance distribution and a feasible state that the program minimises: a divergence penalty plus one entropy term per key-live acceptance value. & \mref{Definition}{def:feasible}\\
\addlinespace[2pt]
\exttermref{../reductions/contract.pdf}{term-output-registers}{output registers} & The tuple of registers on which the full outputs of a protocol presentation live. & \extref{../reductions/contract.pdf}{reductions/def:opreduction}\\
\addlinespace[2pt]
\exttermref{../protocol-compiler/contract.pdf}{term-party-partition}{party partition} & The readout maps of the two parties, with a party label for Alice's system and for each tensor factor of every sector. & \extref{../protocol-compiler/contract.pdf}{protocol-compiler/def:declaration}\\
\addlinespace[2pt]
\exttermref{../protocol-compiler/contract.pdf}{term-party-local}{party-local} & Said of a scoped constraint entry, tagged so, whose condition holds against the marginal of the state on the systems other than its factor. & \extref{../protocol-compiler/contract.pdf}{protocol-compiler/def:homog}\\
\addlinespace[2pt]
\exttermref{../reductions/contract.pdf}{term-party-local-form}{party-local form} & The form of a constraint entry scoped to one factor of a product measured system: tagged party-local and read as an operator inequality against the marginal actually present. & \extref{../reductions/contract.pdf}{reductions/def:licensed}\\
\addlinespace[2pt]
\termref{term-pass-class}{pass class} & An announcement class of an acceptance value that the sifting predicate passes. & \mref{Definition}{def:branch}\\
\addlinespace[2pt]
\exttermref{../protocol-compiler/contract.pdf}{term-povm-form}{POVM form} & The form of Alice's measurement slot that holds the POVM itself. & \extref{../protocol-compiler/contract.pdf}{protocol-compiler/def:declaration}\\
\addlinespace[2pt]
\exttermref{../foundations/contract.pdf}{term-producer}{producer} & a module other than a supplier that emits numerical data, constructed from declared data it read, to another module as its own declaration & \extref{../foundations/contract.pdf}{foundations/def:declared-data}\\
\addlinespace[2pt]
\exttermref{../protocol-compiler/contract.pdf}{term-protocol-declaration}{protocol declaration} & The finite object the compiler emits and the key-rate engine consumes: systems, sector map, marginal slot, measurement families, public maps, raw-key map and constraint list. & \extref{../protocol-compiler/contract.pdf}{protocol-compiler/def:declaration}\\
\addlinespace[2pt]
\exttermref{../protocol-compiler/contract.pdf}{term-readout-map}{readout map} & A total map from the outcome pairs to one party's declared finite alphabet, saying what that party reads. & \extref{../protocol-compiler/contract.pdf}{protocol-compiler/def:declaration}\\
\addlinespace[2pt]
\exttermref{../reductions/contract.pdf}{term-receiver-segment}{receiver segment} & The list of receiver-stage typed links of a reduction chain. & \extref{../reductions/contract.pdf}{reductions/def:chain}\\
\addlinespace[2pt]
\exttermref{../reductions/contract.pdf}{term-reduction-chain}{reduction chain} & A pair of a source segment and a receiver segment of typed links. & \extref{../reductions/contract.pdf}{reductions/def:chain}\\
\addlinespace[2pt]
\exttermref{../reductions/contract.pdf}{term-reduction-instance}{reduction instance} & A handed model with a finite datum and hypotheses about the rest, a target built from the datum and scalar values, a squashing relation with a transported set, and certified entries, all established for every admissible datum. & \extref{../reductions/contract.pdf}{reductions/def:reduction-instance}\\
\addlinespace[2pt]
\exttermref{../foundations/contract.pdf}{term-reference}{reference} & a family of Hermitian positive semidefinite block lists complete at the constant on every block, read together with an entrywise enclosure of every block of every element & \extref{../foundations/contract.pdf}{foundations/def:exact-reference}\\
\addlinespace[2pt]
\termref{term-relaxed-feasible-set}{relaxed feasible set} & The feasible set read at a window: positivity and the sector-block structure kept exact, the marginal or coupling condition, the declared constraint rows and the normalisation each relaxed at the window. & \mref{Definition}{def:relaxed-feasible-set}\\
\addlinespace[2pt]
\termref{term-removable}{removable} & Said of a sector for which the hypotheses of the sector-removal theorem hold, as certified by a witness meeting the removal checklist. & \mref{Definition}{def:classes}\\
\addlinespace[2pt]
\exttermref{../foundations/contract.pdf}{term-rigorously-enclosed}{rigorously enclosed} & said of a computation each of whose operations returns a set of numbers, such as an interval or a ball, containing the exact result of the operation at every value its inputs' sets contain & \extref{../foundations/contract.pdf}{foundations/def:tiers}\\
\addlinespace[2pt]
\termref{term-round-channel}{round channel} & The channel of round j: the attack followed by the single-round protocol map, with the announcement placed in the adversary's register. & \mref{Definition}{def:roundmodel}\\
\addlinespace[2pt]
\exttermref{../foundations/contract.pdf}{term-round-model}{round model} & n rounds of a per-round system S, with the attack sequences and attack sets below & \extref{../foundations/contract.pdf}{foundations/def:round-model}\\
\addlinespace[2pt]
\exttermref{../foundations/contract.pdf}{term-rounds-nowhere}{rounds nowhere} & said of an accumulation every one of whose binary64 additions is exact & \extref{../foundations/contract.pdf}{foundations/def:declared-order}\\
\addlinespace[2pt]
\termref{term-scalar-list}{scalar list} & A constraint list of untrusted-source mode whose right-hand sides are all real numbers, each entry a comparison of real numbers. & \mref{Definition}{def:homog}\\
\addlinespace[2pt]
\exttermref{../protocol-compiler/contract.pdf}{term-scope}{scope} & The measured factor that an entry of the constraint list is declared on. & \extref{../protocol-compiler/contract.pdf}{protocol-compiler/def:declaration}\\
\addlinespace[2pt]
\exttermref{../foundations/contract.pdf}{term-sector-map}{sector map} & a block kind together with a non-empty, ordered list of block identifiers, each with a dimension & \extref{../foundations/contract.pdf}{foundations/def:sector-map}\\
\addlinespace[2pt]
\exttermref{../reductions/contract.pdf}{term-sector-set}{sector set} & The finite ordered list of a reduction instance's target sectors, data and scalar, each with a declared dimension. & \extref{../reductions/contract.pdf}{reductions/def:reduction-instance}\\
\addlinespace[2pt]
\exttermref{../foundations/contract.pdf}{term-sequential-sum}{sequential sum} & the exact sum of a finite sequence, accumulated left to right from zero & \extref{../foundations/contract.pdf}{foundations/def:declared-order}\\
\addlinespace[2pt]
\termref{term-single-round-protocol-map}{single-round protocol map} & The map from Alice's and the measured system to the key, acceptance and announcement registers that the declaration generates for one round. & \mref{Definition}{def:maps}\\
\addlinespace[2pt]
\exttermref{../foundations/contract.pdf}{term-snap}{snap} & at a lattice exponent k, data of the same shape as the given data, over the same alphabet in the same order, with every entry a lattice value at k, that satisfy exactly the conditions the declaration of the given data states & \extref{../foundations/contract.pdf}{foundations/def:snap}\\
\addlinespace[2pt]
\exttermref{../transmitter/contract.pdf}{term-source-model}{source model} & A submodule's named definition of the signal ensemble as mathematical objects, possibly infinite dimensional, rather than as matrices. & \extref{../transmitter/contract.pdf}{transmitter/def:source-spec}\\
\addlinespace[2pt]
\exttermref{../transmitter/contract.pdf}{term-source-segment}{source segment} & The ordered, possibly empty list of source links leading from the virtual ensemble to the specified one. & \extref{../transmitter/contract.pdf}{transmitter/def:source-segment}\\
\addlinespace[2pt]
\exttermref{../transmitter/contract.pdf}{term-source-specification}{source specification} & The declared description of Alice's side of a protocol: alphabet, retained system with sectors, measurement, signal ensemble and source mode. & \extref{../transmitter/contract.pdf}{transmitter/def:source-spec}\\
\addlinespace[2pt]
\exttermref{../transmitter/contract.pdf}{term-source-replacement-certificate}{source-replacement certificate} & An operator on Alice's retained system and the virtual transmitted system whose marginal on her side is her declared marginal and which, measured on her side by her POVM, leaves the weighted virtual members, with retention attested. & \extref{../transmitter/contract.pdf}{transmitter/def:srcert}\\
\addlinespace[2pt]
\exttermref{../reductions/contract.pdf}{term-squashing-channel}{squashing channel} & A channel from the physical system to the squashed system whose adjoint realises a squashing relation. & \extref{../reductions/contract.pdf}{reductions/def:squashing}\\
\addlinespace[2pt]
\termref{term-state-blind}{state-blind} & Said of a sector on which every measurement operator of the measured system, and every constraint operator when constraints are declared, is a scalar multiple of the identity. & \mref{Definition}{def:classes}\\
\addlinespace[2pt]
\termref{term-statistics-map}{statistics map} & The map sending a state on Alice's and the measured system to the probabilities of the acceptance values. & \mref{Definition}{def:maps}\\
\addlinespace[2pt]
\exttermref{../foundations/contract.pdf}{term-sum-in-the-declared-order}{sum in the declared order} & the sequential sum of a family indexed by an alphabet, taken in the alphabet's declared order & \extref{../foundations/contract.pdf}{foundations/def:declared-order}\\
\addlinespace[2pt]
\exttermref{../foundations/contract.pdf}{term-supplier}{supplier} & a module whose own claims decide their conditions on the data it reads or emits in floating point, at slacks an instance declares to it, rather than exactly & \extref{../foundations/contract.pdf}{foundations/def:declared-data}\\
\addlinespace[2pt]
\termref{term-support-frame}{support frame} & The range of an isometry from the Choi input register into Alice's system that contains the support of her declared marginal. & \mref{Definition}{def:support-frame}\\
\addlinespace[2pt]
\termref{term-system-metadata}{system metadata} & The party label of Alice's system and, per sector, the declared tensor factorisation with a party label per factor. & \mref{Definition}{def:declaration}\\
\addlinespace[2pt]
\exttermref{../reductions/contract.pdf}{term-target-family}{target family} & The target POVM of a reduction instance at a datum, assigned one sector block at a time from the datum and the scalar values. & \extref{../reductions/contract.pdf}{reductions/def:reduction-instance}\\
\addlinespace[2pt]
\termref{term-tau-a-homogeneous-list}{$\tau_A$-homogeneous list} & A constraint list of trusted-source mode whose entries are all declared homogeneous, each right-hand side constructed as its scalar times Alice's declared marginal. & \mref{Definition}{def:homog}\\
\addlinespace[2pt]
\exttermref{../foundations/contract.pdf}{term-tier}{tier} & one of three labels, estimate, first-order or certified, that a reported key-rate number carries & \extref{../foundations/contract.pdf}{foundations/def:tiers}\\
\addlinespace[2pt]
\termref{term-tight}{tight} & Said of a support frame equal to the support of Alice's marginal, equivalently of dimension equal to its rank. & \mref{Definition}{def:support-frame}\\
\addlinespace[2pt]
\exttermref{../protocol-compiler/contract.pdf}{term-transcript-protocol}{transcript protocol} & The announcement slot of the general format: an ordered list of messages, each a function of its sender's readout and of the earlier messages. & \extref{../protocol-compiler/contract.pdf}{protocol-compiler/def:declaration}\\
\addlinespace[2pt]
\exttermref{../transmitter/contract.pdf}{term-transmitted-system}{transmitted system} & The system Alice sends into the channel, carrying the specified signal ensemble. & \extref{../transmitter/contract.pdf}{transmitter/def:source-spec}\\
\addlinespace[2pt]
\exttermref{../transmitter/contract.pdf}{term-trusted-source-mode}{trusted-source mode} & The source mode in which Alice's retained system never reaches the adversary, with a declared marginal and, above dimension one, a certificate. & \extref{../transmitter/contract.pdf}{transmitter/def:source-spec}\\
\addlinespace[2pt]
\exttermref{../transmitter/contract.pdf}{term-untrusted-source-mode}{untrusted-source mode} & The source mode with no certificate and no marginal, in which the adversary may have prepared Alice's retained system. & \extref{../transmitter/contract.pdf}{transmitter/def:source-spec}\\
\addlinespace[2pt]
\exttermref{../transmitter/contract.pdf}{term-virtual-ensemble}{virtual ensemble} & The finite-dimensional ensemble the source segment starts from and the certificate is written against. & \extref{../transmitter/contract.pdf}{transmitter/def:source-segment}\\
\bottomrule
\end{xltabular}}

\clearpage
\part{Mathematics}

\section{The declaration and the protocol channel}\label{group-meat-core/def:declaration}
\sectionintro{This section fixes the data every later result is computed from and the channel those data generate. The protocol declaration is the compiler's, read here with the conditions and names this module's claims rest on (\hyperlink{claim-meat-core/def:declaration}{Definition~\ref*{def:declaration}}); an ensemble-form declaration denotes an ordinary one (\hyperlink{claim-meat-core/lem:ensemble-form}{Lemma~\ref*{lem:ensemble-form}}). From the declared measurements, announcements and key map the module generates a single-round protocol channel whose acceptance marginal is the \termref{term-statistics-map}{statistics map} (\hyperlink{claim-meat-core/thm:protocol-channel}{Theorem~\ref*{thm:protocol-channel}}), and \termref{term-branch-map}{branch maps} that split the announcement classes of each acceptance value into pass and \termref{term-fail-class}{fail classes} (\hyperlink{claim-meat-core/def:branch}{Definition~\ref*{def:branch}}, \hyperlink{claim-meat-core/lem:branch-maps}{Lemma~\ref*{lem:branch-maps}}). The branch maps carry the orthogonality the entropy calculation uses, and a component of a register other than the raw-key register that is constant across a branch map's Kraus terms can be removed from it without changing any quantity of this module (\hyperlink{claim-meat-core/rem:inert}{Remark~\ref*{rem:inert}}).}
\sectionimports{\extref{../foundations/contract.pdf}{foundations/def:declared-order}, \extref{../protocol-compiler/contract.pdf}{protocol-compiler/def:declaration}}

% claim meat-core/def:declaration
\claimsummary{meat-core/def:declaration}{The declaration is the compiler's (\xref{protocol-compiler/def:declaration}), in either format, admissible (\hyperlink{claim-meat-core/def:admissible-declaration}{Definition~\ref*{def:admissible-declaration}}), sifting by the announcements and a raw key its holder can compute. This module adds the key-live acceptance values, its reading of the party partition, the summation order and exact admissibility.}
\claimbookmark{meat-core/def:declaration}{Definition}{def:declaration}{The protocol declaration as the MEAT core reads it}
\claimleadin{%
The operator, state and channel conventions of this contract, among them
that every Hilbert space is finite-dimensional, are those of
\xref{foundations/def:pinching}, and its divergences and conditional
entropies are those of \xref{foundations/def:kl} and
\xref{foundations/def:renyi}. In the tuple below, which is the compiler's,
$A$ is Alice's system and $\tau_A$ the content of its \exttermref{../protocol-compiler/contract.pdf}{term-marginal-slot}{marginal slot};
$\Qsys$ is the system the family $\{F_v\}$ measures, with sectors
$\Hs{\Qsys}=\bigoplus_{s\in\SecSet}\Hs{\Qsys_s}$; $\{M_x\}$ is the
measurement on $A$; $\Party$ is the \exttermref{../protocol-compiler/contract.pdf}{term-party-partition}{party partition}, with
\exttermref{../protocol-compiler/contract.pdf}{term-readout-map}{readout maps} $\rd{\mathsf A},\rd{\mathsf B}$ onto the
alphabets $\Oset{\mathsf A},\Oset{\mathsf B}$; the
\exttermref{../protocol-compiler/contract.pdf}{term-announcement-map}{announcement maps} $a_A,a_B$, or in the general
format the messages of (DX5), give each outcome pair $(x,v)$ its joint
announcement $a(x,v)$, an announcement class; $\phi$ maps
announcement classes to acceptance values; $r$ is the raw-key map, $\hold$
the key holder and $\CList$ the constraint list. The definition adds two
conditions on $r$ and names the acceptance values that can carry key, then
says what this module reads of $\Party$, reads the summation order off the
declaration, states exact admissibility, and gives the general format with its reading of
exact admissibility. The
compiler's conditions on the declaration it emits
(\xref{protocol-compiler/def:target}) are the counterpart of
Definition~\ref{def:admissible-declaration}, and the two lists are checked against each other
at the interface between the two modules.
Requirement~\ref{req:exact-admission} says what becomes of a declaration that
is not \exttermref{../foundations/contract.pdf}{term-exactly-admissible}{exactly admissible}, and that restricting the states of the
program to lists of blocks along Alice's sectors is sound is the subject of
\xref{protocol-compiler/rem:block-restriction}, which Theorem~\ref{thm:reduction}
discharges. The two source modes that the two values of the marginal slot
name are compared on the program in Remark~\ref{rem:modes-da1}.
}
\begin{definition}[The protocol declaration as the MEAT core reads it]\label{def:declaration}
In this module a \exttermref{../protocol-compiler/contract.pdf}{term-protocol-declaration}{protocol declaration} is a tuple
\begin{equation*}
  \Decl=\Big(A,\ \tau_A;\ \Qsys,\ \SecSet;\ \{M_x\}_{x\in\Xset};\
  \{F_v\}_{v\in\Vset};\ \Party;\ a_A,\ a_B;\ \phi;\ r,\ \hold;\
  \CList\Big)
\end{equation*}
of \xref{protocol-compiler/def:declaration}, in its \exttermref{../protocol-compiler/contract.pdf}{term-basic-format}{basic format} or in
its general format, (DX1)--(DX6) there, whose components satisfy the
well-formedness conditions of Definition~\ref{def:admissible-declaration} and the two
conditions on its raw-key map below. The two values of its
\exttermref{../protocol-compiler/contract.pdf}{term-marginal-slot}{marginal slot}, $\tau_A\in\Seq{A}$ and $\tau_A=\nomarg$, are the two
source modes, \exttermref{../transmitter/contract.pdf}{term-trusted-source-mode}{trusted-source mode} and
\exttermref{../transmitter/contract.pdf}{term-untrusted-source-mode}{untrusted-source mode}. Write $\Kann$ for the announcement-class
alphabet of that definition, $\CAset\times\CBset$ in the basic format and
the product of the message alphabets of (DX5) in the general one,
$a(x,v)\in\Kann$ for the joint announcement it assigns to an outcome pair,
and $\kappa$ for a class. \emph{Sifting is announcement-determined}: there is
$\varsigma:\Kann\to\{\mathrm{pass},\mathrm{fail}\}$ with
$r(x,v)=\perp\iff\varsigma(a(x,v))=\mathrm{fail}$. \emph{The raw key is
party-computable}: for the declared key holder
$\hold\in\{\mathsf A,\mathsf B\}$ there is
$r_{\hold}:\Oset{\hold}\times\Kann\to\Zset$ with
\begin{equation}\label{eq:partycomp}
  r(x,v)=r_{\hold}\big(\rd{\hold}(x,v),\,a(x,v)\big)
  \qquad\text{for all }(x,v)\in\Xset\times\Vset
\end{equation}
(Definition~\ref{def:admissible-declaration}\,(viii)). We write
\begin{equation}\label{eq:chkey}
  \Chkey\coloneqq\big\{\hat c\in\Chatset:\ \exists\,\kappa\in\phi^{-1}(\hat c)
  \text{ with }\varsigma(\kappa)=\mathrm{pass}\big\}
\end{equation}
for the set of
\termdef{term-key-live-acceptance-value}{key-live acceptance value}s,
$\varsigma$ being the predicate of the sifting condition above. That
condition fixes $\varsigma$ on every class some outcome pair realises, and
on a class $\kappa\notin a(\Xset\times\Vset)$, which no outcome pair
realises, we take $\varsigma(\kappa)=\mathrm{fail}$; so $\Chkey$ is the set
of acceptance values to which some realised class passing sifting maps,
and it is fixed by the declaration.

\emph{What this module reads of $\Party$.} The \exttermref{../protocol-compiler/contract.pdf}{term-party-partition}{party partition}
$\Party$ of that definition consists of (a) its two
\exttermref{../protocol-compiler/contract.pdf}{term-readout-map}{readout maps} $\rd{\mathsf A}$ and $\rd{\mathsf B}$,
naming what each party physically observes in a round, and (b)
\termdef{term-system-metadata}{system metadata}:
the party label for $A$ and, for each $s\in\SecSet$, the declared tensor
factorization $\Hs{\Qsys_s}=\bigotimes_{i}\Hs{\Qsys_{s,i}}$ with a party
label per factor (a comma-indexed \emph{tensor} factor $\Qsys_{s,i}$, never
the parenthesized superscript $X^{(s)}$, which this module reserves for
\emph{sector blocks} of an operator). The \emph{default partition} is that of
\xref{protocol-compiler/def:target}: $\rd{\mathsf A}(x,v)=x$,
$\rd{\mathsf B}(x,v)=v$ with $A$ labelled $\mathsf A$ and every $\Qsys_s$
labelled $\mathsf B$. Part~(a) is consumed by the derivations, through
Definition~\ref{def:admissible-declaration}\,(viii). Part~(b) is consumed by no
derivation; it carries an obligation of the compiler: the declaration's POVMs
must be realizable by each party measuring the factors it is labelled with,
so that $\rd{\mathsf A}$ and $\rd{\mathsf B}$ are what the parties can in fact
read. Part~(b) is read by the premises that place the parties'
laboratories: Assumption~\ref{ass:lab} takes all its labels as naming the
parties' laboratories, and A-SRC (\xref{transmitter/ass:src}) reads the label
of $A$; so a run attests that obligation there, and the claims that assume
those premises rest on it through them. This module never verifies the
obligation.

\emph{The declared summation order.} The declaration fixes, as part of the
declared alphabet $\Vset$ of outcomes, a \exttermref{../foundations/contract.pdf}{term-declared-order}{declared order}
$v_1\prec v_2\prec\dots\prec v_{|\Vset|}$ on it. For each sector $s$ the
completeness sum is the \exttermref{../foundations/contract.pdf}{term-sum-in-the-declared-order}{sum in the declared order} of the blocks
$F^{(s)}_v$, the \exttermref{../foundations/contract.pdf}{term-sequential-sum}{sequential sum}
\begin{equation}\label{eq:declorder}
  \Sigma^{(s)}_0\coloneqq0,\qquad
  \Sigma^{(s)}_i\coloneqq\Sigma^{(s)}_{i-1}+F^{(s)}_{v_i}
  \quad(1\le i\le|\Vset|),\qquad
  \Sigma^{(s)}\coloneqq\Sigma^{(s)}_{|\Vset|},
\end{equation}
formed left to right, entrywise, with no re-ordering and no
re-association. In exact arithmetic \eqref{eq:declorder} is the unordered
sum and the order is immaterial; in binary64 it is not
(\xref{foundations/def:declared-order}), so the order is part of the datum and
this module reads it off the declared alphabet rather than choosing one.

\emph{Exact admissibility.} A declaration reaches this module with its
operator data \exttermref{../foundations/contract.pdf}{term-exactly-admissible}{exactly admissible}: every $F_v$ exactly Hermitian and
exactly positive semidefinite per sector block, every $\hat O_k$ exactly
Hermitian per sector block (positivity is asked of no constraint operator
anywhere in this module), every $M_x$ and the declared $\tau_A$ exactly
Hermitian and exactly positive semidefinite, and the completeness sum
\eqref{eq:declorder} of each sector exactly the identity with every addition
\exttermref{../foundations/contract.pdf}{term-rounds-nowhere}{rounding nowhere}. These are the conditions the compiler decides on the
numbers it emits (\xref{protocol-compiler/def:exact}).

\emph{The general format.} The general format is that of
\xref{protocol-compiler/def:declaration}, (DX1)--(DX6) there, and the basic
format is its case of one Alice sector, one \exttermref{../protocol-compiler/contract.pdf}{term-measured-factor}{measured factor},
\exttermref{../protocol-compiler/contract.pdf}{term-povm-form}{POVM form}, two messages neither of which reads the other, and no
party-local entry. This contract refers to five of that definition's
displays under its own labels, each read as it is defined there: the
product \exttermref{../foundations/contract.pdf}{term-sector-map}{sector map} of (DX1), over the ordered set $\mathbb{F}$ of
measured factors, factor $f$ having the sectors $s_f\in\SecSet^{f}$ with
projectors $\Pi^{f}_{s_f}$ and block dimensions $d_{f,s_f}$,
\begin{equation}\label{eq:productmap}
\begin{gathered}
  \Qsys=\bigotimes_{f\in\mathbb{F}}\Qsys_f,\qquad
  \SecSet=\prod_{f\in\mathbb{F}}\SecSet^{f},\qquad
  \Hs{\Qsys_s}=\bigotimes_{f}\Hs{\Qsys_{f,s_f}},\\
  \Pi_s=\bigotimes_f\Pi^{f}_{s_f},\qquad d_s=\prod_fd_{f,s_f};
\end{gathered}
\end{equation}
the product family of (DX3),
\begin{equation}\label{eq:productfamily}
  F_v\coloneqq\bigotimes_{f\in\mathbb{F}}F^{f}_{v_f},\qquad
  F^{(s)}_v=\bigotimes_{f\in\mathbb{F}}F^{f,(s_f)}_{v_f}
  \qquad\big(v=(v_f)_f,\ s=(s_f)_f\big);
\end{equation}
the two conditions on the declared $\{N_x\}$ of the \exttermref{../protocol-compiler/contract.pdf}{term-ensemble-form}{ensemble form} of
(DX4), whose blocks lie along the index set $\mathbb{T}$ of Alice's sectors,
$\Hs{A}=\bigoplus_{t\in\mathbb{T}}\Hs{A_t}$, and the family of the POVM-form
declaration an ensemble-form declaration denotes, with $P_t$ the projector
onto the support of $\tau^{(t)}_A$ and the inverse square root taken on that
support and zero off it,
\begin{equation}\label{eq:ensembleform}
  \sum_{x\in\Xset}N^{(t)}_x=\tau^{(t)}_A
  \quad\text{and}\quad\tau^{(t)}_A\ne0
  \qquad\text{for every }t\in\mathbb{T},
\end{equation}
\begin{equation}\label{eq:ensemblepovm}
  M^{(t)}_x\coloneqq\big(\tau^{(t)}_A\big)^{-1/2}N^{(t)}_x
  \big(\tau^{(t)}_A\big)^{-1/2}
  +\frac{\tr N^{(t)}_x}{\tr\tau^{(t)}_A}\,\big(\idon{A_t}-P_t\big);
\end{equation}
and the joint announcement function a \exttermref{../protocol-compiler/contract.pdf}{term-transcript-protocol}{transcript protocol} of (DX5)
derives by composition from its $L$ messages, message $m$ having the sender
$\pi_m$, whose readout map is $\rd{\pi_m}$, the alphabet $\mathcal{C}_m$ and
the table $a_m$,
\begin{equation}\label{eq:transcript}
\begin{gathered}
  \kappa_m(x,v)\coloneqq a_m\big(\rd{\pi_m}(x,v),\,\kappa_1(x,v),\dots,
  \kappa_{m-1}(x,v)\big),\\
  a(x,v)\coloneqq\big(\kappa_m(x,v)\big)_{m=1}^{L}\in
  \Kann\coloneqq\prod_{m=1}^{L}\mathcal{C}_m .
\end{gathered}
\end{equation}

\emph{Exact admissibility in the general format.} The exact-admissibility
paragraph above holds with these readings: the conditions on $\{F_v\}$ are
decided per measured factor, on the family $\{F^{f}_{v_f}\}$, its blocks and
its declared order, and on no product; the conditions on $\tau_A$, on every
$M_x$ and on the right-hand side $\Theta_k$ of a \exttermref{../protocol-compiler/contract.pdf}{term-general}{general} entry are
decided block by block along $\mathbb{T}$, the trace of $\tau_A$
accumulated over the blocks in their declared order and rounding nowhere;
and in ensemble form the conditions on $\{M_x\}$ are replaced by those of
(DX4) on $\{N_x\}$, the denoted $\{M_x\}$ being a POVM by
Lemma~\ref{lem:ensemble-form} and decided by nothing.
\end{definition}
\claimmeta{\metalabel{Uses}\extref{../foundations/contract.pdf}{foundations/def:declared-order}, \extref{../protocol-compiler/contract.pdf}{protocol-compiler/def:declaration}
\metasep \metalabel{Used by}\mref{Def.}{def:admissible-declaration}, \mref{Lem.}{lem:ensemble-form}, \mref{Lem.}{lem:branch-maps}, \mref{Def.}{def:homog}, \mref{Def.}{def:roundmodel}, \mref{Def.}{def:joint-sectors}, \mref{Lem.}{lem:joint-sectors}, \mref{Def.}{def:feasible}, \mref{Thm.}{thm:reduction}, \mref{Prop.}{prop:a-neutral}, \mref{Def.}{def:support-frame}, \mref{Lem.}{lem:choi-factor}, \mref{Thm.}{thm:choi}, \mref{Lem.}{lem:rank-one-sector}, \mref{Def.}{def:dropped}, \mref{Lem.}{lem:objective-homogeneous}, \mref{Def.}{def:postprocessing-model}, \mref{Rem.}{rem:cgate}, \mref{Prop.}{prop:optol-nobound}, \mref{Rem.}{rem:representability-ceiling}, \mref{Rem.}{rem:declared-measurement}, \mref{Prop.}{prop:reference-companion}, \mref{Ass.}{ass:keyreg}, \mref{Ass.}{ass:lab}, \mref{Req.}{req:exact-admission}, \mref{Req.}{req:companion-admission}, \mref{Thm.}{thm:firstorder}, \mref{Lem.}{lem:feasible-side-dual}, \mref{Lem.}{lem:certified-bracket}, \mref{Lem.}{lem:certified-bracket-enclosed}
\metasep newly derived
\metasep implemented
\metasep \metalabel{Code}\coderef{https://github.com/fraunhofersg/certiqs-key/blob/v0.1.0/src/certiqs_key/meat_core/declaration/declaration.py}{certiqs_key.meat_core.declaration.declaration:ProtocolDeclaration}}

% claim meat-core/def:admissible-declaration
\claimsummary{meat-core/def:admissible-declaration}{An admissible declaration is finite with total maps, a well-formed marginal slot, block-diagonal POVMs, announcement-determined sifting, constraints typed by the source mode and a non-empty feasible set. Each party announces from its own readout and the key holder computes the raw key from its own.}
\claimbookmark{meat-core/def:admissible-declaration}{Definition}{def:admissible-declaration}{Admissible declaration}
\claimleadin{%
Condition~(viii) is stated for a declared \exttermref{../protocol-compiler/contract.pdf}{term-party-partition}{party partition} and not
only for the default partition, because a declaration that
places both parties' data in $v$, as the $d_A=1$ embedding of
Remark~\ref{rem:modes-da1}\,(b) does, satisfies the default-partition
reading of~(viii) vacuously, with $\hold=\mathsf B$ and $r$ any function of
$v$ at all, while guaranteeing nothing about what the key holder can in
fact compute; (viii) is what makes the error-correction and verification
steps of Theorem~\ref{thm:keylength} physically executable. How these
conditions compare with the compiler's target conditions is
\xref{protocol-compiler/rem:engine-admission}. Conditions (i)--(viii) are
decided on the \exttermref{../foundations/contract.pdf}{term-declared-data}{declared data} by the compiler's validator
(\xref{protocol-compiler/def:validator}, checks V1--V8), condition~(vii)
one-sidedly (\xref{protocol-compiler/thm:feasibility}). This module's
admission decides every condition but~(vii) again, exactly
(Requirement~\ref{req:declaration-admission}), and it reports a key length
only on a declaration the compiler emitted
(Requirement~\ref{req:compiled-declaration-only}), so that every condition
is decided wherever a key length is reported.
}
\begin{definition}[Admissible declaration]\label{def:admissible-declaration}
The tuple $\Decl$ of Definition~\ref{def:declaration} is an
\termdef{term-admissible-declaration}{admissible declaration}
when it satisfies: (i) all systems finite-dimensional, all
alphabets finite, all maps total; (ii) the \exttermref{../protocol-compiler/contract.pdf}{term-marginal-slot}{marginal slot} holds
$\tau_A\in\Seq{A}$ (\exttermref{../transmitter/contract.pdf}{term-trusted-source-mode}{trusted-source mode}) or the value $\nomarg$
(\exttermref{../transmitter/contract.pdf}{term-untrusted-source-mode}{untrusted-source mode}); in particular $\tr\tau_A=1$ whenever a
marginal is declared, so $\tau_A=0$ is inadmissible;
(iii) $\{M_x\}$ and $\{F_v\}$ are POVMs and every $F_v$ and every
$\hat O_k$ is sector-block-diagonal; (iv) $r$ maps into
$\Zval\sqcup\{\perp\}$; (v) sifting is announcement-determined via
$\varsigma$; (vi) $\bowtie_k\in\{\preceq,=,\succeq\}$ and each entry
carries one of the two declared forms of Definition~\ref{def:homog}, so
that $\Theta_k$ is Hermitian on $A$, in trusted-source mode, respectively
$\bowtie_k\in\{\le,=,\ge\}$ and $\Theta_k=c_k\in\mathbb R$ in
untrusted-source mode; (vii) $\Feas\neq\varnothing$; (viii)
\emph{party computability:} the declaration names a party partition
$\Party$ (Definition~\ref{def:declaration}) and
\begin{itemize}[leftmargin=1.6em]
\item each party announces a function of its own readout: $a_A$ is
defined on $\Oset{\mathsf A}$ and $a_B$ on $\Oset{\mathsf B}$, with
$a(x,v)=(a_A(\rd{\mathsf A}(x,v)),a_B(\rd{\mathsf B}(x,v)))$;
\item the declared key holder $\hold\in\{\mathsf A,\mathsf B\}$
computes the raw key from \emph{its own readout} and the announcements,
$r(x,v)=r_{\hold}(\rd{\hold}(x,v),a(x,v))$ with $r_{\hold}$ a declared
total function, Eq.~\eqref{eq:partycomp}.
\end{itemize}
Under the default partition ($\rd{\mathsf A}(x,v)=x$,
$\rd{\mathsf B}(x,v)=v$) these conditions read:
$a_A:\Xset\to\CAset$, $a_B:\Vset\to\CBset$, and
$r(x,v)=r_{\mathsf A}(x,a(x,v))$ resp.\ $r_{\mathsf B}(v,a(x,v))$.

\emph{In the general format} of Definition~\ref{def:declaration}\,(DX1)--(DX6)
the same eight conditions read as follows, and at one Alice sector, one
\exttermref{../protocol-compiler/contract.pdf}{term-measured-factor}{measured factor}, \exttermref{../protocol-compiler/contract.pdf}{term-povm-form}{POVM form} and two messages they are the
conditions above. (iii): every $M_x$, the declared $\tau_A$ and the
right-hand side of every \exttermref{../protocol-compiler/contract.pdf}{term-general}{general} entry are block diagonal along
Alice's sectors $\mathbb{T}$, \emph{by their type as
\exttermref{../foundations/contract.pdf}{term-block-list}{block lists} against a declared factor tagged $A$}, and
$\{M_x\}$ is a POVM; in \exttermref{../protocol-compiler/contract.pdf}{term-ensemble-form}{ensemble form} the declared $\{N_x\}$ satisfy
\eqref{eq:ensembleform} and $\{M_x\}$ is the family \eqref{eq:ensemblepovm},
a POVM block diagonal along $\mathbb{T}$ by
Lemma~\ref{lem:ensemble-form}\,(i); $\{F_v\}$ is the product family
\eqref{eq:productfamily} of per-factor POVMs, each block diagonal for its
factor's \exttermref{../foundations/contract.pdf}{term-sector-map}{sector map}, and every $\hat O_k$ is block diagonal for the
sector map of its \exttermref{../protocol-compiler/contract.pdf}{term-scope}{scope}. (vi): an entry may in addition carry the
\exttermref{../protocol-compiler/contract.pdf}{term-party-local}{party-local} form of Definition~\ref{def:homog}, in either mode, with
$\bowtie_k\in\{\preceq,=,\succeq\}$ and $c_k\in\mathbb{R}$. (viii): each
message of the transcript is a function of its sender's own readout and of
the earlier messages, $a$ is the composition \eqref{eq:transcript}, and
\eqref{eq:partycomp} holds with that $a$.
\end{definition}
\claimmeta{\metalabel{Uses}\mref{Def.}{def:declaration}, \mref{Def.}{def:homog}
\metasep \metalabel{Used by}\mref{Lem.}{lem:ensemble-form}, \mref{Thm.}{thm:protocol-channel}, \mref{Lem.}{lem:branch-maps}, \mref{Lem.}{lem:purifier-indep}, \mref{Thm.}{thm:accumulated-entropy}, \mref{Ex.}{ex:fss}, \mref{Rem.}{rem:homog}, \mref{Lem.}{lem:cone-concavity}, \mref{Lem.}{lem:joint-sectors}, \mref{Prop.}{prop:trivial-alice-system}, \mref{Rem.}{rem:modes-da1}, \mref{Thm.}{thm:reduction}, \mref{Prop.}{prop:a-neutral}, \mref{Cor.}{cor:reduction-constrained}, \mref{Thm.}{thm:convexity}, \mref{Prop.}{prop:retention}, \mref{Prop.}{prop:removal-compose}, \mref{Lem.}{lem:admissible-alice-subspace}, \mref{Lem.}{lem:choi-marginal}, \mref{Lem.}{lem:choi-factor}, \mref{Lem.}{lem:estimate}, \mref{Lem.}{lem:branch-support}, \mref{Lem.}{lem:b-derivative}, \mref{Thm.}{thm:oracle}, \mref{Lem.}{lem:rank-one-sector}, \mref{Prop.}{prop:drop-bound}, \mref{Cor.}{cor:drop-program}, \mref{Prop.}{prop:minimal-divergence}, \mref{Lem.}{lem:objective-homogeneous}, \mref{Thm.}{thm:keylength}, \mref{Prop.}{prop:model-pass-bound}, \mref{Rem.}{rem:cgate}, \mref{Prop.}{prop:optol-nobound}, \mref{Prop.}{prop:reference-companion}, \mref{Req.}{req:declaration-admission}, \mref{Req.}{req:companion-admission}, \mref{Req.}{req:diagonal-frame-test}, \mref{Req.}{req:compiled-declaration-only}, \mref{Prop.}{prop:relaxed-set-properties}, \mref{Rem.}{rem:feas-tol}, \mref{Thm.}{thm:firstorder}, \mref{Cor.}{cor:fo-dropped}, \mref{Cor.}{cor:fo-reference}, \mref{Lem.}{lem:certified-bracket}, \mref{Lem.}{lem:certified-bracket-enclosed}, \mref{Cor.}{cor:certified-keylength}, \mref{Cor.}{cor:certified-keylength-reference}
\metasep newly derived
\metasep accepted}

% claim meat-core/lem:ensemble-form
\claimsummary{meat-core/lem:ensemble-form}{Positive operators summing to Alice's marginal sector by sector denote a POVM, block diagonal along her sectors, that returns them between square roots of the marginal and vanishes on a sector exactly where they do. Conversely a POVM and a marginal, both block diagonal, give such operators.}
\claimbookmark{meat-core/lem:ensemble-form}{Lemma}{lem:ensemble-form}{The measurement an ensemble-form declaration denotes}
\claimleadin{%
A declaration in \exttermref{../protocol-compiler/contract.pdf}{term-ensemble-form}{ensemble form} holds, in place of Alice's POVM,
positive operators $N_x$ that sum to her marginal sector by sector,
\eqref{eq:ensembleform}. The lemma shows that the operators
\eqref{eq:ensemblepovm} such a declaration denotes are a POVM, so that every
claim of this module applies to it as it stands; by part~(iii), every
predicate asking which outcomes can occur on a sector is decided on the
declared numbers. Part~(iv) is the converse; the operators it gives need not
be exactly representable.
}
\begin{lemma}[The measurement an ensemble-form declaration denotes]\label{lem:ensemble-form}
Let $\Hs{A}=\bigoplus_{t\in\mathbb{T}}\Hs{A_t}$, and for every $t$ let
$\tau^{(t)}_A\in\Pos{A_t}$ be non-zero and $N^{(t)}_x\in\Pos{A_t}$,
$x\in\Xset$, with $\sum_{x}N^{(t)}_x=\tau^{(t)}_A$. Let $P_t$ be the
projector onto the support of $\tau^{(t)}_A$, let
$(\tau^{(t)}_A)^{-1/2}$ be the inverse square root on that support, zero off
it, and let $M^{(t)}_x$ be the operator of \eqref{eq:ensemblepovm},
$M_x\coloneqq\bigoplus_tM^{(t)}_x$. Then:
\begin{enumerate}[label=(\roman*),leftmargin=2.2em]
\item $M^{(t)}_x\succeq0$ and $\sum_xM^{(t)}_x=\idon{A_t}$ for every $t$, so
$\{M_x\}_{x\in\Xset}$ is a POVM on $A$ and every element is block diagonal
along $\mathbb{T}$;
\item $(\tau^{(t)}_A)^{1/2}M^{(t)}_x(\tau^{(t)}_A)^{1/2}=N^{(t)}_x$, and in
particular $\tr[M_x\tau_A]=\sum_t\tr N^{(t)}_x$;
\item $M^{(t)}_x=0$ if and only if $N^{(t)}_x=0$;
\item conversely, if $\tau_A\in\Pos{A}$ and $\{M_x\}$ is a POVM on $A$, both
block diagonal along $\mathbb{T}$, then
$N_x\coloneqq\tau_A^{1/2}M_x\tau_A^{1/2}$ is positive semidefinite and block
diagonal along $\mathbb{T}$ and $\sum_xN_x=\tau_A$.
\end{enumerate}
\end{lemma}
\claimmeta{\metalabel{Uses}\mref{Def.}{def:declaration}, \mref{Def.}{def:admissible-declaration}
\metasep \metalabel{Used by}\mref{Def.}{def:joint-sectors}, \mref{Lem.}{lem:choi-factor}, \mref{Lem.}{lem:rank-one-sector}
\metasep newly derived
\metasep implemented (effective: accepted)
\metasep \metalabel{Code}\coderef{https://github.com/fraunhofersg/certiqs-key/blob/v0.1.0/src/certiqs_key/meat_core/declaration/general_validation.py}{certiqs_key.meat_core.declaration.general_validation:validate_general_declaration}}

\begin{proof}
Fix $t$ and drop it from the notation; write $S\coloneqq\tau_A^{-1/2}$, a
positive semidefinite operator with $S\tau_AS=P$ and
$\tau_A^{1/2}S=S\tau_A^{1/2}=P$, and $w_x\coloneqq\tr N_x/\tr\tau_A$, which
is defined because $\tau_A\ne0$ is positive semidefinite, hence of positive
trace, and satisfies $w_x\ge0$ and $\sum_xw_x=\tr\tau_A/\tr\tau_A=1$.

(i) $SN_xS\succeq0$ as a conjugation of a positive semidefinite operator by
a Hermitian one, and $w_x(\idon{}-P)\succeq0$; so $M_x\succeq0$. By
linearity $\sum_xM_x=S\tau_AS+(\idon{}-P)=P+(\idon{}-P)=\idon{}$. Block
diagonality of $M_x$ along $\mathbb{T}$ is its definition as a direct sum.

(ii) From $0\preceq N_x\preceq\sum_{x'}N_{x'}=\tau_A$ it follows that
$N_x$ vanishes on the kernel of $\tau_A$: for $\xi$ in that kernel
$0\le\bra{\xi}N_x\ket{\xi}\le\bra{\xi}\tau_A\ket{\xi}=0$, and a positive
semidefinite operator with $\bra{\xi}N_x\ket{\xi}=0$ has $N_x\ket{\xi}=0$.
Hence $PN_xP=N_x$. Then
$\tau_A^{1/2}M_x\tau_A^{1/2}=PN_xP+w_x\,\tau_A^{1/2}(\idon{}-P)\tau_A^{1/2}
=N_x$, because $\tau_A^{1/2}(\idon{}-P)=0$. Taking the trace,
$\tr[M_x\tau_A]=\tr N_x$ on the sector, and the sum over $t$ is the second
statement.

(iii) If $N_x=0$ then $w_x=0$ and both terms of $M_x$ vanish. If $N_x\ne0$
then $N_x=PN_xP$ is a non-zero operator on the support of $\tau_A$, on which
$S$ is invertible, so $SN_xS\ne0$; and $SN_xS$ is supported on the range of
$P$ while $w_x(\idon{}-P)$ is supported on its complement, so their sum is
non-zero.

(iv) $N_x$ is a conjugation of $M_x\succeq0$ by the Hermitian operator
$\tau_A^{1/2}$, which is block diagonal along $\mathbb{T}$ because
$\tau_A$ is, the square root acting block by block; and
$\sum_xN_x=\tau_A^{1/2}\idon{A}\tau_A^{1/2}=\tau_A$.
\end{proof}

% claim meat-core/def:maps
\claimsummary{meat-core/def:maps}{The declaration generates coherent measurement, announcement, statistic and key isometries, the single-round protocol map, the key pinching and the statistics map. They give the entropy proof a precise representation of the protocol's classical outputs.}
\claimbookmark{meat-core/def:maps}{Definition}{def:maps}{Generated registers, isometries and maps}
\begin{definition}[Generated registers, isometries and maps]\label{def:maps}
The classical registers are $X$ (alphabet $\Xset$), $Y$ (alphabet $\Vset$),
the announcement register $C_AC_B$ and $I$ (alphabet $\Kann$ each), $\Chat$
(alphabet $\Chatset$), $R$ and $Z$ (alphabet $\Zset$). The announcement
register is one register, named after the two messages of the basic format;
in the general format it holds the whole tuple $a(x,v)$ of messages. Define
\begin{align}
  \Vmsm &\coloneqq \sum_{x\in\Xset}\sum_{v\in\Vset}
    \ket{x,v}_{XY}\otimes\big(\sqrt{M_x}\otimes\sqrt{F_v}\big),
    &&\Hs{A\Qsys}\to\Hs{XY\,A\Qsys},\label{eq:Vmsm}\\
  \Vann &\coloneqq \sum_{x,v}\ket{a(x,v)}_{C_AC_B}\otimes\dyad{x,v}_{XY},
    &&\text{controlled on }XY,\label{eq:Vann}\\
  \VI &\coloneqq \sum_{\kappa\in\Kann}\ket{\kappa}_{I}\otimes
    \dyad{\kappa}_{C_AC_B},
    &&\text{controlled on }C_AC_B,\label{eq:VI}\\
  \VChat &\coloneqq \sum_{\kappa\in\Kann}\ket{\phi(\kappa)}_{\Chat}\otimes
    \dyad{\kappa}_{C_AC_B},
    &&\text{controlled on }C_AC_B,\label{eq:VChat}\\
  \VR &\coloneqq \sum_{x,v}\ket{r(x,v)}_{R}\otimes\dyad{x,v}_{XY},
    &&\text{controlled on }XY,\label{eq:VR}\\
  \VZ &\coloneqq \sum_{j\in\Zset}\ket{j}_{Z}\otimes Z_j,
    \qquad Z_j\coloneqq\dyad{j}_{R},
    &&\text{controlled on }R,\label{eq:VZ}
\end{align}
each acting as the identity on all registers not named. The full protocol
isometry and the
\termdef{term-single-round-protocol-map}{single-round protocol map}
are
\begin{equation}\label{eq:Mprot}
\begin{gathered}
  \Wprot\coloneqq\VR\,\VChat\,\VI\,\Vann\,\Vmsm,\\
  \Mprot(\omega)\coloneqq
  \tr_{A\Qsys\,XY\,C_AC_B\,R}\!\left[\VZ\,\Wprot\,\omega\,\Wprot^{\dagger}
  \VZ^{\dagger}\right],
  \qquad \Mprot: A\Qsys\to Z\Chat I .
\end{gathered}
\end{equation}
The
\termdef{term-key-pinching}{key pinching}
is $\cZ(\tau)\coloneqq\sum_{j\in\Zset}Z_j\,\tau\,Z_j$ on $R$, the pinching
of \xref{foundations/def:pinching} by the projectors $Z_j$ (extended by the
identity to any ambient system), and the
\termdef{term-statistics-map}{statistics map}
is
\begin{equation}\label{eq:Phimap}
  \Phimap[\omega]\coloneqq\sum_{\hat c\in\Chatset}
  \Bigg(\sum_{(x,v):\,\phi(a(x,v))=\hat c}
  \tr\!\big[(M_x\otimes F_v)\,\omega\big]\Bigg)\,\hat e_{\hat c},
  \qquad \Phimap:\Lin{A\Qsys}\to\mathbb{R}^{\Chatset},
\end{equation}
with $\hat e_{\hat c}$ the unit vector at position $\hat c$.
\end{definition}
\claimmeta{\metalabel{Used by}\mref{Thm.}{thm:protocol-channel}, \mref{Lem.}{lem:branch-maps}, \mref{Thm.}{thm:branch-formula}, \mref{Lem.}{lem:cone-concavity}, \mref{Lem.}{lem:joint-sectors}, \mref{Prop.}{prop:a-neutral}, \mref{Thm.}{thm:removal}, \mref{Lem.}{lem:choi-factor}, \mref{Cor.}{cor:full-gradient}, \mref{Lem.}{lem:kraus-freedom}, \mref{Lem.}{lem:rank-one-sector}, \mref{Cor.}{cor:drop-program}, \mref{Prop.}{prop:minimal-divergence}, \mref{Lem.}{lem:objective-homogeneous}, \mref{Prop.}{prop:model-pass-bound}, \mref{Req.}{req:companion-admission}
\metasep newly derived
\metasep implemented
\metasep \metalabel{Code}\coderef{https://github.com/fraunhofersg/certiqs-key/blob/v0.1.0/src/certiqs_key/meat_core/objective/branch.py}{certiqs_key.meat_core.objective.branch:generate_maps}
\metasep \metalabel{Oracles}\coderef{https://github.com/fraunhofersg/certiqs-key/blob/v0.1.0/tests/oracles/registry.py}{registry.py}}

% claim meat-core/thm:protocol-channel
\claimsummary{meat-core/thm:protocol-channel}{For an admissible declaration (\hyperlink{claim-meat-core/def:admissible-declaration}{Definition~\ref*{def:admissible-declaration}}) the single-round protocol map is a channel with classical key, acceptance and announcement registers, given by an explicit measurement formula. Its acceptance marginal is exactly the statistics map.}
\claimbookmark{meat-core/thm:protocol-channel}{Theorem}{thm:protocol-channel}{The declaration generates a protocol channel}
\begin{theorem}[The declaration generates a protocol channel]
\label{thm:protocol-channel}
Under Definition~\ref{def:admissible-declaration}:
\begin{enumerate}[label=(\roman*),leftmargin=2.2em]
\item Each of \eqref{eq:Vmsm}--\eqref{eq:VZ}, and hence $\Wprot$, is an
isometry, and $\Mprot$ is a channel $A\Qsys\to Z\Chat I$ with $\Chat$, $I$,
$Z$ classical.
\item For every system $E$ and every $\omega\in\Sleq{A\Qsys E}$,
\begin{equation}\label{eq:povm-form}
\begin{split}
  (\Mprot\otimes\idmap_E)(\omega)
  ={}&\sum_{x,v}
  \dyad{r(x,v)}_{Z}\otimes\dyad{\phi(a(x,v))}_{\Chat}\otimes{}\\
  &\qquad\dyad{a(x,v)}_{I}\otimes
  \tr_{A\Qsys}\!\big[(M_x\otimes F_v\otimes\idon{E})\,\omega\big].
\end{split}
\end{equation}
\item The $\Chat$-marginal of $(\Mprot\otimes\idmap_E)(\omega)$ is the
probability vector $\Phimap[\omega_{A\Qsys}]$; $\Phimap$ is affine, and
$\Phimap[\omega]$ is a probability distribution on $\Chatset$ for every
$\omega\in\Seq{A\Qsys}$.
\end{enumerate}
\end{theorem}
\claimmeta{\metalabel{Uses}\mref{Def.}{def:admissible-declaration}, \mref{Def.}{def:maps}
\metasep \metalabel{Used by}\mref{Lem.}{lem:branch-maps}, \mref{Def.}{def:roundmodel}, \mref{Thm.}{thm:accumulated-entropy}, \mref{Thm.}{thm:branch-formula}, \mref{Thm.}{thm:reduction}, \mref{Rem.}{rem:hlow-range}, \mref{Thm.}{thm:convexity}
\metasep newly derived
\metasep implemented (effective: accepted)
\metasep \metalabel{Code}\coderef{https://github.com/fraunhofersg/certiqs-key/blob/v0.1.0/src/certiqs_key/meat_core/objective/branch.py}{certiqs_key.meat_core.objective.branch:generate_maps}
\metasep \metalabel{Oracles}\coderef{https://github.com/fraunhofersg/certiqs-key/blob/v0.1.0/tests/oracles/registry.py}{registry.py}}

\begin{proof}
(i) With $\{\ket{x,v}\}$ orthonormal,
$\Vmsm^{\dagger}\Vmsm
=\sum_{x,v}\sqrt{M_x}\sqrt{M_x}\otimes\sqrt{F_v}\sqrt{F_v}
=\Big(\sum_xM_x\Big)\otimes\Big(\sum_vF_v\Big)
=\idon{A}\otimes\idon{\Qsys}$
by POVM completeness. Each of \eqref{eq:Vann}--\eqref{eq:VR} has the form
$\sum_{\mu}\ket{g(\mu)}\otimes\dyad{\mu}$ for a function $g$ on the control
alphabet, so $V^{\dagger}V=\sum_{\mu}\dyad{\mu}=\idon{}$ on the control
register; $\VZ^{\dagger}\VZ=\sum_jZ_j=\idon{R}$. Compositions of isometries
are isometries; conjugation by an isometry followed by a partial trace is
CPTP. $Z$, $\Chat$, $I$ are classical because \eqref{eq:povm-form} exhibits
the output diagonal in their computational bases.

(ii) Write $K_{x,v}\coloneqq\sqrt{M_x}\otimes\sqrt{F_v}$, so that
$\Vmsm\,\omega\,\Vmsm^{\dagger}
=\sum_{x,v,x',v'}\ketbra{x,v}{x',v'}_{XY}\otimes
(K_{x,v}\otimes\idon{E})\,\omega\,(K_{x',v'}\otimes\idon{E})^{\dagger}$.
Every subsequent isometry in $\Wprot$, and $\VZ$, is controlled on $XY$
directly or on a register whose content is a function of $XY$, so
conjugation by them appends the labels
$a(x,v)$, $a(x,v)$ (in $I$), $\phi(a(x,v))$, $r(x,v)$, $r(x,v)$ (in $Z$) to
the ket side and the primed labels to the bra side of each term. The partial
trace in \eqref{eq:Mprot} includes $\tr_{XY}$, and
$\tr_{XY}\ketbra{x,v}{x',v'}=\delta_{xx'}\delta_{vv'}$, so only diagonal
terms survive; on those, $\tr_{A\Qsys}[K_{x,v}\,\omega\,K_{x,v}^{\dagger}]
=\tr_{A\Qsys}[(M_x\otimes F_v)\,\omega]$ by cyclicity of the partial trace
over $A\Qsys$. This is \eqref{eq:povm-form}.

(iii) Take the $\Chat$-marginal of \eqref{eq:povm-form} and group terms by
$\hat c=\phi(a(x,v))$; tracing $E$ replaces the last factor by
$\tr[(M_x\otimes F_v)\,\omega_{A\Qsys}]$. Affinity is manifest;
normalization follows from
$\sum_{x,v}\tr[(M_x\otimes F_v)\omega]=\tr[\omega]=1$.
\end{proof}

% claim meat-core/def:branch
\claimsummary{meat-core/def:branch}{Each acceptance value splits its announcement classes into pass and fail classes, and the two branch maps sum the Kraus terms of each, sector by sector as well. They keep the coherence the entropy calculation needs while separating key-bearing rounds from discarded ones.}
\claimbookmark{meat-core/def:branch}{Definition}{def:branch}{Branch maps}
\begin{definition}[Branch maps]\label{def:branch}
For $\hat c\in\Chatset$, a class $\kappa\in\phi^{-1}(\hat c)$ with
$\varsigma(\kappa)=\mathrm{pass}$ is a
\termdef{term-pass-class}{pass class}
of $\hat c$, and one with $\varsigma(\kappa)=\mathrm{fail}$ a
\termdef{term-fail-class}{fail class};
let
$\Kann_{\mathrm{pass}}(\hat c)\coloneqq
\{\kappa\in\phi^{-1}(\hat c):\varsigma(\kappa)=\mathrm{pass}\}$ and
$\Kann_{\mathrm{fail}}(\hat c)\coloneqq
\{\kappa\in\phi^{-1}(\hat c):\varsigma(\kappa)=\mathrm{fail}\}$. For each
class $\kappa$ define the operator
\begin{equation}\label{eq:kraus}
\begin{split}
  K_{\kappa}\coloneqq{}&\sum_{(x,v)\in a^{-1}(\kappa)}
  \ket{r(x,v)}_{R}\otimes\ket{x,v}_{XY}\otimes\ket{\kappa}_{C_AC_B}
  \otimes\big(\sqrt{M_x}\otimes\sqrt{F_v}\big)\\
  &\;:\;\Hs{A\Qsys}\to\Hs{R\,XY\,C_AC_B\,A\Qsys},
\end{split}
\end{equation}
and the CP
\termdef{term-branch-map}{branch map}s,
the pass map $\cG_{\hat c}$ and the fail map $\Theta_{\hat c}$,
\begin{equation}\label{eq:GTheta}
  \cG_{\hat c}(\omega)\coloneqq
  \sum_{\kappa\in\Kann_{\mathrm{pass}}(\hat c)}K_{\kappa}\,\omega\,
  K_{\kappa}^{\dagger},
  \qquad
  \Theta_{\hat c}(\omega)\coloneqq
  \sum_{\kappa\in\Kann_{\mathrm{fail}}(\hat c)}K_{\kappa}\,\omega\,
  K_{\kappa}^{\dagger}.
\end{equation}
For $s\in\SecSet$, the sector versions $K_{\kappa}^{(s)}$, $\cG_{\hat c,s}$,
$\Theta_{\hat c,s}$ are \eqref{eq:kraus}--\eqref{eq:GTheta} with
$\sqrt{F_v}$ replaced by $\sqrt{F_v^{(s)}}$, acting
$\Hs{A\Qsys_s}\to\Hs{R\,XY\,C_AC_B\,A\Qsys_s}$.
\end{definition}
\claimmeta{\metalabel{Used by}\mref{Lem.}{lem:branch-maps}, \mref{Thm.}{thm:branch-formula}, \mref{Prop.}{prop:exact-cancellation}, \mref{Def.}{def:joint-sectors}, \mref{Lem.}{lem:joint-sectors}, \mref{Prop.}{prop:a-neutral}, \mref{Prop.}{prop:retention}, \mref{Lem.}{lem:choi-factor}, \mref{Def.}{def:branch-support}, \mref{Lem.}{lem:branch-support}, \mref{Lem.}{lem:kraus-freedom}, \mref{Lem.}{lem:rank-one-sector}, \mref{Req.}{req:companion-admission}, \mref{Req.}{req:branch-register-components}, \mref{Lem.}{lem:certified-bracket-enclosed}
\metasep newly derived
\metasep implemented
\metasep \metalabel{Code}\coderef{https://github.com/fraunhofersg/certiqs-key/blob/v0.1.0/src/certiqs_key/meat_core/objective/branch.py}{certiqs_key.meat_core.objective.branch:ClassBranch}
\metasep \metalabel{Oracles}\coderef{https://github.com/fraunhofersg/certiqs-key/blob/v0.1.0/tests/oracles/registry.py}{registry.py}}

% claim meat-core/lem:branch-maps
\claimsummary{meat-core/lem:branch-maps}{The pass and fail maps are trace-non-increasing, their traces add up to the acceptance value's probability, their outputs have orthogonal supports and the fail branch is fixed by the key pinching. They split by sector on block-diagonal inputs, and a value that is not key-live has a zero pass map.}
\claimbookmark{meat-core/lem:branch-maps}{Lemma}{lem:branch-maps}{Properties of the pass and fail branch maps}
\begin{lemma}[Properties of the pass and fail branch maps]\label{lem:branch-maps}
Under Definition~\ref{def:admissible-declaration}, for every $\hat c\in\Chatset$:
\begin{enumerate}[label=(\roman*),leftmargin=2.2em]
\item $\cG_{\hat c}$ and $\Theta_{\hat c}$ are CP and trace-non-increasing,
with
\begin{equation}\label{eq:traces}
  \tr\cG_{\hat c}(\omega)+\tr\Theta_{\hat c}(\omega)
  =\Phimap[\omega](\hat c)
  \qquad\text{for all }\omega\in\Pos{A\Qsys}.
\end{equation}
\item $\cG_{\hat c}(\omega)$ is supported on
$\operatorname{span}\{\ket{j}_R:j\in\Zval\}$ and
$\Theta_{\hat c}(\omega)=\dyad{\perp}_{R}\otimes\theta_{\hat c}(\omega)$ for
a CP map $\theta_{\hat c}$; in particular their supports are orthogonal and
$\cZ(\Theta_{\hat c}(\omega)^{p})=\Theta_{\hat c}(\omega)^{p}$ for every
$p>0$.
\item (Sector decomposition.) If
$\omega=\bigoplus_{s\in\SecSet}\omega_s$ is block-diagonal with respect to
the decomposition $\Hs{A\Qsys}=\bigoplus_s\Hs{A\Qsys_s}$, then
\begin{equation}\label{eq:G-sector}
  \cG_{\hat c}(\omega)=\bigoplus_{s\in\SecSet}\cG_{\hat c,s}(\omega_s),
  \qquad
  \Theta_{\hat c}(\omega)=\bigoplus_{s\in\SecSet}\Theta_{\hat c,s}(\omega_s),
  \qquad
  \Phimap[\omega]=\sum_{s\in\SecSet}\Phimap_s[\omega_s],
\end{equation}
where the direct sums refer to the (mutually orthogonal) subspaces
$\Hs{R\,XY\,C_AC_B}\otimes\Hs{A\Qsys_s}$ and
$\Phimap_s$ is \eqref{eq:Phimap} with $F_v^{(s)}$ in place of $F_v$.
\item If $\hat c\notin\Chkey$ then $\cG_{\hat c}=0$.
\end{enumerate}
The complete positivity of $\cG_{\hat c}$ and $\Theta_{\hat c}$ in (i),
and all of (iii), hold as well when $\{F_v\}$ is any family of
sector-block-diagonal positive semidefinite operators, complete or not,
the other components satisfying Definition~\ref{def:admissible-declaration}.
\end{lemma}
\claimmeta{\metalabel{Uses}\mref{Def.}{def:declaration}, \mref{Def.}{def:admissible-declaration}, \mref{Def.}{def:maps}, \mref{Thm.}{thm:protocol-channel}, \mref{Def.}{def:branch}
\metasep \metalabel{Used by}\mref{Thm.}{thm:accumulated-entropy}, \mref{Thm.}{thm:branch-formula}, \mref{Cor.}{cor:concavity}, \mref{Lem.}{lem:cone-concavity}, \mref{Cor.}{cor:f-nonpos}, \mref{Prop.}{prop:exact-cancellation}, \mref{Lem.}{lem:joint-sectors}, \mref{Prop.}{prop:a-neutral}, \mref{Lem.}{lem:choi-factor}, \mref{Thm.}{thm:gradient}, \mref{Prop.}{prop:relaxed-set-properties}, \mref{Lem.}{lem:fo-face}
\metasep newly derived
\metasep implemented (effective: accepted)
\metasep \metalabel{Code}\coderef{https://github.com/fraunhofersg/certiqs-key/blob/v0.1.0/src/certiqs_key/meat_core/objective/branch.py}{certiqs_key.meat_core.objective.branch:evaluate_branches}
\metasep \metalabel{Oracles}\coderef{https://github.com/fraunhofersg/certiqs-key/blob/v0.1.0/tests/oracles/registry.py}{registry.py}}

\begin{proof}
(i) CP is Kraus form. For the trace, note that for a single class $\kappa$,
the register kets $\ket{r(x,v),x,v,\kappa}$ appearing in
$K_{\kappa}$ are mutually orthonormal as $(x,v)$ ranges over
$a^{-1}(\kappa)$ (they differ in $XY$), so
$\tr[K_{\kappa}\omega K_{\kappa}^{\dagger}]
=\sum_{(x,v)\in a^{-1}(\kappa)}\tr[(M_x\otimes F_v)\,\omega]$ by cyclicity.
Summing over $\kappa\in\phi^{-1}(\hat c)$ (pass and fail together) and
comparing with \eqref{eq:Phimap} gives \eqref{eq:traces}; each summand of
\eqref{eq:traces} is non-negative, so
$\tr\cG_{\hat c}(\omega)$ and $\tr\Theta_{\hat c}(\omega)$ are each at
most $\Phimap[\omega](\hat c)$. The remaining step,
$\Phimap[\omega](\hat c)\le1$ for a state $\omega$, is not a property of the
branch maps: it is
Theorem~\ref{thm:protocol-channel}\,(iii), that $\Phimap[\omega]$ is a probability
distribution on $\Chatset$, whose proof consumes the completeness of the
declared POVMs, $\sum_vF_v=\idon{\Qsys}$ and $\sum_xM_x=\idon{A}$
(Definition~\ref{def:admissible-declaration}). Hence both maps are trace-non-increasing on
states.

(ii) By announcement-determined sifting
(Definition~\ref{def:admissible-declaration}\,(v)), $r(x,v)\in\Zval$ for every $(x,v)$ in a
pass class and $r(x,v)=\perp$ for every $(x,v)$ in a fail class. Hence every
Kraus operator of $\cG_{\hat c}$ carries an $R$-ket in $\Zval$ and every
Kraus operator of $\Theta_{\hat c}$ carries $\ket{\perp}_R$; the $\perp$
value factors out of the sum defining $\Theta_{\hat c}$. The pinching $\cZ$
acts as the identity on any operator supported on a single $R$-eigenspace,
and powers preserve supports.

(iii) By blockwise functional calculus,
$\sqrt{F_v}=\bigoplus_s\sqrt{F_v^{(s)}}$, so
$K_{\kappa}=\sum_s K_{\kappa}^{(s)}$ with
$K_{\kappa}^{(s)}=K_{\kappa}(\idon{A}\otimes\Pi_s)$ mapping
$\Hs{A\Qsys_s}$ into $\Hs{R\,XY\,C_AC_B}\otimes\Hs{A\Qsys_s}$ (the operators
$\sqrt{M_x}\otimes\sqrt{F_v^{(s)}}$ do not leave the sector). For
block-diagonal $\omega$,
$K_{\kappa}\,\omega\,K_{\kappa}^{\dagger}
=\sum_{s,s'}K_{\kappa}^{(s)}\,\omega\,K_{\kappa}^{(s')\dagger}$
and $K_{\kappa}^{(s)}\omega\,K_{\kappa}^{(s')\dagger}
=K_{\kappa}^{(s)}(\idon{}\otimes\Pi_s)\,\omega\,(\idon{}\otimes\Pi_{s'})
K_{\kappa}^{(s')\dagger}$ vanishes for $s\ne s'$ because
$(\idon{}\otimes\Pi_s)\,\omega\,(\idon{}\otimes\Pi_{s'})=0$. The images for
different $s$ lie in orthogonal subspaces, giving the direct sums; the
$\Phimap$ statement follows from linearity of the trace and block-diagonality
of $F_v$.

(iv) If $\hat c\notin\Chkey$ there are no pass classes in $\phi^{-1}(\hat c)$
by \eqref{eq:chkey}, so the sum \eqref{eq:GTheta} defining
$\cG_{\hat c}$ is empty.

\emph{Any positive family.} The Kraus-form argument for complete
positivity in (i) and the argument for (iii) read of $\{F_v\}$ only that
each $F_v$ is positive semidefinite, so that $\sqrt{F_v}$ is defined, and
sector-block-diagonal; neither reads $\sum_vF_v=\idon{\Qsys}$, which
enters (i) only through Theorem~\ref{thm:protocol-channel}\,(iii) for the
bound on the traces.
\end{proof}

% claim meat-core/rem:inert
\claimsummary{meat-core/rem:inert}{An isometry on registers other than the raw-key register leaves the pinched trace power unchanged, so no quantity of this module depends on a component of another register that is constant across a branch map's Kraus terms. Varying components and the raw-key register carry what the proofs use.}
\claimbookmark{meat-core/rem:inert}{Remark}{rem:inert}{Inert register components of a branch map}
\claimleadin{%
Theorem~\ref{thm:branch-formula} writes the entropy on an acceptance value
through the pinched trace power
$\tr[(\cZ(\cG_{\hat c}(\omega)^{1/\alpha}))^{\alpha}]$ of the pass map's
output, the trace of that output and the probability of the value. The
remark shows that such a trace power does not see an isometry on registers
other than the raw-key register $R$.
}
\begin{remark}[Inert register components of a branch map]\label{rem:inert}
If $U$ is an isometry acting only on registers disjoint from $R$, then
$(U\sigma U^{\dagger})^{p}=U\sigma^{p}U^{\dagger}$ and
$\cZ(U\sigma^{p}U^{\dagger})=U\cZ(\sigma^{p})U^{\dagger}$, so
$\tr[(\cZ((U\sigma U^{\dagger})^{1/\alpha}))^{\alpha}]
=\tr[(\cZ(\sigma^{1/\alpha}))^{\alpha}]$. Consequently a register component
on a register other than $R$ whose value is \emph{constant across all Kraus
terms of a given branch map} --- for example the $C_AC_B$ factor when a
branch has a single announcement class --- can be removed from that branch
map without changing any quantity in this module: its constant ket
$\ket{c}$ enters every Kraus operator through the isometry
$\ket{\psi}\mapsto\ket{c}\otimes\ket{\psi}$, which acts only on registers
disjoint from $R$. Components that vary across Kraus terms (the class label
$\kappa$ across classes, the outcome labels within a class) are not of this
kind: they carry the orthogonality that the derivations of this module use
(Lemma~\ref{lem:branch-maps}). Nor is the raw-key register $R$ itself, even
where its value is constant, as it is ($\perp$) across every Kraus term of
a fail map (Lemma~\ref{lem:branch-maps}\,(ii)): the pinching $\cZ$ acts on
$R$, and $R$ is what separates the pass branch from the fail branch.
\end{remark}
\claimmeta{\metalabel{Used by}\mref{Req.}{req:branch-register-components}
\metasep newly derived
\metasep accepted}

\section{Accumulation}\label{group-meat-core/fact:meat}
\sectionintro{The imported sequential accumulation theorem bounds the entropy of an accepted multi-round state from below by the round count times a single-round rate, less a penalty in the acceptance probability (\hyperlink{claim-meat-core/fact:meat}{Fact~\ref*{fact:meat}}). Instantiated on the generated \termref{term-round-channel}{round channels}, with every announcement handed to the adversary and her quantum memory carried between rounds, it gives the section's main result, \hyperlink{claim-meat-core/thm:accumulated-entropy}{Theorem~\ref*{thm:accumulated-entropy}}: on acceptance, the entropy of the raw key given the adversary's information is at least n times the single-round rate less (\ensuremath{\alpha}/(\ensuremath{\alpha}\ensuremath{-}1))\ensuremath{\cdot}log\ensuremath{_{2}}(1/Pr[accept]), with the product marginal of Alice's systems in trusted-source mode derived from the source's premises, A-SEED (v) and A-LAB at the party holding A (\hyperlink{claim-meat-core/lem:product-marginal}{Lemma~\ref*{lem:product-marginal}}), and with no source assumption in untrusted-source mode. Licensed constraint lists, \termref{term-licensed-attack-set}{licensed attack sets} and the \termref{term-constrained-single-round-rate}{constrained single-round rate} are defined here; a constraint list may raise the rate only through a licensed restriction of the attack model (\hyperlink{claim-meat-core/rem:no-constraints}{Remark~\ref*{rem:no-constraints}}).}
\sectionimports{\extref{../foundations/contract.pdf}{foundations/def:kl}, \extref{../foundations/contract.pdf}{foundations/def:renyi}, \extref{../foundations/contract.pdf}{foundations/def:round-model}, \extref{../protocol-compiler/contract.pdf}{protocol-compiler/def:homog}, \extref{../transmitter/contract.pdf}{transmitter/def:source-segment}, \extref{../transmitter/contract.pdf}{transmitter/def:source-spec}, \extref{../transmitter/contract.pdf}{transmitter/def:srcert}, \extref{../transmitter/contract.pdf}{transmitter/thm:certificate-sufficiency}}

% claim meat-core/fact:meat
\claimsummary{meat-core/fact:meat}{On an accepted event the R\'enyi entropy of n rounds of sequential channels with quantum memory, under a product marginal constraint, is at least n times a single-round rate less (\ensuremath{\alpha}/(\ensuremath{\alpha}\ensuremath{-}1))\ensuremath{\cdot}log\ensuremath{_{2}}(1/p\_\ensuremath{\Omega}). The rate minimises a KL penalty plus conditional entropies over the channels' convex range.}
\claimbookmark{meat-core/fact:meat}{Fact}{fact:meat}{Marginal-constrained entropy accumulation, sequential form}
\claimleadin{%
The imported theorem is stated below in this module's notation. The locator
notes after the statement tie each display to its place in the source, and
note~(vii) checks that a one-dimensional marginal register is an admissible
instance, on which the \exttermref{../transmitter/contract.pdf}{term-untrusted-source-mode}{untrusted-source mode} case of
Theorem~\ref{thm:accumulated-entropy} rests. Untrusted-source mode on the
program is the subject of Remark~\ref{rem:modes-da1}, whose part (b) quotes
the sentence following \cite[Eq.~(1)]{arqand2025meat} that note~(vii)
cites. The Fact is the source's theorem in its published form, on
$\alpha<\infty$ (note~(vi)), quoted once for both source modes. The source states no separate result for trivial
marginal registers: that its hypotheses admit them is note~(vii)'s reading
of the source at the locators it quotes, and the claim that instantiates the
Fact there, Theorem~\ref{thm:accumulated-entropy}, is newly derived.
}
\begin{fact}[Marginal-constrained entropy accumulation, sequential form]\label{fact:meat}
\emph{(\cite[Thm.~4.2a, Eqs.~(100)--(101)]{arqand2025meat}.)}
For each $j\in\{1,\dots,n\}$ take a state
$\sigma^{(j-1)}\in\Seq{A_{j-1}}$ and a channel
$\Mprot_j\in\mathrm{CPTP}(A_{j-1}E_{j-1},\,S_j\Chat_jE_j)$ with
$\Chat_j$ classical. Let
\begin{equation*}
  \rho_{S_1^n\Chat_1^nE_n}
  =\big(\Mprot_n\circ\dots\circ\Mprot_1\big)\big[\omega^{0}\big],
  \qquad
  \omega^{0}\in\Seq{A_0^{n-1}E_0},
\end{equation*}
subject to the product marginal constraint
$\omega^{0}_{A_0^{n-1}}=\sigma^{(0)}\otimes\dots\otimes\sigma^{(n-1)}$;
$E_0$ is unconstrained, and $\omega^{0}$ ranges over \emph{all}
normalized states compatible with the marginal constraint (purity is
not assumed). Suppose
$\rho=p_{\Omega}\,\rho_{|\Omega}+(1-p_{\Omega})\rho_{|\bar\Omega}$
with $p_{\Omega}\in(0,1]$ and $\rho_{|\Omega},\rho_{|\bar\Omega}$
normalized, all $\Chat_j$ isomorphic to one register $\Chat$ with
alphabet $\Chatset$, and let $S_{\Omega}$ be a convex set of
probability distributions on $\Chatset$ such that
$\operatorname{freq}_{\hat c_1^n}\in S_{\Omega}$ for every frequency
vector with nonzero probability in $\rho_{|\Omega}$. Then, for every
$\alpha\in(1,\infty)$ (locator note (vi): the source also covers the
endpoint $\alpha=\infty$, which no claim of this module uses),
\begin{equation}\label{eq:meat}
  \Hsandup\big(S_1^n\big|\Chat_1^nE_n\big)_{\rho_{|\Omega}}
  \;\ge\;n\,\hup-\frac{\alpha}{\alpha-1}\log_2\frac{1}{p_{\Omega}},
\end{equation}
where
\begin{equation}\label{eq:meat-rate}
  \hup=\inf_{\mathbf q\in S_{\Omega}}\ \inf_{\nu\in\Sigma}
  \left(\frac{\alpha}{\alpha-1}\,
  \sfD\!\left(\mathbf q\,\middle\|\,\boldsymbol{\nu}_{\Chat}\right)
  +\sum_{\hat c\in\operatorname{supp}(\boldsymbol{\nu}_{\Chat})}
  q(\hat c)\,\Hsandup(S|E\Etil)_{\nu_{|\hat c}}\right),
\end{equation}
with $\sfD$ the classical (base-2) KL divergence of
\xref{foundations/def:kl} and $\Sigma=\Sigma_{S\Chat E\Etil}$ the
\emph{marginal-constrained convex range}
(\cite[Def.~4.2, Eq.~(98)]{arqand2025meat}) of the channels
$\Mprot_j\otimes\idmap_{\Etil}$ and states $\sigma^{(j-1)}$:
\begin{equation}\label{eq:meat-range}
  \Sigma\coloneqq\operatorname{conv}\!\Big(\bigcup_{j=1}^{n}\Sigma_j\Big),
  \qquad
  \Sigma_j\coloneqq\big\{(\Mprot_j\otimes\idmap_{\Etil})[\omega]:
  \omega\in\Seq{A_{j-1}E_{j-1}\Etil},\
  \omega_{A_{j-1}}=\sigma^{(j-1)}\big\},
\end{equation}
where every output is embedded in common registers $S\Chat E\Etil$ (the
results hold for any choice of embeddings;
\cite[footnote~12]{arqand2025meat}) and $\Etil$ is a register of large
enough dimension to serve as a purifying register for any of the
$A_{j-1}E_{j-1}$ registers (enlargeable without loss of generality;
\cite[footnote~17]{arqand2025meat}).

\emph{Locator notes.} (i) The statement above is the source's
Thm.~4.2a with its registers written in this module's fonts; its
Eq.~(100) is \eqref{eq:meat}, its Eq.~(101) is \eqref{eq:meat-rate},
and its Def.~4.2, Eq.~(98) is \eqref{eq:meat-range}. Thm.~4.2a is the
special case of the source's Thm.~4.2b in which the per-round secret
test registers (the unbarred $C_j$ of the source) are trivial, which is
the case Theorem~\ref{thm:protocol-channel} generates. The source states
its entropies in base~2
(bits). (ii) The divergence in the penalty term is the \emph{KL}
divergence $\sfD$: the source obtains the rate objective as the convex
conjugate of a log-sum-exponential
(\cite[Lemma~4.11, Eqs.~(113)--(114)]{arqand2025meat}) and evaluates
it in its Eq.~(118), stating there explicitly that the first term ``is
the KL divergence''.
(iii) An alternative import is the tensor-product
specialization \cite[Cor.~4.2, Eq.~(137)]{arqand2025meat}, in which
the round channels carry no $E_j$ output and Eve's entire action is
the choice of one global input state. The source presents that
corollary as a register substitution of the theorems above, and its
Sec.~5 applies the resulting analysis when every public announcement
is made only after \emph{all} measurements are complete; for per-round
(``on-the-fly'') announcements --- the model
Theorem~\ref{thm:protocol-channel} generates, whose round channel
announces $I_j$ in every round --- the source's own
route (its Sec.~5, Protocols~1--2 and the reduction between them) is
the sequential form quoted here, in which Eve carries a quantum memory
$E_{j-1}$ between rounds and may adapt each round's action to all
earlier announcements. The present import therefore covers a strictly
larger attack model than that specialization would. (iv) The theorem has no
announcement register of its own: announcements enter the sequential
model inside the $E_j$ outputs (the source's Sec.~5: Eve
``immediately incorporates'' each $I_j$ into her side information).
(v) Indexing is off by one between this module's rounds and the
source's registers (\cite[Rem.~5.1]{arqand2025meat}): round $j$ here
uses the source's $A_{j-1}E_{j-1}$. (vi) The source states the
theorem for $\alpha\in(1,\infty]$, and its Def.~2.6 defines the
entropies at the endpoint; \xref{foundations/def:renyi}
defines the sandwiched family on $(1,\infty)$ only, so the Fact is
quoted, and Theorems \ref{thm:accumulated-entropy} and \ref{thm:constrained-accumulation} are stated, on
$(1,\infty)$. No claim uses the endpoint: every key-length statement
is confined to $(1,2)$.
(vii) \emph{Trivial marginal registers are admissible} --- the import
check for \exttermref{../transmitter/contract.pdf}{term-untrusted-source-mode}{untrusted-source mode}, read off the source at the
locators named. The hypotheses that
\cite[Thm.~4.2a]{arqand2025meat} places on the marginal registers are,
in full, that one ``take a state $\sigma^{(j-1)}\in S_=(A_{j-1})$, and
a channel $M_j\in\mathrm{CPTP}(A_{j-1}E_{j-1},S_j\widehat C_jE_j)$,
such that $\widehat C_j$ are classical'', together with the product
condition
$\omega_{A_0^{n-1}}=\sigma^{(0)}_{A_0}\otimes\dots\otimes
\sigma^{(n-1)}_{A_{n-1}}$ on the input state; and
\cite[Def.~4.2, Eq.~(98)]{arqand2025meat} likewise only ``consider[s]
channels $E_j\in\mathrm{CPTP}(A_{j-1}\bar A_{j-1},Z_j)$ and states
$\sigma^{(j-1)}\in S_=(A_{j-1})$'' (in the source's symbols, where $E_j$ is a
channel and $Z_j$ its output register, not this module's register $E_j$ and
raw key $Z_j$). \emph{No dimension bound and no
rank condition is imposed on $A_{j-1}$ at either locator}, and the
source's standing convention is explicit that none is imposed globally
either: ``we will assume that all systems are finite-dimensional, but
we will not impose any bounds on the system dimensions unless otherwise
specified'' (\cite[Sec.~2, the convention stated below its Table~1]%
{arqand2025meat}). A one-dimensional $A_{j-1}$ is therefore an
admissible instantiation, and at it the product marginal condition is
satisfied by every normalized $\omega$, since $\Seq{A_0^{n-1}}$ is then
a singleton --- the hypothesis discharges vacuously rather than
failing. Independently, the source reads the same equivalence into its
own definition of channel conditional entropy: see the sentence
following \cite[Eq.~(1)]{arqand2025meat}. Nothing else in Fact~\ref{fact:meat}
mentions $A_{j-1}$ beyond the purifier requirement, which asks $\Etil$
to purify $A_{j-1}E_{j-1}$ and is only weakened by shrinking
$A_{j-1}$.
\end{fact}
\claimmeta{\metalabel{Uses}\extref{../foundations/contract.pdf}{foundations/def:kl}, \extref{../foundations/contract.pdf}{foundations/def:renyi}
\metasep \metalabel{Used by}\mref{Def.}{def:roundmodel}, \mref{Lem.}{lem:range}, \mref{Thm.}{thm:accumulated-entropy}, \mref{Rem.}{rem:no-constraints}, \mref{Thm.}{thm:constrained-accumulation}, \mref{Lem.}{lem:b-derivative}
\metasep published in \cite[{Thm. 4.2a, Eqs. (100)-(101)}]{arqand2025meat}; \cite[{Def. 4.2, Eq. (98)}]{arqand2025meat}; \cite[{footnote 12}]{arqand2025meat}; \cite[{footnote 17}]{arqand2025meat}; \cite[{Lemma 4.11, Eqs. (113)-(114)}]{arqand2025meat}; \cite[{Cor. 4.2, Eq. (137)}]{arqand2025meat}; \cite[{Rem. 5.1}]{arqand2025meat}; \cite[{Thm. 4.2a}]{arqand2025meat}; \cite[{Eq. (1)}]{arqand2025meat}
\metasep accepted
\metasep \metalabel{Oracles}\coderef{https://github.com/fraunhofersg/certiqs-key/blob/v0.1.0/tests/oracles/registry.py}{registry.py}}

% claim meat-core/def:homog
\claimsummary{meat-core/def:homog}{A constraint list is licensed when it is licensed for the declared source mode in the sense of \xref{protocol-compiler/def:homog}, a type check that compares no numbers. Its classes are \ensuremath{\tau}\_A-homogeneous lists in trusted mode and scalar lists in untrusted mode; the general format adds party-local entries.}
\claimbookmark{meat-core/def:homog}{Definition}{def:homog}{Licensed constraint list and its two classes}
\claimleadin{%
The declared forms of a constraint entry and the class of licensed lists are
the compiler's (\xref{protocol-compiler/def:homog}); this definition names the
two classes as this module's claims read them. The security chain of this
module, Theorems~\ref{thm:constrained-accumulation} and~\ref{thm:keylength},
admits only licensed lists; every structural and numerical claim is stated
and proved for arbitrary Hermitian $\Theta_k$ (trusted mode) or arbitrary real
$c_k$ (untrusted mode). Remark~\ref{rem:homog}\,(a) exhibits, in
\exttermref{../transmitter/contract.pdf}{term-untrusted-source-mode}{untrusted-source mode}, a \termref{term-licensed-attack-set}{licensed attack set} for which imposing the operator form
$\Theta_k=c_k\idon{A}$ would empty the \exttermref{../protocol-compiler/contract.pdf}{term-feasible-set}{feasible set}; what licenses a
party-local entry is \xref{protocol-compiler/rem:party-local-soundness}, and
what the constrained accumulation theorem says of it is
Remark~\ref{rem:tagged-entries-and-accumulation}.
}
\begin{definition}[Licensed constraint list and its two classes]\label{def:homog}
The declared forms of an entry of $\CList$ --- \exttermref{../protocol-compiler/contract.pdf}{term-homogeneous}{homogeneous}, carrying
a real scalar $c_k$, \exttermref{../protocol-compiler/contract.pdf}{term-general}{general}, carrying a Hermitian $\Theta_k$ on $A$,
and, in the general format, \exttermref{../protocol-compiler/contract.pdf}{term-party-local}{party-local} --- are those of
\xref{protocol-compiler/def:homog}, and the list is \emph{licensed} when it
is \exttermref{../protocol-compiler/contract.pdf}{term-licensed-for-the-declared-source-mode}{licensed for the declared source mode} there. Its two classes
for an untagged entry are the following; in the general format an untagged
entry is read for its \exttermref{../protocol-compiler/contract.pdf}{term-lifted-operator}{lifted operator}, and the cited class also
licenses a party-local entry in either mode (the paragraph on party-local
entries).
\begin{enumerate}[label=(\alph*),leftmargin=2.2em]
\item \emph{\exttermref{../transmitter/contract.pdf}{term-trusted-source-mode}{trusted-source mode}:} every entry of $\CList$ is declared
in the homogeneous form, carrying $c_k\in\mathbb{R}$; the right-hand side it
denotes is then \emph{constructed} as $\Theta_k\coloneqq c_k\,\tau_A$ from the
declared $c_k$ and the declared $\tau_A$, and such a list is a
\termdef{term-tau-a-homogeneous-list}{$\tau_A$-homogeneous list}
by construction;
\item \emph{\exttermref{../transmitter/contract.pdf}{term-untrusted-source-mode}{untrusted-source mode}:} $\CList$ is a
\termdef{term-scalar-list}{scalar list},
i.e.\ every right-hand side is a real number $\Theta_k=c_k$ and the entry
reads $\tr[(\idon{A}\otimes\hat O_k)\,\omega]\bowtie_k c_k$ --- a
comparison of real numbers, with $\bowtie_k\in\{\le,=,\ge\}$.
\end{enumerate}
Licensing is therefore a \emph{type} check in both modes and compares no
numbers. A general entry remains admissible input, every structural and
numerical claim of this module being proved for it, but it is never licensed,
so a list containing one carries no key-length claim, which is the safe
direction. The two classes agree at $d_A=1$
(Proposition~\ref{prop:trivial-alice-system}\,(iii)), and (b) is what the
licensing condition \eqref{eq:attset-un} transports in untrusted-source
mode, through \eqref{eq:transport-un}:
\emph{the operator form $\Theta_k=c_k\idon{A}$ is not licensed}.

\emph{Party-local entries.} In the general format an entry may be declared
in the \exttermref{../reductions/contract.pdf}{term-party-local-form}{party-local form} (\xref{protocol-compiler/def:homog}); it is
written $(\hat O_k,\bowtie_k,c_k;f)$, and its state-level reading, in either
source mode, is
\begin{equation}\label{eq:partylocal-entry}
  \tr_{\Qsys_f}\!\Big[\big(\idon{}\otimes
  (\hat O_k-c_k\,\idon{\Qsys_f})\big)\,\omega\Big]\;\bowtie_k\;0
  \qquad\text{as an operator on the systems of }\omega
  \text{ other than }\Qsys_f ,
\end{equation}
those systems being $A$ and the \exttermref{../protocol-compiler/contract.pdf}{term-measured-factor}{measured factors}
$\Qsys_{f'}$, $f'\ne f$.
\end{definition}
\claimmeta{\metalabel{Uses}\mref{Def.}{def:declaration}, \extref{../protocol-compiler/contract.pdf}{protocol-compiler/def:homog}
\metasep \metalabel{Used by}\mref{Def.}{def:admissible-declaration}, \mref{Def.}{def:roundmodel}, \mref{Def.}{def:hqkd-list}, \mref{Lem.}{lem:range}, \mref{Rem.}{rem:no-constraints}, \mref{Thm.}{thm:constrained-accumulation}, \mref{Ex.}{ex:fss}, \mref{Rem.}{rem:homog}, \mref{Rem.}{rem:scalar-declaration}, \mref{Rem.}{rem:not-licensed}, \mref{Def.}{def:feasible}, \mref{Prop.}{prop:trivial-alice-system}, \mref{Prop.}{prop:a-neutral}, \mref{Cor.}{cor:reduction-constrained}, \mref{Prop.}{prop:frame-free-rows}, \mref{Rem.}{rem:basis-invariance}, \mref{Thm.}{thm:keylength}, \mref{Rem.}{rem:cgate}, \mref{Prop.}{prop:reference-companion}, \mref{Req.}{req:constrained-key-gate}
\metasep newly derived
\metasep implemented (effective: accepted)
\metasep \metalabel{Code}\coderef{https://github.com/fraunhofersg/certiqs-key/blob/v0.1.0/src/certiqs_key/meat_core/finitekey/keylength.py}{certiqs_key.meat_core.finitekey.keylength:licensed_class_check}}

% claim meat-core/def:roundmodel
\claimsummary{meat-core/def:roundmodel}{The round model is that of foundations at the squashed system, with round channels that hand every public announcement to the adversary's memory. A licensed attack set bounds each attack's adjoint by the declared constraints; whether a protocol reduces to it is the attack-set licence.}
\claimbookmark{meat-core/def:roundmodel}{Definition}{def:roundmodel}{Round model and licensed attack sets}
\claimleadin{%
The \exttermref{../foundations/contract.pdf}{term-round-model}{round model} is the one the foundations module, which holds the
quantum-information preliminaries of this contract, states for a generic
per-round system, taken at the squashed system $\Qsys$: the system the
declared family $\{F_v\}$ measures, onto which the
\exttermref{../reductions/contract.pdf}{term-squashing-channel}{squashing channels} of the reductions module map
the physical system. The reductions module takes the same model at the
physical system for its transfer theorems and states its
\exttermref{../foundations/contract.pdf}{term-attack-set}{attack sets} at $\Qsys$
(\xref{reductions/def:roundmodel}, \xref{reductions/def:attset}), so the
attack sets of the two modules are the same sets. The definition below adds
what this module's accumulation theorems consume: the round channels, the
retention of the announcements and the licensed list.
}
\begin{definition}[Round model and licensed attack sets]
\label{def:roundmodel}
Fix a declaration $\Decl$ with single-round protocol channel
$\Mprot\in\mathrm{CPTP}(A\Qsys\to Z\Chat I)$ (Theorem~\ref{thm:protocol-channel}).
The round model of this contract is the round model of
\xref{foundations/def:round-model} with per-round system $\Qsys$, Alice's
system $A$ and the declared source mode, the round index of $\Qsys$ and $A$
dropped. An \exttermref{../foundations/contract.pdf}{term-attack-sequence}{attack sequence} is thus a family $(\Aev_j)_{j=1}^{n}$ of
channels whose output carries the systems the adversary supplies in the
declared source mode, the display being that definition's with its
per-round system $S$ taken to be $\Qsys$:
\begin{equation}\label{eq:attackchannel}
  \Aev_j\in
  \begin{cases}
    \mathrm{CPTP}(E_{j-1},\ \Qsys\,E'_j), &
      \text{trusted-source mode},\\[0.3ex]
    \mathrm{CPTP}(E_{j-1},\ A\,\Qsys\,E'_j), &
      \text{untrusted-source mode},
  \end{cases}
\end{equation}
where $E_0$ is Eve's initial register and $E'_j$ her outgoing memory. In
trusted-source mode $A$ is \emph{not} an output of $\Aev_j$: it is
supplied by Alice, inaccessible to the adversary, and carried into the
round channel as a separate input --- that, and nothing else, is what
A-SRC (\xref{transmitter/ass:src}) licenses. The generated
\termdef{term-round-channel}{round channel}s
are
\begin{equation}\label{eq:roundchannel}
\begin{gathered}
  \Mprot_j\coloneqq
  \begin{cases}
    (\Mprot\otimes\idmap_{E'_j})\circ(\idmap_{A}\otimes\Aev_j)
      \;\in\;\mathrm{CPTP}(A\,E_{j-1},\;Z\,\Chat\,E_j), &
      \text{trusted mode},\\[0.5ex]
    (\Mprot\otimes\idmap_{E'_j})\circ\Aev_j
      \;\in\;\mathrm{CPTP}(E_{j-1},\;Z\,\Chat\,E_j), &
      \text{untrusted mode},
  \end{cases}\\
  E_j\coloneqq E'_j\,I_j:
\end{gathered}
\end{equation}
the attack delivers the round input on $\Qsys$ --- and, in
untrusted-source mode, on $A$ as well --- together with Eve's
new memory, the protocol channel produces $(Z_j,\Chat_j,I_j)$, and the
announcement $I_j$ --- classical and public --- is placed in Eve's
register. Content of $E_{j-1}$ that Eve wishes to retain (earlier
announcements included) she carries inside $E'_j$; every physical
attack is represented by a sequence that retains \emph{all}
announcements, and Theorem~\ref{thm:accumulated-entropy} uses that representation, so
that $E_n=I_1^n\Ehat$ with $\Ehat$ Eve's retained quantum memory. The
untrusted-mode round channel has no $A$ input at all, which is what
makes it an instance of Fact~\ref{fact:meat} with the marginal
register $A_{j-1}$ trivial (Fact~\ref{fact:meat}, locator note (vii);
Theorem~\ref{thm:accumulated-entropy}).

For a licensed list $\CList$ (Definition~\ref{def:homog}) and
finite-dimensional registers $E',E''$, the
\termdef{term-licensed-attack-set}{licensed attack set}
is the attack set of $\CList$ of \xref{foundations/def:round-model}
at $S=\Qsys$; its displays at that system are, in trusted-source mode,
\begin{equation}\label{eq:attset}
  \AttSet_{\CList}(E',\,\Qsys E'')\coloneqq
  \Big\{\Aev\in\mathrm{CPTP}(E',\Qsys E''):\
  \Aev^{\dagger}\big(\hat O_k\otimes\idon{E''}\big)
  \bowtie_k c_k\,\idon{E'}\ \ \forall k\Big\},
\end{equation}
and, with the constraint operator lifted to the larger supplied system,
\begin{equation}\label{eq:attset-un}
  \AttSet_{\CList}(E',\,A\Qsys E'')\coloneqq
  \Big\{\Aev\in\mathrm{CPTP}(E',A\Qsys E''):\
  \Aev^{\dagger}\big(\idon{A}\otimes\hat O_k\otimes\idon{E''}\big)
  \bowtie_k c_k\,\idon{E'}\ \ \forall k\Big\}
\end{equation}
in untrusted-source mode; the two are distinguished by the output
system named in the argument. An attack sequence is
$\CList$-admissible (an \exttermref{../foundations/contract.pdf}{term-admissible-attack-sequence}{admissible attack sequence} for $\CList$) if
$\Aev_j$ lies in the set of
\eqref{eq:attset}, resp.\ \eqref{eq:attset-un}, with $E'=E_{j-1}$ and
$E''=E'_j$, for every $j$. Where $\CList$ carries entries declared
\exttermref{../protocol-compiler/contract.pdf}{term-party-local}{party-local} (Definition~\ref{def:homog}), each such entry $(\hat O_k,\bowtie_k,c_k;f)$
contributes to the licensed attack set, in place of the scalar condition of
\eqref{eq:attset} or \eqref{eq:attset-un}, the condition of an entry tagged
with the factor $f$ in \xref{foundations/def:round-model}\,(iii) at
$S=\Qsys$. With $\Qsys_{f^c}$ the product of the measured factors other than
$f$, and $R$ equal to $A$ in untrusted-source mode and trivial in
trusted-source mode, its display at that system is
\begin{equation}\label{eq:attset-pl}
\begin{gathered}
  \big(\idmap_{R'}\otimes\Aev\big)^{\dagger}
  \Big(X\otimes\big(\hat O_k-c_k\,\idon{\Qsys_f}\big)\otimes\idon{E''}\Big)
  \;\bowtie_k\;0\\
  \text{for every finite-dimensional }R'\text{ and every }
  X\in\Pos{R'R\,\Qsys_{f^c}} ,
\end{gathered}
\end{equation}
as an operator on $R'E'$, the tensor factors of the argument written in the
order $R'R\,\Qsys_{f^c}$, $\Qsys_f$, $E''$; untagged entries, those not
declared party-local, keep their conditions, read for their
\exttermref{../protocol-compiler/contract.pdf}{term-lifted-operator}{lifted operators}
(\xref{protocol-compiler/def:declaration}\,(DX6)). The reference system $R'$ is what
carries Alice's retained system in trusted-source mode, where $A$ is no
output of $\Aev$. With $R'$ trivial and $X=\idon{}$ the condition is the
scalar one for the lifted operator, so the set with tagged entries is
contained in the set with the same entries untagged. For
$\CList=\varnothing$ the condition is empty:
$\AttSet_{\varnothing}=\mathrm{CPTP}$, every attack sequence is
admissible, and the model is unrestricted. The adjoint condition in
\eqref{eq:attset} is exactly the published squashed-attack-set format
of \cite[Eq.~(151)]{tupkary2026decoy}, and
restricting attention to a declared subset of channels is the standard
modelling step of \cite[Def.~8.3, Rem.~8.4]{tupkary2026decoy}. Whether
a given physical protocol reduces to this model with a given
$\AttSet_{\CList}$ is \emph{not} decided here: it is what
A-ATK (\xref{protocol-compiler/ass:atk}) certifies.
\end{definition}
\claimmeta{\metalabel{Uses}\mref{Def.}{def:declaration}, \mref{Thm.}{thm:protocol-channel}, \mref{Fact}{fact:meat}, \mref{Def.}{def:homog}, \extref{../foundations/contract.pdf}{foundations/def:round-model}
\metasep \metalabel{Used by}\mref{Lem.}{lem:range}, \mref{Lem.}{lem:product-marginal}, \mref{Thm.}{thm:accumulated-entropy}, \mref{Thm.}{thm:constrained-accumulation}, \mref{Rem.}{rem:tagged-entries-and-accumulation}, \mref{Ex.}{ex:fss}, \mref{Rem.}{rem:homog}, \mref{Rem.}{rem:not-licensed}, \mref{Rem.}{rem:hlow-range}, \mref{Def.}{def:postprocessing-model}, \mref{Thm.}{thm:keylength}, \mref{Cor.}{cor:separate-executions}, \mref{Ass.}{ass:lab}
\metasep newly derived
\metasep implemented (effective: accepted)
\metasep \metalabel{Code}\coderef{https://github.com/fraunhofersg/certiqs-key/blob/v0.1.0/src/certiqs_key/meat_core/declaration/validation.py}{certiqs_key.meat_core.declaration.validation:REQUIRED_ASSUMPTIONS}
\metasep \metalabel{Oracles}\coderef{https://github.com/fraunhofersg/certiqs-key/blob/v0.1.0/tests/oracles/registry.py}{registry.py}}

% claim meat-core/def:hqkd-list
\claimsummary{meat-core/def:hqkd-list}{The constrained single-round rate minimises a KL penalty plus the key's conditional entropy over acceptance distributions and states meeting the declared marginal and constraint list; at the empty list it is the unconstrained rate. Constraints only raise it, so using it needs a licence.}
\claimbookmark{meat-core/def:hqkd-list}{Definition}{def:hqkd-list}{Constrained target set and single-round rate}
\begin{definition}[Constrained target set and single-round rate]
\label{def:hqkd-list}
For a register $\Etil'$ and a (possibly empty) declared list $\CList$,
let
\begin{equation}\label{eq:sigL}
  \Sigma^{\CList}(\Etil')\coloneqq
  \big\{(\Mprot\otimes\idmap_{\Etil'})(\sigma):\
  \sigma\in\Seq{A\Qsys\Etil'},\ \sigma\models\Decl,\
  \sigma_{A\Qsys}\models\CList\big\},
\end{equation}
where $\sigma\models\Decl$ abbreviates the declaration's marginal
condition,
\begin{equation}\label{eq:margcond}
  \sigma\models\Decl
  \quad:\Longleftrightarrow\quad
  \begin{cases}
    \sigma_A=\tau_A, & \text{trusted-source mode},\\
    \text{no condition beyond }\tr\sigma=1,
      & \text{untrusted-source mode},
  \end{cases}
\end{equation}
and $\sigma_{A\Qsys}\models\CList$ abbreviates ``$\sigma_{A\Qsys}$
satisfies every entry of $\CList$'' in the form
Definition~\ref{def:homog} gives it for the declared mode, i.e.\
$\tr_{\Qsys}[(\idon{A}\otimes\hat O_k)\,\sigma_{A\Qsys}]
\bowtie_k\Theta_k$ in trusted-source mode and
$\tr[(\idon{A}\otimes\hat O_k)\,\sigma_{A\Qsys}]\bowtie_k c_k$ in
untrusted-source mode, for all $k$; an entry declared \exttermref{../protocol-compiler/contract.pdf}{term-party-local}{party-local}
is read, in either mode, in the form \eqref{eq:partylocal-entry} of
Definition~\ref{def:homog}, an affine semidefinite condition on
$\sigma_{A\Qsys}$. $\Sigma^{\CList}(\Etil')$ is
convex in either mode: the state set is the intersection of the convex
$\Seq{A\Qsys\Etil'}$ with affine equalities and operator- or
scalar-inequality conditions, each convex, and
$\Mprot\otimes\idmap_{\Etil'}$ is affine.
Define the
\termdef{term-constrained-single-round-rate}{constrained single-round rate}
\begin{equation}\label{eq:hqkd-list}
  \hqkd(\CList)\coloneqq\inf_{\mathbf q\in\acSet}\
  \inf_{\substack{\sigma\in\Seq{A\Qsys\Etil}\\
  \sigma\models\Decl,\ \sigma_{A\Qsys}\models\CList}}
  \left(\frac{\alpha}{\alpha-1}\,
  \sfD\!\left(\mathbf q\,\middle\|\,\Phimap[\sigma_{A\Qsys}]\right)
  +\sum_{\hat c\in\Chkey}
  q(\hat c)\,\Hsandup(Z|I\Etil)_{\nu^{\sigma}_{|\hat c}}\right),
\end{equation}
with $\nu^{\sigma}\coloneqq(\Mprot\otimes\idmap_{\Etil})(\sigma)$ and
$\dim\Hs{\Etil}\ge\dim\Hs{A\Qsys}$; by
Lemma~\ref{lem:purifier-indep} the value does not depend on that
choice. For $\CList=\varnothing$ we write
$\hqkd\coloneqq\hqkd(\varnothing)$, the unconstrained single-round rate,
which Theorem~\ref{thm:accumulated-entropy} writes out as
\eqref{eq:hqkd-def}. Since a nonempty
$\CList$ shrinks the inner feasible set,
$\hqkd(\CList)\ge\hqkd(\varnothing)$ always.
\end{definition}
\claimmeta{\metalabel{Uses}\mref{Def.}{def:homog}
\metasep \metalabel{Used by}\mref{Lem.}{lem:purifier-indep}, \mref{Lem.}{lem:range}, \mref{Thm.}{thm:accumulated-entropy}, \mref{Thm.}{thm:constrained-accumulation}, \mref{Rem.}{rem:tagged-entries-and-accumulation}, \mref{Rem.}{rem:not-licensed}, \mref{Def.}{def:feasible}, \mref{Thm.}{thm:reduction}, \mref{Cor.}{cor:reduction-constrained}
\metasep newly derived
\metasep implemented (effective: accepted)
\metasep \metalabel{Code}\coderef{https://github.com/fraunhofersg/certiqs-key/blob/v0.1.0/src/certiqs_key/meat_core/objective/objective.py}{certiqs_key.meat_core.objective.objective:objective_value}
\metasep \metalabel{Oracles}\coderef{https://github.com/fraunhofersg/certiqs-key/blob/v0.1.0/tests/oracles/registry.py}{registry.py}}

% claim meat-core/def:admissible-acceptance-set
\claimsummary{meat-core/def:admissible-acceptance-set}{An acceptance set is admissible when it is a closed convex subset of the probability simplex on the acceptance values. Accumulation, the attainment and limit results and the first-order bound all read that geometry.}
\claimbookmark{meat-core/def:admissible-acceptance-set}{Definition}{def:admissible-acceptance-set}{Admissible acceptance set}
\claimleadin{%
The acceptance set is the set of acceptance-value distributions on which the
acceptance test accepts. Which sets the engine admits, and how it decides
that a declared one is admissible, is Requirement~\ref{req:acceptance-set-admission}.
}
\begin{definition}[Admissible acceptance set]\label{def:admissible-acceptance-set}
An \termdef{term-admissible-acceptance-set}{admissible acceptance set}
is a closed convex subset $\acSet$ of the probability simplex on $\Chatset$.
\end{definition}
\claimmeta{\metalabel{Used by}\mref{Thm.}{thm:accumulated-entropy}, \mref{Rem.}{rem:homog}, \mref{Thm.}{thm:convexity}, \mref{Thm.}{thm:oracle}, \mref{Cor.}{cor:drop-program}, \mref{Prop.}{prop:minimal-divergence}, \mref{Prop.}{prop:model-pass-bound}, \mref{Req.}{req:acceptance-set-admission}, \mref{Lem.}{lem:fo-khyp}, \mref{Thm.}{thm:firstorder}, \mref{Cor.}{cor:fo-dropped}, \mref{Cor.}{cor:fo-reference}, \mref{Lem.}{lem:certified-bracket}, \mref{Rem.}{rem:acc-side}
\metasep newly derived
\metasep accepted}

% claim meat-core/lem:cond-more
\claimsummary{meat-core/lem:cond-more}{The sandwiched conditional R\'enyi entropy cannot increase when the conditioning system grows, by data processing under the partial trace. The conservative steps that hand the adversary extra information rest on it.}
\claimbookmark{meat-core/lem:cond-more}{Lemma}{lem:cond-more}{Conditioning on more side information}
\begin{technicalclaim}
\begin{lemma}[Conditioning on more side information]\label{lem:cond-more}
For $\alpha\in(1,\infty)$ and any $\rho_{ABC}\in\Seq{ABC}$:
$\Hsandup(A|BC)_{\rho}\le\Hsandup(A|B)_{\rho}$.
\end{lemma}
\end{technicalclaim}
\claimmeta{\metalabel{Used by}\mref{Lem.}{lem:purifier-indep}, \mref{Thm.}{thm:reduction}, \mref{Prop.}{prop:value-range}
\metasep newly derived
\metasep accepted}

\begin{proof}
For any $\sigma_{BC}\in\Seq{BC}$, data processing for the sandwiched
divergence under the channel $\tr_C$
(\cite[Thm.~4.1, Eq.~(4.82)]{tomamichel2016finite}; $\alpha>1\ge\frac12$)
gives
$\Dsand(\rho_{ABC}\|\idon{A}\otimes\sigma_{BC})
\ge\Dsand(\rho_{AB}\|\idon{A}\otimes\sigma_{B})$. Taking $-\inf$ over
$\sigma_{BC}$ on the left and noting that every $\sigma_B$ arises as a
marginal on the right,
$\Hsandup(A|BC)_{\rho}
=-\inf_{\sigma_{BC}}\Dsand(\rho_{ABC}\|\idon{A}\otimes\sigma_{BC})
\le-\inf_{\sigma_{B}}\Dsand(\rho_{AB}\|\idon{A}\otimes\sigma_{B})
=\Hsandup(A|B)_{\rho}$.
\end{proof}

% claim meat-core/lem:purifier-indep
\claimsummary{meat-core/lem:purifier-indep}{Once the adversary's register is at least as large as Alice's and the measured system together, enlarging it leaves the constrained single-round rate unchanged. The inner infimum may be restricted to pure states, which the branch formula uses.}
\claimbookmark{meat-core/lem:purifier-indep}{Lemma}{lem:purifier-indep}{Independence of the purifier size}
\begin{lemma}[Independence of the purifier size]\label{lem:purifier-indep}
Fix $\CList$ (possibly empty) and write $h(\Etil')$ for the value of
\eqref{eq:hqkd-list} computed with a register $\Etil'$ of
$\dim\Hs{\Etil'}\ge\dim\Hs{A\Qsys}$ in place of $\Etil$. Then
$h(\Etil')$ is the same for every such $\Etil'$, and the inner infimum
is unchanged if restricted to pure $\sigma$.
\end{lemma}
\claimmeta{\metalabel{Uses}\mref{Def.}{def:admissible-declaration}, \mref{Def.}{def:hqkd-list}, \mref{Lem.}{lem:cond-more}, \extref{../foundations/contract.pdf}{foundations/def:renyi}
\metasep \metalabel{Used by}\mref{Thm.}{thm:accumulated-entropy}, \mref{Thm.}{thm:constrained-accumulation}, \mref{Rem.}{rem:not-licensed}
\metasep newly derived
\metasep accepted}

\begin{proof}
Let $\dim\Hs{\Etil_2}\ge\dim\Hs{\Etil_1}\ge\dim\Hs{A\Qsys}$.

\emph{($h(\Etil_2)\le h(\Etil_1)$.)} An isometry
$V:\Hs{\Etil_1}\to\Hs{\Etil_2}$ maps every feasible $\sigma$ for
$\Etil_1$ to the feasible
$(\idon{A\Qsys}\otimes V)\sigma(\idon{A\Qsys}\otimes V)^{\dagger}$ for
$\Etil_2$ with the same $A\Qsys$-marginal, so the divergence term and
the constraints are unchanged; and $\Hsandup(Z|I\Etil)$ is invariant
under an isometry on the conditioning register: data processing
(\cite[Thm.~4.1, Eq.~(4.82)]{tomamichel2016finite}) under the
isometric channel $V(\cdot)V^{\dagger}$ on $\Etil_1$ and under a
reversal channel $\mathcal{R}$ with $\mathcal{R}(V(\cdot)V^{\dagger})=\idmap$
(e.g.\ $\mathcal{R}(X)=V^{\dagger}XV+\tr[(\idon{}-VV^{\dagger})X]\,\xi$ for a
fixed state $\xi$), applied to the optimisation defining $\Hsandup$ in
\xref{foundations/def:renyi}, gives both inequality directions.
Hence every $\Etil_1$-bracket value is attained
by an $\Etil_2$-feasible state.

\emph{($h(\Etil_2)\ge h(\Etil_1)$; purifications suffice.)} Let
$\sigma$ be feasible for $\Etil_2$ and let
$\hat\sigma\in\Seq{A\Qsys\Etil_1}$ be a purification of
$\sigma_{A\Qsys}$ (it exists: $\dim\Hs{\Etil_1}\ge\dim\Hs{A\Qsys}$).
Every extension of $\sigma_{A\Qsys}$ arises from any fixed
purification by a channel on the purifying register: purify $\sigma$
itself by a further register, relate the two purifications of
$\sigma_{A\Qsys}$ by an isometry on the purifier side, and trace ---
so $\sigma=(\idmap_{A\Qsys}\otimes N)(\hat\sigma)$ for some channel
$N:\Etil_1\to\Etil_2$. Then
$\nu^{\sigma}=(\idmap_{ZI}\otimes N)(\nu^{\hat\sigma})$, and data
processing on the conditioning register
(\cite[Thm.~4.1, Eq.~(4.82)]{tomamichel2016finite}; the general-channel
form of the argument in Lemma~\ref{lem:cond-more}) gives
$\Hsandup(Z|I\Etil_2)_{\nu^{\sigma}_{|\hat c}}
\ge\Hsandup(Z|I\Etil_1)_{\nu^{\hat\sigma}_{|\hat c}}$ for every
$\hat c$ (conditioning is on a classical value of $\Chat$, which
commutes with the channel on $\Etil$). The divergence term and the
constraints agree (same $A\Qsys$-marginal). Hence the
$\Etil_2$-bracket at $\sigma$ dominates the $\Etil_1$-bracket at the
pure $\hat\sigma$, so $h(\Etil_2)\ge h(\Etil_1)$, and the infimum is
attained along purifications. Combining the two directions, all values
coincide.
\end{proof}

% claim meat-core/lem:range
\claimsummary{meat-core/lem:range}{An admissible attack carries the licensed bounds on its adjoint over to the constraint conditions on the round's state, so every round output lies in the constrained target set. A larger output set can only lower the rate, which is the conservative comparison accumulation needs.}
\claimbookmark{meat-core/lem:range}{Lemma}{lem:range}{Constraint transport and range inclusion}
\begin{lemma}[Constraint transport and range inclusion]\label{lem:range}
Let $\CList$ be licensed (Definition~\ref{def:homog}) or empty, and let
$\Aev_j$ be $\CList$-admissible for the declared source mode
(Definition~\ref{def:roundmodel}).
\begin{enumerate}[label=(\roman*),leftmargin=2.2em]
\item \emph{(Transport.)} \emph{Trusted-source mode:} for every
$\omega\in\Seq{A\,E_{j-1}\,\Etil}$ with $\omega_A=\tau_A$, the state
$\varsigma\coloneqq(\idmap_A\otimes\Aev_j\otimes\idmap_{\Etil})(\omega)
\in\Seq{A\Qsys E'_j\Etil}$ satisfies $\varsigma_A=\tau_A$ and
$\varsigma_{A\Qsys}\models\CList$.
\emph{Untrusted-source mode:} for every
$\omega\in\Seq{E_{j-1}\,\Etil}$, the state
$\varsigma\coloneqq(\Aev_j\otimes\idmap_{\Etil})(\omega)
\in\Seq{A\Qsys E'_j\Etil}$ satisfies $\varsigma\models\Decl$ --- which
in that mode is only $\tr\varsigma=1$ --- and
$\varsigma_{A\Qsys}\models\CList$ in the scalar form of
Definition~\ref{def:homog}\,(b). \emph{In either mode}, every entry of
$\CList$ declared \exttermref{../protocol-compiler/contract.pdf}{term-party-local}{party-local} holds of $\varsigma_{A\Qsys}$ in
the form \eqref{eq:partylocal-entry}, the attack being admissible for it in
the sense of \eqref{eq:attset-pl}.
\item \emph{(Range inclusion.)} With $\Mprot_j$ the round channel
\eqref{eq:roundchannel} and $\Sigma_j$ its output set
\eqref{eq:meat-range},
$\Sigma_j\subseteq\Sigma^{\CList}(E'_j\Etil)$ up to the regrouping of the
registers of the output, $(Z\Chat)(E'_jI_j)\Etil\mapsto(Z\Chat I)(E'_j\Etil)$;
hence, embedding the finitely many $E'_j\Etil$ into a common register
$\Etil^{*}$ with $\dim\Hs{\Etil^{*}}=\max_j\dim\Hs{E'_j\Etil}$,
\begin{equation*}
  \Sigma=\operatorname{conv}\Big(\bigcup_{j}\Sigma_j\Big)
  \;\subseteq\;\Sigma^{\CList}(\Etil^{*}).
\end{equation*}
\item \emph{(Antitonicity.)} If $\Sigma\subseteq\Sigma'$ then the
value of \eqref{eq:meat-rate} computed over $\Sigma$ is at least its
value computed over $\Sigma'$.
\end{enumerate}
\end{lemma}
\claimmeta{\metalabel{Uses}\mref{Fact}{fact:meat}, \mref{Def.}{def:homog}, \mref{Def.}{def:roundmodel}, \mref{Def.}{def:hqkd-list}
\metasep \metalabel{Assumes}\extaref{../protocol-compiler/contract.pdf}{protocol-compiler/ass:atk}{A-ATK}, \extaref{../transmitter/contract.pdf}{transmitter/ass:src}{A-SRC}
\metasep \metalabel{Used by}\mref{Thm.}{thm:accumulated-entropy}, \mref{Thm.}{thm:constrained-accumulation}, \mref{Rem.}{rem:tagged-entries-and-accumulation}, \mref{Rem.}{rem:homog}, \mref{Rem.}{rem:scalar-declaration}, \mref{Rem.}{rem:hlow-range}
\metasep newly derived
\metasep accepted}

\begin{proof}
(i) $\varsigma_A=\omega_A=\tau_A$:
$\idmap_A\otimes\Aev_j\otimes\idmap_{\Etil}$ is trace-preserving and
acts trivially on $A$. For the constraints, test against an arbitrary
$X\in\Lin{A}$:
\begin{equation*}
  \tr\big[X\,\tr_{\Qsys}[(\idon{A}\otimes\hat O_k)\,
  \varsigma_{A\Qsys}]\big]
  =\tr\big[(X\otimes\hat O_k\otimes\idon{E'_j})\,
  (\idmap_A\otimes\Aev_j)(\omega_{AE_{j-1}})\big]
  =\tr\big[\big(X\otimes
  \Aev_j^{\dagger}(\hat O_k\otimes\idon{E'_j})\big)\,
  \omega_{AE_{j-1}}\big]
\end{equation*}
by definition of the adjoint. Hence, as operators on $A$,
\begin{equation}\label{eq:transport}
  \tr_{\Qsys}\big[(\idon{A}\otimes\hat O_k)\,\varsigma_{A\Qsys}\big]
  =\tr_{E_{j-1}}\big[\big(\idon{A}\otimes
  \Aev_j^{\dagger}(\hat O_k\otimes\idon{E'_j})\big)\,
  \omega_{AE_{j-1}}\big].
\end{equation}
Write $X_k\coloneqq\Aev_j^{\dagger}(\hat O_k\otimes\idon{E'_j})
-c_k\idon{E_{j-1}}$. For $\bowtie_k\,=\,\succeq$, admissibility
\eqref{eq:attset} says $X_k\succeq0$, and
$\tr_{E_{j-1}}[(\idon{A}\otimes X_k)\,\omega]
=\tr_{E_{j-1}}[(\idon{A}\otimes X_k^{1/2})\,\omega\,
(\idon{A}\otimes X_k^{1/2})]\succeq0$ --- cyclicity under the partial
trace for operators on the traced factor, and the middle expression is
the partial trace of a PSD operator --- so \eqref{eq:transport} equals
$c_k\,\omega_A+\tr_{E_{j-1}}[(\idon{A}\otimes X_k)\,\omega]
\succeq c_k\,\tau_A$. For $\bowtie_k\,=\,\preceq$, apply the same
argument to $-X_k\succeq0$; for $\bowtie_k\,=\,=$, admissibility
forces $X_k=0$, and \eqref{eq:transport} equals $c_k\,\tau_A$ exactly.
(For $\CList=\varnothing$ the claim is only $\varsigma_A=\tau_A$.)

\emph{(i), untrusted-source mode.} Here $A$ is an \emph{output} of
$\Aev_j$, so the computation above has no operator-valued analogue:
there is no $A$ factor on the input side to carry an operator
inequality through, and the identity \eqref{eq:transport} --- whose
left-hand side is an operator on $A$ and whose right-hand side is built
from $\omega_{AE_{j-1}}$ --- has no counterpart. What does survive is
its trace. Normalization is immediate: $\Aev_j\otimes\idmap_{\Etil}$ is
trace-preserving, so $\tr\varsigma=\tr\omega=1$. For the constraints,
by definition of the adjoint,
\begin{equation}\label{eq:transport-un}
  \tr\big[(\idon{A}\otimes\hat O_k)\,\varsigma_{A\Qsys}\big]
  =\tr\big[(\idon{A}\otimes\hat O_k\otimes\idon{E'_j})\,
  \Aev_j(\omega_{E_{j-1}})\big]
  =\tr\big[\Aev_j^{\dagger}
  \big(\idon{A}\otimes\hat O_k\otimes\idon{E'_j}\big)\,
  \omega_{E_{j-1}}\big].
\end{equation}
Write $X_k\coloneqq
\Aev_j^{\dagger}(\idon{A}\otimes\hat O_k\otimes\idon{E'_j})
-c_k\idon{E_{j-1}}$. For $\bowtie_k\,=\,\succeq$ admissibility
\eqref{eq:attset-un} says $X_k\succeq0$, so
$\tr[X_k\,\omega_{E_{j-1}}]\ge0$ (the trace of a product of two PSD
operators is non-negative), and \eqref{eq:transport-un} equals
$c_k\tr[\omega_{E_{j-1}}]+\tr[X_k\,\omega_{E_{j-1}}]\ge c_k$, using
$\tr\omega_{E_{j-1}}=1$. For $\bowtie_k\,=\,\preceq$ apply the same
argument to $-X_k\succeq0$; for $\bowtie_k\,=\,=$ admissibility forces
$X_k=0$ and \eqref{eq:transport-un} equals $c_k$ exactly. This is
precisely the scalar entry format of
Definition~\ref{def:homog}\,(b), and it is the \emph{only} conclusion
the licence supports in this mode; the operator form
$\Theta_k=c_k\idon{A}$ is not available (Remark~\ref{rem:homog}\,(a)).
(For $\CList=\varnothing$ the claim is only $\tr\varsigma=1$.)

\emph{(i), party-local entries.} Let $(\hat O_k,\bowtie_k,c_k;f)$ be
declared party-local, write
$\hat D\coloneqq\hat O_k-c_k\idon{\Qsys_f}$ and
$Y\coloneqq\tr_{\Qsys_f}[(\idon{}\otimes\hat D)\,\varsigma_{A\Qsys}]$, a
Hermitian operator on $A\Qsys_{f^c}$, and let $X\in\Pos{A\Qsys_{f^c}}$. In
trusted-source mode, by the definition of the adjoint,
\begin{equation*}
  \tr[XY]=\tr\big[(X\otimes\hat D\otimes\idon{E'_j})\,
  (\idmap_A\otimes\Aev_j)(\omega_{AE_{j-1}})\big]
  =\tr\Big[(\idmap_A\otimes\Aev_j)^{\dagger}
  \big(X\otimes\hat D\otimes\idon{E'_j}\big)\;\omega_{AE_{j-1}}\Big],
\end{equation*}
and \eqref{eq:attset-pl} with $R'=A$ and $R$ trivial says that the operator
in the last trace is $\bowtie_k0$ on $AE_{j-1}$; since
$\omega_{AE_{j-1}}\succeq0$ the trace is $\ge0$ for $\succeq$, $\le0$ for
$\preceq$ and $0$ for $=$. In untrusted-source mode the same computation
with $\varsigma_{A\Qsys E'_j}=\Aev_j(\omega_{E_{j-1}})$ gives
$\tr[XY]=\tr[\Aev_j^{\dagger}(X\otimes\hat D\otimes\idon{E'_j})\,
\omega_{E_{j-1}}]$, and \eqref{eq:attset-pl} with $R'$ trivial and $R=A$
gives the same signs. A Hermitian $Y$ with $\tr[XY]\ge0$ for every
$X\succeq0$ is positive semidefinite, and one with $\tr[XY]=0$ for every
$X\succeq0$ vanishes, positive semidefinite operators spanning all
operators; so $Y\bowtie_k0$, which is \eqref{eq:partylocal-entry}. The
tensor factors are written in the order of \eqref{eq:attset-pl}; reordering
them is a fixed unitary relabelling on both sides.

(ii) \emph{Trusted-source mode.} Fix
$\omega\in\Seq{A\,E_{j-1}\,\Etil}$ with $\omega_A=\tau_A$ and
set
$\varsigma=(\idmap_A\otimes\Aev_j\otimes\idmap_{\Etil})(\omega)$. By
\eqref{eq:roundchannel},
$(\Mprot_j\otimes\idmap_{\Etil})(\omega)
=(\Mprot\otimes\idmap_{E'_j\Etil})(\varsigma)$ --- the same operator,
with the registers of the round-$j$ output regrouped: the announcement $I_j$,
emitted by $\Mprot$, is assigned to $E_j=E'_jI_j$ on the left and to
the $I$-output on the right. By (i), $\varsigma$ is feasible for
\eqref{eq:sigL} with purifier $E'_j\Etil$, so the output lies in
$\Sigma^{\CList}(E'_j\Etil)$. Isometric embedding of each $E'_j\Etil$
into the common $\Etil^{*}$ preserves feasibility and the
$A\Qsys$-marginal, so every embedded $\Sigma_j$ lies in the convex set
$\Sigma^{\CList}(\Etil^{*})$ (Definition~\ref{def:hqkd-list}), and so
does the convex hull. \emph{Untrusted-source mode.} The same argument
with $\omega\in\Seq{E_{j-1}\Etil}$ and
$\varsigma=(\Aev_j\otimes\idmap_{\Etil})(\omega)$: the second line of
\eqref{eq:roundchannel} gives the same regrouping identity, (i) gives
feasibility for \eqref{eq:sigL} in the untrusted reading of
\eqref{eq:margcond}, and the embedding step is unchanged. Note that
for $\CList=\varnothing$ the inclusion is then immediate for a
different reason as well: $\Sigma^{\varnothing}(\Etil^{*})$ is the set
of \emph{all} $(\Mprot\otimes\idmap)(\sigma)$ with $\sigma$ normalized,
and every round-channel output is of that form.

(iii) An infimum over a superset is no larger; the bracket of
\eqref{eq:meat-rate} is evaluated pointwise.
\end{proof}

% claim meat-core/lem:product-marginal
\claimsummary{meat-core/lem:product-marginal}{In trusted-source mode, under retention, emission of the specified ensemble, sequential rounds, private source randomness and a sealed laboratory, the n rounds start with Alice's systems in the declared marginal in product, and each round finds her system in that marginal.}
\claimbookmark{meat-core/lem:product-marginal}{Lemma}{lem:product-marginal}{Product marginal of Alice's retained systems}
\claimleadin{%
In \exttermref{../transmitter/contract.pdf}{term-trusted-source-mode}{trusted-source mode} the \exttermref{../foundations/contract.pdf}{term-round-model}{round model} takes Alice's retained
systems as inputs of the round channels, and the accumulation theorem
(Theorem~\ref{thm:accumulated-entropy}) asks that they start in the declared
marginal, in product over the rounds. The lemma derives that from the
transmitter's premises. The \exttermref{../transmitter/contract.pdf}{term-source-replacement-certificate}{source-replacement certificate} describes
one round's preparation as Alice's measurement of her half of the certificate
(\xref{transmitter/thm:certificate-sufficiency}\,(iii)); A-ENS
(\xref{transmitter/ass:ens}) makes the $n$ preparations a product, and
A-SEED (Assumption~\ref{ass:seed}\,(v)) makes them independent of what the
adversary holds at the start; and retention, A-SRC
(\xref{transmitter/ass:src}), the sealed laboratory of the party that holds
$A$, A-LAB (Assumption~\ref{ass:lab}\,(i), (ii)), and the sequential round structure,
A-SEQ (\xref{transmitter/ass:seq}), let each measurement be moved into its
round. Retention alone does not give the product: it fixes what may
act on each round's system, and the independence of the source's rounds gives
the product. A specification declared by a \exttermref{../transmitter/contract.pdf}{term-source-model}{source model} carries, in
the place of A-ENS, the premise its submodule declares, which refines A-ENS:
for weak coherent pulses it is A-WCP (\xref{transmitter/decoy/ass:wcp}), whose
last three clauses are A-ENS's for the model's states. At $d_A=1$ the conclusion holds of every state
(Proposition~\ref{prop:trivial-alice-system}\,(ii)).
}
\begin{lemma}[Product marginal of Alice's retained systems]\label{lem:product-marginal}
Let the declaration be in \exttermref{../transmitter/contract.pdf}{term-trusted-source-mode}{trusted-source mode} with $d_A\ge2$, and let
its $A$, $\tau_A$ and $\{M_x\}$ be those of a \exttermref{../transmitter/contract.pdf}{term-source-specification}{source specification}
(\xref{transmitter/def:source-spec}) that carries a source-replacement
certificate $\Psi$ for $(A,\mathbb{T},\tau_A,\{M_x\})$ against its
\exttermref{../transmitter/contract.pdf}{term-virtual-ensemble}{virtual ensemble} $\{p_x,\rvirt_x\}$ on $\Avirt$, in any of the forms of
\xref{transmitter/def:srcert}, an \exttermref{../transmitter/contract.pdf}{term-ensemble-presentation}{ensemble presentation} being read as
the block-form certificate it defines; let its \exttermref{../transmitter/contract.pdf}{term-source-segment}{source segment} satisfy
the source identity $\Xi(\rvirt_x)=\rho_x$ for every $x\in\Xset$, with $\Xi$
the composite channel of the segment and $\Xi=\idmap$ when the segment is
empty (\xref{transmitter/def:source-segment}\,(S3)). Let A-SRC
(\xref{transmitter/ass:src}), A-ENS (\xref{transmitter/ass:ens}), A-SEQ
(\xref{transmitter/ass:seq}), A-SEED (Assumption~\ref{ass:seed})
and A-LAB (Assumption~\ref{ass:lab}) at the party to which the
\exttermref{../protocol-compiler/contract.pdf}{term-party-partition}{party partition} assigns $A$ hold of the physical protocol, let the
adversary's initial register $E$ be finite-dimensional, and let $\sigma_E$
be the state in which it starts. Put
\begin{equation}\label{eq:product-input}
  \omega^{0}\coloneqq\Big(\bigotimes_{j=1}^{n}\Psi_{A_j\Avirt_j}\Big)
  \otimes\sigma_E\ \in\ \Seq{A_1^nE_0},
  \qquad E_0\coloneqq(\Avirt)_1^n\,E .
\end{equation}
Then:
\begin{enumerate}[label=(\roman*),leftmargin=2.2em]
\item \emph{(source replacement over the $n$ rounds)} against every attack,
the physical protocol and the protocol that starts from $\omega^{0}$, applies
$\Xi$ to $\Avirt_j$ in round $j$ in the place of the source's emission, and
applies $\pcNmsm$, Alice's measurement as the channel of
\xref{transmitter/thm:certificate-sufficiency}\,(iii), to $A_j$ at the point of
round $j$ at which the declared protocol's measurements are made, produce
the same state after every round, on every register that both then hold; in
the second protocol $A_j$ is an input of that measurement and of no other
operation, which is how the round model of Definition~\ref{def:roundmodel}
takes Alice's system in trusted-source mode;
\item $\omega^{0}_{A_1^n}=\tau_A^{\otimes n}$;
\item for every \exttermref{../foundations/contract.pdf}{term-attack-sequence}{attack sequence} of Definition~\ref{def:roundmodel}
and every $j$, the state on which the round channel $\Mprot_j$ acts, in a run
of the round model from $\omega^{0}$, has the marginal $\tau_A$ on $A_j$.
\end{enumerate}
At $d_A=1$, where $\tau_A=[1]$ and no certificate is declared, (ii) and (iii)
hold of every state on $A_1^nE_0$.
\end{lemma}
\claimmeta{\metalabel{Uses}\mref{Def.}{def:roundmodel}, \extref{../transmitter/contract.pdf}{transmitter/def:source-segment}, \extref{../transmitter/contract.pdf}{transmitter/def:source-spec}, \extref{../transmitter/contract.pdf}{transmitter/def:srcert}, \extref{../transmitter/contract.pdf}{transmitter/thm:certificate-sufficiency}
\metasep \metalabel{Assumes}\extaref{../transmitter/contract.pdf}{transmitter/ass:src}{A-SRC}, \extaref{../transmitter/contract.pdf}{transmitter/ass:ens}{A-ENS}, \extaref{../transmitter/contract.pdf}{transmitter/ass:seq}{A-SEQ}, \aref{A-SEED}{ass:seed}, \aref{A-LAB}{ass:lab}
\metasep \metalabel{Used by}\mref{Thm.}{thm:accumulated-entropy}
\metasep newly derived
\metasep accepted}

\begin{proof}
\emph{Idea.} The physical preparation is a product over the rounds of one
round's classical-quantum source state, which the certificate writes as
Alice's measurement of her half of $\Psi$ followed by $\Xi$; nothing touches
her half or her letter before the round's measurements, so the interchange
law moves her measurement there; and the marginal of a product of copies of
$\Psi$ on Alice's side is the product of their marginals.

\emph{At $d_A=1$.} $\Seq{A}=\{[1]\}$, so $\Seq{A_1^n}$ is a singleton and
every state on $A_1^nE_0$, and every state the round channels act on, has the
marginal $[1]^{\otimes n}=\tau_A^{\otimes n}$, respectively $[1]=\tau_A$.

\emph{Step 1 (the physical source state).} Let $X_j$ hold the letter of
round $j$, and write
$\rho_{X\Asig}\coloneqq\sum_{x\in\Xset}p_x\dyad{x}_X\otimes\rho_x$ for one
round's classical-quantum source state on the \exttermref{../transmitter/contract.pdf}{term-transmitted-system}{transmitted system}. By
A-ENS (i) and (ii) the letters of the $n$ rounds are independent with the
weights $p_x$ and, given them, the emitted systems are in
$\rho_{x_1}\otimes\dots\otimes\rho_{x_n}$, so the letters and the emitted
systems are jointly in
$\sum_{x_1^n}\prod_jp_{x_j}\dyad{x_1^n}\otimes\rho_{x_1}\otimes\dots\otimes
\rho_{x_n}=\rho_{X\Asig}^{\otimes n}$. By A-ENS (iii) neither the letter of
round $j$ nor the state emitted in it depends on anything that happens in the
other rounds, so preparing all $n$ letters and systems at the start in that
state, and releasing the round-$j$ system only in round $j$, changes nothing
the protocol produces. By A-SEED (v) the letters and the prepared systems
are jointly independent of the adversary's initial register, and by its last
clause no purification of a prepared system is made available to the
adversary during or after the protocol, so that no register beyond the
letters and the emitted systems carries the preparation; the protocol
starts from $\rho_{X\Asig}^{\otimes n}\otimes\sigma_E$.

\emph{Step 2 (the certificate).} By the source identity,
$\rho_{X\Asig}=(\idmap_X\otimes\Xi)(\rho_{X\Avirt})$ with
$\rho_{X\Avirt}\coloneqq\sum_xp_x\dyad{x}_X\otimes\rvirt_x$. By (SR2), in the
form of \xref{transmitter/thm:certificate-sufficiency}\,(iii)(b) with the
identity channel in the place of $\mathcal{E}$,
$(\pcNmsm\otimes\idmap_{\Avirt})(\Psi)=\rho_{X\Avirt}$. Hence
$\rho_{X\Asig}=(\pcNmsm\otimes\Xi)(\Psi)$, and the initial state of Step~1 is
$\big(\bigotimes_j(\pcNmsm\otimes\Xi)_{A_j\Avirt_j}\big)(\omega^{0})$: the
physical protocol is the protocol that starts from $\omega^{0}$ and applies,
at the start, $\pcNmsm$ to every $A_j$ and $\Xi$ to every $\Avirt_j$, the
round-$j$ output of $\Xi$ being released in round $j$. Applying $\Xi$ to
$\Avirt_j$ at the start or in round $j$ produces the same state, since
nothing else acts on $\Avirt_j$ before round $j$.

\emph{Step 3 (moving the measurement into the round).} Fix $j$. Let
$\mathcal{E}_j$ be the composite of every operation of the protocol, of the
parties and of the adversary, from the start to the point of round $j$ at
which the declared protocol's measurements are made. It acts on neither $A_j$
nor $X_j$; the measurements of the other rounds, wherever they stand, act on
$A_k$ and $X_k$ with $k\ne j$. On $A_j$: by A-SRC (ii) no operation of the
protocol or of the adversary acts on $A$ except the declared measurement on it.
On $X_j$, for the adversary's operations: the letters are records of the
party to which the party partition assigns $A$, and by A-LAB (i) at that party
no information about them leaves its laboratory except through the systems
handed out of it, which in trusted-source mode are the systems its source
emits, and the public messages of the protocol, and by A-LAB (ii) the
adversary acts on its laboratory only through the systems delivered into it,
as its declared measurements receive them, and the messages the protocol
delivers to it, so that nothing she does reads or writes the letter record;
no operation of the adversary acts on $X_j$
itself: what the emitted systems carry about the
letter is already in the state of Step~2, and the messages are treated next. On $X_j$,
for the protocol's own operations: the letter of round $j$ is a record of
that same party, which the declared protocol
reads only in the round-$j$ protocol map, at or after the round's
measurements, and in the post-processing; by A-SEQ (i) the messages of round
$j$ are sent after every measurement of the round, by A-SEQ (iv) no operation
on a quantum system of round $j$ depends on them, and by A-SEQ (ii) the
post-processing follows all $n$ rounds. So $\pcNmsm$, a map out of $A_j$
alone into $X_j$, and $\mathcal{E}_j$ act on complementary registers, and by
the interchange law of \xref{transmitter/thm:certificate-sufficiency}\,(iii)(a)
$\mathcal{E}_j\circ(\pcNmsm)_{A_j}=(\pcNmsm)_{A_j}\circ\mathcal{E}_j$ as maps on
the whole input. Applying this for $j=1,\dots,n$ moves each measurement from
the start to its round without changing any state the protocol produces after
the round, which is~(i); in the protocol so obtained $A_j$ enters no operation
but the round-$j$ measurement.

\emph{Step 4 (the marginals).} (ii): $\sigma_E$ is a state, so
$\omega^{0}_{A_1^n}=\bigotimes_j\tr_{\Avirt_j}\Psi_{A_j\Avirt_j}
=\tau_A^{\otimes n}$ by (SR1). (iii): in a run of the round model from
$\omega^{0}$, the state on which $\Mprot_j$ acts is produced from $\omega^{0}$
by $\Mprot_1,\dots,\Mprot_{j-1}$ and the attack channel $\Aev_j$
(Definition~\ref{def:roundmodel}, \eqref{eq:roundchannel}), and in
trusted-source mode each $\Mprot_k$ acts on $A_kE_{k-1}$ and $\Aev_j$ on
$E_{j-1}$, none of them on $A_j$. A channel on a register complementary to
$A_j$ leaves the $A_j$-marginal unchanged, the partial trace over that
register being invariant under it, so the marginal is $\omega^{0}_{A_j}$, which
is $\tau_A$ by~(ii).
\end{proof}

% claim meat-core/thm:accumulated-entropy
\claimsummary{meat-core/thm:accumulated-entropy}{On acceptance, the raw key's R\'enyi entropy given the adversary's information is at least n times the single-round rate less \ensuremath{\alpha}/(\ensuremath{\alpha}\ensuremath{-}1)\ensuremath{\cdot}log\ensuremath{_{2}}(1/Pr[accept]), for any adaptive attack with a finite-dimensional memory. Trusted-source mode reads the product marginal of \hyperlink{claim-meat-core/lem:product-marginal}{Lemma~\ref*{lem:product-marginal}}; untrusted mode needs none.}
\claimbookmark{meat-core/thm:accumulated-entropy}{Theorem}{thm:accumulated-entropy}{Accumulated entropy bound under adaptive attacks}
\claimleadin{%
The imported accumulation theorem, applied to the round channels the
declaration generates, bounds the entropy of the accepted raw key from below
by $n$ times the single-round rate \eqref{eq:hqkd-def}, the rate of
Definition~\ref{def:hqkd-list} at the empty constraint list, less a penalty in
the acceptance probability. The theorem is a statement about attack
sequences of the round model, whose registers are finite-dimensional. That
every physical attack with finite-dimensional registers whose first channel
receives every transmitted signal is such a sequence, with $E_n$ the
adversary's whole information about the rounds, is the mapping of the
physical protocol onto the model; it rests on the \exttermref{../reductions/contract.pdf}{term-reduction-chain}{reduction chain} from which
the declaration was compiled, hypothesis (H-CHAIN) of
Theorem~\ref{thm:keylength}, and, for the round structure, on A-SEQ
(\xref{transmitter/ass:seq})\,(i), (ii) and~(iv), Assumption~\ref{ass:lab}
at each of the two parties and Assumption~\ref{ass:conform}\,(ii), and is made
where the key length is stated (Theorem~\ref{thm:keylength}). The security
of the key extends from these attacks to every physical attack, whatever the
dimension of the adversary's memory, by Corollary~\ref{cor:composed-security}
through Lemma~\ref{lem:memory-lift}.
Remark~\ref{rem:accumulation-scope} says which attacks the bound covers and
what untrusted-source mode costs.
}
\begin{mainclaim}
\begin{theorem}[Accumulated entropy bound under adaptive attacks]\label{thm:accumulated-entropy}
Let $\Decl$ satisfy Definition~\ref{def:admissible-declaration}, let
$\acSet$ satisfy Definition~\ref{def:admissible-acceptance-set}, and let
Assumption~A-RHOM (\xref{protocol-compiler/ass:rhom}) hold; in
trusted-source mode let A-SRC
(\xref{transmitter/ass:src}), A-ENS (\xref{transmitter/ass:ens}), A-SEQ
(\xref{transmitter/ass:seq}), \ref{ass:seed} (A-SEED) and
\ref{ass:lab} (A-LAB) at the party to which the \exttermref{../protocol-compiler/contract.pdf}{term-party-partition}{party partition}
assigns $A$ hold as well, and, at $d_A\ge2$, let the declaration's $A$,
$\tau_A$ and $\{M_x\}$ be those of a \exttermref{../transmitter/contract.pdf}{term-source-specification}{source specification} carrying a
\exttermref{../transmitter/contract.pdf}{term-source-replacement-certificate}{source-replacement certificate} against its \exttermref{../transmitter/contract.pdf}{term-virtual-ensemble}{virtual ensemble},
with a \exttermref{../transmitter/contract.pdf}{term-source-segment}{source segment} satisfying the source identity, as
Lemma~\ref{lem:product-marginal} states them; together with these, the
adversary's initial register being finite-dimensional as every register of
the round model is, the five assumptions are the hypotheses of
Lemma~\ref{lem:product-marginal}, so that by that lemma the
protocol's $n$-round input in that mode has the marginal below (in
untrusted-source mode the lemma is not invoked and none of the hypotheses
named for that mode is used here). Let $(\Aev_j)_{j=1}^{n}$ be an
\emph{arbitrary} attack sequence (Definition~\ref{def:roundmodel}; no
admissibility restriction; its registers $E_0$ and $E'_j$ finite-dimensional,
of any dimensions, as every register of that model is), let
\begin{equation*}
  \rho_{Z_1^n\Chat_1^nI_1^n\Ehat}
  =\big(\Mprot_n\circ\dots\circ\Mprot_1\big)\big[\omega^{0}\big],
  \qquad
  \omega^{0}\in\Seq{A_1^nE_0}\ \text{with}\
  \omega^{0}_{A_1^n}=\tau_A^{\otimes n}
  \ \ \text{(trusted mode)},
\end{equation*}
respectively $\omega^{0}\in\Seq{E_0}$ arbitrary (untrusted mode, where
the round channels take no $A$ input),
with $\Mprot_j$ the round channels \eqref{eq:roundchannel} and Eve's
final register written $E_n=I_1^n\Ehat$
(Definition~\ref{def:roundmodel}), and let $\OmAT$, the acceptance event, be
the event that the frequency vector $\operatorname{freq}_{\hat c_1^n}$ of the
acceptance values of the $n$ rounds lies in $\acSet$, with $\Pr[\OmAT]>0$.
Then for
every $\alpha\in(1,\infty)$,
\begin{equation}\label{eq:c3}
  \Hsandup\big(Z_1^n\big|\Chat_1^nI_1^n\Ehat\big)_{\rho_{|\OmAT}}
  \;\ge\;n\,\hqkd-\frac{\alpha}{\alpha-1}\log_2\frac{1}{\Pr[\OmAT]},
\end{equation}
with $\hqkd=\hqkd(\varnothing)$ the single-round rate of
Definition~\ref{def:hqkd-list} at the empty list, which reads
\begin{equation}\label{eq:hqkd-def}
\begin{gathered}
  \hqkd=\inf_{\mathbf q\in\acSet}\
  \inf_{\substack{\omega\in\Seq{A\Qsys\Etil}\\ \omega\models\Decl}}
  \left(\frac{\alpha}{\alpha-1}\,
  \sfD\!\left(\mathbf q\,\middle\|\,\Phimap[\omega_{A\Qsys}]\right)
  +\sum_{\hat c\in\Chkey}
  q(\hat c)\,\Hsandup(Z|I\Etil)_{\nu^{\omega}_{|\hat c}}\right),\\
  \nu^{\omega}\coloneqq(\Mprot\otimes\idmap_{\Etil})(\omega),
\end{gathered}
\end{equation}
with $\omega\models\Decl$ as in \eqref{eq:margcond} --- so the inner
infimum is over states with $\omega_A=\tau_A$ in trusted-source mode
and over \emph{all} normalized states on $A\Qsys\Etil$ in
untrusted-source mode. It ranges over all (not necessarily pure) such
states, with
$\dim\Hs{\Etil}\ge\dim\Hs{A}\cdot\dim\Hs{\Qsys}$
(Lemma~\ref{lem:purifier-indep}: the value is independent of that
choice).
\end{theorem}
\end{mainclaim}
\claimmeta{\metalabel{Uses}\mref{Def.}{def:admissible-declaration}, \mref{Thm.}{thm:protocol-channel}, \mref{Lem.}{lem:branch-maps}, \mref{Fact}{fact:meat}, \mref{Def.}{def:roundmodel}, \mref{Def.}{def:hqkd-list}, \mref{Def.}{def:admissible-acceptance-set}, \mref{Lem.}{lem:purifier-indep}, \mref{Lem.}{lem:range}, \mref{Lem.}{lem:product-marginal}, \extref{../foundations/contract.pdf}{foundations/def:kl}, \extref{../foundations/contract.pdf}{foundations/def:renyi}
\metasep \metalabel{Assumes}\extaref{../transmitter/contract.pdf}{transmitter/ass:src}{A-SRC}, \extaref{../transmitter/contract.pdf}{transmitter/ass:ens}{A-ENS}, \extaref{../transmitter/contract.pdf}{transmitter/ass:seq}{A-SEQ}, \aref{A-SEED}{ass:seed}, \aref{A-LAB}{ass:lab}, \extaref{../protocol-compiler/contract.pdf}{protocol-compiler/ass:rhom}{A-RHOM}
\metasep \metalabel{Used by}\mref{Rem.}{rem:accumulation-scope}, \mref{Rem.}{rem:no-constraints}, \mref{Thm.}{thm:constrained-accumulation}, \mref{Rem.}{rem:losses}, \mref{Rem.}{rem:hlow-range}, \mref{Thm.}{thm:keylength}, \mref{Cor.}{cor:separate-executions}, \mref{Lem.}{lem:fo-khyp}
\metasep newly derived
\metasep implemented (effective: accepted)
\metasep \metalabel{Code}\coderef{https://github.com/fraunhofersg/certiqs-key/blob/v0.1.0/src/certiqs_key/meat_core/finitekey/keylength.py}{certiqs_key.meat_core.finitekey.keylength:key_length}
\metasep \metalabel{Oracles}\coderef{https://github.com/fraunhofersg/certiqs-key/blob/v0.1.0/tests/oracles/registry.py}{registry.py}}

\begin{proof}
\emph{Idea.} The bound is Fact~\ref{fact:meat} applied to the round
channels \eqref{eq:roundchannel}. Step~1 instantiates the source's marginal
register by the declared source mode, Step~2 checks the hypotheses of the
import, and Steps~3 and~4 bound the imported rate \eqref{eq:meat-rate} from
below by $\hqkd$: Step~3 through the range inclusion of
Lemma~\ref{lem:range}, Step~4 by reducing the sum over acceptance values to
the key-live ones.

\emph{Step 1 (the marginal register).} Apply Fact~\ref{fact:meat} to the
round channels \eqref{eq:roundchannel}, with the identifications
$S_j\to Z_j$, $S_{\Omega}\to\acSet$, and the source's marginal register
$A_{j-1}$ instantiated by the declared source mode:
\begin{itemize}[leftmargin=1.6em]
\item[] \emph{Trusted-source mode:} $A_{j-1}$ is the round-$j$ copy of
$A$ and $\sigma^{(j-1)}\to\tau_A$ (off-by-one indexing;
Fact~\ref{fact:meat}, locator note (v));
$\Mprot_j\in\mathrm{CPTP}(A\,E_{j-1},\,Z\,\Chat\,E_j)$.
\item[] \emph{Untrusted-source mode:} $A_{j-1}$ is \emph{trivial}
($\dim=1$) and $\sigma^{(j-1)}\to[1]$, the unique state on it;
$\Mprot_j\in\mathrm{CPTP}(E_{j-1},\,Z\,\Chat\,E_j)$, which is the same
map read with a trivial tensor factor adjoined to its input. That
instantiation is admissible: neither
\cite[Thm.~4.2a]{arqand2025meat} nor
\cite[Def.~4.2, Eq.~(98)]{arqand2025meat} places any dimension or rank
hypothesis on $A_{j-1}$, and the source's standing convention imposes
no bound on system dimensions
(Fact~\ref{fact:meat}, locator note (vii), where all three are quoted).
\end{itemize}

\emph{Step 2 (the hypotheses of the import).}
\begin{itemize}[leftmargin=1.6em]
\item Each $\Mprot_j$ is CPTP with $\Chat_j$ classical: composition and
tensoring of the CPTP maps $\Mprot$ (Theorem~\ref{thm:protocol-channel}\,(i), which
also makes $\Chat$ classical) and $\Aev_j$, followed by the output
regrouping $E_j=E'_jI_j$. The import takes every system
finite-dimensional and places no bound and no other structural
requirement on the $E_{j-1}$ registers (Fact~\ref{fact:meat}, locator
note (vii)); the round model of Definition~\ref{def:roundmodel} supplies
exactly that for the adversary's registers, with no bound on their
dimensions, and the bound uses none, which is what lets
Corollary~\ref{cor:composed-security} reach a memory of any dimension through
Lemma~\ref{lem:memory-lift}; no generality is lost by $\Qsys$ carrying a
sector decomposition.
\item Product marginal constraint on $A_1^n$: in trusted-source mode
this is the condition on $\omega^{0}$ in the statement, which
Lemma~\ref{lem:product-marginal}\,(ii) establishes for the protocol's
$n$-round input under its hypotheses, which are those the statement names
for that mode with the finite dimension of the adversary's initial
register, and at $d_A=1$ for every state. In
untrusted-source mode
$A_1^n$ is one-dimensional, $\Seq{A_1^n}$ is a singleton, and the
condition
$\omega^{0}_{A_1^n}=\sigma^{(0)}\otimes\dots\otimes\sigma^{(n-1)}$
holds for every normalized $\omega^{0}$ --- it discharges
\emph{vacuously}, which is a satisfied hypothesis and not a waived one.
Correspondingly
Lemma~\ref{lem:product-marginal} is not invoked in that mode and nothing
in this proof appeals to it.
\item Conditioning registers: the imported bound \eqref{eq:meat}
conditions on $\Chat_1^nE_n$, and the physical attack is represented
so that $E_n=I_1^n\Ehat$ (Definition~\ref{def:roundmodel}: the
announcements are classical and public, so the sequence that retains
all of them is available for every physical attack and produces the
same protocol run). This yields exactly the left-hand side of
\eqref{eq:c3}, the form Theorem~\ref{thm:keylength} uses, where Eve is
granted the announcement transcript and the $\Chat$ copies.
\item Round structure: all rounds use the same declaration, hence the
same $\Mprot$, the same source mode, and (in trusted-source mode) the
same $\tau_A$
(A-RHOM (\xref{protocol-compiler/ass:rhom})); Eve's channels $\Aev_j$ are \emph{not}
assumed identical across rounds. No $\Sigma=\Sigma_1$ collapse is
available or used --- the round channels $\Mprot_j$ are genuinely
non-isomorphic (the registers $E_{j-1}$ differ across $j$) --- and the
homogeneity of the declaration enters only through the common target
set of Lemma~\ref{lem:range}\,(ii). It is
$\Mprot$, not the trustedness of $\tau_A$, that supplies the sharing:
in untrusted-source mode the target sets
$\Sigma_j\subseteq\Sigma^{\CList}(\Etil^{*})$ are contained in a common
set for the same reason, with the marginal clause of
\eqref{eq:margcond} empty.
\item Purifying register: the import requires $\Etil$ large enough to
purify $A_{j-1}\,E_{j-1}$ --- a proof-internal requirement, discharged
by enlargement (\cite[footnote~17]{arqand2025meat}) and only weakened
when $A_{j-1}$ is trivial; the defining
optimization \eqref{eq:hqkd-def} needs only
$\dim\Hs{\Etil}\ge\dim\Hs{A\Qsys}$
(the dimension condition in the statement, with
Lemma~\ref{lem:purifier-indep}).
\item Acceptance: $\acSet$ is convex (Definition~\ref{def:admissible-acceptance-set}) and by
definition of $\OmAT$ every frequency with nonzero probability in
$\rho_{|\OmAT}$ lies in $\acSet$; $p_{\Omega}=\Pr[\OmAT]>0$.
\item $\alpha$ domain: the import supplies $(1,\infty]$; this
module invokes it on $(1,\infty)$, where
\xref{foundations/def:renyi} defines the entropies
(Fact~\ref{fact:meat}, locator note (vi)); no other import is used in
this claim.
\end{itemize}

\emph{Step 3 (the imported rate is at least $\hqkd$).} It remains to bound
the rate \eqref{eq:meat-rate} from below by $\hqkd$. By
Lemma~\ref{lem:range}\,(ii) with $\CList=\varnothing$ --- in either source
mode ---
$\Sigma\subseteq\Sigma^{\varnothing}(\Etil^{*})$ for a common register
$\Etil^{*}$, so by Lemma~\ref{lem:range}\,(iii) the rate is bounded
below by the value of \eqref{eq:meat-rate} over
$\Sigma^{\varnothing}(\Etil^{*})$, whose elements are
$\nu^{\sigma}=(\Mprot\otimes\idmap_{\Etil^{*}})(\sigma)$ with
$\sigma\models\Decl$; for these,
$\boldsymbol{\nu}_{\Chat}=\Phimap[\sigma_{A\Qsys}]$ and the
conditioning register regroups as $E\Etil=I\,(E'\Etil)$
(Lemma~\ref{lem:range}\,(ii)), so the bracket is bounded from below by
that of \eqref{eq:hqkd-list} at $\CList=\varnothing$, the $\hat c$-sum
reduced to $\Chkey$ (Step~4), and Lemma~\ref{lem:purifier-indep} converts the purifier
$\Etil^{*}$ to the $\Etil$ of \eqref{eq:hqkd-def} --- choosing $\Etil$ in the
application of Fact~\ref{fact:meat} large enough that
$\dim\Hs{E'_j\Etil}\ge\dim\Hs{A\Qsys}$ for every $j$, so that the
lemma's hypothesis holds for $\Etil^{*}$;
\cite[footnote~17]{arqand2025meat} permits the enlargement, which
only grows $\Sigma$ and weakens the bound in the safe direction.

\emph{Step 4 (reduction of the sum to the key-live values).} Fix
$\hat c\notin\Chkey$ with $q(\hat c)>0$ and
$\boldsymbol{\nu}_{\Chat}(\hat c)>0$. By Lemma~\ref{lem:branch-maps}\,(iv) and
\eqref{eq:povm-form}, every outcome pair contributing to $\hat c$ has
$r(x,v)=\perp$, so $Z=\perp$ deterministically in
$\nu^{\sigma}_{|\hat c}$: the conditional state is
$\dyad{\perp}_Z\otimes\nu'$, with $\nu'$ normalized. For a deterministic
classical $Z$, $\Hsandup(Z|\cdot)\ge0$: the divergence in the definition of
$\Hsandup$ (\xref{foundations/def:renyi}) evaluates on
$\dyad{\perp}\otimes\nu'$ against $\idon{Z}\otimes\sigma'$, and
$\Dsand(\dyad{\perp}\otimes\nu'\|\idon{Z}\otimes\sigma')
=\Dsand(\nu'\|\sigma')$, since the sandwich is
$\dyad{\perp}\otimes\sigma'^{\gamma}\nu'\sigma'^{\gamma}$ with
$\gamma=\frac{1-\alpha}{2\alpha}$; the state $\sigma'=\nu'$ is admissible in
the infimum and gives
$\Dsand(\nu'\|\nu')=\frac{1}{\alpha-1}\log_2\tr\nu'=0$, so the infimum is at
most $0$. Hence those terms are non-negative, and dropping them lowers the
sum: the sum over $\operatorname{supp}(\boldsymbol{\nu}_{\Chat})$ is at least
the sum over $\Chkey$, in the extended sense of \xref{foundations/def:kl}: a
term with
$\hat c\in\Chkey\setminus\operatorname{supp}(\boldsymbol{\nu}_{\Chat})$
and $q(\hat c)=0$ is $0$ by the $0\log_2\frac0v=0$ convention (the
undefined conditional state is multiplied by zero), while $q(\hat c)>0$
for such a $\hat c$ means $\mathbf q\not\ll\boldsymbol{\nu}_{\Chat}$, so
the divergence term is $+\infty$ and the bracket equals $+\infty$ on
both sides regardless of the convention adopted for the undefined
entropy term. Since Step~3 bounds the imported rate from below, a sum
bounded from below by the one over $\Chkey$ is all it needs.
\end{proof}

% claim meat-core/rem:accumulation-scope
\claimsummary{meat-core/rem:accumulation-scope}{The bound of \hyperlink{claim-meat-core/thm:accumulated-entropy}{Theorem~\ref*{thm:accumulated-entropy}} covers adaptive attacks with a finite-dimensional memory of any size, the tensor-product state model being a special case. In untrusted-source mode it holds with a smaller single-round rate, the price of not assuming a trusted source.}
\claimbookmark{meat-core/rem:accumulation-scope}{Remark}{rem:accumulation-scope}{Scope of the accumulated entropy bound}
\begin{remark}[Scope of the accumulated entropy bound]\label{rem:accumulation-scope}
The bound \eqref{eq:c3} of Theorem~\ref{thm:accumulated-entropy} covers
every attack in which Eve adapts each round to all earlier announcements and
keeps a quantum memory between rounds, finite-dimensional and of any size as
every register of the round model is; the security of the key, which this
bound feeds, extends to a memory of any dimension
(Corollary~\ref{cor:composed-security}, through
Lemma~\ref{lem:memory-lift}). The tensor-product state
model ($\rho=\Mprot^{\otimes n}[\omega]$,
$\omega_{A_1^n}=\tau_A^{\otimes n}$) is the special case in which each
$\Aev_j$ routes a pre-prepared register to $\Qsys$. In
untrusted-source mode the left-hand side and the constant
$\hqkd$ are the same expressions read with $\omega$ unconstrained
beyond normalization; the bound is then weaker (a larger inner feasible
set gives a smaller $\hqkd$), which is the direction of conservatism
and the price of not assuming a trusted source.
\end{remark}
\claimmeta{\metalabel{Uses}\mref{Thm.}{thm:accumulated-entropy}
\metasep newly derived
\metasep accepted}

% claim meat-core/rem:no-constraints
\claimsummary{meat-core/rem:no-constraints}{The imported theorem fixes the inner infimum over the outputs of all attacks, so a constraint list may raise the rate only through a licensed restriction of the attack model (\hyperlink{claim-meat-core/thm:constrained-accumulation}{Theorem~\ref*{thm:constrained-accumulation}}). Device knowledge without such a licence enters through the declared measurement.}
\claimbookmark{meat-core/rem:no-constraints}{Remark}{rem:no-constraints}{Where declared constraints may enter the rate}
\begin{remark}[Where declared constraints may enter the rate]
\label{rem:no-constraints}
The inner infimum of \eqref{eq:meat-rate} is fixed by the imported
theorem: for the unrestricted attack model of
Theorem~\ref{thm:accumulated-entropy} it ranges over the outputs of
\emph{all} attacks, and restricting it by a nonempty $\CList$ is not
licensed by anything in this section. What \emph{is} licensed is a
restriction of the attack model itself: when $\CList$ is
licensed for the declared source mode (Definition~\ref{def:homog}) and the
\exttermref{../reductions/contract.pdf}{term-reduction-chain}{reduction chain} that produced the declaration certifies that every
physical attack with finite registers whose first channel receives every
transmitted system is
represented by a $\CList$-admissible sequence (A-ATK
(\xref{protocol-compiler/ass:atk})), the single-round rate tightens to the
constrained $\hqkd(\CList)$ --- that is
Theorem~\ref{thm:constrained-accumulation}, and it is the only constrained
use this module licenses. Device and squashing knowledge that is not backed
by such a certificate enters soundly through the declared POVM $\{F_v\}$
itself, which needs no licence.
\end{remark}
\claimmeta{\metalabel{Uses}\mref{Fact}{fact:meat}, \mref{Def.}{def:homog}, \mref{Thm.}{thm:accumulated-entropy}
\metasep \metalabel{Assumes}\extaref{../protocol-compiler/contract.pdf}{protocol-compiler/ass:atk}{A-ATK}
\metasep newly derived
\metasep accepted}

\section{Constrained accumulation}\label{group-meat-core/thm:constrained-accumulation}
\sectionintro{When the reduction chain behind the declaration restricts every attack to a \termref{term-licensed-attack-set}{licensed attack set}, the accumulated entropy bound holds with the \termref{term-constrained-single-round-rate}{constrained single-round rate} of the declared list, with the same penalty (\hyperlink{claim-meat-core/thm:constrained-accumulation}{Theorem~\ref*{thm:constrained-accumulation}}). The licensed class is the one the licence transports and not an arbitrary choice: multiples of \ensuremath{\tau}\_A in trusted-source mode and scalars in untrusted-source mode (\hyperlink{claim-meat-core/rem:homog}{Remark~\ref*{rem:homog}}), a trusted-mode entry declaring the scalar so that no numerical test is needed (\hyperlink{claim-meat-core/rem:scalar-declaration}{Remark~\ref*{rem:scalar-declaration}}). For a party-local entry of the general format the theorem stays true, though that reading is not shown canonical (\hyperlink{claim-meat-core/rem:tagged-entries-and-accumulation}{Remark~\ref*{rem:tagged-entries-and-accumulation}}). The published flag-state squasher is an instance (\hyperlink{claim-meat-core/ex:fss}{Example~\ref*{ex:fss}}), and \hyperlink{claim-meat-core/rem:not-licensed}{Remark~\ref*{rem:not-licensed}} records what the licence does not cover.}
\sectionimports{\extref{../protocol-compiler/contract.pdf}{protocol-compiler/def:homog}}

% claim meat-core/thm:constrained-accumulation
\claimsummary{meat-core/thm:constrained-accumulation}{Under the attack-set licence \hyperref{../protocol-compiler/contract.pdf}{}{claim-protocol-compiler/ass:atk}{A-ATK}, the accumulated entropy bound holds with the constrained single-round rate of a licensed list in place of the unconstrained rate, with the same acceptance-probability penalty. The constrained rate is never below the unconstrained one.}
\claimbookmark{meat-core/thm:constrained-accumulation}{Theorem}{thm:constrained-accumulation}{Constrained accumulation under the attack-set licence}
\claimleadin{%
Whether an accumulation theorem admits per-round receiver-side constraints is
settled by re-importing the accumulation at the sequential level
(Fact~\ref{fact:meat}): the imported theorem is stated \emph{per channel
sequence}; which attack channels Eve is granted is a modelling input outside
the theorem (\cite[Def.~8.3, Rem.~8.4]{tupkary2026decoy}); and a declared
constraint list that is \exttermref{../protocol-compiler/contract.pdf}{term-licensed-for-the-declared-source-mode}{licensed for the declared source mode} is precisely
the state-level shadow of an attack-set restriction in the published format
\cite[Eq.~(151)]{tupkary2026decoy} (Definition~\ref{def:roundmodel},
Lemma~\ref{lem:range}\,(i)) --- the $\tau_A$-homogeneous entries of
Definition~\ref{def:homog}\,(a) transported by \eqref{eq:transport}, the
scalar entries of Definition~\ref{def:homog}\,(b) by \eqref{eq:transport-un}.
What remains an assumption, and is never asserted here, is that the
\exttermref{../reductions/contract.pdf}{term-reduction-chain}{reduction chain} behind the declaration actually delivers such a
restriction: A-ATK (\xref{protocol-compiler/ass:atk}), under which a
$\CList$-admissible \exttermref{../foundations/contract.pdf}{term-attack-sequence}{attack sequence} covers every attack available against the
physical protocol with finite registers whose first channel receives every
\exttermref{../transmitter/contract.pdf}{term-transmitted-system}{transmitted system}; the others are reached by a limit
(Corollary~\ref{cor:composed-security}).
}
\begin{theorem}[Constrained accumulation under the attack-set licence]\label{thm:constrained-accumulation}
Let the hypotheses of Theorem~\ref{thm:accumulated-entropy} hold in either
source mode, let $\CList$ be
\exttermref{../protocol-compiler/contract.pdf}{term-licensed-for-the-declared-source-mode}{licensed for that mode}
(Definition~\ref{def:homog}), let A-ATK
(\xref{protocol-compiler/ass:atk}) hold, and let the \exttermref{../foundations/contract.pdf}{term-attack-sequence}{attack sequence} be
$\CList$-admissible (Definition~\ref{def:roundmodel}). Then for every
$\alpha\in(1,\infty)$,
\begin{equation}\label{eq:c17}
  \Hsandup\big(Z_1^n\big|\Chat_1^nI_1^n\Ehat\big)_{\rho_{|\OmAT}}
  \;\ge\;n\,\hqkd(\CList)
  -\frac{\alpha}{\alpha-1}\log_2\frac{1}{\Pr[\OmAT]},
\end{equation}
with $\hqkd(\CList)$ the \termref{term-constrained-single-round-rate}{constrained single-round rate} of
Definition~\ref{def:hqkd-list}, and
$\hqkd(\CList)\ge\hqkd(\varnothing)=\hqkd$.
\end{theorem}
\claimmeta{\metalabel{Uses}\mref{Fact}{fact:meat}, \mref{Def.}{def:homog}, \mref{Def.}{def:roundmodel}, \mref{Def.}{def:hqkd-list}, \mref{Lem.}{lem:purifier-indep}, \mref{Lem.}{lem:range}, \mref{Thm.}{thm:accumulated-entropy}
\metasep \metalabel{Assumes}\extaref{../protocol-compiler/contract.pdf}{protocol-compiler/ass:atk}{A-ATK}
\metasep \metalabel{Used by}\mref{Rem.}{rem:tagged-entries-and-accumulation}, \mref{Rem.}{rem:homog}, \mref{Rem.}{rem:not-licensed}, \mref{Rem.}{rem:hlow-range}, \mref{Thm.}{thm:keylength}, \mref{Cor.}{cor:separate-executions}, \mref{Cor.}{cor:constrained-key-length}
\metasep newly derived
\metasep implemented (effective: accepted)
\metasep \metalabel{Code}\coderef{https://github.com/fraunhofersg/certiqs-key/blob/v0.1.0/src/certiqs_key/meat_core/finitekey/keylength.py}{certiqs_key.meat_core.finitekey.keylength:check_key_length_permitted}
\metasep \metalabel{Oracles}\coderef{https://github.com/fraunhofersg/certiqs-key/blob/v0.1.0/tests/oracles/registry.py}{registry.py}}

\begin{proof}
The proof of Theorem~\ref{thm:accumulated-entropy} applies unchanged, with one
refinement: Lemma~\ref{lem:range}\,(ii), invoked for the
$\CList$-admissible sequence, gives
$\Sigma\subseteq\Sigma^{\CList}(\Etil^{*})$, so
Lemma~\ref{lem:range}\,(iii), the $\Chkey$ reduction of the sum (verbatim
from the proof of Theorem~\ref{thm:accumulated-entropy}), and
Lemma~\ref{lem:purifier-indep} bound the rate of \eqref{eq:meat-rate}
below by $\hqkd(\CList)$. Every other hypothesis check is identical.
No inequality of the unconstrained proof is reversed anywhere: the
improvement over Theorem~\ref{thm:accumulated-entropy} is a restriction of
the admissible attack model, purchased entirely by Assumption~A-ATK. The
final inequality $\hqkd(\CList)\ge\hqkd(\varnothing)$ is
Definition~\ref{def:hqkd-list} (the constrained inner set is smaller).
\end{proof}

% claim meat-core/rem:tagged-entries-and-accumulation
\claimsummary{meat-core/rem:tagged-entries-and-accumulation}{No statement of the constrained accumulation theorem becomes false for a list with a party-local entry: its hypotheses are satisfiable, its conclusion holds and its monotonicity survives. The tagged reading is not shown canonical, and the paragraph before the theorem describes only untagged forms.}
\claimbookmark{meat-core/rem:tagged-entries-and-accumulation}{Remark}{rem:tagged-entries-and-accumulation}{Party-local entries and constrained accumulation}
\begin{remark}[Party-local entries and constrained accumulation]\label{rem:tagged-entries-and-accumulation}
Theorem~\ref{thm:constrained-accumulation}, which turns a licensed list into a
key-length licence, is stated for a list licensed for the declared mode and
for an \exttermref{../foundations/contract.pdf}{term-attack-sequence}{attack sequence} admissible for that list. Both readings carry
the \exttermref{../reductions/contract.pdf}{term-party-local-form}{party-local form} of \xref{protocol-compiler/def:homog}, the
licensed class being that definition's and the per-round admissibility format
of Definition~\ref{def:roundmodel} carrying the party-local admissibility set
beside the two of the untagged forms, with
Lemma~\ref{lem:range}\,(i) transporting it to the state level in either
mode. \emph{So no statement of that theorem becomes false of a list carrying
a tagged entry}: its hypotheses are satisfiable, its conclusion holds by that
transport, and its monotonicity $\hqkd(\CList)\ge\hqkd(\varnothing)$, where
$\hqkd(\CList)$ is that theorem's single-round QKD rate at order $\alpha$
under the constraints of $\CList$, holds because a tagged entry shrinks the
\exttermref{../protocol-compiler/contract.pdf}{term-feasible-set}{feasible set} further than its untagged reading does.

\emph{What it does not say} is that the tagged reading is canonical: the
argument that the licensed class is not an arbitrary choice is stated for the
two untagged forms (Remark~\ref{rem:homog}) and is not extended to the third
here, and the third form's soundness rests on the certificate on the scoped
factor's own channel (\xref{protocol-compiler/rem:party-local-soundness}) and
on nothing this contract proves.

\emph{The paragraph motivating the theorem describes the untagged forms.} That
paragraph, before the theorem's statement, rests it on three points: the
imported accumulation theorem is stated per sequence of attack channels;
which attack channels the adversary is granted is a modelling input outside
that theorem; and a constraint list licensed for the declared mode is
``precisely the state-level shadow of an attack-set restriction in the
published format'', the homogeneous entries of form (a) and the scalar
entries of form (b) of \xref{protocol-compiler/def:homog} each being carried
to the state level by a transport of its own. \emph{The third point is stated
of the untagged forms in its main clause, not only in the two transports it
lists}: the published format it cites is the single comparison of a channel
adjoint against a multiple of the identity, whereas the admissibility set of
a tagged entry is a \emph{family} of operator inequalities, one per reference
system and per positive semidefinite operator on it
(Definition~\ref{def:roundmodel}). \emph{What covers the third form is the
second point}, which makes restricting attention to a declared subset of
attack channels a modelling step and is stated for any declared subset, not
for a shape of right-hand side. So the theorem's hypotheses, conclusion and
proof are unaffected, and only the description in its motivating paragraph
does not reach the third form.
\end{remark}
\claimmeta{\metalabel{Uses}\mref{Def.}{def:roundmodel}, \mref{Def.}{def:hqkd-list}, \mref{Lem.}{lem:range}, \mref{Thm.}{thm:constrained-accumulation}, \extref{../protocol-compiler/contract.pdf}{protocol-compiler/def:homog}
\metasep newly derived
\metasep accepted}

% claim meat-core/ex:fss
\claimsummary{meat-core/ex:fss}{The restricted attack set of the published flag-state squasher is exactly the attack set of one licensed entry, a designated outcome's element minus \ensuremath{\lambda}\_min times the flag projector bounded below by 0, in either source mode. At least one published reduction thus delivers a licensed list.}
\claimbookmark{meat-core/ex:fss}{Example}{ex:fss}{The flag-state squasher as a licensed constraint list}
\begin{example}[The flag-state squasher as a licensed constraint list]
\label{ex:fss}
\cite[Lemma~9.4, Eqs.~(151)--(152)]{tupkary2026decoy} constructs, for
a receiver whose physical POVM is block-diagonal with respect to a
photon-number cutoff, a squashed model whose receiver space carries a
preserved subspace and a flag sector --- matching this module's
$\SecSet$ with
$\Hs{\Qsys}=\Hs{\Qsys_{\mathrm{pres}}}\oplus\Hs{\Qsys_{\mathrm{flag}}}$
--- together with the restricted \exttermref{../foundations/contract.pdf}{term-attack-set}{attack set}
\begin{equation*}
  \big\{\Aev:\ \Aev^{\dagger}(F_o\otimes\idon{E''})\;\succeq\;
  \lambda_{\min}\,\Aev^{\dagger}(\Pi_{\mathrm{flag}}\otimes\idon{E''})
  \big\},
\end{equation*}
its Eq.~(151). The source writes the flag projector $\Pi_{Sq}$, defined there
as ``the projector onto the flag subspace''; here it is
$\Pi_{\mathrm{flag}}$, the projector onto $\Hs{\Qsys_{\mathrm{flag}}}$, a
renaming into this module's sector vocabulary. $F_o$ is the squashed POVM
element of a designated outcome $o$ and $\lambda_{\min}>0$ a receiver
constant. By linearity of the adjoint this is exactly
$\AttSet_{\CList}$ of \eqref{eq:attset} for the single
$\tau_A$-homogeneous entry
\begin{equation*}
  \hat O=F_o-\lambda_{\min}\,\Pi_{\mathrm{flag}},\qquad
  \bowtie\ =\ \succeq,\qquad c=0,
\end{equation*}
whose operators are sector-block-diagonal as
Definition~\ref{def:admissible-declaration}\,(iii) requires ($F_o$ by declaration,
$\Pi_{\mathrm{flag}}$ a sector projector); the dimension-independent
marginal form of the same set is the source's Eq.~(152).

The instance carries over verbatim to \exttermref{../transmitter/contract.pdf}{term-untrusted-source-mode}{untrusted-source mode}: the
source's condition is on operators of $\Qsys$ only, so lifting it as
\eqref{eq:attset-un} prescribes gives
$\Aev^{\dagger}(\idon{A}\otimes F_o\otimes\idon{E''})\succeq
\lambda_{\min}\Aev^{\dagger}(\idon{A}\otimes\Pi_{\mathrm{flag}}
\otimes\idon{E''})$, i.e.\ the same single entry
$\hat O=F_o-\lambda_{\min}\Pi_{\mathrm{flag}}$, $\bowtie\,=\,\succeq$,
$c=0$, now read with the scalar right-hand side of
Definition~\ref{def:homog}\,(b); and $c=0$ is a scalar in either
reading, so no wording of the source is stretched.

The instance shows that A-ATK (\xref{protocol-compiler/ass:atk})
is not vacuous: at least one published reduction delivers a licensed pair
$(\CList,\AttSet_{\CList})$. Determining $\lambda_{\min}$ for a given
receiver, and certifying the reduction itself, are obligations of the
receiver and reductions modules; nothing here depends on how they are met.
\end{example}
\claimmeta{\metalabel{Uses}\mref{Def.}{def:admissible-declaration}, \mref{Def.}{def:homog}, \mref{Def.}{def:roundmodel}
\metasep \metalabel{Assumes}\extaref{../protocol-compiler/contract.pdf}{protocol-compiler/ass:atk}{A-ATK}, \extaref{../transmitter/contract.pdf}{transmitter/ass:src}{A-SRC}
\metasep newly derived
\metasep accepted}

% claim meat-core/rem:homog
\claimsummary{meat-core/rem:homog}{The licence transports exactly \ensuremath{\tau}\_A-multiples in trusted mode and scalars in untrusted mode, where the operator form c\_k\ensuremath{\cdot}1\_A can empty the feasible set. On directions the licence can probe, sound trusted-mode right-hand sides are dominated by homogeneous ones; data stay out of the declaration.}
\claimbookmark{meat-core/rem:homog}{Remark}{rem:homog}{Canonicity of the licensed class}
\begin{remark}[Canonicity of the licensed class]\label{rem:homog}
Three observations make Definition~\ref{def:homog} the natural
licensed class rather than an arbitrary restriction.

\emph{(a) The licence transports exactly the licensed right-hand sides, and
nothing wider.} In \exttermref{../transmitter/contract.pdf}{term-trusted-source-mode}{trusted-source mode} the admissibility format
\eqref{eq:attset} compares $\Aev^{\dagger}(\hat O_k\otimes\idon{})$ against
multiples of the identity, and Lemma~\ref{lem:range}\,(i) turns that into the
state-level right-hand side $c_k\,\tau_A$; the trusted marginal is the only
operator on $A$ the per-round model supplies.

In untrusted-source mode there is no such operator, and the difference
is structural rather than a matter of taste: $A$ is an \emph{output}
of $\Aev_j$ (Definition~\ref{def:roundmodel}), so the input side of the
adjoint identity carries no $A$ factor at all and the operator-valued
identity \eqref{eq:transport} has no analogue. Only its trace
survives, giving \eqref{eq:transport-un} and the scalar class of
Definition~\ref{def:homog}\,(b). \emph{The operator form
$\Theta_k=c_k\idon{A}$ is therefore not licensed}, and it is not
merely unlicensed but wrong: take $d_A=2$, $\hat O=\idon{\Qsys}$,
$\bowtie\,=\,\succeq$ and $c=1$. Every CPTP $\Aev$ is admissible, since
$\Aev^{\dagger}(\idon{A}\otimes\idon{\Qsys}\otimes\idon{E''})
=\idon{E'}\succeq1\cdot\idon{E'}$ by trace preservation, so
$\AttSet_{\CList}$ is the whole of $\mathrm{CPTP}$ and the declaration
restricts nothing. The transported scalar constraint is
$\tr[(\idon{A}\otimes\idon{\Qsys})\,\omega]=\tr\omega=1\ge c=1$, which
every feasible $\omega$ satisfies --- correct and vacuous, as it should
be. The operator form would instead demand
$\tr_{\Qsys}[(\idon{A}\otimes\idon{\Qsys})\,\omega]
=\omega_A\succeq c\,\idon{A}=\idon{A}$, which \emph{no} state
satisfies at $d_A\ge2$ ($\tr\omega_A=1<2=\tr\idon{A}$). It would empty
$\Feas$, violating Definition~\ref{def:admissible-declaration}\,(vii), and would do so
for a declaration whose \termref{term-licensed-attack-set}{licensed attack set} is everything. Writing
$\Theta_k=c_k\idon{A}$ in untrusted-source mode is thus a defect, not
a conservative simplification.

\emph{The sentence ``only its trace survives'' is scoped to the per-round
model's own operator on $A$, and the general format shows why the scope
matters.} What has no analogue in untrusted-source mode is the
\emph{operator-valued identity on $A$ against $\tau_A$}, because there is no
declared $\tau_A$ to compare against; what does survive, beside the trace, is
an operator-valued statement on $A$ \emph{tensored with the other
\exttermref{../protocol-compiler/contract.pdf}{term-measured-factor}{measured factors}}, and it survives not by weakening
the obstruction but by a strictly stronger channel-level condition on the
declaring side: the party-local admissibility format of
Definition~\ref{def:roundmodel}, which asks
$\Aev_j^{\dagger}(X\otimes\hat D\otimes\idon{})\bowtie_k0$ for \emph{every}
positive semidefinite $X$ on the retained and unscoped systems, and which
Lemma~\ref{lem:range}\,(i) transports to
$\tr_{\Qsys_f}[(\idon{}\otimes\hat D)\,\varsigma]\bowtie_k0$ as an operator,
in either mode. \emph{So ``only its trace survives'' is true of what the
\emph{scalar} admissibility format \eqref{eq:attset} supplies and is not true
of the licensed class as a whole once the party-local format is declared};
the difference is that the scalar format tests one operator inequality on
Eve's register and the party-local format tests a family of them, one per
$X$, which is more to license and not less. \emph{And the pathology of the
operator-form counterexample does not recur there}: the untrusted-mode
counterexample $\hat O=\idon{\Qsys}$, $\bowtie\,=\,\succeq$, $c=1$ read
party-locally is not admitted with an empty \exttermref{../protocol-compiler/contract.pdf}{term-feasible-set}{feasible set}, because the
same channel-level condition it would need is
$\Aev_j^{\dagger}(X\otimes0\otimes\idon{})\succeq0$ for every $X$ with
$\hat D=0$, which is satisfied, while the state-level reading is
$\tr_{\Qsys}[(\idon{}\otimes0)\,\varsigma]\succeq0$, which every state
satisfies --- correct and vacuous, as the scalar reading is. \emph{The
canonicity argument of (a) is not extended to the third form}: the argument
that the licensed class is not an arbitrary choice is made here for the two
untagged forms, and the third form's soundness rests on the certificate on
the scoped factor's own channel
(\xref{protocol-compiler/rem:party-local-soundness}) and on nothing this
remark proves.

\emph{(b) In trusted-source mode, sound right-hand sides are dominated
by homogeneous ones on the certifiable directions.} Restrict attention
to the discard-and-prepare channels the licensed set \emph{actually
contains}: for every
$\Aev=\varrho\,\tr_{E'}[\,\cdot\,]\in\AttSet_{\CList}(E',\Qsys E'')$
(a state $\varrho$ on $\Qsys E''$ prepared regardless of the input),
the reachable marginals include $\tau_A\otimes\varrho_{\Qsys}$, whose
constraint value $\tr[\hat O_k\,\varrho_{\Qsys}]\,\tau_A$ is a scalar
multiple of $\tau_A$. A sound right-hand side $\Theta_k$ for
$\bowtie_k\,=\,\succeq$ must therefore satisfy
$\Theta_k\preceq\big(\inf\{\tr[\hat O_k\,\varrho_{\Qsys}]:
\varrho\,\tr_{E'}[\,\cdot\,]\in\AttSet_{\CList}\}\big)\,\tau_A$
whenever that set is nonempty: on the directions the licensed model
can probe, $\Theta_k$ is dominated by a homogeneous right-hand side.
\emph{Which} discard-and-prepare channels are licensed is decided by the
reduction chain behind the declaration, not by this module: Eve may
always discard her memory and inject a fresh \emph{physical} state, but what
reaches $\Qsys$ after the reduction is its squashed image, which in
general ranges over a proper subset of the states on $\Qsys$. In
Example~\ref{ex:fss} the licensed set deliberately excludes some
preparations --- a fresh flag state $\dyad{o'}$ with $o'\ne o$
violates \eqref{eq:attset} --- and exactly that exclusion is what
makes the attack set restrictive and
Theorem~\ref{thm:constrained-accumulation} non-vacuous. The canonical status
of the licensed class therefore rests on (a) --- the licence transports only
$\tau_A$-multiples in trusted-source mode, and only scalars in
untrusted-source mode, in the first place --- with (b) a domination statement
over the licensed subset in trusted-source mode, not a universal necessity
claim.

\emph{(c) Data stays out of the declaration.} $\CList$ is
part of $\Decl$, fixed before any run; an observed number on a
right-hand side would make the declaration data-dependent. The only
sanctioned carrier of observed data is the acceptance set $\acSet$
(Definition~\ref{def:admissible-acceptance-set}); a constraint whose $\Theta_k$ encodes
statistics of the current run is unsound here regardless of its
provenance.
\end{remark}
\claimmeta{\metalabel{Uses}\mref{Def.}{def:admissible-declaration}, \mref{Def.}{def:homog}, \mref{Def.}{def:roundmodel}, \mref{Def.}{def:admissible-acceptance-set}, \mref{Lem.}{lem:range}, \mref{Thm.}{thm:constrained-accumulation}
\metasep \metalabel{Used by}\mref{Rem.}{rem:not-licensed}, \mref{Rem.}{rem:cgate}
\metasep newly derived
\metasep accepted}

% claim meat-core/rem:scalar-declaration
\claimsummary{meat-core/rem:scalar-declaration}{A trusted-mode entry declares the scalar c\_k, not an operator \ensuremath{\Theta}\_k said to equal c\_k\ensuremath{\cdot}\ensuremath{\tau}\_A, so its right-hand side and the transported bound are one expression. An operator only near c\_k\ensuremath{\cdot}\ensuremath{\tau}\_A would need a numerical test, and no relaxation of the program covers admitting it.}
\claimbookmark{meat-core/rem:scalar-declaration}{Remark}{rem:scalar-declaration}{Scalar declaration of a homogeneous right-hand side}
\begin{remark}[Scalar declaration of a homogeneous right-hand side]\label{rem:scalar-declaration}
Licensing is a type check in both modes and compares no numbers
(Definition~\ref{def:homog}). This is deliberate, and it is the reason the
trusted-mode entry declares ``a declared $c_k$'' and not ``a declared
$\Theta_k$ satisfying $\Theta_k=c_k\tau_A$''. Under the latter form, a
declaration carrying an operator $\Theta_k$ merely \emph{near} $c_k\tau_A$
would have to be admitted or refused by a numerical test, and admitting one
would be unsafe in a way no relaxation of this module covers: the transport
Lemma~\ref{lem:range}\,(i) and \eqref{eq:transport} deliver exactly
$c_k\tau_A$ on the right-hand side, while \eqref{eq:feas-modes} would impose
the declared $\Theta_k$, and the two feasible sets are then neither nested nor
related by the relaxation of Definition~\ref{def:relaxed-feasible-set}, whose
window is a spectral floor, or an entrywise ball, on the \emph{relation}, and
not an entrywise-relative window on the right-hand side. With the scalar
declaration, $\Theta_k$ and the transported bound are the same expression in
the same $c_k$ and the same $\tau_A$, and the two agree by construction on the
declared numbers.
\end{remark}
\claimmeta{\metalabel{Uses}\mref{Def.}{def:homog}, \mref{Lem.}{lem:range}
\metasep newly derived
\metasep accepted}

% claim meat-core/rem:not-licensed
\claimsummary{meat-core/rem:not-licensed}{The constrained accumulation theorem licenses only restricting the single-round infimum to marginals meeting a licensed list. Limits on the purifying register, data-dependent or multi-round constraints, a trusted marginal, other proof techniques, the reduction and the test statistics lie beyond it.}
\claimbookmark{meat-core/rem:not-licensed}{Remark}{rem:not-licensed}{Limits of the attack-set licence}
\begin{remark}[Limits of the attack-set licence]\label{rem:not-licensed}
Theorem~\ref{thm:constrained-accumulation} licenses exactly one thing:
restricting the single-round infimum to marginals satisfying a declared list
that is licensed for the declared source mode (Definition~\ref{def:homog}; in
the \exttermref{../protocol-compiler/contract.pdf}{term-basic-format}{basic format}, $\tau_A$-homogeneous in trusted-source mode and
scalar in untrusted-source mode), under A-ATK
(\xref{protocol-compiler/ass:atk}). In the general format of
Definition~\ref{def:declaration} the licensed class carries a third form, the
entries declared \exttermref{../protocol-compiler/contract.pdf}{term-party-local}{party-local}, in either mode, whose right-hand side is
a real $c_k$ and whose state-level reading is the operator relation of
Definition~\ref{def:homog}. What that adds to the following list is
\emph{nothing}: every item (a) to (g) is a statement about what is \emph{not}
licensed, each holds of the third form as it holds of the other two, and
item (b)'s right-hand sides read $c_k\tau_A$, $c_k$, or $c_k$ times the
marginal that is actually there, all three with \emph{declared} constants.
The theorem does \emph{not} license:
\begin{enumerate}[label=(\alph*),leftmargin=2.2em]
\item \emph{Constraints on the purifying register.} $\CList$'s syntax
touches $\sigma_{A\Qsys}$ only; the inner infimum of
\eqref{eq:hqkd-list} still ranges over \emph{all} extensions of a
feasible $A\Qsys$-marginal (Lemma~\ref{lem:purifier-indep}), and no
constraint on $\Etil$ --- equivalently, on Eve's side information ---
is expressible or admitted.
\item \emph{Data-dependent constraints.} $\Theta_k=c_k\tau_A$
(respectively $\Theta_k=c_k$) with \emph{declared} constants $c_k$;
observed statistics enter only through
$\acSet$ (Remark~\ref{rem:homog}\,(c)).
\item \emph{A trusted marginal.} Nothing in
Theorem~\ref{thm:constrained-accumulation} licenses declaring $\tau_A$. That
is a separate input with a separate attestation, A-SRC
(\xref{transmitter/ass:src}), and in particular a licensed $\CList$ never
substitutes for it: a constraint list restricts what \emph{Eve's
channel} may do, whereas $\tau_A$ asserts that a register never
reached her at all.
\item \emph{Joint multi-round constraints.} Admissibility is per round
(Definition~\ref{def:roundmodel}); no constraint on the joint
$n$-round state is licensed, and none is expressible in $\CList$.
\item \emph{Reuse under other proof techniques.} The restriction is
compatible with the accumulation route imported here; it is \emph{not}
usable with postselection-based reductions
(\cite[Rem.~9.4]{tupkary2026decoy}). A module importing $\CList$ into
a different proof technique must re-establish the licence there.
\item \emph{The reduction itself.} Nothing here establishes a
squashing reduction, the existence of a licensed attack set for a
given receiver, or the value of a receiver constant such as
$\lambda_{\min}$ (Example~\ref{ex:fss}); these are obligations of the
reductions and receiver modules, assumed through Assumption~A-ATK and never
verified here.
\item \emph{Test-statistics changes.} Nothing here licenses a
fixed-fraction test-statistics model
(\xref{protocol-compiler/req:no-fixed-fraction}) or any modification of
$\acSet$ (A-TSM (\xref{reductions/ass:tsm}) stands unchanged).
\end{enumerate}
\end{remark}
\claimmeta{\metalabel{Uses}\mref{Def.}{def:homog}, \mref{Def.}{def:roundmodel}, \mref{Def.}{def:hqkd-list}, \mref{Lem.}{lem:purifier-indep}, \mref{Thm.}{thm:constrained-accumulation}, \mref{Rem.}{rem:homog}
\metasep \metalabel{Assumes}\extaref{../protocol-compiler/contract.pdf}{protocol-compiler/ass:atk}{A-ATK}, \extaref{../transmitter/contract.pdf}{transmitter/ass:src}{A-SRC}, \extaref{../reductions/contract.pdf}{reductions/ass:tsm}{A-TSM}
\metasep newly derived
\metasep accepted}

\section{Branch entropy and sectors}\label{group-meat-core/thm:branch-formula}
\sectionintro{A pure input's conditional entropy on one acceptance value is a function of the protocol marginal alone, through the ratio of a pinched trace-power quantity to the outcome probability, and that quantity splits into sector deficits on block-diagonal inputs (\hyperlink{claim-meat-core/thm:branch-formula}{Theorem~\ref*{thm:branch-formula}}). The pinched trace power is concave on the whole positive cone (\hyperlink{claim-meat-core/lem:cone-concavity}{Lemma~\ref*{lem:cone-concavity}}), and every sector deficit is non-positive (\hyperlink{claim-meat-core/cor:f-nonpos}{Corollary~\ref*{cor:f-nonpos}}) and vanishes on a key-inert pair, whose variables still enter the statistics (\hyperlink{claim-meat-core/prop:exact-cancellation}{Proposition~\ref*{prop:exact-cancellation}}). With Alice sectors the same holds per joint sector (\hyperlink{claim-meat-core/lem:joint-sectors}{Lemma~\ref*{lem:joint-sectors}}).}
\sectionimports{\extref{../foundations/contract.pdf}{foundations/fact:concavity}, \extref{../foundations/contract.pdf}{foundations/fact:pinch-id}, \extref{../foundations/contract.pdf}{foundations/lem:pinchmono}}

% claim meat-core/thm:branch-formula
\claimsummary{meat-core/thm:branch-formula}{A pure input's conditional entropy on one acceptance value is 1/(1\ensuremath{-}\ensuremath{\alpha}) times log\ensuremath{_{2}} of K/P, the ratio of a pinched trace-power quantity K to the outcome probability P, and K splits into sector deficits on block-diagonal inputs. The bounds 0 \ensuremath{\le} K \ensuremath{\le} P hold on the whole positive cone.}
\claimbookmark{meat-core/thm:branch-formula}{Theorem}{thm:branch-formula}{Branch entropy from the protocol marginal}
\claimleadin{%
The right-hand side of the pinching identity of \xref{foundations/fact:pinch-id}
depends on the state only through $\rho_Q$; this is what makes the objective of
this module a function of the $A\Qsys$-marginal alone, and it is why purity of
the global state must be arranged before the identity is invoked
(Theorem~\ref{thm:reduction}).
}
\begin{theorem}[Branch entropy from the protocol marginal]\label{thm:branch-formula}
For $\alpha\in(1,\infty)$, $\omega\in\Pos{A\Qsys}$ and $\hat c\in\Chatset$
define
\begin{equation}\label{eq:fK}
  P_{\hat c}(\omega)\coloneqq\Phimap[\omega](\hat c),\qquad
  T_{\hat c}(\omega)\coloneqq
  \tr\!\left[\left(\cZ\big(\cG_{\hat c}(\omega)^{1/\alpha}\big)
  \right)^{\alpha}\right],\qquad
  K_{\hat c}(\omega)\coloneqq
  T_{\hat c}(\omega)+P_{\hat c}(\omega)-\tr\cG_{\hat c}(\omega),
\end{equation}
and, for block-diagonal $\omega=\bigoplus_s\omega_s$, the sector deficits
\begin{equation}\label{eq:f-def}
  f_{\hat c,s}(\omega_s)\coloneqq
  \tr\!\left[\left(\cZ\big(\cG_{\hat c,s}(\omega_s)^{1/\alpha}\big)
  \right)^{\alpha}\right]-\tr\cG_{\hat c,s}(\omega_s).
\end{equation}
Let $\alpha\in(1,\infty)$, let $\omega\in\Seq{A\Qsys\Etil}$ be \emph{pure},
let $\nu^{\omega}=(\Mprot\otimes\idmap_{\Etil})(\omega)$, and let
$\hat c\in\Chatset$ with $P_{\hat c}\coloneqq
P_{\hat c}(\omega_{A\Qsys})>0$. Then:
\begin{enumerate}[label=(\roman*),leftmargin=2.2em]
\item \emph{(the identity)}
\begin{equation}\label{eq:c5}
  \Hsanddown(Z|I\Etil)_{\nu^{\omega}_{|\hat c}}
  =\frac{1}{1-\alpha}\log_2\frac{K_{\hat c}(\omega_{A\Qsys})}{P_{\hat c}};
\end{equation}
\item \emph{(the bounds)} $0\le K_{\hat c}(\omega_{A\Qsys})\le P_{\hat c}$, so
the right-hand side of \eqref{eq:c5} is non-negative;
\item \emph{(the sector form, on the whole cone)} for every
$\hat c\in\Chatset$ and \emph{every} block-diagonal
$\varrho=\bigoplus_s\varrho_s\in\Pos{A\Qsys}$, with none of the hypotheses
above on $\omega$ and $P_{\hat c}$ (no purity, no normalization and no
hypothesis $P_{\hat c}(\varrho)>0$),
\begin{equation}\label{eq:c5-sector}
  K_{\hat c}(\varrho)
  =P_{\hat c}(\varrho)+\sum_{s\in\SecSet}f_{\hat c,s}(\varrho_s).
\end{equation}
\end{enumerate}
The two bounds of (ii) hold on the whole cone, as a statement separate from
the identity (i): for every $\alpha\in(1,\infty)$, every
$\hat c\in\Chatset$ and \emph{every} $\varrho\in\Pos{A\Qsys}$, with no
purity, no normalization, no block-diagonality and no hypothesis
$P_{\hat c}(\varrho)>0$, one has $0\le K_{\hat c}(\varrho)\le
P_{\hat c}(\varrho)$. We call this the theorem's \emph{cone clause}.
\end{theorem}
\claimmeta{\metalabel{Uses}\mref{Def.}{def:maps}, \mref{Thm.}{thm:protocol-channel}, \mref{Def.}{def:branch}, \mref{Lem.}{lem:branch-maps}, \extref{../foundations/contract.pdf}{foundations/fact:pinch-id}, \extref{../foundations/contract.pdf}{foundations/lem:pinchmono}
\metasep \metalabel{Used by}\mref{Cor.}{cor:concavity}, \mref{Lem.}{lem:cone-concavity}, \mref{Cor.}{cor:f-nonpos}, \mref{Prop.}{prop:exact-cancellation}, \mref{Def.}{def:joint-sectors}, \mref{Lem.}{lem:joint-sectors}, \mref{Thm.}{thm:reduction}, \mref{Prop.}{prop:a-neutral}, \mref{Rem.}{rem:losses}, \mref{Prop.}{prop:value-range}, \mref{Thm.}{thm:convexity}, \mref{Lem.}{lem:choi-factor}, \mref{Def.}{def:untrusted-choi}, \mref{Lem.}{lem:kraus-freedom}, \mref{Thm.}{thm:oracle}, \mref{Lem.}{lem:rank-one-sector}, \mref{Def.}{def:dropped}, \mref{Cor.}{cor:drop-program}, \mref{Lem.}{lem:objective-homogeneous}, \mref{Prop.}{prop:relaxed-set-properties}, \mref{Lem.}{lem:fo-khyp}, \mref{Thm.}{thm:firstorder}, \mref{Rem.}{rem:fo-danger}, \mref{Lem.}{lem:certified-bracket}, \mref{Rem.}{rem:fo-negative-values}
\metasep newly derived
\metasep implemented (effective: accepted)
\metasep \metalabel{Code}\coderef{https://github.com/fraunhofersg/certiqs-key/blob/v0.1.0/src/certiqs_key/meat_core/objective/branch.py}{certiqs_key.meat_core.objective.branch:evaluate_branch}
\metasep \metalabel{Oracles}\coderef{https://github.com/fraunhofersg/certiqs-key/blob/v0.1.0/tests/oracles/registry.py}{registry.py}}

\begin{proof}
\emph{The identity (i).} Write $Q_{\mathrm{tot}}\coloneqq R\,XY\,C_AC_B\,A\Qsys$ and
$\ket{\psi}\coloneqq\Wprot\ket{\omega}_{A\Qsys\Etil}$, a pure state on
$Q_{\mathrm{tot}}\,\Chat\,I\,\Etil$ (the isometry $\Wprot$ of
\eqref{eq:Mprot} acts trivially on $\Etil$; $\Chat$ and $I$ are classical
in $\psi$ in the sense that $\psi$ is invariant under their computational
dephasings, both being controlled copies). Let
$\ket{\psi_{\hat c}}\coloneqq
P_{\hat c}^{-1/2}\,\bra{\hat c}_{\Chat}\ket{\psi}$, a normalized pure state
on $Q_{\mathrm{tot}}\,I\,\Etil$; normalization holds because
$\|\bra{\hat c}\psi\|^2=\Pr[\Chat=\hat c]=P_{\hat c}$ by
Theorem~\ref{thm:protocol-channel}\,(iii). By construction
$\nu^{\omega}_{|\hat c}$ on $Z\,I\,\Etil$ equals
$\tr_{Q_{\mathrm{tot}}}[\VZ\,\dyad{\psi_{\hat c}}\,\VZ^{\dagger}]$ up to the
retained $\Chat$ value, which is the fixed classical symbol $\hat c$ and may
be dropped from the conditioning without change.

Apply \xref{foundations/fact:pinch-id} with $Q\to Q_{\mathrm{tot}}$,
$E\to I\Etil$, projectors $Z_j=\dyad{j}_R\otimes\idon{}$
($j\in\Zset$; orthogonal, complete), and the pure state
$\dyad{\psi_{\hat c}}$: hypotheses hold verbatim, and
\begin{equation}\label{eq:c5-step1}
  \Hsanddown(Z|I\Etil)_{\nu^{\omega}_{|\hat c}}
  =\frac{1}{1-\alpha}\log_2
  \tr\!\left[\left(\cZ\big(\varrho^{1/\alpha}\big)\right)^{\alpha}\right],
  \qquad
  \varrho\coloneqq\tr_{I\Etil}\dyad{\psi_{\hat c}}.
\end{equation}
Compute $\varrho$. Expanding $\Wprot$ as in the proof of
Theorem~\ref{thm:protocol-channel}\,(ii) and projecting $\Chat$ on $\hat c$ leaves the sum
over pairs $(x,v)$, $(x',v')$ with
$\phi(a(x,v))=\phi(a(x',v'))=\hat c$. The register $I$ holds
$\ket{a(x,v)}\bra{a(x',v')}$; tracing it out annihilates every term with
$a(x,v)\ne a(x',v')$, that is, all cross-\emph{class} terms. Tracing $\Etil$
replaces the rank-one $\dyad{\omega}$ by $\omega_{A\Qsys}$ inside each
remaining term, since the operators $K_{\kappa}$ of \eqref{eq:kraus} act
only on $A\Qsys$:
$\tr_{\Etil}\big[(K_{\kappa}\otimes\idon{\Etil})\dyad{\omega}
(K_{\kappa'}\otimes\idon{\Etil})^{\dagger}\big]
=K_{\kappa}\,\omega_{A\Qsys}\,K_{\kappa'}^{\dagger}$, and only
$\kappa=\kappa'$ survives. Hence
\begin{equation}\label{eq:varrho}
  \varrho=\frac{1}{P_{\hat c}}
  \sum_{\kappa\in\phi^{-1}(\hat c)}K_{\kappa}\,\omega_{A\Qsys}\,
  K_{\kappa}^{\dagger}
  =\frac{1}{P_{\hat c}}\big(\cG_{\hat c}(\omega_{A\Qsys})
  +\Theta_{\hat c}(\omega_{A\Qsys})\big).
\end{equation}
By Lemma~\ref{lem:branch-maps}\,(ii) the two summands have orthogonal supports
($R\in\Zval$ versus $R=\perp$), so
$\varrho^{1/\alpha}=P_{\hat c}^{-1/\alpha}
\big(\cG_{\hat c}(\omega)^{1/\alpha}+\Theta_{\hat c}(\omega)^{1/\alpha}\big)$,
and $\cZ$ preserves both parts, acting as the identity on the second.
Therefore
\begin{equation*}
  \left(\cZ\big(\varrho^{1/\alpha}\big)\right)^{\alpha}
  =\frac{1}{P_{\hat c}}\left[
  \left(\cZ\big(\cG_{\hat c}(\omega)^{1/\alpha}\big)\right)^{\alpha}
  +\Theta_{\hat c}(\omega)\right],
\end{equation*}
again by orthogonality of supports (the two $\alpha$-th powers act
blockwise) and
$(\Theta^{1/\alpha})^{\alpha}=\Theta$. Taking the trace and using
$\tr\Theta_{\hat c}(\omega)=P_{\hat c}-\tr\cG_{\hat c}(\omega)$
(Lemma~\ref{lem:branch-maps}\,(i)) gives
$\tr[(\cZ(\varrho^{1/\alpha}))^{\alpha}]
=K_{\hat c}(\omega)/P_{\hat c}$, which with \eqref{eq:c5-step1} is
\eqref{eq:c5}.

\emph{The bounds (ii) and the cone clause.} $T_{\hat c}\ge0$
because it is the trace of a positive semidefinite operator, and
$P_{\hat c}-\tr\cG_{\hat c}=\tr\Theta_{\hat c}\ge0$ by \eqref{eq:traces},
which Lemma~\ref{lem:branch-maps}\,(i) states for \emph{all}
$\omega\in\Pos{A\Qsys}$; together these give $K_{\hat c}\ge0$. And
$T_{\hat c}\le\tr\cG_{\hat c}$ is
\xref{foundations/lem:pinchmono}, applied at
$\tau=\cG_{\hat c}(\varrho)$ (that lemma is stated on all of $\Pos{}$
and consumes no trace hypothesis), giving $K_{\hat c}\le P_{\hat c}$.
Neither step reads purity, normalization, $P_{\hat c}>0$ or
block-diagonality; both read $\varrho\succeq0$ and nothing else, which is
the cone clause.

\emph{The sector form (iii).} By Lemma~\ref{lem:branch-maps}\,(iii),
$\cG_{\hat c}(\varrho)=\bigoplus_s\cG_{\hat c,s}(\varrho_s)$; powers and
$\cZ$ act blockwise on this decomposition ($\cZ$ pinches $R$, which is a
tensor factor of every sector subspace), so
$T_{\hat c}(\varrho)=\sum_sT_{\hat c,s}(\varrho_s)$, with
$T_{\hat c,s}(\varrho_s)\coloneqq
\tr[(\cZ(\cG_{\hat c,s}(\varrho_s)^{1/\alpha}))^{\alpha}]$, and
$\tr\cG_{\hat c}(\varrho)=\sum_s\tr\cG_{\hat c,s}(\varrho_s)$; insert into
\eqref{eq:fK} and \eqref{eq:f-def}. This reads $\varrho\succeq0$ and its
block-diagonality and nothing else: neither purity nor $P_{\hat c}>0$.
\end{proof}

% claim meat-core/cor:concavity
\claimsummary{meat-core/cor:concavity}{For \ensuremath{\alpha} \textgreater{} 1 the pinched trace power is concave on subnormalised operators, so every sector deficit is concave on subnormalised sector states and each K is concave on the feasible set. It is the imported concavity result read at the parameter 1/\ensuremath{\alpha}.}
\claimbookmark{meat-core/cor:concavity}{Corollary}{cor:concavity}{Concavity of the branch quantities on subnormalised states}
\begin{corollary}[Concavity of the branch quantities on subnormalised states]
\label{cor:concavity}
For $\alpha\in(1,\infty)$, the map
$X\mapsto\tr[(\cZ(X^{1/\alpha}))^{\alpha}]=Q_{1/\alpha}(X)$ is concave on
$\Sleq{}$ of its space; consequently, for every $\hat c$, $s$, the deficit
$f_{\hat c,s}$ of \eqref{eq:f-def} is concave on
$\{\sigma\in\Pos{A\Qsys_s}:\tr\sigma\le1\}$, and
$K_{\hat c}$ is concave on $\Feas$.
\end{corollary}
\claimmeta{\metalabel{Uses}\mref{Lem.}{lem:branch-maps}, \mref{Thm.}{thm:branch-formula}, \extref{../foundations/contract.pdf}{foundations/fact:concavity}
\metasep \metalabel{Used by}\mref{Lem.}{lem:cone-concavity}, \mref{Thm.}{thm:convexity}, \mref{Thm.}{thm:gradient}
\metasep newly derived
\metasep accepted
\metasep \metalabel{Oracles}\coderef{https://github.com/fraunhofersg/certiqs-key/blob/v0.1.0/tests/oracles/registry.py}{registry.py}}

\begin{proof}
$1/\alpha\in(0,1)$, so \xref{foundations/fact:concavity} applies at parameter
$\alpha'=1/\alpha$ and gives concavity of $Q_{1/\alpha}$ (this is
exactly the substitution the source itself performs in its Cor.~27, which
states concavity of $Q_{1/\alpha}\circ\mathcal{G}_{\delta}$ for
$\alpha\in[1,\infty)$). The map $\cG_{\hat c,s}$ is linear and maps the
stated domain into $\Sleq{}$ (Lemma~\ref{lem:branch-maps}\,(i)), and precomposition
of a concave function with a linear map preserves concavity; subtracting
the linear term $\tr\cG_{\hat c,s}(\cdot)$ preserves concavity. For
$K_{\hat c}=P_{\hat c}+\sum_sf_{\hat c,s}$: $P_{\hat c}$ is affine and a
sum of concave functions is concave. The hypothesis check for
\xref{foundations/fact:concavity}: our $\cZ$ is a pinching from the complete
orthogonal family $\{Z_j\}_{j\in\Zset}$ extended by identity factors, which
is again such a family; the arguments $\cG_{\hat c,s}(\sigma)$ are
subnormalized PSD operators, inside the source's stated domain
$\Sleq{Q}$.
\end{proof}

% claim meat-core/lem:cone-concavity
\claimsummary{meat-core/lem:cone-concavity}{The pinched trace-power functional T is positively homogeneous and concave on the whole positive cone. Hence the sector deficits are concave on the cone and the branch quantities K on its sector-block-diagonal part, for every block-diagonal positive measurement family, complete or not.}
\claimbookmark{meat-core/lem:cone-concavity}{Lemma}{lem:cone-concavity}{Concavity on the positive cone}
\claimleadin{%
Corollary~\ref{cor:concavity} gives concavity only on subnormalised
operators $\Sleq{}$, a domain restriction it inherits from
Lemma~\ref{lem:branch-maps}\,(i). Positive homogeneity makes the restriction
dispensable: two positive operators, divided by the larger of their
traces, lie in $\Sleq{}$, and $T$ scales back linearly. Dispensing with it
is what part~(c) needs, and what every use of $T$ off the trace-one slice
needs, the relaxed iterate set of the Frank--Wolfe submodule
(Proposition~\ref{prop:relaxed-set-properties}\,(c)) among them. Part~(c) reads no
completeness relation of $\{F_v\}$, which is what makes it available off
that slice as well. This lemma is the only place where the positive
homogeneity of $T$ is established, and every later use of it, the
objective of Theorem~\ref{thm:firstorder} included, reads it from here.
}
\begin{lemma}[Concavity on the positive cone]\label{lem:cone-concavity}
Let $\alpha\in(1,\infty)$, put $p\coloneqq1/\alpha\in(0,1)$, and let $\cZ$
be a pinching built from a complete family of mutually orthogonal
projectors. On the cone $\Pos{}$ of the space $\cZ$ acts on, consider the
functional
\begin{equation}\label{eq:Tcone}
  T(X)\coloneqq\tr\!\big[\big(\cZ(X^{p})\big)^{\alpha}\big].
\end{equation}
\begin{enumerate}[label=(\alph*),leftmargin=2.4em]
\item $T$ is positively $1$-homogeneous: $T(tX)=t\,T(X)$ for every
$t\ge0$.
\item $T$ is concave on all of $\Pos{}$.
\item Consider a declaration whose components other than $\{F_v\}$
satisfy Definition~\ref{def:admissible-declaration} and whose $\{F_v\}$ is any family of
sector-block-diagonal positive semidefinite operators, complete or not.
Then each deficit $f_{\hat c,s}$ of \eqref{eq:f-def} is concave on all of
$\Pos{A\Qsys_s}$, and $K_{\hat c}$ of \eqref{eq:fK} is concave on the
sector-block-diagonal part of $\Pos{A\Qsys}$, in particular on $\Feas$.
\end{enumerate}
\end{lemma}
\claimmeta{\metalabel{Uses}\mref{Def.}{def:admissible-declaration}, \mref{Def.}{def:maps}, \mref{Lem.}{lem:branch-maps}, \mref{Thm.}{thm:branch-formula}, \mref{Cor.}{cor:concavity}, \extref{../foundations/contract.pdf}{foundations/fact:concavity}
\metasep \metalabel{Used by}\mref{Lem.}{lem:joint-sectors}, \mref{Prop.}{prop:a-neutral}, \mref{Lem.}{lem:choi-factor}, \mref{Thm.}{thm:gradient}, \mref{Lem.}{lem:kraus-freedom}, \mref{Lem.}{lem:objective-homogeneous}, \mref{Prop.}{prop:relaxed-set-properties}
\metasep newly derived
\metasep accepted
\metasep \metalabel{Oracles}\coderef{https://github.com/fraunhofersg/certiqs-key/blob/v0.1.0/tests/oracles/registry.py}{registry.py}}

\begin{proof}
(a) Let $t>0$. First, $(tX)^{p}=t^{p}X^{p}$: the map $X\mapsto X^{p}$ is
defined by the functional calculus of $x\mapsto x^{p}$, and
$(t\lambda)^{p}=t^{p}\lambda^{p}$ on each eigenvalue, the eigenprojectors
being unchanged by the scaling. Second, $\cZ$ is a pinching, hence
linear, so $\cZ\big((tX)^{p}\big)=\cZ(t^{p}X^{p})=t^{p}\cZ(X^{p})$, and
$\tr[(t^{p}Y)^{\alpha}]=t^{p\alpha}\tr[Y^{\alpha}]$ by the same functional
calculus applied to $x\mapsto x^{\alpha}$. Together,
$T(tX)=t^{p\alpha}T(X)=t\,T(X)$ because $p\alpha=1$. At $t=0$ both sides
vanish.

(b) Let $X,Y\in\Pos{}$ and $\lambda\in(0,1)$. If $\tr X=\tr Y=0$ then
$X=Y=0$, a positive semidefinite operator of zero trace being zero, and
the inequality reads $0\ge0$. Otherwise put
$c\coloneqq\max\{\tr X,\tr Y\}>0$, so that $X/c$ and $Y/c$ lie in
$\Sleq{}$ and so does their convex combination. \xref{foundations/fact:concavity}
applies at $\alpha'=p\in(0,1]$, the substitution
Corollary~\ref{cor:concavity} makes, and on the source's own stated
domain $\Sleq{}$, so no hypothesis of the import is stretched; it gives
\begin{equation*}
  T\Big(\lambda\frac{X}{c}+(1-\lambda)\frac{Y}{c}\Big)
  \;\ge\;\lambda\,T\Big(\frac{X}{c}\Big)
  +(1-\lambda)\,T\Big(\frac{Y}{c}\Big).
\end{equation*}
Multiply by $c>0$ and apply (a) once to each of the three arguments:
$T(\lambda X+(1-\lambda)Y)\ge\lambda T(X)+(1-\lambda)T(Y)$. No trace
hypothesis on $X$ or $Y$ is used: the scaling is by their own traces.

(c) $\cG_{\hat c,s}$ is completely positive by Lemma~\ref{lem:branch-maps}\,(i) in
its Kraus-form clause, which needs $F^{(s)}_v\succeq0$ and no completeness
relation; so it is linear and maps $\Pos{A\Qsys_s}$ into $\Pos{}$.
Precomposition of the concave $T$ of (b) with a linear map preserves
concavity wherever the image lies in $\Pos{}$, which is on all of
$\Pos{A\Qsys_s}$, and subtracting the linear
$\sigma\mapsto\tr\cG_{\hat c,s}(\sigma)$ preserves it; by
\eqref{eq:f-def} that is concavity of $f_{\hat c,s}$. For $K_{\hat c}$:
on sector-block-diagonal $\omega$, Lemma~\ref{lem:branch-maps}\,(iii) gives
$\cG_{\hat c}(\omega)=\bigoplus_s\cG_{\hat c,s}(\omega_s)$ with mutually
orthogonal summands, and powers and $\cZ$ act blockwise on that
decomposition, so $T_{\hat c}(\omega)=\sum_sT_{\hat c,s}(\omega_s)$ and
$\tr\cG_{\hat c}(\omega)=\sum_s\tr\cG_{\hat c,s}(\omega_s)$. Inserting
both into \eqref{eq:fK} gives
$K_{\hat c}(\omega)=P_{\hat c}(\omega)+\sum_sf_{\hat c,s}(\omega_s)$,
which is \eqref{eq:c5-sector} obtained from the definitions and
Lemma~\ref{lem:branch-maps}\,(iii) alone, without Theorem~\ref{thm:branch-formula}'s entropy
statement. $P_{\hat c}$ is affine by \eqref{eq:Phimap}, and a sum of
concave functions is concave.
\end{proof}

% claim meat-core/cor:f-nonpos
\claimsummary{meat-core/cor:f-nonpos}{Every sector deficit is non-positive on the whole positive cone, with no trace hypothesis, and vanishes exactly when its passing-branch output is block-diagonal for the key pinching. \hyperlink{claim-meat-core/prop:exact-cancellation}{Proposition~\ref*{prop:exact-cancellation}} reads the equality case.}
\claimbookmark{meat-core/cor:f-nonpos}{Corollary}{cor:f-nonpos}{Non-positivity of the sector deficits}
\claimleadin{%
The corollary is stated on the whole cone. It imposes no trace bound
$\tr\sigma\le1$: such a bound would be inherited from the domain of
Corollary~\ref{cor:concavity} and is used nowhere, and
Proposition~\ref{prop:exact-cancellation} invokes the equality case for
arbitrary $\sigma\in\Pos{A\Qsys_s}$.
}
\begin{corollary}[Non-positivity of the sector deficits]\label{cor:f-nonpos}
For every $\hat c$, $s$ and \emph{every} $\sigma\in\Pos{A\Qsys_s}$, with no
trace hypothesis (in particular not $\tr\sigma\le1$),
$f_{\hat c,s}(\sigma)\le0$, and
$f_{\hat c,s}(\sigma)=0$ if and only if $\cG_{\hat c,s}(\sigma)$ is
block-diagonal for the \termref{term-key-pinching}{key pinching} $\cZ$.
\end{corollary}
\claimmeta{\metalabel{Uses}\mref{Lem.}{lem:branch-maps}, \mref{Thm.}{thm:branch-formula}, \extref{../foundations/contract.pdf}{foundations/lem:pinchmono}
\metasep \metalabel{Used by}\mref{Prop.}{prop:exact-cancellation}, \mref{Lem.}{lem:joint-sectors}, \mref{Thm.}{thm:removal}, \mref{Lem.}{lem:objective-homogeneous}
\metasep newly derived
\metasep implemented (effective: accepted)
\metasep \metalabel{Code}\coderef{https://github.com/fraunhofersg/certiqs-key/blob/v0.1.0/src/certiqs_key/meat_core/objective/branch.py}{certiqs_key.meat_core.objective.branch:evaluate_branch}
\metasep \metalabel{Oracles}\coderef{https://github.com/fraunhofersg/certiqs-key/blob/v0.1.0/tests/oracles/registry.py}{registry.py}}

\begin{proof}
Apply \xref{foundations/lem:pinchmono} with $\tau=\cG_{\hat c,s}(\sigma)$ and
recall \eqref{eq:f-def}. That lemma is stated on all of $\Pos{}$ and
consumes no trace hypothesis, and $\cG_{\hat c,s}$ maps
$\Pos{A\Qsys_s}$ into $\Pos{}$ by Lemma~\ref{lem:branch-maps}\,(i), so the
statement is on the whole cone.
\end{proof}

% claim meat-core/def:keyinert
\claimsummary{meat-core/def:keyinert}{A key-live acceptance value and a sector form a key-inert pair when the raw key is constant on the effective outcomes of each passing announcement class. The predicate is finite and read off the declaration, and it marks the pairs whose deficit cancels exactly (\hyperlink{claim-meat-core/prop:exact-cancellation}{Proposition~\ref*{prop:exact-cancellation}}).}
\claimbookmark{meat-core/def:keyinert}{Definition}{def:keyinert}{Key-inert pair}
\begin{definition}[Key-inert pair]\label{def:keyinert}
For $\hat c\in\Chkey$ and $s\in\SecSet$ define the
\termdef{term-effective-outcome-set}{effective outcome set}
of a \termref{term-pass-class}{pass class} $\kappa\in\Kann_{\mathrm{pass}}(\hat c)$ on sector $s$ as
\begin{equation*}
  E^{(s)}_{\kappa}\coloneqq
  \big\{(x,v)\in a^{-1}(\kappa):\ M_x\ne0,\ F^{(s)}_v\ne0\big\}.
\end{equation*}
The pair $(\hat c,s)$ is
\termdef{term-key-inert}{key-inert}
if for every $\kappa\in\Kann_{\mathrm{pass}}(\hat c)$ the raw-key map $r$ is
constant on $E^{(s)}_{\kappa}$. This is a finite, mechanically checkable
predicate on the declaration.
\end{definition}
\claimmeta{\metalabel{Used by}\mref{Def.}{def:joint-sectors}, \mref{Lem.}{lem:joint-sectors}, \mref{Def.}{def:classes}, \mref{Prop.}{prop:removal-compose}, \mref{Lem.}{lem:certified-bracket}, \mref{Lem.}{lem:certified-bracket-enclosed}
\metasep newly derived
\metasep implemented
\metasep \metalabel{Code}\coderef{https://github.com/fraunhofersg/certiqs-key/blob/v0.1.0/src/certiqs_key/meat_core/objective/branch.py}{certiqs_key.meat_core.objective.branch:is_key_inert}
\metasep \metalabel{Oracles}\coderef{https://github.com/fraunhofersg/certiqs-key/blob/v0.1.0/tests/oracles/registry.py}{registry.py}}

% claim meat-core/prop:exact-cancellation
\claimsummary{meat-core/prop:exact-cancellation}{On a key-inert pair the sector deficit vanishes for every positive input, so the sector contributes nothing to the entropy term. Its variables still enter the outcome probabilities, the statistics and the marginal constraint.}
\claimbookmark{meat-core/prop:exact-cancellation}{Proposition}{prop:exact-cancellation}{Exact cancellation of key-inert sectors}
\claimleadin{%
On a key-inert pair every pass class writes a single raw-key value, so the
branch output never mixes key values and the key pinching leaves it
unchanged; the equality case of Corollary~\ref{cor:f-nonpos} does the rest.
}
\begin{proposition}[Exact cancellation of key-inert sectors]\label{prop:exact-cancellation}
If $(\hat c,s)$ is key-inert, then for \emph{every}
$\sigma\in\Pos{A\Qsys_s}$ the operator $\cG_{\hat c,s}(\sigma)$ is
block-diagonal for the key pinching $\cZ$, and consequently
$f_{\hat c,s}(\sigma)=0$. In particular such sectors contribute nothing to
$K_{\hat c}$ in \eqref{eq:c5-sector}, for every input; they remain
optimization variables, entering through $\Phimap$, $P_{\hat c}$, and the
marginal constraint.
\end{proposition}
\claimmeta{\metalabel{Uses}\mref{Def.}{def:branch}, \mref{Lem.}{lem:branch-maps}, \mref{Thm.}{thm:branch-formula}, \mref{Cor.}{cor:f-nonpos}
\metasep \metalabel{Used by}\mref{Def.}{def:keyactive}, \mref{Def.}{def:classes}, \mref{Thm.}{thm:removal}, \mref{Thm.}{thm:choi}, \mref{Thm.}{thm:oracle}
\metasep newly derived
\metasep implemented (effective: accepted)
\metasep \metalabel{Code}\coderef{https://github.com/fraunhofersg/certiqs-key/blob/v0.1.0/src/certiqs_key/meat_core/declaration/sectors.py}{certiqs_key.meat_core.declaration.sectors:classify_sector}
\metasep \metalabel{Oracles}\coderef{https://github.com/fraunhofersg/certiqs-key/blob/v0.1.0/tests/oracles/registry.py}{registry.py}}

\begin{proof}
Fix a pass class $\kappa$ and let $r_{\kappa}$ be the constant value of $r$
on $E^{(s)}_{\kappa}$. Every term of $K^{(s)}_{\kappa}$
(Definition~\ref{def:branch}) with $(x,v)\notin E^{(s)}_{\kappa}$ vanishes,
because $\sqrt{M_x}\otimes\sqrt{F^{(s)}_v}=0$ there. Hence
$K^{(s)}_{\kappa}=\ket{r_{\kappa}}_{R}\otimes\widetilde K_{\kappa}$ for an
operator $\widetilde K_{\kappa}$ carrying no $R$ factor, and
$K^{(s)}_{\kappa}\sigma K^{(s)\dagger}_{\kappa}
=\dyad{r_{\kappa}}_{R}\otimes
\widetilde K_{\kappa}\sigma\widetilde K_{\kappa}^{\dagger}$ is supported in
a single $R$-eigenspace. A sum of such operators over
$\kappa$ is $R$-block-diagonal (each term lies in one block, possibly
different blocks for different $\kappa$), hence a fixed point of $\cZ$.
By Corollary~\ref{cor:f-nonpos} (equality case), $f_{\hat c,s}(\sigma)=0$.
The last sentence of the statement restates
Lemma~\ref{lem:branch-maps}\,(i),(iii): the sector still carries statistics
weight, so removing it from the \emph{optimization} (rather than from the
objective sum) is a different operation, governed by the sector classes of
Definition~\ref{def:classes} and Theorem~\ref{thm:removal}.
\end{proof}

% claim meat-core/def:keyactive
\claimsummary{meat-core/def:keyactive}{The key-active index set of an acceptance value holds the sectors that are not key-inert with it. Restricting the deficit sum to it leaves the objective exactly equal, not bounded, and removes no state variable.}
\claimbookmark{meat-core/def:keyactive}{Definition}{def:keyactive}{Key-active index set}
\begin{definition}[Key-active index set]\label{def:keyactive}
The
\termdef{term-key-active-index-set}{key-active index set}
of $\hat c$ is
$\SecSet_{K}(\hat c)\coloneqq\{s\in\SecSet:(\hat c,s)\text{ is not
key-inert}\}$. By Proposition~\ref{prop:exact-cancellation},
$K_{\hat c}(\omega)=P_{\hat c}(\omega)
+\sum_{s\in\SecSet_{K}(\hat c)}f_{\hat c,s}(\omega_s)$ for block-diagonal
$\omega$; the objective computed with the sum restricted to
$\SecSet_K(\hat c)$ is \emph{exactly equal} to the full sum, not a bound.
\end{definition}
\claimmeta{\metalabel{Uses}\mref{Prop.}{prop:exact-cancellation}
\metasep \metalabel{Used by}\mref{Def.}{def:joint-sectors}, \mref{Thm.}{thm:choi}, \mref{Req.}{req:structure-operations}, \mref{Lem.}{lem:certified-bracket}
\metasep newly derived
\metasep implemented (effective: accepted)
\metasep \metalabel{Code}\coderef{https://github.com/fraunhofersg/certiqs-key/blob/v0.1.0/src/certiqs_key/meat_core/objective/branch.py}{certiqs_key.meat_core.objective.branch:key_live_values}
\metasep \metalabel{Oracles}\coderef{https://github.com/fraunhofersg/certiqs-key/blob/v0.1.0/tests/oracles/registry.py}{registry.py}}

% claim meat-core/def:joint-sectors
\claimsummary{meat-core/def:joint-sectors}{With Alice sectors, the branch maps, statistics and deficit restrict to each joint sector, and a joint pair is key-inert when the raw key is constant on the outcomes that can occur on it, decided on declared data. With one Alice sector these are the sector notions.}
\claimbookmark{meat-core/def:joint-sectors}{Definition}{def:joint-sectors}{Joint-sector deficits and the joint key-active index set}
\claimleadin{%
Every sector notion of this section is stated on the cone of one measured
sector tensored with all of $A$. Along Alice's declared sectors each of them
restricts to a subcone, and the joint objects below are those restrictions;
nothing about the objective is defined anew.
}
\begin{definition}[Joint-sector deficits and the joint key-active index set]\label{def:joint-sectors}
Let $\Decl$ be a declaration in the general format of
Definition~\ref{def:declaration}, with Alice sectors $t\in\mathbb{T}$ and
\exttermref{../protocol-compiler/contract.pdf}{term-joint-sector}{joint sectors} $(t,s)\in\mathbb{T}\times\SecSet$. Regard
$\Pos{A_t\Qsys_s}$ as the subcone of $\Pos{A\Qsys_s}$ of operators supported
on $\Hs{A_t}\otimes\Hs{\Qsys_s}$. For $\hat c\in\Chatset$ the
\emph{joint-sector} maps, statistics and deficit are the restrictions of the
sector objects of Definition~\ref{def:branch} and of \eqref{eq:f-def} to that
subcone,
\begin{equation}\label{eq:joint-deficit}
  \cG_{\hat c,(t,s)}\coloneqq\cG_{\hat c,s}\big|_{\Pos{A_t\Qsys_s}},\qquad
  P_{\hat c,(t,s)}(\sigma)\coloneqq P_{\hat c}(\sigma),\qquad
  f_{\hat c,(t,s)}(\sigma)\coloneqq f_{\hat c,s}(\sigma)
  \qquad(\sigma\in\Pos{A_t\Qsys_s}),
\end{equation}
so that on such a $\sigma$ the Kraus operators \eqref{eq:kraus} act through
$\sqrt{M^{(t)}_x}\otimes\sqrt{F^{(s)}_v}$. For
$\hat c\in\Chkey$ and a pass class $\kappa\in\Kann_{\mathrm{pass}}(\hat c)$
the \emph{effective outcome set on the joint sector} is
\begin{equation}\label{eq:joint-effective}
  E^{(t,s)}_{\kappa}\coloneqq\big\{(x,v)\in a^{-1}(\kappa):\
  M^{(t)}_x\ne0,\ F^{(s)}_v\ne0\big\}\subseteq E^{(s)}_{\kappa},
\end{equation}
and the pair $(\hat c,(t,s))$ is \emph{key-inert} if $r$ is constant on
$E^{(t,s)}_{\kappa}$ for every $\kappa\in\Kann_{\mathrm{pass}}(\hat c)$. The
\termdef{term-joint-key-active-index-set}{joint key-active index set}
is $\mathbb{J}_K(\hat c)\coloneqq\{(t,s):(\hat c,(t,s))\text{ is not
key-inert}\}$.

Key-inertness of a joint pair is decided on \exttermref{../foundations/contract.pdf}{term-declared-data}{declared data}. In
\exttermref{../protocol-compiler/contract.pdf}{term-povm-form}{POVM form} $M^{(t)}_x\ne0$ is a comparison of the declared block with
the zero block. In \exttermref{../protocol-compiler/contract.pdf}{term-ensemble-form}{ensemble form} it is $N^{(t)}_x\ne0$, by
Lemma~\ref{lem:ensemble-form}\,(iii), again a comparison of a declared block
with zero. For a factorised family
$F^{(s)}_v=\bigotimes_fF^{f,(s_f)}_{v_f}$ is non-zero exactly when every
factor block $F^{f,(s_f)}_{v_f}$ is, a tensor product of matrices vanishing
exactly when one of them does; so the predicate is decided per factor and no
product is formed. With $|\mathbb{T}|=1$ one has
$E^{(t_0,s)}_{\kappa}=E^{(s)}_{\kappa}$, and the pair $(\hat c,(t_0,s))$ is
key-inert exactly when $(\hat c,s)$ is in the sense of
Definition~\ref{def:keyinert}, with
$\mathbb{J}_K(\hat c)=\{t_0\}\times\SecSet_K(\hat c)$.
\end{definition}
\claimmeta{\metalabel{Uses}\mref{Def.}{def:declaration}, \mref{Lem.}{lem:ensemble-form}, \mref{Def.}{def:branch}, \mref{Thm.}{thm:branch-formula}, \mref{Def.}{def:keyinert}, \mref{Def.}{def:keyactive}
\metasep \metalabel{Used by}\mref{Lem.}{lem:joint-sectors}, \mref{Lem.}{lem:choi-factor}, \mref{Thm.}{thm:choi}, \mref{Lem.}{lem:kraus-freedom}, \mref{Lem.}{lem:rank-one-sector}, \mref{Def.}{def:dropped}, \mref{Prop.}{prop:drop-bound}, \mref{Cor.}{cor:drop-program}, \mref{Prop.}{prop:minimal-divergence}, \mref{Lem.}{lem:objective-homogeneous}, \mref{Prop.}{prop:reference-companion}, \mref{Req.}{req:companion-admission}, \mref{Cor.}{cor:fo-dropped}, \mref{Cor.}{cor:fo-reference}, \mref{Lem.}{lem:feasible-side-dual}, \mref{Lem.}{lem:certified-bracket}, \mref{Lem.}{lem:certified-bracket-enclosed}
\metasep newly derived
\metasep implemented (effective: accepted)
\metasep \metalabel{Code}\coderef{https://github.com/fraunhofersg/certiqs-key/blob/v0.1.0/src/certiqs_key/meat_core/objective/joint.py}{certiqs_key.meat_core.objective.joint:is_key_inert}}

% claim meat-core/lem:joint-sectors
\claimsummary{meat-core/lem:joint-sectors}{On states block-diagonal along Alice's and the measured sectors, the branch maps, the statistics and K split per joint sector; each deficit is non-positive, concave and positively homogeneous. Key-inert joint pairs add zero, so the sum may be restricted to the joint key-active index set.}
\claimbookmark{meat-core/lem:joint-sectors}{Lemma}{lem:joint-sectors}{The objective split over joint sectors}
\claimleadin{%
With Alice sectors, the sector form of Theorem~\ref{thm:branch-formula}\,(iii)
refines once more: a state block diagonal along the joint sectors splits the
objective into one deficit per joint sector, and key-inertness is read per
joint pair (Definition~\ref{def:joint-sectors}).
}
\begin{lemma}[The objective split over joint sectors]\label{lem:joint-sectors}
Let $\Decl$ satisfy Definition~\ref{def:admissible-declaration} in the general format of
Definition~\ref{def:declaration}, let $\alpha\in(1,\infty)$ and
$\hat c\in\Chatset$.
\begin{enumerate}[label=(\roman*),leftmargin=2.2em]
\item \emph{(decomposition)} For $s\in\SecSet$ and
$\sigma=\bigoplus_{t}\sigma_t\in\Pos{A\Qsys_s}$ with
$\sigma_t\in\Pos{A_t\Qsys_s}$,
\begin{equation}\label{eq:G-joint}
  \cG_{\hat c,s}(\sigma)=\bigoplus_{t\in\mathbb{T}}
  \cG_{\hat c,(t,s)}(\sigma_t),\qquad
  \Theta_{\hat c,s}(\sigma)=\bigoplus_{t\in\mathbb{T}}
  \Theta_{\hat c,s}(\sigma_t),\qquad
  \Phimap_s[\sigma]=\sum_{t}\Phimap_s[\sigma_t],
\end{equation}
the direct sums along the mutually orthogonal subspaces
$\Hs{R\,XY\,C_AC_B}\otimes\Hs{A_t\Qsys_s}$.
\item \emph{(additivity)} For such $\sigma$,
$f_{\hat c,s}(\sigma)=\sum_tf_{\hat c,(t,s)}(\sigma_t)$. Consequently, for
every $\omega=\bigoplus_{(t,s)}\omega^{(t,s)}$ block diagonal along the joint
sectors,
\begin{equation}\label{eq:K-joint}
  P_{\hat c}(\omega)=\sum_{(t,s)}P_{\hat c,(t,s)}\big(\omega^{(t,s)}\big),
  \qquad
  K_{\hat c}(\omega)=P_{\hat c}(\omega)
  +\sum_{(t,s)}f_{\hat c,(t,s)}\big(\omega^{(t,s)}\big),
\end{equation}
and for every $\sigma\in\Pos{A_t\Qsys_s}$: $f_{\hat c,(t,s)}(\sigma)\le0$,
$0\le P_{\hat c,(t,s)}(\sigma)+f_{\hat c,(t,s)}(\sigma)$, and
$f_{\hat c,(t,s)}$ is concave and positively $1$-homogeneous on
$\Pos{A_t\Qsys_s}$.
\item \emph{(key-inert joint pairs)} If $\hat c\in\Chkey$ and
$(\hat c,(t,s))$ is key-inert (Definition~\ref{def:joint-sectors}), then
$f_{\hat c,(t,s)}(\sigma)=0$ for \emph{every} $\sigma\in\Pos{A_t\Qsys_s}$.
Hence the second sum of \eqref{eq:K-joint} may be restricted to
$(t,s)\in\mathbb{J}_K(\hat c)$, as an identity and not as a bound.
\item \emph{(relation to the sector notions)} If $(\hat c,s)$ is key-inert in
the sense of Definition~\ref{def:keyinert} then $(\hat c,(t,s))$ is key-inert
for every $t$; the converse fails in general, and holds when
$|\mathbb{T}|=1$.
\end{enumerate}
\end{lemma}
\claimmeta{\metalabel{Uses}\mref{Def.}{def:declaration}, \mref{Def.}{def:admissible-declaration}, \mref{Def.}{def:maps}, \mref{Def.}{def:branch}, \mref{Lem.}{lem:branch-maps}, \mref{Thm.}{thm:branch-formula}, \mref{Lem.}{lem:cone-concavity}, \mref{Cor.}{cor:f-nonpos}, \mref{Def.}{def:keyinert}, \mref{Def.}{def:joint-sectors}
\metasep \metalabel{Used by}\mref{Thm.}{thm:choi}, \mref{Prop.}{prop:drop-bound}, \mref{Cor.}{cor:drop-program}, \mref{Lem.}{lem:objective-homogeneous}, \mref{Prop.}{prop:relaxed-set-properties}, \mref{Thm.}{thm:firstorder}, \mref{Cor.}{cor:fo-dropped}, \mref{Lem.}{lem:certified-bracket}
\metasep newly derived
\metasep implemented (effective: accepted)
\metasep \metalabel{Code}\coderef{https://github.com/fraunhofersg/certiqs-key/blob/v0.1.0/src/certiqs_key/meat_core/objective/joint.py}{certiqs_key.meat_core.objective.joint:evaluate_joint_branch}}

\begin{proof}
(i) By Definition~\ref{def:admissible-declaration}\,(iii) in the general format every $M_x$ is
block diagonal along $\mathbb{T}$, and by blockwise functional calculus so
is $\sqrt{M_x}=\bigoplus_t\sqrt{M^{(t)}_x}$. Write
$\widetilde\Pi_t$ for $\Pi^{A}_t$ acting on the $A$ factor of
$\Hs{R\,XY\,C_AC_B\,A\Qsys_s}$. Each term of $K^{(s)}_{\kappa}$ in
\eqref{eq:kraus} carries the operator $\sqrt{M_x}\otimes\sqrt{F^{(s)}_v}$ on
$A\Qsys_s$, which commutes with $\Pi^{A}_t\otimes\idon{}$, so
$K^{(s)}_{\kappa}(\Pi^{A}_t\otimes\idon{})=\widetilde\Pi_tK^{(s)}_{\kappa}$.
Since $\sigma=\sum_t(\Pi^{A}_t\otimes\idon{})\sigma(\Pi^{A}_t\otimes\idon{})$
and $\sigma_t=(\Pi^{A}_t\otimes\idon{})\sigma(\Pi^{A}_t\otimes\idon{})$,
\begin{equation*}
  K^{(s)}_{\kappa}\,\sigma\,K^{(s)\dagger}_{\kappa}
  =\sum_t\widetilde\Pi_t\,K^{(s)}_{\kappa}\,\sigma_t\,
  K^{(s)\dagger}_{\kappa}\,\widetilde\Pi_t ,
\end{equation*}
and the $t$-th term is $K^{(s)}_{\kappa}\sigma_tK^{(s)\dagger}_{\kappa}$,
supported on the range of $\widetilde\Pi_t$. Summing over the pass, resp.\
fail, classes of $\phi^{-1}(\hat c)$ gives the two direct sums; the
statistics statement is linearity of \eqref{eq:Phimap} in its argument.

(ii) The subspaces of (i) are invariant under $\cZ$, which acts on $R$
alone, and under powers; so for an operator
$X=\bigoplus_tX_t$ along them, $X^{1/\alpha}=\bigoplus_tX_t^{1/\alpha}$,
$\cZ(X^{1/\alpha})=\bigoplus_t\cZ(X_t^{1/\alpha})$, the $\alpha$-th power is
again blockwise, and the trace adds. With $X=\cG_{\hat c,s}(\sigma)$ and (i),
\eqref{eq:f-def} gives $f_{\hat c,s}(\sigma)=\sum_tf_{\hat c,s}(\sigma_t)$,
which is the claim by \eqref{eq:joint-deficit}. For $\omega$ block diagonal
along the joint sectors, put $\omega_s\coloneqq\bigoplus_t\omega^{(t,s)}$;
then $\omega=\bigoplus_s\omega_s$ is block diagonal along $\SecSet$,
Lemma~\ref{lem:branch-maps}\,(iii) gives
$\cG_{\hat c}(\omega)=\bigoplus_s\cG_{\hat c,s}(\omega_s)$ and
$\Phimap[\omega]=\sum_s\Phimap_s[\omega_s]$, and by the same blockwise
argument along $\SecSet$ together with \eqref{eq:fK},
$K_{\hat c}(\omega)=P_{\hat c}(\omega)+\sum_sf_{\hat c,s}(\omega_s)$; insert
the additivity just proved, and the additivity of $\Phimap_s$ from (i), to
obtain \eqref{eq:K-joint}. For the bounds let $\sigma\in\Pos{A_t\Qsys_s}$.
$f_{\hat c,(t,s)}(\sigma)=f_{\hat c,s}(\sigma)\le0$ by
Corollary~\ref{cor:f-nonpos}, stated on all of $\Pos{A\Qsys_s}$. Regard
$\sigma$ as an element $\varrho$ of $\Pos{A\Qsys}$ supported on one joint
sector: then $\cG_{\hat c}(\varrho)=\cG_{\hat c,s}(\sigma)$ by
Lemma~\ref{lem:branch-maps}\,(iii), so
$K_{\hat c}(\varrho)=P_{\hat c}(\sigma)+f_{\hat c,s}(\sigma)$ by
\eqref{eq:fK} and \eqref{eq:f-def}, and the cone clause of
Theorem~\ref{thm:branch-formula} gives $K_{\hat c}(\varrho)\ge0$. Concavity on
$\Pos{A_t\Qsys_s}$ is the restriction to a convex subcone of the concavity
of $f_{\hat c,s}$ on $\Pos{A\Qsys_s}$, Lemma~\ref{lem:cone-concavity}\,(c);
positive $1$-homogeneity is part~(a) of that lemma applied to
$T\circ\cG_{\hat c,s}$, the map $\cG_{\hat c,s}$ and the trace being linear.

(iii) Fix a pass class $\kappa$ and let $r_{\kappa}$ be the constant value
of $r$ on $E^{(t,s)}_{\kappa}$. For $\sigma\in\Pos{A_t\Qsys_s}$ one has
$\sigma=(\Pi^{A}_t\otimes\idon{})\sigma(\Pi^{A}_t\otimes\idon{})$, so only
$K^{(s)}_{\kappa}(\Pi^{A}_t\otimes\idon{})$ enters
$K^{(s)}_{\kappa}\sigma K^{(s)\dagger}_{\kappa}$, and its term for $(x,v)$
carries $\sqrt{M^{(t)}_x}\otimes\sqrt{F^{(s)}_v}$, which vanishes unless
$(x,v)\in E^{(t,s)}_{\kappa}$. Hence
$K^{(s)}_{\kappa}(\Pi^{A}_t\otimes\idon{})
=\ket{r_{\kappa}}_R\otimes\widetilde K_{\kappa}$ with $\widetilde K_{\kappa}$
carrying no $R$ factor, and
$K^{(s)}_{\kappa}\sigma K^{(s)\dagger}_{\kappa}
=\dyad{r_{\kappa}}_R\otimes\widetilde K_{\kappa}\sigma
\widetilde K_{\kappa}^{\dagger}$ lies in one eigenspace of the key pinching.
A sum of such operators over $\kappa$ is a fixed point of $\cZ$, so
$\cG_{\hat c,(t,s)}(\sigma)$ is block diagonal for $\cZ$ and
$f_{\hat c,(t,s)}(\sigma)=0$ by the equality case of
Corollary~\ref{cor:f-nonpos}. The restriction of the sum follows.

(iv) $E^{(t,s)}_{\kappa}\subseteq E^{(s)}_{\kappa}$, since $M^{(t)}_x\ne0$
implies $M_x\ne0$; a map constant on a set is constant on a subset. For the
failure of the converse take two Alice sectors, $\Xset=\{0,1\}$ with
$M_0=\Pi^{A}_{1}$, $M_1=\Pi^{A}_{2}$, one announcement class, and
$r(x,v)=x$: on each joint sector one $x$ occurs and $r$ is constant, on the
$\Qsys$ sector both occur and it is not. With $|\mathbb{T}|=1$ the two sets
coincide.
\end{proof}

\section{The program and its reduction}\label{group-meat-core/def:feasible}
\sectionintro{The feasible set, the objective and its infimum h\_low over the acceptance set form the finite convex program of this module; at d\_A = 1 the two source modes give the same program (\hyperlink{claim-meat-core/prop:trivial-alice-system}{Proposition~\ref*{prop:trivial-alice-system}}), and untrusted-source mode never gives a larger value. The section's main result, \hyperlink{claim-meat-core/thm:reduction}{Theorem~\ref*{thm:reduction}}, shows that h\_low bounds the single-round rate from below for an unconstrained declaration in either source mode, every step being an equality or conservative, and \hyperlink{claim-meat-core/cor:reduction-constrained}{Corollary~\ref*{cor:reduction-constrained}} extends it to arbitrary constraint lists. Restricting the variable along Alice's sectors changes no value (\hyperlink{claim-meat-core/prop:a-neutral}{Proposition~\ref*{prop:a-neutral}}). The value lies in [0, +\ensuremath{\infty}], each entropy term being at most log\ensuremath{_{2}}(d\_Z + 1) while the divergence penalty is unbounded (\hyperlink{claim-meat-core/prop:value-range}{Proposition~\ref*{prop:value-range}}), and an infinite value means that no attack the declared list admits is accepted (\hyperlink{claim-meat-core/rem:hlow-range}{Remark~\ref*{rem:hlow-range}}). The objective is jointly convex and lower semicontinuous, so a finite infimum is attained (\hyperlink{claim-meat-core/thm:convexity}{Theorem~\ref*{thm:convexity}}).}
\sectionimports{\extref{../foundations/contract.pdf}{foundations/def:kl}, \extref{../foundations/contract.pdf}{foundations/def:pinching}, \extref{../foundations/contract.pdf}{foundations/def:renyi}, \extref{../foundations/contract.pdf}{foundations/lem:kernel}, \extref{../protocol-compiler/contract.pdf}{protocol-compiler/def:feasible}}

% claim meat-core/def:feasible
\claimsummary{meat-core/def:feasible}{The program minimises an entropy objective over the compiler's feasible set (\xref{protocol-compiler/def:feasible}), read as normalised block-diagonal states meeting the declared marginal and constraints. Its infimum over the acceptance set, h\_low, is the value a key length is computed from.}
\claimbookmark{meat-core/def:feasible}{Definition}{def:feasible}{Feasible set and computable objective}
\claimleadin{%
The set over which this module's program minimises is the compiler's
\exttermref{../protocol-compiler/contract.pdf}{term-feasible-set}{feasible set} (\xref{protocol-compiler/def:feasible}); the objective and
$\hlow$ are this module's alone.
}
\begin{definition}[Feasible set and computable objective]\label{def:feasible}
For a declaration $\Decl$, $\Feas$ is the \exttermref{../protocol-compiler/contract.pdf}{term-feasible-set}{feasible set} of
\xref{protocol-compiler/def:feasible}. Read as a set
of operators on $A\Qsys$, a list of sector blocks being the block-diagonal
operator it denotes, it is
\begin{equation}\label{eq:feas}
  \Feas\coloneqq\Big\{\!\omega\in\Pos{A\Qsys}:\
  (\idon{A}\otimes\Pi_s)\,\omega\,(\idon{A}\otimes\Pi_{s'})=0
  \ \forall s\ne s',\quad
  \tr\omega=1,\quad
  \omega\models\Decl,\quad
  \omega\models\CList\Big\},
\end{equation}
writing $\omega_s$ for the sector blocks, with $\omega\models\Decl$ the
marginal condition of \eqref{eq:margcond} and $\omega\models\CList$ the
constraint conditions in the form
Definition~\ref{def:homog} gives them for the declared mode. Written
out, $\Feas$ is
\begin{equation}\label{eq:feas-modes}
\begin{aligned}
  \text{trusted:}\quad&
  \omega_A=\tau_A,\quad
  \tr_{\Qsys}[(\idon{A}\otimes\hat O_k)\,\omega]\bowtie_k\Theta_k\
  \ \forall k;\\
  \text{untrusted:}\quad&
  \text{(no marginal condition)},\quad
  \tr[(\idon{A}\otimes\hat O_k)\,\omega]\bowtie_k c_k\ \ \forall k .
\end{aligned}
\end{equation}
The normalization $\tr\omega=1$ is written \emph{explicitly} in
\eqref{eq:feas}, and is not derived from anything: in
\exttermref{../transmitter/contract.pdf}{term-trusted-source-mode}{trusted-source mode} it also follows from $\omega_A=\tau_A$ (whence
$\tr\omega=\tr\tau_A=1$), but that derivation is unavailable in
\exttermref{../transmitter/contract.pdf}{term-untrusted-source-mode}{untrusted-source mode}, where normalization is the \emph{only}
unconditional constraint the feasible set carries. Leaving it implied
would empty the untrusted-mode program of its one defining condition.
$\Feas$ is convex and compact in either mode and in either format
(\xref{protocol-compiler/def:feasible}). For
$\alpha\in(1,\infty)$,
$\mathbf q$ a distribution on $\Chatset$ and $\omega\in\Feas$, define the
\termdef{term-objective}{objective}
\begin{equation}\label{eq:objective}
  O_{\alpha}(\mathbf q,\omega)\coloneqq
  \frac{\alpha}{\alpha-1}\sum_{\hat c\notin\Chkey}
  q(\hat c)\log_2\frac{q(\hat c)}{P_{\hat c}(\omega)}
  +\sum_{\hat c\in\Chkey}\left[
  q(\hat c)\log_2\frac{q(\hat c)}{P_{\hat c}(\omega)}
  +\frac{q(\hat c)}{\alpha-1}\log_2\frac{q(\hat c)}{K_{\hat c}(\omega)}
  \right],
\end{equation}
with the conventions of \xref{foundations/def:kl} at vanishing arguments
($0\log_2\frac0v=0$; $u\log_2\frac u0=+\infty$ for $u>0$), and
\begin{equation}\label{eq:hlow-def}
  \hlow\coloneqq\inf_{\mathbf q\in\acSet}\ \inf_{\omega\in\Feas}\
  O_{\alpha}(\mathbf q,\omega).
\end{equation}

\emph{The general format.} For a declaration in the general format of
Definition~\ref{def:declaration}, $\Feas$ is the set of lists over the
\exttermref{../protocol-compiler/contract.pdf}{term-joint-sector}{joint sectors} of \xref{protocol-compiler/def:feasible},
$\omega=\bigoplus_{(t,s)}\omega^{(t,s)}$, with the $s$-block
$\hat O^{\uparrow(s)}_k$ of the \exttermref{../protocol-compiler/contract.pdf}{term-lifted-operator}{lifted operator} of an entry scoped to $f_k$
and, for a party-local entry scoped to $f$, the split $s=(s_{f^c},s_f)$ and
the product $\Qsys_{f^c,s_{f^c}}$ of the other factors' blocks written
there, with every block positive semidefinite; its conditions read, block by
block,
\begin{equation}\label{eq:feas-general}
\begin{aligned}
  &\textstyle\sum_{(t,s)}\tr\omega^{(t,s)}=1
  &&\text{(normalization)};\\
  &\textstyle\sum_{s}\tr_{\Qsys_s}\omega^{(t,s)}=\tau^{(t)}_A
  \quad\forall t
  &&\text{(marginal, trusted-source mode)};\\
  &\textstyle\sum_{s}\tr_{\Qsys_s}\!\big[(\idon{A_t}\otimes
  \hat O^{\uparrow(s)}_k)\,\omega^{(t,s)}\big]\bowtie_k\Theta^{(t)}_k
  \quad\forall t
  &&\text{(trusted; }\Theta^{(t)}_k=c_k\tau^{(t)}_A
  \text{ if homogeneous)};\\
  &\textstyle\sum_{(t,s)}\tr\!\big[(\idon{A_t}\otimes
  \hat O^{\uparrow(s)}_k)\,\omega^{(t,s)}\big]\bowtie_kc_k
  &&\text{(scalar entry, untrusted-source mode)};\\
  &\textstyle\sum_{s_f\in\SecSet^{f}}\tr_{\Qsys_{f,s_f}}\!\Big[\big(\idon{}
  \otimes(\hat O^{(s_f)}_k-c_k\idon{})\big)\,
  \omega^{(t,(s_{f^c},s_f))}\Big]\bowtie_k0
  \quad\forall(t,s_{f^c})
  &&\text{(party-local, either mode)},
\end{aligned}
\end{equation}
the last as an operator on $\Hs{A_t}\otimes\Hs{\Qsys_{f^c,s_{f^c}}}$. The
objective \eqref{eq:objective} and $\hlow$ of \eqref{eq:hlow-def} are
defined on this $\Feas$ by the same formulas.
\end{definition}
\claimmeta{\metalabel{Uses}\mref{Def.}{def:declaration}, \mref{Def.}{def:homog}, \mref{Def.}{def:hqkd-list}, \extref{../foundations/contract.pdf}{foundations/def:kl}, \extref{../protocol-compiler/contract.pdf}{protocol-compiler/def:feasible}
\metasep \metalabel{Used by}\mref{Prop.}{prop:trivial-alice-system}, \mref{Prop.}{prop:a-neutral}, \mref{Cor.}{cor:reduction-constrained}, \mref{Prop.}{prop:value-range}, \mref{Thm.}{thm:convexity}, \mref{Def.}{def:classes}, \mref{Thm.}{thm:removal}, \mref{Prop.}{prop:removal-compose}, \mref{Rem.}{rem:classes}, \mref{Lem.}{lem:admissible-alice-subspace}, \mref{Lem.}{lem:choi-factor}, \mref{Thm.}{thm:choi}, \mref{Prop.}{prop:frame-free-rows}, \mref{Rem.}{rem:basis-invariance}, \mref{Def.}{def:untrusted-choi}, \mref{Lem.}{lem:estimate}, \mref{Lem.}{lem:b-derivative}, \mref{Thm.}{thm:oracle}, \mref{Lem.}{lem:rank-one-sector}, \mref{Def.}{def:dropped}, \mref{Prop.}{prop:drop-bound}, \mref{Cor.}{cor:drop-program}, \mref{Prop.}{prop:minimal-divergence}, \mref{Lem.}{lem:objective-homogeneous}, \mref{Thm.}{thm:keylength}, \mref{Prop.}{prop:model-pass-bound}, \mref{Rem.}{rem:cgate}, \mref{Prop.}{prop:optol-nobound}, \mref{Rem.}{rem:optol-direction}, \mref{Prop.}{prop:reference-companion}, \mref{Req.}{req:objective-alpha-domain}, \mref{Def.}{def:relaxed-feasible-set}, \mref{Prop.}{prop:relaxed-set-properties}, \mref{Thm.}{thm:firstorder}, \mref{Cor.}{cor:fo-dropped}, \mref{Lem.}{lem:feasible-side-dual}
\metasep newly derived
\metasep implemented (effective: accepted)
\metasep \metalabel{Code}\coderef{https://github.com/fraunhofersg/certiqs-key/blob/v0.1.0/src/certiqs_key/meat_core/objective/objective.py}{certiqs_key.meat_core.objective.objective:check_feasible}
\metasep \metalabel{Oracles}\coderef{https://github.com/fraunhofersg/certiqs-key/blob/v0.1.0/tests/oracles/registry.py}{registry.py}}

% claim meat-core/prop:trivial-alice-system
\claimsummary{meat-core/prop:trivial-alice-system}{At d\_A = 1 the two source modes have the same feasible set, the n-round product-marginal condition holds for every normalised state, the two licensed constraint classes coincide and the source licence deletes no attack. The two modes are then one declaration.}
\claimbookmark{meat-core/prop:trivial-alice-system}{Proposition}{prop:trivial-alice-system}{The two source modes coincide at a trivial Alice system}
\claimleadin{%
When Alice's system is one-dimensional, $d_A=\dim\Hs{A}=1$, it admits a
single state, and the two values of the \exttermref{../protocol-compiler/contract.pdf}{term-marginal-slot}{marginal slot} differ only in
name. The proposition says so for four objects the source mode touches: the
feasible set, the product-marginal condition of
Lemma~\ref{lem:product-marginal}, the licensed constraint classes and
Assumption~A-SRC; for the Choi parametrisation it is
Definition~\ref{def:untrusted-choi}\,(d). A condition that every
normalized state satisfies is satisfied vacuously, which is not the same as
violated; Remark~\ref{rem:modes-da1}\,(b) reads the proposition as an
embedding of protocols with no trusted source.
}
\begin{proposition}[The two source modes coincide at a trivial Alice system]\label{prop:trivial-alice-system}
Let $d_A=1$. Then trusted-source mode admits exactly one marginal,
$\tau_A=[1]$, and:
\begin{enumerate}[label=(\roman*),leftmargin=2.2em]
\item the two feasible sets of \eqref{eq:feas} coincide;
\item the $n$-round condition $\omega_{A_1^n}=\tau_A^{\otimes n}$ of
Lemma~\ref{lem:product-marginal}\,(ii) holds for every normalized state; it is
satisfied vacuously;
\item the licensed constraint classes of Definition~\ref{def:homog}
coincide: trusted-mode $\tau_A$-homogeneity is exactly the scalar typing
of untrusted-source mode;
\item A-SRC (\xref{transmitter/ass:src}) is vacuous: what it
licenses is the condition $\tr\omega=1$, satisfied by every normalized
state, so it deletes no attack.
\end{enumerate}
\end{proposition}
\claimmeta{\metalabel{Uses}\mref{Def.}{def:admissible-declaration}, \mref{Def.}{def:homog}, \mref{Def.}{def:feasible}
\metasep \metalabel{Used by}\mref{Prop.}{prop:retention}, \mref{Def.}{def:untrusted-choi}, \mref{Rem.}{rem:cgate}, \mref{Prop.}{prop:optol-nobound}, \mref{Req.}{req:marginal-licence-gate}
\metasep newly derived
\metasep accepted}

\begin{proof}
Since $d_A=1$, $\Seq{A}=\{[1]\}$, so trusted-source mode admits exactly one
marginal, $\tau_A=[1]$. For $\omega\in\Pos{A\Qsys}$ the $A$-marginal is the
$1\times1$ matrix $\tr\omega$, so $\omega_A=\tau_A$ and $\tr\omega=1$ are
the same condition.

(i) The marginal lines of \eqref{eq:feas-modes} agree: that is the
equality of conditions just shown. The constraint lines agree as well: at
$d_A=1$ the partial trace $\tr_{\Qsys}$ of an operator on $A\Qsys$ is its
trace, a $1\times1$ Hermitian $\Theta_k$ is $c_k[1]=c_k\tau_A$ with
$c_k\in\mathbb{R}$ its one entry, and $\preceq$, $=$, $\succeq$ between
$1\times1$ matrices are $\le$, $=$, $\ge$, so the trusted line
$\tr_{\Qsys}[(\idon{A}\otimes\hat O_k)\,\omega]\bowtie_k\Theta_k$ is the
untrusted line $\tr[(\idon{A}\otimes\hat O_k)\,\omega]\bowtie_kc_k$. In
the general format the lines of \eqref{eq:feas-general} agree in the same
way, the \exttermref{../protocol-compiler/contract.pdf}{term-party-local}{party-local} line being the same in either mode.

(ii) $A_1^n$ is one-dimensional, so $\Seq{A_1^n}$ is a singleton and the
$A_1^n$-marginal of every normalized state is its one element
$\tau_A^{\otimes n}$.

(iii) In trusted-source mode an untagged entry of a licensed list is
declared homogeneous, carrying $c_k\in\mathbb{R}$, and its right-hand side
is constructed as $\Theta_k=c_k\tau_A=c_k[1]$
(Definition~\ref{def:homog}\,(a)); in untrusted-source mode an untagged
entry of a licensed list carries a real $c_k$
(Definition~\ref{def:homog}\,(b)). So the two classes consist of the same
lists of entries carrying a real scalar $c_k$, and by (i) the line of such
an entry is in both modes the one comparison of real numbers
$\tr[(\idon{A}\otimes\hat O_k)\,\omega]\bowtie_kc_k$. An entry declared
\exttermref{../protocol-compiler/contract.pdf}{term-general}{general} lies in neither class: its $1\times1$ Hermitian
$\Theta_k$ equals $c_k\tau_A$ with $c_k=\Theta_k\in\mathbb{R}$, so it
could have been declared homogeneous with the same line, but licensing
reads the declared form (Definition~\ref{def:homog}).

(iv) What Assumption~A-SRC licenses is the marginal condition
$\omega_A=\tau_A$, which by the equality of conditions above is
$\tr\omega=1$; every normalized state satisfies it.
\end{proof}

% claim meat-core/rem:modes-da1
\claimsummary{meat-core/rem:modes-da1}{A declared marginal deletes every feasible state with another marginal, so untrusted-source mode is always available and never gives a larger value. At d\_A = 1 the two modes coincide, which lets a protocol with no trusted source be declared there; \ensuremath{\tau}\_A = 0 is inadmissible.}
\claimbookmark{meat-core/rem:modes-da1}{Remark}{rem:modes-da1}{The two source modes on the program}
\begin{remark}[The two source modes on the program]\label{rem:modes-da1}
\emph{(a) What each mode asserts} is the transmitter's statement
(\xref{transmitter/rem:source-mode}\,(a): in trusted-source mode the
declaration asserts, through A-SRC
(\xref{transmitter/ass:src}), that Alice's register $A$ never passes
through the adversary's hands, so that $\omega_A=\tau_A$ in every round;
in untrusted-source mode it asserts nothing about $A$). What that does to
this module's program is stated here. In untrusted-source mode
($\tau_A=\nomarg$) the adversary prepares $A\Qsys$ jointly
(Definition~\ref{def:roundmodel}), the \exttermref{../protocol-compiler/contract.pdf}{term-feasible-set}{feasible set} carries only
$\tr\omega=1$ (Eq.~\eqref{eq:feas}), and Lemma~\ref{lem:product-marginal} is
not invoked at all. A declared $\tau_A$ \emph{deletes} from the feasible set
every state whose $A$-marginal differs, hence every attack that would have
produced one. That is a different \emph{kind} of input from an entry of
$\CList$, and a more demanding one to justify: an entry of $\CList$
restricts what Eve's channel may do to a register she holds, whereas
$\tau_A$ asserts that a register never reached her at all. Declaring one
without the attestation of Assumption~A-SRC overstates the key in the
unsafe direction. Untrusted-source mode is always available and always
more conservative. Precisely: fix the declaration data
$(\hat O_k,\bowtie_k,c_k)$ and compare the feasible set \eqref{eq:feas}
with $\tau_A$ declared against the one with $\tau_A=\nomarg$. Every
$\omega$ feasible in trusted-source mode is feasible in untrusted-source
mode --- from $\omega_A=\tau_A$ we get $\tr\omega=\tr\tau_A=1$, and taking
the trace of the operator entry
$\tr_{\Qsys}[(\idon{A}\otimes\hat O_k)\,\omega]\bowtie_k c_k\tau_A$
gives the scalar entry
$\tr[(\idon{A}\otimes\hat O_k)\,\omega]\bowtie_k c_k$ (using
$\tr\tau_A=1$, and that $X\succeq Y$ implies $\tr X\ge\tr Y$). Hence
the untrusted feasible set is the larger, its infimum the smaller, and
$\hlow$ computed in untrusted-source mode never exceeds the
trusted-mode value. Dropping the marginal therefore loses rate; it
never gains any, and it can never turn an unsound declaration into a
sound-looking one.

\emph{(b) At $d_A=1$ the two modes are the same declaration}
(Proposition~\ref{prop:trivial-alice-system}). This is the
\emph{embedding}: a protocol with no trusted source may be declared at
$d_A=1$ by folding the system the sender measures into the sectors of
$\Qsys$ and her outcome components into $v$, with the
\exttermref{../protocol-compiler/contract.pdf}{term-party-partition}{party partition} naming which components she holds. The
empty-marginal case is therefore the \emph{safe} direction --- it removes
an assumption rather than adding one.

The equivalence is also the accumulation paper's own reading of its
channel conditional entropy: the sentence immediately following
\cite[Eq.~(1)]{arqand2025meat} records that the definitions of earlier
works are recovered ``by setting the $C$ register to be trivial and
omitting the marginal constraint (equivalently, setting $A$ to be
trivial)''. That sentence is about the source's Eq.~(1), not about its
Theorem~4.2a; the derivation of Proposition~\ref{prop:trivial-alice-system}
is this module's own and does not rest on it. The corresponding check on the
import side is Fact~\ref{fact:meat}, locator note~(vii).

\emph{(c) $\tau_A=0$ is inadmissible}, for the two reasons
\xref{transmitter/rem:source-mode}\,(c) gives; on this module's side
Definition~\ref{def:admissible-declaration}\,(ii) refuses it outright, and were the
normalization dropped, $\omega_A=0$ would give $\tr\omega=0$ and hence
$\omega=0$, so $\Feas=\varnothing$ in violation of
Definition~\ref{def:admissible-declaration}\,(vii) and \eqref{eq:hlow-def} would be an
infimum over the empty set. The absent marginal is $\tau_A=\nomarg$; the
trivial marginal is $d_A=1$ with $\tau_A=[1]$; $\tau_A=0$ is neither.
\end{remark}
\claimmeta{\metalabel{Uses}\mref{Def.}{def:admissible-declaration}
\metasep \metalabel{Assumes}\extaref{../transmitter/contract.pdf}{transmitter/ass:src}{A-SRC}
\metasep newly derived
\metasep accepted}

% claim meat-core/thm:reduction
\claimsummary{meat-core/thm:reduction}{For a declaration that adds no constraints on the state (an empty constraint list), in either source mode, the value h\_low of the finite convex program bounds the single-round rate from below. Every step of the reduction is an equality or lowers the bound, so the computable program is conservative.}
\claimbookmark{meat-core/thm:reduction}{Theorem}{thm:reduction}{Reduction of the rate to the convex program}
\claimleadin{%
Step~1 of the proof is where the block restriction is discharged. The
compiler states, as a consumer obligation it does not discharge, that this
module's minimisation over the sector-block-diagonal feasible set is sound
(\xref{protocol-compiler/rem:block-restriction}); Step~1, the dephasing of an arbitrary
feasible state by the \emph{joint} pinching along the measured sectors and
along Alice's declared sectors, is that statement's proof. It is proved under
exactly the typing of Definition~\ref{def:declaration}\,(DX2), which types
Alice's measurement slot, $\tau_A$ and every $\Theta_k$ against a declared
factor. The restriction along $\mathbb{T}$ is licensed by this step, and
Proposition~\ref{prop:a-neutral} adds that it costs nothing. Corollary~\ref{cor:reduction-constrained} extends the theorem to
non-empty constraint lists.
}
\begin{mainclaim}
\begin{theorem}[Reduction of the rate to the convex program]\label{thm:reduction}
Let $\Decl$ satisfy Definition~\ref{def:admissible-declaration} with $\CList=\varnothing$,
in either source mode, and let $\alpha\in(1,\infty)$. Then, with
$\hqkd=\hqkd(\varnothing)$ the single-round rate of
Definition~\ref{def:hqkd-list} at the empty list,
\begin{equation}\label{eq:c8}
  \hqkd\;\ge\;\hlow .
\end{equation}
\end{theorem}
\end{mainclaim}
\claimmeta{\metalabel{Uses}\mref{Def.}{def:declaration}, \mref{Def.}{def:admissible-declaration}, \mref{Thm.}{thm:protocol-channel}, \mref{Def.}{def:hqkd-list}, \mref{Lem.}{lem:cond-more}, \mref{Thm.}{thm:branch-formula}, \extref{../foundations/contract.pdf}{foundations/def:renyi}
\metasep \metalabel{Used by}\mref{Cor.}{cor:reduction-constrained}, \mref{Rem.}{rem:losses}, \mref{Prop.}{prop:value-range}, \mref{Rem.}{rem:hlow-range}, \mref{Thm.}{thm:oracle}, \mref{Thm.}{thm:keylength}, \mref{Prop.}{prop:relaxed-set-properties}
\metasep newly derived
\metasep implemented (effective: accepted)
\metasep \metalabel{Code}\coderef{https://github.com/fraunhofersg/certiqs-key/blob/v0.1.0/src/certiqs_key/meat_core/objective/objective.py}{certiqs_key.meat_core.objective.objective:objective_value}
\metasep \metalabel{Oracles}\coderef{https://github.com/fraunhofersg/certiqs-key/blob/v0.1.0/tests/oracles/registry.py}{registry.py}}

\begin{proof}
\emph{Idea.} Five steps take the bracket of \eqref{eq:hqkd-list} at
$\CList=\varnothing$, the rate that Theorem~\ref{thm:accumulated-entropy}
writes out as \eqref{eq:hqkd-def}, at a state of its inner infimum to the objective at a point of $\Feas$: a joint dephasing that
changes no output (Step~1), a purification (Step~2), the replacement of the
up-arrow entropy by the down-arrow one (Step~3), the branch formula of
Theorem~\ref{thm:branch-formula} (Step~4) and a rearrangement of the
logarithms (Step~5). Every step of the proof is an equality or an inequality
whose direction lowers the bound (conservative); no step uses $\CList$.

Fix $\mathbf q\in\acSet$ and a feasible $\omega\in\Seq{A\Qsys\Etil}$ with
$\omega\models\Decl$ (\eqref{eq:margcond}); abbreviate
$\Phi\coloneqq\Phimap[\omega_{A\Qsys}]$. We
lower-bound the bracket of \eqref{eq:hqkd-list} at $\CList=\varnothing$ by
$O_{\alpha}(\mathbf q,\tilde\omega)$ for some $\tilde\omega\in\Feas$; taking
infima then yields \eqref{eq:c8}. The divergence term of
\eqref{eq:hqkd-list} is already the KL divergence
$\frac{\alpha}{\alpha-1}\sfD(\mathbf q\|\Phi)$ of the import, so no
divergence relaxation occurs anywhere in this chain.

\emph{Step 1 (joint sector dephasing; exact).} Let
$\mathcal{P}_{\mathbb{T}}(X)\coloneqq\sum_{t\in\mathbb{T}}\Pi^{A}_tX\Pi^{A}_t$ be
the pinching along Alice's declared sectors
(Definition~\ref{def:declaration}\,(DX1); the identity map when
$|\mathbb{T}|=1$), $\PinS$ the pinching along $\SecSet$, for a product
\exttermref{../foundations/contract.pdf}{term-sector-map}{sector map} the tensor product of the
\exttermref{../protocol-compiler/contract.pdf}{term-measured-factor}{measured factors}' pinchings, and
$\omega'\coloneqq(\mathcal{P}_{\mathbb{T}}\otimes\PinS\otimes\idmap_{\Etil})(\omega)$.
Then $\omega'$ is feasible: both pinchings are CPTP, so $\omega'$ is a
state, and in trusted-source mode
$\omega'_A=\mathcal{P}_{\mathbb{T}}(\omega_A)
=\mathcal{P}_{\mathbb{T}}(\tau_A)=\tau_A$ because $\tau_A$ is block diagonal
along $\mathbb{T}$ (Definition~\ref{def:admissible-declaration}\,(iii)). Its
$A\Qsys$-marginal is block diagonal along the joint sectors, and
$\nu^{\omega'}=\nu^{\omega}$: by \eqref{eq:povm-form} the output depends on
$\omega$ only through the operators
$\tr_{A\Qsys}[(M_x\otimes F_v\otimes\idon{\Etil})\,\omega]$, and
$\tr_{A\Qsys}[(M_x\otimes F_v\otimes\idon{\Etil})\,\omega']
=\tr_{A\Qsys}[(\mathcal{P}_{\mathbb{T}}(M_x)\otimes\PinS(F_v)\otimes
\idon{\Etil})\,\omega]
=\tr_{A\Qsys}[(M_x\otimes F_v\otimes\idon{\Etil})\,\omega]$, because both
pinchings are self-adjoint for the trace pairing and act on the traced
systems, every $F_v$ is block diagonal along $\SecSet$ and every $M_x$ along
$\mathbb{T}$ (Definition~\ref{def:admissible-declaration}\,(iii)). \emph{The step is an
equality}: $\omega'$ is another feasible point with literally the same
output state, Eve's register included; no dilation of a pinching is handed
to her here, and the purification is Step~2. Hence every entropy term and $\Phi$ are unchanged.
(If $\CList\neq\varnothing$, every entry still holds at $\omega'$; this is
Corollary~\ref{cor:reduction-constrained}.)

\emph{Step 2 (purification; conservative).} Let
$\ket{\hat\omega}\in\Hs{A\Qsys\Etil\Etil_2}$ purify $\omega'$, with
$\Etil_2$ of dimension $\dim\Hs{A\Qsys\Etil}$. Enlarging the purifying
register is within the scope of the import
(\cite[footnote~17]{arqand2025meat}).
For each $\hat c\in\Chkey$ with $q(\hat c)>0$: if $P_{\hat c}(\omega')=0$
then $\mathbf q\not\ll\Phi$ and the divergence term of the bracket is
already $+\infty$, so the claimed lower bound holds trivially; assume
$P_{\hat c}>0$. Since
$\tr_{\Etil_2}$ maps $\nu^{\hat\omega}_{|\hat c}$ to
$\nu^{\omega'}_{|\hat c}$ (conditioning on the classical $\Chat$ commutes
with the partial trace), Lemma~\ref{lem:cond-more} gives
\begin{equation*}
  \Hsandup(Z|I\Etil)_{\nu^{\omega'}_{|\hat c}}
  \;\ge\;\Hsandup(Z|I\Etil\Etil_2)_{\nu^{\hat\omega}_{|\hat c}}.
\end{equation*}

\emph{Step 3 (arrow replacement; conservative).} By the arrow order of
\xref{foundations/def:renyi},
$\Hsandup(Z|I\Etil\Etil_2)_{\nu^{\hat\omega}_{|\hat c}}
\ge\Hsanddown(Z|I\Etil\Etil_2)_{\nu^{\hat\omega}_{|\hat c}}$.

\emph{Step 4 (branch formula; exact).} $\hat\omega$ is pure with
$\hat\omega_{A\Qsys}=\omega'_{A\Qsys}\eqqcolon\tilde\omega$
block diagonal along the joint sectors, and $\tilde\omega\in\Feas$: normalization is
preserved by every step ($\tr\tilde\omega=\tr\omega=1$), and in
trusted-source mode additionally
$\tilde\omega_A=\omega_A=\tau_A$, none of Steps 1--3 having touched the
$A\Qsys$-marginal except through the joint pinching of Step~1, which
fixes $\tau_A$. Theorem~\ref{thm:branch-formula}, applied with
purifying register $\Etil\Etil_2$, gives
\begin{equation*}
  \Hsanddown(Z|I\Etil\Etil_2)_{\nu^{\hat\omega}_{|\hat c}}
  =\frac{1}{1-\alpha}\log_2
  \frac{K_{\hat c}(\tilde\omega)}{P_{\hat c}(\tilde\omega)}
  =\frac{1}{1-\alpha}\log_2\frac{K_{\hat c}(\tilde\omega)}{P_{\hat c}},
\end{equation*}
using $P_{\hat c}(\tilde\omega)=P_{\hat c}(\omega')=P_{\hat c}$.

\emph{Step 5 (rearrangement; exact).} Collecting Steps 1--4, the bracket of
\eqref{eq:hqkd-list} is bounded from below by
\begin{equation*}
  \frac{\alpha}{\alpha-1}\sfD(\mathbf q\|\Phi)
  +\sum_{\hat c\in\Chkey}
  \frac{q(\hat c)}{1-\alpha}\log_2
  \frac{K_{\hat c}(\tilde\omega)}{P_{\hat c}} .
\end{equation*}
For each $\hat c\in\Chkey$ with $q(\hat c)>0$ and
$P_{\hat c},K_{\hat c}>0$,
\begin{equation}\label{eq:rearrange}
\begin{split}
  \frac{\alpha}{\alpha-1}\,q(\hat c)\log_2\frac{q(\hat c)}{P_{\hat c}}
  &+\frac{q(\hat c)}{1-\alpha}\log_2\frac{K_{\hat c}}{P_{\hat c}}\\
  &=q(\hat c)\log_2\frac{q(\hat c)}{P_{\hat c}}
  +\frac{q(\hat c)}{\alpha-1}
  \left[\log_2\frac{q(\hat c)}{P_{\hat c}}
  -\log_2\frac{K_{\hat c}}{P_{\hat c}}\right]\\
  &=q(\hat c)\log_2\frac{q(\hat c)}{P_{\hat c}}
  +\frac{q(\hat c)}{\alpha-1}\log_2\frac{q(\hat c)}{K_{\hat c}},
\end{split}
\end{equation}
an identity of real numbers. At the degenerate values the identity persists
in the extended sense: $q(\hat c)=0$ makes all terms $0$; $P_{\hat c}=0$
with $q(\hat c)>0$ makes both sides $+\infty$; $K_{\hat c}=0$ with
$q(\hat c)>0$ makes both sides $+\infty$ (using $K_{\hat c}\le P_{\hat c}$
from Theorem~\ref{thm:branch-formula}'s \emph{cone clause}, which is stated without the
hypothesis $P_{\hat c}>0$ that the degenerate cases here violate, so the
$+\infty$ appears on both sides through
the $K$ term on the right and through the combination on the left, where
$\frac{1}{1-\alpha}\log_2\frac{K}{P}\to+\infty$). Summing
\eqref{eq:rearrange} over $\hat c\in\Chkey$ and adding the
$\hat c\notin\Chkey$ part of $\sfD$ yields exactly
$O_{\alpha}(\mathbf q,\tilde\omega)\ge\hlow$ after taking the infimum over
$\Feas$ and then over $\acSet$.
\end{proof}

% claim meat-core/prop:a-neutral
\claimsummary{meat-core/prop:a-neutral}{Pinching a feasible state along Alice's declared sectors keeps it feasible with the same statistics and an objective no larger. So restricting the variable to states block diagonal there changes no value, and the pinching carries minimisers to minimisers.}
\claimbookmark{meat-core/prop:a-neutral}{Proposition}{prop:a-neutral}{Value neutrality along Alice's sectors}
\claimleadin{%
Step~1 of the proof of Theorem~\ref{thm:reduction} shows that restricting the
variable to states block diagonal along Alice's declared sectors is sound; the
proposition below shows, independently, that the restriction also loses no
value.
}
\begin{proposition}[Value neutrality along Alice's sectors]
\label{prop:a-neutral}
Let $\Decl$ satisfy Definition~\ref{def:admissible-declaration} in the general format of
Definition~\ref{def:declaration}, in either source mode and with an arbitrary
list $\CList$, \exttermref{../protocol-compiler/contract.pdf}{term-general}{general} and \exttermref{../protocol-compiler/contract.pdf}{term-party-local}{party-local} entries included.
Let $\Feas^{\circ}$ be the set defined as $\Feas$ is in
Definition~\ref{def:feasible} but with block diagonality required along
$\SecSet$ only --- the feasible set of the same declaration with the factor
tagged $A$ forgotten --- and let
$\mathcal{P}_{\mathbb{T}}(X)\coloneqq\sum_{t}\Pi^{A}_tX\Pi^{A}_t$. Then for
every $\alpha\in(1,\infty)$, every distribution $\mathbf q$ on $\Chatset$ and
every $\omega\in\Feas^{\circ}$, the operator
$\omega'\coloneqq(\mathcal{P}_{\mathbb{T}}\otimes\idmap_{\Qsys})(\omega)$
satisfies
\begin{equation}\label{eq:a-neutral}
  \omega'\in\Feas,\qquad
  P_{\hat c}(\omega')=P_{\hat c}(\omega),\qquad
  K_{\hat c}(\omega')\ge K_{\hat c}(\omega)\quad\forall\hat c,\qquad
  O_{\alpha}(\mathbf q,\omega')\le O_{\alpha}(\mathbf q,\omega),
\end{equation}
the last in the extended-real sense. Consequently
$\inf_{\acSet\times\Feas}O_{\alpha}=\inf_{\acSet\times\Feas^{\circ}}O_{\alpha}$,
and $\omega\mapsto\omega'$ maps minimisers over $\Feas^{\circ}$ to minimisers
over $\Feas$.
\end{proposition}
\claimmeta{\metalabel{Uses}\mref{Def.}{def:declaration}, \mref{Def.}{def:admissible-declaration}, \mref{Def.}{def:maps}, \mref{Def.}{def:branch}, \mref{Lem.}{lem:branch-maps}, \mref{Def.}{def:homog}, \mref{Thm.}{thm:branch-formula}, \mref{Lem.}{lem:cone-concavity}, \mref{Def.}{def:feasible}, \extref{../foundations/contract.pdf}{foundations/def:pinching}, \extref{../foundations/contract.pdf}{foundations/lem:kernel}
\metasep \metalabel{Used by}\mref{Rem.}{rem:fo-general-format}
\metasep newly derived
\metasep accepted}

\begin{proof}
$\mathcal{P}_{\mathbb{T}}$ is a pinching (\xref{foundations/def:pinching}):
completely positive, trace preserving, unital and self-adjoint for the trace
pairing.

\emph{Membership.} $\omega'\succeq0$ and $\tr\omega'=\tr\omega=1$.
$\mathcal{P}_{\mathbb{T}}\otimes\idmap_{\Qsys}$ commutes with
$X\mapsto(\idon{A}\otimes\Pi_s)X(\idon{A}\otimes\Pi_{s'})$, so $\omega'$ is
block diagonal along $\SecSet$ as $\omega$ is, and along $\mathbb{T}$ by
construction; hence along the joint sectors. In trusted-source mode
$\omega'_A=\mathcal{P}_{\mathbb{T}}(\omega_A)=\mathcal{P}_{\mathbb{T}}(\tau_A)
=\tau_A$, because $\tau_A$ is block diagonal along $\mathbb{T}$
(Definition~\ref{def:admissible-declaration}\,(iii)). For an entry read through its lifted
operator, the partial trace over $\Qsys$ commutes with a map acting on $A$
alone, so
$\tr_{\Qsys}[(\idon{A}\otimes\hat O_k)\omega']
=\mathcal{P}_{\mathbb{T}}\big(\tr_{\Qsys}[(\idon{A}\otimes\hat O_k)\omega]\big)$.
In trusted-source mode $\mathcal{P}_{\mathbb{T}}$ is positive and linear, so
$X\bowtie_k\Theta_k$ implies
$\mathcal{P}_{\mathbb{T}}(X)\bowtie_k\mathcal{P}_{\mathbb{T}}(\Theta_k)
=\Theta_k$, every $\Theta_k$ being block diagonal along $\mathbb{T}$: for a
\exttermref{../protocol-compiler/contract.pdf}{term-homogeneous}{homogeneous} entry $\Theta_k=c_k\tau_A$, for a general one by
its type (Definition~\ref{def:declaration}\,(DX2)). In untrusted-source mode
the entry value is the trace of the same operator, which
$\mathcal{P}_{\mathbb{T}}$ preserves. For a party-local entry
scoped to $f$, with $\hat D\coloneqq\hat O_k-c_k\idon{\Qsys_f}$, the partial
trace over $\Qsys_f$ commutes with a map acting on $A$ alone, so
$\tr_{\Qsys_f}[(\idon{}\otimes\hat D)\omega']
=(\mathcal{P}_{\mathbb{T}}\otimes\idmap)\big(
\tr_{\Qsys_f}[(\idon{}\otimes\hat D)\omega]\big)$ with the identity map on
the other measured factors; a pinching tensored with an identity map is a
pinching, hence positive, and $Y\bowtie_k0$ implies
$(\mathcal{P}_{\mathbb{T}}\otimes\idmap)(Y)\bowtie_k0$. So
$\omega'\in\Feas$.

\emph{Linear quantities.} For all $(x,v)$, by self-adjointness of the
pinching and block diagonality of $M_x$ along $\mathbb{T}$
(Definition~\ref{def:admissible-declaration}\,(iii)),
$\tr[(M_x\otimes F_v)\omega']
=\tr[(\mathcal{P}_{\mathbb{T}}(M_x)\otimes F_v)\omega]
=\tr[(M_x\otimes F_v)\omega]$. Hence $\Phimap[\omega']=\Phimap[\omega]$ by
\eqref{eq:Phimap}, so $P_{\hat c}$ is unchanged, and
$\tr\cG_{\hat c}(\omega')=\tr\cG_{\hat c}(\omega)$, that trace being
$\sum_{\kappa\in\Kann_{\mathrm{pass}}(\hat c)}\sum_{(x,v)\in a^{-1}(\kappa)}
\tr[(M_x\otimes F_v)\,\cdot\,]$ by the computation in the proof of
Lemma~\ref{lem:branch-maps}\,(i).

\emph{The non-linear term.} With $m\coloneqq|\mathbb{T}|$, an enumeration
$t\mapsto j_t\in\{1,\dots,m\}$ and $\zeta\coloneqq e^{2\pi i/m}$, the
operators $U_k\coloneqq\sum_t\zeta^{j_tk}\Pi^{A}_t$, $k=0,\dots,m-1$, are
unitaries on $A$ with
$\mathcal{P}_{\mathbb{T}}(X)=\frac1m\sum_kU_kXU_k^{\dagger}$, by the
unitary-mixture form of a pinching in \xref{foundations/def:pinching}, with
the projectors $\Pi^{A}_t$ enumerated by $j_t$. Each $U_k$ is a function of the
projectors $\Pi^{A}_t$ and therefore commutes with every operator block
diagonal along $\mathbb{T}$, in particular with every $\sqrt{M_x}$. Let
$\widetilde U_k$ be $U_k$ on the $A$ factor of the output space
$\Hs{R\,XY\,C_AC_B\,A\Qsys}$ of \eqref{eq:kraus} and the identity on the other
factors. Every term of $K_{\kappa}$ carries
$\sqrt{M_x}\otimes\sqrt{F_v}$ on $A\Qsys$, so
$K_{\kappa}(U_k\otimes\idon{\Qsys})=\widetilde U_kK_{\kappa}$ and, $\cG_{\hat
c}$ being linear,
\begin{equation*}
  \cG_{\hat c}(\omega')=\frac1m\sum_{k=0}^{m-1}
  \widetilde U_k\,\cG_{\hat c}(\omega)\,\widetilde U_k^{\dagger}.
\end{equation*}
For a unitary $U$ that is the identity on $R$ and any $X\succeq0$:
$(UXU^{\dagger})^{1/\alpha}=UX^{1/\alpha}U^{\dagger}$ by functional calculus;
$\cZ(UYU^{\dagger})=U\cZ(Y)U^{\dagger}$ because $U$ commutes with every
$Z_j=\dyad{j}_R$; and
$\tr[(U\cZ(Y)U^{\dagger})^{\alpha}]=\tr[\cZ(Y)^{\alpha}]$. So the functional
$T$ of \eqref{eq:Tcone} satisfies $T(UXU^{\dagger})=T(X)$, and by its
concavity on the whole cone (Lemma~\ref{lem:cone-concavity}\,(b))
\begin{equation*}
  T_{\hat c}(\omega')=T\big(\cG_{\hat c}(\omega')\big)
  \ge\frac1m\sum_kT\big(\widetilde U_k\cG_{\hat c}(\omega)
  \widetilde U_k^{\dagger}\big)=T_{\hat c}(\omega).
\end{equation*}
With \eqref{eq:fK},
$K_{\hat c}=T_{\hat c}+P_{\hat c}-\tr\cG_{\hat c}$, and the two linear terms
unchanged, $K_{\hat c}(\omega')\ge K_{\hat c}(\omega)$.

\emph{The objective and the infima.} In \eqref{eq:objective} the state
enters through $P_{\hat c}$, unchanged, and through
$\frac{1}{\alpha-1}d(q(\hat c),K_{\hat c})$ with $d$ the kernel of
\xref{foundations/lem:kernel}, non-increasing in its second argument; the
conventions at vanishing arguments are those of that kernel, so the
inequality holds in the extended-real sense. Finally
$\Feas\subseteq\Feas^{\circ}$, so the infimum over $\Feas$ is no smaller,
and by \eqref{eq:a-neutral} it is no larger; if $(\mathbf q,\omega)$ attains
the infimum over $\Feas^{\circ}$ then $(\mathbf q,\omega')$ attains the same
value over $\Feas$.
\end{proof}

% claim meat-core/cor:reduction-constrained
\claimsummary{meat-core/cor:reduction-constrained}{For any declared constraint list, licensed or not, the constrained value bounds the constrained single-round rate from below, with the same list on both sides. A key length from the constrained value still needs the attack-set licence of \hyperlink{claim-meat-core/thm:constrained-accumulation}{Theorem~\ref*{thm:constrained-accumulation}}.}
\claimbookmark{meat-core/cor:reduction-constrained}{Corollary}{cor:reduction-constrained}{The reduction for arbitrary constraint lists}
\claimleadin{%
No step of the reduction reads the constraint list beyond checking that the
joint dephasing keeps every entry, so the comparison holds for any list.
Whether the constrained value may enter a key length is a separate question,
answered by the attack-set licence (Theorem~\ref{thm:constrained-accumulation}).
}
\begin{corollary}[The reduction for arbitrary constraint lists]\label{cor:reduction-constrained}
Let $\CList$ be arbitrary: $\Theta_k$ an arbitrary Hermitian operator on $A$
in \exttermref{../transmitter/contract.pdf}{term-trusted-source-mode}{trusted-source mode}, $\Theta_k=c_k$ an arbitrary real number in
\exttermref{../transmitter/contract.pdf}{term-untrusted-source-mode}{untrusted-source mode}, and entries declared \exttermref{../protocol-compiler/contract.pdf}{term-party-local}{party-local} in
either mode (Definition~\ref{def:admissible-declaration}\,(vi)); no licensing restriction
(Definition~\ref{def:homog}) is needed in either mode. Then
\begin{equation}\label{eq:c8-general}
  \hqkd(\CList)\;\ge\;\hlow(\CList),
\end{equation}
where $\hqkd(\CList)$ is the \termref{term-constrained-single-round-rate}{constrained single-round rate} of
Definition~\ref{def:hqkd-list} and $\hlow(\CList)$ denotes
\eqref{eq:hlow-def} over the feasible set $\Feas$ of \eqref{eq:feas},
which carries the same list ($\hlow(\CList)$ and the unadorned
$\hlow$ are the same quantity; the argument is written when the
dependence on the declared list matters, as in
$\hlow(\CList)\ge\hlow(\varnothing)$).
\end{corollary}
\claimmeta{\metalabel{Uses}\mref{Def.}{def:admissible-declaration}, \mref{Def.}{def:homog}, \mref{Def.}{def:hqkd-list}, \mref{Def.}{def:feasible}, \mref{Thm.}{thm:reduction}
\metasep \metalabel{Used by}\mref{Rem.}{rem:hlow-range}, \mref{Prop.}{prop:drop-bound}, \mref{Thm.}{thm:keylength}, \mref{Cor.}{cor:constrained-key-length}, \mref{Prop.}{prop:relaxed-set-properties}
\metasep newly derived
\metasep implemented (effective: accepted)
\metasep \metalabel{Code}\coderef{https://github.com/fraunhofersg/certiqs-key/blob/v0.1.0/src/certiqs_key/meat_core/objective/objective.py}{certiqs_key.meat_core.objective.objective:constraint_value}
\metasep \metalabel{Oracles}\coderef{https://github.com/fraunhofersg/certiqs-key/blob/v0.1.0/tests/oracles/registry.py}{registry.py}}

\begin{proof}
Steps 1--5 of the proof of Theorem~\ref{thm:reduction} apply verbatim
to a feasible point of \eqref{eq:hqkd-list}: Step 1 preserves every
entry. For an entry read through its lifted operator, block diagonality
of $\hat O_k$ along $\SecSet$ and self-adjointness of $\PinS$ give
$\tr_{\Qsys}[(\idon{A}\otimes\hat O_k)\,\omega'_{A\Qsys}]
=\mathcal{P}_{\mathbb{T}}\big(\tr_{\Qsys}[(\idon{A}\otimes\hat O_k)\,
\omega_{A\Qsys}]\big)$; in trusted-source mode $\mathcal{P}_{\mathbb{T}}$ is
positive and linear and fixes $\Theta_k$, which is block diagonal along
$\mathbb{T}$ ($c_k\tau_A$ for a \exttermref{../protocol-compiler/contract.pdf}{term-homogeneous}{homogeneous} entry, by type for
a \exttermref{../protocol-compiler/contract.pdf}{term-general}{general} one, Definition~\ref{def:admissible-declaration}\,(iii)), so
$X\bowtie_k\Theta_k$ implies
$\mathcal{P}_{\mathbb{T}}(X)\bowtie_k\Theta_k$; in untrusted-source mode the
entry value is the trace of the same operator, which the pinching
preserves. For a party-local entry scoped to $f$, with
$\hat D=\hat O_k-c_k\idon{\Qsys_f}$ block diagonal for the pinching
$\mathcal{P}_f$ of that factor and $\mathcal{P}_{f^c}$ the product of the
other factors' pinchings,
$\tr_{\Qsys_f}[(\idon{}\otimes\hat D)\,\omega'_{A\Qsys}]
=(\mathcal{P}_{\mathbb{T}}\otimes\mathcal{P}_{f^c})\big(
\tr_{\Qsys_f}[(\idon{}\otimes\hat D)\,\omega_{A\Qsys}]\big)$, and a positive
map preserves $\bowtie_k0$; Steps 2--3 do not change
the $A\Qsys$-marginal, and the state
$\tilde\omega=\omega'_{A\Qsys}$ of Step 4 therefore satisfies every
entry of $\CList$ in addition to $\tr\tilde\omega=1$, block diagonality along the joint
sectors
and (in trusted-source mode) $\tilde\omega_A=\tau_A$, that is,
$\tilde\omega\in\Feas$. Taking infima over
the constrained sets on both sides yields \eqref{eq:c8-general}.
\end{proof}

% claim meat-core/rem:losses
\claimsummary{meat-core/rem:losses}{The divergence term passes through the reduction as an identity, purification loses nothing when the optimum is attained on pure states, and the arrow replacement's gap closes as \ensuremath{\alpha} \ensuremath{\to} 1. No step reverses direction.}
\claimbookmark{meat-core/rem:losses}{Remark}{rem:losses}{Where each step of the reduction loses}
\begin{remark}[Where each step of the reduction loses]\label{rem:losses}
The divergence term is passed through as an identity: the imported rate
\eqref{eq:hqkd-def} already carries the KL divergence $\sfD$, so nothing
is lost there. Step 2 of the proof of Theorem~\ref{thm:reduction} loses
nothing when the feasible optimum is attained on pure states; Step 3 is the
standard down-arrow relaxation (\xref{foundations/def:renyi}), whose gap
closes at $\alpha\to1^+$ (both variants converge to the von~Neumann
conditional entropy). No step reverses direction. The bounds
$0\le K_{\hat c}\le P_{\hat c}$ that Step~4 also supplies are read at points
that are not pure, which is why Theorem~\ref{thm:branch-formula} states them
as its \emph{cone clause}.
\end{remark}
\claimmeta{\metalabel{Uses}\mref{Thm.}{thm:accumulated-entropy}, \mref{Thm.}{thm:branch-formula}, \mref{Thm.}{thm:reduction}, \extref{../foundations/contract.pdf}{foundations/def:renyi}
\metasep newly derived
\metasep accepted}

% claim meat-core/prop:value-range
\claimsummary{meat-core/prop:value-range}{The objective is non-negative on the feasible set, so h\_low lies in [0, +\ensuremath{\infty}]; each entropy term is at most log\ensuremath{_{2}}(d\_Z + 1), while the divergence penalty has no finite bound. If h\_low = +\ensuremath{\infty}, no accepted distribution is absolutely continuous with respect to the statistics of a feasible state.}
\claimbookmark{meat-core/prop:value-range}{Proposition}{prop:value-range}{Range of the value and its finite entropy part}
\claimleadin{%
The \termref{term-objective}{objective} is a divergence penalty plus one entropy term per
\termref{term-key-live-acceptance-value}{key-live acceptance value} (Definition~\ref{def:feasible}). The two parts
behave differently: every entropy term is at most $\log_2(d_Z+1)$, while the
penalty is not bounded at all.
}
\begin{proposition}[Range of the value and its finite entropy part]\label{prop:value-range}
Let $\alpha\in(1,\infty)$.
\begin{enumerate}[label=(\roman*),leftmargin=2.2em]
\item \emph{(lower bound)} $O_{\alpha}(\mathbf q,\omega)\ge0$ for every
distribution $\mathbf q$ on $\Chatset$ and every $\omega\in\Feas$; hence
$\hlow\in[0,+\infty]$.
\item \emph{(entropy-part ceiling)} For every $\omega\in\Feas$ and every
$\hat c\in\Chkey$ with $P_{\hat c}(\omega)>0$, the entropy term
$e_{\hat c}(\omega)\coloneqq\frac{1}{1-\alpha}
\log_2\frac{K_{\hat c}(\omega)}{P_{\hat c}(\omega)}$ satisfies
$e_{\hat c}(\omega)\le\log_2|\Zset|=\log_2(d_Z+1)$.
\item \emph{(no overall ceiling)} No finite a-priori upper bound on $\hlow$
exists: the divergence penalty is unbounded already at fixed $\alpha$.
\item \emph{(infinite value)} $\hlow=+\infty$ whenever no
$(\mathbf q,\omega)$ with $\mathbf q\ll\Phimap[\omega]$ is feasible. If
$\hlow=+\infty$, then no $\mathbf q\in\acSet$ is absolutely continuous with
respect to $\Phimap[\omega]$ for any $\omega\in\Feas$.
\end{enumerate}
\end{proposition}
\claimmeta{\metalabel{Uses}\mref{Lem.}{lem:cond-more}, \mref{Thm.}{thm:branch-formula}, \mref{Def.}{def:feasible}, \mref{Thm.}{thm:reduction}, \extref{../foundations/contract.pdf}{foundations/def:kl}, \extref{../foundations/contract.pdf}{foundations/def:renyi}
\metasep \metalabel{Used by}\mref{Rem.}{rem:hlow-range}, \mref{Rem.}{rem:classes}, \mref{Prop.}{prop:minimal-divergence}, \mref{Cor.}{cor:keylength-bound}, \mref{Prop.}{prop:optol-nobound}, \mref{Rem.}{rem:representability-ceiling}, \mref{Rem.}{rem:optol-direction}, \mref{Rem.}{rem:declared-measurement}, \mref{Req.}{req:infinite-value}, \mref{Req.}{req:unreachable-acceptance-set}, \mref{Prop.}{prop:relaxed-set-properties}, \mref{Lem.}{lem:fo-khyp}, \mref{Thm.}{thm:firstorder}, \mref{Rem.}{rem:fo-oracle-repair}, \mref{Rem.}{rem:fo-negative-values}, \mref{Req.}{req:negative-estimate-refused}
\metasep newly derived
\metasep accepted}

\begin{proof}
\emph{(i)} By \xref{foundations/def:kl} and
\eqref{eq:objective}--\eqref{eq:hlow-def},
$O_{\alpha}(\mathbf q,\omega)
=\frac{\alpha}{\alpha-1}\sfD(\mathbf q\|\Phimap[\omega])
+\sum_{\hat c\in\Chkey}q(\hat c)\,e_{\hat c}(\omega)$ with
$e_{\hat c}\coloneqq\frac{1}{1-\alpha}\log_2\frac{K_{\hat c}}{P_{\hat c}}$;
the KL term is non-negative by Gibbs' inequality of
\xref{foundations/def:kl}, $\mathbf q$ and $\Phimap[\omega]$ being
probability vectors, and $e_{\hat c}\ge0$ since
$K_{\hat c}\le P_{\hat c}$ (Theorem~\ref{thm:branch-formula}, \emph{cone
clause}: a point of $\Feas$ is not pure, so the identity \eqref{eq:c5} is not
what is read here) and $\frac{1}{1-\alpha}<0$. Both parts are non-negative at
every $(\mathbf q,\omega)$, so $O_{\alpha}(\mathbf q,\omega)\ge0$ there and
$\hlow\ge0$.

\emph{(ii)} (If $P_{\hat c}=0$ and $q(\hat c)>0$, the KL term is already
$+\infty$.) Let $P_{\hat c}(\omega)>0$. Purifying $\omega$ as in
Theorem~\ref{thm:reduction}, by Theorem~\ref{thm:branch-formula}
$e_{\hat c}$ equals $\Hsanddown(Z|I\Etil\Etil_2)$ of a normalized branch
state, and by the arrow order $\Hsanddown\le\Hsandup$ of
\xref{foundations/def:renyi} and Lemma~\ref{lem:cond-more} (dropping all
conditioning) it is at most the R\'enyi entropy
$H_{\alpha}(Z)=\frac{1}{1-\alpha}\log_2\sum_zp(z)^{\alpha}
\le\log_2|\Zset|=\log_2(d_Z+1)$, using that
$\sum_zp(z)^{\alpha}\ge|S|^{1-\alpha}$ for $\alpha>1$ and
$S=\operatorname{supp}p$ (convexity of $t\mapsto t^{\alpha}$:
$\frac{1}{|S|}\sum_{z\in S}p(z)^{\alpha}
\ge\big(\frac{1}{|S|}\big)^{\alpha}$), so that
$\frac{1}{1-\alpha}\log_2\sum_zp(z)^{\alpha}\le\log_2|S|$.

\emph{(iii)} Take $d_A=d_Q=1$ (single Alice outcome, trivial
announcement), $\Vset=\{0,1\}$ with $F_0=\epsilon$, $F_1=1-\epsilon$ for
$\epsilon\in(0,\tfrac12)$ and $a_B=\mathrm{id}$, $\phi(\kappa)=c_B$,
all-fail sifting ($\varsigma\equiv\mathrm{fail}$, $r\equiv\perp$, so
$\Chkey=\varnothing$), and a singleton acceptance set
$\acSet=\{(q_0,1-q_0)\}$ with $q_0\in(0,1]$ fixed. Then
$\Phimap\equiv(\epsilon,1-\epsilon)$ and the objective is the constant
$\frac{\alpha}{\alpha-1}\sfD\big((q_0,1-q_0)\|(\epsilon,1-\epsilon)\big)
\ge\frac{\alpha}{\alpha-1}\,\big(q_0\log_2\frac{q_0}{\epsilon}
+(1-q_0)\log_2(1-q_0)\big)\to+\infty$ as $\epsilon\to0^{+}$, at every
fixed $\alpha\in(1,\infty)$.

\emph{(iv)} If no $(\mathbf q,\omega)\in\acSet\times\Feas$ has
$\mathbf q\ll\Phimap[\omega]$, the KL term is $+\infty$ at every such pair
(the convention $u\log_2\frac u0=+\infty$ for $u>0$ of
\xref{foundations/def:kl}) and the entropy terms are non-negative by (i), so
$\hlow=+\infty$. Conversely, if some $(\mathbf q,\omega)\in\acSet\times\Feas$
has $\mathbf q\ll\Phimap[\omega]$, its KL term is finite, and every
$\hat c\in\Chkey$ with $q(\hat c)>0$ has $P_{\hat c}(\omega)>0$ and an
entropy term in $[0,\log_2|\Zset|]$ by (i) and (ii), so
$O_{\alpha}(\mathbf q,\omega)<+\infty$ and $\hlow<+\infty$. Hence
$\hlow=+\infty$ means that no $\mathbf q\in\acSet$ is absolutely continuous
with respect to $\Phimap[\omega]$ for any $\omega\in\Feas$.
\end{proof}

% claim meat-core/rem:hlow-range
\claimsummary{meat-core/rem:hlow-range}{An infinite value of the program is the infeasibility of the declared acceptance set against the declaration: no attack the declared list admits is then accepted and no key is produced. The entropy part of the objective is capped, the divergence penalty is not.}
\claimbookmark{meat-core/rem:hlow-range}{Remark}{rem:hlow-range}{Reading an infinite value}
\begin{remark}[Reading an infinite value]\label{rem:hlow-range}
An infinite value is not a large rate. By
Proposition~\ref{prop:value-range}\,(iv), $\hlow=+\infty$ is the
infeasibility of the declared acceptance set against the declaration: no
$\mathbf q\in\acSet$ is absolutely continuous with respect to
$\Phimap[\omega]$ for any $\omega\in\Feas$. For the attacks the declared
list admits this means $\Pr[\OmAT]=0$: for every \exttermref{../foundations/contract.pdf}{term-attack-sequence}{attack sequence} of
Definition~\ref{def:roundmodel} when $\CList=\varnothing$, and for every
$\CList$-admissible one when $\CList$ is
\exttermref{../protocol-compiler/contract.pdf}{term-licensed-for-the-declared-source-mode}{licensed for the declared source mode}, the attacks
Theorem~\ref{thm:constrained-accumulation} ranges over, the initial state
being in either case that of Theorem~\ref{thm:accumulated-entropy}.
Indeed, let $\sigma_j\in\Seq{A\Qsys}$ be the unconditioned state the
protocol channel acts on in round $j$, the $A\Qsys$-marginal of the output
of the round-$j$ attack; in trusted-source mode the attack acts on the
$A_jE_{j-1}$-marginal of the state before round $j$, whose $A_j$-marginal
is $\tau_A$, the earlier rounds acting on other copies of $A$. By
Lemma~\ref{lem:range}\,(i), $\sigma_j$ satisfies the declaration's marginal
condition and every entry of $\CList$, and by
Theorem~\ref{thm:protocol-channel}\,(iii) the acceptance value of round $j$
is distributed as $\Phimap[\sigma_j]$. The average
$\bar\sigma\coloneqq\frac1n\sum_j\sigma_j$ satisfies the same conditions,
each of them being convex, and
$\Phimap[\bar\sigma]=\frac1n\sum_j\Phimap[\sigma_j]$. Pinching
$\bar\sigma$ along the joint sectors, as in Step~1 of the proof of
Theorem~\ref{thm:reduction}, keeps its statistics and its marginal and, by
the proof of Corollary~\ref{cor:reduction-constrained}, every entry of
$\CList$, and gives a point $\bar\omega\in\Feas$ with
$\Phimap[\bar\omega]=\Phimap[\bar\sigma]$. If $\Pr[\OmAT]>0$, some sequence
$\hat c_1^n$ of positive probability has its frequency vector $\mathbf f$
in $\acSet$; every $\hat c$ with $f(\hat c)>0$ is the value of some round
$j$ of it, so $\Phimap[\sigma_j](\hat c)>0$ and
$\Phimap[\bar\omega](\hat c)\ge\frac1n\Phimap[\sigma_j](\hat c)>0$. Then
$\mathbf f\ll\Phimap[\bar\omega]$ with $\mathbf f\in\acSet$ and
$\bar\omega\in\Feas$, which Proposition~\ref{prop:value-range}\,(iv) rules
out. So the hypothesis $\Pr[\OmAT]>0$ of
Theorem~\ref{thm:accumulated-entropy}, respectively
Theorem~\ref{thm:constrained-accumulation}, fails, and the accepted state
is null: there is no run on which a key would be produced.

For a nonempty list the conclusion reaches no further than the
$\CList$-admissible attacks. Take $d_A=1$, $\Qsys=\mathbb{C}^2$ with one
sector, $F_v=\dyad{v}$ for $v\in\{0,1\}$, the acceptance value the outcome
$v$, $\acSet=\{(\frac12,\frac12)\}$ and for $\CList$ the one homogeneous
entry $\tr[(\idon{A}\otimes\dyad{1})\,\omega]=0$. Then
$\Feas=\{\dyad{0}\}$ and $\hlow=+\infty$, while the attack that prepares
$\ket{+}$ in every round, which is not $\CList$-admissible, is accepted with
positive probability at every even $n$.

The two parts of the objective differ here: the entropy part is capped, each
term at $\log_2(d_Z+1)$ (Proposition~\ref{prop:value-range}\,(ii)), while the
divergence penalty is not (Proposition~\ref{prop:value-range}\,(iii)). What
the engine does with an infinite value is
Requirement~\ref{req:infinite-value}.
\end{remark}
\claimmeta{\metalabel{Uses}\mref{Thm.}{thm:protocol-channel}, \mref{Def.}{def:roundmodel}, \mref{Lem.}{lem:range}, \mref{Thm.}{thm:accumulated-entropy}, \mref{Thm.}{thm:constrained-accumulation}, \mref{Thm.}{thm:reduction}, \mref{Cor.}{cor:reduction-constrained}, \mref{Prop.}{prop:value-range}
\metasep newly derived
\metasep accepted}

% claim meat-core/thm:convexity
\claimsummary{meat-core/thm:convexity}{The objective is jointly convex and lower semicontinuous in the acceptance distribution and the feasible state, so h\_low is a convex minimisation over a compact convex set. A finite infimum is attained.}
\claimbookmark{meat-core/thm:convexity}{Theorem}{thm:convexity}{Joint convexity of the program}
\claimleadin{%
The objective is a positive combination of the divergence kernel of
\xref{foundations/lem:kernel}, evaluated at the affine probabilities $P_{\hat c}$
and at the concave quantities $K_{\hat c}$ (Corollary~\ref{cor:concavity});
that is what makes the program convex.
}
\begin{theorem}[Joint convexity of the program]\label{thm:convexity}
For every $\alpha\in(1,\infty)$ and every declaration satisfying
Definition~\ref{def:admissible-declaration} (any $\CList$), the extended-real objective
$O_{\alpha}$ of \eqref{eq:objective} is jointly convex and lower
semicontinuous on $\Delta(\Chatset)\times\Feas$, where $\Delta(\Chatset)$
is the probability simplex. Consequently \eqref{eq:hlow-def} is a convex
minimization over the compact convex set $\acSet\times\Feas$, and the
infimum is attained whenever it is finite.
\end{theorem}
\claimmeta{\metalabel{Uses}\mref{Def.}{def:admissible-declaration}, \mref{Thm.}{thm:protocol-channel}, \mref{Def.}{def:admissible-acceptance-set}, \mref{Thm.}{thm:branch-formula}, \mref{Cor.}{cor:concavity}, \mref{Def.}{def:feasible}, \extref{../foundations/contract.pdf}{foundations/lem:kernel}
\metasep \metalabel{Used by}\mref{Thm.}{thm:choi}, \mref{Def.}{def:untrusted-choi}, \mref{Rem.}{rem:cgate}, \mref{Prop.}{prop:relaxed-set-properties}, \mref{Thm.}{thm:firstorder}
\metasep newly derived
\metasep implemented (effective: accepted)
\metasep \metalabel{Code}\coderef{https://github.com/fraunhofersg/certiqs-key/blob/v0.1.0/src/certiqs_key/meat_core/objective/objective.py}{certiqs_key.meat_core.objective.objective:objective_from_branches}
\metasep \metalabel{Oracles}\coderef{https://github.com/fraunhofersg/certiqs-key/blob/v0.1.0/tests/oracles/registry.py}{registry.py}}

\begin{proof}
Write $O_{\alpha}$ through the kernel $d$ of \xref{foundations/lem:kernel}:
\begin{equation*}
  O_{\alpha}(\mathbf q,\omega)
  =\frac{\alpha}{\alpha-1}\sum_{\hat c\notin\Chkey}
  d\big(q(\hat c),P_{\hat c}(\omega)\big)
  +\sum_{\hat c\in\Chkey}\Big[
  d\big(q(\hat c),P_{\hat c}(\omega)\big)
  +\frac{1}{\alpha-1}\,d\big(q(\hat c),K_{\hat c}(\omega)\big)\Big].
\end{equation*}
All prefactors are positive for $\alpha>1$. Each
$P_{\hat c}$ is affine in $\omega$ (Theorem~\ref{thm:protocol-channel}\,(iii)), so
$(\mathbf q,\omega)\mapsto d(q(\hat c),P_{\hat c}(\omega))$ is jointly
convex and lsc as the composition of the jointly convex lsc $d$ with an
affine map. For the $K$-terms: $K_{\hat c}$ is concave, non-negative
(Theorem~\ref{thm:branch-formula}, \emph{cone clause} --- a point of $\Feas$ is not
pure, and $\Feas\subseteq\Pos{A\Qsys}$ is all that clause asks), and
continuous on $\Feas$
(Corollary~\ref{cor:concavity}; continuity holds because $K_{\hat c}$ is a
composition of continuous maps on a finite-dimensional compact domain).
For feasible pairs $(\mathbf q_i,\omega_i)$ and $\lambda\in(0,1)$,
concavity of $K_{\hat c}$ and monotonicity of $d$ in its second argument
give
\begin{equation*}
  d\big(q_{\lambda}(\hat c),K_{\hat c}(\omega_{\lambda})\big)
  \le d\big(q_{\lambda}(\hat c),
  \lambda K_{\hat c}(\omega_1)+(1-\lambda)K_{\hat c}(\omega_2)\big)
  \le\lambda\,d\big(q_1(\hat c),K_{\hat c}(\omega_1)\big)
  +(1-\lambda)\,d\big(q_2(\hat c),K_{\hat c}(\omega_2)\big),
\end{equation*}
the second step by joint convexity of $d$. Lower semicontinuity of the
$K$-terms follows from lsc of $d$, continuity of $K_{\hat c}$, and
monotonicity of $d$ in $v$. A finite non-negative combination of jointly
convex lsc functions is jointly convex and lsc. $\acSet\times\Feas$ is
compact ($\acSet$ closed in the simplex by Definition~\ref{def:admissible-acceptance-set};
$\Feas$ closed and bounded); an lsc function on a compact set attains its
infimum when finite.
\end{proof}

\section{Sector classes and sector removal}\label{group-meat-core/def:classes}
\sectionintro{Sectors are classified by exact predicates on the declared data as key-active, constraint-active, state-blind and removable, removability being certified, for a list without party-local entries, by an exhibited witness that meets a six-condition checklist (\hyperlink{claim-meat-core/def:classes}{Definition~\ref*{def:classes}}). A state-blind sector may be deleted when a mixture of states on distinct other key-inert sectors reproduces its responses and constraint values (\hyperlink{claim-meat-core/thm:removal}{Theorem~\ref*{thm:removal}}), and a set of such sectors when every witness targets only sectors that stay (\hyperlink{claim-meat-core/prop:removal-compose}{Proposition~\ref*{prop:removal-compose}}). Dropping a constraint-active sector (\hyperlink{claim-meat-core/prop:retention}{Proposition~\ref*{prop:retention}}) or averaging sectors (\hyperlink{claim-meat-core/prop:no-averaging}{Proposition~\ref*{prop:no-averaging}}) can overstate the key. The remarks show why the index conditions of the checklist are needed, why the witness must be exact, and that a tolerance may only refuse a removal.}

% claim meat-core/def:classes
\claimsummary{meat-core/def:classes}{Exact predicates on the declared data make a sector key-active, constraint-active, state-blind or removable; for a list without party-local entries, removability is certified by a witness meeting a six-condition checklist. Deleting several removable sectors needs a further condition on the set.}
\claimbookmark{meat-core/def:classes}{Definition}{def:classes}{Sector classes and the removal checklist}
\claimleadin{%
Two operations on the declared sectors cost nothing: evaluating the
objective over the sectors that carry a deficit, and deleting a sector whose
weight other key-inert sectors can take over. The classes below name the
sectors each operation applies to, and the checklist says when a sector may
be deleted. Neither half of (W2) is redundant
(Remark~\ref{rem:target-conditions}); why (W3) and (W4) must hold exactly,
and why (W4) names one form, is Remark~\ref{rem:witness-exactness}; and
Remark~\ref{rem:removal-sets} exhibits two removable sectors whose joint
deletion overstates the value.
}
\begin{definition}[Sector classes and the removal checklist]\label{def:classes}
Let $\Decl$ satisfy Definition~\ref{def:admissible-declaration}. For $s\in\SecSet$:
\begin{enumerate}[label=(\alph*),leftmargin=2.2em]
\item $s$ is \termdef{term-key-active}{key-active} if $(\hat c,s)$ fails to be key-inert
(Definition~\ref{def:keyinert}) for at least one $\hat c\in\Chkey$; these
sectors carry the deficits $f_{\hat c,s}$ in the objective.
\item $s$ is \termdef{term-constraint-active}{constraint-active} if it is key-inert for every
$\hat c\in\Chkey$; by Proposition~\ref{prop:exact-cancellation} it contributes nothing to
any $K_{\hat c}$ but remains an optimization variable (it enters
$\Phimap$, every $P_{\hat c}$, the marginal constraint, and $\CList$).
\item $s$ is \termdef{term-state-blind}{state-blind} if $F^{(s)}_v=c_{v,s}\,\idon{\Qsys_s}$ for
scalars $c_{v,s}\ge0$, for every $v\in\Vset$ (then
$\sum_vc_{v,s}=1$ by POVM completeness restricted to the sector); if
$\CList\neq\varnothing$ we additionally require each
$\hat O_k^{(s)}=o_{k,s}\,\idon{\Qsys_s}$ for scalars $o_{k,s}$.
\item $s$ is \termdef{term-removable}{removable} if the hypotheses of
Theorem~\ref{thm:removal} hold for it.
\end{enumerate}
All four predicates are decided by finitely many linear-algebra checks on
the declaration. Items (a) and (b) reduce to key-inertness
(Definition~\ref{def:keyinert}), a finite combinatorial predicate on the
declared function tables and on which of the $M_x$, $F^{(s)}_v$ are the zero
operator, decided exactly, at no window. Items (c) and (d) are decided
exactly as well: (c) is a condition on the
\exttermref{../foundations/contract.pdf}{term-declared-value}{declared values} themselves
(\xref{foundations/def:binary64}), so that a block that is only
\emph{nearly} scalar fails it and $s$ is then not state-blind, and (d) is
decided by the checklist below.

\emph{The removal checklist.} A \emph{candidate witness} for a sector
$s^{*}$ is a target list $s_1,\dots,s_L$ as index data, states $\chi_i$ as
binary64 sector-block data, weights $\mu_i$ as binary64 data, and the
evaluation order of (W6). It is a \emph{witness} for $s^{*}$ in the sense
of~(d) when $\CList$ carries no entry declared \exttermref{../protocol-compiler/contract.pdf}{term-party-local}{party-local} and all six
of the following hold:
\begin{enumerate}[label=(W\arabic*),leftmargin=3.0em]
\item \emph{$s^{*}$ is state-blind} in the sense of (c), the
$\hat O^{(s^{*})}_k$ half included whenever $\CList\neq\varnothing$;
\item \emph{the witness does not name $s^{*}$, and names no sector
twice}: $L\ge1$, $s_i\in\SecSet\setminus\{s^{*}\}$ for every $i$, and the
$s_i$ pairwise distinct, all compared as index data at
no window. A candidate naming $s^{*}$ fails (W2), and so does one that
repeats a target;
\item \emph{weights}: $\mu_i\ge0$ for every $i$ --- a comparison of an
exhibited binary64 datum against zero --- and $\sum_i\mu_i=1$, the sum
formed in the evaluation order the witness exhibits and verified to
\exttermref{../foundations/contract.pdf}{term-rounds-nowhere}{round nowhere}, exactly as (W6) requires of its
residuals, so that the computed sum \emph{is} the real sum;
\item \emph{states}: $\chi_i\succeq0$ and $\tr\chi_i=1$ for every $i$,
decided on the \emph{exhibited form} and never on a computed spectrum:
the witness presents each $\chi_i$ diagonal in the declared basis of
$\Qsys_{s_i}$ --- every off-diagonal entry the datum $0$ in the sense of
(c), every diagonal entry a datum $\ge0$, and their sum formed in the
exhibited order, rounding nowhere, and equal to the datum $1$ --- so that
both conditions hold of the declared numbers by construction. Any other
presentation fails (W4), which can only \emph{refuse} a removal;
\item \emph{key-inert targets}, which is (RM2): every pair
$(\hat c,s_i)$ with $\hat c\in\Chkey$ and $i\in\{1,\dots,L\}$ is key-inert
(Definition~\ref{def:keyinert});
\item \emph{both equality families close at residual exactly zero}: (RM1)
for every $v\in\Vset$, and (RM3) for every $k$ (vacuous when
$\CList=\varnothing$).
\end{enumerate}
The condition on $\CList$ and (W1)--(W6) are the hypotheses of
Theorem~\ref{thm:removal} transcribed, all of them and in the theorem's own
words --- the condition on $\CList$ from its first hypothesis, (W1) and
(W2) from its opening sentence, (W3) and (W4) from its quantifiers over
$\mu_i$ and $\chi_i$, (W5) from (RM2), (W6) from (RM1) and (RM3) --- which
is what (d) means by ``the hypotheses of Theorem~\ref{thm:removal}
hold''.
Condition (W6) is read on the residual itself: each residual is formed in
an evaluation order that the witness exhibits alongside its data, every
elementary operation of which rounds nothing, so that the computed residual
\emph{is} the real residual and its comparison against zero decides the
real identity.

\emph{Deleting a set of removable sectors needs one condition more, and it is
not part of this definition.} Removability is a property of one sector, and
Theorem~\ref{thm:removal} licenses the deletion of one sector; that two
sectors are each removable does not license deleting both.
Proposition~\ref{prop:removal-compose} states the further hypothesis on the
set $\SecSet_{\mathrm{del}}$ a run deletes.
Nothing in (W1)--(W6) implies that hypothesis.

\emph{In the general format} of Definition~\ref{def:declaration} the four
classes remain predicates of the sectors $s\in\SecSet$ of $\Qsys$, decided
as above, with $\hat O^{(s)}_k$ the block of the \exttermref{../protocol-compiler/contract.pdf}{term-lifted-operator}{lifted operator} of an
entry. The finer notion along Alice's sectors --- which joint
pairs $(\hat c,(t,s))$ carry a deficit --- is
Definition~\ref{def:joint-sectors}, and a sector that is key-active here may
have key-inert joint pairs. Theorem~\ref{thm:removal} is stated for a list
carrying no entry declared party-local, so for a list carrying one,
clause~(d) is unsatisfiable and no sector carries a witness in its sense.
\end{definition}
\claimmeta{\metalabel{Uses}\mref{Def.}{def:keyinert}, \mref{Prop.}{prop:exact-cancellation}, \mref{Def.}{def:feasible}
\metasep \metalabel{Used by}\mref{Thm.}{thm:removal}, \mref{Prop.}{prop:removal-compose}, \mref{Rem.}{rem:target-conditions}, \mref{Rem.}{rem:witness-exactness}, \mref{Rem.}{rem:removal-general-format}, \mref{Rem.}{rem:removal-sets}, \mref{Rem.}{rem:classes}, \mref{Req.}{req:sector-report}, \mref{Req.}{req:state-blind-test}, \mref{Req.}{req:removal-checklist}, \mref{Req.}{req:single-sector-deletion}
\metasep newly derived
\metasep implemented (effective: accepted)
\metasep \metalabel{Code}\coderef{https://github.com/fraunhofersg/certiqs-key/blob/v0.1.0/src/certiqs_key/meat_core/declaration/sectors.py}{certiqs_key.meat_core.declaration.sectors:classify_sector}
\metasep \metalabel{Oracles}\coderef{https://github.com/fraunhofersg/certiqs-key/blob/v0.1.0/tests/oracles/registry.py}{registry.py}}

% claim meat-core/thm:removal
\claimsummary{meat-core/thm:removal}{A state-blind sector may be deleted when a mixture of states on distinct other sectors, all key-inert, reproduces its responses and its constraint values. Deleting it changes neither the value of the program nor its attainment.}
\claimbookmark{meat-core/thm:removal}{Theorem}{thm:removal}{Sound removal of a state-blind sector}
\claimleadin{%
A state-blind sector answers every state alike, so its weight can be moved
onto other sectors that reproduce its answers. When those sectors carry no
key entropy, the move removes one non-positive deficit and changes nothing
else.
}
\begin{theorem}[Sound removal of a state-blind sector]\label{thm:removal}
Let $\CList$ carry no entry declared \exttermref{../protocol-compiler/contract.pdf}{term-party-local}{party-local}, let
$s^{*}\in\SecSet$ be state-blind, and suppose there exist
\emph{pairwise distinct} sectors
$s_1,\dots,s_L\in\SecSet\setminus\{s^{*}\}$, states
$\chi_i\in\Seq{\Qsys_{s_i}}$, and weights $\mu_i\ge0$ with
$\sum_{i=1}^{L}\mu_i=1$, such that
\begin{enumerate}[label=(RM\arabic*),leftmargin=3.0em]
\item \emph{response match:}\quad
$c_{v,s^{*}}=\sum_{i}\mu_i\tr\big[F^{(s_i)}_v\chi_i\big]$ for every
$v\in\Vset$;
\item \emph{key-inert targets:}\quad every pair $(\hat c,s_i)$ with
$\hat c\in\Chkey$ is key-inert;
\item \emph{constraint match} (vacuous for $\CList=\varnothing$):\quad
$o_{k,s^{*}}=\sum_i\mu_i\tr\big[\hat O^{(s_i)}_k\chi_i\big]$ for every
$k$.
\end{enumerate}
Then restricting the infimum \eqref{eq:hlow-def} to
$\{\omega\in\Feas:\omega_{s^{*}}=0\}$ changes neither its value nor its
attainment: the sector $s^{*}$ may be deleted from the variable list, and
its rows deleted from the generated maps, without loss.
\end{theorem}
\claimmeta{\metalabel{Uses}\mref{Def.}{def:maps}, \mref{Cor.}{cor:f-nonpos}, \mref{Prop.}{prop:exact-cancellation}, \mref{Def.}{def:feasible}, \mref{Def.}{def:classes}
\metasep \metalabel{Used by}\mref{Prop.}{prop:no-averaging}, \mref{Prop.}{prop:removal-compose}, \mref{Rem.}{rem:target-conditions}, \mref{Rem.}{rem:witness-exactness}, \mref{Rem.}{rem:removal-general-format}, \mref{Rem.}{rem:removal-sets}, \mref{Rem.}{rem:classes}, \mref{Rem.}{rem:cgate}, \mref{Req.}{req:removal-checklist}
\metasep newly derived
\metasep implemented (effective: accepted)
\metasep \metalabel{Code}\coderef{https://github.com/fraunhofersg/certiqs-key/blob/v0.1.0/src/certiqs_key/meat_core/declaration/sectors.py}{certiqs_key.meat_core.declaration.sectors:verify_removal_witness}
\metasep \metalabel{Oracles}\coderef{https://github.com/fraunhofersg/certiqs-key/blob/v0.1.0/tests/oracles/registry.py}{registry.py}}

\begin{proof}
Given $\omega\in\Feas$ and $\mathbf q\in\acSet$, let
$m\coloneqq\tr_{\Qsys_{s^{*}}}\omega_{s^{*}}\in\Pos{A}$ and define
$\hat\omega$ by
\begin{equation*}
  \hat\omega_{s^{*}}\coloneqq0,\qquad
  \hat\omega_{s_i}\coloneqq\omega_{s_i}+\mu_i\,m\otimes\chi_i,\qquad
  \hat\omega_{s}\coloneqq\omega_{s}\ \ (s\notin\{s^{*},s_1,\dots,s_L\}).
\end{equation*}
The three assignments are mutually consistent, which the two index
conditions of the hypothesis ensure: no $s_i$ equals $s^{*}$, so
the first two do not collide, and the $s_i$ are pairwise distinct, so the
second assigns one value to each block it touches.
\emph{Feasibility.} Each block is PSD; sector-block-diagonality is
manifest. With Alice sectors, $\omega_s=\bigoplus_t\omega^{(t,s)}$ is block
diagonal along $\mathbb{T}$, hence so is $m=\bigoplus_tm_t$ with
$m_t=\tr_{\Qsys_{s^{*}}}\omega^{(t,s^{*})}$, and so is every
$m\otimes\chi_i$: the modified point is block diagonal along the
\exttermref{../protocol-compiler/contract.pdf}{term-joint-sector}{joint sectors}, and the move takes place inside each
Alice sector separately. \emph{Normalization} --- the condition that holds
in both source modes, and the only one in
\exttermref{../transmitter/contract.pdf}{term-untrusted-source-mode}{untrusted-source mode}:
$\tr\hat\omega=\tr\omega-\tr m+\sum_i\mu_i\,\tr m\,\tr\chi_i
=\tr\omega=1$, using $\tr\chi_i=1$, $\sum_i\mu_i=1$. \emph{Alice's
marginal}, in \exttermref{../transmitter/contract.pdf}{term-trusted-source-mode}{trusted-source mode} only: the same computation at the
level of the partial trace,
$\tr_{\Qsys}\hat\omega
=\tr_{\Qsys}\omega-m+\sum_i\mu_i\,m\,\tr\chi_i
=\tr_{\Qsys}\omega=\tau_A$; the normalization statement is its trace,
so replacing the marginal by normalization is exactly the step that
carries this theorem into untrusted-source mode.
\emph{Constraints.} The $s^{*}$ contribution to
$\tr_{\Qsys}[(\idon{A}\otimes\hat O_k)\omega]$ is
$\tr_{\Qsys_{s^{*}}}[(\idon{A}\otimes o_{k,s^{*}}\idon{})\,\omega_{s^{*}}]
=o_{k,s^{*}}m$, and the added contribution is
$\sum_i\mu_i\tr[\hat O^{(s_i)}_k\chi_i]\,m=o_{k,s^{*}}m$ by (RM3); all
other sectors unchanged. In untrusted-source mode the constraint value
is the trace of that same operator, so the identity gives
$o_{k,s^{*}}\tr m$ on both sides.
\emph{Statistics.} For every $(x,v)$, the $s^{*}$ term of
\eqref{eq:Phimap} is
$\tr[(M_x\otimes c_{v,s^{*}}\idon{})\,\omega_{s^{*}}]
=c_{v,s^{*}}\tr[M_xm]$, and the added terms give
$\sum_i\mu_i\tr[M_xm]\tr[F^{(s_i)}_v\chi_i]=c_{v,s^{*}}\tr[M_xm]$ by
(RM1). Hence $\Phimap[\hat\omega]=\Phimap[\omega]$, resolved per outcome
pair, so every $P_{\hat c}$ is unchanged.
\emph{Objective.} The divergence terms depend on
$(\mathbf q,\Phimap[\cdot])$ only, hence are unchanged. For
$K_{\hat c}$: by (RM2) and Proposition~\ref{prop:exact-cancellation},
$f_{\hat c,s_i}\equiv0$ on \emph{every} input, before and after the
modification; $f_{\hat c,s}$ for other retained $s$ is unchanged; and the
term $f_{\hat c,s^{*}}(\omega_{s^{*}})\le0$
(Corollary~\ref{cor:f-nonpos}) is deleted. Hence
$K_{\hat c}(\hat\omega)\ge K_{\hat c}(\omega)$, and by monotonicity of the
kernel $d$ in its second argument,
$O_{\alpha}(\mathbf q,\hat\omega)\le O_{\alpha}(\mathbf q,\omega)$.
Therefore the restricted infimum is no larger than the unrestricted one;
being a restriction it is also no smaller; and the map
$\omega\mapsto\hat\omega$ carries minimizers to restricted minimizers.
\end{proof}

% claim meat-core/prop:retention
\claimsummary{meat-core/prop:retention}{A constraint-active sector adds no entropy deficit yet shapes the statistics: in a two-sector example, forcing it to zero raises the value from 0 to 1 bit per round. Dropping such a sector from the variable list can therefore overstate the key.}
\claimbookmark{meat-core/prop:retention}{Proposition}{prop:retention}{Unsoundness of dropping a constraint-active sector}
\claimleadin{%
A constraint-active sector contributes nothing to any $K_{\hat c}$, which
might suggest deleting it along with its variable. Its weight, however,
still enters the statistics, and an adversary may place it there.
}
\begin{proposition}[Unsoundness of dropping a constraint-active sector]
\label{prop:retention}
There exists a declaration satisfying Definition~\ref{def:admissible-declaration} and a
constraint-active sector $s_0$ such that forcing $\omega_{s_0}=0$ strictly
\emph{increases} $\hlow$. Dropping a constraint-active sector from the
variable list is therefore unsound: it can overstate the extractable key.
\end{proposition}
\claimmeta{\metalabel{Uses}\mref{Def.}{def:admissible-declaration}, \mref{Def.}{def:branch}, \mref{Prop.}{prop:trivial-alice-system}
\metasep \metalabel{Used by}\mref{Prop.}{prop:no-averaging}, \mref{Rem.}{rem:target-conditions}, \mref{Rem.}{rem:removal-sets}, \mref{Rem.}{rem:classes}
\metasep newly derived
\metasep implemented (effective: accepted)
\metasep \metalabel{Code}\coderef{https://github.com/fraunhofersg/certiqs-key/blob/v0.1.0/src/certiqs_key/meat_core/declaration/sectors.py}{certiqs_key.meat_core.declaration.sectors:build_sector_report}
\metasep \metalabel{Oracles}\coderef{https://github.com/fraunhofersg/certiqs-key/blob/v0.1.0/tests/oracles/registry.py}{registry.py}}

\begin{proof}
Take $d_A=1$ ($\tau_A=[1]$, equivalently \exttermref{../transmitter/contract.pdf}{term-untrusted-source-mode}{untrusted-source mode} by
Proposition~\ref{prop:trivial-alice-system}; single Alice outcome with
$M=1$, announcements
trivial, one announcement class, $\varsigma\equiv\mathrm{pass}$,
$\phi\equiv\hat c_0$, $\Chatset=\Chkey=\{\hat c_0\}$,
$\acSet=\{\delta_{\hat c_0}\}$), $\Zval=\{0,1\}$, and
$\Qsys=\Qsys_1\oplus\Qsys_2$ with both sectors one-dimensional. Declare
$\Vset=\{0,1\}$, $r(v)=v$, and
\begin{equation*}
  F^{(1)}_0=F^{(1)}_1=\tfrac12,
  \qquad
  F^{(2)}_0=1,\quad F^{(2)}_1=0 .
\end{equation*}
Sector 2 is key-inert (its \termref{term-effective-outcome-set}{effective outcome set} is $\{v=0\}$, a
singleton), hence constraint-active; sector 1 is key-active. A feasible
state is $\omega=(w_1,w_2)$, $w_i\ge0$, $w_1+w_2=1$, and
$P_{\hat c_0}=1$. A direct computation of \eqref{eq:GTheta} for sector~1
gives, for weight $w_1$, the rank-one operator
$\cG_{\hat c_0,1}(w_1)=w_1\dyad{\eta}$ with
$\ket{\eta}=\tfrac1{\sqrt2}\sum_{v}\ket{v}_R\otimes\ket{x_0,v}_{XY}
\otimes\ket{\kappa_0}$, so
$\cG^{1/\alpha}=w_1^{1/\alpha}\dyad{\eta}$,
$\cZ(\dyad\eta)=\tfrac12\sum_v\dyad{\xi_v}$ for orthonormal
$\ket{\xi_v}$, and
$T=w_1\,\tr[(\tfrac12\sum_v\dyad{\xi_v})^{\alpha}]=w_1\,2^{1-\alpha}$;
hence $f_{\hat c_0,1}(w_1)=w_1(2^{1-\alpha}-1)$ and
$f_{\hat c_0,2}\equiv0$, so
$K=1+w_1(2^{1-\alpha}-1)$ and
$O_{\alpha}=\frac{1}{\alpha-1}\log_2\frac{1}{K}$. The infimum over
$w_1\in[0,1]$ is attained at $w_1=0$ with value $0$; forcing
$\omega_{s_0}=\omega_2=0$ (i.e.\ $w_1=1$) gives the strictly larger value
$\frac{1}{\alpha-1}\log_2 2^{\alpha-1}=1$. The value $0$ is also the
correct one: an adversary that always populates sector~2 makes the raw key
deterministic. Hence the drop overstates the key by one bit per round.
\end{proof}

% claim meat-core/prop:no-averaging
\claimsummary{meat-core/prop:no-averaging}{Replacing the sectors of a declaration by one sector with averaged measurement operators can raise the value even where the averaged declaration reproduces the statistics, in the example from 0 to 1 bit per round. Merging sectors is therefore no sound simplification.}
\claimbookmark{meat-core/prop:no-averaging}{Proposition}{prop:no-averaging}{Unsoundness of averaging sectors}
\claimleadin{%
Theorem~\ref{thm:removal} deletes a sector by moving its weight onto other
sectors that reproduce its responses, and the sector index survives the
move. Averaging sectors into one erases that index, and the following
proposition shows that this is not sound.
}
\begin{proposition}[Unsoundness of averaging sectors]\label{prop:no-averaging}
Coarse-graining the sector decomposition of a declaration --- replacing
$\Qsys=\bigoplus_s\Qsys_s$ and $\{F^{(s)}_v\}$ by a single sector carrying
averaged POVM elements --- is unsound: there exists a declaration
$\Decl$ and an averaged declaration $\bar\Decl$ reproducing the same
outcome statistics on the averaged state space for which
$\hlow(\bar\Decl)>\hlow(\Decl)$.
\end{proposition}
\claimmeta{\metalabel{Uses}\mref{Thm.}{thm:removal}, \mref{Prop.}{prop:retention}
\metasep \metalabel{Used by}\mref{Rem.}{rem:classes}, \mref{Req.}{req:structure-operations}
\metasep newly derived
\metasep implemented (effective: accepted)
\metasep \metalabel{Code}\coderef{https://github.com/fraunhofersg/certiqs-key/blob/v0.1.0/src/certiqs_key/meat_core/declaration/sectors.py}{certiqs_key.meat_core.declaration.sectors:delete_sectors}
\metasep \metalabel{Oracles}\coderef{https://github.com/fraunhofersg/certiqs-key/blob/v0.1.0/tests/oracles/registry.py}{registry.py}}

\begin{proof}
Take the declaration of Proposition~\ref{prop:retention} with the two
sectors replaced by projective responses: sector $s$ responds
deterministically with outcome $v=s-1$, i.e.
\begin{equation*}
  F^{(1)}_0=1,\quad F^{(1)}_1=0,\qquad
  F^{(2)}_0=0,\quad F^{(2)}_1=1
\end{equation*}
(so $\sum_vF^{(s)}_v=1=\idon{\Qsys_s}$ in each sector; all other data
unchanged, and the raw key $r(v)=v$ takes the value $s-1$ on sector
$s$). Every pair
$(\hat c_0,s)$ is key-inert (each sector's effective outcome set is the
singleton $\{v=s-1\}$), so $f\equiv0$, $K=P=1$, and
$\hlow(\Decl)=0$ --- correctly: the sector index determines the raw key,
and the purifying register knows the sector. Now form $\bar\Decl$ on the
one-dimensional $\bar\Qsys$ with
$\bar F_v\coloneqq\tfrac12(F^{(1)}_v+F^{(2)}_v)=\tfrac12$ for both $v$
(a valid POVM, and the statistics of $\bar\Decl$ on its feasible state
$\bar\omega=1$ equal those of $\Decl$ at the uniform sector weights).
The computation of Proposition~\ref{prop:retention} with $w_1=1$ applies
verbatim to $\bar\Decl$ and gives $\hlow(\bar\Decl)=1$. An engine that
averaged the sectors would thus certify one secret bit per round for a
protocol whose key the adversary can know exactly. The obstruction is
structural: averaging erases the sector index from the state space while
the adversary retains it, so the averaged \exttermref{../protocol-compiler/contract.pdf}{term-feasible-set}{feasible set} no longer
contains (an image of) every reachable attack.
\end{proof}

% claim meat-core/prop:removal-compose
\claimsummary{meat-core/prop:removal-compose}{Several removable sectors may be deleted together when every deleted sector's witness targets only sectors that stay. The whole set may then be deleted without changing the value of the program or its attainment.}
\claimbookmark{meat-core/prop:removal-compose}{Proposition}{prop:removal-compose}{Deleting a set of removable sectors}
\claimleadin{%
Theorem~\ref{thm:removal} deletes one sector and moves its weight onto the
targets of its witness. Several deletions compose when no deletion removes a
sector that another deleted sector's weight is moved onto.
}
\begin{proposition}[Deleting a set of removable sectors]
\label{prop:removal-compose}
Let $\Decl$ satisfy Definition~\ref{def:admissible-declaration} and let
$\SecSet_{\mathrm{del}}\subseteq\SecSet$ be the set of sectors a run
deletes. Suppose \emph{every $s\in\SecSet_{\mathrm{del}}$
carries a witness in the sense of Definition~\ref{def:classes}\,(d) all of
whose targets lie outside $\SecSet_{\mathrm{del}}$}. Then restricting the
infimum \eqref{eq:hlow-def} to
$\{\omega\in\Feas:\omega_s=0\ \text{for all }s\in\SecSet_{\mathrm{del}}\}$
changes neither its value nor its attainment: the whole set
$\SecSet_{\mathrm{del}}$ may be deleted from the variable list, and its
rows from the generated maps, without loss.
\end{proposition}
\claimmeta{\metalabel{Uses}\mref{Def.}{def:admissible-declaration}, \mref{Def.}{def:keyinert}, \mref{Def.}{def:feasible}, \mref{Def.}{def:classes}, \mref{Thm.}{thm:removal}
\metasep \metalabel{Used by}\mref{Rem.}{rem:removal-general-format}, \mref{Rem.}{rem:removal-sets}, \mref{Rem.}{rem:classes}, \mref{Req.}{req:sector-report}, \mref{Req.}{req:structure-operations}, \mref{Req.}{req:single-sector-deletion}
\metasep newly derived
\metasep implemented (effective: accepted)
\metasep \metalabel{Code}\coderef{https://github.com/fraunhofersg/certiqs-key/blob/v0.1.0/src/certiqs_key/meat_core/declaration/sectors.py}{certiqs_key.meat_core.declaration.sectors:delete_sectors}}

\begin{proof}
By induction on $|\SecSet_{\mathrm{del}}|$. The case $0$ is vacuous and
the case $1$ is Theorem~\ref{thm:removal}. For the step, delete the
members one at a time, in any order, and check at each step that the
theorem applies to the \emph{reduced} declaration
$\Decl'$ obtained by deleting the sectors removed so far.
\emph{The sector being deleted is still state-blind, and every
key-inertness predicate is unchanged}: both are per-sector properties of
that sector's own declared blocks $F^{(s)}_v$, $\hat O^{(s)}_k$ and of
$M_x$, $a$, $r$, $\varphi$, $\varsigma$ (Definitions~\ref{def:keyinert}
and~\ref{def:classes}\,(c)), none of which deleting \emph{another} sector
touches. \emph{Its targets are still present}, never having been in
$\SecSet_{\mathrm{del}}$, so (W2) and (W5) read the same rows and (RM1),
(RM3) the same numbers. \emph{And $\Decl'$ satisfies
Definition~\ref{def:admissible-declaration}}: (iii) because $\{F_v\}$ is sector-block-diagonal
with $\sum_vF^{(s)}_v=\idon{\Qsys_s}$ per surviving sector, so its
restriction is a POVM on $\bigoplus_{s\notin\text{deleted}}\Hs{\Qsys_s}$,
and likewise for $\{M_x\}$, which is untouched; (vi) because $\Theta_k$
and $c_k$ are untouched and each $\hat O_k$ restricts to a
sector-block-diagonal operator; (vii) because Theorem~\ref{thm:removal}'s
own map $\omega\mapsto\hat\omega$ carries a point of $\Feas$ --- nonempty
by (vii) for the declaration before the step --- to a feasible point with
$\omega_{s^{*}}=0$, whose restriction is a point of the reduced feasible
set; and (i), (ii), (iv), (v), (viii) read no sector at all. So
Theorem~\ref{thm:removal} applies at each step and leaves the infimum and
its attainment unchanged. Finally the theorem's modified point adds mass
only to the targets, hence only outside $\SecSet_{\mathrm{del}}$, so no
already deleted sector is repopulated by a later step, and the objective
on the reduced declaration equals the restricted objective, the deleted
sector contributing $f_{\hat c,s^{*}}(0)=0$ and
$\Phimap_{s^{*}}[0]=0$.
\end{proof}

% claim meat-core/rem:target-conditions
\claimsummary{meat-core/rem:target-conditions}{Without (W2) every state-blind constraint-active sector witnesses its own removal, and deleting it can overstate the value by a bit per round. Distinct targets cost no removal in exact arithmetic: repeated targets merge into one witness certifying the same identities.}
\claimbookmark{meat-core/rem:target-conditions}{Remark}{rem:target-conditions}{The two index conditions of the removal checklist}
\claimleadin{%
Condition (W2) of Definition~\ref{def:classes} asks two things of a witness
for $s^{*}$: that it does not name $s^{*}$, and that it names no sector
twice. Both come from the construction in the proof of
Theorem~\ref{thm:removal}.
}
\begin{remark}[The two index conditions of the removal checklist]\label{rem:target-conditions}
\emph{Without (W2), every state-blind constraint-active sector witnesses
its own removal.} Let $s^{*}$ be state-blind and key-inert for every
$\hat c\in\Chkey$ --- that is, state-blind and constraint-active in the
sense of Definition~\ref{def:classes}\,(b) --- and take the
\emph{self-witness} $L=1$, $s_1\coloneqq s^{*}$, $\mu_1\coloneqq1$, and
$\chi_1$ any state on $\Qsys_{s^{*}}$ in the form (W4) requires --- a
diagonal projector $\dyad{i}$ will do. Then (W1) holds by hypothesis, (W3)
and (W4) hold by construction, (W5) holds because $s^{*}$ is itself
key-inert for every $\hat c\in\Chkey$, and \emph{both equality families
close exactly}: state-blindness gives
$\tr[F^{(s^{*})}_v\chi_1]=c_{v,s^{*}}\tr\chi_1=c_{v,s^{*}}$ and
$\tr[\hat O^{(s^{*})}_k\chi_1]=o_{k,s^{*}}\tr\chi_1=o_{k,s^{*}}$, and the
evaluation is the product $c_{v,s^{*}}\cdot1$, which rounds nothing, so
each residual is the datum $0$ and (W6) passes as well. Every condition
but (W2) is met, and a checklist without (W2) would flag the sector
removable and delete a constraint-active sector --- the one operation
Proposition~\ref{prop:retention} proves unsound. Its counterexample is of
exactly this shape: its sector~2 has $F^{(2)}_0=1$, $F^{(2)}_1=0$ on a
one-dimensional block, hence is state-blind, and is key-inert, hence
constraint-active; deleting it \emph{overstates $\hlow$ by one bit per
round}. What forbids the self-witness is the construction in the proof of
Theorem~\ref{thm:removal}: it sets $\hat\omega_{s^{*}}\coloneqq0$ and
$\hat\omega_{s_i}\coloneqq\omega_{s_i}+\mu_i\,m\otimes\chi_i$
\emph{simultaneously}, and if some $s_i=s^{*}$ those two assignments
contradict each other, so $\hat\omega$ is not defined and no step of the
proof is available. The theorem quantifies over
$s_1,\dots,s_L\in\SecSet\setminus\{s^{*}\}$ for that reason, and the
first half of (W2) is that quantifier and nothing else.

\emph{Distinct targets, at no cost.} With $s_i=s_j$ for $i\ne j$ the
proof's second assignment
$\hat\omega_{s_i}\coloneqq\omega_{s_i}+\mu_i\,m\otimes\chi_i$ would assign
two values to one block, the same collision the self-witness produces at
$s_i=s^{*}$, and $\hat\omega$ would again not be defined. The theorem
therefore quantifies over pairwise distinct targets and (W2) is that
quantifier as well. \emph{No removal is lost by it.} A candidate witness
with repeated targets \emph{merges}: drop every term with $\mu_i=0$, and
replace the terms at a repeated target $s$ by the single term
$(\bar\mu_s,\bar\chi_s)$ with $\bar\mu_s\coloneqq\sum_{i:s_i=s}\mu_i$ and
$\bar\chi_s\coloneqq\bar\mu_s^{-1}\sum_{i:s_i=s}\mu_i\chi_i$, defined
because $\bar\mu_s>0$ after the drop. Then $\bar\mu_s\ge0$ and
$\sum_s\bar\mu_s=1$, so (W3) holds; $\bar\chi_s\succeq0$ with
$\tr\bar\chi_s=1$, a convex combination of states, so (W4) holds; and
$\bar\mu_s\tr[F^{(s)}_v\bar\chi_s]=\sum_{i:s_i=s}\mu_i\tr[F^{(s)}_v\chi_i]$
by linearity of $X\mapsto\tr[F^{(s)}_vX]$, and likewise for
$\hat O^{(s)}_k$, so (RM1) and (RM3) take the same values and (W6)'s
identities are the same identities. \emph{In exact arithmetic the merged
witness certifies exactly what the original would have.} The merged
witness and its evaluation order must still be \emph{exhibited}, because
(W3), (W4) and (W6) are conditions on the exhibited numbers: a merge whose
arithmetic rounds sets no flag, which is the safe direction
(Remark~\ref{rem:classes}).
\end{remark}
\claimmeta{\metalabel{Uses}\mref{Def.}{def:classes}, \mref{Thm.}{thm:removal}, \mref{Prop.}{prop:retention}
\metasep newly derived
\metasep accepted}

% claim meat-core/rem:witness-exactness
\claimsummary{meat-core/rem:witness-exactness}{A witness is computed, not declared, so its weights and states can be made exact by construction; a naively summed weight vector can read 1 while its real sum exceeds 1. The diagonal form of (W4) is one exact presentation of a state, not the only one binary64 admits.}
\claimbookmark{meat-core/rem:witness-exactness}{Remark}{rem:witness-exactness}{Exactness of the exhibited weights and states}
\claimleadin{%
Conditions (W3) and (W4) of Definition~\ref{def:classes} read exhibited
\exttermref{../foundations/contract.pdf}{term-binary64-numbers}{binary64 numbers}, and the normalization step of Theorem~\ref{thm:removal}
uses $\sum_i\mu_i=1$ and $\tr\chi_i=1$ as identities of reals.
}
\begin{remark}[Exactness of the exhibited weights and states]\label{rem:witness-exactness}
\emph{Exact by construction.} (W3) and (W4) are decided at zero defect,
but not in the way the declared marginal $\tau_A$ is checked at admission.
That check decides \exttermref{../foundations/contract.pdf}{term-declared-data}{declared data}, which an \exttermref{../foundations/contract.pdf}{term-instance}{instance}
\exttermref{../foundations/contract.pdf}{term-snap}{snaps} under a \exttermref{../foundations/contract.pdf}{term-producer}{producer}'s obligation and a
\exttermref{../foundations/contract.pdf}{term-consumer}{consumer} must take as it finds them. A witness is instead a
\exttermref{../foundations/contract.pdf}{term-computed-object}{computed object} that the implementation exhibits, and the
implementation can make both conditions exact by construction --- which
(W3) and (W4) as stated require it to. The difference is on the false-positive side, which is the
unsafe one. A naively summed $\mu=(1,2^{-60})$ returns the datum $1.0$
while $\sum_i\mu_i>1$ as reals, and then the normalization step of
Theorem~\ref{thm:removal} reads
$\tr\hat\omega=\tr\omega+(\sum_i\mu_i-1)\tr m\ne1$: the modified point is
not feasible and the removal would have been licensed on an identity that
does not hold. Requiring the sum in the exhibited order and verified to
round nowhere closes that window exactly as (W6) closes it for the
residuals, and requiring the exhibited form of (W4) closes the
corresponding window on a computed $\lmin(\chi_i)$ or a computed
$\tr\chi_i$. \emph{Neither leaves a residual exposure}: on the forms those
two conditions admit, both hold of the real numbers the data denote.

\emph{The diagonal form is one exact presentation among others.} Any
presentation deciding $\chi_i\succeq0$ and $\tr\chi_i=1$ exactly by
construction would serve, and the diagonal form is \emph{one} such
presentation, not the only one binary64 admits. An exhibited unit vector's
outer product $\dyad{\psi_i}$ with $\braket{\psi_i}{\psi_i}$ verified to
round nowhere is a genuine second: on a four-dimensional $\Qsys_{s_i}$ the
vector $\psi_i=(\tfrac12,\tfrac12,\tfrac12,\tfrac12)$ has exactly
representable entries, its norm accumulates $0.25\to0.5\to0.75\to1.0$
rounding nowhere, and its outer product --- every entry $\tfrac14$ --- is a
rank-one state that is \emph{not} diagonal. (W4) does not admit it, and
excluding an admissible presentation can only \emph{refuse} a removal. No
impossibility is claimed here.

\emph{The evaluation order of (W6).} The order in which a witness's
residuals are formed is exhibited with the witness. This is the one
difference from the \exttermref{../foundations/contract.pdf}{term-declared-order}{declared order} of \eqref{eq:declorder}, which is
part of the declaration because the datum it governs is declared.
\end{remark}
\claimmeta{\metalabel{Uses}\mref{Def.}{def:classes}, \mref{Thm.}{thm:removal}
\metasep newly derived
\metasep accepted}

% claim meat-core/rem:removal-general-format
\claimsummary{meat-core/rem:removal-general-format}{A party-local entry makes clause (d) unsatisfiable, since the removal theorem excludes it; several measured factors leave the checklist no declared value to read, though the theorem can hold. Either way no sector is flagged, so no set of sectors is deleted.}
\claimbookmark{meat-core/rem:removal-general-format}{Remark}{rem:removal-general-format}{Removal in the general format}
\claimleadin{%
Requirement~\ref{req:removal-checklist} sets the removable flag for no
sector of a declaration whose list carries an entry declared
\exttermref{../protocol-compiler/contract.pdf}{term-party-local}{party-local}, nor of one with more than one \exttermref{../protocol-compiler/contract.pdf}{term-measured-factor}{measured factor}.
The two grounds differ in kind.
}
\begin{remark}[Removal in the general format]\label{rem:removal-general-format}
\emph{A party-local entry is excluded by the theorem.}
Theorem~\ref{thm:removal} is stated for a list carrying no entry declared
party-local, so for a list carrying one, clause~(d) of
Definition~\ref{def:classes} is unsatisfiable and no sector carries a
witness in its sense.

\emph{Several measured factors are excluded by the checklist, not by the
theorem.} The theorem's proof reads the sector blocks as operators and
needs no factorisation of them, so its hypotheses can hold of a declaration
with more than one measured factor. What fails there is the decision
procedure: the blocks $F^{(s)}_v$ are products \eqref{eq:productfamily}
that are never formed as data, so the bit-level test of (W1) and the
residuals of (W6) have no \exttermref{../foundations/contract.pdf}{term-declared-value}{declared value} to read. This ground is a
refusal of the flag and not a soundness ground.

\emph{Both close the deletion of a set.} Both refusals retain sectors,
which is unconditionally sound, and both are stated of the flag. The flag is
what closes the composite deletion: the hypothesis of
Proposition~\ref{prop:removal-compose} quantifies over sectors carrying a
witness in the sense of~(d), and a run deletes a set of sectors only on the
strength of the flag, which is set for none of them in either case. So no
composite deletion is available for a declaration carrying a party-local
entry or measuring more than one factor, and the proposition is not reached
there.
\end{remark}
\claimmeta{\metalabel{Uses}\mref{Def.}{def:classes}, \mref{Thm.}{thm:removal}, \mref{Prop.}{prop:removal-compose}
\metasep newly derived
\metasep accepted}

% claim meat-core/rem:removal-sets
\claimsummary{meat-core/rem:removal-sets}{Per-sector removal flags do not compose: two sectors can witness each other's removal, and deleting both raises the value by a bit per round. The set hypothesis of \hyperlink{claim-meat-core/prop:removal-compose}{Proposition~\ref*{prop:removal-compose}} is sufficient, not necessary, and is read on the recorded witnesses.}
\claimbookmark{meat-core/rem:removal-sets}{Remark}{rem:removal-sets}{The set hypothesis of the composition rule}
\begin{remark}[The set hypothesis of the composition rule]\label{rem:removal-sets}
\emph{The hypothesis of Proposition~\ref{prop:removal-compose} is on the
\emph{set}, not on any one sector, and without it the per-sector flags do
not compose.} The flag \emph{removable} is a per-sector flag and
Theorem~\ref{thm:removal} licenses the deletion of \emph{one} sector;
neither says that deleting every flagged sector at once is sound, and
\emph{in general it is not}. Two state-blind
constraint-active sectors with identical declared
responses each witness the other: take the declaration of
Proposition~\ref{prop:retention} and split its sector~2 into two
one-dimensional sectors $2a$ and $2b$ carrying the same $F_0=1$, $F_1=0$
(all other data unchanged, so $f_{\hat c_0,2a}=f_{\hat c_0,2b}\equiv0$,
$K=1+w_1(2^{1-\alpha}-1)$ and the infimum is still $0$, at $w_1=0$). Then
$\{2b\}$ with $\mu_1=1$ and $\chi_1$ the state on $\Qsys_{2b}$ is a
witness for $2a$ meeting (W1)--(W6) exactly, and symmetrically for $2b$,
so both flags are set with (W2) satisfied at every step --- and deleting
both forces $w_1=1$ and raises the reported value from $0$ to $1$, the
same one bit per round by the same mechanism as
Proposition~\ref{prop:retention}. The hypothesis of
Proposition~\ref{prop:removal-compose} is what fails there, and it is the
only thing that fails.

\emph{Three readings of that hypothesis.} \emph{First, it is sufficient and
not necessary.} With
recorded single-target witnesses forming a chain
$s_a\to s_b\to s_c\to s_d$ and
$\SecSet_{\mathrm{del}}=\{s_a,s_b,s_c\}$ ($s_d$ retained), the hypothesis
fails and the deletion is refused --- although deleting all three is in
fact sound: $s_c$ may be deleted first, and $s_b$ is then witnessed by
$s_c$'s own witness, since $s_c$ is state-blind by (W1) and therefore
$\tr[F^{(s_c)}_v\chi]=c_{v,s_c}$ for every state $\chi$, so
$c_{v,s_b}=c_{v,s_c}=\sum_j\nu_j\tr[F^{(s_d)}_v\xi_j]$ and (RM1) closes
with $(\nu_j,\xi_j)$ in place of the original targets, and likewise
(RM3). Refusal is the safe side and this module claims sufficiency
only. \emph{Second, the hypothesis is decided on the \emph{recorded}
witnesses}, so it constrains the record and not the search: a run may
re-derive each surviving sector's witness in the reduced declaration after
every deletion and record \emph{those}. \emph{Nothing here forbids
sequential re-witnessing}; what is forbidden is deleting a set whose
recorded witnesses point into it. \emph{But the strategy is
order-sensitive, and the order is part of it}: re-derivation runs on the
\emph{surviving} sectors, so it can drop a recorded reference to a sector
already deleted and never to one deleted \emph{later}. What a run must
therefore achieve is that every recorded witness names only sectors it
\emph{retains}, and the way to achieve it is to delete leaf-first --- a
sector only when its current witness names retained sectors alone. On the
chain $s_a\to s_b\to s_c\to s_d$ above with
$\SecSet_{\mathrm{del}}=\{s_a,s_b,s_c\}$ the order $s_c,s_b,s_a$ does it:
after $s_c$ goes, $s_b$ re-derives onto $s_d$ by the argument just given,
then $s_a$ likewise, and the final record is
$s_a\to s_d$, $s_b\to s_d$, $s_c\to s_d$, which meets the hypothesis and
deletes all three. \emph{The order $s_a,s_b,s_c$ does not}: re-derivation
never touches $s_a$'s or $s_b$'s already-recorded witnesses, the record
ends as $s_a\to s_b$, $s_b\to s_c$ with both targets inside
$\SecSet_{\mathrm{del}}$, and the deletion of the set is refused as a
whole. On that record, with $s_c\to s_d$ beside them, a subset of
$\SecSet_{\mathrm{del}}$ meets the hypothesis exactly when it does not
contain $s_b$ together with $s_a$ or $s_c$; the largest is
$\{s_a,s_c\}$, whose recorded targets $s_b$ and $s_d$ it retains, so the
run keeps at most the deletion of $s_a$ and $s_c$, taken as a set of its
own, and retains $s_b$. \emph{So re-witnessing loses nothing
\emph{when such an order exists}}, and where none does --- the two
mutually witnessing sectors of the counterexample above, whose witness
relation is cyclic --- no order helps and the refusal is the correct
outcome, that being precisely the case the hypothesis of
Proposition~\ref{prop:removal-compose} exists to refuse. \emph{Third, the
comparison is of index data and is decided at no window}, exactly like
(W2), and a singleton $\SecSet_{\mathrm{del}}$ satisfies the hypothesis
automatically, by (W2).
\end{remark}
\claimmeta{\metalabel{Uses}\mref{Def.}{def:classes}, \mref{Thm.}{thm:removal}, \mref{Prop.}{prop:retention}, \mref{Prop.}{prop:removal-compose}
\metasep newly derived
\metasep accepted}

% claim meat-core/rem:classes
\claimsummary{meat-core/rem:classes}{Retaining every sector computes the value as defined; deleting one on an approximate match computes a restriction whose excess no function of a tolerance bounds, so a tolerance may only refuse a removal. Dropping a constraint-active sector and averaging sectors can both overstate the key.}
\claimbookmark{meat-core/rem:classes}{Remark}{rem:classes}{Retention and deletion of sectors}
\begin{remark}[Retention and deletion of sectors]\label{rem:classes}
\emph{The direction of every removal decision.} Removability of
Definition~\ref{def:classes}\,(d) is a licence to delete, and every outcome
other than all six of (W1)--(W6) --- a nonzero residual, a rounding
operation, an absent witness, a witness that names $s^{*}$, a target that
is not key-inert, a block that is not exactly scalar --- retains the
sector. Retention is unconditionally sound and costs nothing: $\hlow$ is
\emph{defined} by \eqref{eq:hlow-def} as the infimum over the whole of
$\Feas$, and Theorem~\ref{thm:removal} says only that one particular
restriction of that infimum does not change it. A run that retains every
sector computes the definition. A run that deletes one on an
\emph{approximate} match computes instead the infimum over
$\{\omega\in\Feas:\omega_{s^{*}}=0\}$, a restriction, whose value is
$\ge\hlow$ and which would be reported under the name $\hlow$ --- an
overstatement, and one no function of a \exttermref{../foundations/contract.pdf}{term-declared-tolerance}{declared tolerance} bounds:
the mechanism is a $P_{\hat c}$ or $K_{\hat c}$ driven towards zero inside a
logarithm, which Proposition~\ref{prop:value-range}\,(iii) exhibits driving
the objective up without limit and which
Proposition~\ref{prop:optol-nobound}\,(i) turns, in the neighbouring
setting of a perturbed declared POVM, into two declarations at distance
$\epsilon$ whose values differ by tens of bits per round --- by
$\frac{\alpha}{\alpha-1}N$ with $N$ up to $52$, that ceiling being the
number model's and not the mechanism's. So a tolerance may only ever
\emph{refuse} a removal, never license one.

\emph{Two unsound operations.} Dropping a constraint-active sector
(Proposition~\ref{prop:retention}) and averaging sectors
(Proposition~\ref{prop:no-averaging}) are both unsound, and both failure
modes overstate the key --- the dangerous direction. The first is reachable
in two ways from a checklist that looks complete: by dropping (W2), which
lets a sector witness its own removal (Remark~\ref{rem:target-conditions}),
and by deleting every removable sector at once, which lets two sectors
witness each other's (Remark~\ref{rem:removal-sets}). (W2) refuses the
first; the hypothesis of Proposition~\ref{prop:removal-compose} refuses the
second, and is the reason removability being a property of one sector does
not by itself license the deletion of a set.
\end{remark}
\claimmeta{\metalabel{Uses}\mref{Def.}{def:feasible}, \mref{Prop.}{prop:value-range}, \mref{Def.}{def:classes}, \mref{Thm.}{thm:removal}, \mref{Prop.}{prop:retention}, \mref{Prop.}{prop:no-averaging}, \mref{Prop.}{prop:removal-compose}
\metasep newly derived
\metasep accepted}

\section{The Choi program}\label{group-meat-core/lem:admissible-alice-subspace}
\sectionintro{Every feasible state lives on the admissible Alice subspace, the support of the declared marginal (all of Alice's system in untrusted-source mode), tensored with the measured system (\hyperlink{claim-meat-core/lem:admissible-alice-subspace}{Lemma~\ref*{lem:admissible-alice-subspace}}). In trusted-source mode every state with Alice's marginal is the image of a positive Choi operator on a \termref{term-support-frame}{support frame} that need only contain the marginal's support, so no rank of the marginal is decided (\hyperlink{claim-meat-core/lem:choi-marginal}{Lemma~\ref*{lem:choi-marginal}}). \hyperlink{claim-meat-core/thm:choi}{Theorem~\ref*{thm:choi}} rewrites the program exactly in sector Choi blocks, with one coupling constraint per Alice sector, and on a tight frame a homogeneous constraint row may drop its conjugation (\hyperlink{claim-meat-core/prop:frame-free-rows}{Proposition~\ref*{prop:frame-free-rows}}); in untrusted-source mode the parametrisation is the identity (\hyperlink{claim-meat-core/def:untrusted-choi}{Definition~\ref*{def:untrusted-choi}}). The section's main result, \hyperlink{claim-meat-core/lem:estimate}{Lemma~\ref*{lem:estimate}}, states that the objective at any exactly feasible point is at least h\_low: an upper estimate of the value, never a bound on it.}
\sectionimports{\extref{../foundations/contract.pdf}{foundations/def:sector-map}}

% claim meat-core/lem:admissible-alice-subspace
\claimsummary{meat-core/lem:admissible-alice-subspace}{Every feasible state is supported in the admissible Alice subspace, the support of the declared marginal or all of Alice's system, tensored with the measured system; with no constraints the bound is tight. The subspace depends on the marginal slot alone.}
\claimbookmark{meat-core/lem:admissible-alice-subspace}{Lemma}{lem:admissible-alice-subspace}{The admissible Alice subspace}
\claimleadin{%
This is where the declared marginal enters the support and face rules of the
gradients: the \termref{term-maximal-branch-support}{maximal branch support} of
Definition~\ref{def:branch-support} is formed from $\Sadm$, and for a
nonempty $\CList$ the inclusion~(i) is all that the bound
\eqref{eq:Wcs-bound} of Lemma~\ref{lem:branch-support} uses.
The marginal also enters the gradient formulas, through the $W_{\tau}$
conjugation in \eqref{eq:grad-pieces} of
Corollary~\ref{cor:full-gradient}, which \eqref{eq:grad-pieces-un} drops.
}
\begin{lemma}[The admissible Alice subspace]\label{lem:admissible-alice-subspace}
Let $\Decl$ be a declaration, and write
\begin{equation}\label{eq:SA}
  \Sadm\coloneqq
  \begin{cases}
    \operatorname{supp}\tau_A, & \text{trusted-source mode},\\[0.3ex]
    \Hs{A}, & \text{untrusted-source mode},
  \end{cases}
\end{equation}
the \emph{admissible Alice subspace}; it is determined by the
\exttermref{../protocol-compiler/contract.pdf}{term-marginal-slot}{marginal slot} alone and never by $\CList$ or by a particular feasible
point. Then:
\begin{enumerate}[label=(\roman*),leftmargin=2.2em]
\item every $\omega\in\Feas$ satisfies
$\operatorname{supp}\omega\subseteq\Sadm\otimes\Hs{\Qsys}$;
\item the bound~(i) is tight for $\CList=\varnothing$, and for a nonempty
$\CList$ it remains a valid upper bound.
\end{enumerate}
\end{lemma}
\claimmeta{\metalabel{Uses}\mref{Def.}{def:admissible-declaration}, \mref{Def.}{def:feasible}
\metasep \metalabel{Used by}\mref{Thm.}{thm:choi}, \mref{Def.}{def:branch-support}, \mref{Lem.}{lem:branch-support}
\metasep newly derived
\metasep accepted}

\begin{proof}
(i) In \exttermref{../transmitter/contract.pdf}{term-untrusted-source-mode}{untrusted-source mode} this is vacuous. In
\exttermref{../transmitter/contract.pdf}{term-trusted-source-mode}{trusted-source mode}, with
$\Pi^{\perp}\coloneqq\idon{A}-\Pi_{\operatorname{supp}\tau_A}$,
$\tr[(\Pi^{\perp}\otimes\idon{\Qsys})\,\omega]=\tr[\Pi^{\perp}\tau_A]=0$
forces $(\Pi^{\perp}\otimes\idon{})\,\omega=0$ by positivity.

(ii) For $\CList=\varnothing$ take $\omega=\tau_A\otimes\chi$ with $\chi$ any
full-rank sector-block-diagonal state on $\Qsys$, for instance
$\idon{\Qsys}/d_{\Qsys}$, respectively any full-rank sector-block-diagonal
state $\omega$ on $A\Qsys$, for instance $\idon{A\Qsys}/(d_Ad_{\Qsys})$.
Such an $\omega$ lies in $\Feas$, and
$\operatorname{supp}\omega=\operatorname{supp}\tau_A\otimes\Hs{\Qsys}$,
respectively $\operatorname{supp}\omega=\Hs{A\Qsys}$, which is
$\Sadm\otimes\Hs{\Qsys}$. For a nonempty $\CList$, (i) holds as proved.
\end{proof}

% claim meat-core/def:support-frame
\claimsummary{meat-core/def:support-frame}{In trusted-source mode a support frame is the range of an isometry \ensuremath{\iota} from a Choi input register A\ensuremath{_{0}} into A that contains the support of \ensuremath{\tau}\_A, not necessarily equal to it. Two exact constructions are named: the full frame, always available, and the diagonal frame, tight by construction.}
\claimbookmark{meat-core/def:support-frame}{Definition}{def:support-frame}{Support frame and declared isometry}
\claimleadin{%
The Choi program parametrises the states with Alice's declared marginal
$\tau_A$ by positive operators on a register $A_0$ mapped into $A$. The
frame is asked only to contain the support of $\tau_A$, not to equal
it, so that no run has to decide $\operatorname{rank}\tau_A$ from a computed
spectrum at a window. In \exttermref{../transmitter/contract.pdf}{term-untrusted-source-mode}{untrusted-source mode} there is no marginal, and no
$A_0$, $\iota$ or $W_{\tau}$ is defined or needed
(Definition~\ref{def:untrusted-choi}); at $d_A=1$ the two modes give the
same program (Definition~\ref{def:untrusted-choi}\,(d)).
}
\begin{definition}[Support frame and declared isometry]\label{def:support-frame}
Let the declaration be in \exttermref{../transmitter/contract.pdf}{term-trusted-source-mode}{trusted-source mode}, with declared marginal
$\tau_A$. Let $A_0$ be a register of dimension $d_{A_0}$ and let
$\iota\in\Linn{A_0}{A}$ satisfy
\begin{equation}\label{eq:iota}
  \iota^{\dagger}\iota=\idon{A_0},
  \qquad
  \iota\,\iota^{\dagger}\,\tau_A=\tau_A .
\end{equation}
Write $\Pi_{\iota}\coloneqq\iota\iota^{\dagger}$, which is the orthogonal
projector onto $\operatorname{ran}\iota$
($\Pi_{\iota}^{\dagger}=\Pi_{\iota}$ and
$\Pi_{\iota}^{2}=\iota(\iota^{\dagger}\iota)\iota^{\dagger}=\Pi_{\iota}$
by the first condition). The first condition says $\iota$ is an isometry;
the second says its range \emph{contains} the support of $\tau_A$, since
$\Pi_{\iota}\tau_A=\tau_A$ holds exactly when every column of $\tau_A$
lies in $\operatorname{ran}\iota$, i.e.\ when
$\operatorname{ran}\tau_A\subseteq\operatorname{ran}\iota$. We call
$\operatorname{ran}\iota$ the
\termdef{term-support-frame}{support frame}
of the run and $\iota$ the
\termdef{term-declared-isometry}{declared isometry};
a support frame is given by any $\iota$ satisfying \eqref{eq:iota}. The
frame is
\termdef{term-tight}{tight}
when $\operatorname{ran}\iota=\operatorname{supp}\tau_A$, equivalently when
$d_{A_0}=\operatorname{rank}\tau_A$; in general
$\operatorname{rank}\tau_A\le d_{A_0}\le d_A$. The images under $\iota$
of a fixed orthonormal basis of $A_0$ are the
\termdef{term-declared-basis-of-a-0}{declared basis of $A_0$}.

Two constructions satisfy \eqref{eq:iota} exactly:
\begin{enumerate}[label=(\alph*),leftmargin=2.2em]
\item \emph{The
  \termdef{term-full-frame}{full frame}.}
  $A_0=A$, $d_{A_0}=d_A$, $\iota=\idon{A}$.
  Then $\iota^{\dagger}\iota=\idon{A}$ and
  $\Pi_{\iota}\tau_A=\idon{A}\,\tau_A=\tau_A$: both conditions of
  \eqref{eq:iota} hold for \emph{any} admissible $\tau_A$, and both hold
  \emph{exactly} as binary64 data, a product with the identity rounding
  nothing. This construction is always available and uses no property of
  $\tau_A$ beyond admission. It is tight exactly when $\tau_A$ is
  nonsingular.
\item \emph{The
  \termdef{term-diagonal-frame}{diagonal frame}.}
  If the declared $\tau_A$ is diagonal in
  the declared basis of $A$, $\tau_A=\sum_it_i\dyad{i}$, then
  $\operatorname{supp}\tau_A=\operatorname{span}\{\ket i:t_i\neq0\}$, and
  one may take $d_{A_0}=\#\{i:t_i\neq0\}$ and
  $\iota=\sum_j\ket{i_j}\bra{j}$, the coordinate injection onto those
  indices in increasing order. Both conditions of \eqref{eq:iota} are then
  arithmetic in $\{0,1\}$ and hold exactly, and the frame is tight.
  Whether $t_i\neq0$ is a bit-level comparison of a declared binary64
  datum against zero and not a threshold: $\tau_A$ is admitted only with
  $\tau_A\succeq0$ and $\tr\tau_A=1$ decided exactly
  (Definition~\ref{def:declaration}, the exact-admissibility paragraph), so
  its diagonal entries \emph{are}
  its eigenvalues and \emph{are} the declared numbers.
  The construction is available only for a $\tau_A$ that is diagonal in the
  declared basis, every off-diagonal entry being the datum $0$.
  Construction~(a) is always available
  and exhibits no tightness at all.
\end{enumerate}
Set
\begin{equation}\label{eq:Wtau}
  W_{\tau}\coloneqq\tau_A^{1/2}\,\iota\ \in\ \Linn{A_0}{A}.
\end{equation}
\end{definition}
\claimmeta{\metalabel{Uses}\mref{Def.}{def:declaration}
\metasep \metalabel{Used by}\mref{Lem.}{lem:choi-marginal}, \mref{Lem.}{lem:choi-factor}, \mref{Thm.}{thm:choi}, \mref{Prop.}{prop:frame-free-rows}, \mref{Rem.}{rem:tightness-by-construction}, \mref{Rem.}{rem:basis-invariance}, \mref{Def.}{def:branch-support}, \mref{Lem.}{lem:kraus-freedom}, \mref{Req.}{req:frame-constructed}, \mref{Req.}{req:diagonal-frame-test}, \mref{Req.}{req:frame-recorded}, \mref{Req.}{req:row-form-by-construction}, \mref{Prop.}{prop:relaxed-set-properties}, \mref{Req.}{req:verified-feasible-side}, \mref{Req.}{req:certified-evaluation}
\metasep newly derived
\metasep implemented
\metasep \metalabel{Code}\coderef{https://github.com/fraunhofersg/certiqs-key/blob/v0.1.0/src/certiqs_key/meat_core/objective/joint.py}{certiqs_key.meat_core.objective.joint:marginal_frame}}

% claim meat-core/lem:choi-marginal
\claimsummary{meat-core/lem:choi-marginal}{On every support frame, W\_\ensuremath{\tau} times its adjoint is \ensuremath{\tau}\_A, and the Choi map carries the positive operators with identity input marginal onto the states with A-marginal \ensuremath{\tau}\_A. It is injective on a tight frame, and only there unless the second system is trivial; being onto, it needs no rank decided.}
\claimbookmark{meat-core/lem:choi-marginal}{Lemma}{lem:choi-marginal}{Fixed-marginal parametrisation on a support frame}
\claimleadin{%
Whatever \termref{term-support-frame}{support frame} a run uses, every state with Alice's marginal is the
image of a positive Choi operator on it; a tight frame makes the
correspondence one to one. The identities (i)--(iii) are what the
parametrisation (iv) and the Choi program of Theorem~\ref{thm:choi} use of
$W_{\tau}$.
}
\begin{lemma}[Fixed-marginal parametrisation on a support frame]\label{lem:choi-marginal}
Let the declaration be in \exttermref{../transmitter/contract.pdf}{term-trusted-source-mode}{trusted-source mode}, let $\iota$ be the
\termref{term-declared-isometry}{declared isometry} of a \termref{term-support-frame}{support frame}, satisfying
\eqref{eq:iota}, and let $W_{\tau}$ be \eqref{eq:Wtau}
(Definition~\ref{def:support-frame}). Then:
\begin{enumerate}[label=(\roman*),leftmargin=2.2em]
\item $\Pi_{\iota}\tau_A^{1/2}=\tau_A^{1/2}$.
\item For every support frame, tight or not,
\begin{equation}\label{eq:WWtau}
  W_{\tau}W_{\tau}^{\dagger}
  =\tau_A^{1/2}\,\Pi_{\iota}\,\tau_A^{1/2}
  =\tau_A^{1/2}\,\tau_A^{1/2}=\tau_A .
\end{equation}
\item $\operatorname{ran}W_{\tau}=\operatorname{supp}\tau_A$, and the
Moore--Penrose inverse $W_{\tau}^{+}$ satisfies
\begin{equation}\label{eq:Wplus}
  W_{\tau}W_{\tau}^{+}=\Pi_{\operatorname{supp}\tau_A},
  \qquad
  W_{\tau}^{+}W_{\tau}=P_{\iota}
  \coloneqq\Pi_{(\ker W_{\tau})^{\perp}} ,
\end{equation}
where $\ker W_{\tau}=\{\xi\in\Hs{A_0}:
\iota\xi\in(\operatorname{supp}\tau_A)^{\perp}\}$ has dimension
$d_{A_0}-\operatorname{rank}\tau_A$. The frame is tight if and only if
$\ker W_{\tau}=0$, i.e.\ $P_{\iota}=\idon{A_0}$, and then $W_{\tau}$ has
full column rank and
$W_{\tau}^{+}=(W_{\tau}^{\dagger}W_{\tau})^{-1}W_{\tau}^{\dagger}$.
\item For any finite-dimensional system $B'$, the map
$\chi_{B'}(J)\coloneqq(W_{\tau}\otimes\idon{B'})\,J\,
(W_{\tau}^{\dagger}\otimes\idon{B'})$ maps
$\{J\in\Pos{A_0B'}:\tr_{B'}J=\idon{A_0}\}$ \emph{onto}
$\{\omega\in\Pos{AB'}:\tr_{B'}\omega=\tau_A\}$. It is \emph{injective},
hence a bijection with inverse
$\omega\mapsto(W_{\tau}^{+}\otimes\idon{B'})\,\omega\,
(W_{\tau}^{+\dagger}\otimes\idon{B'})$, when the support frame is tight,
and, if $\dim\Hs{B'}\ge2$, only then; if $\dim\Hs{B'}=1$ its domain is the
single point $\idon{A_0}$. In general $\chi_{B'}(J)=\chi_{B'}(J')$ if and only if
$(P_{\iota}\otimes\idon{B'})J(P_{\iota}\otimes\idon{B'})
=(P_{\iota}\otimes\idon{B'})J'(P_{\iota}\otimes\idon{B'})$.
\end{enumerate}
\end{lemma}
\claimmeta{\metalabel{Uses}\mref{Def.}{def:admissible-declaration}, \mref{Def.}{def:support-frame}
\metasep \metalabel{Used by}\mref{Lem.}{lem:choi-factor}, \mref{Thm.}{thm:choi}, \mref{Prop.}{prop:frame-free-rows}, \mref{Rem.}{rem:tightness-by-construction}, \mref{Rem.}{rem:ordering}, \mref{Rem.}{rem:basis-invariance}, \mref{Lem.}{lem:branch-support}, \mref{Req.}{req:frame-constructed}
\metasep newly derived
\metasep implemented (effective: accepted)
\metasep \metalabel{Code}\coderef{https://github.com/fraunhofersg/certiqs-key/blob/v0.1.0/src/certiqs_key/meat_core/objective/joint.py}{certiqs_key.meat_core.objective.joint:marginal_frame}
\metasep \metalabel{Oracles}\coderef{https://github.com/fraunhofersg/certiqs-key/blob/v0.1.0/tests/oracles/registry.py}{registry.py}}

\begin{proof}
\emph{Idea.} (i)--(iii) are linear algebra of the two conditions of
\eqref{eq:iota}. For (iv), the image of an operator with input marginal
$\idon{A_0}$ has $A$-marginal $W_{\tau}W_{\tau}^{\dagger}=\tau_A$;
conversely $W_{\tau}^{+}$ pulls a state with marginal $\tau_A$ back to an
operator with input marginal $P_{\iota}$, which a block annihilated by
$W_{\tau}$ completes to $\idon{A_0}$.

\emph{Step 1 (identity (i)).} For a positive semidefinite $X$ one has
$\operatorname{ran}X^{1/2}=(\ker X)^{\perp}=\operatorname{ran}X$, so
$\operatorname{ran}\tau_A^{1/2}=\operatorname{ran}\tau_A
\subseteq\operatorname{ran}\Pi_{\iota}$ by \eqref{eq:iota}, and a
projector acts as the identity on a subspace of its range.

\emph{Step 2 (identity (ii)).} The first equality of \eqref{eq:WWtau} is
$W_{\tau}=\tau_A^{1/2}\iota$ with $\Pi_{\iota}=\iota\iota^{\dagger}$, and
the second is (i).

\emph{Step 3 (identity (iii)).} By (ii),
$\operatorname{ran}W_{\tau}
=\operatorname{ran}(W_{\tau}W_{\tau}^{\dagger})
=\operatorname{supp}\tau_A$, so the Moore--Penrose inverse $W_{\tau}^{+}$
satisfies \eqref{eq:Wplus}. The kernel $\ker W_{\tau}$ has dimension
$d_{A_0}-\operatorname{rank}\tau_A$ because $\iota$ is injective and
$\operatorname{ran}\iota=\operatorname{supp}\tau_A\oplus
\big(\operatorname{ran}\iota\cap(\operatorname{supp}\tau_A)^{\perp}\big)$.
So $\ker W_{\tau}=0$ exactly when $d_{A_0}=\operatorname{rank}\tau_A$,
which is tightness (Definition~\ref{def:support-frame}), and then
$W_{\tau}$ has full column rank and
$W_{\tau}^{+}=(W_{\tau}^{\dagger}W_{\tau})^{-1}W_{\tau}^{\dagger}$.

\emph{Step 4 (the map (iv): forward).} $\chi_{B'}(J)\succeq0$, and
$\tr_{B'}\chi_{B'}(J)=W_{\tau}(\tr_{B'}J)W_{\tau}^{\dagger}
=W_{\tau}W_{\tau}^{\dagger}=\tau_A$ by \eqref{eq:WWtau}.

\emph{Step 5 (surjectivity).} Let $\omega\succeq0$ with
$\tr_{B'}\omega=\tau_A$ and let
$\Pi^{\perp}\coloneqq\idon{A}-\Pi_{\operatorname{supp}\tau_A}$. Then
$\tr[(\Pi^{\perp}\otimes\idon{B'})\,\omega]
=\tr[\Pi^{\perp}\tau_A]=0$, and positivity forces
$(\Pi^{\perp}\otimes\idon{})\,\omega=0$, i.e.\ $\omega$ is supported on
$\operatorname{supp}\tau_A\otimes\Hs{B'}$. Put
$J_0\coloneqq(W_{\tau}^{+}\otimes\idon{})\,\omega\,
(W_{\tau}^{+\dagger}\otimes\idon{})\succeq0$; by \eqref{eq:WWtau} and
\eqref{eq:Wplus},
\begin{equation*}
  \tr_{B'}J_0=W_{\tau}^{+}\tau_AW_{\tau}^{+\dagger}
  =W_{\tau}^{+}W_{\tau}\,W_{\tau}^{\dagger}W_{\tau}^{+\dagger}
  =P_{\iota}P_{\iota}^{\dagger}=P_{\iota},
\end{equation*}
and $\chi_{B'}(J_0)
=(W_{\tau}W_{\tau}^{+}\otimes\idon{})\,\omega\,
(W_{\tau}W_{\tau}^{+}\otimes\idon{})^{\dagger}
=(\Pi_{\operatorname{supp}\tau_A}\otimes\idon{})\,\omega\,
(\Pi_{\operatorname{supp}\tau_A}\otimes\idon{})=\omega$. If the frame is
tight, $P_{\iota}=\idon{A_0}$ and $J_0$ is already in the domain. In
general fix any state $\varphi$ on $\Hs{B'}$ and put
$J\coloneqq J_0+(\idon{A_0}-P_{\iota})\otimes\varphi\succeq0$. Then
$\tr_{B'}J=P_{\iota}+(\idon{A_0}-P_{\iota})=\idon{A_0}$, and the added
term is annihilated by $\chi_{B'}$, because
$W_{\tau}(\idon{A_0}-P_{\iota})
=W_{\tau}-W_{\tau}W_{\tau}^{+}W_{\tau}=0$ by the Moore--Penrose identity;
so $\chi_{B'}(J)=\chi_{B'}(J_0)=\omega$.

\emph{Step 6 (the fibres).} $W_{\tau}=W_{\tau}P_{\iota}$ by the same
identity, so for any $J$ in the domain
$\chi_{B'}(J)=(W_{\tau}\otimes\idon{})\big[(P_{\iota}\otimes\idon{})J
(P_{\iota}\otimes\idon{})\big](W_{\tau}^{\dagger}\otimes\idon{})$, and
$W_{\tau}$ restricted to $\operatorname{ran}P_{\iota}
=(\ker W_{\tau})^{\perp}$ is injective; hence $\chi_{B'}(J)$ determines,
and is determined by, the compression
$(P_{\iota}\otimes\idon{})J(P_{\iota}\otimes\idon{})$. That is the stated
fibre description, and it gives injectivity when $P_{\iota}=\idon{A_0}$,
i.e.\ when the frame is tight; the displayed inverse is then $J_0$.
Conversely, let the frame not be tight, so that
$\idon{A_0}-P_{\iota}\ne0$, and let $\dim\Hs{B'}\ge2$, so that $B'$ carries
two states $\varphi\ne\varphi'$. With any state $\varphi_0$ on $B'$, the
operators $J\coloneqq P_{\iota}\otimes\varphi_0
+(\idon{A_0}-P_{\iota})\otimes\varphi$ and
$J'\coloneqq P_{\iota}\otimes\varphi_0
+(\idon{A_0}-P_{\iota})\otimes\varphi'$ are positive semidefinite with
$\tr_{B'}J=\tr_{B'}J'=\idon{A_0}$, have the same compression
$P_{\iota}\otimes\varphi_0$, hence the same image, and differ by
$(\idon{A_0}-P_{\iota})\otimes(\varphi-\varphi')\ne0$; so $\chi_{B'}$ is
not injective. If $\dim\Hs{B'}=1$, then $\tr_{B'}J=J$ and the domain is
the single point $\idon{A_0}$.
\end{proof}

% claim meat-core/lem:choi-factor
\claimsummary{meat-core/lem:choi-factor}{Any factor W of Alice's marginal parametrises the fixed-marginal states per Alice sector, and the objective reads her side only through her measurement sandwiched by W. On the full frame these are the declared operators of an ensemble-form declaration, so no rank is computed.}
\claimbookmark{meat-core/lem:choi-factor}{Lemma}{lem:choi-factor}{Fixed-marginal parametrisation with any factor}
\claimleadin{%
Of $W_{\tau}$, the parametrisation of Lemma~\ref{lem:choi-marginal} reads
only that it is a factor of $\tau_A$. Stated for any factor, and per Alice
sector, it shows that the objective reads Alice's side only through the
operators $W^{\dagger}M_xW$, which on the \termref{term-full-frame}{full frame} are the operators an
ensemble-form declaration declares.
}
\begin{lemma}[Fixed-marginal parametrisation with any factor]\label{lem:choi-factor}
Let $\Decl$ satisfy Definition~\ref{def:admissible-declaration} in trusted-source mode, in the
general format of Definition~\ref{def:declaration}.
\begin{enumerate}[label=(\roman*),leftmargin=2.2em]
\item \emph{(any factor serves)} Let $\tau\in\Pos{A'}$ on a
finite-dimensional system $A'$, of any trace, and let $W\in\Linn{A_0}{A'}$
satisfy $WW^{\dagger}=\tau$. Let $P_W$ be the projector onto
$(\ker W)^{\perp}$. For every finite-dimensional $B'$ the map
$\chi^{W}_{B'}(J)\coloneqq(W\otimes\idon{B'})J(W^{\dagger}\otimes\idon{B'})$
maps $\{J\in\Pos{A_0B'}:\tr_{B'}J=\idon{A_0}\}$ onto
$\{\omega\in\Pos{A'B'}:\tr_{B'}\omega=\tau\}$; it is injective when
$\ker W=0$, and, if $\dim\Hs{B'}\ge2$, only then; and
$\chi^{W}_{B'}(J)=\chi^{W}_{B'}(J')$ exactly when
$(P_W\otimes\idon{})J(P_W\otimes\idon{})=(P_W\otimes\idon{})J'(P_W\otimes\idon{})$.
The operator $W_{\tau}$ of \eqref{eq:Wtau} is such a factor of $\tau_A$ for
every admissible frame.
\item \emph{(per Alice sector)} For every $t\in\mathbb{T}$ let $W_t\in
\Linn{A_{0,t}}{A_t}$ satisfy $W_tW_t^{\dagger}=\tau^{(t)}_A$, and put
$A_0\coloneqq\bigoplus_tA_{0,t}$, $W\coloneqq\bigoplus_tW_t$. Then
$WW^{\dagger}=\tau_A$, and the assignment
$(J_{t,s})\mapsto\big(\omega^{(t,s)}\big)$,
$\omega^{(t,s)}\coloneqq(W_t\otimes\idon{\Qsys_s})J_{t,s}
(W_t^{\dagger}\otimes\idon{\Qsys_s})$, maps
\begin{equation}\label{eq:choi-coupling-t}
  \Big\{(J_{t,s})_{(t,s)}:\ J_{t,s}\in\Pos{A_{0,t}\Qsys_s},\quad
  \sum_{s\in\SecSet}\tr_{\Qsys_s}J_{t,s}=\idon{A_{0,t}}
  \ \ \forall t\in\mathbb{T}\Big\}
\end{equation}
onto the set of lists over the \exttermref{../protocol-compiler/contract.pdf}{term-joint-sector}{joint sectors} with positive semidefinite
blocks and $\sum_s\tr_{\Qsys_s}\omega^{(t,s)}=\tau^{(t)}_A$ for every $t$,
bijectively when $\ker W_t=0$ for every $t$, and, if
$\dim\Hs{\Qsys}\ge2$, only then. The set
\eqref{eq:choi-coupling-t} is a Cartesian product over $t$. If every
$\iota_t$ satisfies \eqref{eq:iota} for $\tau^{(t)}_A$ on $A_t$, then
$\iota\coloneqq\bigoplus_t\iota_t$ satisfies \eqref{eq:iota} for $\tau_A$ and
$W_{\tau}=\bigoplus_t(\tau^{(t)}_A)^{1/2}\iota_t$.
\item \emph{(the objective reads Alice's side through $N_x$)} For a factor
$W=\bigoplus_tW_t$ as in (ii) put
\begin{equation}\label{eq:NfromW}
  N^{W}_x\coloneqq W^{\dagger}M_xW=\bigoplus_tW_t^{\dagger}M^{(t)}_xW_t
  \in\Pos{A_0},
\end{equation}
let $K^{N}_{\kappa}$ be \eqref{eq:kraus} with $\sqrt{N^{W}_x}$ in place of
$\sqrt{M_x}$, an operator
$\Hs{A_0\Qsys}\to\Hs{R\,XY\,C_AC_B\,A_0\Qsys}$, and let $\cG^{N}_{\hat c}$ be
\eqref{eq:GTheta} with $K^{N}_{\kappa}$. Then for every $J\in\Pos{A_0\Qsys}$,
every $(x,v)$ and every $\hat c$, with
$\omega\coloneqq\chi^{W}_{\Qsys}(J)$,
\begin{equation}\label{eq:choi-N}
  \tr[(M_x\otimes F_v)\,\omega]=\tr[(N^{W}_x\otimes F_v)\,J],\qquad
  \tr\cG_{\hat c}(\omega)=\tr\cG^{N}_{\hat c}(J),\qquad
  T\big(\cG_{\hat c}(\omega)\big)=T\big(\cG^{N}_{\hat c}(J)\big),
\end{equation}
$T$ the functional \eqref{eq:Tcone}; hence $P_{\hat c}$, $K_{\hat c}$ and
$O_{\alpha}$ at $\omega$ are functions of $J$, $\{N^{W}_x\}$, $\{F_v\}$ and
the classical maps alone, and the same holds sector by sector for the
joint-sector quantities of Definition~\ref{def:joint-sectors}.
\item \emph{(\exttermref{../protocol-compiler/contract.pdf}{term-ensemble-form}{ensemble form})} On the full frame, $A_{0,t}=A_t$ and
$W_t=(\tau^{(t)}_A)^{1/2}$, one has
$N^{W}_x=\tau_A^{1/2}M_x\tau_A^{1/2}=N_x$, the operators of
Definition~\ref{def:declaration}\,(DX4): for an ensemble-form declaration
they are the \emph{declared} numbers. So the program of an ensemble-form
declaration is the program \eqref{eq:choi-coupling-t}, (iii) in the declared
$\{N_x\}$ and $\{F_v\}$, with every constraint row in the conjugated form of
\eqref{eq:choi-feas}; it is defined, and equals $\hlow$ of the denoted
declaration, at every rank of every $\tau^{(t)}_A$, and no rank is computed.
\end{enumerate}
\end{lemma}
\claimmeta{\metalabel{Uses}\mref{Def.}{def:declaration}, \mref{Def.}{def:admissible-declaration}, \mref{Lem.}{lem:ensemble-form}, \mref{Def.}{def:maps}, \mref{Def.}{def:branch}, \mref{Lem.}{lem:branch-maps}, \mref{Thm.}{thm:branch-formula}, \mref{Lem.}{lem:cone-concavity}, \mref{Def.}{def:joint-sectors}, \mref{Def.}{def:feasible}, \mref{Def.}{def:support-frame}, \mref{Lem.}{lem:choi-marginal}
\metasep \metalabel{Used by}\mref{Thm.}{thm:choi}, \mref{Prop.}{prop:relaxed-set-properties}, \mref{Cor.}{cor:fo-dropped}
\metasep newly derived
\metasep implemented (effective: accepted)
\metasep \metalabel{Code}\coderef{https://github.com/fraunhofersg/certiqs-key/blob/v0.1.0/src/certiqs_key/meat_core/objective/joint.py}{certiqs_key.meat_core.objective.joint:joint_choi_parametrization}}

\begin{proof}
(i) Three facts about $W$. (a) $WW^{\dagger}=\tau$ is the hypothesis. (b)
$\operatorname{ran}W=\operatorname{ran}(WW^{\dagger})=\operatorname{supp}\tau$:
for any operator, $\operatorname{ran}W\supseteq\operatorname{ran}WW^{\dagger}$
and $\ker W^{\dagger}=\ker WW^{\dagger}$ (if $WW^{\dagger}\xi=0$ then
$\|W^{\dagger}\xi\|^2=\bra{\xi}WW^{\dagger}\ket{\xi}=0$), so the two ranges,
being the orthogonal complements of equal kernels, agree. (c) Therefore the
Moore--Penrose inverse satisfies
$WW^{+}=\Pi_{\operatorname{supp}\tau}$ and $W^{+}W=P_W$, and
$W=WW^{+}W=WP_W$. The proof of Lemma~\ref{lem:choi-marginal} uses of
$W_{\tau}$ exactly \eqref{eq:WWtau}, which is (a), and \eqref{eq:Wplus} with
the identity $W_{\tau}(\idon{A_0}-P_{\iota})=0$, which are (c); it uses
neither $\tr\tau_A=1$ nor the isometry $\iota$ otherwise. We repeat it for
$W$. \emph{Forward:} $\chi^{W}(J)\succeq0$ and
$\tr_{B'}\chi^{W}(J)=W(\tr_{B'}J)W^{\dagger}=WW^{\dagger}=\tau$.
\emph{Onto:} let $\omega\succeq0$ with $\tr_{B'}\omega=\tau$ and
$\Pi^{\perp}\coloneqq\idon{}-\Pi_{\operatorname{supp}\tau}$; then
$\tr[(\Pi^{\perp}\otimes\idon{})\omega]=\tr[\Pi^{\perp}\tau]=0$, and
positivity gives $(\Pi^{\perp}\otimes\idon{})\omega=0$. Put
$J_0\coloneqq(W^{+}\otimes\idon{})\omega(W^{+\dagger}\otimes\idon{})\succeq0$.
Then $\tr_{B'}J_0=W^{+}WW^{\dagger}W^{+\dagger}=P_WP_W^{\dagger}=P_W$ and
$\chi^{W}(J_0)=(WW^{+}\otimes\idon{})\omega(WW^{+}\otimes\idon{})^{\dagger}
=\omega$. With any state $\varphi$ on $B'$,
$J\coloneqq J_0+(\idon{A_0}-P_W)\otimes\varphi\succeq0$ has
$\tr_{B'}J=\idon{A_0}$ and $\chi^{W}(J)=\omega$, because
$W(\idon{}-P_W)=0$. \emph{Fibres:} $\chi^{W}(J)=(W\otimes\idon{})
[(P_W\otimes\idon{})J(P_W\otimes\idon{})](W^{\dagger}\otimes\idon{})$ and $W$
is injective on $\operatorname{ran}P_W$, so $\chi^{W}(J)$ determines and is
determined by the compression; injectivity holds when $P_W=\idon{A_0}$. If
$P_W\ne\idon{A_0}$ and $\dim\Hs{B'}\ge2$, two states $\varphi\ne\varphi'$
on $B'$ and any state $\varphi_0$ give the two operators
$P_W\otimes\varphi_0+(\idon{A_0}-P_W)\otimes\varphi$ and
$P_W\otimes\varphi_0+(\idon{A_0}-P_W)\otimes\varphi'$ of the domain, which
differ and have the same compression, hence the same image, as in the proof
of Lemma~\ref{lem:choi-marginal}. Finally $W_{\tau}W_{\tau}^{\dagger}=\tau_A$ is
\eqref{eq:WWtau}.

(ii) $WW^{\dagger}=\bigoplus_tW_tW_t^{\dagger}=\tau_A$. Fix $t$ and apply (i)
with $A'=A_t$, $\tau=\tau^{(t)}_A$, $B'=\Qsys$. The map $\chi^{W_t}_{\Qsys}$
and the section $\omega\mapsto J_0$ act on the first tensor factor alone, so
they commute with $X\mapsto(\idon{}\otimes\Pi_s)X(\idon{}\otimes\Pi_{s'})$
and preserve block diagonality along $\SecSet$; choosing the state $\varphi$
of (i) supported on one $\Hs{\Qsys_{s_0}}$ keeps the completed $J$ block
diagonal as well. For operators block diagonal along $\SecSet$,
$\tr_{\Qsys}J_t=\sum_s\tr_{\Qsys_s}J_{t,s}$, and likewise for $\omega$. So
for each $t$ the map is onto, and injective when $\ker W_t=0$; if
$\ker W_t\ne0$ and $\dim\Hs{\Qsys}\ge2$, the two operators of (i) for
$W_t$, built with a sector-block-diagonal $\varphi_0$ and two different
sector-block-diagonal states $\varphi\ne\varphi'$ on $\Qsys$ (states on two
sectors, or two states on one sector of dimension at least two), are two
different block-diagonal preimages of one point. The domain and the
codomain are products over $t$ of the per-sector sets, and the map is the
product of the per-sector maps, which gives the statement.
For the frames: $\iota^{\dagger}\iota=\bigoplus_t\iota_t^{\dagger}\iota_t
=\idon{A_0}$ and $\iota\iota^{\dagger}\tau_A=\bigoplus_t\iota_t
\iota_t^{\dagger}\tau^{(t)}_A=\tau_A$, and
$\tau_A^{1/2}=\bigoplus_t(\tau^{(t)}_A)^{1/2}$ by blockwise functional
calculus.

(iii) The first identity is cyclicity of the trace:
$\tr[(M_x\otimes F_v)(W\otimes\idon{})J(W^{\dagger}\otimes\idon{})]
=\tr[(W^{\dagger}M_xW\otimes F_v)J]$. The second follows from it and the
trace formula in the proof of Lemma~\ref{lem:branch-maps}\,(i), which reads of the
operators in the slot of $\sqrt{M_x}$ only that their squares' traces against
the input are $\tr[(M_x\otimes F_v)\,\cdot\,]$, resp.\
$\tr[(N^{W}_x\otimes F_v)\,\cdot\,]$. For the third, fix $x$. The operator
$Y_x\coloneqq\sqrt{M_x}\,W\in\Linn{A_0}{A}$ satisfies
$Y_x^{\dagger}Y_x=W^{\dagger}M_xW=N^{W}_x$, so by the polar decomposition
$Y_x=V_x\sqrt{N^{W}_x}$ with $V_x\in\Linn{A_0}{A}$ a partial isometry whose
initial space is $\mathcal{I}_x\coloneqq\overline{\operatorname{ran}}\,N^{W}_x$:
$V_x^{\dagger}V_x$ is the projector onto $\mathcal{I}_x$. Put
$\widetilde V\coloneqq\sum_{x\in\Xset}\dyad{x}_X\otimes V_x$, acting as the
identity on $R$, $Y$, $C_AC_B$ and $\Qsys$; it is a partial isometry with
initial projector
$\widetilde E\coloneqq\sum_x\dyad{x}_X\otimes V_x^{\dagger}V_x$. Every term
of $K_{\kappa}(W\otimes\idon{\Qsys})$ is
$\ket{r(x,v)}\ket{x,v}\ket{\kappa}\otimes(V_x\sqrt{N^{W}_x}\otimes\sqrt{F_v})$,
so
\begin{equation*}
  K_{\kappa}(W\otimes\idon{\Qsys})=\widetilde V\,K^{N}_{\kappa},\qquad
  \widetilde E\,K^{N}_{\kappa}=K^{N}_{\kappa},
\end{equation*}
the second because the term of $K^{N}_{\kappa}$ at $x$ has its $A_0$ range in
$\operatorname{ran}\sqrt{N^{W}_x}\subseteq\mathcal{I}_x$. Hence
$\cG_{\hat c}(\omega)=\widetilde V\,X\,\widetilde V^{\dagger}$ with
$X\coloneqq\cG^{N}_{\hat c}(J)\succeq0$ and $\widetilde EX\widetilde E=X$.
For such a pair, $\widetilde V$ is isometric on the support of $X$, so
$(\widetilde VX\widetilde V^{\dagger})^{1/\alpha}
=\widetilde VX^{1/\alpha}\widetilde V^{\dagger}$; $\widetilde V$ and
$\widetilde E$ commute with every $Z_j=\dyad{j}_R$, so
$\cZ(\widetilde VX^{1/\alpha}\widetilde V^{\dagger})
=\widetilde V\cZ(X^{1/\alpha})\widetilde V^{\dagger}$ and
$\widetilde E\cZ(X^{1/\alpha})\widetilde E=\cZ(X^{1/\alpha})$; therefore
$(\widetilde V\cZ(X^{1/\alpha})\widetilde V^{\dagger})^{\alpha}
=\widetilde V\cZ(X^{1/\alpha})^{\alpha}\widetilde V^{\dagger}$ and its trace
is $\tr[\cZ(X^{1/\alpha})^{\alpha}\widetilde E]=\tr[\cZ(X^{1/\alpha})^{\alpha}]$.
That is $T(\cG_{\hat c}(\omega))=T(X)$. The consequences: $P_{\hat c}$ is a
sum of the first quantities, $K_{\hat c}=T+P-\tr\cG$ by \eqref{eq:fK}, and
$O_{\alpha}$ is a function of these by \eqref{eq:objective}. The argument
used no property of $J$ beyond $J\succeq0$; applied to a $J$ supported on
one joint sector, where $W$, $M_x$ and $N^{W}_x$ act through their $t$-blocks
and $F_v$ through its $s$-block, it gives the sector-by-sector statement.

(iv) $W_t=(\tau^{(t)}_A)^{1/2}$ is Hermitian with $W_tW_t^{\dagger}
=\tau^{(t)}_A$, a factor in the sense of (ii) and the operator
\eqref{eq:Wtau} of the full frame. $N^{W}_x=\tau_A^{1/2}M_x\tau_A^{1/2}$ is
the definition \eqref{eq:NfromW}; for the declaration an ensemble form
denotes it is the declared $N_x$ by Lemma~\ref{lem:ensemble-form}\,(ii). By
(ii) the assignment is onto the marginal-constrained lists at every rank, a
constraint row evaluated at $\omega=\chi^{W}(J)$ is the conjugated row by
substitution, exactly as in the proof of Theorem~\ref{thm:choi}, and by (iii)
the objective is evaluated from the declared operators; so the infimum over
$\acSet$ and the feasible lists of that program is the infimum of $O_{\alpha}$ over
$\acSet\times\Feas$ for the denoted declaration, which is $\hlow$. Nothing in
(i)--(iii) reads a rank, an inverse or the operators
\eqref{eq:ensemblepovm}.
\end{proof}

% claim meat-core/thm:choi
\claimsummary{meat-core/thm:choi}{In trusted-source mode the program is exactly a jointly convex minimisation over sector Choi blocks, with one coupling constraint per Alice sector and affine or conic rows. Every support frame gives the value h\_low, since only surjectivity of the parametrisation enters it.}
\claimbookmark{meat-core/thm:choi}{Theorem}{thm:choi}{The Choi program}
\claimleadin{%
Applied block by block along the sectors, the parametrisation of
Lemma~\ref{lem:choi-marginal} rewrites the program of
Definition~\ref{def:feasible} in Choi blocks, the variables a solver works
in, without changing its value on any support frame.
}
\begin{theorem}[The Choi program]\label{thm:choi}
Let the declaration be in \exttermref{../transmitter/contract.pdf}{term-trusted-source-mode}{trusted-source mode}, fix a
\termref{term-support-frame}{support frame} (Definition~\ref{def:support-frame}), apply
Lemma~\ref{lem:choi-marginal} with $B'=\Qsys$, and restrict to
sector-block-diagonal operators. Then the assignment
$J=\bigoplus_{s}J_s\mapsto\omega=\chi_{\Qsys}(J)$, with
$\omega_s=\chi_{\Qsys_s}(J_s)\coloneqq
(W_{\tau}\otimes\idon{\Qsys_s})J_s(W_{\tau}^{\dagger}\otimes\idon{\Qsys_s})$,
maps
\begin{equation}\label{eq:choi-feas}
  \Feas^{J}\coloneqq\Big\{\!(J_s)_{s\in\SecSet}:\ J_s\in\Pos{A_0\Qsys_s},\
  \sum_{s\in\SecSet}\tr_{\Qsys_s}J_s=\idon{A_0},\
  W_{\tau}\Big(\textstyle\sum_s
  \tr_{\Qsys_s}[(\idon{A_0}\otimes\hat O^{(s)}_k)J_s]\Big)
  W_{\tau}^{\dagger}\bowtie_k\Theta_k\ \forall k\!\Big\}
\end{equation}
\emph{onto} $\Feas$ of \eqref{eq:feas}, bijectively when the support frame
is tight (Lemma~\ref{lem:choi-marginal}), and, if $\Feas\ne\varnothing$ and
$\dim\Hs{\Qsys}\ge2$, only then. Writing
$\Phimap_{J}[J]\coloneqq\Phimap[\chi_{\Qsys}(J)]$,
$\cG^{J}_{\hat c,s}\coloneqq\cG_{\hat c,s}\circ\chi_{\Qsys_s}$,
$f^{J}_{\hat c,s}\coloneqq f_{\hat c,s}\circ\chi_{\Qsys_s}$,
$P^{J}_{\hat c}(J)\coloneqq\Phimap_J[J](\hat c)$ and
$K^{J}_{\hat c}(J)\coloneqq P^{J}_{\hat c}(J)
+\sum_{s\in\SecSet_K(\hat c)}f^{J}_{\hat c,s}(J_s)$, the value
\eqref{eq:hlow-def} equals
\begin{equation}\label{eq:choi-program}
\boxed{\ \
\begin{aligned}
  \hlow=\inf_{\mathbf q,\ (J_s)_s}\quad
  &\frac{\alpha}{\alpha-1}\sum_{\hat c\notin\Chkey}
  q(\hat c)\log_2\frac{q(\hat c)}{P^{J}_{\hat c}(J)}
  +\sum_{\hat c\in\Chkey}\left[
  q(\hat c)\log_2\frac{q(\hat c)}{P^{J}_{\hat c}(J)}
  +\frac{q(\hat c)}{\alpha-1}
  \log_2\frac{q(\hat c)}{K^{J}_{\hat c}(J)}\right]\\
  \text{s.t.}\quad
  &\mathbf q\in\acSet,\qquad (J_s)_s\in\Feas^{J},
\end{aligned}\ \ }
\end{equation}
an \emph{exact} reformulation (not a relaxation), jointly convex in
$(\mathbf q,(J_s)_s)$, with all constraints affine or conic.

\emph{In the general format} of Definition~\ref{def:declaration}, let the
frame be a direct sum $\iota=\bigoplus_{t\in\mathbb{T}}\iota_t$ with every
$\iota_t\in\Linn{A_{0,t}}{A_t}$ satisfying \eqref{eq:iota} for
$\tau^{(t)}_A$ --- either construction of Definition~\ref{def:support-frame}
applied sector by sector --- and put $W_t\coloneqq(\tau^{(t)}_A)^{1/2}\iota_t$. Let $\Feas^{J}$
be the set of lists $(J_{t,s})$ over the \exttermref{../protocol-compiler/contract.pdf}{term-joint-sector}{joint sectors} with
$J_{t,s}\in\Pos{A_{0,t}\Qsys_s}$, the coupling constraints
$\sum_s\tr_{\Qsys_s}J_{t,s}=\idon{A_{0,t}}$ for every $t$, and for every
entry the row of \eqref{eq:feas-general} evaluated at
$\omega^{(t,s)}=(W_t\otimes\idon{})J_{t,s}(W_t^{\dagger}\otimes\idon{})$: for
an entry read through its \exttermref{../protocol-compiler/contract.pdf}{term-lifted-operator}{lifted operator},
$W_t\big(\sum_s\tr_{\Qsys_s}[(\idon{A_{0,t}}\otimes\hat O^{\uparrow(s)}_k)
J_{t,s}]\big)W_t^{\dagger}\bowtie_k\Theta^{(t)}_k$ for every $t$, and for a
\exttermref{../protocol-compiler/contract.pdf}{term-party-local}{party-local} entry the last line of \eqref{eq:feas-general} at $J$,
conjugated by $W_t\otimes\idon{}$. Then $(J_{t,s})\mapsto(\omega^{(t,s)})$
maps $\Feas^{J}$ onto $\Feas$ of \eqref{eq:feas-general}, bijectively when
every sector's frame is tight, and, if $\Feas\ne\varnothing$ and
$\dim\Hs{\Qsys}\ge2$, only then, and \eqref{eq:choi-program} holds with
$K^{J}_{\hat c}(J)=P^{J}_{\hat c}(J)+\sum_{(t,s)\in\mathbb{J}_K(\hat c)}
f_{\hat c,(t,s)}(\omega^{(t,s)})$. Every constraint of $\Feas^{J}$ involves
the blocks of one Alice sector only, so $\Feas^{J}=\prod_t\Feas^{J}_t$, and a
linear function of $(J_{t,s})$ is minimised over $\Feas^{J}$ by minimising,
independently for each $t$, its part on the blocks of $t$ over
$\Feas^{J}_t$.
\end{theorem}
\claimmeta{\metalabel{Uses}\mref{Def.}{def:declaration}, \mref{Prop.}{prop:exact-cancellation}, \mref{Def.}{def:keyactive}, \mref{Def.}{def:joint-sectors}, \mref{Lem.}{lem:joint-sectors}, \mref{Def.}{def:feasible}, \mref{Thm.}{thm:convexity}, \mref{Lem.}{lem:admissible-alice-subspace}, \mref{Def.}{def:support-frame}, \mref{Lem.}{lem:choi-marginal}, \mref{Lem.}{lem:choi-factor}, \extref{../foundations/contract.pdf}{foundations/def:sector-map}
\metasep \metalabel{Used by}\mref{Prop.}{prop:frame-free-rows}, \mref{Rem.}{rem:tightness-by-construction}, \mref{Rem.}{rem:ordering}, \mref{Rem.}{rem:basis-invariance}, \mref{Def.}{def:untrusted-choi}, \mref{Lem.}{lem:estimate}, \mref{Cor.}{cor:full-gradient}, \mref{Lem.}{lem:kraus-freedom}, \mref{Rem.}{rem:oracle}, \mref{Lem.}{lem:objective-homogeneous}, \mref{Rem.}{rem:cgate}, \mref{Req.}{req:choi-convention}, \mref{Def.}{def:relaxed-feasible-set}, \mref{Prop.}{prop:relaxed-set-properties}, \mref{Def.}{def:first-order-setting}, \mref{Lem.}{lem:fo-khyp}, \mref{Lem.}{lem:fo-face}, \mref{Thm.}{thm:firstorder}, \mref{Cor.}{cor:fo-dropped}, \mref{Rem.}{rem:fo-general-format}, \mref{Lem.}{lem:feasible-side-dual}, \mref{Lem.}{lem:certified-bracket}, \mref{Lem.}{lem:certified-bracket-enclosed}, \mref{Rem.}{rem:fo-iterate-checks}, \mref{Req.}{req:certified-evaluation}
\metasep newly derived
\metasep implemented (effective: accepted)
\metasep \metalabel{Code}\coderef{https://github.com/fraunhofersg/certiqs-key/blob/v0.1.0/src/certiqs_key/meat_core/objective/choi.py}{certiqs_key.meat_core.objective.choi:choi_parametrization}
\metasep \metalabel{Oracles}\coderef{https://github.com/fraunhofersg/certiqs-key/blob/v0.1.0/tests/oracles/registry.py}{registry.py}}

\begin{proof}
\emph{Idea.} The map $\chi_{\Qsys}$ acts on the $A$ factor alone, so it
respects the sector blocks and carries every constraint row over by
substitution. Surjectivity is that of Lemma~\ref{lem:choi-marginal},
completed on one sector by a block that $W_{\tau}$ annihilates; the value
follows by substituting into the objective, and injectivity is not needed.

\emph{The map.} Lemma~\ref{lem:choi-marginal} restricts to
sector-block-diagonal operators because $\chi_{\Qsys}$ acts only on the
$A$ factor and therefore commutes with the sector projectors
$\idon{}\otimes\Pi_s$; the global trace-preservation condition
$\tr_{\Qsys}J=\idon{A_0}$ is the single coupling constraint
$\sum_s\tr_{\Qsys_s}J_s=\idon{A_0}$, which by the lemma \emph{implies}
$\omega_A=\tau_A$ and is implied by it up to the fibre the lemma
describes. The $\CList$ rows transfer by substitution, using
$\tr_{\Qsys}[(\idon{A}\otimes\hat O_k)\chi_{\Qsys}(J)]
=W_{\tau}(\tr_{\Qsys}[(\idon{A_0}\otimes\hat O_k)J])W_{\tau}^{\dagger}$
--- which holds because
$(\idon{A}\otimes\hat O_k)(W_{\tau}\otimes\idon{})
=(W_{\tau}\otimes\idon{})(\idon{A_0}\otimes\hat O_k)$ --- and
block-diagonality of $\hat O_k$: the $k$-th row of \eqref{eq:choi-feas}
\emph{is} the $k$-th row of \eqref{eq:feas-modes} evaluated at
$\omega=\chi_{\Qsys}(J)$, so it holds for $J$ exactly when it holds for
its image. Normalization $\tr\omega=1$ follows from the coupling
constraint and \eqref{eq:iota}, since
$\tr\omega=\tr[W_{\tau}^{\dagger}W_{\tau}\cdot\tr_{\Qsys}J]
=\tr[W_{\tau}^{\dagger}W_{\tau}]
=\tr[\iota^{\dagger}\tau_A\iota]=\tr[\tau_A\,\iota\iota^{\dagger}]
=\tr[\Pi_{\iota}\tau_A]=\tr\tau_A=1$. Together with positivity and
block-diagonality, which $\chi_{\Qsys_s}$ preserves, this gives
$\chi_{\Qsys}(\Feas^{J})\subseteq\Feas$.

\emph{Onto.} Fix $\omega\in\Feas$. Put
$J_{0,s}\coloneqq(W_{\tau}^{+}\otimes\idon{})\,\omega_s\,
(W_{\tau}^{+\dagger}\otimes\idon{})\succeq0$ for every $s$; then
$\chi_{\Qsys_s}(J_{0,s})=\omega_s$ exactly as in
Lemma~\ref{lem:choi-marginal} (each $\omega_s$ is supported on
$\operatorname{supp}\tau_A\otimes\Hs{\Qsys_s}$, which is
Lemma~\ref{lem:admissible-alice-subspace} with $\Sadm$ of \eqref{eq:SA}), and
$\sum_s\tr_{\Qsys_s}J_{0,s}
=W_{\tau}^{+}\big(\sum_s\tr_{\Qsys_s}\omega_s\big)W_{\tau}^{+\dagger}
=W_{\tau}^{+}\tau_AW_{\tau}^{+\dagger}=P_{\iota}$. Choose one
$s_0\in\SecSet$ ($\SecSet$ is non-empty, being the index set of a \exttermref{../foundations/contract.pdf}{term-sector-map}{sector map},
\xref{foundations/def:sector-map}) and any
state $\varphi$ on $\Hs{\Qsys_{s_0}}$, and set
$J_s\coloneqq J_{0,s}+\delta_{s,s_0}\,(\idon{A_0}-P_{\iota})
\otimes\varphi$. Each $J_s\succeq0$;
$\sum_s\tr_{\Qsys_s}J_s=P_{\iota}+(\idon{A_0}-P_{\iota})=\idon{A_0}$; and
$W_{\tau}(\idon{A_0}-P_{\iota})=0$ kills the added block under both
$\chi_{\Qsys_{s_0}}$ and the conjugated constraint rows, so
$\chi_{\Qsys}(J)=\omega$ and every row of \eqref{eq:choi-feas} reads the
corresponding row of \eqref{eq:feas-modes} at $\omega$, which holds.
Hence $(J_s)_s\in\Feas^{J}$ and $\chi_{\Qsys}(\Feas^{J})=\Feas$.
Injectivity, when the frame is tight, is the lemma's. If the frame is not
tight, $\idon{A_0}-P_{\iota}\ne0$; if moreover $\Feas\ne\varnothing$ and
$\dim\Hs{\Qsys}\ge2$, there are two different choices of $(s_0,\varphi)$ ---
two sectors, or two states on one sector of dimension at least two --- and
for a fixed $\omega\in\Feas$ they give two different lists $(J_s)_s$ in
$\Feas^{J}$, differing in the added block, with the same image $\omega$.

\emph{The value.} $O^{J}_{\alpha}=O_{\alpha}\circ(\idmap\times
\chi_{\Qsys})$ by construction --- the objective transfers by
substitution of $\omega_s=\chi_{\Qsys_s}(J_s)$ into
\eqref{eq:objective}, using Definition~\ref{def:keyactive} for the
restriction of the $K$-sum to key-active sectors (an identity, by
Proposition~\ref{prop:exact-cancellation}) --- so
$\inf_{\acSet\times\Feas^{J}}O^{J}_{\alpha}
=\inf_{\mathbf q\in\acSet}\inf_{\omega\in\chi_{\Qsys}(\Feas^{J})}
O_{\alpha}(\mathbf q,\omega)
=\inf_{\mathbf q\in\acSet}\inf_{\omega\in\Feas}O_{\alpha}(\mathbf q,\omega)
=\hlow$, which is \eqref{eq:choi-program}. Only surjectivity enters this
chain and injectivity does not, so a non-tight support frame changes no
value.

\emph{Convexity.} Each $\chi_{\Qsys_s}$ is linear, so
Theorem~\ref{thm:convexity} applies composed with a linear map; the
constraint set \eqref{eq:choi-feas} is an intersection of the PSD cones
with affine conditions (for $\bowtie_k\in\{\preceq,\succeq\}$, a
conjugated affine-semidefinite condition), hence convex and compact ---
compact because $J_s\succeq0$ with
$\sum_s\tr J_s=\tr\idon{A_0}=d_{A_0}$.

\emph{The general format.} By Lemma~\ref{lem:choi-factor}\,(ii), applied with
the factors $W_t$ (factors of $\tau^{(t)}_A$ by \eqref{eq:WWtau} for each
sector), the assignment maps the lists satisfying the coupling constraints
onto the lists over the joint sectors with positive semidefinite blocks and
marginal $\tau_A$, injectively when every $\ker W_t=0$, which is
tightness of every sector's frame. Where some $\ker W_t\ne0$ and
$\dim\Hs{\Qsys}\ge2$, the onto step of that lemma's proof, run for $W_t$
with two different sector-block-diagonal states $\varphi\ne\varphi'$ (states
on two sectors, or two states on one sector of dimension at least two), gives
two different preimages of any point, differing by
$(\idon{A_{0,t}}-P_{W_t})\otimes(\varphi-\varphi')$; having the same
image, they satisfy the same rows, which $\Feas^{J}$ reads at the image
point, so at a point of a non-empty $\Feas$ both lie in $\Feas^{J}$.
Normalization follows from the marginal
condition and $\tr\tau_A=1$. The rows of $\Feas^{J}$ are by definition the
rows of \eqref{eq:feas-general} at the image point, so a list lies in
$\Feas^{J}$ exactly when it satisfies the coupling constraints and its image
lies in $\Feas$; for the lifted-operator rows the displayed conjugated form
is the substitution made in the paragraph \emph{The map}, sector by sector. Hence the image of
$\Feas^{J}$ is $\Feas$. The value: substitute into \eqref{eq:objective} and
use Lemma~\ref{lem:joint-sectors}\,(ii), (iii) for $K_{\hat c}$ and the
restriction of its sum to $\mathbb{J}_K(\hat c)$; surjectivity gives the
equality of the infima as before. Convexity and compactness are as before,
each coupling constraint bounding $\sum_s\tr J_{t,s}=d_{A_{0,t}}$. The
product structure: the coupling constraint of $t$, the lifted-operator rows
of $t$ and the party-local rows at $(t,s_{f^c})$ each involve only blocks
$J_{t,s}$ with that $t$, and a linear function on a Cartesian product is a
sum of linear functions of the factors, minimised factor by factor.
\end{proof}

% claim meat-core/prop:frame-free-rows
\claimsummary{meat-core/prop:frame-free-rows}{On a tight support frame a homogeneous row of the Choi program may drop its conjugation by W\_\ensuremath{\tau} without changing the feasible set; on any frame the frame-free row only restricts the set, so the infimum over it is at least h\_low. The same holds per Alice sector and for party-local rows.}
\claimbookmark{meat-core/prop:frame-free-rows}{Proposition}{prop:frame-free-rows}{Frame-free constraint rows on a tight frame}
\claimleadin{%
In the Choi program of Theorem~\ref{thm:choi} every constraint row is
imposed conjugated by $W_{\tau}$. For a homogeneous entry the right-hand
side $c_k\tau_A$ is itself $c_kW_{\tau}W_{\tau}^{\dagger}$
\eqref{eq:WWtau}, so the conjugation can be taken off both sides. The
proposition says when that changes the set $\Feas^{J}$ and in which
direction.
}
\begin{proposition}[Frame-free constraint rows on a tight frame]\label{prop:frame-free-rows}
In (i)--(iii) let the declaration be in trusted-source mode, let $W_{\tau}$ be
\eqref{eq:Wtau} for a support frame (Definition~\ref{def:support-frame}),
and for an entry $k$ of $\CList$ and $(J_s)_{s\in\SecSet}$ with
$J_s\in\Pos{A_0\Qsys_s}$ write
$Y_k\coloneqq\sum_s\tr_{\Qsys_s}[(\idon{A_0}\otimes\hat O^{(s)}_k)J_s]$.
For a \exttermref{../protocol-compiler/contract.pdf}{term-homogeneous}{homogeneous} entry, whose right-hand side is
$\Theta_k=c_k\tau_A$ (Definition~\ref{def:homog}\,(a)), call the $k$-th row
of \eqref{eq:choi-feas}, $W_{\tau}Y_kW_{\tau}^{\dagger}\bowtie_kc_k\tau_A$,
its \emph{conjugated form}, and
\begin{equation}\label{eq:frame-free-row}
  Y_k\ \bowtie_k\ c_k\,\idon{A_0}
\end{equation}
its \emph{frame-free form}.
\begin{enumerate}[label=(\roman*),leftmargin=2.2em]
\item If $\ker W_{\tau}=0$, that is, if the frame is \termref{term-tight}{tight}, the
frame-free form of a homogeneous row holds for $(J_s)_s$ if and only if its
conjugated form does; so imposing some homogeneous rows in the frame-free
form leaves $\Feas^{J}$ unchanged.
\item On every support frame the frame-free form of a homogeneous row
implies its conjugated form. So the set obtained from $\Feas^{J}$ by
imposing some homogeneous rows in the frame-free form is contained in
$\Feas^{J}$, and the infimum in \eqref{eq:choi-program} taken over that set
in place of $\Feas^{J}$ is at least $\hlow$.
\item In the general format, with $\iota=\bigoplus_t\iota_t$ and $W_t$ as in
Theorem~\ref{thm:choi}:
\begin{enumerate}[label=(\alph*),leftmargin=2.2em]
\item for a homogeneous entry read through its \exttermref{../protocol-compiler/contract.pdf}{term-lifted-operator}{lifted operator} and
every $t\in\mathbb{T}$, (i) and (ii) hold for the row of sector $t$, with
$W_t$ and $A_{0,t}$ in place of $W_{\tau}$ and $A_0$ and with the
frame-free form
$\sum_s\tr_{\Qsys_s}[(\idon{A_{0,t}}\otimes\hat O^{\uparrow(s)}_k)
J_{t,s}]\bowtie_kc_k\,\idon{A_{0,t}}$;
\item for a \exttermref{../protocol-compiler/contract.pdf}{term-party-local}{party-local} entry, with $Y_k$ the expression of the last
line of \eqref{eq:feas-general} at $J$ in place of $\omega$, the
unconjugated row $Y_k\bowtie_k0$ is equivalent to the conjugated row
$(W_t\otimes\idon{})\,Y_k\,(W_t^{\dagger}\otimes\idon{})\bowtie_k0$ if
$\ker W_t=0$, and implies it on every frame.
\end{enumerate}
\item In untrusted-source mode there is no conjugation to remove: every
row is a condition on the state variable as written.
\end{enumerate}
\end{proposition}
\claimmeta{\metalabel{Uses}\mref{Def.}{def:homog}, \mref{Def.}{def:feasible}, \mref{Def.}{def:support-frame}, \mref{Lem.}{lem:choi-marginal}, \mref{Thm.}{thm:choi}, \mref{Def.}{def:untrusted-choi}
\metasep \metalabel{Used by}\mref{Rem.}{rem:tightness-by-construction}, \mref{Rem.}{rem:basis-invariance}, \mref{Req.}{req:row-form-by-construction}, \mref{Req.}{req:row-form-record}, \mref{Lem.}{lem:feasible-side-dual}
\metasep newly derived
\metasep accepted}

\begin{proof}
\emph{Idea.} By \eqref{eq:WWtau} the right-hand side $c_k\tau_A$ is
$c_kW_{\tau}W_{\tau}^{\dagger}$, so the conjugated form is a congruence of
the frame-free one. A congruence preserves each of $\preceq$, $=$ and
$\succeq$, and on a tight frame $W_{\tau}$ has a left inverse that undoes
it.

\emph{Step 1 (the conjugated form as a congruence).} By \eqref{eq:WWtau}
(Lemma~\ref{lem:choi-marginal}\,(ii)),
$\Theta_k=c_k\tau_A=c_kW_{\tau}W_{\tau}^{\dagger}$, so the conjugated form
of the $k$-th row is $W_{\tau}XW_{\tau}^{\dagger}\bowtie_k0$ with
$X\coloneqq Y_k-c_k\idon{A_0}$, and the frame-free form is
$X\bowtie_k0$.

\emph{Step 2 (a congruence keeps the relation; part (ii)).} For any linear
map $V$, $X\succeq0$ implies $VXV^{\dagger}\succeq0$, since
$\langle v,VXV^{\dagger}v\rangle=\langle V^{\dagger}v,XV^{\dagger}v\rangle
\ge0$ for every $v$; $X\preceq0$ is the same statement for $-X$, and $X=0$
gives $VXV^{\dagger}=0$. With $V=W_{\tau}$ the frame-free form implies the
conjugated form, on every support frame. A list satisfying the frame-free
forms of some homogeneous rows therefore satisfies their conjugated forms,
and so lies in $\Feas^{J}$. An infimum over a subset is at least the
infimum over the set, and the latter is $\hlow$ by \eqref{eq:choi-program}
(Theorem~\ref{thm:choi}).

\emph{Step 3 (part (i)).} If $\ker W_{\tau}=0$, then $W_{\tau}$ is
injective and has a left inverse $L$, $LW_{\tau}=\idon{A_0}$; for instance
$L=W_{\tau}^{+}$, by \eqref{eq:Wplus} with $P_{\iota}=\idon{A_0}$
(Lemma~\ref{lem:choi-marginal}\,(iii)). Then
$X=L\,(W_{\tau}XW_{\tau}^{\dagger})\,L^{\dagger}$, and Step~2 with $V=L$
shows that the conjugated form implies the frame-free one. With Step~2 the
two forms are equivalent row by row, so imposing either gives the same
set.

\emph{Step 4 (part (iii)).} (a) In the general format the row of sector $t$
of a lifted-operator entry is, by Theorem~\ref{thm:choi},
$W_t\big(\sum_s\tr_{\Qsys_s}[(\idon{A_{0,t}}\otimes\hat O^{\uparrow(s)}_k)
J_{t,s}]\big)W_t^{\dagger}\bowtie_k\Theta^{(t)}_k$, with
$\Theta^{(t)}_k=c_k\tau^{(t)}_A$ for a homogeneous entry
\eqref{eq:feas-general}, and $W_tW_t^{\dagger}=\tau^{(t)}_A$ by
\eqref{eq:WWtau} for each sector. Steps~1 to~3 apply with $W_t$ and
$A_{0,t}$ in place of $W_{\tau}$ and $A_0$, and the value of the general
program is $\hlow$ by the same theorem. (b) For a party-local row, Step~2
with $V=W_t\otimes\idon{}$ shows that $Y_k\bowtie_k0$ implies the conjugated
row on every frame. If $\ker W_t=0$, then $W_t$ has a left inverse $L_t$,
and $L_t\otimes\idon{}$ is a left inverse of $W_t\otimes\idon{}$; Step~3
with it gives the converse.

\emph{Step 5 (part (iv)).} In untrusted-source mode the parametrisation is
the identity and no $W_{\tau}$ is introduced
(Definition~\ref{def:untrusted-choi}\,(a), (b)): the rows of
\eqref{eq:choi-feas-un}, and in the general format those of
\eqref{eq:feas-general}, are conditions on the blocks of $\omega$ as
written.
\end{proof}

% claim meat-core/rem:tightness-by-construction
\claimsummary{meat-core/rem:tightness-by-construction}{Deciding tightness of a given \ensuremath{\tau}\_A is a rank decision on a computed spectrum; the diagonal frame exhibits it by construction, the full frame exhibits nothing. Read at a positive window, the diagonality test would admit the frame-free rows on a frame that is not tight.}
\claimbookmark{meat-core/rem:tightness-by-construction}{Remark}{rem:tightness-by-construction}{Tightness read off the construction of the frame}
\begin{remark}[Tightness read off the construction of the frame]\label{rem:tightness-by-construction}
Tightness of a support frame is $\ker W_{\tau}=0$, equivalently
$d_{A_0}=\operatorname{rank}\tau_A$ (Lemma~\ref{lem:choi-marginal}\,(iii)),
so establishing it \emph{of a given} $\tau_A$ is a rank decision on a
computed spectrum: the decision that \eqref{eq:iota} removes, since
Lemma~\ref{lem:choi-marginal} and Theorem~\ref{thm:choi} hold on every
support frame. What can be known without it is whether the construction
that produced the frame exhibits $\ker W_{\tau}=0$ without computing a rank.

Of the two constructions of Definition~\ref{def:support-frame}, the
diagonal frame~(b) does: there
$\operatorname{ran}\iota=\operatorname{supp}\tau_A$ holds by construction,
the index set $\{i:t_i\neq0\}$ being fixed by a bit-level comparison of
declared binary64 data, which is exact because~(b) is available only for a
$\tau_A$ diagonal in the declared basis, where its diagonal entries
\emph{are} its eigenvalues. This covers the case where every $t_i\neq0$
and construction~(b) returns $\iota=\idon{A}$. The full frame taken as
construction~(a) is built from $\idon{A}$ and from no property of $\tau_A$
at all, and therefore exhibits nothing about $\operatorname{rank}\tau_A$,
however nonsingular $\tau_A$ looks.

The diagonality on which~(b) rests is itself decided by comparing declared
entries with zero. Read at a positive window instead, a singular $\tau_A$
with small off-diagonal entries and every $t_i\ne0$ would pass as
diagonal, and~(b) would return the full frame, which is not tight. The
frame-free form of Proposition~\ref{prop:frame-free-rows} read off that
construction would then be imposed on a non-tight frame, where it imposes
the row on the directions of $\ker W_{\tau}$ as well, cutting $\Feas^{J}$
down to a subset, in general proper, whose image under $\chi_{\Qsys}$ need
not be all of $\Feas$; by Proposition~\ref{prop:frame-free-rows}\,(ii) an
infimum over that subset can only come out higher, the unsafe direction for
a key length. For a party-local row of the general format, the
unconjugated form $Y_k\bowtie_k0$ of
Proposition~\ref{prop:frame-free-rows}\,(iii)(b) is equivalent to the
conjugated row when $\ker W_t=0$, and otherwise implies it and in general
restricts further.

Declining the frame-free form costs nothing: the conjugated form is exact
on every support frame, tight or not (Theorem~\ref{thm:choi}), so imposing
it changes no value and forgoes only the redundant directions of
$\ker W_{\tau}$. That is why the condition may be resolved conservatively
rather than decided numerically.
\end{remark}
\claimmeta{\metalabel{Uses}\mref{Def.}{def:support-frame}, \mref{Lem.}{lem:choi-marginal}, \mref{Thm.}{thm:choi}, \mref{Prop.}{prop:frame-free-rows}
\metasep newly derived
\metasep accepted}

% claim meat-core/rem:ordering
\claimsummary{meat-core/rem:ordering}{In trusted-source mode each Choi block J\_s is the Choi operator of a completely positive branch from A\ensuremath{_{0}} to the sector, input factor first, and only the branches' sum preserves trace. The reading is fixed only up to a unitary relabelling of the input, set by the declared isometry.}
\claimbookmark{meat-core/rem:ordering}{Remark}{rem:ordering}{Ordering convention and channel reading}
\begin{remark}[Ordering convention and channel reading]\label{rem:ordering}
In trusted-source mode each $J_s$ is the (unnormalized) Choi operator of a
CP branch $\mathcal{E}_s:A_0\to\Qsys_s$ in the convention with the
\emph{input} factor $A_0$ first and trace over the \emph{output} equal to
the identity on the input, $\sum_s\tr_{\Qsys_s}J_s=\idon{A_0}$; the induced
action is
$\mathcal E_s(\varrho)=\tr_{A_0}[(\varrho^{\transpose_{A_0}}
\otimes\idon{\Qsys_s})J_s]$ with the transpose taken in the declared basis
of $A_0$, i.e.\ in the columns of the declared isometry $\iota$ of
\eqref{eq:iota}. This matches the convention of
\cite[Eqs.~(82)--(83) and (85)--(86)]{kamin2025renyi}. The reading is
interpretive only: the parametrisation of Lemma~\ref{lem:choi-marginal} and
the program of Theorem~\ref{thm:choi} do not use it. It is relative
to the declared $\iota$: replacing $\iota$ by $\iota U$ for a unitary $U$
on $A_0$ replaces the block representing a given $\omega_s$ by
$J'_s=(U^{\dagger}\otimes\idon{})J_s(U\otimes\idon{})$
(Remark~\ref{rem:basis-invariance}\,(b)), and since
$\tr_{A_0}[(\varrho^{\transpose}\otimes\idon{})J'_s]
=\tr_{A_0}[((U\varrho^{\transpose}U^{\dagger})\otimes\idon{})J_s]$ by
cyclicity in the traced factor, while
$(U\varrho^{\transpose}U^{\dagger})^{\transpose}
=\overline{U}\,\varrho\,\overline{U}^{\dagger}$, the branch reads
$\mathcal E'_s(\varrho)
=\mathcal E_s(\overline{U}\,\varrho\,\overline{U}^{\dagger})$ with
$\overline{U}$ the entrywise complex conjugate of $U$ in the declared
basis. The channel is therefore pinned down only up to that unitary
relabelling of its input, while every quantity this module outputs is
invariant (Remark~\ref{rem:basis-invariance}\,(d)); the declared $\iota$ is
recorded with any reported $J_s$ (Requirement~\ref{req:frame-recorded}),
and Requirement~\ref{req:choi-convention} fixes the factor order and the
transpose basis. The branches are generally
trace-non-increasing; only their sum is trace-preserving, which is what
couples the sector blocks through the single constraint in
\eqref{eq:choi-feas}. All declared sectors, including one-dimensional
classical (flag-type) sectors, are treated uniformly as output sectors of
this parametrized channel.
\end{remark}
\claimmeta{\metalabel{Uses}\mref{Lem.}{lem:choi-marginal}, \mref{Thm.}{thm:choi}
\metasep \metalabel{Used by}\mref{Rem.}{rem:basis-invariance}, \mref{Req.}{req:frame-recorded}, \mref{Req.}{req:choi-convention}
\metasep newly derived
\metasep accepted}

% claim meat-core/rem:basis-invariance
\claimsummary{meat-core/rem:basis-invariance}{Within one support frame the declared isometry is fixed up to a unitary on A\ensuremath{_{0}}, which carries the Choi program onto itself, so the reported estimate, h\_lb and key length are invariant. Across frames only h\_low and the sector report are frame-free, and h\_lb \ensuremath{\le} h\_low holds in each frame.}
\claimbookmark{meat-core/rem:basis-invariance}{Remark}{rem:basis-invariance}{Invariance under the declared basis of the Choi input}
\begin{remark}[Invariance under the declared basis of the Choi input]
\label{rem:basis-invariance}
In trusted-source mode the isometry $\iota$ of \eqref{eq:iota} is not
canonical --- neither its range, the support frame, nor, within a fixed
range, its columns --- and this module declares no canonical choice. It
needs none, and the two freedoms are settled separately. \emph{Within one
range} $\iota$ is canonical up to a unitary on $A_0$, and every quantity
this module computes is invariant under that freedom; that is (a)--(d).
\emph{Across ranges} the programs genuinely differ --- they have
different numbers of variables --- but their value does not: by
Theorem~\ref{thm:choi} each maps \emph{onto} the one set $\Feas$, whose
infimum is $\hlow$ by \eqref{eq:hlow-def}, a definition that never
mentions $A_0$. What is frame-free across ranges is therefore the
\emph{mathematical} $\hlow$ of \eqref{eq:hlow-def}, and not the reported
numbers: the reported \exttermref{../foundations/contract.pdf}{term-estimate}{estimate} is the objective at the returned
\emph{point}, and the reported key length is computed from
$\max\{0,\hlb\}$, so both are functions of the iterates and neither
\emph{is} $\hlow$. Across ranges the programs carry different numbers of
variables, a run in one frame is not the image of a run in the other, and
those three reported numbers --- the estimate, $\hlb$ and the key
length --- may all differ. What (a)--(d) establish is the
\emph{within-range} statement, and there they are invariant, the two runs
corresponding under $\mathrm{Ad}_{U}$. What holds in every frame
regardless is the inequality: $\hlb\le\hlow$ by
Theorem~\ref{thm:firstorder}\,(ii), respectively~(iiib), applied in that
frame, against the one frame-free $\hlow$.

\emph{(a) Within one range the freedom is exactly a unitary.} Let
$\iota,\iota'$ both satisfy \eqref{eq:iota} \emph{with the same range},
write $\Pi\coloneqq\iota\iota^{\dagger}=\iota'\iota'^{\dagger}$ and put
$U\coloneqq\iota^{\dagger}\iota'$. Since
$\operatorname{ran}\iota'=\operatorname{ran}\Pi$ we have
$\Pi\iota'=\iota'$, so
$U^{\dagger}U=\iota'^{\dagger}\iota\,\iota^{\dagger}\iota'
=\iota'^{\dagger}\Pi\iota'=\iota'^{\dagger}\iota'=\idon{A_0}$, and
likewise $UU^{\dagger}=\iota^{\dagger}\Pi\iota=\idon{A_0}$: $U$ is unitary
on $A_0$, and
$\iota U=\iota\,\iota^{\dagger}\iota'=\Pi\iota'=\iota'$. Conversely
$\iota U$ satisfies \eqref{eq:iota}, with the same range, for every
unitary $U$ on $A_0$. The two admissible Choi factors are therefore
related by $W'_{\tau}\coloneqq\tau_A^{1/2}\iota'=W_{\tau}U$, and every
identity of Lemma~\ref{lem:choi-marginal}\,(i)--(iii) holds verbatim for
either:
$W'_{\tau}W'^{\dagger}_{\tau}
=W_{\tau}UU^{\dagger}W_{\tau}^{\dagger}=\tau_A$, and
$W'^{+}_{\tau}=(W_{\tau}U)^{+}=U^{\dagger}W^{+}_{\tau}$ (the
Moore--Penrose inverse of a product with a unitary right factor), so
$W'_{\tau}W'^{+}_{\tau}=W_{\tau}W^{+}_{\tau}
=\Pi_{\operatorname{supp}\tau_A}$ and
$W'^{+}_{\tau}W'_{\tau}=U^{\dagger}P_{\iota}U$, which is $\idon{A_0}$
exactly when the frame is tight. In particular $\ker W'_{\tau}
=U^{\dagger}(\ker W_{\tau})$ has the same dimension, so tightness is a
property of the range and not of the columns.

\emph{(b) The two parametrizations differ by a relabelling.} For a unitary
$U$ on $A_0$ write
$\mathrm{Ad}_{U}(J)\coloneqq(U\otimes\idon{})\,J\,
(U^{\dagger}\otimes\idon{})$ on $\Lin{A_0B'}$. From
$W'_{\tau}=W_{\tau}U$,
\begin{equation}\label{eq:chi-relabel}
  \chi'_{B'}=\chi_{B'}\circ\mathrm{Ad}_{U}.
\end{equation}
Moreover $\mathrm{Ad}_{U}$ carries
$\{J\in\Pos{A_0B'}:\tr_{B'}J=\idon{A_0}\}$ onto itself, with inverse
$\mathrm{Ad}_{U^{\dagger}}$: it preserves positivity, and
$\tr_{B'}\mathrm{Ad}_{U}(J)=U(\tr_{B'}J)U^{\dagger}
=U\idon{A_0}U^{\dagger}=\idon{A_0}$. Hence
Lemma~\ref{lem:choi-marginal} holds for either choice \emph{with the same
image} $\{\omega\in\Pos{AB'}:\tr_{B'}\omega=\tau_A\}$, and the two
parametrizations of one and the same $\omega$ are related by
$J'=\mathrm{Ad}_{U^{\dagger}}(J)$.

\emph{(c) The Choi program is carried onto itself.} Apply
$\mathrm{Ad}_{U}$ blockwise to $(J_s)_{s}$. Positivity and the coupling
constraint $\sum_s\tr_{\Qsys_s}J_s=\idon{A_0}$ are preserved by (b). For a
constraint entry, $(\idon{A_0}\otimes\hat O^{(s)}_k)$ commutes with
$U\otimes\idon{}$, so
\begin{equation*}
  \sum_s\tr_{\Qsys_s}\big[(\idon{A_0}\otimes\hat O^{(s)}_k)\,
  \mathrm{Ad}_{U}(J_s)\big]
  =U\Big(\sum_s\tr_{\Qsys_s}\big[(\idon{A_0}\otimes\hat O^{(s)}_k)\,
  J_s\big]\Big)U^{\dagger},
\end{equation*}
and conjugating by $W_{\tau}$ gives
$W_{\tau}U(\cdot)U^{\dagger}W^{\dagger}_{\tau}
=W'_{\tau}(\cdot)W'^{\dagger}_{\tau}$: the constraint of
\eqref{eq:choi-feas} built with $\iota'$ and imposed on $(J_s)_s$ is the
constraint built with $\iota$ and imposed on
$(\mathrm{Ad}_{U}(J_s))_s$. So $\mathrm{Ad}_{U}$ is a bijection from
$\Feas^{J}$ built with $\iota'$ onto $\Feas^{J}$ built with $\iota$. It
preserves the objective as well, because every ingredient of
\eqref{eq:choi-program} is a pullback along $\chi$: by
\eqref{eq:chi-relabel},
$P'^{J}_{\hat c}(J)=P^{J}_{\hat c}(\mathrm{Ad}_{U}(J))$,
$f'^{J}_{\hat c,s}(J_s)=f^{J}_{\hat c,s}(\mathrm{Ad}_{U}(J_s))$ and
$K'^{J}_{\hat c}(J)=K^{J}_{\hat c}(\mathrm{Ad}_{U}(J))$. The two programs
therefore have the same value, and $\hlow$ is unchanged. (Independently:
$\hlow$ is \emph{defined} by \eqref{eq:hlow-def} as an infimum over
$\Feas$ in the state variable $\omega$, which never mentions $A_0$;
Theorem~\ref{thm:choi} asserts that the Choi program equals that number
for any admissible $\iota$, so the invariance is already forced by the
theorem.)

\emph{(d) Everything downstream is invariant or covariant.} The
frame-free form \eqref{eq:frame-free-row} of a homogeneous row
(Proposition~\ref{prop:frame-free-rows}) is preserved, since
$c_k\idon{A_0}=U(c_k\idon{A_0})U^{\dagger}$ and conjugation by a unitary
preserves each of $\preceq$, $=$, $\succeq$. The branch operators
$X_{\hat c,s}=\cG_{\hat c,s}(\omega_s)$, the
\termref{term-maximal-branch-support}{maximal branch supports}
$W_{\hat c,s}$ of \eqref{eq:Wcs} and every value the objective takes are
\emph{literally} unchanged, being functions of
$\omega_s=\chi_{\Qsys_s}(J_s)$ and of $\operatorname{supp}\tau_A$ alone.
The gradients of Corollary~\ref{cor:full-gradient} are covariant:
$W'^{\dagger}_{\tau}M_xW'_{\tau}
=U^{\dagger}(W^{\dagger}_{\tau}M_xW_{\tau})U$ and
$\chi'^{\dagger}_{\Qsys_{s}}
=\mathrm{Ad}_{U^{\dagger}}\circ\chi^{\dagger}_{\Qsys_{s}}$, so both
$\Gamma^{(s')}_{\hat c}$ and $\Delta^{(s')}_{\hat c}$ transform as
$G\mapsto(U^{\dagger}\otimes\idon{})G(U\otimes\idon{})$, and a
Frank--Wolfe run in one declared basis is the
$\mathrm{Ad}_{U^{\dagger}}$-image of the same run in the other.
Consequently, \emph{within one range}, the reported estimate, $\hlb$ and
the reported key length are independent of which columns were chosen: the
two runs correspond under $\mathrm{Ad}_{U}$ and the objective, the
gradients and the pairings are carried along with them. Only the reported
$J_s$ and the channel reading of Remark~\ref{rem:ordering} are
basis-relative, and they are what the frame is recorded for
(Requirement~\ref{req:frame-recorded}). The sector report is frame-free in
the stronger sense, across ranges as well: its predicates read the declared
blocks $F^{(s)}_v$, $\hat O^{(s)}_k$ and $\{M_x\}$ and the declared tables,
and no $A_0$ enters any of them. \emph{Across ranges the reported numbers
are not claimed invariant}; what survives there is the mathematical
$\hlow$, and $\hlb\le\hlow$ in each frame, as the opening paragraph says.
\end{remark}
\claimmeta{\metalabel{Uses}\mref{Def.}{def:homog}, \mref{Def.}{def:feasible}, \mref{Def.}{def:support-frame}, \mref{Lem.}{lem:choi-marginal}, \mref{Thm.}{thm:choi}, \mref{Prop.}{prop:frame-free-rows}, \mref{Rem.}{rem:ordering}
\metasep newly derived
\metasep accepted}

% claim meat-core/def:untrusted-choi
\claimsummary{meat-core/def:untrusted-choi}{In untrusted-source mode the Choi step is the identity: the variables are the sector blocks of the state, coupled by the single scalar condition that their traces sum to one. At d\_A = 1 this is the trusted-mode Choi program.}
\claimbookmark{meat-core/def:untrusted-choi}{Definition}{def:untrusted-choi}{The program in untrusted-source mode}
\claimleadin{%
In \exttermref{../transmitter/contract.pdf}{term-untrusted-source-mode}{untrusted-source mode} there is no marginal to fix, so there is
nothing to reparametrize. The program is not obtained by substituting
$W_{\tau}=\idon{A}$ into \eqref{eq:choi-feas}: that substitution would give
$\sum_s\tr_{\Qsys_s}\omega_s=\idon{A}$, i.e.\ $\omega_A$ fixed to the
\emph{unnormalized} identity --- an infeasible condition, and the operator
form the licence does not supply (Remark~\ref{rem:homog}\,(a)). The
degeneration runs the other way: $W_{\tau}$ is simply never introduced.
Since part (a) below introduces no $A_0$ and no $\iota$, the declared-basis
question of Remark~\ref{rem:basis-invariance} does not arise at all, and no
support frame is recorded in this mode
(Requirement~\ref{req:frame-recorded}).
}
\begin{definition}[The program in untrusted-source mode]
\label{def:untrusted-choi}
In untrusted-source mode the Choi step is the identity. This is the
degenerate case of Theorem~\ref{thm:choi}, not a second program:
\begin{enumerate}[label=(\alph*),leftmargin=2.2em]
\item \emph{The parametrization.} $\chi_{B'}=\idmap$, $J=\omega$,
$J_s=\omega_s$; no $A_0$, no $\iota$, no $W_{\tau}$.
\item \emph{The variables and constraints.} $\Feas$ of
\eqref{eq:feas} is already the \exttermref{../protocol-compiler/contract.pdf}{term-feasible-set}{feasible set}, and by
\eqref{eq:feas-modes} it reads
\begin{equation}\label{eq:choi-feas-un}
  \Feas=\Big\{(\omega_s)_{s\in\SecSet}:\
  \omega_s\in\Pos{A\Qsys_s},\
  \sum_{s\in\SecSet}\tr\omega_s=1,\
  \sum_{s\in\SecSet}
  \tr\big[(\idon{A}\otimes\hat O^{(s)}_k)\,\omega_s\big]
  \bowtie_k c_k\ \forall k\Big\},
\end{equation}
so the coupling constraint that binds the sector blocks is the single
\emph{scalar} equation $\sum_s\tr\omega_s=1$, in place of the operator
equation $\sum_s\tr_{\Qsys_s}J_s=\idon{A_0}$ of
\eqref{eq:choi-feas}.
\item \emph{The objective.} \eqref{eq:choi-program} holds verbatim with
$J_s$ replaced by $\omega_s$ and $P^{J}_{\hat c}$, $f^J_{\hat c,s}$,
$K^{J}_{\hat c}$ by their state-form originals
$P_{\hat c}$, $f_{\hat c,s}$, $K_{\hat c}$ of \eqref{eq:fK}: the
substitution $\omega_s=\chi_{\Qsys_s}(J_s)$ that
Theorem~\ref{thm:choi} performs is the identity here. Hence the value
is $\hlow$ of \eqref{eq:hlow-def} by definition rather than by a
theorem, the program is jointly convex by
Theorem~\ref{thm:convexity} directly, and the feasible set is convex
and compact by Definition~\ref{def:feasible}.
\item \emph{Consistency at $d_A=1$.} There trusted-source mode admits only
$\tau_A=[1]$, the one state on a one-dimensional $A$, and its Choi
parametrization degenerates: $d_{A_0}=\operatorname{rank}\tau_A=1$,
$\iota=[1]$, $W_{\tau}=[1]$, $\chi_{B'}=\idmap$ and $J_s=\omega_s$. So
$\idon{A_0}=[1]$ and $\sum_s\tr_{\Qsys_s}J_s=\idon{A_0}$ \emph{is}
$\sum_s\tr\omega_s=1$, the constraint rows agree because every $1\times1$
Hermitian $\Theta_k$ is $c_k\tau_A$ with $c_k=\Theta_k$
(Proposition~\ref{prop:trivial-alice-system}\,(iii)), and the two feasible
sets of \eqref{eq:feas} coincide
(Proposition~\ref{prop:trivial-alice-system}\,(i)): \eqref{eq:choi-feas}
and \eqref{eq:choi-feas-un} are the same set, and the two modes give
literally the same program.
\end{enumerate}
\end{definition}
\claimmeta{\metalabel{Uses}\mref{Thm.}{thm:branch-formula}, \mref{Def.}{def:feasible}, \mref{Prop.}{prop:trivial-alice-system}, \mref{Thm.}{thm:convexity}, \mref{Thm.}{thm:choi}
\metasep \metalabel{Used by}\mref{Prop.}{prop:frame-free-rows}, \mref{Lem.}{lem:branch-support}, \mref{Cor.}{cor:full-gradient}, \mref{Lem.}{lem:kraus-freedom}, \mref{Lem.}{lem:objective-homogeneous}, \mref{Def.}{def:first-order-setting}, \mref{Cor.}{cor:fo-dropped}, \mref{Lem.}{lem:feasible-side-dual}
\metasep newly derived
\metasep accepted}

% claim meat-core/lem:estimate
\claimsummary{meat-core/lem:estimate}{The objective at any exactly feasible point, a state or the Choi operators of \hyperlink{claim-meat-core/thm:choi}{Theorem~\ref*{thm:choi}} that parametrise one, is at least the value h\_low of the program. It warrants the estimate every run reports, a solver diagnostic and never the lower bound a key length needs.}
\claimbookmark{meat-core/lem:estimate}{Lemma}{lem:estimate}{Upper estimate of the value at any feasible point}
\claimleadin{%
The number a run reports as its \exttermref{../foundations/contract.pdf}{term-estimate}{estimate} is the objective at the
point its optimizer returned. What that number is worth is one line, and it
is a property of the program, not of any method that searches it. The
membership $x\in\acSet\times\Feas^{J}$ is exactly what must be checked for
this warrant to hold, and which check establishes it is the solver
submodule's claim, not the core's. Definition~\ref{def:admissible-acceptance-set} is not needed:
nothing in the lemma reads convexity or closedness, and at an empty
$\acSet$ the statement is vacuous. In either form the number is an \emph{upper}
estimate of $\hlow$ and never a bound on it: nothing here says how far above
$\hlow$ it lies.
}
\begin{lemma}[Upper estimate of the value at any feasible point]\label{lem:estimate}
Let $\Decl$ satisfy Definition~\ref{def:admissible-declaration}, let $\acSet$ be any set
of probability vectors on the acceptance alphabet, and let
$\alpha\in(1,\infty)$.
\begin{enumerate}[label=(\roman*),leftmargin=2.2em]
\item \emph{(State form.)} For every $(\mathbf q,\omega)\in\acSet\times\Feas$,
\begin{equation}\label{eq:estimate}
  O_{\alpha}(\mathbf q,\omega)\;\ge\;\hlow .
\end{equation}
\item \emph{(Choi form.)} Write
$O^{J}_{\alpha}(\mathbf q,J)\coloneqq O_{\alpha}(\mathbf q,\chi_{\Qsys}(J))$
for the objective of the Choi program \eqref{eq:choi-program} of
Theorem~\ref{thm:choi}, with $\chi_{\Qsys}$ the map of that theorem. For
every $x=(\mathbf q,J)\in\acSet\times\Feas^{J}$,
\begin{equation}\label{eq:estimate-choi}
  O^{J}_{\alpha}(x)\;\ge\;\hlow .
\end{equation}
\end{enumerate}
\end{lemma}
\claimmeta{\metalabel{Uses}\mref{Def.}{def:admissible-declaration}, \mref{Def.}{def:feasible}, \mref{Thm.}{thm:choi}
\metasep \metalabel{Used by}\mref{Thm.}{thm:firstorder}
\metasep newly derived
\metasep implemented (effective: accepted)
\metasep \metalabel{Code}\coderef{https://github.com/fraunhofersg/certiqs-key/blob/v0.1.0/src/certiqs_key/meat_core/objective/objective.py}{certiqs_key.meat_core.objective.objective:objective_value}}

\begin{proof}
(i) $\hlow$ is by \eqref{eq:hlow-def} the infimum of $O_{\alpha}$ over
$\acSet\times\Feas$, and $(\mathbf q,\omega)$ is a point of that set; an
infimum is a lower bound on every value it is taken over.

(ii) By Theorem~\ref{thm:choi}, $\chi_{\Qsys}$ maps $\Feas^{J}$ onto
$\Feas$ and the value of the Choi program is $\hlow$; so for
$x=(\mathbf q,J)\in\acSet\times\Feas^{J}$ the pair
$(\mathbf q,\chi_{\Qsys}(J))$ lies in $\acSet\times\Feas$ and
$O^{J}_{\alpha}(x)=O_{\alpha}(\mathbf q,\chi_{\Qsys}(J))\ge\hlow$ by (i).
\end{proof}

\section{Gradients and the limit of order one}\label{group-meat-core/def:branch-support}
\sectionintro{The \termref{term-maximal-branch-support}{maximal branch support} is the span of the pass-class Kraus images of the admissible Alice subspace tensored with the sector; every branch operator the program forms lies in it, with equality at a positive definite Choi block (\hyperlink{claim-meat-core/lem:branch-support}{Lemma~\ref*{lem:branch-support}}). On the face of a singular operator the pinched trace power has a gradient in split form, with divided differences over the positive spectrum and pinched operators on the ambient space (\hyperlink{claim-meat-core/thm:gradient}{Theorem~\ref*{thm:gradient}}). The full objective is differentiable along every feasible direction at a point of maximal support with positive probabilities (\hyperlink{claim-meat-core/cor:full-gradient}{Corollary~\ref*{cor:full-gradient}}); its branch values and pulled-back gradient depend on the Kraus factors only through their Gram matrices, so exact factors may replace the square roots (\hyperlink{claim-meat-core/lem:kraus-freedom}{Lemma~\ref*{lem:kraus-freedom}}). As the R\'enyi order tends to one the value converges to a blockwise relative-entropy program when every acceptable distribution gives each key-live value probability at least some p\ensuremath{_{0}} \textgreater{} 0 and some feasible state reproduces an acceptable distribution (\hyperlink{claim-meat-core/thm:oracle}{Theorem~\ref*{thm:oracle}}), a validation reference and never a bound.}
\sectionimports{\extref{../foundations/contract.pdf}{foundations/def:kl}, \extref{../foundations/contract.pdf}{foundations/def:pinching}, \extref{../foundations/contract.pdf}{foundations/fact:concavity}, \extref{../foundations/contract.pdf}{foundations/lem:kernel}, \extref{../foundations/contract.pdf}{foundations/lem:matrix-function-perturbation}}

% claim meat-core/def:branch-support
\claimsummary{meat-core/def:branch-support}{The maximal branch support of a key-live acceptance value and a sector is the span of the pass-class Kraus images of the admissible Alice subspace tensored with the sector. The side-by-side matrix of those Kraus operators depends on the declaration and the support frame alone.}
\claimbookmark{meat-core/def:branch-support}{Definition}{def:branch-support}{Maximal branch support}
\claimleadin{%
A branch operator of a \termref{term-key-live-acceptance-value}{key-live acceptance value} and a sector cannot
reach beyond a subspace that the declaration fixes. That subspace is the face
on which Theorem~\ref{thm:gradient} differentiates the pinched trace power,
and Lemma~\ref{lem:branch-support} records how the branch operators the
program forms sit inside it.
}
\begin{definition}[Maximal branch support]\label{def:branch-support}
For $\hat c\in\Chkey$ and $s\in\SecSet$ the
\termdef{term-maximal-branch-support}{maximal branch support}
is the subspace
\begin{equation}\label{eq:Wcs}
  W_{\hat c,s}\coloneqq
  \sum_{\kappa\in\Kann_{\mathrm{pass}}(\hat c)}
  K^{(s)}_{\kappa}\big(\Sadm\otimes\Hs{\Qsys_s}\big)
  \ \subseteq\ \Hs{R\,XY\,C_AC_B\,A\Qsys_s},
\end{equation}
the span of the images, under the Kraus operators of $\cG_{\hat c,s}$
(Definition~\ref{def:branch}), of the admissible Alice subspace
$\Sadm$ of \eqref{eq:SA} tensored with the sector, that is, of the
subspace the declaration admits: $\operatorname{supp}\tau_A$ in
\exttermref{../transmitter/contract.pdf}{term-trusted-source-mode}{trusted-source mode}, all of $\Hs{A}$ in
\exttermref{../transmitter/contract.pdf}{term-untrusted-source-mode}{untrusted-source mode}.

Listing the operators of \eqref{eq:Wcs} side by side, with $W_{\tau}$ of
\eqref{eq:Wtau} on the \termref{term-support-frame}{support frame} of
Definition~\ref{def:support-frame}, put
\begin{equation}\label{eq:Kstack}
  \mathbf K_{\hat c,s}\coloneqq
  \Big[\ K^{(s)}_{\kappa}\,(W_{\tau}\otimes\idon{\Qsys_s})\ \Big]
  _{\kappa\in\Kann_{\mathrm{pass}}(\hat c)}
  \;:\;\bigoplus_{\kappa\in\Kann_{\mathrm{pass}}(\hat c)}\Hs{A_0\Qsys_s}
  \;\longrightarrow\;\Hs{R\,XY\,C_AC_B\,A\Qsys_s},
\end{equation}
with $W_{\tau}\otimes\idon{\Qsys_s}$ replaced by $\idon{A\Qsys_s}$, and
the domain by $\bigoplus_{\kappa}\Hs{A\Qsys_s}$, in untrusted-source mode.
The matrix $\mathbf K_{\hat c,s}$ depends on the declaration and on
$W_{\tau}$ alone and on no iterate.
\end{definition}
\claimmeta{\metalabel{Uses}\mref{Def.}{def:branch}, \mref{Lem.}{lem:admissible-alice-subspace}, \mref{Def.}{def:support-frame}
\metasep \metalabel{Used by}\mref{Lem.}{lem:branch-support}, \mref{Thm.}{thm:gradient}, \mref{Cor.}{cor:full-gradient}, \mref{Req.}{req:split-form-gradient}, \mref{Lem.}{lem:fo-face}, \mref{Thm.}{thm:firstorder}, \mref{Cor.}{cor:fo-dropped}, \mref{Rem.}{rem:fo-one-ruler}
\metasep newly derived
\metasep implemented (effective: accepted)
\metasep \metalabel{Code}\coderef{https://github.com/fraunhofersg/certiqs-key/blob/v0.1.0/src/certiqs_key/meat_core/objective/branch.py}{certiqs_key.meat_core.objective.branch:branch_support_of}}

% claim meat-core/lem:branch-support
\claimsummary{meat-core/lem:branch-support}{Every branch operator the program forms is supported in the maximal branch support, with equality at a positive definite Choi block; the support is the range of the stacked Kraus matrix. It is always smaller than the ambient space and in general not invariant under the key pinching.}
\claimbookmark{meat-core/lem:branch-support}{Lemma}{lem:branch-support}{Properties of the maximal branch support}
\claimleadin{%
Four facts about the \termref{term-maximal-branch-support}{maximal branch support} are used by the gradient
of Theorem~\ref{thm:gradient} and by the first-order argument of the
Frank--Wolfe submodule: it bounds the support of every branch operator the
program forms, it is the range of the stacked Kraus matrix, it is always
strictly smaller than the ambient space, and it is in general not invariant
under the \termref{term-key-pinching}{key pinching}. The bound in~(ii) is not a number a run computes:
$\operatorname{rank}\tau_A$ enters it only as a mathematical quantity.
Hypothesis (H-SUPP) of Theorem~\ref{thm:firstorder} compares a rank with
$\dim W_{\hat c,s}$, and how that comparison is decided is
Requirement~\ref{req:support-rank-check}.
}
\begin{lemma}[Properties of the maximal branch support]\label{lem:branch-support}
Let $\hat c\in\Chkey$ and $s\in\SecSet$, and let $W_{\hat c,s}$ and
$\mathbf K_{\hat c,s}$ be as in Definition~\ref{def:branch-support}.
\begin{enumerate}[label=(\roman*),leftmargin=2.2em]
\item In trusted-source mode every $\omega_s=\chi_{\Qsys_s}(\Delta)$ with
  $\Delta\in\Pos{A_0\Qsys_s}$, in particular every feasible $\omega_s$
  and the image of the positive and the negative part of every Hermitian
  direction, satisfies
  $\operatorname{supp}\omega_s
  =(W_{\tau}\otimes\idon{})(\operatorname{supp}\Delta)
  \subseteq\operatorname{supp}\tau_A\otimes\Hs{\Qsys_s}
  =\Sadm\otimes\Hs{\Qsys_s}$; in untrusted-source mode
  $\Sadm=\Hs{A}$ and the inclusion holds for every
  $\omega_s\in\Pos{A\Qsys_s}$. In either mode, therefore,
  \begin{equation}\label{eq:Wcs-bound}
    \operatorname{supp}\cG_{\hat c,s}(\omega_s)\subseteq W_{\hat c,s},
  \end{equation}
  with equality whenever $\omega_s$ attains that maximal support,
  $\operatorname{supp}\omega_s=\Sadm\otimes\Hs{\Qsys_s}$. In
  trusted-source mode $J_s\succ0$ for the Choi block suffices for it, and
  on a tight support frame is equivalent to it; in untrusted-source mode it
  is $\omega_s\succ0$ on $\Hs{A\Qsys_s}$, which is $J_s\succ0$, as
  $J_s=\omega_s$ there (Definition~\ref{def:untrusted-choi}\,(a)).
\item $W_{\hat c,s}$ is the range of one matrix, and its dimension is
  bounded:
  \begin{equation}\label{eq:Kstack-rank}
    W_{\hat c,s}=\operatorname{ran}\mathbf K_{\hat c,s},
    \qquad
    \dim W_{\hat c,s}=\operatorname{rank}\mathbf K_{\hat c,s}
    \;\le\;\big|\Kann_{\mathrm{pass}}(\hat c)\big|\cdot
    \dim\Sadm\cdot d_s ,
  \end{equation}
  where $\dim\Sadm=\operatorname{rank}\tau_A$ in trusted-source mode and
  $d_A$ in untrusted-source mode.
\item $\dim W_{\hat c,s}$ is \emph{strictly} smaller than
  $\dim\Hs{R\,XY\,C_AC_B\,A\Qsys_s}$, so $0$ always lies in the ambient
  spectrum of a branch operator $X_{\hat c,s}=\cG_{\hat c,s}(\omega_s)$ of~(i)
  and ``maximal support'' is never the statement that $X_{\hat c,s}\succ0$
  on the ambient space.
\item $W_{\hat c,s}$ is \emph{not} in general $\cZ$-invariant.
\end{enumerate}
\end{lemma}
\claimmeta{\metalabel{Uses}\mref{Def.}{def:admissible-declaration}, \mref{Def.}{def:branch}, \mref{Lem.}{lem:admissible-alice-subspace}, \mref{Lem.}{lem:choi-marginal}, \mref{Def.}{def:untrusted-choi}, \mref{Def.}{def:branch-support}
\metasep \metalabel{Used by}\mref{Thm.}{thm:gradient}, \mref{Cor.}{cor:full-gradient}, \mref{Req.}{req:gradient-at-maximal-support}, \mref{Lem.}{lem:fo-face}, \mref{Thm.}{thm:firstorder}, \mref{Cor.}{cor:fo-dropped}, \mref{Cor.}{cor:fo-reference}, \mref{Rem.}{rem:fo-general-format}, \mref{Lem.}{lem:certified-bracket}, \mref{Rem.}{rem:certified-scope}, \mref{Rem.}{rem:fo-one-ruler}, \mref{Rem.}{rem:fo-reachability}, \mref{Req.}{req:support-rank-check}, \mref{Req.}{req:certified-evaluation}
\metasep newly derived
\metasep implemented (effective: accepted)
\metasep \metalabel{Code}\coderef{https://github.com/fraunhofersg/certiqs-key/blob/v0.1.0/src/certiqs_key/meat_core/objective/gradient.py}{certiqs_key.meat_core.objective.gradient:is_support_maximal}}

\begin{proof}
\emph{Support of a sum.} For $\varrho\succeq0$ the support of
$\sum_{\kappa}K_{\kappa}\varrho K_{\kappa}^{\dagger}$ is
$\sum_{\kappa}K_{\kappa}(\operatorname{supp}\varrho)$: a vector $\ket w$
in the kernel of the sum has
$\bra wK_{\kappa}\varrho K_{\kappa}^{\dagger}\ket w=0$, hence
$\varrho^{1/2}K^{\dagger}_{\kappa}\ket w=0$, for every $\kappa$, and
conversely each summand has range in $K_{\kappa}(\operatorname{supp}
\varrho)$.

(i) In trusted-source mode,
$\operatorname{supp}\omega_s=(W_{\tau}\otimes\idon{})(\operatorname{supp}\Delta)$
by the support of a sum with the single operator $W_{\tau}\otimes\idon{}$;
it lies in $\operatorname{ran}W_{\tau}\otimes\Hs{\Qsys_s}$, and
$\operatorname{ran}W_{\tau}=\operatorname{supp}\tau_A$
(Lemma~\ref{lem:choi-marginal}\,(iii)), which is $\Sadm$ by \eqref{eq:SA}. In
untrusted-source mode $\Sadm=\Hs{A}$ and the inclusion is trivial. The
support of a sum then gives
$\operatorname{supp}\cG_{\hat c,s}(\omega_s)
=\sum_{\kappa}K^{(s)}_{\kappa}(\operatorname{supp}\omega_s)
\subseteq\sum_{\kappa}K^{(s)}_{\kappa}(\Sadm\otimes\Hs{\Qsys_s})
=W_{\hat c,s}$, with equality where
$\operatorname{supp}\omega_s=\Sadm\otimes\Hs{\Qsys_s}$. That is the case
when $J_s\succ0$, since then
$\operatorname{supp}\omega_s=\operatorname{ran}W_{\tau}\otimes\Hs{\Qsys_s}$,
and when $\omega_s\succ0$ on $\Hs{A\Qsys_s}$ in untrusted-source mode. On a
tight frame the converse holds as well: $W_{\tau}$ is then injective
(Lemma~\ref{lem:choi-marginal}\,(iii)), so $W_{\tau}\otimes\idon{}$ maps
$\operatorname{supp}J_s$ injectively onto $\operatorname{supp}\omega_s$,
and $\operatorname{supp}\omega_s=\operatorname{ran}W_{\tau}\otimes
\Hs{\Qsys_s}$, of dimension $d_{A_0}d_s$, forces
$\operatorname{supp}J_s=\Hs{A_0\Qsys_s}$, that is $J_s\succ0$.

(ii) The range of a side-by-side listing is the sum of the ranges of its
blocks, and
$\operatorname{ran}(W_{\tau}\otimes\idon{\Qsys_s})
=\operatorname{ran}W_{\tau}\otimes\Hs{\Qsys_s}
=\Sadm\otimes\Hs{\Qsys_s}$ by
$\operatorname{ran}W_{\tau}=\operatorname{supp}\tau_A$
(Lemma~\ref{lem:choi-marginal}\,(iii)), so
$\operatorname{ran}\mathbf K_{\hat c,s}
=\sum_{\kappa}K^{(s)}_{\kappa}(\Sadm\otimes\Hs{\Qsys_s})$, which is
\eqref{eq:Wcs}. The rank of a side-by-side listing is at most the sum of
the ranks of its blocks, and each block has rank at most
$\operatorname{rank}(W_{\tau}\otimes\idon{\Qsys_s})=\dim\Sadm\cdot d_s$,
which gives the inequality. In untrusted-source mode the same two steps
apply with $\idon{A\Qsys_s}$ in place of $W_{\tau}\otimes\idon{\Qsys_s}$.

(iii) For a \termref{term-pass-class}{pass class} $\kappa$ of $\hat c$ every
$(x,v)\in a^{-1}(\kappa)$ has $r(x,v)\in\Zval$, by announcement-determined
sifting and Definition~\ref{def:admissible-declaration}\,(iv), (v); so every Kraus operator
$K^{(s)}_{\kappa}$ of \eqref{eq:Wcs} carries an $R$-ket $\ket{j}_R$ with
$j\in\Zval$, and $W_{\hat c,s}\subseteq\operatorname{span}\{\ket{j}_R:
j\in\Zval\}\otimes\Hs{XY\,C_AC_B\,A\Qsys_s}$. The nonzero vectors
$\ket{\perp}_R\otimes\xi$ are orthogonal to that subspace, so
$\dim W_{\hat c,s}$ is smaller than the ambient dimension. By~(i),
$\operatorname{supp}X_{\hat c,s}\subseteq W_{\hat c,s}$, so those vectors
lie in the ambient kernel of $X_{\hat c,s}$, and $0$ is an eigenvalue of
$X_{\hat c,s}$ on the ambient space.

(iv) The declaration is free to let the raw-key value $r(x,v)$ vary with
$x$ inside one announcement class, since
Definition~\ref{def:admissible-declaration}\,(viii) constrains only that the key holder can
compute it, and then a single Kraus operator \eqref{eq:kraus} has range
vectors spread over several $R$-eigenvalues. $Z_j$ applied to such a
vector $K^{(s)}_{\kappa}\ket\psi$ retains only the terms with
$r(x,v)=j$, which is in general not $K^{(s)}_{\kappa}\ket{\psi'}$ for any
$\ket{\psi'}$, nor in the range of any other Kraus operator (the classes
$\kappa$ occupy orthogonal $C_AC_B$ sectors). An instance: $d_A=1$, one
Alice outcome $x$, $\Qsys=\mathbb{C}$ with one sector $s$, two outcomes
$v_0,v_1$ with $F_{v_0}=F_{v_1}=\frac12$, both outcome pairs in one
announcement class $\kappa$, which passes sifting, and the key held by
Bob with $r(x,v_0)=0$ and $r(x,v_1)=1$, a function of his readout; let
$\hat c\coloneqq\phi(\kappa)$ have no other pass class. Then
$W_{\hat c,s}$ is the line spanned by
$\big(\ket{0}_R\ket{x,v_0}+\ket{1}_R\ket{x,v_1}\big)\ket{\kappa}$, and
$Z_0$ maps that vector to $\ket{0}_R\ket{x,v_0}\ket{\kappa}$, which is not
on the line.
\end{proof}

% claim meat-core/lem:dk
\claimsummary{meat-core/lem:dk}{At a Hermitian operator, a matrix function continuously differentiable around its spectrum has as derivative its divided-difference kernel between spectral projectors, self-adjoint under the trace pairing. At trace level only the diagonal kernel survives.}
\claimbookmark{meat-core/lem:dk}{Lemma}{lem:dk}{Divided-difference derivative of a matrix function}
\claimleadin{%
Theorem~\ref{thm:gradient} differentiates $X\mapsto X^{p}$ inside the
support of a branch operator with the kernel below, and
Lemma~\ref{lem:b-derivative} uses its trace form.
}
\begin{technicalclaim}
\begin{lemma}[Divided-difference derivative of a matrix function]\label{lem:dk}
Write $p\coloneqq1/\alpha\in(0,1)$ and $T$ for the functional
\eqref{eq:Tcone} of Lemma~\ref{lem:cone-concavity} on the positive
semidefinite cone of a branch space; $\log_2\mathrm{e}=1/\ln2$.

Let $X$ be Hermitian with spectral decomposition $X=\sum_i\lambda_iP_i$, and
let $f$ be continuously differentiable on an open interval
$I\supseteq\operatorname{spec}X$. Then $X'\mapsto f(X')$, defined on the
Hermitian operators with spectrum in $I$, is Fr\'echet differentiable at $X$
with
\begin{equation}\label{eq:dk}
  Df(X)[H]=\sum_{i,j}f^{[1]}(\lambda_i,\lambda_j)\,P_iHP_j,
  \qquad
  f^{[1]}(\lambda,\mu)\coloneqq
  \begin{cases}\dfrac{f(\lambda)-f(\mu)}{\lambda-\mu},&\lambda\ne\mu,\\[1ex]
  f'(\lambda),&\lambda=\mu,\end{cases}
\end{equation}
and for Hermitian $A$, $H$:
$\tr\!\big[A\,Df(X)[H]\big]=\tr\!\big[Df(X)[A]\,H\big]$. In particular
this applies to $f(x)=x^{p}$ at every $X\succ0$, with $I=(0,\infty)$; and,
at the trace level, for every Hermitian $Y$ and every $g$ continuously
differentiable on an open interval containing $\operatorname{spec}Y$, the
map $Y'\mapsto\tr g(Y')$ is Fr\'echet differentiable at $Y$ with
\begin{equation}\label{eq:trace-diff}
  \frac{\mathrm d}{\mathrm dt}\Big|_{t=0}\tr\,g(Y+tK)=\tr\big[g'(Y)K\big]
  \qquad\text{for every Hermitian }K.
\end{equation}
For $\alpha>1$ the function $g(y)\coloneqq\max\{y,0\}^{\alpha}$ is
continuously differentiable on $\mathbb{R}$, with
$g'(y)=\alpha\max\{y,0\}^{\alpha-1}$, and $\tr g(Y)=\tr[Y^{\alpha}]$ for
every $Y\succeq0$.
\end{lemma}
\end{technicalclaim}
\claimmeta{\metalabel{Used by}\mref{Thm.}{thm:gradient}, \mref{Lem.}{lem:b-derivative}
\metasep newly derived
\metasep implemented
\metasep \metalabel{Code}\coderef{https://github.com/fraunhofersg/certiqs-key/blob/v0.1.0/src/certiqs_key/meat_core/objective/linalg.py}{certiqs_key.meat_core.objective.linalg:divided_difference}
\metasep \metalabel{Oracles}\coderef{https://github.com/fraunhofersg/certiqs-key/blob/v0.1.0/tests/oracles/registry.py}{registry.py}}

\begin{proof}
\begin{sloppypar}
\emph{Polynomials.} For a monomial $f(x)=x^{m}$,
$\frac{\mathrm d}{\mathrm dt}(X+tH)^{m}\big|_{0}
=\sum_{a+b=m-1}X^{a}HX^{b}$, and
$\bra{i}\sum_{a+b=m-1}X^aHX^b\ket{j}
=H_{ij}\sum_{a+b=m-1}\lambda_i^a\lambda_j^b
=H_{ij}\,f^{[1]}(\lambda_i,\lambda_j)$, which is \eqref{eq:dk}; linearity
extends it to polynomials. A polynomial map $X'\mapsto q(X')$ is a finite
sum of products of $X'$ with itself, hence Fr\'echet differentiable at every
$X$, and its Fr\'echet derivative is the directional derivative just
computed.

\emph{A Hilbert--Schmidt bound.} Let $h$ be continuously differentiable on a
compact interval $[c,C]$, and let $A=\sum_ia_iP_i$ and $B=\sum_jb_jP'_j$ be
Hermitian with spectra in $[c,C]$. Then
$P_i\big(h(A)-h(B)\big)P'_j=(h(a_i)-h(b_j))P_iP'_j$ and
$P_i(A-B)P'_j=(a_i-b_j)P_iP'_j$, so
$h(A)-h(B)=\sum_{i,j}h^{[1]}(a_i,b_j)\,P_i(A-B)P'_j$, the terms with
$a_i=b_j$ being $0$ on both sides. The maps $Z\mapsto P_iZP'_j$ are
orthogonal projections of the Hilbert--Schmidt space onto mutually
orthogonal subspaces that sum to it, and $|h^{[1]}|\le\|h'\|_{\infty}$ on
$[c,C]^2$ by the mean value theorem, so
$\|h(A)-h(B)\|_2\le\|h'\|_{\infty}\,\|A-B\|_2$; in the same way
$\|\sum_{i,j}h^{[1]}(\lambda_i,\lambda_j)P_iHP_j\|_2
\le\|h'\|_{\infty}\,\|H\|_2$.

\emph{General $f$.} Choose a compact $[c,C]\subset I$ and $\delta>0$ with
$\lambda_i\pm\delta\in[c,C]$ for every $i$. By Weyl's inequality every
eigenvalue of $X+H$ lies within $\|H\|$ of one of $X$, so
$\operatorname{spec}(X+H)\subseteq[c,C]$ whenever $\|H\|_2\le\delta$. Pick
polynomials $q_n\to f'$ uniformly on $[c,C]$ (Weierstrass) and let $Q_n$ be
primitives with $Q_n(c)=f(c)$, so that
$\varepsilon_n\coloneqq\sup_{[c,C]}|f'-Q_n'|\to0$. Write
$L(H)\coloneqq\sum_{i,j}f^{[1]}(\lambda_i,\lambda_j)P_iHP_j$ and $L_n$ for
the same expression with $Q_n$, which is $DQ_n(X)$ by the polynomial step.
For $\|H\|_2\le\delta$ the Hilbert--Schmidt bound, applied to $h=f-Q_n$,
whose divided differences are those of $f$ less those of $Q_n$, gives
\begin{equation*}
  \big\|\big(f(X+H)-f(X)-L(H)\big)
  -\big(Q_n(X+H)-Q_n(X)-L_n(H)\big)\big\|_2
  \;\le\;2\varepsilon_n\|H\|_2 .
\end{equation*}
Given $\eta>0$, fix $n$ with $2\varepsilon_n\le\eta/2$ and then
$\delta'\in(0,\delta]$ with
$\|Q_n(X+H)-Q_n(X)-L_n(H)\|_2\le\frac{\eta}{2}\|H\|_2$ for
$\|H\|_2\le\delta'$; then $\|f(X+H)-f(X)-L(H)\|_2\le\eta\|H\|_2$ for
$\|H\|_2\le\delta'$. The estimate holds uniformly in $H$ on that ball, so
$f$ is Fr\'echet differentiable at $X$ with $Df(X)=L$, which is
\eqref{eq:dk}.

\emph{Self-adjointness.} With $k_{ij}\coloneqq
f^{[1]}(\lambda_i,\lambda_j)=k_{ji}$ real,
$\tr[A\sum_{ij}k_{ij}P_iHP_j]=\sum_{ij}k_{ij}\tr[P_jAP_iH]
=\tr\big[\big(\sum_{ij}k_{ji}P_iAP_j\big)H\big]
=\tr[Df(X)[A]\,H]$.

\emph{Trace level.} The general step applies to $g$ at $Y=\sum_i\mu_iP_i$,
so $Y'\mapsto g(Y')$ is Fr\'echet differentiable at $Y$ with derivative
$\sum_{i,j}g^{[1]}(\mu_i,\mu_j)P_iKP_j$, and so is its composition with the
linear trace. Only the diagonal kernel survives the trace:
$\tr[P_iKP_j]=\delta_{ij}\tr[P_iK]$, so
$\frac{\mathrm d}{\mathrm dt}\tr g(Y+tK)\big|_0
=\sum_ig'(\mu_i)\tr[P_iK]=\tr[g'(Y)K]$. For $\alpha>1$ the function
$g(y)=\max\{y,0\}^{\alpha}$ has derivative $\alpha y^{\alpha-1}$ at $y>0$,
$0$ at $y<0$, and $0$ at $y=0$, since $|g(y)|/|y|\le|y|^{\alpha-1}\to0$;
that derivative is continuous because $\alpha-1>0$; and $g$ is $y^{\alpha}$
on $[0,\infty)$, so $\tr g(Y)=\tr[Y^{\alpha}]$ for $Y\succeq0$.
\end{sloppypar}
\end{proof}

% claim meat-core/thm:gradient
\claimsummary{meat-core/thm:gradient}{The pinched trace power is differentiable at a positive definite operator, and on the face of a singular one in split form: divided differences over the positive spectrum, pinched operators on the ambient space. A one-sided derivative is +\ensuremath{\infty} only along a direction that enlarges the support.}
\claimbookmark{meat-core/thm:gradient}{Theorem}{thm:gradient}{Gradient of the pinched trace power on a face}
\claimleadin{%
The key-active branch operators are in general singular, and their
\termref{term-maximal-branch-support}{maximal branch support} is in general not invariant under the
\termref{term-key-pinching}{key pinching} (Lemma~\ref{lem:branch-support}\,(iv)); part~(i) is therefore
stated so that no such invariance is needed, and
Requirement~\ref{req:split-form-gradient} is the rule it gives an
implementation. Part~(ii) is why gradients are taken at points of maximal
support (Requirement~\ref{req:gradient-at-maximal-support}). At maximal
support the split form gives the true one-sided derivative in every
feasible direction, because there no feasible direction enlarges the
support; that is Lemma~\ref{lem:fo-face} of the Frank--Wolfe submodule,
where Theorem~\ref{thm:firstorder} consumes it, and
Remark~\ref{rem:fo-nopert} explains why the explicit-perturbation route has
no accounted instance.
}
\begin{theorem}[Gradient of the pinched trace power on a face]\label{thm:gradient}
Let $\alpha\in(1,\infty)$, $p=1/\alpha$, let $T$ be the pinched trace
power \eqref{eq:Tcone} on the positive semidefinite cone of a branch space
(Lemma~\ref{lem:dk}), and let $X\succ0$. Then $T$ is Fr\'echet
differentiable at $X$ with gradient
\begin{equation}\label{eq:gradT}
  \nabla T(X)
  =\alpha\,D X^{p}\Big[\cZ\big(Y^{\alpha-1}\big)\Big],
  \qquad Y\coloneqq\cZ(X^{p}),
\end{equation}
i.e.\ $\frac{\mathrm d}{\mathrm dt}T(X+tH)\big|_0=\tr[\nabla T(X)H]$,
where $DX^{p}[\cdot]$ is the self-adjoint map \eqref{eq:dk} with
$f(x)=x^{p}$ and $Y^{\alpha-1}$ is defined by functional calculus
($0^{\alpha-1}\coloneqq0$; $Y$ may be singular). Explicitly, in the
eigenbasis $X=\sum_i\lambda_iP_i$:
\begin{equation}\label{eq:gradT-kernel}
  \nabla T(X)=\alpha\sum_{i,j}
  \frac{\lambda_i^{p}-\lambda_j^{p}}{\lambda_i-\lambda_j}\,
  P_i\,\cZ\big(Y^{\alpha-1}\big)\,P_j
  \qquad(\text{diagonal kernel }p\lambda_i^{p-1}).
\end{equation}
Moreover:
\begin{enumerate}[label=(\roman*),leftmargin=2.2em]
\item \emph{(faces)} Let $X\succeq0$, write
$S\coloneqq\operatorname{supp}X$, and let $H$ be Hermitian with
$\operatorname{supp}H\subseteq S$. Then $t\mapsto T(X+tH)$ is
differentiable at $t=0$, and \eqref{eq:gradT}--\eqref{eq:gradT-kernel}
hold in the following \emph{split} form:
\begin{enumerate}[label=(\alph*),leftmargin=2.2em]
\item the divided-difference map $DX^{p}[\cdot]$ is compressed to $S$,
equivalently, the sum in \eqref{eq:gradT-kernel} runs over the strictly
positive eigenvalues of $X$ only, where $X$ is positive definite and the
kernel $f^{[1]}$ is finite;
\item $Y=\cZ(X^{p})$, $Y^{\alpha-1}$ and $\cZ(Y^{\alpha-1})$ are formed on
the \emph{ambient} branch space and are \emph{not} compressed to $S$.
\end{enumerate}
The resulting $\nabla T(X)$ determines the derivative in every direction
supported in $S$; its components outside $S$ are neither determined by
this nor needed, support-enlarging directions being the subject of (ii).
Compressing the \emph{whole} computation to a subspace $S'$ is licensed
when $\sum_{j\in\Zset}Z_jS\subseteq S'$, in particular when $S\subseteq S'$
and $S'$ is $\cZ$-invariant ($Z_jS'\subseteq S'$ for every $j\in\Zset$):
then the ambient complement contributes $0$ to $T$, and the compressed and
ambient computations agree. Compressing to $S$ itself is licensed
\emph{only} when $S$ is $\cZ$-invariant; otherwise $\cZ$ transports weight
out of $S$ and the compressed value of $T$ is strictly smaller (the proof
gives the argument and a rank-one witness). For the branch operators $X=\cG_{\hat c,s}(\omega_s)$, whose
faces are in general not $\cZ$-invariant, the relevant face is
$S=W_{\hat c,s}$ of \eqref{eq:Wcs} at maximal support, and (a)--(b) is
what makes the gradient available there.
\item \emph{(one-sided derivatives)} For $X\succeq0$ and any Hermitian
$H$ with $X+tH\succeq0$ for small $t>0$, the one-sided directional
derivative $\lim_{t\downarrow0}\frac{T(X+tH)-T(X)}{t}$ exists in
$(-\infty,+\infty]$, and can equal $+\infty$ only when $H$ strictly
enlarges the support of $X$.
\end{enumerate}
\end{theorem}
\claimmeta{\metalabel{Uses}\mref{Lem.}{lem:branch-maps}, \mref{Cor.}{cor:concavity}, \mref{Lem.}{lem:cone-concavity}, \mref{Def.}{def:branch-support}, \mref{Lem.}{lem:branch-support}, \mref{Lem.}{lem:dk}, \extref{../foundations/contract.pdf}{foundations/fact:concavity}
\metasep \metalabel{Used by}\mref{Cor.}{cor:full-gradient}, \mref{Lem.}{lem:kraus-freedom}, \mref{Rem.}{rem:cgate}, \mref{Req.}{req:split-form-gradient}, \mref{Req.}{req:gradient-at-maximal-support}, \mref{Lem.}{lem:fo-face}, \mref{Thm.}{thm:firstorder}, \mref{Rem.}{rem:fo-danger}, \mref{Rem.}{rem:fo-nopert}, \mref{Cor.}{cor:fo-dropped}, \mref{Lem.}{lem:certified-bracket}
\metasep newly derived
\metasep implemented (effective: accepted)
\metasep \metalabel{Code}\coderef{https://github.com/fraunhofersg/certiqs-key/blob/v0.1.0/src/certiqs_key/meat_core/objective/gradient.py}{certiqs_key.meat_core.objective.gradient:state_gradient}
\metasep \metalabel{Oracles}\coderef{https://github.com/fraunhofersg/certiqs-key/blob/v0.1.0/tests/oracles/registry.py}{registry.py}}

\begin{proof}
\emph{Idea.} At $X\succ0$ the gradient is the chain rule through
$X\mapsto X^{p}$ (Lemma~\ref{lem:dk}), the linear pinching and
$Y\mapsto\tr[Y^{\alpha}]$. On a face the inner map is differentiated
inside the support, where $X$ is positive definite, while the outer map is
differentiable at every $Y\succeq0$ and so stays on the ambient space.
Concavity of $T$ locates the only infinite one-sided derivatives.

\emph{Step 1 (differentiability at $X\succ0$).} $X\mapsto X^{p}$ is
Fr\'echet differentiable at $X\succ0$ (Lemma~\ref{lem:dk} with
$f=x^{p}$ on $I=(0,\infty)$); $\cZ$ is linear; and with
$g(y)=\max\{y,0\}^{\alpha}$ of Lemma~\ref{lem:dk}, continuously
differentiable on $\mathbb{R}$ because $\alpha>1$, the map
$Y\mapsto\tr g(Y)$, which is $\tr[Y^{\alpha}]$ on the positive
semidefinite operators, is Fr\'echet differentiable at every $Y\succeq0$
with derivative $K\mapsto\alpha\tr[Y^{\alpha-1}K]$ by
\eqref{eq:trace-diff}. For $|t|$ small $X+tH\succ0$, so
$T(X+tH)=\tr g\big(\cZ((X+tH)^{p})\big)$. The chain rule gives
\begin{equation*}
  \frac{\mathrm d}{\mathrm dt}T(X+tH)\Big|_0
  =\alpha\tr\big[Y^{\alpha-1}\,\cZ\big(DX^{p}[H]\big)\big]
  =\alpha\tr\big[\cZ\big(Y^{\alpha-1}\big)\,DX^{p}[H]\big]
  =\tr\big[\alpha\,DX^{p}\big[\cZ(Y^{\alpha-1})\big]\,H\big],
\end{equation*}
using self-adjointness of the pinching and then of $DX^{p}$
(Lemma~\ref{lem:dk}). \eqref{eq:gradT-kernel} is \eqref{eq:dk} written
out.

\emph{Step 2 (faces, part (i)).} Write $S=\operatorname{supp}X$ and let
$\Pi_{S}$ be its projector. For Hermitian $H$ with
$\operatorname{supp}H\subseteq S$ we have
$X+tH=\Pi_{S}(X+tH)\Pi_{S}$, and $X$ is positive definite on $S$, so for
$|t|$ small the compression $(X+tH)|_{S}$ is positive definite and
$(X+tH)^{p}=\big((X+tH)|_{S}\big)^{p}\oplus0$ by the convention $0^{p}=0$
of \xref{foundations/def:pinching}. Lemma~\ref{lem:dk}, applied inside
$\Lin{S}$ with $f(x)=x^{p}$ on $I=(0,\infty)$, which contains
$\operatorname{spec}(X|_{S})$, gives
\begin{equation}\label{eq:face-dpow}
  \frac{\mathrm d}{\mathrm dt}(X+tH)^{p}\Big|_{0}
  =\!\!\sum_{i,j:\ \lambda_i,\lambda_j>0}\!\!
  \frac{\lambda_i^{p}-\lambda_j^{p}}{\lambda_i-\lambda_j}\,P_iHP_j
  \eqqcolon D_{S}X^{p}[H],
\end{equation}
a self-adjoint map on the Hermitian operators supported in $S$. No term
with $\lambda_i=0$ or $\lambda_j=0$ occurs, which removes the divergence
of the kernel at the origin; this is (a). The outer factor is untouched by
that compression: with $g(y)=\max\{y,0\}^{\alpha}$ as in Step~1,
\eqref{eq:trace-diff} makes $Y\mapsto\tr g(Y)$, which is $\tr[Y^{\alpha}]$
on the positive semidefinite operators, differentiable at every
$Y\succeq0$, in particular at the \emph{ambient}
$Y=\cZ(X^{p})$, which is not supported in $S$ unless $S$ is
$\cZ$-invariant. This is (b). Composing the two,
\begin{equation*}
\begin{aligned}
  \frac{\mathrm d}{\mathrm dt}T(X+tH)\Big|_{0}
  &=\alpha\tr\big[Y^{\alpha-1}\,\cZ\big(D_{S}X^{p}[H]\big)\big]
  =\alpha\tr\big[\cZ\big(Y^{\alpha-1}\big)\,D_{S}X^{p}[H]\big]\\
  &=\tr\big[\alpha\,D_{S}X^{p}\big[\cZ(Y^{\alpha-1})\big]\,H\big],
\end{aligned}
\end{equation*}
by self-adjointness of the pinching and then of $D_{S}X^{p}$
(Lemma~\ref{lem:dk} inside $\Lin{S}$; only
$\Pi_{S}\cZ(Y^{\alpha-1})\Pi_{S}$ enters, because $H$ is supported in
$S$).

\emph{Step 3 (compression).} Put $S_{\cZ}\coloneqq\sum_{j\in\Zset}Z_jS$.
It contains $S$, because $\sum_jZ_j=\idon{}$, and it is $\cZ$-invariant,
because $Z_iZ_j=\delta_{ij}Z_i$; and for every $W$ supported in $S_{\cZ}$,
$\cZ(W)=\sum_jZ_jWZ_j$ is supported in $\sum_jZ_jS_{\cZ}=S_{\cZ}$. With
$X^{p}$ supported in $S$, therefore, $Y=\cZ(X^{p})$, $Y^{\alpha-1}$ and
$\cZ(Y^{\alpha-1})$ are supported in $S_{\cZ}$. Let
$S_{\cZ}\subseteq S'$, which holds in particular when $S\subseteq S'$ and
$Z_jS'\subseteq S'$ for all $j$. Compressing the pinching to $S'$,
$W\mapsto\Pi_{S'}\cZ(W)\Pi_{S'}$, then changes none of these operators, so
every trace above may be taken inside $S'$ and the ambient complement
contributes $0$: the compressed and ambient computations agree.

If $S$ is not $\cZ$-invariant, compressing to $S$ itself changes $T$
already. Then $S_{\cZ}\ne S$, and $\operatorname{supp}Y=S_{\cZ}$: a vector
$\ket{w}$ with $\bra{w}Z_jX^{p}Z_j\ket{w}=0$ for every $j$ has
$X^{p/2}Z_j\ket{w}=0$, hence $Z_j\ket{w}\perp S$, for every $j$, and
conversely each $Z_jX^{p}Z_j$ has range in $Z_jS$. So $Y$ is not supported
in $S$, and the compressed value $\tr[(\Pi_SY\Pi_S)^{\alpha}]$ is strictly
smaller than $T(X)=\tr[Y^{\alpha}]$: for a unit eigenvector $u\in S$ of
$\Pi_SY\Pi_S$, with eigenvalue $\mu=\bra{u}Y\ket{u}$, Jensen's inequality
for the convex $x\mapsto x^{\alpha}$ and the spectral distribution of $Y$
in $u$ gives $\mu^{\alpha}\le\bra{u}Y^{\alpha}\ket{u}$; summing over an
orthonormal eigenbasis of $S$ gives
$\tr[(\Pi_SY\Pi_S)^{\alpha}]\le\tr[\Pi_SY^{\alpha}]$, the operator
$\Pi_SY\Pi_S$ vanishing on $S^{\perp}$; and
$\tr[\Pi_SY^{\alpha}]<\tr[Y^{\alpha}]$, because $Y^{\alpha}\succeq0$ has the
support of $Y$, which is not contained in $S$. Witness: let
$\ket\eta=a\ket{j}_R\ket{u}+b\ket{j'}_R\ket{u'}$ with $j\ne j'$,
$|a|^{2}+|b|^{2}=1$, and $X=\dyad{\eta}$, so that $X^{p}=X$ and
$\cZ(X^{p})=|a|^{2}\dyad{j,u}+|b|^{2}\dyad{j',u'}$ has rank two while $X$
has rank one. The ambient value is
$T(X)=|a|^{2\alpha}+|b|^{2\alpha}$, whereas compressing to
$S=\mathbb{C}\ket\eta$ would return
$\big(\bra\eta\cZ(X^{p})\ket\eta\big)^{\alpha}
=(|a|^{4}+|b|^{4})^{\alpha}$; at $|a|^{2}=|b|^{2}=\tfrac12$ these are
$2^{1-\alpha}$ and $2^{-\alpha}$.

\emph{Step 4 (branch operators).} For $X=\cG_{\hat c,s}(\omega_s)$,
\eqref{eq:Wcs-bound} gives $S\subseteq W_{\hat c,s}$, with equality at
maximal support, and $W_{\hat c,s}$ need not be $\cZ$-invariant
(Lemma~\ref{lem:branch-support}\,(iv)). The split form (a)--(b) is
therefore the operative one: it needs no $\cZ$-invariance, only
$\operatorname{supp}H\subseteq S$. The whole computation may be
compressed to a $\cZ$-invariant $S'\supseteq W_{\hat c,s}$, for instance
$\sum_{j\in\Zset}Z_jW_{\hat c,s}$, which is $\cZ$-invariant because
$Z_iZ_j=\delta_{ij}Z_i$.

\emph{Step 5 (one-sided derivatives, part (ii)).} $T$ is concave on the
positive semidefinite cone by Lemma~\ref{lem:cone-concavity}\,(b), which
carries the concavity of \xref{foundations/fact:concavity}, in the form of
Corollary~\ref{cor:concavity} on $\Sleq{}$, to the whole cone by the
positive $1$-homogeneity $T(cX)=c\,T(X)$ for $c\ge0$. For a concave
function, $t\mapsto\frac{T(X+tH)-T(X)}{t}$ is non-increasing on the
feasible ray, so the limit exists in $(-\infty,+\infty]$. If $H$
preserves the support, (i) gives a finite value. The scalar example
$x\mapsto x^{p}$ at $x=0$ shows the support-enlarging derivative can be
$+\infty$.
\end{proof}

% claim meat-core/cor:full-gradient
\claimsummary{meat-core/cor:full-gradient}{At a point with positive probabilities and maximal branch supports, the objective of the Choi program is differentiable along every feasible direction, with an explicit gradient. The point need only be positive semidefinite block by block, so the corollary applies at relaxed iterates.}
\claimbookmark{meat-core/cor:full-gradient}{Corollary}{cor:full-gradient}{Gradient of the objective along feasible directions}
\claimleadin{%
The gradient of the pinched trace power on a face,
Theorem~\ref{thm:gradient}\,(i), combines with the affine maps of the Choi
program into the gradient of the whole objective. Of the $J$ factor of the
point the hypothesis asks only per-block positive semidefiniteness, and no
normalization, coupling condition or declared row, which is what lets
Theorem~\ref{thm:firstorder}\,(iiib) apply this corollary at an iterate of
the \termref{term-relaxed-feasible-set}{relaxed feasible set} $\Feas^{J}(\varepsilon)$, the only set a
run establishes membership in (Remark~\ref{rem:fo-iterate-checks}).
Lemma~\ref{lem:fo-khyp}\,(iii) weakens the positivity hypothesis to the
\termref{term-acceptance-support}{acceptance support} $\Chacc$ of \eqref{eq:chacc-def}, and to
$\Chkey\cap\Chacc$ for the $K$-part, which is the form
Theorem~\ref{thm:firstorder} uses and the only one satisfiable when
$\Chacc\subsetneq\Chatset$. Remark~\ref{rem:directional-differentiability}
says why nothing beyond directional differentiability is claimed.
}
\begin{corollary}[Gradient of the objective along feasible directions]\label{cor:full-gradient}
Let $\alpha\in(1,\infty)$ and let $x=(\mathbf q,(J_s)_s)$ be a point with
$\mathbf q\in\acSet$ and $J_s\in\Pos{A_0\Qsys_s}$ for every
$s\in\SecSet$ --- per-block positive semidefiniteness of the $J$ factor
and nothing else: no normalization, no coupling condition
$\sum_s\tr_{\Qsys_s}J_s=\idon{A_0}$ and no declared row $\bowtie_k$ is
assumed of it, so that every point of $\acSet\times\Feas^{J}(\varepsilon)$
qualifies, at every $\varepsilon\ge0$ and the exact $\Feas^{J}$ included
(Definition~\ref{def:relaxed-feasible-set}: per-block positive
semidefiniteness is the one condition that family does not relax) ---
and such that $q(\hat c)>0$, $P^{J}_{\hat c}>0$ and $K^{J}_{\hat c}>0$
for all $\hat c$, and such that each branch operator
$X_{\hat c,s}\coloneqq\cG^{J}_{\hat c,s}(J_s)$
($\hat c\in\Chkey$, $s\in\SecSet_K(\hat c)$) has \emph{maximal support},
$\operatorname{supp}X_{\hat c,s}=W_{\hat c,s}$ of \eqref{eq:Wcs}, for
which $J_s\succ0$ suffices, so that Theorem~\ref{thm:gradient}\,(i)
applies on that face. The displayed formulas are written for
trusted-source mode; the untrusted-source-mode statement is at the end of
this corollary.

Then the objective $O^{J}_{\alpha}$ of \eqref{eq:choi-program} is
differentiable at that point \emph{along every feasible direction}, and
that is the whole of the conclusion: for every
$y\in\acSet\times\Feas^{J}$ the one-variable map
$t'\mapsto O^{J}_{\alpha}\big(x+t'(y-x)\big)$ is differentiable at
$t'=0$, with derivative the pairing of $y-x$ against the gradient
\begin{align}
  \frac{\partial O^{J}_{\alpha}}{\partial q(\hat c)}
  &=\begin{cases}
  \dfrac{\alpha}{\alpha-1}\Big(\log_2\dfrac{q(\hat c)}{P^{J}_{\hat c}}
  +\log_2\mathrm e\Big), & \hat c\notin\Chkey,\\[1.2ex]
  \log_2\dfrac{q(\hat c)}{P^{J}_{\hat c}}+\log_2\mathrm e
  +\dfrac{1}{\alpha-1}\Big(\log_2\dfrac{q(\hat c)}{K^{J}_{\hat c}}
  +\log_2\mathrm e\Big), & \hat c\in\Chkey,
  \end{cases}\label{eq:grad-q}\\[0.6ex]
  \nabla_{J_{s'}}O^{J}_{\alpha}
  &=-\log_2\mathrm e\!\!\sum_{\hat c\notin\Chkey}\!
  \frac{\alpha}{\alpha-1}\,\frac{q(\hat c)}{P^{J}_{\hat c}}\,
  \Gamma^{(s')}_{\hat c}
  \;-\;\log_2\mathrm e\!\!\sum_{\hat c\in\Chkey}\!\left[
  \frac{q(\hat c)}{P^{J}_{\hat c}}\,\Gamma^{(s')}_{\hat c}
  +\frac{1}{\alpha-1}\,\frac{q(\hat c)}{K^{J}_{\hat c}}
  \Big(\Gamma^{(s')}_{\hat c}+\Delta^{(s')}_{\hat c}\Big)\right],
  \label{eq:grad-J}
\end{align}
where
\begin{equation}\label{eq:grad-pieces}
  \Gamma^{(s')}_{\hat c}\coloneqq
  \sum_{(x,v):\,\phi(a(x,v))=\hat c}
  \big(W_{\tau}^{\dagger}M_xW_{\tau}\big)\otimes F^{(s')}_v,
  \qquad
  \Delta^{(s')}_{\hat c}\coloneqq
  \mathbb 1[s'\in\SecSet_K(\hat c)]\;
  \chi_{\Qsys_{s'}}^{\dagger}\Big(
  \cG_{\hat c,s'}^{\dagger}\big[\nabla T(X_{\hat c,s'})
  -\idon{}\big]\Big),
\end{equation}
with $\cG^{\dagger}_{\hat c,s'}(Y)=\sum_{\kappa}K^{(s')\dagger}_{\kappa}
YK^{(s')}_{\kappa}$ the Kraus adjoint and
$\chi^{\dagger}_{\Qsys_{s'}}(A)=(W_{\tau}^{\dagger}\otimes\idon{})A
(W_{\tau}\otimes\idon{})$. Each $\nabla T(X_{\hat c,s'})$ is read in the
split form of Theorem~\ref{thm:gradient}\,(i): divided differences over
the strictly positive spectrum of $X_{\hat c,s'}$, the pinched operators
$\cZ(X^{p}_{\hat c,s'})$ and $\cZ(Y^{\alpha-1})$ on the ambient branch
space.

\emph{Untrusted-source mode.} There the \emph{state form} is primary,
the Choi form being its pullback rather than the other way round: the
variables are the sector blocks $\omega_{s}$ themselves
(Definition~\ref{def:untrusted-choi}), $\chi_{\Qsys_{s'}}=\idmap$, and
\eqref{eq:grad-q}--\eqref{eq:grad-J} hold verbatim with
$J_{s'}\to\omega_{s'}$, $P^{J}_{\hat c}\to P_{\hat c}$,
$K^{J}_{\hat c}\to K_{\hat c}$, $X_{\hat c,s}=\cG_{\hat c,s}(\omega_s)$
and the pieces \eqref{eq:grad-pieces} losing their $W_{\tau}$
conjugation:
\begin{equation}\label{eq:grad-pieces-un}
  \Gamma^{(s')}_{\hat c}=
  \sum_{(x,v):\,\phi(a(x,v))=\hat c} M_x\otimes F^{(s')}_v,
  \qquad
  \Delta^{(s')}_{\hat c}=
  \mathbb 1[s'\in\SecSet_K(\hat c)]\;
  \cG_{\hat c,s'}^{\dagger}\big[\nabla T(X_{\hat c,s'})-\idon{}\big],
\end{equation}
Hermitian operators on $A\Qsys_{s'}$. The derivation is the proof
below with $\chi_{\Qsys_{s'}}$ deleted; the positivity hypothesis on the
point reads $\omega_{s}\in\Pos{A\Qsys_s}$ per block, there being no $A_0$
in that mode; and the maximal-support hypothesis
reads $\operatorname{supp}X_{\hat c,s'}=W_{\hat c,s'}$ with $\Sadm$
now all of $\Hs{A}$, for which $\omega_{s'}\succ0$ suffices. In
trusted-source mode the same formulas conjugated by $\chi$ give the
state-form gradients, and at $d_A=1$ the two agree because
$W_{\tau}=[1]$.
\end{corollary}
\claimmeta{\metalabel{Uses}\mref{Def.}{def:maps}, \mref{Thm.}{thm:choi}, \mref{Def.}{def:untrusted-choi}, \mref{Def.}{def:branch-support}, \mref{Lem.}{lem:branch-support}, \mref{Thm.}{thm:gradient}
\metasep \metalabel{Used by}\mref{Lem.}{lem:kraus-freedom}, \mref{Rem.}{rem:directional-differentiability}, \mref{Lem.}{lem:fo-khyp}, \mref{Lem.}{lem:fo-face}, \mref{Thm.}{thm:firstorder}, \mref{Cor.}{cor:fo-dropped}, \mref{Lem.}{lem:certified-bracket}, \mref{Lem.}{lem:certified-bracket-enclosed}, \mref{Rem.}{rem:fo-oracle-repair}, \mref{Req.}{req:certified-evaluation}
\metasep newly derived
\metasep implemented (effective: accepted)
\metasep \metalabel{Code}\coderef{https://github.com/fraunhofersg/certiqs-key/blob/v0.1.0/src/certiqs_key/meat_core/objective/gradient.py}{certiqs_key.meat_core.objective.gradient:choi_gradient}
\metasep \metalabel{Oracles}\coderef{https://github.com/fraunhofersg/certiqs-key/blob/v0.1.0/tests/oracles/registry.py}{registry.py}}

\begin{proof}
\emph{Idea.} The objective is differentiated term by term, through affine
maps and, for the deficits, through the chain rule of
Theorem~\ref{thm:gradient}. Of the point $x$ the argument consumes three
things. (a) The strict positivity clause, which keeps every logarithm and
every denominator of \eqref{eq:grad-q}--\eqref{eq:grad-J} away from its
singularity. (b) Maximal support of each key-active branch operator
$X_{\hat c,s}$, which makes Theorem~\ref{thm:gradient}\,(i) applicable on
that face; the ambient factors $\cZ(X^{p}_{\hat c,s})$ and
$\cZ(Y^{\alpha-1})$ it forms need only $X_{\hat c,s}\succeq0$, which
$J_s\succeq0$ and complete positivity of $\cG^{J}_{\hat c,s}$ supply.
(c) That the $J$-component of the direction $y-x$ is Hermitian, which
per-block positive semidefiniteness of $x$ and of $y$ supplies, a
difference of two positive semidefinite operators being Hermitian; this
is the hypothesis Lemma~\ref{lem:fo-face} is stated on, and it names no
\exttermref{../protocol-compiler/contract.pdf}{term-feasible-set}{feasible set}, exact or relaxed. Normalization, the coupling condition and the declared rows
$\bowtie_k$ are consumed nowhere: not in \eqref{eq:grad-q}, not in
\eqref{eq:grad-J}, not in \eqref{eq:grad-pieces}, and at no step below,
which is why the hypothesis asks for per-block positive semidefiniteness
and not for feasibility.

\emph{Step 1 (the $\mathbf q$-gradient).} Differentiate
\eqref{eq:choi-program} term by term. In $\mathbf q$:
$\frac{\mathrm d}{\mathrm du}\,u\log_2\frac uv
=\log_2\frac uv+\log_2\mathrm e$.

\emph{Step 2 (the divergence terms in $J$).}
$\frac{\partial}{\partial v}\,u\log_2\frac uv=-\frac{u}{v}\log_2\mathrm
e$ composed with the affine maps $J\mapsto P^{J}_{\hat c}$, whose
gradients are the constant Hermitian operators
$\Gamma^{(s')}_{\hat c}$ (expand \eqref{eq:Phimap} through
$\chi$ as in the proof of Theorem~\ref{thm:choi} and read off the linear
coefficient of $J_{s'}$).

\emph{Step 3 (the $K$-terms).}
$\nabla_{J_{s'}}K^{J}_{\hat c}
=\Gamma^{(s')}_{\hat c}+\Delta^{(s')}_{\hat c}$: the $P$ part is as
before, and the deficit part is the chain rule through the linear maps
$\chi_{\Qsys_{s'}}$ and $\cG_{\hat c,s'}$ and
Theorem~\ref{thm:gradient}, together with
$\nabla_{X}\tr[X]=\idon{}$; the adjoints implement
$\tr[Y\,\mathcal L(H)]=\tr[\mathcal L^{\dagger}(Y)H]$ for each linear
factor.

\emph{Step 4 (the chain rule at a singular branch operator).} The
maximal-support hypothesis is what licenses that chain rule at a singular
$X_{\hat c,s'}$: an admissible direction
$\Delta J_{s'}=J'_{s'}-J_{s'}$ is Hermitian, both endpoints being
per-block positive semidefinite, and it produces
$H=\cG^{J}_{\hat c,s'}(\Delta J_{s'})$ with
$\operatorname{supp}H\subseteq W_{\hat c,s'}
=\operatorname{supp}X_{\hat c,s'}$ by \eqref{eq:Wcs-bound} applied to the
positive and negative parts of $\Delta J_{s'}$ separately (the fact the
submodule states as Lemma~\ref{lem:fo-face}, derived here and reading no
feasible set), so Theorem~\ref{thm:gradient}\,(i) applies; and the
components of $\nabla T(X_{\hat c,s'})$ outside that face never enter
\eqref{eq:grad-pieces}, because
$\chi^{\dagger}_{\Qsys_{s'}}\circ\cG^{\dagger}_{\hat c,s'}$ sandwiches its
argument between the operators
$K^{(s')}_{\kappa}(W_{\tau}\otimes\idon{})$ and their adjoints, whose
ranges lie in $W_{\hat c,s'}$ by \eqref{eq:Wcs}. All compositions are
differentiable at the stated point, so the sum rule applies.
\end{proof}

% claim meat-core/lem:kraus-freedom
\claimsummary{meat-core/lem:kraus-freedom}{Kraus factors with the same Gram matrices give branch operators related by a partial isometry commuting with the key pinching, so the trace power, deficit and pulled-back gradient agree. A subspace containing the support is the support once the compression to it is positive definite.}
\claimbookmark{meat-core/lem:kraus-freedom}{Lemma}{lem:kraus-freedom}{Kraus freedom in the branch maps}
\claimleadin{%
The Kraus operators of a \termref{term-branch-map}{branch map} carry the square roots
$\sqrt{M_x}\otimes\sqrt{F_v}$ (Definition~\ref{def:branch}), and in the Choi
form the frame $W_{\tau}$ as well. A square root of an exact positive
semidefinite matrix is irrational in general, and is computed only
approximately. The lemma shows that nothing the objective or its gradient
reads depends on that choice: any factor with the same Gram matrix gives the
same branch values and the same pulled-back gradient, through a partial
isometry that the key pinching cannot see. The lemma is stated for one pair
of factor families at a time, so it holds for every element an enclosure
leaves possible, each with its own factors. A certified evaluation may
therefore build the branch maps from factors whose rank is fixed: exact
factors, or, for an element $LWL^{\dagger}$ known through an enclosed factor
$L$ and an exact positive diagonal core $W$, the factor $W^{1/2}L^{\dagger}$,
whose Gram matrix is that element for every $L$ the enclosure allows, though
$L$ itself is not known (Requirement~\ref{req:certified-evaluation}).
Part~(iv) then
compresses a branch operator, or a pinched block of one, to its support,
where it is positive definite and the divided-difference map of
\xref{foundations/lem:matrix-function-perturbation} applies: a subspace known
to contain the support, such as the range of the factors, is the support as
soon as the compression to it is shown positive definite.
}
\begin{lemma}[Kraus freedom in the branch maps]\label{lem:kraus-freedom}
Let $\alpha\in(1,\infty)$, $p\coloneqq1/\alpha$ and $\hat c\in\Chkey$, and let
$\mathcal{H}_{\mathrm{in}}$ be a finite-dimensional Hilbert space. For every
outcome pair $(x,v)$ of a \termref{term-pass-class}{pass class} of $\hat c$, that is
$a(x,v)\in\Kann_{\mathrm{pass}}(\hat c)$, let $S_{xv}$ and $S'_{xv}$ be linear
maps from $\mathcal{H}_{\mathrm{in}}$ into finite-dimensional Hilbert spaces
$\mathcal{K}$ and $\mathcal{K}'$ with
\begin{equation}\label{eq:kraus-gram}
  S'^{\dagger}_{xv}S'_{xv}\;=\;S^{\dagger}_{xv}S_{xv} .
\end{equation}
For $\kappa\in\Kann_{\mathrm{pass}}(\hat c)$ put
\begin{equation}\label{eq:kraus-factored}
  K_{\kappa}\coloneqq\sum_{(x,v)\in a^{-1}(\kappa)}\ket{r(x,v)}_R\ket{x,v}_{XY}
  \ket{\kappa}_{C_AC_B}\otimes S_{xv},
\end{equation}
a map into $\mathcal{B}\coloneqq\Hs{R\,XY\,C_AC_B}\otimes\mathcal{K}$, and
$K'_{\kappa}$ likewise with $S'_{xv}$, into
$\mathcal{B}'\coloneqq\Hs{R\,XY\,C_AC_B}\otimes\mathcal{K}'$. For every
positive semidefinite operator $\omega$ on $\mathcal{H}_{\mathrm{in}}$ put
$X(\omega)\coloneqq\sum_{\kappa}K_{\kappa}\omega K_{\kappa}^{\dagger}$ and
$X'(\omega)\coloneqq\sum_{\kappa}K'_{\kappa}\omega K'^{\dagger}_{\kappa}$,
the sums over $\Kann_{\mathrm{pass}}(\hat c)$. Let $\cZ$ and $\cZ'$ be the
\termref{term-key-pinching}{key pinchings} on $\mathcal{B}$ and $\mathcal{B}'$, by
the projectors $Z_j=\dyad{j}_R\otimes\id$, $j\in\Zset$, and $T$ and $T'$ the
functionals \eqref{eq:Tcone} with these pinchings. For $X\succeq0$ write
$\nabla T(X)$ for the operator of \eqref{eq:gradT-kernel} with the sum over
the strictly positive eigenvalues of $X$ only, $Y=\cZ(X^{p})$ and
$0^{\alpha-1}\coloneqq0$, which is the split form of
Theorem~\ref{thm:gradient}\,(i), and $\nabla T'$ likewise.
\begin{enumerate}[label=(\roman*),leftmargin=2.4em]
\item \emph{(A partial isometry.)} Let $\mathcal{W}\subseteq\mathcal{B}$ be
  the span of the subspaces
  $\ket{r(x,v)}_R\ket{x,v}_{XY}\ket{a(x,v)}_{C_AC_B}\otimes
  \operatorname{ran}S_{xv}$. There is a linear map
  $\mathbf U:\mathcal{B}\to\mathcal{B}'$ that is isometric on $\mathcal{W}$
  and vanishes on $\mathcal{W}^{\perp}$, with $\mathbf UK_{\kappa}=K'_{\kappa}$
  for every $\kappa$ and $(Z_j\otimes\idon{\mathcal{K}'})\,\mathbf U
  =\mathbf U\,(Z_j\otimes\idon{\mathcal{K}})$ for every $j$; and
  $\mathcal{W}$ is invariant under every $Z_j\otimes\idon{\mathcal{K}}$.
  Hence $X'(\omega)=\mathbf UX(\omega)\mathbf U^{\dagger}$, with
  $\operatorname{supp}X(\omega)\subseteq\mathcal{W}$, for every
  $\omega\succeq0$.
\item \emph{(Values.)} For every $\omega\succeq0$:
  $T'(X'(\omega))=T(X(\omega))$, $\tr X'(\omega)=\tr X(\omega)$,
  $\operatorname{supp}X'(\omega)=\mathbf U\operatorname{supp}X(\omega)$ and
  $\operatorname{rank}X'(\omega)=\operatorname{rank}X(\omega)$.
\item \emph{(Gradients.)} For every $\omega\succeq0$, with $X\coloneqq
  X(\omega)$ and $X'\coloneqq X'(\omega)$,
  \begin{equation}\label{eq:kraus-gradient}
    \nabla T'(X')=\mathbf U\,\nabla T(X)\,\mathbf U^{\dagger},
    \qquad
    \sum_{\kappa}K'^{\dagger}_{\kappa}\big(\nabla T'(X')-\id\big)K'_{\kappa}
    =\sum_{\kappa}K^{\dagger}_{\kappa}\big(\nabla T(X)-\id\big)K_{\kappa}.
  \end{equation}
\item \emph{(Compression to the support.)} Let $A\succeq0$ be an operator
  on $\mathcal{B}$.
  \begin{enumerate}[label=(\alph*),leftmargin=2.2em]
  \item \emph{(Containment.)} For every $\omega\succeq0$,
    $\operatorname{supp}X(\omega)\subseteq\sum_{\kappa}\operatorname{ran}K_{\kappa}$;
    and $\operatorname{supp}\big((Z_j\otimes\idon{\mathcal{K}})A
    (Z_j\otimes\idon{\mathcal{K}})\big)\subseteq(Z_j\otimes\idon{\mathcal{K}})
    \operatorname{supp}A$ for every $j$; and
    $\operatorname{supp}A^{q}=\operatorname{supp}A$ for every $q>0$.
  \item \emph{(A floor proves equality.)} If $Q:\mathbb{C}^{m}\to\mathcal{B}$
    is an isometry with $\operatorname{ran}Q\supseteq\operatorname{supp}A$
    and $Q^{\dagger}AQ$ is positive definite, then
    $\operatorname{ran}Q=\operatorname{supp}A$.
  \item \emph{(The compression.)} If $Q:\mathbb{C}^{m}\to\mathcal{B}$ is an
    isometry with $\operatorname{ran}Q=\operatorname{supp}A$ and $m\ge1$, then
    $N\coloneqq Q^{\dagger}AQ$ is positive definite, $A=QNQ^{\dagger}$,
    $A^{q}=QN^{q}Q^{\dagger}$ for every $q>0$, and for every operator $H$ on
    $\mathcal{B}$ the divided-difference sum of \eqref{eq:face-dpow} at $A$,
    over the strictly positive eigenvalues of $A$, equals
    $Q\,\mathrm{D}_N[Q^{\dagger}HQ]\,Q^{\dagger}$, with $\mathrm{D}_N$ the
    map of \xref{foundations/lem:matrix-function-perturbation}\,(iii) for
    $f(x)=x^{p}$.
  \end{enumerate}
\item \emph{(The branch maps.)} With $\mathcal{H}_{\mathrm{in}}=\Hs{A\Qsys_s}$
  and $S_{xv}=\sqrt{M_x}\otimes\sqrt{F^{(s)}_v}$, the $K_{\kappa}$ are the
  $K^{(s)}_{\kappa}$ of Definition~\ref{def:branch} and
  $X(\omega)=\cG_{\hat c,s}(\omega)$. In \exttermref{../transmitter/contract.pdf}{term-trusted-source-mode}{trusted-source mode}, with
  $\mathcal{H}_{\mathrm{in}}=\Hs{A_0\Qsys_s}$, a support frame and
  $S_{xv}=\sqrt{M_x}\,W_{\tau}\otimes\sqrt{F^{(s)}_v}$, the $K_{\kappa}$ are
  the $K^{(s)}_{\kappa}(W_{\tau}\otimes\idon{\Qsys_s})$,
  $X(J_s)=\cG^{J}_{\hat c,s}(J_s)$, and
  $\sum_{\kappa}K^{\dagger}_{\kappa}(\nabla T(X)-\id)K_{\kappa}
  =\chi^{\dagger}_{\Qsys_s}\big(\cG^{\dagger}_{\hat c,s}[\nabla T(X)-\id]\big)$,
  the operator of \eqref{eq:grad-pieces}; the same holds on a
  \exttermref{../protocol-compiler/contract.pdf}{term-joint-sector}{joint sector}
  $(t,s)$ with $M^{(t)}_x$ and the frame of the Alice sector $t$
  (Definition~\ref{def:joint-sectors}). In each case
  $S^{\dagger}_{xv}S_{xv}$ is $M_x\otimes F^{(s)}_v$, respectively
  $(W_{\tau}^{\dagger}M_xW_{\tau})\otimes F^{(s)}_v$. So for any $S'_{xv}$
  with that Gram matrix, the deficit $f_{\hat c,s}$, the branch operator's
  rank and the operator $\Delta^{(s)}_{\hat c}$ of \eqref{eq:grad-pieces}
  are given, by (ii) and (iii), by $X'$ in place of $X$.
\end{enumerate}
\end{lemma}
\claimmeta{\metalabel{Uses}\mref{Def.}{def:maps}, \mref{Def.}{def:branch}, \mref{Thm.}{thm:branch-formula}, \mref{Lem.}{lem:cone-concavity}, \mref{Def.}{def:joint-sectors}, \mref{Def.}{def:support-frame}, \mref{Thm.}{thm:choi}, \mref{Def.}{def:untrusted-choi}, \mref{Thm.}{thm:gradient}, \mref{Cor.}{cor:full-gradient}, \extref{../foundations/contract.pdf}{foundations/def:pinching}, \extref{../foundations/contract.pdf}{foundations/lem:matrix-function-perturbation}
\metasep \metalabel{Used by}\mref{Req.}{req:certified-evaluation}
\metasep newly derived
\metasep implemented (effective: accepted)
\metasep \metalabel{Code}\coderef{https://github.com/fraunhofersg/certiqs-key/blob/v0.1.0/src/certiqs_key/solvers/frank_wolfe/certified.py}{certiqs_key.solvers.frank_wolfe.certified:_deficit}}

\begin{proof}
\emph{Isometric conjugation.} Let $\mathbf V$ be isometric on a subspace
$\mathcal{S}$ and vanish on $\mathcal{S}^{\perp}$, or be an isometry and
$\mathcal{S}$ its whole domain. Then $\mathbf V^{\dagger}\mathbf V$ is the
orthogonal projector onto $\mathcal{S}$: for $\phi$, $\psi$ with components
$\phi_{\mathcal{S}}$, $\psi_{\mathcal{S}}$ in $\mathcal{S}$,
$\langle\phi,\mathbf V^{\dagger}\mathbf V\psi\rangle
=\langle\mathbf V\phi_{\mathcal{S}},\mathbf V\psi_{\mathcal{S}}\rangle
=\langle\phi_{\mathcal{S}},\psi_{\mathcal{S}}\rangle$, an isometry preserving
inner products by polarisation. If $A\succeq0$ has
$\operatorname{supp}A\subseteq\mathcal{S}$,
write $A=\sum_ka_ku_ku_k^{\dagger}$ with $a_k>0$ and orthonormal
$u_k\in\mathcal{S}$; the $\mathbf Vu_k$ are orthonormal, so
$\mathbf VA\mathbf V^{\dagger}=\sum_ka_k(\mathbf Vu_k)(\mathbf Vu_k)^{\dagger}$
is a spectral decomposition with the same positive eigenvalues, its spectral
projector for an eigenvalue $\lambda>0$ is $\mathbf VP\mathbf V^{\dagger}$
with $P$ that of $A$, its support is $\mathbf V\operatorname{supp}A$, and
$(\mathbf VA\mathbf V^{\dagger})^{q}=\mathbf VA^{q}\mathbf V^{\dagger}$ for
every $q>0$, with $0^{q}\coloneqq0$ as in \xref{foundations/def:pinching};
moreover $\tr[\mathbf VA\mathbf V^{\dagger}]=\tr[A\,\mathbf V^{\dagger}
\mathbf V]=\tr A$, since $\mathbf V^{\dagger}\mathbf V$ is the identity
on $\mathcal{S}\supseteq\operatorname{supp}A$.

\emph{(i).} For each pair $(x,v)$ define $V_{xv}$ on
$\operatorname{ran}S_{xv}$ by $V_{xv}(S_{xv}\phi)\coloneqq S'_{xv}\phi$.
By \eqref{eq:kraus-gram},
$\|S_{xv}\phi\|^{2}=\langle\phi,S^{\dagger}_{xv}S_{xv}\phi\rangle
=\|S'_{xv}\phi\|^{2}$, so $S_{xv}\phi=S_{xv}\phi'$ forces
$S'_{xv}(\phi-\phi')=0$: $V_{xv}$ is well defined, linear and isometric.
Extend it by $0$ on $(\operatorname{ran}S_{xv})^{\perp}$. Write
$e_{xv}\coloneqq\ket{r(x,v)}_R\ket{x,v}_{XY}\ket{a(x,v)}_{C_AC_B}$; these
unit vectors are mutually orthogonal for distinct $(x,v)$, their $XY$
components being. Put $\mathbf U\coloneqq\sum_{(x,v)}e_{xv}e_{xv}^{\dagger}
\otimes V_{xv}$. On $\mathcal{W}=\bigoplus_{(x,v)}e_{xv}\otimes
\operatorname{ran}S_{xv}$ it acts as $V_{xv}$ on each orthogonal summand,
with orthogonal images, so it is isometric there, and it vanishes on
$\mathcal{W}^{\perp}$. Then $\mathbf UK_{\kappa}
=\sum_{(x,v)\in a^{-1}(\kappa)}e_{xv}\otimes V_{xv}S_{xv}=K'_{\kappa}$. Since
$Z_je_{xv}e_{xv}^{\dagger}=e_{xv}e_{xv}^{\dagger}Z_j$, both being
$e_{xv}e_{xv}^{\dagger}$ if $r(x,v)=j$ and $0$ otherwise, $\mathbf U$
commutes with the pinching projectors as stated, and each $Z_j$ maps every
summand of $\mathcal{W}$ to itself or to $0$. Finally
$X'(\omega)=\sum_{\kappa}\mathbf UK_{\kappa}\omega K_{\kappa}^{\dagger}
\mathbf U^{\dagger}=\mathbf UX(\omega)\mathbf U^{\dagger}$, and
$\operatorname{supp}X(\omega)\subseteq\sum_{\kappa}\operatorname{ran}K_{\kappa}
\subseteq\mathcal{W}$.

\emph{(ii).} Put $X\coloneqq X(\omega)$. By isometric conjugation,
$X'^{p}=\mathbf UX^{p}\mathbf U^{\dagger}$, $\tr X'=\tr X$, and the support
and rank statements hold. By (i),
$\cZ'(\mathbf UA\mathbf U^{\dagger})=\sum_j(Z_j\otimes\id)\mathbf UA
\mathbf U^{\dagger}(Z_j\otimes\id)=\mathbf U\cZ(A)\mathbf U^{\dagger}$ for
every $A$, the adjoint of the commutation relation giving
$\mathbf U^{\dagger}(Z_j\otimes\id)=(Z_j\otimes\id)\mathbf U^{\dagger}$.
So $Y'\coloneqq\cZ'(X'^{p})=\mathbf UY\mathbf U^{\dagger}$ with
$Y\coloneqq\cZ(X^{p})\succeq0$, and $\operatorname{supp}Y\subseteq\mathcal{W}$,
$\operatorname{supp}X^{p}$ lying in $\mathcal{W}$ and $\mathcal{W}$ being
invariant under every $Z_j$. Isometric conjugation again gives
$T'(X')=\tr[(\mathbf UY\mathbf U^{\dagger})^{\alpha}]
=\tr[\mathbf UY^{\alpha}\mathbf U^{\dagger}]=\tr[Y^{\alpha}]=T(X)$.

\emph{(iii).} With $Y$ as in (ii), $Y'^{\alpha-1}=\mathbf UY^{\alpha-1}
\mathbf U^{\dagger}$ and so $\cZ'(Y'^{\alpha-1})=\mathbf U\cZ(Y^{\alpha-1})
\mathbf U^{\dagger}$. Let $X=\sum_i\lambda_iP_i$ over the distinct strictly
positive eigenvalues of $X$; the strictly positive eigenvalues of $X'$ are the
same, with spectral projectors $\mathbf UP_i\mathbf U^{\dagger}$. As
$\operatorname{ran}P_i\subseteq\mathcal{W}$, on which
$\mathbf U^{\dagger}\mathbf U$ acts as the identity,
$\mathbf UP_i\mathbf U^{\dagger}\,\mathbf U\cZ(Y^{\alpha-1})\mathbf U^{\dagger}
\,\mathbf UP_j\mathbf U^{\dagger}=\mathbf UP_i\cZ(Y^{\alpha-1})P_j
\mathbf U^{\dagger}$, and summing with the coefficients
$\alpha\,(\lambda_i^{p}-\lambda_j^{p})/(\lambda_i-\lambda_j)$, respectively
$\alpha p\lambda_i^{p-1}$, gives the first identity of
\eqref{eq:kraus-gradient}. For the second,
$K'^{\dagger}_{\kappa}AK'_{\kappa}=K_{\kappa}^{\dagger}\mathbf U^{\dagger}A
\mathbf UK_{\kappa}$; with $A=\mathbf U\nabla T(X)\mathbf U^{\dagger}$ and
with $A=\id$, and $\mathbf U^{\dagger}\mathbf U$ acting as the identity on
$\mathcal{W}\supseteq\operatorname{ran}K_{\kappa}\cup\operatorname{supp}\nabla T(X)$
(the latter inside $\operatorname{supp}X$), both terms reduce to those with
$K_{\kappa}$ and $\nabla T(X)$.

\emph{(iv)(a).} The support of a positive semidefinite operator is its
range. $\operatorname{ran}X(\omega)\subseteq\sum_{\kappa}
\operatorname{ran}(K_{\kappa}\omega K_{\kappa}^{\dagger})\subseteq
\sum_{\kappa}\operatorname{ran}K_{\kappa}$. With
$Z\coloneqq Z_j\otimes\idon{\mathcal{K}}$, $\operatorname{ran}(ZAZ)\subseteq
Z\operatorname{ran}(A)$, since $ZAZw=Z(AZw)$. With
$A=\sum_ka_ku_ku_k^{\dagger}$, $a_k>0$ and orthonormal $u_k$, both $A$ and
$A^{q}=\sum_ka_k^{q}u_ku_k^{\dagger}$ have range the span of the $u_k$.

\emph{(iv)(b).} For $y\ne0$,
$(Qy)^{\dagger}A(Qy)=y^{\dagger}(Q^{\dagger}AQ)y>0$, so $Qy\notin\ker A$.
Hence $\operatorname{ran}Q\cap\ker A=\{0\}$ and
$m=\dim\operatorname{ran}Q\le\dim\mathcal{B}-\dim\ker A
=\operatorname{rank}A$. The inclusion
$\operatorname{supp}A\subseteq\operatorname{ran}Q$ gives
$\operatorname{rank}A\le m$, so the two subspaces have the same dimension and
are equal.

\emph{(iv)(c).} $QQ^{\dagger}$ is the orthogonal projector onto
$\operatorname{supp}A$, so $A=QQ^{\dagger}AQQ^{\dagger}=QNQ^{\dagger}$. For
$y\ne0$, $Qy\ne0$ lies in $\operatorname{supp}A$, where $A$ is positive
definite, so $y^{\dagger}Ny=(Qy)^{\dagger}A(Qy)>0$. Isometric conjugation
with $\mathbf V=Q$ and $N$ in the role of $A$ there gives
$A^{q}=QN^{q}Q^{\dagger}$ and shows that
the strictly positive eigenvalues of $A$ are the eigenvalues $\mu_a$ of $N$,
with spectral projectors $Q\Pi_aQ^{\dagger}$; so the sum of
\eqref{eq:face-dpow} is $\sum_{a,b}f^{[1]}(\mu_a,\mu_b)\,Q\Pi_aQ^{\dagger}H
Q\Pi_bQ^{\dagger}=Q\,\mathrm{D}_N[Q^{\dagger}HQ]\,Q^{\dagger}$ by the
definition of $\mathrm{D}_N$ in
\xref{foundations/lem:matrix-function-perturbation}\,(iii).

\emph{(v).} Substituting $S_{xv}$ in the definition of $K_{\kappa}$ gives
\eqref{eq:kraus} of Definition~\ref{def:branch}, and $X(\omega)$ is then
\eqref{eq:GTheta} for the sector. In the Choi form,
$(\sqrt{M_x}\otimes\sqrt{F^{(s)}_v})(W_{\tau}\otimes\id)
=\sqrt{M_x}W_{\tau}\otimes\sqrt{F^{(s)}_v}$, so the $K_{\kappa}$ are the
$K^{(s)}_{\kappa}(W_{\tau}\otimes\id)$, and $X(J_s)=\cG_{\hat c,s}
(\chi_{\Qsys_s}(J_s))=\cG^{J}_{\hat c,s}(J_s)$ by the definition of
$\chi_{\Qsys_s}$ in Theorem~\ref{thm:choi}. The Kraus adjoint
$\cG^{\dagger}_{\hat c,s}(A)=\sum_{\kappa}K^{(s)\dagger}_{\kappa}AK^{(s)}_{\kappa}$
and $\chi^{\dagger}_{\Qsys_s}(B)=(W_{\tau}^{\dagger}\otimes\id)B
(W_{\tau}\otimes\id)$ (Corollary~\ref{cor:full-gradient}) compose to
$\sum_{\kappa}K_{\kappa}^{\dagger}AK_{\kappa}$. On a joint sector the Kraus
operators act through $\sqrt{M^{(t)}_x}\otimes\sqrt{F^{(s)}_v}$
(\eqref{eq:joint-deficit}), and the same substitution applies. The Gram
matrices are $(\sqrt{M_x}\otimes\sqrt{F^{(s)}_v})^{2}=M_x\otimes F^{(s)}_v$
and $(\sqrt{M_x}W_{\tau}\otimes\sqrt{F^{(s)}_v})^{\dagger}
(\sqrt{M_x}W_{\tau}\otimes\sqrt{F^{(s)}_v})
=W_{\tau}^{\dagger}M_xW_{\tau}\otimes F^{(s)}_v$. The deficit
$f_{\hat c,s}=T\circ\cG_{\hat c,s}-\tr\circ\,\cG_{\hat c,s}$
(\eqref{eq:f-def}) is unchanged by (ii), and $\Delta^{(s)}_{\hat c}$, the
pulled-back operator of \eqref{eq:grad-pieces} with $\nabla T$ read in the
split form, is unchanged by (iii).
\end{proof}

% claim meat-core/rem:directional-differentiability
\claimsummary{meat-core/rem:directional-differentiability}{The objective of the Choi program is undefined off the positive cone and the sets \hyperlink{claim-meat-core/cor:full-gradient}{Corollary~\ref*{cor:full-gradient}} quantifies over have empty interior, so only directional derivatives are well posed. Even within the affine hull, a direction off the acceptance support can have one-sided derivative \ensuremath{-}\ensuremath{\infty}.}
\claimbookmark{meat-core/rem:directional-differentiability}{Remark}{rem:directional-differentiability}{Scope of the directional differentiability of the objective}
\begin{remark}[Scope of the directional differentiability of the objective]\label{rem:directional-differentiability}
Corollary~\ref{cor:full-gradient} claims no ambient Fr\'echet
differentiability, and the set on which it differentiates is part of its
statement. There are two grounds, and the first is decisive on its own.

\emph{First}, $O^{J}_{\alpha}$ is not defined on an ambient neighbourhood
of the corollary's point $x$ at all, for either of two independent causes: the objective
factors through $X\mapsto X^{p}$, which is undefined off the positive
semidefinite cone, so an ambient neighbourhood of a singular $J_s$ leaves
the domain; and the sets the corollary quantifies over have empty
interior in the ambient real vector space: the affine equality
$\sum_{\hat c}q(\hat c)=1$ that $\acSet$ carries forces that on its own,
and on the exact $\Feas^{J}$ the coupling equality
$\sum_s\tr_{\Qsys_s}J_s=\idon{A_0}$ forces it again. So the ambient notion
is not well posed here, whereas the directional one is, and it is the one
\eqref{eq:fo-subgrad} consumes.

\emph{Second}, the restriction cannot be lifted even to directions inside
that affine hull. Where the hypothesis is taken in the weakened form of
Lemma~\ref{lem:fo-khyp}\,(iii), a coordinate $q(\hat c)$ with
$\hat c\notin\Chacc$ sits at $0$ throughout $\acSet$, and
$\partial_u\big(u\log_2\tfrac uv\big)=\log_2\tfrac uv+\log_2\mathrm e
\to-\infty$ as $u\downarrow0$ for $v>0$, so a direction that raises such a
coordinate off $0$, one that leaves $\acSet$ while preserving
$\sum_{\hat c}q(\hat c)=1$, has one-sided derivative $-\infty$; inside
$\acSet$ the coordinate stays $0$ identically and the term is deleted
rather than bounded (Lemma~\ref{lem:fo-khyp}\,(ii)). Feasible directions
are all that Step~1 of Theorem~\ref{thm:firstorder} consumes, and no
statement of this module reads more.
\end{remark}
\claimmeta{\metalabel{Uses}\mref{Cor.}{cor:full-gradient}
\metasep newly derived
\metasep accepted}

% claim meat-core/lem:b-derivative
\claimsummary{meat-core/lem:b-derivative}{At R\'enyi order one the derivative of the pinched trace power of a subnormalised operator is minus ln 2 times its relative entropy to its pinching. A jointly continuous majorant bounds the derivative just above order one, which carries the limit of order one through the infimum.}
\claimbookmark{meat-core/lem:b-derivative}{Lemma}{lem:b-derivative}{Derivative of the pinched trace power at order one}
\claimleadin{%
The limit of order one, Theorem~\ref{thm:oracle}, needs the derivative of
the pinched trace power at $\alpha=1$ and a bound on it just above $1$
that is continuous in the operator. In the proof, $\Xi(\alpha)$ is the
derivative of $\sigma^{1/\alpha}$; the letter $\Sigma$ is taken by the
marginal-constrained range of Fact~\ref{fact:meat}.
}
\begin{technicalclaim}
\begin{lemma}[Derivative of the pinched trace power at order one]\label{lem:b-derivative}
We write $\hlow(\alpha)$ for \eqref{eq:hlow-def} with its
$\alpha$-dependence explicit, and, for $\sigma\in\Pos{}$ with
$\tr\sigma\le1$,
$B_{\sigma}(\alpha)\coloneqq\tr[(\cZ(\sigma^{1/\alpha}))^{\alpha}]$.
Let $\delta>0$, one number for all $\sigma$, and let $\sigma\succeq0$ with
$\tr\sigma\le1$ (so $\|\sigma\|\le1$). Then
$B_{\sigma}\in C^{1}([1,1+\delta])$, with
$B_{\sigma}(1)=\tr\sigma$ and
\begin{equation}\label{eq:Bprime}
  B_{\sigma}'(1)=\tr[\cZ(\sigma)\ln\cZ(\sigma)]-\tr[\sigma\ln\sigma]
  =-\ln2\;\Dtwo\big(\sigma\,\big\|\,\cZ(\sigma)\big)
  \;\le\;0,
\end{equation}
with the convention $0\ln0=0$. Moreover, for all
$\alpha\in(1,1+\delta]$,
\begin{equation}\label{eq:majorant}
  B_{\sigma}'(\alpha)\;\le\;c_{\sigma}(\alpha)\coloneqq
  \tr\big[A(\alpha)^{\alpha}\ln A(\alpha)\big]
  -\frac1{\alpha}\sum_i\lambda_i^{1/\alpha}\ln\lambda_i ,
  \qquad A(\alpha)\coloneqq\cZ\big(\sigma^{1/\alpha}\big),
\end{equation}
where $\lambda_i$ are the eigenvalues of $\sigma$; and
$(\alpha,\sigma)\mapsto c_{\sigma}(\alpha)$ is jointly continuous on
$[1,1+\delta]\times\{\sigma\succeq0:\tr\sigma\le1\}$ with
$c_{\sigma}(1)=B_{\sigma}'(1)$.
\end{lemma}
\end{technicalclaim}
\claimmeta{\metalabel{Uses}\mref{Def.}{def:admissible-declaration}, \mref{Fact}{fact:meat}, \mref{Def.}{def:feasible}, \mref{Lem.}{lem:dk}, \extref{../foundations/contract.pdf}{foundations/def:pinching}
\metasep \metalabel{Used by}\mref{Thm.}{thm:oracle}, \mref{Cor.}{cor:drop-program}
\metasep newly derived
\metasep implemented (effective: accepted)
\metasep \metalabel{Code}\coderef{https://github.com/fraunhofersg/certiqs-key/blob/v0.1.0/src/certiqs_key/meat_core/objective/limit.py}{certiqs_key.meat_core.objective.limit:relative_entropy_to_pinching}
\metasep \metalabel{Oracles}\coderef{https://github.com/fraunhofersg/certiqs-key/blob/v0.1.0/tests/oracles/registry.py}{registry.py}}

\begin{proof}
\emph{Finiteness.} $\Dtwo(\sigma\|\cZ(\sigma))<\infty$ always: a vector
in the kernel of $\cZ(\sigma)$ satisfies $\sigma^{1/2}Z_j\ket w=0$ for all
$j$, hence lies in the kernel of $\sigma$, so
$\operatorname{supp}\sigma\subseteq\operatorname{supp}\cZ(\sigma)$.

\emph{Fixed supports.} $\operatorname{supp}\sigma^{1/\alpha}
=\operatorname{supp}\sigma$ for every $\alpha$, and the kernel argument
above shows $\operatorname{supp}\cZ(\tau)$ depends on $\tau$ only through
$\operatorname{supp}\tau$; hence
$S\coloneqq\operatorname{supp}A(\alpha)$ is the same subspace for all
$\alpha$, and $S$ is $\cZ$-compatible (its projector is a fixed point of
$\cZ$, being the support projector of a $\cZ$-output, which is
block-diagonal). On $S$, $A(\alpha)$ is positive definite with
eigenvalues bounded away from $0$ uniformly on compact
$\alpha$-intervals (eigenvalues are continuous in $\alpha$ and positive).

\emph{Smoothness.} $\alpha\mapsto\sigma^{1/\alpha}
=\sum_i\lambda_i^{1/\alpha}P_i$ is $C^{1}$ with
$\frac{\mathrm d}{\mathrm d\alpha}\sigma^{1/\alpha}
=-\frac1{\alpha^{2}}\sum_i\lambda_i^{1/\alpha}\ln\lambda_i\,P_i
\eqqcolon\Xi(\alpha)\succeq0$ (each $-\lambda^{1/\alpha}\ln\lambda\ge0$
for $\lambda\in[0,1]$; the $\lambda=0$ terms vanish). Hence
$A'(\alpha)=\cZ(\Xi(\alpha))\succeq0$ is continuous. Restricting to
$S$, where $A$ is smooth and positive definite,
$B_{\sigma}(\alpha)=\tr[A(\alpha)^{\alpha}]$ is $C^1$ with
\begin{equation}\label{eq:Bprime-general}
  B_{\sigma}'(\alpha)
  =\tr\big[A^{\alpha}\ln A\big]
  +\alpha\,\tr\big[A^{\alpha-1}A'\big],
\end{equation}
by \eqref{eq:trace-diff} (for the second term) and the scalar chain rule
in the exponent (for the first; both computed in the joint eigenbasis
data of $A(\alpha)$, using that
$\frac{\partial}{\partial\alpha}a^{\alpha}=a^{\alpha}\ln a$ and the
$C^1$ dependence of eigenvalues and eigenprojectors along the smooth
positive-definite path on $S$; concretely, write
$B_{\sigma}(\alpha)=\tr\exp(\alpha\ln A(\alpha))$ on $S$ and
differentiate the two slots).

\emph{At $\alpha=1$.} $A(1)=\cZ(\sigma)$,
$A'(1)=\cZ(-\sigma\ln\sigma)$, and $A^{\alpha-1}\to\Pi_S$ as
$\alpha\downarrow1$ (eigenvalues of $A$ on $S$ bounded below). Hence
\begin{equation*}
  B_{\sigma}'(1)=\tr[\cZ(\sigma)\ln\cZ(\sigma)]
  +\tr[\Pi_S\,\cZ(-\sigma\ln\sigma)]
  =\tr[\cZ(\sigma)\ln\cZ(\sigma)]-\tr[\sigma\ln\sigma],
\end{equation*}
using self-adjointness of $\cZ$, $\cZ(\Pi_S)=\Pi_S$, and
$\operatorname{supp}(\sigma\ln\sigma)\subseteq\operatorname{supp}\sigma
\subseteq S$. The identification with the relative entropy uses
$\tr[\sigma\ln\cZ(\sigma)]=\tr[\cZ(\sigma)\ln\cZ(\sigma)]$ (again
self-adjointness and block-diagonality of $\ln\cZ(\sigma)$ on $S$), and
non-positivity is monotonicity of $\Dtwo\ge0$.

\emph{Majorant.} In \eqref{eq:Bprime-general},
$\|A(\alpha)\|\le\|\sigma^{1/\alpha}\|\le1$ (a pinching is an average of
unitary conjugates, \xref{foundations/def:pinching}, so it cannot increase
the operator norm), hence
$A^{\alpha-1}\preceq\Pi_S\preceq\idon{}$, and $A'\succeq0$ gives
$\tr[A^{\alpha-1}A']\le\tr[A']=\tr[\Xi(\alpha)]
=-\frac1{\alpha^{2}}\sum_i\lambda_i^{1/\alpha}\ln\lambda_i$; multiplying
by $\alpha$ yields \eqref{eq:majorant}. Joint continuity: the map
$(\alpha,\sigma)\mapsto A(\alpha)$ is jointly continuous; the scalar
functions $(a,\alpha)\mapsto a^{\alpha}\ln a$ (value $0$ at $a=0$) and
$(\lambda,\alpha)\mapsto\lambda^{1/\alpha}\ln\lambda$ are jointly
continuous on $[0,1]\times[1,1+\delta]$; and a trace of a jointly
continuous spectral function of a jointly continuous operator family is
jointly continuous. No step restricts $\delta$ or makes it depend on
$\sigma$: the eigenvalue bound on $S$ holds on every compact
$\alpha$-interval, and the norm bound and the two scalar functions hold for
every $\alpha\ge1$. At $\alpha=1$ the two expressions
\eqref{eq:majorant} and \eqref{eq:Bprime-general} coincide, both equal to
\eqref{eq:Bprime}.
\end{proof}

% claim meat-core/thm:oracle
\claimsummary{meat-core/thm:oracle}{If every acceptable distribution gives each key-live value probability at least some p\ensuremath{_{0}} \textgreater{} 0 and some feasible state reproduces an acceptable distribution, the value tends, as \ensuremath{\alpha} \ensuremath{\to} 1, to a blockwise relative-entropy program. The limit is a validation reference and not a bound.}
\claimbookmark{meat-core/thm:oracle}{Theorem}{thm:oracle}{Limit of the value as the Rényi order tends to one}
\claimleadin{%
As the R\'enyi order tends to one, the weight $\alpha/(\alpha-1)$ of the
divergence term of the objective diverges, and the entropy part of each
branch tends to a relative entropy to its pinching. Remark~\ref{rem:oracle} says what the
limit validates and on which declarations it is available.
}
\begin{theorem}[Limit of the value as the Rényi order tends to one]\label{thm:oracle}
Let $\Decl$ satisfy Definition~\ref{def:admissible-declaration} and let
$\acSet$ satisfy Definition~\ref{def:admissible-acceptance-set}. Assume that
there is $p_0>0$ with $q(\hat c)\ge p_0$ for every $\mathbf q\in\acSet$ and
$\hat c\in\Chkey$, and that the asymptotic feasible set
$\Feas_1\coloneqq\{\omega\in\Feas:\Phimap[\omega]\in\acSet\}$ is
nonempty. Then
\begin{equation}\label{eq:oracle}
  \lim_{\alpha\downarrow1}\hlow(\alpha)
  \;=\;h_1\coloneqq
  \inf_{\omega\in\Feas_1}\ \sum_{\hat c\in\Chkey}\
  \sum_{s\in\SecSet_K(\hat c)}
  \Dtwo\Big(\cG_{\hat c,s}(\omega_s)\,\Big\|\,
  \cZ\big(\cG_{\hat c,s}(\omega_s)\big)\Big).
\end{equation}
\end{theorem}
\claimmeta{\metalabel{Uses}\mref{Def.}{def:admissible-declaration}, \mref{Def.}{def:admissible-acceptance-set}, \mref{Thm.}{thm:branch-formula}, \mref{Prop.}{prop:exact-cancellation}, \mref{Def.}{def:feasible}, \mref{Thm.}{thm:reduction}, \mref{Lem.}{lem:b-derivative}, \extref{../foundations/contract.pdf}{foundations/def:kl}, \extref{../foundations/contract.pdf}{foundations/lem:kernel}
\metasep \metalabel{Used by}\mref{Rem.}{rem:oracle}, \mref{Cor.}{cor:drop-program}, \mref{Req.}{req:limit-not-a-bound}
\metasep newly derived
\metasep implemented (effective: accepted)
\metasep \metalabel{Code}\coderef{https://github.com/fraunhofersg/certiqs-key/blob/v0.1.0/src/certiqs_key/meat_core/objective/limit.py}{certiqs_key.meat_core.objective.limit:check_limit_hypotheses}
\metasep \metalabel{Oracles}\coderef{https://github.com/fraunhofersg/certiqs-key/blob/v0.1.0/tests/oracles/registry.py}{registry.py}}

\begin{proof}
\emph{Idea.} The upper bound evaluates the objective at a state that
reproduces an acceptable distribution exactly, where only the entropy part
survives, and passes to the limit with the derivative of
Lemma~\ref{lem:b-derivative} at $\alpha=1$. The lower bound takes
near-optimal points along $\alpha_n\downarrow1$, uses the divergence part
to force a limit point that reproduces its distribution, and bounds the
entropy part from below by the continuous majorant of that lemma.

\emph{Notation.} Write $\sigma_{\hat c,s}(\omega)\coloneqq\cG_{\hat c,s}(\omega_s)$ and,
for $P_{\hat c}(\omega)>0$,
$e_{\hat c}(\alpha,\omega)\coloneqq\frac{1}{1-\alpha}\log_2
\frac{K_{\hat c}(\omega)}{P_{\hat c}(\omega)}\ge0$.

\emph{Upper bound.} Fix $\omega\in\Feas_1$ and take
$\mathbf q\coloneqq\Phimap[\omega]\in\acSet$. All divergence terms of
$O_{\alpha}$ vanish, and $q(\hat c)=P_{\hat c}(\omega)\ge p_0$ for
$\hat c\in\Chkey$. Then
\begin{equation*}
  O_{\alpha}(\mathbf q,\omega)
  =\sum_{\hat c\in\Chkey}P_{\hat c}\,e_{\hat c}(\alpha,\omega)
  =\sum_{\hat c\in\Chkey}\frac{P_{\hat c}}{(1-\alpha)\ln2}\,
  \ln\frac{K_{\hat c}(\alpha)}{P_{\hat c}} ,
\end{equation*}
and $K_{\hat c}(\alpha)-P_{\hat c}
=\sum_{s}\big(B_{\sigma_{\hat c,s}}(\alpha)
-B_{\sigma_{\hat c,s}}(1)\big)$ by \eqref{eq:c5-sector} and
$B_{\sigma}(1)=\tr\sigma$ together with
$P_{\hat c}-\sum_s\tr\sigma_{\hat c,s}=\tr\Theta_{\hat c}$ cancelling
between $K$ and $P$ (both formulas from \eqref{eq:fK}). By
Lemma~\ref{lem:b-derivative} and l'H\^opital (or the $C^1$ mean value
theorem), as $\alpha\downarrow1$,
\begin{equation*}
  O_{\alpha}(\mathbf q,\omega)\longrightarrow
  -\frac{1}{\ln2}\sum_{\hat c\in\Chkey}\sum_{s}
  B'_{\sigma_{\hat c,s}}(1)
  =\sum_{\hat c\in\Chkey}\sum_{s}
  \Dtwo\big(\sigma_{\hat c,s}\|\cZ(\sigma_{\hat c,s})\big),
\end{equation*}
where the sum may be restricted to $s\in\SecSet_K(\hat c)$ because
key-inert sectors have $\sigma$ already $\cZ$-block-diagonal
(Proposition~\ref{prop:exact-cancellation}) and hence zero relative entropy. Therefore
$\limsup_{\alpha\downarrow1}\hlow(\alpha)\le
h_1$-objective of $\omega$, for every $\omega\in\Feas_1$, so
$\limsup\hlow(\alpha)\le h_1$.

\emph{Lower bound.} Let $\alpha_n\downarrow1$ and pick feasible
$(\mathbf q_n,\omega_n)$ with
$O_{\alpha_n}(\mathbf q_n,\omega_n)\le\hlow(\alpha_n)+\frac1n$; by the
upper bound this sequence of values is bounded above by some $C<\infty$.
Decompose $O_{\alpha}=\frac{\alpha}{\alpha-1}\sfD(\mathbf q\|\Phimap)
+\sum_{\hat c\in\Chkey}q(\hat c)e_{\hat c}$ (undo
\eqref{eq:rearrange}); both parts are non-negative (the divergence by
Gibbs' inequality of \xref{foundations/def:kl}, $\mathbf q_n$ and
$\Phimap[\omega_n]$ being probability vectors, and $e_{\hat c}\ge0$ since $K_{\hat c}\le P_{\hat c}$ by
Theorem~\ref{thm:branch-formula}'s \emph{cone clause}, the points $\omega_n$ being
feasible but not pure). Hence
$\sfD(\mathbf q_n\|\Phimap[\omega_n])\le C\frac{\alpha_n-1}{\alpha_n}
\to0$. By compactness pass to a subsequence with
$(\mathbf q_n,\omega_n)\to(\mathbf q^{*},\omega^{*})
\in\acSet\times\Feas$: $\acSet$ is a closed subset of the simplex
(Definition~\ref{def:admissible-acceptance-set}) and $\Feas$ is compact
(Definition~\ref{def:feasible}). $\sfD(\mathbf q\|\mathbf p)$ is the sum
over $\hat c$ of the kernel $d(q(\hat c),p(\hat c))$ of
\xref{foundations/lem:kernel}, which is lower semicontinuous, and
$\Phimap$ is continuous, so
$\sfD(\mathbf q^{*}\|\Phimap[\omega^{*}])\le
\liminf_n\sfD(\mathbf q_n\|\Phimap[\omega_n])=0$; by the equality case of
Gibbs' inequality (\xref{foundations/def:kl})
$\mathbf q^{*}=\Phimap[\omega^{*}]$ and
$\omega^{*}\in\Feas_1$; in particular
$P_{\hat c}(\omega^{*})=q^{*}(\hat c)\ge p_0$ for $\hat c\in\Chkey$, so
$P_{\hat c}(\omega_n)\ge p_0/2$ eventually. For each
$\hat c\in\Chkey$, using $\ln(1+u)\le u$ with
$u_n\coloneqq\frac{K_{\hat c}(\omega_n)-P_{\hat c}(\omega_n)}
{P_{\hat c}(\omega_n)}\in[-1,0]$ and $\frac{1}{1-\alpha_n}<0$:
\begin{equation*}
  e_{\hat c}(\alpha_n,\omega_n)
  \;\ge\;\frac{u_n}{(1-\alpha_n)\ln2}
  \;=\;\frac{1}{\ln2\,P_{\hat c}(\omega_n)}
  \sum_{s}\Big[-\frac{B_{\sigma_{\hat c,s}(\omega_n)}(\alpha_n)
  -B_{\sigma_{\hat c,s}(\omega_n)}(1)}{\alpha_n-1}\Big].
\end{equation*}
By the mean value theorem there are $\xi_n\in(1,\alpha_n)$ with the
bracket equal to $-B'_{\sigma_{\hat c,s}(\omega_n)}(\xi_n)
\ge-c_{\sigma_{\hat c,s}(\omega_n)}(\xi_n)$ by \eqref{eq:majorant}. Joint
continuity of $(\alpha,\sigma)\mapsto c_{\sigma}(\alpha)$
(Lemma~\ref{lem:b-derivative}) and
$\sigma_{\hat c,s}(\omega_n)\to\sigma_{\hat c,s}(\omega^{*})$ (the maps
$\cG_{\hat c,s}$ are continuous) give
$c_{\sigma(\omega_n)}(\xi_n)\to
c_{\sigma(\omega^{*})}(1)=B'_{\sigma(\omega^{*})}(1)$. Hence
\begin{equation*}
  \liminf_n\ q_n(\hat c)\,e_{\hat c}(\alpha_n,\omega_n)
  \;\ge\;\frac{q^{*}(\hat c)}{\ln2\,P_{\hat c}(\omega^{*})}
  \sum_{s}\big(-B'_{\sigma_{\hat c,s}(\omega^{*})}(1)\big)
  \;=\;\sum_{s}\Dtwo\big(\sigma_{\hat c,s}(\omega^{*})\,\big\|\,
  \cZ(\sigma_{\hat c,s}(\omega^{*}))\big),
\end{equation*}
using $q^{*}(\hat c)=P_{\hat c}(\omega^{*})$. Dropping the non-negative
divergence part and summing over $\hat c$ (a finite sum, so
$\liminf$ passes through),
$\liminf_n\hlow(\alpha_n)\ge h_1$-objective of
$\omega^{*}\ge h_1$. Combining with the upper bound proves
\eqref{eq:oracle}.
\end{proof}

% claim meat-core/rem:oracle
\claimsummary{meat-core/rem:oracle}{The limit of \hyperlink{claim-meat-core/thm:oracle}{Theorem~\ref*{thm:oracle}} is a blockwise relative-entropy program that checks an implementation as \ensuremath{\alpha} \ensuremath{\to} 1; it validates no finite-\ensuremath{\alpha} value and bounds nothing. It is available only where every acceptable distribution is positive on the key-live values.}
\claimbookmark{meat-core/rem:oracle}{Remark}{rem:oracle}{Validation against the order-one limit}
\begin{remark}[Validation against the order-one limit]\label{rem:oracle}
The right-hand side of \eqref{eq:oracle} is a blockwise relative-entropy
objective, of the structure of the conic von~Neumann key-distance program
of \cite[Eqs.~(1)--(6)]{winick2018reliable} with acceptance-set membership
in place of the equality constraints used there. Theorem~\ref{thm:oracle}
is stated over the declared $\Feas$, so it covers arbitrary $\CList$
unchanged, and affine constraints enter both programs natively. The limit
validates the consistency of the $\alpha\to1^{+}$ limit of an
implementation of \eqref{eq:choi-program}; it does not validate
finite-$\alpha$ objective values or the finite-$\alpha$ $K_{\hat c}$, and
it is a bound of no kind, in particular no lower bound on $\hlow$
(Requirement~\ref{req:limit-not-a-bound}).

\emph{On which declarations it is available.}
Theorem~\ref{thm:oracle} assumes a $p_0>0$ with $q(\hat c)\ge p_0$ for
\emph{every} $\mathbf q\in\acSet$ and every $\hat c\in\Chkey$, which asks
that every element of $\acSet$ be strictly positive on $\Chkey$; on a
compact $\acSet$ that is exactly the existence of $p_0$, the minimum being
attained. A singleton $\acSet=\{\mathbf q^{*}\}$ with
$\Chkey\subseteq\operatorname{supp}\mathbf q^{*}$ satisfies it; a polytope
with a vertex vanishing somewhere on $\Chkey$ does not, and there the
limit is unavailable. Theorem~\ref{thm:oracle} is a validation statement:
no key length passes through it, and a declaration outside its hypothesis
loses this check and no claim. The condition is narrower than the
\termref{term-acceptance-support}{acceptance support} $\Chacc$ of \eqref{eq:chacc-def}, which asks
positivity of \emph{some} element and against which (H-POS) of
Theorem~\ref{thm:firstorder} is stated for exactly that reason.
\end{remark}
\claimmeta{\metalabel{Uses}\mref{Thm.}{thm:choi}, \mref{Thm.}{thm:oracle}
\metasep \metalabel{Used by}\mref{Req.}{req:limit-not-a-bound}
\metasep newly derived
\metasep accepted}

\section{The dropped program}\label{group-meat-core/lem:rank-one-sector}
\sectionintro{A run may give up the entropy of chosen pairs of a \termref{term-key-live-acceptance-value}{key-live acceptance value} and a joint sector, setting their deficits to zero while keeping their variables, statistics and constraints (\hyperlink{claim-meat-core/def:dropped}{Definition~\ref*{def:dropped}}). The resulting value bounds h\_low and the single-round rate from below (\hyperlink{claim-meat-core/prop:drop-bound}{Proposition~\ref*{prop:drop-bound}}); it is attained when finite, and a \termref{term-fully-dropped}{fully dropped} joint sector enters only through its statistics, with no matrix function (\hyperlink{claim-meat-core/cor:drop-program}{Corollary~\ref*{cor:drop-program}}). With every pair dropped the value is \ensuremath{\alpha}/(\ensuremath{\alpha}\ensuremath{-}1) times the minimal divergence between the accepted distributions and the statistics of the feasible set, a floor under h\_low that lies within the entropy ceiling of it, and a positive lower bound on it proves that no feasible state has an accepted distribution as its statistics (\hyperlink{claim-meat-core/prop:minimal-divergence}{Proposition~\ref*{prop:minimal-divergence}}). On an Alice sector whose marginal has rank one, with default readouts and Alice holding the key, the deficit is linear with constants from declared numbers (\hyperlink{claim-meat-core/lem:rank-one-sector}{Lemma~\ref*{lem:rank-one-sector}}). The objective and the dropped objective, read on the whole cone, are positively homogeneous of degree one (\hyperlink{claim-meat-core/lem:objective-homogeneous}{Lemma~\ref*{lem:objective-homogeneous}}).}
\sectionimports{\extref{../foundations/contract.pdf}{foundations/def:kl}, \extref{../foundations/contract.pdf}{foundations/lem:kernel}}

% claim meat-core/lem:rank-one-sector
\claimsummary{meat-core/lem:rank-one-sector}{On an Alice sector of rank-one marginal every feasible block is the fixed pure state times an operator on Q, and with default readouts and Alice holding the key its deficit is linear, with constants from declared numbers. A uniform key bit gives one bit per unit of pass weight.}
\claimbookmark{meat-core/lem:rank-one-sector}{Lemma}{lem:rank-one-sector}{Linear deficit on a rank-one Alice sector}
\claimleadin{%
When Alice's marginal on a sector is pure, her part of every feasible block
is fixed. With the default readouts and Alice holding the key, the deficit
of the block then reduces to a Rényi quantity of her key value given the
announcement, weighted linearly by the block's statistics.
}
\begin{lemma}[Linear deficit on a rank-one Alice sector]\label{lem:rank-one-sector}
Let $\Decl$ satisfy Definition~\ref{def:admissible-declaration} in \exttermref{../transmitter/contract.pdf}{term-trusted-source-mode}{trusted-source mode}
and the general format, and let $t\in\mathbb{T}$ with
$\tau^{(t)}_A=p_t\dyad{\hat v}$, $p_t>0$, $\ket{\hat v}\in\Hs{A_t}$ a unit
vector.
\begin{enumerate}[label=(\roman*),leftmargin=2.2em]
\item \emph{(product blocks)} Every list $(\omega^{(t,s)})_{s\in\SecSet}$ of
positive semidefinite blocks with
$\sum_s\tr_{\Qsys_s}\omega^{(t,s)}=\tau^{(t)}_A$, in particular the sector-$t$
part of every $\omega\in\Feas$, has the form
$\omega^{(t,s)}=\dyad{\hat v}\otimes\sigma_s$ with
$\sigma_s\in\Pos{\Qsys_s}$ and $\sum_s\tr\sigma_s=p_t$.
\item \emph{(\termref{term-pass-class}{pass classes} are product sets)} If the
\exttermref{../protocol-compiler/contract.pdf}{term-readout-map}{readout maps} are the default ones,
$\rd{\mathsf A}(x,v)=x$ and $\rd{\mathsf B}(x,v)=v$, then for
every \exttermref{../protocol-compiler/contract.pdf}{term-transcript-protocol}{transcript protocol} and every $\kappa\in\Kann$ there are sets
$\Xset_{\kappa}\subseteq\Xset$, $\Vset_{\kappa}\subseteq\Vset$ with
$a^{-1}(\kappa)=\Xset_{\kappa}\times\Vset_{\kappa}$; if moreover
$\hold=\mathsf A$, then $r(x,v)=r_{\mathsf A}(x,\kappa)$ on $a^{-1}(\kappa)$.
\item \emph{(linear deficit)} Under the hypotheses of (ii), let
$\hat c\in\Chkey$, $s\in\SecSet$ and $\alpha\in(1,\infty)$, put
$\mu_x\coloneqq\bra{\hat v}M^{(t)}_x\ket{\hat v}$ and, for
$\kappa\in\Kann_{\mathrm{pass}}(\hat c)$ and $\mathsf z\in\Zval$,
\begin{equation}\label{eq:rank-one-weights}
  m_{\kappa}\coloneqq\sum_{x\in\Xset_{\kappa}}\mu_x,\qquad
  m_{\kappa,\mathsf z}\coloneqq\sum_{x\in\Xset_{\kappa}:\,
  r_{\mathsf A}(x,\kappa)=\mathsf z}\mu_x,\qquad
  \eta_{\kappa}\coloneqq
  \begin{cases}
   1-\sum_{\mathsf z}\big(m_{\kappa,\mathsf z}/m_{\kappa}\big)^{\alpha},
   & m_{\kappa}>0,\\ 0,& m_{\kappa}=0 .
  \end{cases}
\end{equation}
Then for every $\sigma_s\in\Pos{\Qsys_s}$
\begin{equation}\label{eq:rank-one-deficit}
  f_{\hat c,(t,s)}\big(\dyad{\hat v}\otimes\sigma_s\big)
  =-\sum_{\kappa\in\Kann_{\mathrm{pass}}(\hat c)}\eta_{\kappa}\,m_{\kappa}
  \sum_{v\in\Vset_{\kappa}}\tr\big[F^{(s)}_v\sigma_s\big],
\end{equation}
a \emph{linear} function of $\sigma_s$, with $0\le\eta_{\kappa}<1$. In
\exttermref{../protocol-compiler/contract.pdf}{term-ensemble-form}{ensemble form} $\mu_x=\tr N^{(t)}_x/p_t$ and $p_t=\tr\tau^{(t)}_A$, declared
numbers.
\item \emph{(two values)} If within every pass class of $\hat c$ the key
takes $d$ values of equal weight, $m_{\kappa,\mathsf z}=m_{\kappa}/d$, then
$\eta_{\kappa}=1-d^{1-\alpha}$ and the sector contributes to
$\frac{1}{1-\alpha}\log_2$ of the ratio $K/P$ \emph{as if the conditional
entropy were $\log_2d$ on the pass weight}: one bit per unit of pass weight
for a uniform key bit. The qualifier matters: what the
expression assigns to the pair is its fail weight plus $d^{1-\alpha}$ times
its pass weight, so with a non-zero fail weight in the same joint pair the
contribution is strictly less than that of a sector of conditional entropy
$\log_2d$, and equals it exactly when the pass weight is the whole of the
pair's weight. If the key is
constant on every $\Xset_{\kappa}\cap\{x:\mu_x>0\}$ then $\eta_{\kappa}=0$
and $f_{\hat c,(t,s)}$ vanishes on these blocks.
\end{enumerate}
\end{lemma}
\claimmeta{\metalabel{Uses}\mref{Def.}{def:declaration}, \mref{Def.}{def:admissible-declaration}, \mref{Lem.}{lem:ensemble-form}, \mref{Def.}{def:maps}, \mref{Def.}{def:branch}, \mref{Thm.}{thm:branch-formula}, \mref{Def.}{def:joint-sectors}, \mref{Def.}{def:feasible}
\metasep newly derived
\metasep implemented (effective: accepted)
\metasep \metalabel{Code}\coderef{https://github.com/fraunhofersg/certiqs-key/blob/v0.1.0/src/certiqs_key/meat_core/objective/joint.py}{certiqs_key.meat_core.objective.joint:rank_one_deficit}}

\begin{proof}
(i) Let $\omega_t\coloneqq\bigoplus_s\omega^{(t,s)}\in\Pos{A_t\Qsys}$, so
$\tr_{\Qsys}\omega_t=p_t\dyad{\hat v}$, and let
$\Pi^{\perp}\coloneqq\idon{A_t}-\dyad{\hat v}$. Then
$\tr[(\Pi^{\perp}\otimes\idon{})\omega_t]=p_t\tr[\Pi^{\perp}\dyad{\hat v}]=0$,
and for a positive semidefinite $\omega_t$ this forces
$(\Pi^{\perp}\otimes\idon{})\omega_t=0$. Hence
$\omega_t=(\dyad{\hat v}\otimes\idon{})\omega_t(\dyad{\hat v}\otimes\idon{})
=\dyad{\hat v}\otimes\sigma$ with
$\sigma\coloneqq(\bra{\hat v}\otimes\idon{})\omega_t(\ket{\hat v}\otimes\idon{})
\succeq0$, block diagonal along $\SecSet$ because $\omega_t$ is, with blocks
$\sigma_s$; and $\tr\sigma=\tr\omega_t=p_t$. For $\omega\in\Feas$ the
hypothesis is the marginal line of \eqref{eq:feas-general}.

(ii) By \eqref{eq:transcript}, $a(x,v)=\kappa$ holds exactly when
$a_m(\rd{\pi_m}(x,v),\kappa_1,\dots,\kappa_{m-1})=\kappa_m$ for every $m$:
by induction on $m$, if the first $m-1$ components of $a(x,v)$ are
$\kappa_1,\dots,\kappa_{m-1}$ then its $m$-th component is the left-hand
side. With the default readouts the $m$-th condition is a condition on $x$
alone when $\pi_m=\mathsf A$ and on $v$ alone when $\pi_m=\mathsf B$. So
$a^{-1}(\kappa)=\Xset_{\kappa}\times\Vset_{\kappa}$ with $\Xset_{\kappa}$ the
set of $x$ meeting all conditions with $\pi_m=\mathsf A$ and
$\Vset_{\kappa}$ the set of $v$ meeting all those with $\pi_m=\mathsf B$. The
last statement is \eqref{eq:partycomp} with $\rd{\mathsf A}(x,v)=x$.

(iii) Fix $\kappa\in\Kann_{\mathrm{pass}}(\hat c)$ and write
$\ket{m_x}\coloneqq\sqrt{M^{(t)}_x}\ket{\hat v}$, so
$\braket{m_x}{m_x}=\mu_x$. By (ii) the operator $K^{(s)}_{\kappa}$ of
\eqref{eq:kraus} satisfies, on vectors $\ket{\hat v}\otimes\ket{\xi}$,
\begin{equation*}
  K^{(s)}_{\kappa}\big(\ket{\hat v}\otimes\ket{\xi}\big)
  =\ket{a_{\kappa}}\otimes\ket{\kappa}\otimes B_{\kappa}\ket{\xi},\qquad
  \ket{a_{\kappa}}\coloneqq\sum_{x\in\Xset_{\kappa}}
  \ket{r_{\mathsf A}(x,\kappa)}_R\ket{x}_X\ket{m_x}_A,\qquad
  B_{\kappa}\coloneqq\sum_{v\in\Vset_{\kappa}}\ket{v}_Y\otimes\sqrt{F^{(s)}_v},
\end{equation*}
up to the fixed ordering of tensor factors. Hence
$G_{\kappa}\coloneqq K^{(s)}_{\kappa}(\dyad{\hat v}\otimes\sigma_s)
K^{(s)\dagger}_{\kappa}=\dyad{a_{\kappa}}\otimes\dyad{\kappa}\otimes
\beta_{\kappa}$ with $\beta_{\kappa}\coloneqq B_{\kappa}\sigma_sB_{\kappa}^{\dagger}
\succeq0$, and $\braket{a_{\kappa}}{a_{\kappa}}=m_{\kappa}$, the kets
$\ket{x}_X$ being orthonormal. The operators $G_{\kappa}$ for different
$\kappa$ are supported on the mutually orthogonal subspaces fixed by
$\ket{\kappa}$, which are invariant under $\cZ$ and under powers, so as in
the proof of Lemma~\ref{lem:joint-sectors}\,(ii) the functional $T$ and the
trace add over $\kappa$, and
$f_{\hat c,(t,s)}=\sum_{\kappa}\big(T(G_{\kappa})-\tr G_{\kappa}\big)$. If
$m_{\kappa}=0$ then $G_{\kappa}=0$. Otherwise
$\dyad{a_{\kappa}}=m_{\kappa}\Pi_a$ with $\Pi_a$ a rank-one projector, so
$G_{\kappa}^{1/\alpha}=m_{\kappa}^{1/\alpha-1}\dyad{a_{\kappa}}\otimes
\dyad{\kappa}\otimes\beta_{\kappa}^{1/\alpha}$. The vectors
$\ket{a_{\kappa,\mathsf z}}\coloneqq Z_{\mathsf z}\ket{a_{\kappa}}$ are
mutually orthogonal with squared norms $m_{\kappa,\mathsf z}$, and
$Z_{\perp}\ket{a_{\kappa}}=0$ because $\kappa$ is a pass class, so
$\cZ(G_{\kappa}^{1/\alpha})=m_{\kappa}^{1/\alpha-1}\sum_{\mathsf z}
\dyad{a_{\kappa,\mathsf z}}\otimes\dyad{\kappa}\otimes
\beta_{\kappa}^{1/\alpha}$. The summands have mutually orthogonal supports
and $(\dyad{a_{\kappa,\mathsf z}})^{\alpha}
=m_{\kappa,\mathsf z}^{\alpha-1}\dyad{a_{\kappa,\mathsf z}}$, hence
\begin{equation*}
  T(G_{\kappa})=m_{\kappa}^{1-\alpha}\sum_{\mathsf z}
  m_{\kappa,\mathsf z}^{\alpha}\,\tr\beta_{\kappa}
  =\Big(\sum_{\mathsf z}\big(m_{\kappa,\mathsf z}/m_{\kappa}\big)^{\alpha}\Big)
  \tr G_{\kappa},\qquad
  \tr G_{\kappa}=m_{\kappa}\tr\beta_{\kappa}
  =m_{\kappa}\sum_{v\in\Vset_{\kappa}}\tr\big[F^{(s)}_v\sigma_s\big].
\end{equation*}
Subtracting gives \eqref{eq:rank-one-deficit}. The numbers
$m_{\kappa,\mathsf z}/m_{\kappa}$ form a probability vector, so
$0<\sum_{\mathsf z}(\cdot)^{\alpha}\le1$ and $0\le\eta_{\kappa}<1$. In
ensemble form $\mu_x=\tr[M^{(t)}_x\tau^{(t)}_A]/p_t=\tr N^{(t)}_x/p_t$ by
Lemma~\ref{lem:ensemble-form}\,(ii).

(iv) With $d$ equal weights $\sum_{\mathsf z}(1/d)^{\alpha}=d^{1-\alpha}$. If
every pass class of $\hat c$ has this form then, by (iii), on these blocks
$P_{\hat c,(t,s)}+f_{\hat c,(t,s)}$ equals the fail weight plus
$d^{1-\alpha}=2^{(1-\alpha)\log_2d}$ times the pass weight, which is the
value the branch formula \eqref{eq:c5} assigns to conditional entropy
$\log_2d$ on the pass weight. If the key is constant where $\mu_x>0$, one
$m_{\kappa,\mathsf z}$ equals $m_{\kappa}$ and the others vanish.
\end{proof}

% claim meat-core/def:dropped
\claimsummary{meat-core/def:dropped}{A run may give up the entropy of chosen pairs of a key-live acceptance value and a joint sector by setting their deficits to zero, keeping every variable, statistic and constraint. The result is a program of its own; once any pair is dropped, its value is never reported as the undropped one.}
\claimbookmark{meat-core/def:dropped}{Definition}{def:dropped}{Dropped set and dropped program}
\claimleadin{%
A run may prefer a weaker bound that is cheaper to compute: a
\exttermref{../protocol-compiler/contract.pdf}{term-joint-sector}{joint sector} whose deficits it gives up for every
\termref{term-key-live-acceptance-value}{key-live acceptance value} needs no matrix function (Corollary~\ref{cor:drop-program}\,(iii)). What the run
gives up is only the entropy those deficits carry; the sector stays in the
program.
}
\begin{definition}[Dropped set and dropped program]\label{def:dropped}
Let $\Decl$ be a declaration in the general format of
Definition~\ref{def:declaration}. A
\termdef{term-dropped-set}{dropped set}
is a set
$\mathcal{D}\subseteq\Chkey\times(\mathbb{T}\times\SecSet)$ of pairs of a
\termref{term-key-live-acceptance-value}{key-live acceptance value} and a \exttermref{../protocol-compiler/contract.pdf}{term-joint-sector}{joint sector}. For
$\omega=\bigoplus_{(t,s)}\omega^{(t,s)}$ block diagonal along the joint
sectors put
\begin{equation}\label{eq:KD}
  K^{\mathcal{D}}_{\hat c}(\omega)\coloneqq P_{\hat c}(\omega)
  +\sum_{(t,s):\,(\hat c,(t,s))\notin\mathcal{D}}
  f_{\hat c,(t,s)}\big(\omega^{(t,s)}\big)
  \qquad(\hat c\in\Chkey),
\end{equation}
with $f_{\hat c,(t,s)}$ of \eqref{eq:joint-deficit}; let
$O^{\mathcal{D}}_{\alpha}$ be \eqref{eq:objective} with
$K^{\mathcal{D}}_{\hat c}$ in place of $K_{\hat c}$, with the same
conventions at vanishing arguments; and let
\begin{equation}\label{eq:hlowD}
  \underline{h}^{\mathcal{D}}_{\alpha}\coloneqq
  \inf_{\mathbf q\in\acSet}\ \inf_{\omega\in\Feas}\
  O^{\mathcal{D}}_{\alpha}(\mathbf q,\omega).
\end{equation}
The minimisation in \eqref{eq:hlowD} is the
\termdef{term-dropped-program}{dropped program}.
A joint sector $(t,s)$ is
\termdef{term-fully-dropped}{fully dropped}
when $\Chkey\times\{(t,s)\}\subseteq\mathcal{D}$. For
$\mathcal{D}=\varnothing$ these are $K_{\hat c}$ (\eqref{eq:K-joint}),
$O_{\alpha}$ and $\hlow$.

\emph{What is not changed.} The \exttermref{../protocol-compiler/contract.pdf}{term-feasible-set}{feasible set} is $\Feas$, unchanged: a
dropped sector's block remains a variable, enters $\Phimap$ and every
$P_{\hat c}$ with its full weight --- the term
$P_{\hat c,(t,s)}(\omega^{(t,s)})$ stays inside $P_{\hat c}$ in
\eqref{eq:KD} --- and enters the marginal condition and every constraint
row. Setting a dropped block to zero, omitting its statistics, or removing
its weight from $P_{\hat c}$ in \eqref{eq:KD} are different operations, and
each can overstate the value (Proposition~\ref{prop:retention} for the
first). \emph{The dropped set is a formulation choice of a run and not
\exttermref{../foundations/contract.pdf}{term-declared-data}{declared data}}.
\end{definition}
\claimmeta{\metalabel{Uses}\mref{Def.}{def:declaration}, \mref{Thm.}{thm:branch-formula}, \mref{Def.}{def:joint-sectors}, \mref{Def.}{def:feasible}
\metasep \metalabel{Used by}\mref{Prop.}{prop:drop-bound}, \mref{Cor.}{cor:drop-program}, \mref{Prop.}{prop:minimal-divergence}, \mref{Lem.}{lem:objective-homogeneous}, \mref{Prop.}{prop:model-pass-bound}, \mref{Req.}{req:dropped-set-record}, \mref{Req.}{req:unreachable-acceptance-set}, \mref{Cor.}{cor:fo-dropped}, \mref{Lem.}{lem:certified-bracket}, \mref{Lem.}{lem:certified-bracket-enclosed}
\metasep newly derived
\metasep implemented (effective: accepted)
\metasep \metalabel{Code}\coderef{https://github.com/fraunhofersg/certiqs-key/blob/v0.1.0/src/certiqs_key/meat_core/objective/joint.py}{certiqs_key.meat_core.objective.joint:joint_objective_value}}

% claim meat-core/prop:drop-bound
\claimsummary{meat-core/prop:drop-bound}{Setting the deficits of a dropped set to zero can only lower the objective, so the dropped value lies between 0 and h\_low, hence below the single-round rate, and never grows as the set grows. Dropping only key-inert pairs changes nothing.}
\claimbookmark{meat-core/prop:drop-bound}{Proposition}{prop:drop-bound}{The dropped program bounds the value from below}
\claimleadin{%
Deleting summands that are all non-positive can only raise each $K$, and
the objective falls as $K$ rises. The loss is in the rate and needs no
hypothesis on the key map or the key holder, and it is a bound only because
the dropped sectors stay in the program.
}
\begin{proposition}[The dropped program bounds the value from below]\label{prop:drop-bound}
Let $\Decl$ satisfy Definition~\ref{def:admissible-declaration} in the general format, in
either source mode and with an arbitrary list $\CList$; let
$\alpha\in(1,\infty)$ and let $\mathcal{D}\subseteq\mathcal{D}'$ be dropped
sets (Definition~\ref{def:dropped}). For every $\omega$ block diagonal along
the joint sectors, every $\hat c\in\Chkey$ and every distribution $\mathbf q$
on $\Chatset$:
\begin{enumerate}[label=(\roman*),leftmargin=2.2em]
\item $0\le K_{\hat c}(\omega)\le K^{\mathcal{D}}_{\hat c}(\omega)
\le K^{\mathcal{D}'}_{\hat c}(\omega)\le P_{\hat c}(\omega)$;
\item $O^{\mathcal{D}'}_{\alpha}(\mathbf q,\omega)\le
O^{\mathcal{D}}_{\alpha}(\mathbf q,\omega)\le O_{\alpha}(\mathbf q,\omega)$
in the extended-real sense, and $O^{\mathcal{D}'}_{\alpha}(\mathbf
q,\omega)\ge0$ whenever $\tr\omega=1$;
\item consequently
$0\le\underline{h}^{\mathcal{D}'}_{\alpha}\le
\underline{h}^{\mathcal{D}}_{\alpha}\le\hlow\le\hqkd(\CList)$;
\item if every pair of $\mathcal{D}$ is key-inert
(Definition~\ref{def:joint-sectors}) then
$K^{\mathcal{D}}_{\hat c}=K_{\hat c}$ and
$O^{\mathcal{D}}_{\alpha}=O_{\alpha}$ identically.
\end{enumerate}
\end{proposition}
\claimmeta{\metalabel{Uses}\mref{Def.}{def:admissible-declaration}, \mref{Def.}{def:joint-sectors}, \mref{Lem.}{lem:joint-sectors}, \mref{Def.}{def:feasible}, \mref{Cor.}{cor:reduction-constrained}, \mref{Def.}{def:dropped}, \extref{../foundations/contract.pdf}{foundations/def:kl}, \extref{../foundations/contract.pdf}{foundations/lem:kernel}
\metasep \metalabel{Used by}\mref{Cor.}{cor:drop-program}, \mref{Prop.}{prop:minimal-divergence}, \mref{Req.}{req:dropped-set-record}, \mref{Cor.}{cor:fo-dropped}, \mref{Lem.}{lem:certified-bracket}, \mref{Cor.}{cor:certified-keylength}
\metasep newly derived
\metasep accepted}

\begin{proof}
(i) By \eqref{eq:K-joint} and \eqref{eq:KD}, the three quantities
$K_{\hat c}$, $K^{\mathcal{D}}_{\hat c}$, $K^{\mathcal{D}'}_{\hat c}$ are
$P_{\hat c}(\omega)$ plus the sum of $f_{\hat c,(t,s)}(\omega^{(t,s)})$ over
all joint sectors, over those with $(\hat c,(t,s))\notin\mathcal{D}$, and
over those with $(\hat c,(t,s))\notin\mathcal{D}'$ respectively: three
nested index sets, the last the smallest. Every summand is $\le0$
(Lemma~\ref{lem:joint-sectors}\,(ii)), so deleting summands does not
decrease the sum, which gives the two middle inequalities, and the last
holds because the remaining sum is $\le0$. $K_{\hat c}(\omega)\ge0$: by
Lemma~\ref{lem:joint-sectors}\,(ii),
$P_{\hat c,(t,s)}+f_{\hat c,(t,s)}\ge0$ on every joint sector, and
$K_{\hat c}(\omega)$ is the sum of these over the joint sectors.

(ii) In \eqref{eq:objective} the quantity $K$ enters only through the terms
$\frac{1}{\alpha-1}d\big(q(\hat c),K\big)$, $\hat c\in\Chkey$, with $d$ the
kernel of \xref{foundations/lem:kernel} --- whose conventions at vanishing arguments
are those of \xref{foundations/def:kl} --- and the prefactor is positive. $d$
is non-increasing in its second argument, so (i) gives the two inequalities
termwise, in the extended-real sense. For the lower bound let $\tr\omega=1$,
so that $\mathbf P\coloneqq\Phimap[\omega]$ is a probability distribution. If
some $q(\hat c)>0$ has $P_{\hat c}=0$, or
$K^{\mathcal{D}'}_{\hat c}=0$ with $\hat c\in\Chkey$, the objective is
$+\infty$. Otherwise, by $K^{\mathcal{D}'}_{\hat c}\le P_{\hat c}$ and
monotonicity of $d$,
$d(q(\hat c),K^{\mathcal{D}'}_{\hat c})\ge d(q(\hat c),P_{\hat c})$, so
\begin{equation*}
  O^{\mathcal{D}'}_{\alpha}(\mathbf q,\omega)\;\ge\;
  \frac{\alpha}{\alpha-1}\sum_{\hat c\in\Chatset}
  q(\hat c)\log_2\frac{q(\hat c)}{P_{\hat c}}
  =\frac{\alpha}{\alpha-1}\,\sfD(\mathbf q\,\|\,\mathbf P)\;\ge\;0 ,
\end{equation*}
the terms $\hat c\notin\Chkey$ carrying the prefactor
$\frac{\alpha}{\alpha-1}$ already and the terms $\hat c\in\Chkey$ the sum
$1+\frac{1}{\alpha-1}$ of their two prefactors; the last step is
non-negativity of the divergence between probability distributions,
Gibbs' inequality of \xref{foundations/def:kl}.

(iii) Take infima over $\acSet\times\Feas$ in (ii); every $\omega\in\Feas$
has $\tr\omega=1$. The last inequality is
Corollary~\ref{cor:reduction-constrained}.

(iv) For a key-inert pair $f_{\hat c,(t,s)}\equiv0$
(Lemma~\ref{lem:joint-sectors}\,(iii)), so the summands deleted in
\eqref{eq:KD} are identically zero.
\end{proof}

% claim meat-core/cor:drop-program
\claimsummary{meat-core/cor:drop-program}{The dropped objective is jointly convex and lower semicontinuous, its infimum attained when finite, and on a fully dropped joint sector it reads the block only through its statistics. As the R\'enyi order tends to one the dropped value tends to the relative-entropy program over the kept pairs.}
\claimbookmark{meat-core/cor:drop-program}{Corollary}{cor:drop-program}{Convexity, gradient and limit of the dropped program}
\claimleadin{%
The dropped program keeps the undropped program's feasible set and only
removes deficits from each $K$, so the properties the program's numerics
rest on carry over, and a fully dropped joint sector costs no matrix
function. Part~(iii) is written in state coordinates; in ensemble form the
statistics operators are read in the Choi coordinates of
Lemma~\ref{lem:choi-factor}, Alice's blocks not being declared there.
}
\begin{corollary}[Convexity, gradient and limit of the dropped program]\label{cor:drop-program}
Under the hypotheses of Proposition~\ref{prop:drop-bound}, with $\acSet$
convex and compact (Definition~\ref{def:admissible-acceptance-set}):
\begin{enumerate}[label=(\roman*),leftmargin=2.2em]
\item $K^{\mathcal{D}}_{\hat c}$ is concave and continuous on the cone of
operators block diagonal along the joint sectors;
\item $O^{\mathcal{D}}_{\alpha}$ is jointly convex and lower semicontinuous on
$\Delta(\Chatset)\times\Feas$ with values in $[0,+\infty]$, and the infimum
\eqref{eq:hlowD} is attained whenever it is finite;
\item if $(t,s)$ is fully dropped, then $O^{\mathcal{D}}_{\alpha}$ depends on
the block $\omega^{(t,s)}$ only through the real numbers
$\tr[(M^{(t)}_x\otimes F^{(s)}_v)\,\omega^{(t,s)}]$, $(x,v)\in\Xset\times\Vset$;
at a point where every $P_{\hat c}$ with $q(\hat c)>0$ and every
$K^{\mathcal{D}}_{\hat c}$ with $q(\hat c)>0$, $\hat c\in\Chkey$, is
positive, its derivative in the direction of a Hermitian $\Delta$ supported
on that joint sector is
\begin{equation}\label{eq:drop-gradient}
  -\log_2\!\mathrm{e}\sum_{\hat c\in\Chatset}\beta_{\hat c}\,
  \tr\big[\Gamma^{(t,s)}_{\hat c}\Delta\big],\qquad
  \Gamma^{(t,s)}_{\hat c}\coloneqq\!\!\sum_{(x,v):\,\phi(a(x,v))=\hat c}\!\!
  M^{(t)}_x\otimes F^{(s)}_v,\qquad
  \beta_{\hat c}\coloneqq
  \begin{cases}
   \frac{\alpha}{\alpha-1}\frac{q(\hat c)}{P_{\hat c}},&\hat c\notin\Chkey,\\[0.6ex]
   \frac{q(\hat c)}{P_{\hat c}}+\frac{1}{\alpha-1}
   \frac{q(\hat c)}{K^{\mathcal{D}}_{\hat c}},&\hat c\in\Chkey,
  \end{cases}
\end{equation}
with $\beta_{\hat c}\coloneqq0$ where $q(\hat c)=0$;
\item under the additional hypotheses of Theorem~\ref{thm:oracle},
\begin{equation}\label{eq:drop-oracle}
  \lim_{\alpha\downarrow1}\underline{h}^{\mathcal{D}}_{\alpha}
  =\inf_{\omega\in\Feas_1}\ \sum_{\hat c\in\Chkey}\
  \sum_{\substack{(t,s)\in\mathbb{J}_K(\hat c)\\(\hat c,(t,s))\notin\mathcal{D}}}
  \Dtwo\Big(\cG_{\hat c,(t,s)}\big(\omega^{(t,s)}\big)\,\Big\|\,
  \cZ\big(\cG_{\hat c,(t,s)}\big(\omega^{(t,s)}\big)\big)\Big).
\end{equation}
\end{enumerate}
\end{corollary}
\claimmeta{\metalabel{Uses}\mref{Def.}{def:admissible-declaration}, \mref{Def.}{def:maps}, \mref{Def.}{def:admissible-acceptance-set}, \mref{Thm.}{thm:branch-formula}, \mref{Def.}{def:joint-sectors}, \mref{Lem.}{lem:joint-sectors}, \mref{Def.}{def:feasible}, \mref{Lem.}{lem:b-derivative}, \mref{Thm.}{thm:oracle}, \mref{Def.}{def:dropped}, \mref{Prop.}{prop:drop-bound}, \extref{../foundations/contract.pdf}{foundations/def:kl}, \extref{../foundations/contract.pdf}{foundations/lem:kernel}
\metasep \metalabel{Used by}\mref{Prop.}{prop:minimal-divergence}, \mref{Cor.}{cor:fo-dropped}
\metasep newly derived
\metasep implemented (effective: accepted)
\metasep \metalabel{Code}\coderef{https://github.com/fraunhofersg/certiqs-key/blob/v0.1.0/src/certiqs_key/meat_core/objective/joint_gradient.py}{certiqs_key.meat_core.objective.joint_gradient:joint_gradient}}

\begin{proof}
\emph{Idea.} Each $K^{\mathcal{D}}_{\hat c}$ is $P_{\hat c}$ plus a sum of
concave deficits, so (i) holds summand by summand, and (ii) follows from the
monotonicity and joint convexity of the kernel as for any concave,
continuous, non-negative argument, with attainment from compactness. On a
fully dropped joint sector only the
linear $P_{\hat c}$ read the block, which gives (iii). For (iv) the proof of
Theorem~\ref{thm:oracle} goes through with the kept summands in place of
all of them; we list the changes.

(i) $P_{\hat c}$ is linear. Each $f_{\hat c,(t,s)}$ is concave on its cone
(Lemma~\ref{lem:joint-sectors}\,(ii)), and
$\omega\mapsto\omega^{(t,s)}$ is linear, so \eqref{eq:KD} is a sum of concave
functions. Continuity: $X\mapsto X^{1/\alpha}$ and $Y\mapsto Y^{\alpha}$ are
continuous on the positive semidefinite cone by the continuous functional
calculus, $\cZ$, $\cG_{\hat c,s}$ and the trace are linear, so each
$f_{\hat c,(t,s)}$ of \eqref{eq:f-def} is continuous.

(ii) Write $O^{\mathcal{D}}_{\alpha}$ as the sum over $\hat c$ of
$\frac{\alpha}{\alpha-1}d(q(\hat c),P_{\hat c}(\omega))$ for
$\hat c\notin\Chkey$ and of
$d(q(\hat c),P_{\hat c}(\omega))+\frac{1}{\alpha-1}
d(q(\hat c),K^{\mathcal{D}}_{\hat c}(\omega))$ for $\hat c\in\Chkey$, with
$d$ the kernel of \xref{foundations/lem:kernel}. Let $g$ be concave, continuous and
non-negative on a convex set, as $P_{\hat c}$ and, by (i) and
Proposition~\ref{prop:drop-bound}\,(i), $K^{\mathcal{D}}_{\hat c}$ are on the
convex set $\Feas$. For a convex combination with weights $\lambda$,
$1-\lambda$: $g(\omega_{\lambda})\ge\lambda g(\omega_1)+(1-\lambda)g(\omega_2)$,
so, $d$ being non-increasing in its second argument and jointly convex,
$d(q_{\lambda},g(\omega_{\lambda}))\le
d\big(q_{\lambda},\lambda g(\omega_1)+(1-\lambda)g(\omega_2)\big)\le
\lambda d(q_1,g(\omega_1))+(1-\lambda)d(q_2,g(\omega_2))$. For lower
semicontinuity let $(q_n,\omega_n)\to(q,\omega)$; then
$g(\omega_n)\to g(\omega)$ and lower semicontinuity of $d$ gives
$\liminf_nd(q_n,g(\omega_n))\ge d(q,g(\omega))$. Sums with positive
coefficients preserve both properties. Non-negativity is
Proposition~\ref{prop:drop-bound}\,(ii). $\acSet\times\Feas$ is compact
(Definition~\ref{def:admissible-acceptance-set}, Definition~\ref{def:feasible}), and a lower
semicontinuous function with values in $(-\infty,+\infty]$ attains its
infimum on a non-empty compact set.

(iii) If $(t,s)$ is fully dropped, no $f_{\hat c,(t,s)}$ occurs in any
$K^{\mathcal{D}}_{\hat c}$, so the block enters
$O^{\mathcal{D}}_{\alpha}$ through the values $P_{\hat c}(\omega)$ alone,
and its contribution to $P_{\hat c}$ is
$P_{\hat c,(t,s)}(\omega^{(t,s)})
=\tr[\Gamma^{(t,s)}_{\hat c}\omega^{(t,s)}]$ by \eqref{eq:Phimap}, a sum of
the stated numbers. With every relevant argument positive,
$\frac{\partial}{\partial v}\,u\log_2\frac uv=-\log_2\!\mathrm{e}\,\frac uv$,
and the directional derivative of $P_{\hat c}$ and, the deficits of the
block being absent, of $K^{\mathcal{D}}_{\hat c}$ along $\Delta$ is
$\tr[\Gamma^{(t,s)}_{\hat c}\Delta]$ in both cases. Collecting the prefactors
of (ii) gives \eqref{eq:drop-gradient}; terms with $q(\hat c)=0$ are
identically zero.

(iv) We list the changes to the proof of Theorem~\ref{thm:oracle}, which is
written for $\mathcal{D}=\varnothing$ and the index $s$. Write
$\sigma_{\hat c,(t,s)}(\omega)\coloneqq\cG_{\hat c,(t,s)}(\omega^{(t,s)})$,
an operator of trace at most $1$ for $\omega\in\Feas$, and
$e^{\mathcal{D}}_{\hat c}\coloneqq\frac{1}{1-\alpha}\log_2
\frac{K^{\mathcal{D}}_{\hat c}}{P_{\hat c}}$, which is $\ge0$ by
Proposition~\ref{prop:drop-bound}\,(i). By \eqref{eq:KD}, \eqref{eq:f-def}
and $B_{\sigma}(1)=\tr\sigma$,
\begin{equation*}
  K^{\mathcal{D}}_{\hat c}(\omega)-P_{\hat c}(\omega)
  =\sum_{(t,s):\,(\hat c,(t,s))\notin\mathcal{D}}
  \Big(B_{\sigma_{\hat c,(t,s)}(\omega)}(\alpha)
  -B_{\sigma_{\hat c,(t,s)}(\omega)}(1)\Big),
\end{equation*}
which replaces the identity for $K_{\hat c}-P_{\hat c}$ used there.
\emph{Upper bound:} with $\mathbf q=\Phimap[\omega]$ the divergence terms
vanish and $O^{\mathcal{D}}_{\alpha}(\mathbf q,\omega)
=\sum_{\hat c\in\Chkey}P_{\hat c}\,e^{\mathcal{D}}_{\hat c}$; the limit is
taken summand by summand with Lemma~\ref{lem:b-derivative}, exactly as there,
and gives the sum of relative entropies over the kept pairs; pairs that are
key-inert contribute zero because their $\sigma$ is block diagonal for
$\cZ$ (Lemma~\ref{lem:joint-sectors}\,(iii) and its proof), which restricts
the sum to $\mathbb{J}_K(\hat c)$. \emph{Lower bound:} the decomposition of
the objective into the divergence part and
$\sum_{\hat c}q(\hat c)e^{\mathcal{D}}_{\hat c}$ is the identity
\eqref{eq:rearrange} of real numbers, valid for any $K\in(0,P]$; both parts
are non-negative; compactness of $\acSet\times\Feas$, lower semicontinuity of
the divergence and the conclusion $\omega^{*}\in\Feas_1$ do not involve $K$;
$u_n\coloneqq(K^{\mathcal{D}}_{\hat c}-P_{\hat c})/P_{\hat c}\in[-1,0]$ by
Proposition~\ref{prop:drop-bound}\,(i); and the mean-value step with the
majorant \eqref{eq:majorant} and the joint continuity of $c_{\sigma}(\alpha)$
is applied to each kept summand, the maps
$\omega\mapsto\sigma_{\hat c,(t,s)}(\omega)$ being continuous. The resulting
$\liminf$ is the right-hand side of \eqref{eq:drop-oracle} at $\omega^{*}$.
\end{proof}

% claim meat-core/prop:minimal-divergence
\claimsummary{meat-core/prop:minimal-divergence}{Dropping every pair leaves \ensuremath{\alpha}/(\ensuremath{\alpha}\ensuremath{-}1) times the minimal divergence between accepted distributions and feasible statistics, and h\_low lies between that and it plus log\ensuremath{_{2}}(d\_Z+1) times the key-live mass. A positive lower bound proves that no feasible state has accepted statistics.}
\claimbookmark{meat-core/prop:minimal-divergence}{Proposition}{prop:minimal-divergence}{The minimal divergence as the value with every pair dropped}
\claimleadin{%
The \termref{term-objective}{objective} is a divergence penalty plus one entropy term per
\termref{term-key-live-acceptance-value}{key-live acceptance value}, and only the entropy terms are bounded
(Proposition~\ref{prop:value-range}); an infinite value is the extreme case
(Remark~\ref{rem:hlow-range}). Dropping every pair removes the entropy terms
altogether, so the dropped program then measures the penalty alone: how far
the accepted distributions lie from the statistics of the feasible set.
}
\begin{proposition}[The minimal divergence as the value with every pair dropped]\label{prop:minimal-divergence}
Let $\Decl$ satisfy Definition~\ref{def:admissible-declaration} in the
general format, in either source mode and with an arbitrary list $\CList$,
let $\acSet$ satisfy Definition~\ref{def:admissible-acceptance-set}, and let
$\alpha\in(1,\infty)$. Let
$\mathcal{D}_{\mathrm{all}}\coloneqq\Chkey\times(\mathbb{T}\times\SecSet)$
be the dropped set of every pair (Definition~\ref{def:dropped}), and put
\begin{equation}\label{eq:dmin}
  \sfD_{\min}\coloneqq\inf_{\mathbf q\in\acSet}\ \inf_{\omega\in\Feas}\
  \sfD\big(\mathbf q\,\big\|\,\Phimap[\omega]\big)\in[0,+\infty],
  \qquad
  m_{\mathrm{key}}\coloneqq\sup_{\mathbf q\in\acSet}\
  \sum_{\hat c\in\Chkey}q(\hat c)\in[0,1],
\end{equation}
with $\inf\varnothing=+\infty$ and $\sup\varnothing=0$.
\begin{enumerate}[label=(\roman*),leftmargin=2.2em]
\item \emph{(The value with every pair dropped.)} For every distribution
  $\mathbf q$ on $\Chatset$ and every $\omega\in\Feas$,
  $O^{\mathcal{D}_{\mathrm{all}}}_{\alpha}(\mathbf q,\omega)
  =\frac{\alpha}{\alpha-1}\,\sfD(\mathbf q\|\Phimap[\omega])$; hence
  $\underline{h}^{\mathcal{D}_{\mathrm{all}}}_{\alpha}
  =\frac{\alpha}{\alpha-1}\,\sfD_{\min}$.
\item \emph{(The value between its floor and its ceiling.)}
  \begin{equation}\label{eq:dmin-sandwich}
    \frac{\alpha}{\alpha-1}\,\sfD_{\min}\;\le\;\hlow\;\le\;
    \frac{\alpha}{\alpha-1}\,\sfD_{\min}+\log_2(d_Z+1)\,m_{\mathrm{key}} .
  \end{equation}
\item \emph{(When the floor is positive.)} $\sfD_{\min}=0$ if and only if
  $\mathbf q=\Phimap[\omega]$ for some $\mathbf q\in\acSet$ and some
  $\omega\in\Feas$. Equivalently, $\sfD_{\min}>0$ if and only if no
  distribution of $\acSet$ is the statistics $\Phimap[\omega]$ of a state
  $\omega\in\Feas$.
\item \emph{(Bounds at one order serve every order.)} Let $b$ be a real
  number with $b\le\underline{h}^{\mathcal{D}_{\mathrm{all}}}_{\alpha}$, and
  let $(\mathbf q,\omega)\in\acSet\times\Feas$. Then
  \begin{equation}\label{eq:dmin-bracket}
    \frac{\alpha-1}{\alpha}\,b\;\le\;\sfD_{\min}\;\le\;
    \frac{\alpha-1}{\alpha}\,
    O^{\mathcal{D}_{\mathrm{all}}}_{\alpha}(\mathbf q,\omega),
  \end{equation}
  and for every $\alpha'\in(1,\infty)$ the value $\underline{h}_{\alpha'}$
  of \eqref{eq:hlow-def} at the order $\alpha'$ satisfies
  $\underline{h}_{\alpha'}\ge\frac{\alpha'}{\alpha'-1}\cdot
  \frac{\alpha-1}{\alpha}\,b$. If $b>0$, then $\sfD_{\min}>0$, and no
  distribution of $\acSet$ is the statistics of a state of $\Feas$.
  For every $\alpha'\ge\alpha$,
  \begin{equation}\label{eq:dmin-ceiling}
    \frac{\alpha'}{\alpha'-1}\,\sfD_{\min}\;\le\;
    \frac{\alpha'}{\alpha'-1}\cdot\frac{\alpha-1}{\alpha}\,
    O^{\mathcal{D}_{\mathrm{all}}}_{\alpha}(\mathbf q,\omega)\;\le\;
    O^{\mathcal{D}_{\mathrm{all}}}_{\alpha}(\mathbf q,\omega),
  \end{equation}
  so that $\underline{h}_{\alpha'}\le
  \frac{\alpha'}{\alpha'-1}\cdot\frac{\alpha-1}{\alpha}\,
  O^{\mathcal{D}_{\mathrm{all}}}_{\alpha}(\mathbf q,\omega)
  +\log_2(d_Z+1)\,m_{\mathrm{key}}$.
\end{enumerate}
\end{proposition}
\claimreading{The key-rate program charges a penalty whenever the accepted statistics differ from the per-round statistics of every state the declared devices admit. This result isolates that penalty. If some such state has statistics the acceptance test accepts, the program's value is at most the entropy ceiling, log\ensuremath{_{2}}(d\_Z+1) times the largest key-live probability an accepted distribution gives. If none has, the value has a floor that grows without limit as the R\'enyi order approaches one, so a key length can grow with the mismatch rather than with the secrecy of the key. A positive lower bound on the penalty-only program, which the engine computes with its usual solver, proves that no such state has accepted per-round statistics. It does not prove that a finite run of those devices fails the test: its frequencies can still fall in the acceptance set.}
\claimmeta{\metalabel{Uses}\mref{Def.}{def:admissible-declaration}, \mref{Def.}{def:maps}, \mref{Def.}{def:admissible-acceptance-set}, \mref{Def.}{def:joint-sectors}, \mref{Def.}{def:feasible}, \mref{Prop.}{prop:value-range}, \mref{Def.}{def:dropped}, \mref{Prop.}{prop:drop-bound}, \mref{Cor.}{cor:drop-program}, \extref{../foundations/contract.pdf}{foundations/def:kl}
\metasep \metalabel{Used by}\mref{Prop.}{prop:model-pass-bound}, \mref{Req.}{req:unreachable-acceptance-set}
\metasep newly derived
\metasep implemented (effective: accepted)
\metasep \metalabel{Code}\coderef{https://github.com/fraunhofersg/certiqs-key/blob/v0.1.0/src/certiqs_key/meat_core/finitekey/consistency.py}{certiqs_key.meat_core.finitekey.consistency:proven_floor}}

\begin{proof}
\emph{Idea.} With every pair dropped, $K$ equals $P$ and the two prefactors
of a key-live term add up to that of the other terms, so the objective is
the scaled divergence. The entropy terms are non-negative and each at most
$\log_2(d_Z+1)$, which places $\hlow$ between the scaled minimal divergence
and that plus the entropy ceiling. A divergence vanishes only between equal
distributions, and the minimum is attained.

Throughout, $\omega\in\Feas$ has $\tr\omega=1$ (Definition~\ref{def:feasible}),
so $\Phimap[\omega]$ of \eqref{eq:Phimap} is a probability distribution on
$\Chatset$ whose entry at $\hat c$ is $P_{\hat c}(\omega)$, as in the proof of
Proposition~\ref{prop:drop-bound}\,(ii).

\emph{(i).} Every pair of $\Chkey\times(\mathbb{T}\times\SecSet)$ lies in
$\mathcal{D}_{\mathrm{all}}$, so the sum in \eqref{eq:KD} is empty and
$K^{\mathcal{D}_{\mathrm{all}}}_{\hat c}(\omega)=P_{\hat c}(\omega)$ for every
$\hat c\in\Chkey$. Fix $\hat c$ and write $q=q(\hat c)$, $P=P_{\hat c}(\omega)$.
For $\hat c\notin\Chkey$ the term of \eqref{eq:objective} is
$\frac{\alpha}{\alpha-1}q\log_2\frac qP$. For $\hat c\in\Chkey$, with $K$
replaced by $P$, it is $q\log_2\frac qP+\frac{q}{\alpha-1}\log_2\frac qP$:
if $q=0$ both summands are $0$ by the convention
$0\log_2\frac0v=0$ of \xref{foundations/def:kl}, as is
$\frac{\alpha}{\alpha-1}q\log_2\frac qP$; if $q>0$ and $P=0$ both are
$+\infty$, as is $\frac{\alpha}{\alpha-1}q\log_2\frac qP$; if $q,P>0$ they
add up to $\big(1+\frac{1}{\alpha-1}\big)q\log_2\frac qP
=\frac{\alpha}{\alpha-1}q\log_2\frac qP$. Summing over $\hat c\in\Chatset$
gives $O^{\mathcal{D}_{\mathrm{all}}}_{\alpha}(\mathbf q,\omega)
=\frac{\alpha}{\alpha-1}\sum_{\hat c}q(\hat c)\log_2
\frac{q(\hat c)}{P_{\hat c}(\omega)}
=\frac{\alpha}{\alpha-1}\sfD(\mathbf q\|\Phimap[\omega])$ in the extended-real
sense. Taking the infimum over $\acSet\times\Feas$ in \eqref{eq:hlowD}, with
the positive factor $\frac{\alpha}{\alpha-1}$ taken out (and
$\frac{\alpha}{\alpha-1}\cdot(+\infty)=+\infty$), gives
$\underline{h}^{\mathcal{D}_{\mathrm{all}}}_{\alpha}
=\frac{\alpha}{\alpha-1}\sfD_{\min}$.

\emph{(ii), the lower inequality.} By
Proposition~\ref{prop:drop-bound}\,(iii),
$\underline{h}^{\mathcal{D}_{\mathrm{all}}}_{\alpha}\le\hlow$, and by (i) the
left-hand side is $\frac{\alpha}{\alpha-1}\sfD_{\min}$.

\emph{(ii), the upper inequality.} If $\sfD_{\min}=+\infty$ the right-hand
side is $+\infty$. Otherwise let $(\mathbf q,\omega)\in\acSet\times\Feas$ with
$\sfD(\mathbf q\|\Phimap[\omega])<+\infty$; then $q(\hat c)>0$ implies
$P_{\hat c}(\omega)>0$. For $\hat c\in\Chkey$ with $q(\hat c)>0$ let
$e_{\hat c}(\omega)\coloneqq\frac{1}{1-\alpha}\log_2
\frac{K_{\hat c}(\omega)}{P_{\hat c}(\omega)}$; by
Proposition~\ref{prop:value-range}\,(ii), $P_{\hat c}(\omega)>0$ gives
$e_{\hat c}(\omega)\le\log_2(d_Z+1)$, so $e_{\hat c}(\omega)$ is a real
number and $K_{\hat c}(\omega)>0$. For such $\hat c$, with $q=q(\hat c)$,
$P=P_{\hat c}(\omega)$, $K=K_{\hat c}(\omega)$, the term of
\eqref{eq:objective} is
$q\log_2\frac qP+\frac{q}{\alpha-1}\big(\log_2\frac qP+\log_2\frac PK\big)
=\frac{\alpha}{\alpha-1}q\log_2\frac qP+q\,e_{\hat c}(\omega)$; for
$\hat c\in\Chkey$ with $q(\hat c)=0$ the term is $0
=\frac{\alpha}{\alpha-1}q(\hat c)\log_2\frac{q(\hat c)}{P_{\hat c}(\omega)}$.
Summing,
\begin{equation*}
  O_{\alpha}(\mathbf q,\omega)
  =\frac{\alpha}{\alpha-1}\sfD(\mathbf q\|\Phimap[\omega])
  +\sum_{\hat c\in\Chkey:\,q(\hat c)>0}q(\hat c)\,e_{\hat c}(\omega)
  \le\frac{\alpha}{\alpha-1}\sfD(\mathbf q\|\Phimap[\omega])
  +\log_2(d_Z+1)\,m_{\mathrm{key}} ,
\end{equation*}
using $\sum_{\hat c\in\Chkey}q(\hat c)\le m_{\mathrm{key}}$. For every
$\varepsilon>0$ there is such a pair with
$\sfD(\mathbf q\|\Phimap[\omega])\le\sfD_{\min}+\varepsilon$, so
$\hlow\le O_{\alpha}(\mathbf q,\omega)\le\frac{\alpha}{\alpha-1}
(\sfD_{\min}+\varepsilon)+\log_2(d_Z+1)\,m_{\mathrm{key}}$; letting
$\varepsilon\to0$ gives \eqref{eq:dmin-sandwich}.

\emph{(iii).} If $\mathbf q=\Phimap[\omega]$ with $\mathbf q\in\acSet$ and
$\omega\in\Feas$, then $\sfD(\mathbf q\|\Phimap[\omega])=0$, and since
$\sfD_{\min}\ge0$ ($\underline{h}^{\mathcal{D}_{\mathrm{all}}}_{\alpha}\ge0$ by
Proposition~\ref{prop:drop-bound}\,(iii), and (i)), $\sfD_{\min}=0$.
Conversely let $\sfD_{\min}=0$. Then
$\underline{h}^{\mathcal{D}_{\mathrm{all}}}_{\alpha}=0$ by (i) is finite, so by
Corollary~\ref{cor:drop-program}\,(ii) it is attained at some
$(\mathbf q,\omega)\in\acSet\times\Feas$, where by (i)
$\sfD(\mathbf q\|\mathbf p)=0$ with $\mathbf p\coloneqq\Phimap[\omega]$. Let
$S\coloneqq\{\hat c:q(\hat c)>0\}$. The divergence being finite,
$p(\hat c)>0$ on $S$, and
$0=\sfD(\mathbf q\|\mathbf p)=-\sum_{\hat c\in S}q(\hat c)\log_2
\frac{p(\hat c)}{q(\hat c)}\ge-\log_2\sum_{\hat c\in S}p(\hat c)\ge0$, the
first inequality by concavity of $\log_2$ (Jensen's inequality for the
weights $q(\hat c)$, $\hat c\in S$, which sum to $1$), the second because
$\sum_{\hat c\in S}p(\hat c)\le1$. Both are equalities. The second gives
$\sum_{\hat c\in S}p(\hat c)=1$, so $p(\hat c)=0$ off $S$. In the first,
$\log_2$ is strictly concave, so equality forces
$\frac{p(\hat c)}{q(\hat c)}$ to be one constant $\lambda$ on $S$; summing
$p(\hat c)=\lambda q(\hat c)$ over $S$ gives $\lambda=1$. Hence
$\mathbf q=\mathbf p=\Phimap[\omega]$. The second sentence is the
contrapositive of the first, $\sfD_{\min}$ lying in $[0,+\infty]$.

\emph{(iv).} By (i), $b\le\frac{\alpha}{\alpha-1}\sfD_{\min}$, which is the
lower inequality of \eqref{eq:dmin-bracket}; by (i) again,
$\frac{\alpha-1}{\alpha}O^{\mathcal{D}_{\mathrm{all}}}_{\alpha}(\mathbf q,\omega)
=\sfD(\mathbf q\|\Phimap[\omega])\ge\sfD_{\min}$, the upper one. The lower
inequality of \eqref{eq:dmin-sandwich} at the order $\alpha'$, whose
$\sfD_{\min}$ and $\Feas$ do not depend on the order, gives
$\underline{h}_{\alpha'}\ge\frac{\alpha'}{\alpha'-1}\sfD_{\min}
\ge\frac{\alpha'}{\alpha'-1}\cdot\frac{\alpha-1}{\alpha}\,b$. If $b>0$, the
lower inequality of \eqref{eq:dmin-bracket} gives $\sfD_{\min}>0$, and (iii)
the assertion that follows it. For \eqref{eq:dmin-ceiling}, multiply the upper
inequality of \eqref{eq:dmin-bracket} by $\frac{\alpha'}{\alpha'-1}>0$; the
second inequality holds because $t\mapsto\frac{t}{t-1}=1+\frac{1}{t-1}$ is
decreasing on $(1,\infty)$, so $\frac{\alpha'}{\alpha'-1}\le
\frac{\alpha}{\alpha-1}$ for $\alpha'\ge\alpha$, and
$O^{\mathcal{D}_{\mathrm{all}}}_{\alpha}(\mathbf q,\omega)\ge0$
(Proposition~\ref{prop:drop-bound}\,(ii)). The last bound on
$\underline{h}_{\alpha'}$ is the upper inequality of
\eqref{eq:dmin-sandwich} at the order $\alpha'$ with the first inequality of
\eqref{eq:dmin-ceiling}.
\end{proof}

% claim meat-core/lem:objective-homogeneous
\claimsummary{meat-core/lem:objective-homogeneous}{Read on the whole cone, the objective, its Choi form and the dropped objective are positively homogeneous of degree one: scaling the acceptance vector and the state by c \textgreater{} 0 scales the value by c. This is what lets a first-order bracket be formed without evaluating the objective.}
\claimbookmark{meat-core/lem:objective-homogeneous}{Lemma}{lem:objective-homogeneous}{Positive homogeneity of the objective}
\claimleadin{%
The objective is defined on the acceptance set and the feasible set, but its
formula makes sense on the whole cone, and there it scales linearly: the
divergence terms are jointly homogeneous of degree one in their two
arguments, the statistics are linear, and the pinched trace power is
homogeneous of degree one because $p\alpha=1$. Differentiated along the point
itself, this says that the gradient paired with the point is the value at the
point, which is how the Frank--Wolfe submodule forms a bracket from its two
linear minima alone (Lemma~\ref{lem:certified-bracket}).
}
\begin{technicalclaim}
\begin{lemma}[Positive homogeneity of the objective]\label{lem:objective-homogeneous}
Let $\alpha\in(1,\infty)$ and let $\Decl$ be a tuple of the form of
Definition~\ref{def:declaration}, in the basic or the general format, that
satisfies conditions (i)--(vi) and (viii) of
Definition~\ref{def:admissible-declaration}; its maps are read below. For
$\mathbf q\in[0,\infty)^{\Chatset}$ and
$\omega\in\Pos{A\Qsys}$ block diagonal along the sectors, along the
\exttermref{../protocol-compiler/contract.pdf}{term-joint-sector}{joint sectors} in the general format of
Definition~\ref{def:declaration}, let $O_{\alpha}(\mathbf q,\omega)$ be the
right-hand side of \eqref{eq:objective}, with the conventions of
\xref{foundations/def:kl} at vanishing arguments, and for a
\termref{term-dropped-set}{dropped set} $\mathcal{D}$ let $K^{\mathcal{D}}_{\hat c}$ and
$O^{\mathcal{D}}_{\alpha}$ be those of Definition~\ref{def:dropped} at the
same arguments. These are defined on all such $(\mathbf q,\omega)$, with
values in $(-\infty,+\infty]$: there $P_{\hat c}(\omega)\ge0$ and
$0\le K_{\hat c}(\omega)\le K^{\mathcal{D}}_{\hat c}(\omega)$. Let $c>0$,
and read $c\cdot(+\infty)\coloneqq+\infty$.
\begin{enumerate}[label=(\roman*),leftmargin=2.4em]
\item $P_{\hat c}(c\omega)=c\,P_{\hat c}(\omega)$ for every
  $\hat c\in\Chatset$; every sector deficit and every joint-sector deficit
  satisfies $f(c\sigma)=c\,f(\sigma)$ on its cone; and
  $K_{\hat c}(c\omega)=c\,K_{\hat c}(\omega)$ and
  $K^{\mathcal{D}}_{\hat c}(c\omega)=c\,K^{\mathcal{D}}_{\hat c}(\omega)$ for
  every $\hat c\in\Chkey$.
\item $O_{\alpha}(c\mathbf q,c\omega)=c\,O_{\alpha}(\mathbf q,\omega)$ and
  $O^{\mathcal{D}}_{\alpha}(c\mathbf q,c\omega)
  =c\,O^{\mathcal{D}}_{\alpha}(\mathbf q,\omega)$.
\item In \exttermref{../transmitter/contract.pdf}{term-trusted-source-mode}{trusted-source mode}, with $\chi_{\Qsys}$ the assignment of
  Theorem~\ref{thm:choi} for a \termref{term-support-frame}{support frame}, the Choi forms
  $O^{J}_{\alpha}(\mathbf q,J)\coloneqq O_{\alpha}(\mathbf q,\chi_{\Qsys}(J))$
  and $O^{J,\mathcal{D}}_{\alpha}(\mathbf q,J)\coloneqq
  O^{\mathcal{D}}_{\alpha}(\mathbf q,\chi_{\Qsys}(J))$ satisfy
  $O^{J}_{\alpha}(c\mathbf q,cJ)=c\,O^{J}_{\alpha}(\mathbf q,J)$ and the same
  for $O^{J,\mathcal{D}}_{\alpha}$, at every $\mathbf q\in[0,\infty)^{\Chatset}$
  and every $J$ whose blocks are positive semidefinite; in
  \exttermref{../transmitter/contract.pdf}{term-untrusted-source-mode}{untrusted-source mode} the Choi form is the state form
  (Definition~\ref{def:untrusted-choi}) and (ii) is the statement.
\end{enumerate}
\end{lemma}
\end{technicalclaim}
\claimmeta{\metalabel{Uses}\mref{Def.}{def:declaration}, \mref{Def.}{def:admissible-declaration}, \mref{Def.}{def:maps}, \mref{Thm.}{thm:branch-formula}, \mref{Lem.}{lem:cone-concavity}, \mref{Cor.}{cor:f-nonpos}, \mref{Def.}{def:joint-sectors}, \mref{Lem.}{lem:joint-sectors}, \mref{Def.}{def:feasible}, \mref{Thm.}{thm:choi}, \mref{Def.}{def:untrusted-choi}, \mref{Def.}{def:dropped}, \extref{../foundations/contract.pdf}{foundations/def:kl}
\metasep \metalabel{Used by}\mref{Lem.}{lem:certified-bracket}
\metasep newly derived
\metasep accepted}

\begin{proof}
\emph{The hypothesis.} The facts cited below, the cone clause of
Theorem~\ref{thm:branch-formula}, Lemma~\ref{lem:joint-sectors}\,(ii),
Corollary~\ref{cor:f-nonpos} and Lemma~\ref{lem:cone-concavity}\,(a), are
stated for \termref{term-admissible-declaration}{admissible declarations} or for the
cone of a branch space, and
their proofs read the declaration's maps, positivity and block structure
only, never condition~(vii) of Definition~\ref{def:admissible-declaration},
the non-empty feasible set; so they hold of $\Decl$.

\emph{The domain.} $P_{\hat c}(\omega)=\Phimap[\omega](\hat c)$ is linear in
$\omega$ (\eqref{eq:Phimap}) and non-negative on $\Pos{A\Qsys}$, being a sum
of traces of products of positive semidefinite operators. The cone clause of
Theorem~\ref{thm:branch-formula} gives $0\le K_{\hat c}(\omega)$ on the
whole cone, and in the general format so does Lemma~\ref{lem:joint-sectors}\,(ii)
through \eqref{eq:K-joint}, each $P_{\hat c,(t,s)}+f_{\hat c,(t,s)}$ being
non-negative. $K^{\mathcal{D}}_{\hat c}$ is $K_{\hat c}$ with the deficits
of the pairs of $\mathcal{D}$ removed (\eqref{eq:KD}, \eqref{eq:K-joint}), and
every deficit is non-positive (Corollary~\ref{cor:f-nonpos},
Lemma~\ref{lem:joint-sectors}\,(ii)), so $K^{\mathcal{D}}_{\hat c}\ge
K_{\hat c}\ge0$. So every term $u\log_2\frac uv$ of \eqref{eq:objective} has
$u\ge0$ and $v\ge0$, and by \xref{foundations/def:kl} it is $0$ when $u=0$,
$+\infty$ when $u>0=v$, and real otherwise; a sum of finitely many such terms
with positive coefficients lies in $(-\infty,+\infty]$.

\emph{(i).} Linearity of $\Phimap$ gives the first claim. A sector deficit is
$f_{\hat c,s}(\sigma)=T(\cG_{\hat c,s}(\sigma))-\tr\cG_{\hat c,s}(\sigma)$
(\eqref{eq:f-def}) with $T$ the pinched trace power of \eqref{eq:Tcone};
$\cG_{\hat c,s}$ and the trace are linear, and $T(cX)=c\,T(X)$ by
Lemma~\ref{lem:cone-concavity}\,(a), so $f_{\hat c,s}(c\sigma)
=c\,f_{\hat c,s}(\sigma)$. A joint-sector deficit is a restriction of a
sector deficit (\eqref{eq:joint-deficit}), and
Lemma~\ref{lem:joint-sectors}\,(ii) states its homogeneity as well. By
\eqref{eq:fK},
$K_{\hat c}(\omega)=T(\cG_{\hat c}(\omega))+P_{\hat c}(\omega)
-\tr\cG_{\hat c}(\omega)$, each term homogeneous by the same two facts; in
the general format $K_{\hat c}$ and $K^{\mathcal{D}}_{\hat c}$ are $P_{\hat c}$
plus a sum of joint-sector deficits (\eqref{eq:K-joint}, \eqref{eq:KD}),
each term homogeneous.

\emph{(ii).} For $u,v\ge0$ and $c>0$,
$(cu)\log_2\frac{cu}{cv}=c\cdot u\log_2\frac uv$ under the conventions of
\xref{foundations/def:kl}: both sides are $0$ when $u=0$, both are $+\infty$
when $u>0=v$, and otherwise the logarithm is unchanged by the common factor.
Each term of \eqref{eq:objective} is such an expression, multiplied by a
positive constant ($\frac{\alpha}{\alpha-1}$, $1$ or $\frac1{\alpha-1}$),
with $u=q(\hat c)$ and $v=P_{\hat c}(\omega)$ or $v=K_{\hat c}(\omega)$, and
by (i) the pair $(u,v)$ at $(c\mathbf q,c\omega)$ is $c$ times the pair at
$(\mathbf q,\omega)$; the same holds with $K^{\mathcal{D}}_{\hat c}$. For
$c>0$ and $a,b\in(-\infty,+\infty]$, $c(a+b)=ca+cb$, so the finite sum scales
by $c$.

\emph{(iii).} $\chi_{\Qsys}$ is linear and maps a list of positive
semidefinite blocks to a block-diagonal positive semidefinite operator
(Theorem~\ref{thm:choi}), so
$O^{J}_{\alpha}(c\mathbf q,cJ)=O_{\alpha}(c\mathbf q,c\,\chi_{\Qsys}(J))
=c\,O_{\alpha}(\mathbf q,\chi_{\Qsys}(J))$ by (ii), and likewise for
$O^{J,\mathcal{D}}_{\alpha}$. In untrusted-source mode
$\chi_{\Qsys}$ is the identity (Definition~\ref{def:untrusted-choi}).
\end{proof}

\section{Key length}\label{group-meat-core/def:postprocessing-model}
\sectionintro{The accumulated entropy feeds a \termref{term-fixed-length-post-processing-model}{fixed-length post-processing model} with a fixed-length error-correction transcript, error verification and privacy amplification by ideal two-universal hashing, every message public (\hyperlink{claim-meat-core/def:postprocessing-model}{Definition~\ref*{def:postprocessing-model}}). The section's main results are \hyperlink{claim-meat-core/thm:keylength}{Theorem~\ref*{thm:keylength}}, an \ensuremath{\varepsilon}\_PA-secret key of the stated length on the accepted state for a R\'enyi order in (1, 2), and \hyperlink{claim-meat-core/cor:keylength-bound}{Corollary~\ref*{cor:keylength-bound}}, which lets any lower bound h \ensuremath{\le} h\_low take the place of h\_low through max\{0, h\} and is where every solver obtains its key length. Error verification makes the key \ensuremath{\varepsilon}\_EV-correct (\hyperlink{claim-meat-core/lem:ev-correctness}{Lemma~\ref*{lem:ev-correctness}}), the two compose to (\ensuremath{\varepsilon}\_PA + \ensuremath{\varepsilon}\_EV)-security with honest authentication, and over a real authentication scheme the protocol is (\ensuremath{\varepsilon}\_PA + \ensuremath{\varepsilon}\_EV + 2\ensuremath{\varepsilon}\_auth)-secure with net key gain \ensuremath{\ell} \ensuremath{-} k\_auth (\hyperlink{claim-meat-core/cor:composed-security}{Corollary~\ref*{cor:composed-security}}), all of it against an adversary whose memory has any dimension, by a limit over finite-dimensional ones (\hyperlink{claim-meat-core/lem:memory-lift}{Lemma~\ref*{lem:memory-lift}}) and every adversary of the \termref{term-honest-authentication-setting}{honest-authentication setting} read as such an attack (\hyperlink{claim-meat-core/lem:honest-attacks}{Lemma~\ref*{lem:honest-attacks}}); the parameter bounds one execution, separate executions that share no key add their parameters (\hyperlink{claim-meat-core/cor:separate-executions}{Corollary~\ref*{cor:separate-executions}}), and a reused authentication key enters through \ensuremath{\varepsilon}\_auth, which the earlier execution's parameter does not bound by itself (\hyperlink{claim-meat-core/rem:repeated-executions}{Remark~\ref*{rem:repeated-executions}}). A run whose rounds follow statistics of feasible states passes the acceptance test with probability at most 2\textasciicircum{}(\ensuremath{-}n D\_min), D\_min the minimal divergence (\hyperlink{claim-meat-core/prop:model-pass-bound}{Proposition~\ref*{prop:model-pass-bound}}). The reported \termref{term-maximal-key-length}{maximal key length} is that bound floored and clamped at 0, and more leakage or a tighter error budget never lengthens it (\hyperlink{claim-meat-core/prop:monotone}{Proposition~\ref*{prop:monotone}}). The constrained value enters under the attack-set licence (\hyperlink{claim-meat-core/cor:constrained-key-length}{Corollary~\ref*{cor:constrained-key-length}}), behind a two-part gate whose licensed-class half is vacuous in untrusted-source mode and trivially avoidable at d\_A = 1 (\hyperlink{claim-meat-core/rem:cgate}{Remark~\ref*{rem:cgate}}).}
\sectionimports{\extref{../foundations/contract.pdf}{foundations/def:kl}, \extref{../foundations/contract.pdf}{foundations/def:renyi}, \extref{../foundations/contract.pdf}{foundations/lem:kernel}, \extref{../foundations/contract.pdf}{foundations/lem:register-compression}, \extref{../reductions/contract.pdf}{reductions/def:protocol}}

% claim meat-core/def:postprocessing-model
\claimsummary{meat-core/def:postprocessing-model}{After the n rounds the key holder's raw string passes an acceptance test, one-way error correction with a padded fixed-length transcript, error verification with an announced flag and privacy amplification, each message public and delivered unaltered. Secrecy and correctness are stated in it.}
\claimbookmark{meat-core/def:postprocessing-model}{Definition}{def:postprocessing-model}{Fixed-length post-processing model}
\claimleadin{%
The model is the fixed-length specialisation of the generic post-processing
of \cite[Sec.~4.1, Protocol~1, steps~2--5]{tupkary2026decoy}: the key length
is either $\ell$ or zero, error correction is one-way with a transcript of
fixed length, and the verification outcome is announced. Every message is
public, and the model is stated in the setting in which authentication
behaves honestly \cite[Sec.~3.2]{tupkary2026decoy}; how a run over a real
authentication scheme reduces to it is a separate step.
}
\begin{definition}[Fixed-length post-processing model]
\label{def:postprocessing-model}
Fix a declaration $\Decl$ with key holder $\hold$
(Definition~\ref{def:declaration}), a round count $n$, an acceptance set
$\acSet$, an output length $\ell$, a transcript length $\lEC$ and error
budgets $\epa,\eev\in(0,1]$. A hash family is a set of functions with a
probability law on it \cite[Def.~3.5]{tupkary2026decoy}. The
\termdef{term-fixed-length-post-processing-model}{fixed-length post-processing model}
is the following classical processing, applied after the $n$
\termref{term-round-channel}{round channels} of Definition~\ref{def:roundmodel}.
The raw string $Z_1^n$ is the content of the classical registers
$Z_1,\dots,Z_n$ of the round channels, which the single-round protocol channel
$\Mprot$ retains (Theorem~\ref{thm:protocol-channel}); the key holder computes
it from its own outcome record and the announcements $I_1^n$ through the
function $r_{\hold}$ of \eqref{eq:partycomp}, and we call the other party the
partner. The acceptance values $\Chat_1^n$ are announced. The processing
appends, in order:
\begin{enumerate}[label=(\roman*),leftmargin=2.4em]
\item the \emph{acceptance test}: the event
  $\OmAT\coloneqq\{\operatorname{freq}_{\hat c_1^n}\in\acSet\}$, a function
  of $\Chat_1^n$;
\item \emph{error correction}, one-way from the holder: a public seed
  $S_{\mathrm{pub}}$, drawn with the law the error-correction scheme
  prescribes, and the public transcript
  $L_{\mathrm{EC}}=g_{\mathrm{EC}}(Z_1^n,S_{\mathrm{pub}})$, which occupies
  a fixed-length register of $\lEC$ bits: a shorter transcript is padded to
  length $\lEC$, so that $\dim(L_{\mathrm{EC}})=2^{\lEC}$ exactly; from
  $L_{\mathrm{EC}}$ and its own data the partner forms a guess $\hat Z_1^n$
  of $Z_1^n$;
\item \emph{error verification}: a public seed $S_{\mathrm{EV}}$ selects a
  function $h_{\mathrm{EV}}$, with the law of the error-verification family,
  from a hash family of functions with
  $\lceil\log_2\frac1{\eev}\rceil$-bit values; the holder announces
  $T_{\mathrm{EV}}\coloneqq h_{\mathrm{EV}}(Z_1^n)$, the partner compares it
  with $h_{\mathrm{EV}}(\hat Z_1^n)$ and announces the one-bit outcome
  $F_{\mathrm{acc}}\in\{\mathrm{acc},\mathrm{rej}\}$, which is
  $\mathrm{acc}$ exactly when the two values agree;
\item \emph{privacy amplification}: a function $G$ is drawn, with the law of
  the privacy-amplification family, from a hash family of functions with
  $\ell$-bit values, and announced; the holder outputs
  $K\coloneqq G(Z_1^n)$ and the partner $\hat K\coloneqq G(\hat Z_1^n)$.
\end{enumerate}
We write $K_A$ and $K_B$ for Alice's and Bob's outputs, which are $K$ and
$\hat K$ in the order the key holder fixes, and
\begin{equation}\label{eq:pp-registers}
  S\coloneqq(S_{\mathrm{pub}},S_{\mathrm{EV}}),\qquad
  L\coloneqq(L_{\mathrm{EC}},T_{\mathrm{EV}},F_{\mathrm{acc}}),\qquad
  E_{\mathrm{tot}}\coloneqq\Chat_1^nI_1^n\Ehat\,S\,L\,G,
\end{equation}
where $L$ has at most
$\lambda\coloneqq\lEC+\lceil\log_2\frac1{\eev}\rceil+1$ bits and $\Ehat$ is
the adversary's retained quantum memory of Definition~\ref{def:roundmodel}.
The protocol \emph{accepts} on
$\Omega_{\mathrm{acc}}\coloneqq\OmAT\wedge\{F_{\mathrm{acc}}=\mathrm{acc}\}$;
on its complement both parties abort and output the symbol $\perp$ in place
of a key. The model is stated in the
\termdef{term-honest-authentication-setting}{honest-authentication setting}:
every classical message of the protocol, the round announcements included,
reaches its recipient exactly as it was sent, in the order sent and at some
time after it was sent, the adversary receives a copy of it, and no message
fails authentication.
\end{definition}
\claimmeta{\metalabel{Uses}\mref{Def.}{def:declaration}, \mref{Def.}{def:roundmodel}
\metasep \metalabel{Used by}\mref{Def.}{def:authenticated-channel}, \mref{Def.}{def:execution-end}, \mref{Lem.}{lem:ev-correctness}, \mref{Lem.}{lem:memory-lift}, \mref{Lem.}{lem:honest-attacks}, \mref{Thm.}{thm:keylength}, \mref{Cor.}{cor:composed-security}, \mref{Rem.}{rem:completeness-scope}, \mref{Prop.}{prop:model-pass-bound}, \mref{Ass.}{ass:keyreg}, \mref{Ass.}{ass:hash}, \mref{Ass.}{ass:seed}, \mref{Ass.}{ass:fixed}, \mref{Ass.}{ass:conform}, \mref{Ass.}{ass:lab}
\metasep newly derived
\metasep accepted}

% claim meat-core/def:authenticated-channel
\claimsummary{meat-core/def:authenticated-channel}{The authenticated classical channel delivers each message unaltered or as an authentication failure, never before it was sent, and a scheme realises it within \ensuremath{\varepsilon} from a pre-shared key. The authentication post-processing makes both parties abort when either saw a failure.}
\claimbookmark{meat-core/def:authenticated-channel}{Definition}{def:authenticated-channel}{Authenticated classical channel and authentication post-processing}
\claimleadin{%
A real authentication scheme cannot deliver every message: an adversary who
tampers with a message makes its recipient reject it. The channel below is
the functionality such a scheme can approach
\cite[Sec.~5, Eqs.~(37)--(38)]{tupkary2026decoy}, the approximation is
measured by the distinguishing advantage of the composable framework
\cite[Sec.~II.E, Def.~1; Sec.~VII.A, Eq.~(36)]{portmann2022security}, and
the authentication post-processing is step~6 of
\cite[Sec.~4.1, Protocol~1]{tupkary2026decoy} together with the rule of
that section's opening paragraphs that a party receiving an authentication
failure aborts.
}
\begin{definition}[Authenticated classical channel and authentication post-processing]
\label{def:authenticated-channel}
Label each classical message of a party by its position $i$ in that party's
sequence of sent messages, and the messages a party receives by their
position in its sequence of received messages; times are global times, which
the parties need not know.
\begin{enumerate}[label=(\alph*),leftmargin=2.4em]
\item The
  \termdef{term-authenticated-classical-channel}{authenticated classical channel}
  between the two parties gives the adversary a copy of every message sent
  over it and lets her choose when each message is delivered, subject to:
  (A1) if the $i$-th message a party receives arrives before the other
  party sent its $i$-th message, the received symbol is
  $\mathrm{auth\text{-}abort}$; (A2) otherwise the received symbol is either
  the $i$-th sent message, unaltered, or $\mathrm{auth\text{-}abort}$.
\item An \emph{authentication scheme} is a pair of procedures, one per party,
  that use a pre-shared key and an unauthenticated classical channel, on
  which the adversary may read, alter, inject, delay and suppress messages.
  It \emph{$\varepsilon$-realises} the authenticated classical channel if
  there is a simulator at the adversary's interface such that the
  distinguishing advantage between the scheme run on the key and the
  unauthenticated channel, and the authenticated classical channel with that
  simulator, is at most $\varepsilon$
  \cite[Sec.~II.E, Def.~1]{portmann2022security}. Its \emph{key consumption}
  is the number of bits of the pre-shared key it uses in one run of the
  protocol.
\item The
  \termdef{term-authentication-post-processing}{authentication post-processing}
  appended to the \termref{term-fixed-length-post-processing-model}{fixed-length post-processing model}
  (Definition~\ref{def:postprocessing-model}) is: a party that receives
  $\mathrm{auth\text{-}abort}$ for any message aborts, replaces its key by
  $\perp$ and sends $\mathrm{abort}$ in place of each later message; after
  privacy amplification the partner sends $\mathrm{accept}$ if it holds a key
  and $\mathrm{abort}$ otherwise; the key holder sends $\mathrm{accept}$ if
  it received $\mathrm{accept}$ and holds a key, and otherwise sends
  $\mathrm{abort}$ and replaces its key by $\perp$; the partner keeps its key
  if it receives $\mathrm{accept}$ and otherwise replaces it by $\perp$.
\end{enumerate}
\end{definition}
\claimmeta{\metalabel{Uses}\mref{Def.}{def:postprocessing-model}
\metasep \metalabel{Used by}\mref{Def.}{def:execution-end}, \mref{Cor.}{cor:composed-security}, \mref{Rem.}{rem:completeness-scope}, \mref{Ass.}{ass:fixed}, \mref{Ass.}{ass:conform}, \mref{Ass.}{ass:auth}, \mref{Ass.}{ass:lab}
\metasep newly derived
\metasep accepted}

% claim meat-core/def:execution-end
\claimsummary{meat-core/def:execution-end}{A party has ended an execution once it has sent and received every message the post-processing prescribes to it and fixed its output, a key or an abort; the execution has ended once both parties have. Several executions are compared at that time.}
\claimbookmark{meat-core/def:execution-end}{Definition}{def:execution-end}{The end of an execution}
\claimleadin{%
Under the authentication post-processing the two parties do not finish at
the same time: the key holder finishes when it sends its last decision, the
partner when that decision is delivered, at a time the adversary chooses.
}
\begin{technicalclaim}
\begin{definition}[The end of an execution]\label{def:execution-end}
An \emph{execution} of the protocol is one run of its $n$ rounds, the
\termref{term-fixed-length-post-processing-model}{fixed-length post-processing model}
(Definition~\ref{def:postprocessing-model}) and the
\termref{term-authentication-post-processing}{authentication post-processing}
(Definition~\ref{def:authenticated-channel}\,(c)). A party has ended an
execution once it has sent and received every message that these prescribe
to it and has fixed its output, its key or $\perp$: the key holder once it
has sent its final $\mathrm{accept}$ or $\mathrm{abort}$, the partner once
that message, or $\mathrm{auth\text{-}abort}$ in its place, has been
delivered to it. The
\termdef{term-end-of-an-execution}{end of an execution}
is the time at which both of its parties have ended it.
\end{definition}
\end{technicalclaim}
\claimmeta{\metalabel{Uses}\mref{Def.}{def:postprocessing-model}, \mref{Def.}{def:authenticated-channel}
\metasep \metalabel{Used by}\mref{Cor.}{cor:separate-executions}, \mref{Ass.}{ass:conform}, \mref{Ass.}{ass:separate-executions}
\metasep newly derived
\metasep accepted}

% claim meat-core/lem:ev-correctness
\claimsummary{meat-core/lem:ev-correctness}{In the fixed-length post-processing model the two parties accept with different keys with probability at most \ensuremath{\varepsilon}\_EV, for every attack and every output length. It is the correctness half of the protocol's security.}
\claimbookmark{meat-core/lem:ev-correctness}{Lemma}{lem:ev-correctness}{Correctness from error verification}
\claimleadin{%
The argument reads only the verification step of the model: two different
raw strings pass it only if their short hashes collide, and the hash is
chosen after, and independently of, the strings it compares.
}
\begin{lemma}[Correctness from error verification]\label{lem:ev-correctness}
Let the protocol run the \termref{term-fixed-length-post-processing-model}{fixed-length post-processing model}
(Definition~\ref{def:postprocessing-model}) in the
\termref{term-honest-authentication-setting}{honest-authentication setting}, let Assumptions~\ref{ass:hash},
\ref{ass:seed} and~\ref{ass:conform} hold, and let $\eev\in(0,1]$. Then for every attack and
every output length $\ell$ the protocol is $\eev$-correct on the accepted
state, in the sense of \cite[Sec.~2.3.1, Def.~2.2, Eq.~(4)]{tupkary2025survey}:
\begin{equation}\label{eq:ev-correct}
  \Pr\big[\Omega_{\mathrm{acc}}\wedge K_A\ne K_B\big]\le\eev .
\end{equation}
\end{lemma}
\claimreading{This is the correctness half of security. Before the key is formed, the key holder announces a \ensuremath{\lceil}log\ensuremath{_{2}}(1/\ensuremath{\varepsilon}\_EV)\ensuremath{\rceil}-bit hash of its raw key, computed with a hash function drawn at random independently of the raw keys, and the partner compares it with the hash of its own corrected copy; a mismatch aborts the run. Two different raw keys pass only through a hash collision, so the chance that both parties accept yet hold different keys is at most \ensuremath{\varepsilon}\_EV, whatever the attacker does and whatever key length you choose. It needs a two-universal hash family (\hyperlink{claim-meat-core/ass:hash}{A-HASH}), seeds independent of the raw keys (\hyperlink{claim-meat-core/ass:seed}{A-SEED}) and an implementation that aborts as specified (\hyperlink{claim-meat-core/ass:conform}{A-CONFORM}), and takes messages to arrive unaltered; \hyperlink{claim-meat-core/cor:composed-security}{Corollary~\ref*{cor:composed-security}} adds real authentication. It does not promise that error correction succeeds, only that a failure is caught, except with probability \ensuremath{\varepsilon}\_EV.}
\claimmeta{\metalabel{Uses}\mref{Def.}{def:postprocessing-model}
\metasep \metalabel{Assumes}\aref{A-HASH}{ass:hash}, \aref{A-SEED}{ass:seed}, \aref{A-CONFORM}{ass:conform}
\metasep \metalabel{Used by}\mref{Cor.}{cor:composed-security}, \mref{Cor.}{cor:separate-executions}, \mref{Cor.}{cor:keylength-bound}
\metasep newly derived
\metasep accepted}

\begin{proof}
Both parties abort when the verification flag is $\mathrm{rej}$ and output
the keys the model computes (Assumption~\ref{ass:conform}\,(iii)--(iv)). On
$\Omega_{\mathrm{acc}}\wedge\{K_A\ne K_B\}$ the hashed keys differ, so the
raw strings differ, yet their $\lceil\log_2\frac1{\eev}\rceil$-bit
two-universal hashes agreed. (If the raw strings agree, $K_A=K_B$ because
the same public $G$ is applied: in the honest-authentication setting the
partner receives the holder's $G$ unaltered.) The verification function
$h_{\mathrm{EV}}$ is selected by $S_{\mathrm{EV}}$, which
Assumption~\ref{ass:seed} makes independent of the pair
$(Z_1^n,\hat Z_1^n)$, and its family is two-universal
(Assumption~\ref{ass:hash}\,(ii)). Conditioned on any pair of distinct
strings, the probability that their hashes agree is therefore at most
$2^{-\lceil\log_2(1/\eev)\rceil}\le\eev$, and averaging over the pair gives
\eqref{eq:ev-correct}.
\end{proof}

% claim meat-core/lem:memory-lift
\claimsummary{meat-core/lem:memory-lift}{Every attack on a run, its memory of any dimension and its signals taken in any order, is approximated in trace norm by attacks with finite memory that take every signal at the start. A trace-distance bound that holds for those therefore holds for every attack with the same \ensuremath{\varepsilon}.}
\claimbookmark{meat-core/lem:memory-lift}{Lemma}{lem:memory-lift}{Approximating an attack by all-at-once attacks with finite memory}
\claimleadin{%
The bound of Theorem~\ref{thm:keylength} is proved for the attacks the
\exttermref{../foundations/contract.pdf}{term-round-model}{round model} of Definition~\ref{def:roundmodel} represents: their
registers are finite-dimensional, as every register of that model is, and in
\exttermref{../transmitter/contract.pdf}{term-trusted-source-mode}{trusted-source mode} their first channel receives every transmitted
signal, the type the reductions transfer
(\xref{reductions/lem:source-transfer}) and the timing the published analysis
moves to \cite[Lemma~8.1]{tupkary2026decoy}. A physical adversary may keep a
memory of any dimension and receives the signals as the source releases
them; where a signal lives in an infinite-dimensional space, a finite memory
cannot hold it for a later round. The lemma closes that gap by a limit, the
route of \cite[Def.~8.2, Sec.~8.2, Thm.~A.1, Rem.~A.2]{tupkary2026decoy}:
compressing the adversary's registers, and each signal released after the
first round, onto finite-dimensional subspaces
(\xref{foundations/lem:register-compression}) turns every attack into one of
the analysed type whose output state is as close as wanted, and a bound on a
trace distance, which no dimension enters, passes to the limit.
Corollary~\ref{cor:composed-security} applies it. The compression touches
neither the \exttermref{../reductions/contract.pdf}{term-measured-system}{measured system} nor the \exttermref{../transmitter/contract.pdf}{term-transmitted-system}{transmitted system} in
transit, so no finite dimension of either is needed; that is why the limit is
taken at the physical run rather than at the analysed protocol. Unlike the
transfers of the reductions module, which hold for every criterion because
they are exact (\xref{reductions/rem:why-exact}), the limit needs a criterion
continuous in the output state, and the trace distances of the security
statement are.
}
\begin{lemma}[Approximating an attack by all-at-once attacks with finite memory]\label{lem:memory-lift}
Read a run of the protocol and its attacks as follows.
\begin{enumerate}[label=(\alph*),leftmargin=2.4em]
\item \emph{(the run)} Its quantum phase is an $n$-round protocol of
  \xref{reductions/def:protocol}, whose measured system $\Bsys_j$ of round $j$
  is separable and of any dimension, followed by the
  \termref{term-fixed-length-post-processing-model}{fixed-length post-processing model} of
  Definition~\ref{def:postprocessing-model} in the
  \termref{term-honest-authentication-setting}{honest-authentication setting}. In trusted-source mode its source
  releases, for $i=1,\dots,n$, a separable transmitted system
  $\Asig_i$ of any dimension, prepared together with Alice's letter register
  $X_i$ in a state $\rho_i$ on $X_i\Asig_i$ that depends on nothing else, and
  the run reads $X_i$ from round $i$ on; in \exttermref{../transmitter/contract.pdf}{term-untrusted-source-mode}{untrusted-source mode} there
  are no $X_i$ and no $\Asig_i$. Every system of the run other than the
  $\Bsys_j$ and the $\Asig_i$ is finite-dimensional, and its public messages
  are classical registers it keeps.
\item \emph{(attacks)} An \emph{attack} consists of a release schedule
  $0=k_0\le k_1\le\dots\le k_n=n$ with $k_j\ge j$ for every $j$, a state
  $\zeta$ on the adversary's initial register $E_0$, and channels, in the
  sense of \xref{foundations/lem:register-compression},
  \begin{equation}\label{eq:lift-attack}
    \Aev_j\in
    \begin{cases}
      \mathrm{CPTP}\big(\Asig_{k_{j-1}+1}\cdots\Asig_{k_j}E_{j-1},\
        \Bsys_jE'_j\big), & \text{trusted-source mode},\\[0.3ex]
      \mathrm{CPTP}\big(E_{j-1},\ A_j\Bsys_jE'_j\big),
        & \text{untrusted-source mode},
    \end{cases}
  \end{equation}
  for $j=1,\dots,n$, where $E'_j$ is the adversary's memory after round $j$,
  $E_j$ is $E'_j$ together with a copy of the public messages of round $j$,
  and $\Hs{E_0}$ and every $\Hs{E'_j}$ are Hilbert spaces of any dimension.
  Under the attack the run starts from $\zeta$ and its own registers, prepares
  $\rho_i$ in round $j$ for $k_{j-1}<i\le k_j$ before $\Aev_j$ acts, takes
  $\Bsys_j$, and in untrusted-source mode Alice's system $A_j$, from the
  output of $\Aev_j$, and places a copy of the public messages of round $j$ in
  $E_j$. Its \emph{output state} is the state of the run's
  \exttermref{../reductions/contract.pdf}{term-output-registers}{output registers}, its public messages and $E'_n$. The attack is
  \emph{all-at-once} if $k_1=n$, and has \emph{finite registers} if $E_0$ and
  every $E'_j$ are finite-dimensional.
\end{enumerate}
Then:
\begin{enumerate}[label=(\roman*),leftmargin=2.4em]
\item \emph{(approximation)} for every attack and every $\eta>0$ there is an
  all-at-once attack with finite registers whose memory after round $n$ is a
  subspace of $\Hs{E'_n}$, such that the output states $\rho$ and
  $\rho^{\circ}$ of the two attacks satisfy
  $\|\iota(\rho^{\circ})-\rho\|_1\le\eta$, with $\iota$ the inclusion of that
  subspace into $\Hs{E'_n}$, acting as the identity on every other register;
\item \emph{(trace-distance criteria)} let $L_1$ and $L_2$ be completely
  positive maps, not increasing the trace, from the run's output registers
  and public messages to a finite-dimensional system, put $L\coloneqq
  L_1-L_2$, and for a state $\sigma$ of those registers and the adversary's
  memory after round $n$ put
  $f(\sigma)\coloneqq\frac12\|(L\otimes\idmap)(\sigma)\|_1$; if
  $f(\rho^{\circ})\le\varepsilon$ at the output state $\rho^{\circ}$ of every
  all-at-once attack with finite registers, then $f(\rho)\le\varepsilon$ at
  the output state $\rho$ of every attack;
\item \emph{(processing by the adversary)} let, in (ii), each $L_t$ keep the
  public messages and act on the other registers given their value $m$,
  $L_t(\dyad{m}\otimes x)=\dyad{m}\otimes L_{t,m}(x)$, and let
  $\Phi=\sum_m\dyad{m}\otimes\Phi_m$ act on a state that is classical on the
  public messages by a channel $\Phi_m$ on the adversary's register, given
  their value $m$, keeping them. Then $f((\idmap\otimes\Phi)(\sigma))\le
  f(\sigma)$ for every state $\sigma$ classical on the public messages.
\end{enumerate}
\end{lemma}
\claimmeta{\metalabel{Uses}\mref{Def.}{def:postprocessing-model}, \extref{../foundations/contract.pdf}{foundations/lem:register-compression}, \extref{../reductions/contract.pdf}{reductions/def:protocol}
\metasep \metalabel{Used by}\mref{Lem.}{lem:honest-attacks}, \mref{Thm.}{thm:keylength}, \mref{Cor.}{cor:composed-security}, \mref{Cor.}{cor:keylength-bound}, \mref{Cor.}{cor:fo-keylength}
\metasep newly derived
\metasep accepted}

\begin{proof}
\emph{Idea.} Written as a sequence of channels, the run is compressed at
finitely many points: the adversary's initial register, her memory after each
of her channels, and every signal released after the first round, as soon as
it is released. The compressions cost at most $\eta$ in trace norm
(\xref{foundations/lem:register-compression}\,(ii), (iii)), and the compressed
run is the run under an all-at-once attack with finite registers: the
adversary takes every signal at the start, compresses those she would have
received later and keeps them, compressed, in her memory until their round.
A trace distance of the output state moves by at most the distance of the
states, which gives (ii); a channel on the adversary's side commutes past the
criterion's maps and does not increase a trace norm, which gives (iii).

\emph{Step 1 (the run as a sequence of channels).} Fix an attack. The run's
own registers are finite-dimensional and start in fixed states, which we take
as written by the channels that first write them. Round $j$ is the composite
of three channels: the preparation $\mathcal{Q}_j$, which appends
$\bigotimes_{k_{j-1}<i\le k_j}\rho_i$ (the identity in untrusted-source mode);
$\Aev_j$, extended by the identity on the run's registers; and the round map
$\Grnd_j$ of \xref{reductions/def:protocol}, extended by the identity on
$E'_j$. After round $n$ the run applies $\Kfin$, which includes the
post-processing, extended by the identity on $E'_n$. These extensions are
channels acting on products as the named map tensored with the identity. A
channel extended by the identity on a finite-dimensional register applies it
to every entry of a block matrix, which is a channel by complete positivity.
$\Grnd_j$ is the \exttermref{../reductions/contract.pdf}{term-measurement-channel}{measurement channel} of the measured system followed
by processing of finite-dimensional registers; the measurement channel
$\tau\mapsto\sum_v\tr[M_v\tau]\dyad{v}$ extends to
$\tau\mapsto\sum_v\dyad{v}\otimes\tr_{\Bsys_j}\big[(\sqrt{M_v}\otimes\idon{})
\tau(\sqrt{M_v}\otimes\idon{})\big]$, a sum of conjugations followed by
partial traces and the appending of states, and a completely positive map
$\Psi$ between finite-dimensional systems is a finite sum
$\sum_tK_t(\cdot)K_t^{\dagger}$, which extends to
$\sum_t(K_t\otimes\idon{})(\cdot)(K_t\otimes\idon{})^{\dagger}$. The
output state is the result of these channels applied in order to $\zeta$.

\emph{Step 2 (the compressions).} The compression points are: $E_0$ before
round one; $E'_j$ right after $\Aev_j$, for every $j$; and, for every $i>k_1$,
the transmitted system $\Asig_i$ right after the preparation $\mathcal{Q}_j$
that appends $\rho_i$, before $\Aev_j$ receives it, one signal at a time with
identity channels between. There are at most $2n+1$ of them. At each point,
\xref{foundations/lem:register-compression}\,(ii), applied to the run's state
there under the given attack, gives a non-zero finite-dimensional subspace of
the register at which the bound of that part, $2\sqrt{\delta}+\delta$, is at
most $\eta/(2n+1)$; we write $E_0^{\circ}$, $E_j^{\prime\circ}$ and
${\Asig_i}^{\circ}$ for them. Inserting at each point the compression onto its
subspace, and the identity at every other point, part (iii) of that lemma
gives $\|\rho^{\mathcal{C}}-\rho\|_1\le\eta$ for the output
$\rho^{\mathcal{C}}$ of the run so compressed.

\emph{Step 3 (the compressed run under an attack with finite registers).}
Write each compression as $\iota\circ\mathcal{P}$
(\xref{foundations/lem:register-compression}\,(i)). Each inclusion $\iota$ of
$E'_{j}$ commutes with the maps that follow it up to $\Aev_{j+1}$, the round
map $\Grnd_j$, the preparations of round $j+1$ and the compressions of the
signals they release, all of which act on other registers: for the channels
this is part (iv) of that lemma, and for an appended state $\rho$ it is
$\iota(y)\otimes\rho=(\iota\otimes\idmap)(y\otimes\rho)$, the product form of
part (i). Likewise the inclusion of $E_0$ moves to $\Aev_1$, the inclusion of
each ${\Asig_i}^{\circ}$ to the channel $\Aev_j$ that receives $\Asig_i$, and
the inclusion of $E_n^{\prime\circ}$ past $\Grnd_n$ and $\Kfin$ to the end.
Each projection $\mathcal{P}$ stays where it acts: after $\Aev_j$ it joins
it, and after the preparation of $\rho_i$ it gives
$\rho_i^{\circ}\coloneqq(\idmap_{X_i}\otimes\mathcal{P})(\rho_i)$, by the
product form $\mathcal{P}_{R_1R_2}(x\otimes y)=x\otimes\mathcal{P}_{R_2}(y)$
of part (i) with $R_2=X_i$ and $R_1$ the other systems. So
$\rho^{\mathcal{C}}=\iota(\rho^{\mathrm{fs}})$, where $\rho^{\mathrm{fs}}$ is
the output state of the run that prepares $\rho_i$ for $i\le k_1$ and
$\rho_i^{\circ}$ for $i>k_1$, on $X_i{\Asig_i}^{\circ}$, in the round of the
schedule, against the initial state $\mathcal{P}(\zeta)$ on $E_0^{\circ}$ and
the channels
\begin{equation*}
  \Aev^{\circ}_1\coloneqq\mathcal{P}\circ\Aev_1\circ\iota,
  \qquad
  \Aev^{\circ}_j\coloneqq\mathcal{P}\circ\Aev_j\circ\iota\ \ (j\ge2),
\end{equation*}
the projection being that of $E'_j$, and the inclusion that of $E_0$ for
$j=1$ and of $E'_{j-1}$ and every ${\Asig_i}^{\circ}$ received in round $j$ for
$j\ge2$, each acting as the identity on the other registers. Each
$\Aev^{\circ}_j$ is a channel between finite-dimensional adversary's
registers, with the signals of the schedule and $\Bsys_j$, and $A_j$ in
untrusted-source mode, as before.

\emph{Step 4 (an all-at-once attack with the same output).} Put
$E^{\star}_j\coloneqq E_j^{\prime\circ}{\Asig_{k_j+1}}^{\circ}\cdots
{\Asig_n}^{\circ}$, finite-dimensional, so that $E^{\star}_n=E_n^{\prime\circ}$.
The all-at-once attack has the initial state $\mathcal{P}(\zeta)$, the first
channel
\begin{equation*}
  \mathcal{A}^{\star}_1\coloneqq
  \big(\Aev^{\circ}_1\otimes\idmap_{{\Asig_{k_1+1}}^{\circ}\cdots{\Asig_n}^{\circ}}\big)
  \circ\mathcal{P}_{\Asig_{k_1+1}}\circ\dots\circ\mathcal{P}_{\Asig_n}
  \in\mathrm{CPTP}\big(\Asig_1\cdots\Asig_nE_0^{\circ},\ \Bsys_1E^{\star}_1\big),
\end{equation*}
each $\mathcal{P}_{\Asig_i}$ the projection of $\Asig_i$ onto
${\Asig_i}^{\circ}$ with the identity on the other systems, and, for $j\ge2$,
$\mathcal{A}^{\star}_j\coloneqq\Aev^{\circ}_j\otimes
\idmap_{{\Asig_{k_j+1}}^{\circ}\cdots{\Asig_n}^{\circ}}$, from $E^{\star}_{j-1}$
with the public messages of round $j-1$ to $\Bsys_jE^{\star}_j$; the input
of $\Aev^{\circ}_j$ is a factor of that register. In untrusted-source mode it
is the attack of Step~3 itself. Its run prepares
$\rho_1\otimes\dots\otimes\rho_n$ at the start, and $\mathcal{A}^{\star}_1$
turns each factor $\rho_i$ with $i>k_1$ into $\rho_i^{\circ}$, by the product
form of part (i). From then on no map acts on $X_i{\Asig_i}^{\circ}$ before
the round that reads $X_i$, which is not earlier than the round of the
schedule since $k_i\ge i$, or before the channel that receives
${\Asig_i}^{\circ}$; every map before them acts on other registers, and a map
$\Psi$ extended by the identity on the finite-dimensional
$X_i{\Asig_i}^{\circ}$ takes $y\otimes\rho_i^{\circ}$ to
$\Psi(y)\otimes\rho_i^{\circ}$. Appending $\rho_i^{\circ}$ at the start or in
its round of the schedule therefore leaves every later state the same, and
the output state $\rho^{\circ}$ of the all-at-once attack is
$\rho^{\mathrm{fs}}$. With Steps~2 and~3,
$\|\iota(\rho^{\circ})-\rho\|_1=\|\rho^{\mathcal{C}}-\rho\|_1\le\eta$, which
is (i).

\emph{Step 5 (part (ii)).} The maps $L_1\otimes\idmap$ and $L_2\otimes\idmap$
are completely positive and do not increase the trace, by the Kraus form of
Step~1, so each contracts the trace norm of a Hermitian operator: for
$Z=Z^{+}-Z^{-}$, $\|(L_t\otimes\idmap)(Z)\|_1\le\tr(L_t\otimes\idmap)(Z^{+})
+\tr(L_t\otimes\idmap)(Z^{-})\le\|Z\|_1$. Hence for two states $\sigma$,
$\sigma'$,
\begin{equation*}
  |f(\sigma)-f(\sigma')|\le\tfrac12\big\|(L\otimes\idmap)(\sigma-\sigma')\big\|_1
  \le\tfrac12\sum_{t=1,2}\big\|(L_t\otimes\idmap)(\sigma-\sigma')\big\|_1
  \le\|\sigma-\sigma'\|_1 .
\end{equation*}
Given an attack and $\eta>0$, take the all-at-once attack of (i). By part
(iv) of \xref{foundations/lem:register-compression}, with $G=L$,
$(L\otimes\idmap_{E'_n})\circ\iota=\iota\circ(L\otimes\idmap_{E_n^{\prime\circ}})$,
and $\iota$ preserves the trace norm, so $f(\iota(\rho^{\circ}))=f(\rho^{\circ})
\le\varepsilon$. So $f(\rho)\le\varepsilon+\eta$ for every $\eta>0$.

\emph{Step 6 (part (iii)).} Write $\sigma=\sum_m\dyad{m}\otimes\sigma_m$, with
$\sigma_m$ positive on the other registers and the adversary's. Then
\begin{equation*}
  (L_t\otimes\idmap)\big((\idmap\otimes\Phi)(\sigma)\big)
  =\sum_m\dyad{m}\otimes(L_{t,m}\otimes\Phi_m)(\sigma_m)
  =\Phi'\big((L_t\otimes\idmap)(\sigma)\big),
\end{equation*}
with $\Phi'\coloneqq\sum_m\dyad{m}\otimes\Phi_m$ acting on the output of
$L_t$, which keeps the public messages: both sides apply $L_{t,m}$ to the run's
registers and $\Phi_m$ to the adversary's in each block, which commute by
part (iv) of \xref{foundations/lem:register-compression}. So
$(L\otimes\idmap)(\idmap\otimes\Phi)(\sigma)=\Phi'((L\otimes\idmap)(\sigma))$,
and $\Phi'$ is a positive trace-preserving map on block-diagonal operators,
which contracts the trace norm of the Hermitian $(L\otimes\idmap)(\sigma)$ as
in Step~5. Halving gives the claim.
\end{proof}

% claim meat-core/lem:honest-attacks
\claimsummary{meat-core/lem:honest-attacks}{Whatever an adversary does in the honest-authentication setting, adaptively and with a memory of any dimension, the output state is that of an all-at-once attack of \hyperlink{claim-meat-core/lem:memory-lift}{Lemma~\ref*{lem:memory-lift}} followed by her own processing of her register.}
\claimbookmark{meat-core/lem:honest-attacks}{Lemma}{lem:honest-attacks}{Adversaries of the honest-authentication setting as round attacks}
\claimleadin{%
An adversary of the \termref{term-honest-authentication-setting}{honest-authentication setting} is not bound to
act once per round: she may act at any time, decide adaptively when to deliver
which system, hold the signals she receives, read her copies of the messages
when they reach her, and process her memory after the last round. The lemma
shows that none of this leaves the all-at-once attacks of
Lemma~\ref{lem:memory-lift}\,(b), up to a final processing of her own register,
as the timing modification of \cite[Lemma~8.1]{tupkary2026decoy} does for the
published analysis: she is handed every signal at the start and ignores each
until it would have reached her. That needs every party to measure in round
order, A-SEQ\,(v), the measurement part of the timing condition of
\cite[Sec.~8.1.1, Protocol~2]{tupkary2026decoy}, read for each measuring
party. Corollary~\ref{cor:composed-security} uses the
lemma to state its bound at every output state of the setting.
}
\begin{lemma}[Adversaries of the honest-authentication setting as round attacks]\label{lem:honest-attacks}
Let the run be as in Lemma~\ref{lem:memory-lift}\,(a), with A-ALICE-MEAS
(\xref{protocol-compiler/ass:alice-meas})\,(i) for Alice's measurement where
no \exttermref{../transmitter/contract.pdf}{term-source-replacement-certificate}{source-replacement certificate} is demanded, and let A-SEQ
(\xref{transmitter/ass:seq}), its clause (v) included,
Assumption~\ref{ass:lab} (A-LAB) at each of the two parties and
Assumption~\ref{ass:conform} (A-CONFORM) hold. Let an adversary act on the run in the honest-authentication setting
(Definition~\ref{def:postprocessing-model}) by any operations, at any times,
each chosen given the classical outcomes of her earlier ones, on her own
registers, of any dimension, on each transmitted system once the source has
released it and on her copies of the public messages once she holds them,
delivering the measured system of each round, and in untrusted-source mode
Alice's system, for the round's measurements. Then there are an all-at-once
attack in the sense of Lemma~\ref{lem:memory-lift}\,(b), with output state
$\rho$, and a channel $\Phi=\sum_m\dyad{m}\otimes\Phi_m$ that applies a
channel $\Phi_m$ to her register $E'_n$ given the value $m$ of the run's
public messages and keeps them, such that the state of the run's output
registers, its public messages and the adversary's final register is
$(\idmap\otimes\Phi)(\rho)$.
\end{lemma}
\claimmeta{\metalabel{Uses}\mref{Def.}{def:postprocessing-model}, \mref{Lem.}{lem:memory-lift}, \extref{../foundations/contract.pdf}{foundations/lem:register-compression}, \extref{../reductions/contract.pdf}{reductions/def:protocol}
\metasep \metalabel{Assumes}\extaref{../transmitter/contract.pdf}{transmitter/ass:seq}{A-SEQ}, \aref{A-LAB}{ass:lab}, \aref{A-CONFORM}{ass:conform}, \extaref{../protocol-compiler/contract.pdf}{protocol-compiler/ass:alice-meas}{A-ALICE-MEAS}
\metasep \metalabel{Used by}\mref{Cor.}{cor:composed-security}, \mref{Cor.}{cor:keylength-bound}, \mref{Cor.}{cor:fo-keylength}
\metasep newly derived
\metasep accepted}

\begin{proof}
\emph{Idea.} Every signal is handed to her at the start and passed on
untouched until she would have received it; her operations between the
deliveries that complete consecutive rounds, with their classical branches,
form one channel; every measurement of a round is made, in the
representation, right after the delivery that completes it; and neither the timing of message
delivery nor the time a signal is prepared changes anything.

\emph{Step 1 (what each side acts on).} By Assumption~\ref{ass:lab}\,(ii) at
each party, the adversary acts on a laboratory only through the systems
delivered into it, as its declared measurements receive them, and through the
public messages it is delivered, and by (i) she learns nothing of its records
except through the systems handed out of it and the public messages. So every
operation of hers acts on her registers, on transmitted systems the source has
released, on her copies of public messages, and on a system before she
delivers it; once delivered, a system is acted on only by the measurement
that receives it. Every operation of the run acts on its own registers and on
the systems delivered to it. In the honest-authentication setting every
classical message reaches its recipient exactly as sent, in the order sent and
after it was sent, and the adversary receives a copy. By
Assumption~\ref{ass:conform}\,(ii) the messages after the rounds are those of
Definition~\ref{def:postprocessing-model}, and by (i), (iii) and (iv) the run
applies that model's processing.

\emph{Step 2 (what the run does not depend on).} The run's quantum
operations, the preparations of the states $\rho_i$, which depend on nothing
else, the one measurement channel of every round
(\xref{reductions/def:protocol}) and, in untrusted-source mode, Alice's
measurement of $A_j$ in the round's processing, the same in every round and
on that round's $A_j$ alone by A-ALICE-MEAS\,(i), read no message. Its classical operations
compute functions of its records and of the contents of the messages it has
received, and a message arrives unaltered and in order. So the run's
registers at its end, and the content of every public message, are the same
whenever each message is delivered before the operation that reads it, which
the run waits for, and whenever each $\rho_i$ is prepared before its signal
is touched and its letter read.

\emph{Step 3 (the cuts $s_j$).} The measurements of round $j$ are those of
every party that measures in it, Alice's measurement of $A_j$ in
untrusted-source mode among them (A-SEQ\,(i)). For $j=1,\dots,n$ let $s_j$ be
the latest delivery, among the adversary's operations, of a system that a
measurement of a round $i\le j$ receives. It is an operation of hers, so it is
fixed by the branch of her classical outcomes: each of the finitely many
systems of the rounds $i\le j$ is delivered exactly once on every branch, so
at each of her operations her branch so far decides whether it is the last
such delivery. And $s_{j-1}$ precedes or is $s_j$. Every public message of a round $j'\ge j$ is sent after $s_j$: the
system delivered at $s_j$ is received by the measurement of a round $i\le j$
of some party $\mathsf P$, which completes after the delivery; by A-SEQ\,(v)
$\mathsf P$'s measurement of round $j'$ completes no earlier, $\mathsf P$
measuring in every round (\xref{reductions/def:protocol}); and by
A-SEQ\,(i) every message of round $j'$ follows that measurement. So a message
the adversary holds before $s_j$ is one of a round $j'<j$.

\emph{Step 4 (moving the measurements to the cuts).} Every measurement of
round $j$ receives its system at or before $s_j$, and its outcome is read by
nothing before the messages of round $j$, which follow $s_j$ (Step~3). By
Step~1 the adversary does not touch the system once delivered, and by
Assumption~\ref{ass:lab}\,(i) none of her operations depends on the outcome
before those messages; every operation of the run between the measurement and
$s_j$ acts on other systems and reads nothing it outputs. Moving the
measurement to just after $s_j$, whether it was made before or after,
therefore changes no later state, operations on disjoint systems commuting
(\xref{foundations/lem:register-compression}\,(iv)), with the extensions of
part (v) of that lemma. From now on every measurement of round $j$ is made
right after $s_j$, on the systems delivered for round $j$.

\emph{Step 5 (the attack).} Hand the adversary every transmitted system at the
start: by Step~2 preparing $\rho_i$ there changes nothing, and she passes
each on untouched, her operations extended by the identity on it
(\xref{foundations/lem:register-compression}\,(v)), until the operation at
which she received it. Let $\zeta$ be the state of her registers before the
run starts. For $j=1,\dots,n$, let $\Aev_j$ be the composite of her operations
after $s_{j-1}$ ($s_0$ the start) and up to $s_j$: on each branch of her
classical outcomes the composite of the operations on that branch, the record
of the branch kept in her memory, summed over the branches; it is a channel,
every operation being one and her branching being on recorded classical
outcomes. It acts on her memory, on the signals she holds and on her copies
of messages, which by Step~3 are of rounds before $j$, those she has not yet
received being passed on unread. The time at which a copy reaches her is set
by the run and may depend on its devices' timing; she records each arrival
as a classical outcome of hers, so it is part of her branch. By
Assumption~\ref{ass:lab}\,(ii) at each party, a party's devices, and so the
timing of its messages, are influenced by nothing but the systems and
messages delivered to it, whose delivery times are in her branch; any
further dependence on the quantum systems in its laboratory passes through a
measurement, whose outcome is a record. So, given her branch, the timing is
a function of the parties' records and local randomness, and by
Assumption~\ref{ass:lab}\,(i) at each party, the timing of the parties'
messages carrying no information about their records, the arrival times are
independent, given her branch, of every register other than her branch: her
quantum memory, the signals she holds and the parties' systems and records.
Their law given her branch is therefore a fixed function of that branch;
$\Aev_j$ samples them from that law with randomness of its own, which leaves
the joint law of every register unchanged. It outputs the measured systems of round
$j$, those she delivered before $s_{j-1}$ having been kept untouched in her
memory until then, Alice's system $A_j$ among them in untrusted-source mode
(Step~3), together
with her memory $E'_j$, which holds everything she still passes on. With
$k_1=n$, $(\zeta,(\Aev_j)_j)$ is an all-at-once attack in the sense of
Lemma~\ref{lem:memory-lift}\,(b). By A-SEQ\,(ii) the post-processing follows
the $n$ rounds; let $\Phi$ be the composite of her operations after $s_n$,
which deliver nothing and act on her register and on her copies of the public
messages;
the public messages are classical, so $\Phi$ applies, given their value $m$,
a channel $\Phi_m$ to her register and keeps them, the arrival of each copy
it reads being sampled inside $\Phi_m$ as inside the $\Aev_j$.

\emph{Step 6 (the same output).} The given execution, with its measurements
moved as in Step~4, and the run under the attack of Step~5 followed by $\Phi$
consist of the same channels, of the run and of the adversary. They differ in
the order of channels acting on disjoint systems, which commute
(\xref{foundations/lem:register-compression}\,(iv)); in the time at which
each message reaches the run and each $\rho_i$ is prepared, which changes
nothing by Step~2; and in the time at which the adversary holds a signal, a
delivered system or a copy, which the passing on of Step~5 leaves without
effect, the arrival of each copy being sampled inside the $\Aev_j$ and
$\Phi$ as Step~5 says. So the two output states agree.
\end{proof}

% claim meat-core/thm:keylength
\claimsummary{meat-core/thm:keylength}{For a R\'enyi order \ensuremath{\alpha} in (1, 2), a key of length \ensuremath{\ell} \ensuremath{\le} n\ensuremath{\cdot}h\_low \ensuremath{-} \ensuremath{\lambda}\_EC \ensuremath{-} \ensuremath{\lceil}log\ensuremath{_{2}}(1/\ensuremath{\varepsilon}\_EV)\ensuremath{\rceil} \ensuremath{-} \ensuremath{\alpha}/(\ensuremath{\alpha}\ensuremath{-}1)\ensuremath{\cdot}log\ensuremath{_{2}}(1/\ensuremath{\varepsilon}\_PA) + 1 is \ensuremath{\varepsilon}\_PA-secret on the accepted state at every all-at-once finite-register attack the reduction chain carries onto the round model (\hyperref{../protocol-compiler/contract.pdf}{}{claim-protocol-compiler/ass:link}{A-LINK}). \hyperlink{claim-meat-core/cor:composed-security}{Corollary~\ref*{cor:composed-security}} extends it to every attack.}
\claimbookmark{meat-core/thm:keylength}{Theorem}{thm:keylength}{Secrecy of the key from the accumulated entropy}
\claimleadin{%
The accumulated entropy bound feeds the fixed-length post-processing model:
error correction with a transcript of fixed length, error verification whose
one-bit outcome is announced, and privacy amplification by an ideal
two-universal hash. Announcing the verification outcome costs one bit and
makes the secrecy statement hold on the accepted state
(Remark~\ref{rem:plus2}).

The theorem is stated at the attacks the analysis is carried out for: those
with finite registers whose first channel receives every transmitted signal,
which the \exttermref{../reductions/contract.pdf}{term-reduction-chain}{reduction chain} carries onto the \exttermref{../foundations/contract.pdf}{term-round-model}{round model}, as
hypothesis (H-CHAIN) says. Its link part is proved, not attested, on the
identity link of \exttermref{../protocol-compiler/contract.pdf}{term-band-a}{Band A} (\xref{protocol-compiler/thm:band-a}) and for
a chain the reductions module emits with status exact
(\xref{reductions/thm:two-sided}, \xref{reductions/cor:sequential},
\xref{reductions/thm:licence}), where A-LINK is discharged
(\xref{protocol-compiler/rem:obligations-discharged}); for any other chain
it rests on A-LINK. With constraints, the part that makes the attacks
admissible for the constraint list rests on A-ATK on every route, Band A
included, wherever the list carries an entry no theorem licenses.
Corollary~\ref{cor:composed-security} extends the bound to every attack,
whatever the adversary's memory.
}
\begin{mainclaim}
\begin{theorem}[Secrecy of the key from the accumulated entropy]\label{thm:keylength}
Let the protocol run the \termref{term-fixed-length-post-processing-model}{fixed-length post-processing model}
(Definition~\ref{def:postprocessing-model}) in the
\termref{term-honest-authentication-setting}{honest-authentication setting}, with the registers
\eqref{eq:pp-registers}. Let one of the two accumulation branches hold:
\emph{(unconstrained)} the hypotheses of
Theorem~\ref{thm:accumulated-entropy} with $\CList=\varnothing$; or
\emph{(constrained)} the hypotheses of
Theorem~\ref{thm:constrained-accumulation} ($\CList$
\exttermref{../protocol-compiler/contract.pdf}{term-licensed-for-the-declared-source-mode}{licensed for the declared source mode}, Definition~\ref{def:homog};
A-ATK (\xref{protocol-compiler/ass:atk})), in either branch save the
condition $\Pr[\OmAT]>0$ that the accumulation theorem places on the
attack, which the attack below need not meet. Let the run be as in
Lemma~\ref{lem:memory-lift}\,(a), as A-SPEC (\xref{protocol-compiler/ass:spec})\,(i)
attests for its measurement, with A-ALICE-MEAS
(\xref{protocol-compiler/ass:alice-meas})\,(i) for Alice's measurement where
no \exttermref{../transmitter/contract.pdf}{term-source-replacement-certificate}{source-replacement certificate} is demanded, and A-ENS
(\xref{transmitter/ass:ens}), or the
premise a \exttermref{../transmitter/contract.pdf}{term-source-model}{source model} declares in its place, for its source states;
let it be attacked by an attack in the sense of
Lemma~\ref{lem:memory-lift}\,(b) that is all-at-once and has finite
registers; and let the \exttermref{../reductions/contract.pdf}{term-reduction-chain}{reduction chain} from which the declaration was
compiled carry it onto the round model, as A-LINK
(\xref{protocol-compiler/ass:link})\,(i)(a), (i)(b), (ii)(a) and (iv) attest,
the existence of each link's channel, the pull-back identity, the source
identity and complete positivity, with A-ATK in the constrained branch:
\begin{itemize}[leftmargin=2.4em]
\item[(H-CHAIN)] the marginal on $Z_1^n\Chat_1^nI_1^nE_n$ of the run's state
  after the $n$ rounds is the state the
  \termref{term-round-channel}{round channels} of Definition~\ref{def:roundmodel}
  produce under an \exttermref{../foundations/contract.pdf}{term-attack-sequence}{attack sequence} of that definition,
  $\CList$-admissible in the constrained branch, from an input state as in
  the branch's accumulation theorem; the run's state is an extension of it,
  holding the parties' further records.
\end{itemize}
Let Assumptions~\ref{ass:keyreg}, \ref{ass:hash}, \ref{ass:seed},
\ref{ass:fixed}, \ref{ass:conform} and~\ref{ass:lab}, the last at each of
the two parties, and A-SEQ (\xref{transmitter/ass:seq}) hold, let
$\alpha\in(1,2)$ and $\epa,\eev\in(0,1]$. If the output length satisfies
\begin{equation}\label{eq:keylength}
  \ell\;\le\;n\,\hlow-\lEC-\Big\lceil\log_2\frac1{\eev}\Big\rceil
  -\frac{\alpha}{\alpha-1}\log_2\frac1{\epa}+1,
\end{equation}
then, under that attack, the protocol is $\epa$-secret on the accepted
state:
\begin{equation*}
  \frac12\Big\|\rho_{K\,E_{\mathrm{tot}}\wedge\Omega_{\mathrm{acc}}}
  -\frac{\idon{K}}{2^{\ell}}\otimes
  \rho_{E_{\mathrm{tot}}\wedge\Omega_{\mathrm{acc}}}\Big\|_1\le\epa .
\end{equation*}
\end{theorem}
\end{mainclaim}
\claimreading{Over n rounds the attacker's uncertainty about the raw key accumulates to n times the per-round bound h\_low, less a term for how unlikely the acceptance test is to pass, which cancels over all runs; the bound h\_low holds over every attack compatible with your declared devices and with passing statistics. Subtracted are the error-correction transcript of \ensuremath{\lambda}\_EC bits, the verification hash of \ensuremath{\lceil}log\ensuremath{_{2}}(1/\ensuremath{\varepsilon}\_EV)\ensuremath{\rceil} bits and a finite-size penalty of \ensuremath{\alpha}/(\ensuremath{\alpha}\ensuremath{-}1)\ensuremath{\cdot}log\ensuremath{_{2}}(1/\ensuremath{\varepsilon}\_PA) bits; one bit is added back. The remaining key of \ensuremath{\ell} bits is secret, counted over all runs: it is within the secrecy parameter \ensuremath{\varepsilon}\_PA of a uniform key independent of the attacker's memory and every public message. It is proved for a finite-memory attacker receiving every signal at once; \hyperlink{claim-meat-core/cor:composed-security}{Corollary~\ref*{cor:composed-security}} covers every attacker. The R\'enyi order \ensuremath{\alpha} in (1, 2) is fixed before the run (\hyperlink{claim-meat-core/ass:fixed}{A-FIXED}) and changes the key's length, never its secrecy.}
\claimmeta{\metalabel{Uses}\mref{Def.}{def:admissible-declaration}, \mref{Def.}{def:homog}, \mref{Def.}{def:roundmodel}, \mref{Thm.}{thm:accumulated-entropy}, \mref{Thm.}{thm:constrained-accumulation}, \mref{Def.}{def:feasible}, \mref{Thm.}{thm:reduction}, \mref{Cor.}{cor:reduction-constrained}, \mref{Def.}{def:postprocessing-model}, \mref{Lem.}{lem:memory-lift}, \extref{../foundations/contract.pdf}{foundations/def:renyi}
\metasep \metalabel{Assumes}\aref{A-KEYREG}{ass:keyreg}, \aref{A-HASH}{ass:hash}, \extaref{../protocol-compiler/contract.pdf}{protocol-compiler/ass:atk}{A-ATK}, \aref{A-SEED}{ass:seed}, \aref{A-FIXED}{ass:fixed}, \aref{A-CONFORM}{ass:conform}, \aref{A-LAB}{ass:lab}, \extaref{../transmitter/contract.pdf}{transmitter/ass:seq}{A-SEQ}, \extaref{../protocol-compiler/contract.pdf}{protocol-compiler/ass:link}{A-LINK}, \extaref{../transmitter/contract.pdf}{transmitter/ass:ens}{A-ENS}, \extaref{../protocol-compiler/contract.pdf}{protocol-compiler/ass:spec}{A-SPEC}, \extaref{../protocol-compiler/contract.pdf}{protocol-compiler/ass:alice-meas}{A-ALICE-MEAS}
\metasep \metalabel{Used by}\mref{Cor.}{cor:composed-security}, \mref{Cor.}{cor:separate-executions}, \mref{Def.}{def:max-key-length}, \mref{Rem.}{rem:plus2}, \mref{Rem.}{rem:hash-penalty}, \mref{Prop.}{prop:monotone}, \mref{Cor.}{cor:keylength-bound}, \mref{Cor.}{cor:constrained-key-length}, \mref{Rem.}{rem:cgate}, \mref{Prop.}{prop:optol-nobound}, \mref{Req.}{req:infinite-value}, \mref{Req.}{req:reported-key-length}, \mref{Req.}{req:hash-penalty}, \mref{Req.}{req:key-length-alpha-domain}, \mref{Req.}{req:compiled-declaration-only}, \mref{Req.}{req:security-statement-not-vacuous}, \mref{Req.}{req:constrained-key-gate}
\metasep newly derived
\metasep implemented (effective: accepted)
\metasep \metalabel{Code}\coderef{https://github.com/fraunhofersg/certiqs-key/blob/v0.1.0/src/certiqs_key/meat_core/finitekey/keylength.py}{certiqs_key.meat_core.finitekey.keylength:key_length}
\metasep \metalabel{Oracles}\coderef{https://github.com/fraunhofersg/certiqs-key/blob/v0.1.0/tests/oracles/registry.py}{registry.py}}

\begin{proof}
\emph{Idea.} Secrecy is first proved on the state conditioned on the
acceptance test: the chain rule for a classical register charges the whole
public transcript $L$ against the accumulated entropy, the public seeds cost
nothing because they are independent of everything else, and the R\'enyi
leftover-hashing lemma turns the result into a trace-distance bound, whose
acceptance-probability factor cancels. A contraction then passes from the
acceptance test to the accepted state, on which the announced flag
$F_{\mathrm{acc}}$ is a register of the adversary's side information. An
attack under which the test never passes is disposed of first.

\emph{Step 0 (an attack that never passes the test).} The event $\OmAT$ is
a function of $\Chat_1^n$, so by (H-CHAIN) its probability under the
attack is the one in the round model. If $\Pr[\OmAT]=0$, then
$\Omega_{\mathrm{acc}}=\OmAT\wedge\{F_{\mathrm{acc}}=\mathrm{acc}\}$
(Definition~\ref{def:postprocessing-model}) has probability zero, so
$\rho_{K\,E_{\mathrm{tot}}\wedge\Omega_{\mathrm{acc}}}=0$ and
$\rho_{E_{\mathrm{tot}}\wedge\Omega_{\mathrm{acc}}}=0$, and the left-hand
side of the secrecy criterion is $0\le\epa$. From now on $\Pr[\OmAT]>0$,
and the hypotheses of the branch's accumulation theorem hold in full.

\emph{Step 1 (secrecy on the acceptance test).} Abbreviate
$B\coloneqq\Chat_1^nI_1^n\Ehat$. By (H-CHAIN) the rounds of the run under the
attack leave, on $Z_1^n\Chat_1^nI_1^nE_n$, the state the
\termref{term-round-channel}{round channels} of Definition~\ref{def:roundmodel}
produce under an attack sequence of that definition, to which the
accumulation theorem of the branch applies, the run's further records being
an extension of it, and
the register $E_{\mathrm{tot}}=B\,S\,L\,G$ is what the adversary holds at the
end of the protocol. By A-SEQ (\xref{transmitter/ass:seq})\,(i), (ii)
and~(iv), every message of a round is sent after every measurement of that
round, no operation on a quantum system of a round depends on that round's
messages, and the classical post-processing is applied after the $n$
rounds, so each round is the measurement of its round channel followed by
announcements that the model hands to the adversary, and the post-processing
acts on the state the $n$ round channels produce. By
Assumption~\ref{ass:lab} at each of the two parties, during the protocol
nothing about the raw string, the keys or the parties' local randomness
leaves either laboratory, other than as the protocol's output, except
through the systems the \exttermref{../foundations/contract.pdf}{term-round-model}{round model} hands to the adversary and the
public messages, and the adversary acts on neither laboratory except through
the systems and messages the model routes to it. By
Assumption~\ref{ass:conform}\,(ii) the classical messages after the rounds
are exactly those of the model and of the authentication post-processing
(Definition~\ref{def:authenticated-channel}\,(c)). In the
honest-authentication setting no message fails authentication, so the
messages of the authentication post-processing are accept or abort
decisions that are functions of $\Omega_{\mathrm{acc}}$, hence of
$\Chat_1^n$ and $F_{\mathrm{acc}}$, which $E_{\mathrm{tot}}$ holds;
appending them is a channel acting on $E_{\mathrm{tot}}$ alone, which does
not increase the trace distance of the secrecy criterion and keeps its ideal
term of the same form, so they cost nothing. By
Assumption~\ref{ass:fixed}, $n$, $\acSet$, $\ell$, $\lEC$, $\epa$ and $\eev$
are constants of the run, so the event $\OmAT$, the dimension $2^{\lambda}$
of $L$ and the output length do not depend on the run's data, and by
Assumption~\ref{ass:conform}\,(i) and~(iv) the run applies that test and
outputs the $\ell$-bit key the model computes, or $\perp$. The
post-processing is executable: the key holder computes $Z_1^n$ from its own
readout and the announcements (Definition~\ref{def:admissible-declaration}\,(viii)), and every
object the post-processing forms is a function of registers present in the
state that Theorem~\ref{thm:accumulated-entropy} bounds, together with the
public seeds.

The state after the public phase, conditioned on $\OmAT$, is classical on
$L$; by the chain rule for conditioning on a classical register
\cite[Lemmas~5.3--5.4, Eq.~(5.94)]{tomamichel2016finite}, applied with
the classical system $L$ of dimension $2^{\lambda}$ and the conditioning
system $B\,S$ (the $\alpha$-domain
restriction to $(1,2)$ is this module's, inherited from the
leftover-hashing lemma applied next; the cited chain rule is stated for the
sandwiched family of \cite[Def.~5.2, Eqs.~(5.18)--(5.19)]%
{tomamichel2016finite} on its data-processing range
$\alpha\in[\frac12,\infty]$, the range on which
\cite[Thm.~4.1, Eq.~(4.82)]{tomamichel2016finite} grounds data
processing, and $(1,2)\subset[\frac12,\infty]$),
\begin{equation}\label{eq:chainrule}
\begin{aligned}
  \Hsandup\big(Z_1^n\big|B\,S\,L\big)_{\rho_{|\OmAT}}
  &\;\ge\;\Hsandup\big(Z_1^n\big|B\,S\big)_{\rho_{|\OmAT}}-\lambda
  \;\ge\;\Hsandup\big(Z_1^n\big|B\big)_{\rho_{|\OmAT}}-\lambda\\
  &\;\ge\;n\,\hlow-\frac{\alpha}{\alpha-1}\log_2\frac{1}{\Pr[\OmAT]}
  -\lambda .
\end{aligned}
\end{equation}
The second inequality holds because the seeds cost nothing. By
Assumption~\ref{ass:seed} they are drawn after the rounds, independently of
every other register, so that
$\rho_{Z_1^nBS|\OmAT}=\rho_{Z_1^nB|\OmAT}\otimes\omega_S$ for the law
$\omega_S$ of $S$, $\OmAT$ being a function of $\Chat_1^n$. For every
$\sigma_B\in\Seq{B}$ the state $\sigma_B\otimes\omega_S$ is admissible in the
infimum that defines $\Hsandup(Z_1^n|BS)$ (\xref{foundations/def:renyi}), and
with $\gamma\coloneqq\frac{1-\alpha}{2\alpha}$, powers taken on supports,
\[
  (\idon{Z}\otimes\sigma_B\otimes\omega_S)^{\gamma}
  (\rho_{Z_1^nB}\otimes\omega_S)
  (\idon{Z}\otimes\sigma_B\otimes\omega_S)^{\gamma}
  =\big((\idon{Z}\otimes\sigma_B)^{\gamma}\rho_{Z_1^nB}
  (\idon{Z}\otimes\sigma_B)^{\gamma}\big)\otimes\omega_S^{1/\alpha};
\]
the trace of the $\alpha$-th power of a tensor product is the product of the
traces, $\tr[(\omega_S^{1/\alpha})^{\alpha}]=\tr\omega_S=1$, and the support
condition holds for the product exactly when it holds for the first factor.
Hence
$\Dsand(\rho_{Z_1^nB}\otimes\omega_S\|\idon{Z}\otimes\sigma_B\otimes\omega_S)
=\Dsand(\rho_{Z_1^nB}\|\idon{Z}\otimes\sigma_B)$, and the infimum over
$\sigma_B$ gives $\Hsandup(Z_1^n|BS)\ge\Hsandup(Z_1^n|B)$.

The third inequality of \eqref{eq:chainrule} uses
Theorem~\ref{thm:accumulated-entropy} with
Theorem~\ref{thm:reduction} ($\hqkd\ge\hlow$) in the unconstrained
branch, and Theorem~\ref{thm:constrained-accumulation} with
Corollary~\ref{cor:reduction-constrained} --- that is,
Corollary~\ref{cor:constrained-key-length} --- in the constrained branch;
here and in the rest of the proof $\hlow$ is \eqref{eq:hlow-def} over the
declared $\Feas$, so in the constrained branch it is the constrained value
$\hlow(\CList)$, and the remainder of the proof consumes only the
displayed bound, supplied by either branch. The chain rule holds for an
arbitrary jointly distributed classical $L$ of dimension $2^{\lambda}$ ---
no functional dependence of $L$ on $Z_1^n$ is required --- so it covers
$F_{\mathrm{acc}}$, which depends also on the partner's reconstruction
$\hat Z_1^n$. The third inequality of \eqref{eq:chainrule} applies to the
post-processed state because
$\Hsandup(Z_1^n|B)$ is a functional of the $(Z_1^n,B)$ marginal only, and
the post-processing leaves that marginal unchanged: the registers traced out
by the \termref{term-round-channel}{round channels} (the partner's outcome and
announcement copies;
Definition~\ref{def:roundmodel}) still exist in the physical protocol as an
extension of the modelled state, and every appended object ---
$S$ drawn independently, $L_{\mathrm{EC}}$ and $T_{\mathrm{EV}}$ from
classical copies of $Z_1^n$ and the public seeds, $F_{\mathrm{acc}}$ from
$T_{\mathrm{EV}}$ and the partner's data (Assumption~\ref{ass:keyreg}) --- is
generated by channels acting trivially on $Z_1^nB$. Hence the accumulation
bound of the branch in force, Theorem~\ref{thm:accumulated-entropy} or
Theorem~\ref{thm:constrained-accumulation}, applies verbatim. Now apply the
R\'enyi leftover-hashing lemma \cite[Thm.~4, Eq.~(13)]{kamin2025renyi}
(stated there for $\alpha\in(1,2)$; this is what confines the theorem to
that domain), under the ideal-family hypothesis of
Assumption~\ref{ass:hash}\,(i) (\cite[Lemma~6.5, Rem.~6.1]{tupkary2026decoy};
a merely two-universal family costs one further bit there) and with $G$
drawn with the law of its family independently of $Z_1^nBSL$
(Assumption~\ref{ass:seed}), to the normalized state $\rho_{|\OmAT}$ with
side information $E\coloneqq B\,S\,L$ and hash register $G$:
\begin{equation*}
  \frac12\Big\|\rho_{K\,EG|\OmAT}-\frac{\idon{K}}{2^{\ell}}\otimes
  \rho_{EG|\OmAT}\Big\|_1
  \;\le\;2^{\frac{2(1-\alpha)}{\alpha}}\,
  2^{\frac{1-\alpha}{\alpha}\left(
  \Hsandup(Z_1^n|BSL)_{\rho_{|\OmAT}}-\ell\right)} .
\end{equation*}
Multiply both sides by $\Pr[\OmAT]$ to pass to the subnormalized
$\wedge\OmAT$ state, and insert \eqref{eq:chainrule}: with
$\kappa\coloneqq\frac{\alpha-1}{\alpha}>0$, the factor
$2^{\kappa\log_2(1/\Pr[\OmAT])}=1/\Pr[\OmAT]$ cancels the prefactor
$\Pr[\OmAT]$, leaving
\begin{equation}\label{eq:secrecy-at}
  \frac12\Big\|\rho_{K\,E_{\mathrm{tot}}\wedge\OmAT}
  -\frac{\idon{K}}{2^{\ell}}\otimes
  \rho_{E_{\mathrm{tot}}\wedge\OmAT}\Big\|_1
  \;\le\;2^{-2\kappa}\,2^{-\kappa\left(n\hlow-\lambda-\ell\right)} .
\end{equation}
Requiring the right side to be at most $\epa$ is equivalent to
$n\hlow-\lambda-\ell\ge\frac1{\kappa}\log_2\frac1{\epa}-2$, i.e.\ to
\eqref{eq:keylength} after inserting
$\lambda=\lEC+\lceil\log_2\frac1{\eev}\rceil+1$.

\emph{Step 2 (from $\OmAT$ to $\Omega_{\mathrm{acc}}$).} The flag
$F_{\mathrm{acc}}$ is a register of $E_{\mathrm{tot}}$ (it is part of
$L$). Let $\Pi\coloneqq\dyad{\mathrm{acc}}_{F_{\mathrm{acc}}}$, acting on
$E_{\mathrm{tot}}$ only. The selection map
$X\mapsto(\idon{K}\otimes\Pi)X(\idon{K}\otimes\Pi)$ is trace-norm
contractive ($\|AXB\|_1\le\|A\|_{\infty}\|X\|_1\|B\|_{\infty}$ with
$\|\Pi\|_{\infty}=1$), maps
$\rho_{\wedge\OmAT}\mapsto\rho_{\wedge\Omega_{\mathrm{acc}}}$, and maps
the ideal term to
$\frac{\idon{K}}{2^{\ell}}\otimes
\Pi\rho_{E_{\mathrm{tot}}\wedge\OmAT}\Pi
=\frac{\idon{K}}{2^{\ell}}\otimes
\rho_{E_{\mathrm{tot}}\wedge\Omega_{\mathrm{acc}}}$, preserving the ideal
form because $\Pi$ does not act on $K$. Applying it to
\eqref{eq:secrecy-at} gives the claimed accepted-state secrecy.
\end{proof}

% claim meat-core/cor:composed-security
\claimsummary{meat-core/cor:composed-security}{Secrecy and correctness compose to (\ensuremath{\varepsilon}\_PA + \ensuremath{\varepsilon}\_EV)-security with honest authentication; over an authentication scheme within \ensuremath{\varepsilon}\_auth that consumes k\_auth bits of pre-shared key, the protocol is (\ensuremath{\varepsilon}\_PA + \ensuremath{\varepsilon}\_EV + 2\ensuremath{\varepsilon}\_auth)-secure; when both parties output a key, its net key gain is \ensuremath{\ell} \ensuremath{-} k\_auth bits.}
\claimbookmark{meat-core/cor:composed-security}{Corollary}{cor:composed-security}{Security of the protocol and its net key}
\claimleadin{%
Secrecy (Theorem~\ref{thm:keylength}) and correctness
(Lemma~\ref{lem:ev-correctness}) are the two halves of the security of the
fixed-length protocol in the setting where authentication behaves honestly.
A real authentication scheme can abort asymmetrically and has its own error;
the authentication post-processing reduces the first to the honest setting
\cite[Sec.~5, Thm.~5.1]{tupkary2026decoy}, and the scheme's error enters the
security parameter, here through the triangle inequality, which charges it
twice. That it enters is the motivation of
\cite[Sec.~4, footnote~8]{tupkary2026decoy} and
\cite[Sec.~5.6.1]{tupkary2025survey}; (iii) is derived without them.
}
\begin{mainclaim}
\begin{corollary}[Security of the protocol and its net key]\label{cor:composed-security}
Let the hypotheses of Theorem~\ref{thm:keylength}, the run as in
Lemma~\ref{lem:memory-lift}\,(a) and (H-CHAIN) among them, hold at every
attack in the sense of Lemma~\ref{lem:memory-lift}\,(b) that is all-at-once
and has finite registers, with $\ell\ge1$ satisfying \eqref{eq:keylength},
A-SEQ with its clause (v) among them; let those of
Lemma~\ref{lem:ev-correctness} hold; and put
$\tau_{K_AK_B}\coloneqq2^{-\ell}\sum_{k\in\{0,1\}^{\ell}}
\dyad{k}_{K_A}\otimes\dyad{k}_{K_B}$.
\begin{enumerate}[label=(\roman*),leftmargin=2.4em]
\item In the \termref{term-honest-authentication-setting}{honest-authentication setting}, at every output state the
  setting admits, whatever the dimension of the adversary's registers and the
  order in which the source releases its signals, the protocol is
  $\epa$-secret on the accepted state as Theorem~\ref{thm:keylength} states
  it, and $(\epa+\eev)$-secure in the sense of
  \cite[Sec.~2.3, Def.~2.1, Eq.~(1)]{tupkary2025survey}:
  \begin{equation}\label{eq:composed-honest}
    \frac12\Big\|\rho_{K_AK_B\,E_{\mathrm{tot}}\wedge\Omega_{\mathrm{acc}}}
    -\tau_{K_AK_B}\otimes
    \rho_{E_{\mathrm{tot}}\wedge\Omega_{\mathrm{acc}}}\Big\|_1
    \le\epa+\eev ,
  \end{equation}
  and on the complement of $\Omega_{\mathrm{acc}}$ both parties output
  $\perp$.
\item Run with the \termref{term-authentication-post-processing}{authentication post-processing}
  (Definition~\ref{def:authenticated-channel}\,(c)) over the
  \termref{term-authenticated-classical-channel}{authenticated classical channel}
  (Definition~\ref{def:authenticated-channel}\,(a)), the protocol is
  $(\epa+\eev)$-secure in the sense of
  \cite[Sec.~3.3, Def.~3.8]{tupkary2026decoy}, which admits asymmetric
  aborts, for every output state that channel admits.
\item Let Assumptions~\ref{ass:auth}, \ref{ass:fixed}
  and~\ref{ass:conform} hold, with $\varepsilon_{\mathrm{auth}}$ and
  $k_{\mathrm{auth}}$ the parameters of Assumption~\ref{ass:auth}. Run with
  the authentication post-processing over the classical channel of
  Assumption~\ref{ass:auth}, the protocol is
  $(\epa+\eev+2\varepsilon_{\mathrm{auth}})$-secure in the sense of
  \cite[Sec.~3.3, Def.~3.8]{tupkary2026decoy}: for every attack on that run,
  with output state $\rho'$ on the key registers and the adversary's final
  information,
  \begin{equation}\label{eq:composed-auth}
    \frac12\big\|\rho'-\mathcal R_{\mathrm{ideal}}(\rho')\big\|_1
    \le\epa+\eev+2\varepsilon_{\mathrm{auth}},
  \end{equation}
  with $\mathcal R_{\mathrm{ideal}}$ the map of
  \cite[Sec.~3.3, Eqs.~(13)--(17)]{tupkary2026decoy}. It consumes at most
  $k_{\mathrm{auth}}$ bits of the pre-shared key, and when both parties
  output an $\ell$-bit key its \emph{net key gain} is
  $\ell_{\mathrm{net}}\coloneqq\ell-k_{\mathrm{auth}}$ bits; for
  $\ell_{\mathrm{net}}\le0$ the run yields no net key.
\end{enumerate}
\end{corollary}
\end{mainclaim}
\claimreading{A run either aborts or outputs a key of \ensuremath{\ell} bits that is a perfect key up to the security parameter \ensuremath{\varepsilon}: identical for both parties, uniform and independent of all the attacker holds, public messages included. One party may abort while the other gets a key. No test, whatever the attacker knows, tells the run's output from an ideal protocol's with advantage above \ensuremath{\varepsilon} = \ensuremath{\varepsilon}\_PA + \ensuremath{\varepsilon}\_EV + 2\ensuremath{\varepsilon}\_auth. It holds against every attack, of any memory, if every premise in the run record holds, \hyperlink{claim-meat-core/ass:auth}{A-AUTH}, \hyperlink{claim-meat-core/ass:conform}{A-CONFORM} and \hyperlink{claim-meat-core/ass:fixed}{A-FIXED} included. A first-order run also needs its bound to hold, computed h\_lb \ensuremath{\le} h\_low, which no premise covers (\hyperlink{claim-meat-core/frank-wolfe/rem:fo-arithmetic}{Remark~\ref*{rem:fo-arithmetic}}). \ensuremath{\varepsilon} averages over aborted runs too. The criterion is composable. Separate executions provably add \ensuremath{\varepsilon} (\hyperlink{claim-meat-core/cor:separate-executions}{Corollary~\ref*{cor:separate-executions}}). Not every later use of the key is covered (\hyperlink{claim-meat-core/rem:repeated-executions}{Remark~\ref*{rem:repeated-executions}}). A run where both get a key nets \ensuremath{\ell} \ensuremath{-} k\_auth bits; an abort uses up to k\_auth. Abort rates: \hyperlink{claim-meat-core/rem:completeness-scope}{Remark~\ref*{rem:completeness-scope}}.}
\claimmeta{\metalabel{Uses}\mref{Def.}{def:postprocessing-model}, \mref{Def.}{def:authenticated-channel}, \mref{Lem.}{lem:ev-correctness}, \mref{Lem.}{lem:memory-lift}, \mref{Lem.}{lem:honest-attacks}, \mref{Thm.}{thm:keylength}
\metasep \metalabel{Assumes}\aref{A-AUTH}{ass:auth}, \aref{A-FIXED}{ass:fixed}, \aref{A-CONFORM}{ass:conform}
\metasep \metalabel{Used by}\mref{Cor.}{cor:separate-executions}, \mref{Rem.}{rem:repeated-executions}, \mref{Cor.}{cor:keylength-bound}, \mref{Req.}{req:security-statement-not-vacuous}
\metasep newly derived
\metasep implemented (effective: accepted)
\metasep \metalabel{Code}\coderef{https://github.com/fraunhofersg/certiqs-key/blob/v0.1.0/src/certiqs_key/meat_core/finitekey/keylength.py}{certiqs_key.meat_core.finitekey.keylength:composed_security}}

\begin{proof}
\emph{(i).} We prove the two bounds at an all-at-once attack with finite
registers first, and then at every output state of the setting.

\emph{At an all-at-once attack with finite registers.} Theorem~\ref{thm:keylength}
gives the secrecy there. The rest is the statement that correctness and secrecy together imply
security \cite[Sec.~2.3.1, Thm.~2.1]{tupkary2025survey}, which is
\cite[Sec.~III.B.3, Thm.~2]{portmann2022security}; we give the argument in
this contract's notation. In the honest-authentication setting both parties
learn $\Omega_{\mathrm{acc}}$, a function of the announced $\Chat_1^n$ and
$F_{\mathrm{acc}}$, so on its complement both output $\perp$, where the real
and the ideal state coincide; it remains to bound
\eqref{eq:composed-honest}. Write
$\sigma\coloneqq\rho_{K\hat K E_{\mathrm{tot}}\wedge\Omega_{\mathrm{acc}}}
=\sum_{k,k'}\dyad{k}_K\otimes\dyad{k'}_{\hat K}\otimes\sigma^{k,k'}$, with
$\sigma^{k,k'}\ge0$ on $E_{\mathrm{tot}}$, and let $C$ be the channel on $K$
that measures $K$ in its classical basis and writes the outcome into $K$ and
into a copy $\hat K$, $C(X)\coloneqq\sum_k\bra{k}X\ket{k}\,
\dyad{k}_K\otimes\dyad{k}_{\hat K}$. Put
$\sigma'\coloneqq(C\otimes\idmap_{E_{\mathrm{tot}}})
(\rho_{K E_{\mathrm{tot}}\wedge\Omega_{\mathrm{acc}}})
=\sum_{k}\dyad{k}_K\otimes\dyad{k}_{\hat K}\otimes\sum_{k'}\sigma^{k,k'}$.
The difference $\sigma-\sigma'$ is block diagonal in $k$, the terms with
$k'=k$ cancel, and in block $k$ it is
$\sum_{k'\ne k}\dyad{k'}_{\hat K}\otimes\sigma^{k,k'}
-\dyad{k}_{\hat K}\otimes\sum_{k'\ne k}\sigma^{k,k'}$, whose trace norm is
$2\sum_{k'\ne k}\tr\sigma^{k,k'}$ because every $\sigma^{k,k'}$ is positive.
Hence
$\frac12\|\sigma-\sigma'\|_1=\Pr[\Omega_{\mathrm{acc}}\wedge K\ne\hat K]
\le\eev$ by Lemma~\ref{lem:ev-correctness}. Since $C$ maps
$\idon{K}/2^{\ell}$ to $\tau_{K\hat K}$ and a channel does not increase the
trace norm,
$\frac12\|\sigma'-\tau_{K\hat K}\otimes
\rho_{E_{\mathrm{tot}}\wedge\Omega_{\mathrm{acc}}}\|_1\le\epa$ by
Theorem~\ref{thm:keylength}. The triangle inequality gives
\eqref{eq:composed-honest}, read for $(K_A,K_B)$, which is $(K,\hat K)$ in
the order the key holder fixes; $\tau$ is symmetric under exchanging the two
registers.

\emph{At every output state of the setting.} With $\Pi$ the projection onto
$\Omega_{\mathrm{acc}}$, a function of the public messages, the secrecy
distance of Theorem~\ref{thm:keylength} is the criterion $f$ of
Lemma~\ref{lem:memory-lift}\,(ii) for $L_1(\sigma)$ the selection
$\Pi\sigma\Pi$ followed by the partial trace over every register other than
$K$ and the public registers of $E_{\mathrm{tot}}$, and
$L_2(\sigma)\coloneqq\frac{\idon{K}}{2^{\ell}}\otimes\tr_KL_1(\sigma)$; the
left side of \eqref{eq:composed-honest} is the same with $K_AK_B$ in place of
$K$ and $\tau_{K_AK_B}$ in place of $\idon{K}/2^{\ell}$. Each $L_t$ is
completely positive, does not increase the trace, keeps the public messages
and acts on the other registers given their value, as part (iii) of that
lemma asks, being composed of the selection by $\Pi$, which is diagonal in
the public messages, partial traces over other registers and the appending
of a state. By the first part both bounds hold at every all-at-once attack
with finite registers, so by Lemma~\ref{lem:memory-lift}\,(ii) at every attack
in the sense of its (b). By Lemma~\ref{lem:honest-attacks}, whose premises
A-SEQ, its clause (v) included, A-LAB at both parties, A-CONFORM and
A-ALICE-MEAS\,(i) where no \exttermref{../transmitter/contract.pdf}{term-source-replacement-certificate}{source-replacement certificate} is demanded
are among those of Theorem~\ref{thm:keylength}, every output state of the
honest-authentication setting is $(\idmap\otimes\Phi)(\rho)$ for the output
state $\rho$ of such an attack and a channel $\Phi$ of the adversary given
the public messages, and by Lemma~\ref{lem:memory-lift}\,(iii) neither
criterion increases under $\Phi$. That both parties output $\perp$ off
$\Omega_{\mathrm{acc}}$ holds at every output state, being the run's
processing of its public messages.

\emph{(ii).} Write $\widetilde{\mathcal P}$ for the protocol of the
\termref{term-fixed-length-post-processing-model}{fixed-length post-processing model} and
$\mathcal P\coloneqq\mathcal P_{\mathrm{APP}}\circ\widetilde{\mathcal P}$ for
it followed by the authentication post-processing. The model is an instance
of the protocol of \cite[Sec.~4.1, Protocol~1, steps~2--5]{tupkary2026decoy}
with the key holder as Alice: the length and error-correction decisions are
$\ell(\hat c_1^n)=\ell\cdot\mathbf 1[\operatorname{freq}_{\hat c_1^n}\in
\acSet]$ and the constant $\lEC$, so the messages that carry them are public
functions of $\Chat_1^n$; error correction is one-way with at most
$2^{\lEC}$ transcripts; error verification and privacy amplification are
those steps with the hash lengths of
Definition~\ref{def:postprocessing-model}. By (i),
$\widetilde{\mathcal P}$ is $(\epa+\eev)$-secure, for every output state of
the honest-authentication setting, in the sense of
\cite[Sec.~3.3, Def.~3.9]{tupkary2026decoy} at fixed length
\cite[Eq.~(23)]{tupkary2026decoy}: its key lengths always coincide, and
$E_{\mathrm{tot}}$ holds every public message of $\widetilde{\mathcal P}$
(the announcements, the acceptance values, the seeds, $L$ and $G$) together
with the adversary's quantum memory. That is the hypothesis of
\cite[Sec.~5, Thm.~5.1]{tupkary2026decoy}, whose honest-authentication setting
is this contract's and whose practical authentication setting is the
authenticated classical channel of
Definition~\ref{def:authenticated-channel}\,(a), properties (A1) and (A2)
being its Eqs.~(37)--(38). The theorem concludes $(\epa+\eev)$-security of
$\mathcal P$ in the sense of Def.~3.8 for every output state of that channel;
by the paragraph following it, it applies to an arbitrary protocol followed
by the authentication post-processing of that work, which is
Definition~\ref{def:authenticated-channel}\,(c): step~6 of
\cite[Sec.~4.1, Protocol~1]{tupkary2026decoy} together with the rule of its
preamble that a party receiving an authentication failure aborts, with the
key holder in the role of Alice and the partner in that of Bob.

\emph{(iii).} By Assumption~\ref{ass:conform} the run is $\mathcal P$: it
applies the acceptance test and the verification abort, sends no message
beyond those of the model and of the authentication post-processing,
outputs the $\ell$-bit keys the model computes or $\perp$, uses no key
before the \termref{term-end-of-an-execution}{end of an execution} (Definition~\ref{def:execution-end}),
and aborts on an authentication failure. By
Assumption~\ref{ass:fixed}, $\varepsilon_{\mathrm{auth}}$ and
$k_{\mathrm{auth}}$, with every other parameter, are fixed before the run,
and by Assumption~\ref{ass:auth}\,(i) the run uses one scheme, so the scheme
and the protocol are fixed systems. If the classical channel
is itself the authenticated classical channel
(Assumption~\ref{ass:auth}\,(i), second case), the run is one of those of
(ii), $\varepsilon_{\mathrm{auth}}=0$, and \eqref{eq:composed-auth} is (ii).
Otherwise fix an attack on the run over the scheme, with output state
$\rho'$. By Assumption~\ref{ass:auth}\,(iii) and
Definition~\ref{def:authenticated-channel}\,(b) there is a simulator $\sigma$
at the adversary's interface such that no distinguisher tells the scheme,
run on two identical copies of a uniformly distributed key independent of
everything else and the unauthenticated channel, from the authenticated classical channel
with $\sigma$ with advantage above $\varepsilon_0$
\cite[Sec.~II.E, Def.~1; Sec.~II.F, Eqs.~(4)--(6)]{portmann2022security}.
The run, the attack and any distinguisher after them form a channel applied
to the joint state, before the run, of the two parties' copies of the
pre-shared key with every other register and the adversary's information, together with fresh systems
independent of it; by Assumption~\ref{ass:auth}\,(iv) and because a channel
does not increase the trace distance, every final state against the scheme
run on the actual copies lies within $\delta$ of the one against the scheme
run on the ideal copies. So the scheme run on the actual copies is told from
the
authenticated classical channel with $\sigma$ with advantage at most
$\varepsilon_0+\delta\le\varepsilon_{\mathrm{auth}}$
(Assumption~\ref{ass:auth}\,(v)).
Consider the distinguisher that plays both parties of $\mathcal P$, every
quantum round included, and the given attack, which reaches the classical
channel only through the adversary's interface of the system the
distinguisher is given, and which then holds the parties' key registers and
the adversary's final information. Its final state is $\rho'$ against the
scheme and, against the authenticated channel with $\sigma$, the output
state $\rho$ of $\mathcal P$ under the attack that the given one forms with
$\sigma$; that attack reads messages and controls their delivery only
through the authenticated classical channel, so $\rho$ is an output state
that channel admits. Any measurement on the final state completes this
distinguisher, and the distinguishing advantage between two states is their
trace distance \cite[Sec.~III.B.1]{portmann2022security}, so
$\frac12\|\rho'-\rho\|_1\le\varepsilon_{\mathrm{auth}}$. By (ii),
$\frac12\|\rho-\mathcal R_{\mathrm{ideal}}(\rho)\|_1\le\epa+\eev$. The map
$\mathcal R_{\mathrm{ideal}}$ is a channel on the key registers
\cite[Sec.~3.3, Eqs.~(13), (17)]{tupkary2026decoy}, so it does not increase the
trace norm and
$\frac12\|\mathcal R_{\mathrm{ideal}}(\rho)-\mathcal R_{\mathrm{ideal}}(\rho')\|_1
\le\varepsilon_{\mathrm{auth}}$. The triangle inequality over $\rho'$, $\rho$,
$\mathcal R_{\mathrm{ideal}}(\rho)$ and $\mathcal R_{\mathrm{ideal}}(\rho')$
gives \eqref{eq:composed-auth}. The run consumes at most $k_{\mathrm{auth}}$
bits of pre-shared key (Assumption~\ref{ass:auth}\,(ii)) and, when both
parties output a key, produces an $\ell$-bit key, which is the accounting of
key consumed against key produced in \cite[Sec.~VII.D]{portmann2022security};
the net key gain is then $\ell-k_{\mathrm{auth}}$.
\end{proof}

% claim meat-core/cor:separate-executions
\claimsummary{meat-core/cor:separate-executions}{N executions that share no key and no device state and leave their keys unused until all have ended, run against one adversary, are together (\ensuremath{\varepsilon}\_1 + ... + \ensuremath{\varepsilon}\_N)-secure for the joint ideal map. A key reused from one execution in another is not covered.}
\claimbookmark{meat-core/cor:separate-executions}{Corollary}{cor:separate-executions}{Security parameters of separate executions add}
\claimleadin{%
Composable frameworks add the errors of protocols composed in sequence or
in parallel \cite[Sec.~II.E]{portmann2022security}; that is the motivation
for the trace-distance statement below, which is derived without it. Each
execution's parameter holds against every attack on it, and the other
executions, which touch none of its registers, are part of one such
attack. The per-execution ideal maps are then applied one at a time, each
step costing that execution's parameter.
}
\begin{corollary}[Security parameters of separate executions add]\label{cor:separate-executions}
Let executions $1,\dots,N$ of the protocol run against one adversary, in
sequence, concurrently or interleaved, the adversary carrying information
from each to the others. For each $i$ let the hypotheses of
Corollary~\ref{cor:composed-security}\,(iii) hold for execution~$i$, with
its parameters $\epa^{(i)}$, $\eev^{(i)}$ and
$\varepsilon_{\mathrm{auth}}^{(i)}$. Write
$K^{(i)}\coloneqq K_A^{(i)}K_B^{(i)}$ for its key registers,
$\mathcal R^{(i)}$ for the map $\mathcal R_{\mathrm{ideal}}$ of
\cite[Sec.~3.3, Eqs.~(13)--(17)]{tupkary2026decoy} on $K^{(i)}$, and
$\varepsilon_i\coloneqq\epa^{(i)}+\eev^{(i)}+2\varepsilon_{\mathrm{auth}}^{(i)}$.
Let Assumption~\ref{ass:separate-executions} hold. Then for every attack on
the $N$ executions under which each of them reaches its
\termref{term-end-of-an-execution}{end of an execution} (Definition~\ref{def:execution-end}) on every
branch, with $\rho$ the joint state of $K^{(1)}\cdots K^{(N)}$ and the
adversary's information $E$ at the end of the last of them,
\begin{equation}\label{eq:separate-executions}
  \frac12\Big\|\rho-\big(\mathcal R^{(1)}\otimes\cdots\otimes
  \mathcal R^{(N)}\otimes\idmap_E\big)(\rho)\Big\|_1
  \le\sum_{i=1}^{N}\varepsilon_i .
\end{equation}
\end{corollary}
\claimreading{Running several executions against the same attacker, one after another or at the same time, costs the sum of their security parameters. If each execution i is secure with parameter \ensuremath{\varepsilon}\_i and \hyperlink{claim-meat-core/ass:separate-executions}{A-SEPARATE} holds, whose conditions include that their number and parameters are fixed before the first starts, that they share no key material and no device state, and that their keys stay unused until all have ended, then all their keys together are within \ensuremath{\varepsilon}\_1 + ... + \ensuremath{\varepsilon}\_N of independent ideal keys. Parameters chosen from an earlier execution's data are not covered, nor is a key from one execution reused in another, for example to authenticate it.}
\claimmeta{\metalabel{Uses}\mref{Def.}{def:roundmodel}, \mref{Thm.}{thm:accumulated-entropy}, \mref{Thm.}{thm:constrained-accumulation}, \mref{Def.}{def:execution-end}, \mref{Lem.}{lem:ev-correctness}, \mref{Thm.}{thm:keylength}, \mref{Cor.}{cor:composed-security}
\metasep \metalabel{Assumes}\aref{A-SEPARATE}{ass:separate-executions}
\metasep \metalabel{Used by}\mref{Rem.}{rem:repeated-executions}
\metasep newly derived
\metasep accepted}

\begin{proof}
Fix an attack on the $N$ executions under which each of them reaches its
end on every branch. For each $i$ let $t_i$ be the
\termref{term-end-of-an-execution}{end of an execution} (Definition~\ref{def:execution-end}) for
execution~$i$, and let
$t\coloneqq\max_i t_i$. For a time $s$ and an execution $i$, write
$Y_i(s)$ for the adversary's registers at $s$ together with every register
of the honest parties of every execution $j\ne i$ at $s$. By
Assumption~\ref{ass:separate-executions}\,(iii), at $t$ the register
$K^{(j)}$ of every execution holds that execution's output.

The times $t_i$ and $t$ are random and may depend on the attack. The joint
process is therefore read branch by branch over the classical record of
when each party ends (where the end times are bounded, equivalently with
idle steps that give every branch one schedule); the states named below
are averaged over the branches. That record lies in $Y_i$: the key holder of execution~$i$ ends
by sending a message the adversary sees, and the partner when the
adversary delivers it (Definition~\ref{def:execution-end}). By
Assumption~\ref{ass:separate-executions}\,(ii) and (iii), the evolution
after execution~$i$ has ended depends on that execution's history only
through $Y_i$. The restriction to attacks under which each execution
reaches its end on every branch is the conjunction of the restrictions
that Corollary~\ref{cor:composed-security}\,(iii) already makes for each
execution, whose bound concerns the output state of a run that ends.

\emph{Step 1: one execution against the others.} Fix $i$. Let
$\mathcal E_i$ be the adversary whose initial register holds the given
adversary's initial register together with the initial registers of the
honest parties of every execution $j\ne i$ (their pre-shared keys, their
local randomness and the initial states of their devices), in their joint
initial state, and who carries out, at the times and in the order of the
given attack, the given adversary's operations and every operation of the
honest parties of the executions $j\ne i$. By
Assumption~\ref{ass:separate-executions}\,(ii), no operation of the honest
parties of another execution acts on a register of execution~$i$'s honest
parties, none of theirs acts on a register of another execution, and
execution~$i$ meets the other executions only through systems and
classical channels the adversary has access to. The given adversary's own
operations reach execution~$i$ only through those systems and channels,
its attack interface, as in any attack on it. So
$\mathcal E_i$ is an attack on execution~$i$ alone, whose registers at time
$s$ are $Y_i(s)$, and execution~$i$ run against $\mathcal E_i$ is the same
channel, applied to the same initial state, as the $N$ executions against
the given adversary, with the registers grouped differently.

Corollary~\ref{cor:composed-security}\,(iii) holds for every attack that
its premises admit, whatever the dimension of the adversary's registers, and
it restricts the adversary's initial register and its state by nothing
else. Through Theorem~\ref{thm:keylength} it rests on
Theorem~\ref{thm:accumulated-entropy}, whose initial state $\omega^0$ is
arbitrary on the adversary's initial register $E_0$ of the
\exttermref{../foundations/contract.pdf}{term-round-model}{round model} (Definition~\ref{def:roundmodel}), with Alice's marginal
fixed in \exttermref{../transmitter/contract.pdf}{term-trusted-source-mode}{trusted-source mode}; in the constrained branch of
Theorem~\ref{thm:keylength} it rests on
Theorem~\ref{thm:constrained-accumulation}, whose hypotheses include those
of Theorem~\ref{thm:accumulated-entropy} and so leave $\omega^0$ as
arbitrary, the \exttermref{../foundations/contract.pdf}{term-attack-sequence}{attack sequence} being $\CList$-admissible. It also
rests on Lemma~\ref{lem:ev-correctness}, which holds for every attack. The
initial register of $\mathcal E_i$, however it
is correlated with the given adversary's and the other executions'
systems, is such an $E_0$. The premises are the assumptions on which
Corollary~\ref{cor:composed-security}\,(iii) rests for execution~$i$, and
by Assumption~\ref{ass:separate-executions}\,(iv) each of them holds for
execution~$i$ against $\mathcal E_i$. The registers of $\mathcal E_i$,
which hold the other executions' registers, may be of any dimension, which
the corollary allows. Where a constraint list is declared, A-ATK
(\xref{protocol-compiler/ass:atk}) states its reduction over the
\termref{term-licensed-attack-set}{licensed attack sets} of
Definition~\ref{def:roundmodel}, defined for finite-dimensional registers,
and is asked, with the corollary's other premises, of the all-at-once attacks
with finite registers through which the corollary reaches $\mathcal E_i$.
Clause (iv) asserts each of these, and the proof uses them only through
Corollary~\ref{cor:composed-security}\,(iii). By
Assumption~\ref{ass:separate-executions}\,(i) the parameters of
execution~$i$, and with them $\varepsilon_i$, are fixed before any
execution starts and do not depend on the attack.

Corollary~\ref{cor:composed-security}\,(iii) therefore applies to
execution~$i$ against $\mathcal E_i$, whose information when execution~$i$
has ended is $Y_i(t_i)$. With $\sigma_i$ the state of $K^{(i)}Y_i(t_i)$ at
$t_i$, \eqref{eq:composed-auth} gives
\begin{equation*}
  \frac12\big\|\sigma_i-(\mathcal R^{(i)}\otimes\idmap)(\sigma_i)\big\|_1
  \le\varepsilon_i .
\end{equation*}
From $t_i$ to $t$, nothing acts on $K^{(i)}$ and no information about
another register of execution~$i$'s honest parties leaves their
laboratories (Assumption~\ref{ass:separate-executions}\,(iii)); those
parties, having ended, act on nothing; and no operation of the honest
parties of another execution acts on a register of execution~$i$
(Assumption~\ref{ass:separate-executions}\,(ii)). What happens between
$t_i$ and $t$ therefore acts on $Y_i$ alone and takes no input from a
register of execution~$i$: it is a channel $\Phi_i$ from $Y_i(t_i)$ to
$Y_i(t)$, and the state of $K^{(i)}Y_i(t)$ at $t$ is
$(\idmap_{K^{(i)}}\otimes\Phi_i)(\sigma_i)$. The registers $K^{(j)}$,
$j\ne i$, and $E$ at $t$ belong to $Y_i(t)$, so
$\rho=(\idmap_{K^{(i)}}\otimes\Lambda_i)(\sigma_i)$ with $\Lambda_i$ the
channel $\Phi_i$ followed by the partial trace onto them. This channel acts
on registers other than $K^{(i)}$, so it commutes with
$\mathcal R^{(i)}\otimes\idmap$, and a channel does not increase the trace
norm of a Hermitian operator. Hence
\begin{equation}\label{eq:separate-one}
  \frac12\big\|\rho-(\mathcal R^{(i)}\otimes\idmap)(\rho)\big\|_1
  \le\varepsilon_i\qquad(i=1,\dots,N).
\end{equation}

\emph{Step 2: the hybrid.} Put $\mathcal T_0\coloneqq\idmap$ and
$\mathcal T_i\coloneqq\mathcal R^{(1)}\otimes\cdots\otimes
\mathcal R^{(i)}\otimes\idmap$ for $i=1,\dots,N$, the identity acting on
every other register. The maps $\mathcal R^{(i)}$ act on the disjoint
registers $K^{(i)}$, so
$\mathcal T_i=\mathcal T_{i-1}\circ(\mathcal R^{(i)}\otimes\idmap)$, and
$\mathcal T_N$ is the map of \eqref{eq:separate-executions}. By linearity,
\[
  \rho-\mathcal T_N(\rho)
  =\sum_{i=1}^{N}\big(\mathcal T_{i-1}(\rho)-\mathcal T_i(\rho)\big)
  =\sum_{i=1}^{N}\mathcal T_{i-1}\big(\rho-(\mathcal R^{(i)}\otimes
  \idmap)(\rho)\big).
\]
Each $\mathcal R^{(i)}$ is a channel
\cite[Sec.~3.3, Eqs.~(13), (17)]{tupkary2026decoy}, so each
$\mathcal T_{i-1}$ is a channel and does not increase the trace norm of the
Hermitian operator it is applied to. The triangle inequality and
\eqref{eq:separate-one} give \eqref{eq:separate-executions}.

By \cite[Sec.~3.3, Eqs.~(14), (16), (17)]{tupkary2026decoy},
$\mathcal R^{(i)}$ replaces, for each pair of key lengths of execution~$i$,
its key registers by the ideal keys of that pair, independent of every other
register. Applied for one execution after another, $\mathcal T_N(\rho)$ is
therefore the state in which, for each tuple of the executions' pairs of key
lengths, with its probability under $\rho$, the key registers of every
execution hold the ideal keys of its pair, independent of each other and of
$E$: the ideal of the $N$ executions together.
\end{proof}

% claim meat-core/rem:repeated-executions
\claimsummary{meat-core/rem:repeated-executions}{One execution's parameter covers one run, and separate executions that share no key add their parameters (\hyperlink{claim-meat-core/cor:separate-executions}{Corollary~\ref*{cor:separate-executions}}). A reused authentication key enters through \ensuremath{\varepsilon}\_auth, charged twice, and the earlier run's parameter does not bound its distance in general.}
\claimbookmark{meat-core/rem:repeated-executions}{Remark}{rem:repeated-executions}{One execution's security parameter and a reused key}
\begin{remark}[One execution's security parameter and a reused key]\label{rem:repeated-executions}
The parameter $\epa+\eev+2\varepsilon_{\mathrm{auth}}$ of
Corollary~\ref{cor:composed-security}\,(iii) bounds one execution: one run
of $n$ rounds and its post-processing. It is not a rate per round or per
bit, and it does not shrink when executions are repeated.

A run whose pre-shared authentication key is not fresh, for instance part of
the key of an earlier execution, enters through
Assumption~\ref{ass:auth}\,(iii)--(v). If, before the run starts, the key
jointly with every other register and the adversary's information is within
trace distance $\delta$ of the state in which it is replaced by a uniform and
independent key, then running the scheme, the protocol and any
distinguisher on it is a channel applied to that joint state, together with
fresh systems independent of it; a channel does not increase the trace
distance, so the outputs differ by at most $\delta$ from those on the ideal
key, and by the triangle inequality the scheme realises the authenticated
classical channel within $\varepsilon_0+\delta$, where $\varepsilon_0$ is its
error on an ideal key. The declared $\varepsilon_{\mathrm{auth}}$ includes
$\delta$, and Corollary~\ref{cor:composed-security}\,(iii) charges it twice,
$2\delta$ in all. For a key taken from an earlier execution's output, that
execution's parameter does not bound $\delta$ in general. That parameter is
a sum over the pairs of key lengths
\cite[Sec.~3.3, Def.~3.8, Eqs.~(11), (17)]{tupkary2026decoy}: it bounds the
key's distance from an ideal one where both parties hold a key only
weighted by the probability of that event, which the adversary controls
\cite[Sec.~2.3, Remark~2.1]{tupkary2025survey}. Where one party holds a key
and the other none, as after a blocked final $\mathrm{accept}$ of the
\termref{term-authentication-post-processing}{authentication post-processing}, the ideal state of that definition
also leaves one party without a key, so the parameter does not bound the
probability of that event; there no shared key exists to reuse, and
$\delta$, the distance of the whole state on which the next run starts,
conditioned on its starting, counts that event where the next run starts
on it, and otherwise enters $\delta$ through the conditioning on its
start, which the adversary influences.

For executions that share no key material and no device state, and whose
keys stay unused until the last of them has reached its
\termref{term-end-of-an-execution}{end of an execution} (Assumption~\ref{ass:separate-executions}), the
parameters add:
Corollary~\ref{cor:separate-executions} bounds $N$ such executions, run
against one adversary in any order, by
$\varepsilon_1+\dots+\varepsilon_N$ for the joint ideal map, which is what
composition gives for constructive statements
\cite[Sec.~II.E]{portmann2022security}. Not covered are a key used before
the last execution has ended, since the ideal map applied at the end would
then replace a key whose use the adversary has already seen; parameters
chosen from an earlier execution's data, which would make each
$\varepsilon_i$ depend on the attack; and a chain in which a key of one
execution authenticates the next. For such a chain the hybrid would have
to replace a key by an ideal one at a moment both parties have ended the
earlier execution, which the authentication post-processing does not
provide, and the earlier parameter does not bound $\delta$ (the paragraph
above). That the literature expects such chains to add
\cite[Sec.~2.3, footnote~2]{tupkary2025survey} is cited as motivation
only, and a bound on the failure probability accumulated over every use
\cite[Sec.~II.F]{portmann2022security} lies outside what this contract
proves.
\end{remark}
\claimreading{The security parameter \ensuremath{\varepsilon} covers one execution: one run of n rounds and its post-processing. It is not a per-round or per-bit figure, and it does not shrink when you repeat the protocol. Several executions are together secure with the sum of their \ensuremath{\varepsilon} (\hyperlink{claim-meat-core/cor:separate-executions}{Corollary~\ref*{cor:separate-executions}}) under \hyperlink{claim-meat-core/ass:separate-executions}{A-SEPARATE}. The sum does not cover parameters chosen from an earlier run's data, such as the error-correction budget or the acceptance test. The conditions of \hyperlink{claim-meat-core/ass:separate-executions}{A-SEPARATE} include: (i) their number and parameters are all fixed before the first starts; (ii) they share no key material and no device state; (iii) their keys stay unused until the last has ended. Reusing key from an earlier run to authenticate the next is not covered either. Such a key enters only through the next run's declared authentication error \ensuremath{\varepsilon}\_auth, counted twice, and the earlier run's \ensuremath{\varepsilon} does not by itself bound how far that key is from an ideal one.}
\claimmeta{\metalabel{Uses}\mref{Cor.}{cor:composed-security}, \mref{Cor.}{cor:separate-executions}
\metasep newly derived
\metasep accepted}

% claim meat-core/rem:completeness-scope
\claimsummary{meat-core/rem:completeness-scope}{The completeness parameter \ensuremath{\varepsilon}* of the statistics module bounds an honest run's failure of the acceptance test only. Aborts in error verification and in authentication are bounded by no claim, so \ensuremath{\varepsilon}* does not bound the honest abort probability of the whole protocol.}
\claimbookmark{meat-core/rem:completeness-scope}{Remark}{rem:completeness-scope}{Completeness covers the acceptance test only}
\begin{remark}[Completeness covers the acceptance test only]\label{rem:completeness-scope}
The completeness statements of the statistics module,
\xref{statistics/binomial-box/thm:completeness-bound} with its parameter
$\varepsilon^{\star}$ and \xref{statistics/singleton/thm:singleton-completeness},
bound the probability that an honest run fails the acceptance test, the
complement of the event $\OmAT$ of
Definition~\ref{def:postprocessing-model}\,(i). An honest run of the
fixed-length protocol can also abort later: in error verification, when the
partner's reconstruction $\hat Z_1^n$ differs from $Z_1^n$, and in
authentication
(Definition~\ref{def:authenticated-channel}). That the error-correction scheme lets the partner reconstruct $Z_1^n$
from $L_{\mathrm{EC}}$ and its own data on honest runs is a premise of no
security claim here: it bears only on how often an honest run aborts. No
claim of this contract bounds those abort probabilities, so $\varepsilon^{\star}$ bounds the honest abort
probability of the acceptance test and not that of the whole protocol.
Security is unaffected: Corollary~\ref{cor:composed-security} holds for every
abort pattern.
\end{remark}
\claimreading{The completeness figure \ensuremath{\varepsilon}* bounds one thing only: the chance that an honest run, with no attacker and devices behaving as the declared honest model says, fails the acceptance test on its statistics. An honest run can also abort later: in error verification, when error correction leaves the two raw keys different, and in authentication. No result bounds those aborts, so an honest run's overall abort chance can exceed \ensuremath{\varepsilon}*. Whether your error-correction code reconstructs the key on honest runs is yours to check; it affects only how often you obtain a key. Completeness never affects the security parameter \ensuremath{\varepsilon}, which holds for every abort pattern (\hyperlink{claim-meat-core/cor:composed-security}{Corollary~\ref*{cor:composed-security}}) and so says nothing about how often runs abort.}
\claimmeta{\metalabel{Uses}\mref{Def.}{def:postprocessing-model}, \mref{Def.}{def:authenticated-channel}
\metasep newly derived
\metasep accepted}

% claim meat-core/prop:model-pass-bound
\claimsummary{meat-core/prop:model-pass-bound}{If every round's acceptance value, given the earlier ones, follows the statistics of a feasible state, the run passes the acceptance test with probability at most 2\textasciicircum{}(\ensuremath{-}n D\_min). A lower bound b at one order thus caps that probability at 2\textasciicircum{}(\ensuremath{-}n (\ensuremath{\alpha}\ensuremath{_{0}}\ensuremath{-}1)/\ensuremath{\alpha}\ensuremath{_{0}} \ensuremath{\cdot} b).}
\claimbookmark{meat-core/prop:model-pass-bound}{Proposition}{prop:model-pass-bound}{The pass probability of devices the declared model admits}
\claimleadin{%
Proposition~\ref{prop:minimal-divergence}\,(iii) says that where
$\sfD_{\min}>0$ no single feasible state has its statistics in the acceptance
set; it says nothing about a finite run, whose frequencies can still fall in
$\acSet$. This proposition bounds that event. A state $\omega$ on $A\Qsys$
put through the \termref{term-single-round-protocol-map}{single-round protocol map} yields an acceptance value
distributed as $\Phimap[\omega]$ (Theorem~\ref{thm:protocol-channel}), so the
hypothesis below describes a run without an adversary in which the devices
behave as the declaration says: each round prepares and measures a state of
the \exttermref{../protocol-compiler/contract.pdf}{term-feasible-set}{feasible set}, which may depend on the earlier rounds. That is a
hypothesis about the devices the declaration describes. It is no consequence of the
security model: the \exttermref{../foundations/contract.pdf}{term-round-model}{round model} and the accumulation theorems hold
against every attack and say nothing about how often a run passes, and the
security claims of this contract do not use the hypothesis. Read the other way, it is a statistical test of the model, for an
acceptance set fixed before the run: where the bound is tiny and real
devices pass, that is evidence at the level of the bound, not a proof, that the
declaration does not describe them, which is the reading of
\xref{statistics/ass:reach}. An acceptance set built from the same run's
frequencies passes by construction and supports no such reading. A round state that depends on a
richer record of the past than the earlier acceptance values, with finitely
many values, is covered: given the acceptance values the round's law is then
a convex combination of statistics of feasible states, which is the
statistics of their convex combination, a feasible state, $\Phimap$ being
linear and $\Feas$ convex. The proof is a Chernoff bound whose exponent is
read off a minimising pair of the divergence, with the statistics set mixed
with the uniform distribution so that every outcome has positive probability.
}
\begin{proposition}[The pass probability of devices the declared model admits]\label{prop:model-pass-bound}
Let $\Decl$ and $\acSet$ satisfy the hypotheses of
Proposition~\ref{prop:minimal-divergence}, let $\sfD_{\min}$ be as in
\eqref{eq:dmin} and $\mathcal{D}_{\mathrm{all}}$ the \termref{term-dropped-set}{dropped set} of every
pair there, and let $n\ge1$ be an integer. Let $\Pr$ be a probability
distribution on $\Chatset^{\,n}$, write $\Chat_1,\dots,\Chat_n$ for the
coordinates and $\Chat_1^{t}\coloneqq(\Chat_1,\dots,\Chat_t)$, and suppose:
\begin{description}
\item[(H)] for every $t\in\{1,\dots,n\}$ and every $h\in\Chatset^{\,t-1}$
  with $\Pr[\Chat_1^{t-1}=h]>0$ there is a state $\omega_t(h)\in\Feas$ with
  \begin{equation}\label{eq:pass-hypothesis}
    \Pr\big[\Chat_t=\hat c\ \big|\ \Chat_1^{t-1}=h\big]
    =\Phimap[\omega_t(h)](\hat c)\qquad\text{for every }\hat c\in\Chatset ,
  \end{equation}
  where $\Phimap[\omega](\hat c)$ is the entry of $\Phimap[\omega]$ at
  $\hat c$ and, for $t=1$, $h$ is the empty sequence and the left-hand side
  is $\Pr[\Chat_1=\hat c]$.
\end{description}
The case of independent rounds with one state $\omega\in\Feas$ is
$\omega_t(h)=\omega$ for all $t$ and $h$. Let
$\OmAT=\{\operatorname{freq}_{\Chat_1^n}\in\acSet\}$ be the event of
Definition~\ref{def:postprocessing-model}\,(i), where
$\operatorname{freq}_{\hat c_1^n}(\hat c)\coloneqq
\frac1n\,|\{t:\hat c_t=\hat c\}|$.
\begin{enumerate}[label=(\roman*),leftmargin=2.2em]
\item \emph{(The bound.)} For every real number $\delta\le\sfD_{\min}$,
  \begin{equation}\label{eq:pass-bound}
    \Pr[\OmAT]\;\le\;2^{-n\delta} .
  \end{equation}
  In particular $\Pr[\OmAT]=0$ if $\sfD_{\min}=+\infty$.
\item \emph{(From a bound at one order.)} Let $\alpha_0\in(1,\infty)$ and let
  $b$ be a real number with
  $b\le\underline{h}^{\mathcal{D}_{\mathrm{all}}}_{\alpha_0}$. Then
  $\Pr[\OmAT]\le2^{-x}$ for every real number
  $x\le n\,\frac{\alpha_0-1}{\alpha_0}\,b$. For every
  $\alpha'\in(1,\infty)$ the number
  $f(\alpha')\coloneqq\frac{\alpha'}{\alpha'-1}\cdot\frac{\alpha_0-1}{\alpha_0}\,b$
  satisfies $\frac{\alpha'-1}{\alpha'}\,n\,f(\alpha')
  =n\,\frac{\alpha_0-1}{\alpha_0}\,b$.
\end{enumerate}
\end{proposition}
\claimreading{This result bounds how often devices that behave as your declared model says would pass your acceptance test. Suppose that in every round, given what happened before, the acceptance value is distributed as the statistics of some state the declaration admits as feasible. Then the run passes the test with probability at most 2\textasciicircum{}(\ensuremath{-}n D\_min), where D\_min is the minimal divergence of \hyperlink{claim-meat-core/prop:minimal-divergence}{Proposition~\ref*{prop:minimal-divergence}}, and the engine's proven lower bound on D\_min gives a computable cap. This is a statement about the declared model, not about your real devices and not about an attacker. If the acceptance set was fixed before the run, the cap is tiny and your real devices do pass, that is evidence at the level of the cap, not a proof, that the model does not describe them (\xref{statistics/ass:reach}). A set built from the same run's frequencies passes by construction.}
\claimmeta{\metalabel{Uses}\mref{Def.}{def:admissible-declaration}, \mref{Def.}{def:maps}, \mref{Def.}{def:admissible-acceptance-set}, \mref{Def.}{def:feasible}, \mref{Def.}{def:dropped}, \mref{Prop.}{prop:minimal-divergence}, \mref{Def.}{def:postprocessing-model}, \extref{../foundations/contract.pdf}{foundations/def:kl}, \extref{../foundations/contract.pdf}{foundations/lem:kernel}
\metasep \metalabel{Used by}\mref{Req.}{req:unreachable-acceptance-set}
\metasep newly derived
\metasep implemented (effective: accepted)
\metasep \metalabel{Code}\coderef{https://github.com/fraunhofersg/certiqs-key/blob/v0.1.0/src/certiqs_key/meat_core/finitekey/consistency.py}{certiqs_key.meat_core.finitekey.consistency:model_pass_exponent}}

\begin{proof}
\emph{Idea.} If a weight $w\ge0$ on the acceptance values has mean at most
one under every statistics of a feasible state, the product of the weights
of the rounds has expectation at most one, and on the acceptance event its
logarithm is $n$ times the mean of $\log_2w$ under the observed frequency.
Markov's inequality then bounds the event. The weight is the ratio
$q^{*}/p^{*}$ of a minimising pair of the divergence: the first-order
conditions of the two minimisations give the two properties. Mixing every
statistics with a little of the uniform distribution makes every outcome
positive, so that the derivatives exist; the mixing is removed in a limit.

Write $C\coloneqq\Chatset$, a non-empty finite set, $\Delta$ for the
probability simplex on $C$, $u$ for the uniform distribution on $C$, and
\begin{equation*}
  \mathcal{P}\coloneqq\{\Phimap[\omega]:\omega\in\Feas\}.
\end{equation*}
Every $\Phimap[\omega]$ with $\omega\in\Feas$ lies in $\Delta$, as in the
proof of Proposition~\ref{prop:minimal-divergence}, so $\mathcal{P}\subseteq\Delta$
and, by \eqref{eq:dmin},
$\sfD_{\min}=\inf_{\mathbf q\in\acSet}\inf_{\mathbf p\in\mathcal{P}}
\sfD(\mathbf q\|\mathbf p)$. The map $\Phimap$ of \eqref{eq:Phimap} is
linear (Definition~\ref{def:maps}), hence continuous, and $\Feas$ is convex,
compact and non-empty (Definition~\ref{def:feasible};
Definition~\ref{def:admissible-declaration}), so $\mathcal{P}$ is convex,
compact and non-empty. $\acSet$ is a closed convex subset of $\Delta$
(Definition~\ref{def:admissible-acceptance-set}), hence compact. For
$\mathbf q,\mathbf p\in\Delta$,
$\sfD(\mathbf q\|\mathbf p)=\sum_{c\in C}d(q(c),p(c))$ with $d$ the kernel
of \xref{foundations/lem:kernel}, whose values at vanishing arguments are the
conventions of \xref{foundations/def:kl}; $d$ is lower semicontinuous with values in
$(-\infty,+\infty]$, so $\sfD$, a finite sum of such functions composed with
the continuous coordinate maps, is lower semicontinuous on $\Delta\times\Delta$ with values in
$[0,+\infty]$ (non-negative by \xref{foundations/def:kl}).

If $\acSet=\varnothing$, then $\OmAT=\varnothing$, $\Pr[\OmAT]=0$, and
\eqref{eq:pass-bound} holds for every $\delta$. Assume from now on
$\acSet\ne\varnothing$. Since $2^{-n\delta}\ge2^{-n\sfD_{\min}}$ for $\delta\le\sfD_{\min}$
(with $2^{-\infty}\coloneqq0$), (i) follows from
\begin{equation}\label{eq:pass-bound-dmin}
  \Pr[\OmAT]\le2^{-n\sfD_{\min}} ,
\end{equation}
which steps 1--5 prove.

\emph{Step 1 (the mixed sets).} For $\varepsilon\in(0,1)$ put
$\mathcal{P}_{\varepsilon}\coloneqq\{(1-\varepsilon)\mathbf p+\varepsilon u:
\mathbf p\in\mathcal{P}\}$, the image of $\mathcal{P}$ under an affine map
of $\mathbb{R}^{C}$ into $\Delta$; it is convex, compact and non-empty, and
every $\mathbf p'\in\mathcal{P}_{\varepsilon}$ has $p'(c)\ge\varepsilon/|C|>0$
for every $c$. Put $D_{\varepsilon}\coloneqq\inf_{\mathbf q\in\acSet}
\inf_{\mathbf p'\in\mathcal{P}_{\varepsilon}}\sfD(\mathbf q\|\mathbf p')$.
On $\Delta\times\mathcal{P}_{\varepsilon}$,
$\sfD(\mathbf q\|\mathbf p')=\sum_c\big[\eta(q(c))-q(c)\log_2p'(c)\big]$
with $\eta(x)\coloneqq x\log_2x$ for $x>0$ and $\eta(0)\coloneqq0$,
continuous on $[0,\infty)$,
and $\log_2p'(c)$ continuous as $p'(c)\ge\varepsilon/|C|$; so $\sfD$ is
finite and continuous there, and on the non-empty compact set
$\acSet\times\mathcal{P}_{\varepsilon}$ it attains its minimum at some pair
$(\mathbf q^{*},\mathbf p^{*})$, with $\sfD(\mathbf q^{*}\|\mathbf p^{*})
=D_{\varepsilon}\in[0,\infty)$. Define
\begin{equation*}
  w(c)\coloneqq\frac{q^{*}(c)}{p^{*}(c)}\in[0,\infty),\qquad
  S^{*}\coloneqq\{c\in C:q^{*}(c)>0\}=\{c:w(c)>0\}.
\end{equation*}

\emph{Step 2 (the acceptance side).} We show: every $\mathbf q\in\acSet$
vanishes off $S^{*}$, and
\begin{equation}\label{eq:pass-tilt}
  \sum_{c\in S^{*}}q(c)\log_2w(c)\;\ge\;D_{\varepsilon}
  \qquad\text{for every }\mathbf q\in\acSet .
\end{equation}
Fix $\mathbf q\in\acSet$ and put $\mathbf q_s\coloneqq\mathbf q^{*}
+s(\mathbf q-\mathbf q^{*})$ for $s\in[0,1]$, which lies in $\acSet$ by
convexity, and $g(s)\coloneqq\sfD(\mathbf q_s\|\mathbf p^{*})$. As $p^{*}$
has no zero, $g(s)=\sum_c\varphi_c(q_s(c))$ with
$\varphi_c(x)\coloneqq\eta(x)-x\log_2p^{*}(c)$ on $[0,\infty)$, which is
differentiable on $(0,\infty)$ with
$\varphi_c'(x)=\log_2x+\log_2\mathrm e-\log_2p^{*}(c)$. By minimality of the
pair, $g(s)\ge D_{\varepsilon}=g(0)$ for every $s\in[0,1]$.

Suppose $q(c_0)>0=q^{*}(c_0)$ for some $c_0$. Consider
$\frac{g(s)-g(0)}{s}=\sum_c\frac{\varphi_c(q_s(c))-\varphi_c(q^{*}(c))}{s}$
as $s\downarrow0$. For $c\in S^{*}$, $q_s(c)\to q^{*}(c)>0$ and the quotient tends
to the finite number $\varphi_c'(q^{*}(c))\,(q(c)-q^{*}(c))$. For $c\notin S^{*}$
with $q(c)=0$, $q_s(c)=0$ and the quotient is $0$. For $c\notin S^{*}$ with
$q(c)>0$, $q_s(c)=s\,q(c)$ and the quotient is
$q(c)\log_2\frac{s\,q(c)}{p^{*}(c)}\to-\infty$; $c_0$ is such a $c$. So
$\frac{g(s)-g(0)}{s}\to-\infty$, and $g(s)<g(0)$ for some $s\in(0,1]$, a
contradiction. Hence $q(c)=0$ for every $c\notin S^{*}$.

Then for $c\notin S^{*}$ the summand of $g$ is $\varphi_c(0)=0$ for all $s$, and
for $c\in S^{*}$ and $s\in[0,1)$, $q_s(c)\ge(1-s)q^{*}(c)>0$, so $g$ has the
right derivative
$g'(0^{+})=\sum_{c\in S^{*}}\varphi_c'(q^{*}(c))\,(q(c)-q^{*}(c))$ at $0$, which is
$\ge0$ because $g(s)\ge g(0)$ for $s>0$. Since $\varphi_c'(q^{*}(c))
=\log_2w(c)+\log_2\mathrm e$ and $\sum_{c\in S^{*}}q(c)=1=\sum_{c\in S^{*}}q^{*}(c)$
(both vanish off $S^{*}$), this reads
$\sum_{c\in S^{*}}q(c)\log_2w(c)\ge\sum_{c\in S^{*}}q^{*}(c)\log_2w(c)
=\sfD(\mathbf q^{*}\|\mathbf p^{*})=D_{\varepsilon}$, which is
\eqref{eq:pass-tilt}.

\emph{Step 3 (the model side).} We show
\begin{equation}\label{eq:pass-weight}
  \sum_{c\in C}p(c)\,w(c)\;\le\;\frac{1}{1-\varepsilon}
  \qquad\text{for every }\mathbf p\in\mathcal{P}.
\end{equation}
Let $\mathbf p'\in\mathcal{P}_{\varepsilon}$, put
$\mathbf p_s\coloneqq\mathbf p^{*}+s(\mathbf p'-\mathbf p^{*})\in
\mathcal{P}_{\varepsilon}$ for $s\in[0,1]$ (convexity), so that
$p_s(c)\ge\varepsilon/|C|>0$, and
$\rho(s)\coloneqq\sfD(\mathbf q^{*}\|\mathbf p_s)
=\sum_{c\in S^{*}}q^{*}(c)\big[\log_2q^{*}(c)-\log_2p_s(c)\big]$. Each $p_s(c)$ is
affine in $s$ and positive, so $\rho$ is differentiable on $[0,1]$ with
$\rho'(0)=-\log_2\mathrm e\sum_{c\in S^{*}}q^{*}(c)\frac{p'(c)-p^{*}(c)}{p^{*}(c)}$.
By minimality of the pair, $\rho(s)\ge\rho(0)$ for $s\in[0,1]$, so $\rho'(0)\ge0$,
that is $\sum_{c\in S^{*}}q^{*}(c)\frac{p'(c)}{p^{*}(c)}\le\sum_{c\in S^{*}}q^{*}(c)=1$.
As $w$ vanishes off $S^{*}$, $\sum_cp'(c)w(c)\le1$ for every
$\mathbf p'\in\mathcal{P}_{\varepsilon}$. For $\mathbf p\in\mathcal{P}$ take
$\mathbf p'=(1-\varepsilon)\mathbf p+\varepsilon u$; since
$\sum_cu(c)w(c)\ge0$, $(1-\varepsilon)\sum_cp(c)w(c)\le\sum_cp'(c)w(c)\le1$,
which is \eqref{eq:pass-weight}.

\emph{Step 4 (the process).} Put $Z_0\coloneqq1$ and
$Z_t\coloneqq\prod_{s=1}^{t}w(\Chat_s)\ge0$ for $t=1,\dots,n$, and write
$\mathbb E$ for the expectation under $\Pr$. For $t\ge1$, grouping
$\Chat_1^{t}=(h,\hat c)$ by its first $t-1$ entries $h$, and using
$\Pr[\Chat_1^t=(h,\hat c)]=0$ where $\Pr[\Chat_1^{t-1}=h]=0$,
\begin{equation*}
  \mathbb E[Z_t]
  =\sum_{h:\,\Pr[\Chat_1^{t-1}=h]>0}\Pr[\Chat_1^{t-1}=h]
  \prod_{s=1}^{t-1}w(h_s)\sum_{\hat c\in C}
  \Phimap[\omega_t(h)](\hat c)\,w(\hat c)
  \;\le\;\frac{1}{1-\varepsilon}\,\mathbb E[Z_{t-1}],
\end{equation*}
by \eqref{eq:pass-hypothesis} and by \eqref{eq:pass-weight} for
$\mathbf p=\Phimap[\omega_t(h)]\in\mathcal{P}$. By induction,
$\mathbb E[Z_n]\le(1-\varepsilon)^{-n}$. On $\OmAT$, put
$\mathbf q\coloneqq\operatorname{freq}_{\Chat_1^n}\in\acSet$ and
$N_c\coloneqq nq(c)$. Every value $c=\Chat_t$ has $q(c)\ge\frac1n>0$, so
$c\in S^{*}$ by Step~2 and $w(c)>0$; hence $Z_n>0$ and
\begin{equation*}
  \log_2Z_n=\sum_{c:\,N_c>0}N_c\log_2w(c)
  =n\sum_{c\in S^{*}}q(c)\log_2w(c)\;\ge\;nD_{\varepsilon}
\end{equation*}
by \eqref{eq:pass-tilt}, the summands of $S^{*}$ with $q(c)=0$ being $0$. So
$\OmAT\subseteq\{Z_n\ge2^{nD_{\varepsilon}}\}$, and Markov's inequality for
the non-negative $Z_n$ gives
\begin{equation}\label{eq:pass-eps}
  \Pr[\OmAT]\le2^{-nD_{\varepsilon}}\,\mathbb E[Z_n]
  \le(1-\varepsilon)^{-n}\,2^{-nD_{\varepsilon}}
  \qquad\text{for every }\varepsilon\in(0,1).
\end{equation}

\emph{Step 5 (the limit).} Let $\varepsilon_j\coloneqq\frac{1}{j+1}$. The
sequence $(D_{\varepsilon_j})_j$ in $[0,\infty)$ has a subsequence, which we
relabel, converging to some $L\in[0,+\infty]$. Then
$(1-\varepsilon_j)^{-n}\to1$ and $2^{-nD_{\varepsilon_j}}\to2^{-nL}$, so
\eqref{eq:pass-eps} gives $\Pr[\OmAT]\le2^{-nL}$. We show $L\ge\sfD_{\min}$,
which gives \eqref{eq:pass-bound-dmin}. If $L=+\infty$ there is nothing to
show. Otherwise, for each $j$ pick $\mathbf q_j\in\acSet$ and
$\mathbf p_j\in\mathcal{P}$ with
$\sfD(\mathbf q_j\|\mathbf p'_j)\le D_{\varepsilon_j}+\frac1j$, where
$\mathbf p'_j\coloneqq(1-\varepsilon_j)\mathbf p_j+\varepsilon_ju$. By
compactness of $\acSet\times\mathcal{P}$ a further subsequence has
$(\mathbf q_j,\mathbf p_j)\to(\mathbf q,\mathbf p)\in\acSet\times\mathcal{P}$,
and then $\mathbf p'_j\to\mathbf p$. Lower semicontinuity of $\sfD$ gives
$\sfD_{\min}\le\sfD(\mathbf q\|\mathbf p)\le\liminf_j\sfD(\mathbf q_j\|\mathbf p'_j)
\le\lim_j\big(D_{\varepsilon_j}+\tfrac1j\big)=L$. This proves (i); where
$\sfD_{\min}=+\infty$, \eqref{eq:pass-bound-dmin} reads $\Pr[\OmAT]=0$.

\emph{(ii).} By Proposition~\ref{prop:minimal-divergence}\,(iv),
\eqref{eq:dmin-bracket}, $\frac{\alpha_0-1}{\alpha_0}\,b\le\sfD_{\min}$. For
real $x\le n\frac{\alpha_0-1}{\alpha_0}b$ put $\delta\coloneqq x/n$; then
$\delta\le\sfD_{\min}$, and (i) gives $\Pr[\OmAT]\le2^{-n\delta}=2^{-x}$. The identity
is $\frac{\alpha'-1}{\alpha'}\cdot\frac{\alpha'}{\alpha'-1}=1$.
\end{proof}

% claim meat-core/def:max-key-length
\claimsummary{meat-core/def:max-key-length}{The maximal key length \ensuremath{\ell}\_max is the right-hand side of the key-length bound rounded down to an integer and clamped at 0. The floor and the clamp are both in the safe direction.}
\claimbookmark{meat-core/def:max-key-length}{Definition}{def:max-key-length}{Maximal key length}
\begin{definition}[Maximal key length]\label{def:max-key-length}
The
\termdef{term-maximal-key-length}{maximal key length}
is
\[
  \ell_{\max}\coloneqq\max\{0,\lfloor
  n\hlow-\lEC-\lceil\log_2\tfrac1{\eev}\rceil
  -\tfrac{\alpha}{\alpha-1}\log_2\tfrac1{\epa}+1\rfloor\},
\]
the right-hand side of \eqref{eq:keylength} under a floor and a clamp at
$0$; the floor and the clamp are both in the safe direction: an integer
at most the floor of a number is at most that number, and a key of length
$0$ is trivially secret and correct and claims no key.
\end{definition}
\claimmeta{\metalabel{Uses}\mref{Thm.}{thm:keylength}
\metasep \metalabel{Used by}\mref{Prop.}{prop:monotone}, \mref{Prop.}{prop:optol-nobound}, \mref{Req.}{req:reported-key-length}, \mref{Req.}{req:certified-evaluation}
\metasep newly derived
\metasep accepted}

% claim meat-core/rem:plus2
\claimsummary{meat-core/rem:plus2}{The key-length bound carries the constant +1 where its source derivation and the published fixed-length statement carry +2: one announced bit buys secrecy on the accepted state by an elementary contraction. The source's secrecy definition, conditioned on the acceptance test alone, does not spend it.}
\claimbookmark{meat-core/rem:plus2}{Remark}{rem:plus2}{The constant in the key length and the source's convention}
\begin{remark}[The constant in the key length and the source's convention]
\label{rem:plus2}
The derivation this module's key-length theorem generalises proves
\eqref{eq:secrecy-at} with $\lambda=\lEC+\lceil\log_2\frac1{\eev}\rceil$ and
stops there, i.e.\ its secrecy statement is for the $\wedge\OmAT$ state and
its constant is $+2$; the published fixed-length statement of
\cite[Thm.~7, Eq.~(44)]{kamin2025renyi} reaches the accepted state by a
different route (conditioning on the verification event), with the same
constant. This module spends one deterministic bit (the announced flag) to
reach the accepted state by an elementary contraction argument instead,
avoiding any further import. The cost is exactly one bit in
\eqref{eq:keylength}, which the $\wedge\OmAT$ secrecy definition of that
derivation does not spend; what an instance adopting that definition declares
is Requirement~\ref{req:secrecy-convention}.
\end{remark}
\claimmeta{\metalabel{Uses}\mref{Thm.}{thm:keylength}
\metasep \metalabel{Used by}\mref{Req.}{req:secrecy-convention}
\metasep newly derived
\metasep accepted}

% claim meat-core/rem:hash-penalty
\claimsummary{meat-core/rem:hash-penalty}{With a privacy-amplification family that is only two-universal, collision probability at most 2\textasciicircum{}\ensuremath{-}\ensuremath{\ell}, the key-length bound holds after tightening it by one further bit. Toeplitz hashing is ideal two-universal and needs no such bit.}
\claimbookmark{meat-core/rem:hash-penalty}{Remark}{rem:hash-penalty}{One further bit for a merely two-universal family}
\begin{remark}[One further bit for a merely two-universal family]\label{rem:hash-penalty}
The R\'enyi leftover-hashing step of Theorem~\ref{thm:keylength} is taken
under the reading of \cite[Lemma~6.5]{tupkary2026decoy}, which requires an
ideal two-universal family, collision probability exactly $2^{-\ell}$ for
distinct inputs; that is Assumption~\ref{ass:hash}\,(i).
For a family that is only two-universal, collision probability at most
$2^{-\ell}$, the requirement can be relaxed at the cost of one bit of key
length \cite[Rem.~6.1]{tupkary2026decoy}. Toeplitz hashing is ideal
two-universal \cite[Rem.~6.1]{tupkary2026decoy}.
\end{remark}
\claimmeta{\metalabel{Uses}\mref{Thm.}{thm:keylength}
\metasep \metalabel{Used by}\mref{Req.}{req:hash-penalty}
\metasep published in \cite[{Rem. 6.1}]{tupkary2026decoy}; \cite[{Lemma 6.5}]{tupkary2026decoy}
\metasep accepted}

% claim meat-core/rem:alpha-domain
\claimsummary{meat-core/rem:alpha-domain}{Every key-length statement takes the R\'enyi order in (1, 2), the domain of the R\'enyi leftover-hashing lemma, while the entropy, program and gradient results hold on (1, \ensuremath{\infty}). On (1, 2) the key-length statements and the results they consume apply together.}
\claimbookmark{meat-core/rem:alpha-domain}{Remark}{rem:alpha-domain}{The Rényi order of the key-length statements}
\begin{remark}[The Rényi order of the key-length statements]\label{rem:alpha-domain}
Every key-length statement of this contract, among them
Theorem~\ref{thm:keylength}, Corollary~\ref{cor:keylength-bound} and the
key-length corollary of the first-order results
(Corollary~\ref{cor:fo-keylength}), takes $\alpha\in(1,2)$, because the
R\'enyi leftover-hashing lemma their secrecy argument applies is stated
there \cite[Thm.~4, Eq.~(13)]{kamin2025renyi}. The results on entropy
accumulation, the program, its reduction and its gradients are stated on
$(1,\infty)$, where \xref{foundations/def:renyi} defines the sandwiched
entropies, so the intersection of the two domains, on which the key-length
statements and the results they consume apply together, is $(1,2)$; a
result stated on a narrower range of $\alpha$ holds on that range only. The
accumulation import itself extends to $\alpha=\infty$, an endpoint that
\xref{foundations/def:renyi} does not define and no claim uses
(Fact~\ref{fact:meat}, locator note~(vi)).
\end{remark}
\claimmeta{\metalabel{Uses}\extref{../foundations/contract.pdf}{foundations/def:renyi}
\metasep \metalabel{Used by}\mref{Req.}{req:key-length-alpha-domain}
\metasep newly derived
\metasep accepted}

% claim meat-core/prop:monotone
\claimsummary{meat-core/prop:monotone}{The maximal key length does not increase with the reconciliation leakage, does not decrease as either error budget grows, and is unchanged by any change of the declaration that keeps h\_low. More leakage or a tighter budget therefore never lengthens it.}
\claimbookmark{meat-core/prop:monotone}{Proposition}{prop:monotone}{Monotonicity of the maximal key length}
\begin{proposition}[Monotonicity of the maximal key length]\label{prop:monotone}
The \termref{term-maximal-key-length}{maximal key length} $\ell_{\max}$ of
Definition~\ref{def:max-key-length}
is
\begin{enumerate}[label=(\roman*),leftmargin=2.4em]
\item non-increasing in $\lEC$;
\item non-decreasing in $\eev$;
\item non-decreasing in $\epa$;
\item unchanged under any modification of the declaration that leaves
  $\hlow$ unchanged.
\end{enumerate}
Consequently increasing leakage or tightening either error budget can never
increase $\ell_{\max}$.
\end{proposition}
\claimmeta{\metalabel{Uses}\mref{Thm.}{thm:keylength}, \mref{Def.}{def:max-key-length}
\metasep \metalabel{Used by}\mref{Cor.}{cor:keylength-bound}
\metasep newly derived
\metasep implemented (effective: accepted)
\metasep \metalabel{Code}\coderef{https://github.com/fraunhofersg/certiqs-key/blob/v0.1.0/src/certiqs_key/meat_core/finitekey/keylength.py}{certiqs_key.meat_core.finitekey.keylength:key_length}
\metasep \metalabel{Oracles}\coderef{https://github.com/fraunhofersg/certiqs-key/blob/v0.1.0/tests/oracles/registry.py}{registry.py}}

\begin{proof}
$\hlow$ does not depend on $(\lEC,\eev,\epa)$. In
\eqref{eq:keylength}: $-\lEC$ is decreasing in $\lEC$;
$-\lceil\log_2\frac1{\eev}\rceil$ is non-decreasing in $\eev$ (the
ceiling of a decreasing function of $\eev$ is non-increasing in
$\eev$); $-\frac{\alpha}{\alpha-1}\log_2\frac1{\epa}$ is increasing in
$\epa$ ($\frac{\alpha}{\alpha-1}>0$). The floor and the clamp
$\max\{0,\cdot\}$ are monotone.
\end{proof}

% claim meat-core/cor:keylength-bound
\claimsummary{meat-core/cor:keylength-bound}{Any h \ensuremath{\le} h\_low may replace h\_low in the key-length bound through max\{0, h\}, with the same security statement. This is where every solver obtains its key length; a bound h \ensuremath{\le} 0 yields no positive key when \ensuremath{\varepsilon}\_EV \textless{} 1.}
\claimbookmark{meat-core/cor:keylength-bound}{Corollary}{cor:keylength-bound}{Key length from any lower bound on the value}
\claimleadin{%
This corollary is where a solver obtains its key length, and the only
place: a solver submodule establishes a lower bound $h\le\hlow$ for the $h$
it computes, under its own trust root and in its own claims, and the
Frank--Wolfe linearisation bound, with both linear minima bounded from
below, is the instance this repository carries. The core asserts nothing about how
$h$ was obtained, in what arithmetic, or how close it lies to $\hlow$.
}
\begin{mainclaim}
\begin{corollary}[Key length from any lower bound on the value]\label{cor:keylength-bound}
Let $\alpha\in(1,2)$, let the hypotheses of Theorem~\ref{thm:keylength}
hold in its unconstrained or its constrained branch, the run as in
Lemma~\ref{lem:memory-lift}\,(a) and hypothesis (H-CHAIN) included, at every
all-at-once attack with finite registers, A-SEQ with its clause (v) among
them, and let those of Lemma~\ref{lem:ev-correctness} hold. Let
$h\in\mathbb{R}\cup\{-\infty\}$ satisfy
\begin{equation}\label{eq:lower-bound-h}
  h\;\le\;\hlow ,
\end{equation}
with $\hlow$ the exact value \eqref{eq:hlow-def} of the declaration (for
the constrained branch: with $\Feas$ carrying the declared $\CList$). If the
output length satisfies
\begin{equation}\label{eq:keylength-h}
  \ell\;\le\;n\,\max\{0,h\}-\lEC-\Big\lceil\log_2\frac1{\eev}\Big\rceil
  -\frac{\alpha}{\alpha-1}\log_2\frac1{\epa}+1,
\end{equation}
then the protocol is $\epa$-secret and $\eev$-correct on the accepted
state, hence, for $\ell\ge1$, $(\epa+\eev)$-secure in the
\termref{term-honest-authentication-setting}{honest-authentication setting}; under the further hypotheses of
Corollary~\ref{cor:composed-security}\,(iii) it is
$(\epa+\eev+2\varepsilon_{\mathrm{auth}})$-secure over the classical channel
of that part, with net key gain $\ell-k_{\mathrm{auth}}$ when both parties
output a key. The secrecy and the security statements hold at every output
state, whatever the dimension of the adversary's registers and the order in
which the source releases its signals, as
Corollary~\ref{cor:composed-security} states them, and the correctness
statement at every attack. Moreover:
\begin{enumerate}[label=(\alph*),leftmargin=2.2em]
\item if $h$ depends on none of $(\lEC,\eev,\epa)$, the four conclusions of
  Proposition~\ref{prop:monotone} hold verbatim for the maximal key length
  (Definition~\ref{def:max-key-length}) built from \eqref{eq:keylength-h},
  with $\max\{0,h\}$ in place of $\hlow$;
\item if $h\le0$ and $\eev<1$ then no positive $\ell$ satisfies
  \eqref{eq:keylength-h} and no key is claimed (at $\eev=1$ the correctness
  statement is empty and the corner carries no security claim).
\end{enumerate}
\end{corollary}
\end{mainclaim}
\claimmeta{\metalabel{Uses}\mref{Prop.}{prop:value-range}, \mref{Lem.}{lem:ev-correctness}, \mref{Lem.}{lem:memory-lift}, \mref{Lem.}{lem:honest-attacks}, \mref{Thm.}{thm:keylength}, \mref{Cor.}{cor:composed-security}, \mref{Prop.}{prop:monotone}
\metasep \metalabel{Used by}\mref{Req.}{req:reported-key-length}, \mref{Req.}{req:key-length-alpha-domain}, \mref{Req.}{req:compiled-declaration-only}, \mref{Req.}{req:security-statement-not-vacuous}, \mref{Cor.}{cor:fo-keylength}, \mref{Cor.}{cor:certified-keylength}, \mref{Cor.}{cor:certified-keylength-reference}, \mref{Rem.}{rem:fo-returned-point}
\metasep newly derived
\metasep implemented (effective: accepted)
\metasep \metalabel{Code}\coderef{https://github.com/fraunhofersg/certiqs-key/blob/v0.1.0/src/certiqs_key/meat_core/finitekey/keylength.py}{certiqs_key.meat_core.finitekey.keylength:key_length_from_bound}
\metasep \metalabel{Oracles}\coderef{https://github.com/fraunhofersg/certiqs-key/blob/v0.1.0/tests/oracles/registry.py}{registry.py}}

\begin{proof}
\eqref{eq:keylength-h} is \eqref{eq:keylength} with $\hlow$ replaced by
$\max\{0,h\}$ and \emph{nothing else changed}: no additional subtrahend,
no floor beyond the one $\max\{0,\cdot\}$ already supplies, no arithmetic
requirement of any kind. By \eqref{eq:lower-bound-h} and $\hlow\ge0$
(Proposition~\ref{prop:value-range}\,(i)) one has
$\max\{0,h\}\le\hlow$; the right-hand side of \eqref{eq:keylength} is
non-decreasing in $\hlow$ and $n\ge1$; so \eqref{eq:keylength-h} implies
\eqref{eq:keylength} and Theorem~\ref{thm:keylength} applies unchanged, as
does Corollary~\ref{cor:composed-security}, whose hypotheses are those of
that theorem at an $\ell$ satisfying \eqref{eq:keylength} and those of
Lemma~\ref{lem:ev-correctness}, which holds for every $\ell$ and every
attack; the corollary states the secrecy and its three security statements at
every output state, whatever the adversary's memory.
(a) Under the stated independence, the proof of
Proposition~\ref{prop:monotone} carries over word for word with
$\max\{0,h\}$ in place of $\hlow$: the only property of $\hlow$ that proof
reads is that it does not depend on $(\lEC,\eev,\epa)$. (b) If $h\le0$ then
$\max\{0,h\}=0$ and the right-hand side of \eqref{eq:keylength-h} is
$1-\lEC-\lceil\log_2\frac1{\eev}\rceil-\frac{\alpha}{\alpha-1}\log_2\frac1{\epa}$,
and for $\eev<1$ the ceiling is at least $1$ and the other two subtrahends
are non-negative, so the right-hand side is at most $0$ and no positive
$\ell$ satisfies \eqref{eq:keylength-h} --- the conservative outcome. The
hypothesis $\eev<1$ is not removable: at $\eev=1$, $\epa=1$ and $\lEC=0$,
all inside the hypotheses of Theorem~\ref{thm:keylength}, the right-hand
side equals $1$ and $\ell=1$ satisfies \eqref{eq:keylength-h} under a
$(\epa+\eev)$-security statement that is empty; the corner carries no
security claim, which is why the clause excludes it rather than claiming
it.
\end{proof}

% claim meat-core/cor:constrained-key-length
\claimsummary{meat-core/cor:constrained-key-length}{Under the attack-set licence the accumulated entropy is at least n times the value of the program under the licensed list, less the acceptance-probability penalty. That is the form the key-length theorem reads, so its key length holds with the constrained value.}
\claimbookmark{meat-core/cor:constrained-key-length}{Corollary}{cor:constrained-key-length}{Key length from the constrained value}
\begin{corollary}[Key length from the constrained value]
\label{cor:constrained-key-length}
Under the hypotheses of Theorem~\ref{thm:constrained-accumulation}, with
Corollary~\ref{cor:reduction-constrained},
\begin{equation*}
  \Hsandup\big(Z_1^n\big|\Chat_1^nI_1^n\Ehat\big)_{\rho_{|\OmAT}}
  \;\ge\;n\,\hlow(\CList)
  -\frac{\alpha}{\alpha-1}\log_2\frac{1}{\Pr[\OmAT]},
\end{equation*}
which is the form Theorem~\ref{thm:keylength} consumes; under
A-ATK (\xref{protocol-compiler/ass:atk}) and the remaining
hypotheses of Theorem~\ref{thm:keylength}, the key-length statement
of Theorem~\ref{thm:keylength} therefore holds with the constrained $\hlow$
(its constrained branch).
\end{corollary}
\claimmeta{\metalabel{Uses}\mref{Thm.}{thm:constrained-accumulation}, \mref{Cor.}{cor:reduction-constrained}, \mref{Thm.}{thm:keylength}
\metasep \metalabel{Assumes}\extaref{../protocol-compiler/contract.pdf}{protocol-compiler/ass:atk}{A-ATK}
\metasep newly derived
\metasep implemented (effective: accepted)
\metasep \metalabel{Code}\coderef{https://github.com/fraunhofersg/certiqs-key/blob/v0.1.0/src/certiqs_key/meat_core/finitekey/keylength.py}{certiqs_key.meat_core.finitekey.keylength:key_length_permitted}
\metasep \metalabel{Oracles}\coderef{https://github.com/fraunhofersg/certiqs-key/blob/v0.1.0/tests/oracles/registry.py}{registry.py}}

\begin{proof}
Theorem~\ref{thm:constrained-accumulation} gives
$\Hsandup(Z_1^n|\Chat_1^nI_1^n\Ehat)_{\rho_{|\OmAT}}
\ge n\,\hqkd(\CList)-\frac{\alpha}{\alpha-1}\log_2\frac{1}{\Pr[\OmAT]}$
under its hypotheses, and Corollary~\ref{cor:reduction-constrained} gives
$\hqkd(\CList)\ge\hlow(\CList)$ for the same licensed list; since $n\ge1$,
substituting the second inequality into the first yields the display. That
display is the accumulation bound the secrecy step of
Theorem~\ref{thm:keylength} reads, with $\hlow(\CList)$ in the place of
$\hlow$, so the constrained branch of that theorem applies under the
hypotheses of Theorem~\ref{thm:constrained-accumulation}, which include A-ATK,
together with the remaining hypotheses of Theorem~\ref{thm:keylength}: its
post-processing model and the assumptions it names, $\alpha\in(1,2)$ and
$\epa,\eev\in(0,1]$.
\end{proof}

% claim meat-core/rem:cgate
\claimsummary{meat-core/rem:cgate}{In either format the licensed-class half of the constraint gate is a genuine type check in trusted-source mode at d\_A \ensuremath{\ge} 2, vacuous in untrusted-source mode, and at d\_A = 1 rejects only what a producer can trivially avoid. In the last two cases the key length rests on the attack-set licence alone.}
\claimbookmark{meat-core/rem:cgate}{Remark}{rem:cgate}{Degeneracy of the licensed-class half of the gate}
\claimleadin{%
The numerical machinery --- Definition~\ref{def:feasible},
Theorems~\ref{thm:convexity}, \ref{thm:removal}, \ref{thm:choi} and
\ref{thm:gradient} --- is valid for arbitrary $\CList$, with
$\hlow(\CList)\ge\hlow(\varnothing)$; what the two-part gate on $\CList$
(Requirement~\ref{req:constrained-key-gate}) decides is only whether its
output may reach the key length of Theorem~\ref{thm:keylength}.
}
\begin{remark}[Degeneracy of the licensed-class half of the gate]\label{rem:cgate}
Part (i) of the two-part gate --- ``the list lies in the licensed class
of the declared source mode'' --- is the half a program can check. How
much it checks depends on the source mode and on $d_A$.

\emph{\exttermref{../transmitter/contract.pdf}{term-trusted-source-mode}{Trusted-source mode} with $d_A\ge2$:} the restriction is genuine
--- the multiples of $\tau_A$ are a one-dimensional subspace of the
$d_A^2$-dimensional real space of Hermitian operators on $A$, so a
\exttermref{../protocol-compiler/contract.pdf}{term-general}{general} entry is in general not a licensed one --- and a declaration
can fail part~(i). \emph{The check is nevertheless a \emph{type} check}: an
entry declares either its scalar $c_k$ or an operator $\Theta_k$
(Definition~\ref{def:homog}), and part~(i) inspects which, rather than
comparing the operator against $c_k\tau_A$. What is being rejected is a
declaration that \emph{says} it is general, never one whose numbers came out
near a multiple; the class is the same and the decision is exact.

\emph{$d_A=1$:} it is a type check that rejects nothing a producer cannot
trivially avoid. An entry declared \exttermref{../protocol-compiler/contract.pdf}{term-general}{general} fails it, as at every
$d_A$ (Definition~\ref{def:homog}); but every $1\times1$ Hermitian
$\Theta_k$ equals $\Theta_k\cdot[1]=c_k\tau_A$
(Proposition~\ref{prop:trivial-alice-system}\,(iii)), so such an entry could
always have been declared homogeneous with the same condition, and no
declaration \emph{need} fail part~(i).

\emph{\exttermref{../transmitter/contract.pdf}{term-untrusted-source-mode}{Untrusted-source mode}, any $d_A$:} it is
\emph{vacuously true}. The entry \emph{type} is scalar by
Definition~\ref{def:declaration}, so part~(i) degenerates from a
mathematical condition to a type check that a declaration passes by
construction whenever it is well formed.

In both degenerate cases the constrained branch of
Theorem~\ref{thm:keylength} therefore rests \emph{entirely} on
part~(ii), the attack-set licence of A-ATK
(\xref{protocol-compiler/ass:atk}), which no program can verify. Arguably it
is the correct degeneration: in the no-marginal model the licensed class
\emph{is} the scalar class (Remark~\ref{rem:homog}\,(a)), so there is no
wider class for part~(i) to exclude, and the honest statement is that the
attestation is doing all of the work. What an implementation reports of a
passing part~(i) in these cases is Requirement~\ref{req:degenerate-gate-report}.

\emph{The general format.} The scalar class named for the no-marginal model
is the licensed class of the \exttermref{../protocol-compiler/contract.pdf}{term-basic-format}{basic format}. In the general format of
Definition~\ref{def:declaration} the licensed class of
Definition~\ref{def:homog} contains, in \emph{either} mode, the entries
declared \exttermref{../protocol-compiler/contract.pdf}{term-party-local}{party-local} beside the mode's own form; so in the
no-marginal model the licensed class is the scalar form \emph{together with}
party-local. \emph{And those two are exactly the forms that are
expressible there}: a general right-hand side does not exist in
untrusted-source mode at all (Definition~\ref{def:homog};
Definition~\ref{def:admissible-declaration}\,(vi) in its general-format reading), so it is
refused by well-formedness before part~(i) is reached. \emph{The licensed
class therefore grew by exactly what the expressible class grew by, part~(i)
excludes nothing new, and the degeneration stands}, along the same two
readings: \emph{in the no-marginal model} part~(i) remains a type check that
a declaration passes \emph{by construction} whenever it is well formed, both
expressible forms being licensed --- two forms there, not three; and
\emph{at $d_A=1$} it remains the weaker statement that it rejects nothing a
producer cannot trivially avoid, a general entry being expressible and
unlicensed in trusted-source mode while every $1\times1$ Hermitian
$\Theta_k$ equals $c_k\tau_A$, so that such an entry could always have been
declared homogeneous instead. In both readings the honest statement is still
that the attestation is doing all of the work. In trusted-source mode at
$d_A\ge2$ part~(i) is not degenerate, in either format. \emph{What the third
form does change} is what part~(ii) carries: a tagged entry makes the
attestation attest \emph{more} --- the certificate on the scoped factor's
own channel (\xref{protocol-compiler/rem:party-local-soundness}) --- so the
rules of the gate bind it exactly as they bind the others
(Requirements~\ref{req:constrained-key-gate}
and~\ref{req:degenerate-gate-report}).
\end{remark}
\claimmeta{\metalabel{Uses}\mref{Def.}{def:declaration}, \mref{Def.}{def:admissible-declaration}, \mref{Def.}{def:homog}, \mref{Rem.}{rem:homog}, \mref{Def.}{def:feasible}, \mref{Prop.}{prop:trivial-alice-system}, \mref{Thm.}{thm:convexity}, \mref{Thm.}{thm:removal}, \mref{Thm.}{thm:choi}, \mref{Thm.}{thm:gradient}, \mref{Thm.}{thm:keylength}
\metasep \metalabel{Assumes}\extaref{../protocol-compiler/contract.pdf}{protocol-compiler/ass:atk}{A-ATK}, \extaref{../transmitter/contract.pdf}{transmitter/ass:src}{A-SRC}
\metasep \metalabel{Used by}\mref{Req.}{req:degenerate-gate-report}
\metasep newly derived
\metasep accepted}

\section{Perturbations of the declared measurement}\label{group-meat-core/prop:optol-nobound}
\sectionintro{Declared data are binary64 numbers, and a declaration is exactly admissible when the conditions of \hyperlink{claim-meat-core/def:declaration}{Definition~\ref*{def:declaration}} hold of those numbers themselves, among them operators exactly Hermitian, measurements and marginal exactly positive semidefinite, and completeness sums, formed in the declared order with no addition rounding, exactly the identity. One might hope that two such declarations whose measurements on Q differ by at most \ensuremath{\epsilon} have nearby rates; \hyperlink{claim-meat-core/prop:optol-nobound}{Proposition~\ref*{prop:optol-nobound}} shows that they need not. On data of fixed small dimensions their rate bounds can lie \ensuremath{\alpha}/(\ensuremath{\alpha}\ensuremath{-}1)\ensuremath{\cdot}N bits per round apart, for every positive integer N \ensuremath{\le} 53 + log\ensuremath{_{2}}\ensuremath{\epsilon}, which is up to 52 (\hyperlink{claim-meat-core/rem:representability-ceiling}{Remark~\ref*{rem:representability-ceiling}}). The gap sits in the acceptance-test penalty, can point either way and at the larger N is far above anything a key-length claim could absorb (\hyperlink{claim-meat-core/rem:optol-direction}{Remark~\ref*{rem:optol-direction}}). The engine therefore admits declared data only when they are exactly admissible (\hyperlink{claim-meat-core/req:exact-admission}{Requirement~\ref*{req:exact-admission}}), and the declared measurement model is itself a premise of every reported rate; where a reference of it is held, the first-order bound is taken on the reference and the snap onto the lattice only steers the search, unless a number certified on the snapped declaration carries the key length, and no claim bounds the gap between the model and the device (\hyperlink{claim-meat-core/rem:declared-measurement}{Remark~\ref*{rem:declared-measurement}}).}
\sectionimports{\extref{../foundations/contract.pdf}{foundations/def:evaluated-off-snap}, \extref{../foundations/contract.pdf}{foundations/def:kl}, \extref{../foundations/contract.pdf}{foundations/def:snap}, \extref{../foundations/contract.pdf}{foundations/rem:noise-floor}, \extref{../foundations/contract.pdf}{foundations/req:snap-reference}}

% claim meat-core/prop:optol-nobound
\claimsummary{meat-core/prop:optol-nobound}{Two exactly admissible measurements on Q within \ensuremath{\epsilon} of each other can have rate bounds \ensuremath{\alpha}/(\ensuremath{\alpha}\ensuremath{-}1)\ensuremath{\cdot}N bits per round apart, for every positive integer N \ensuremath{\le} 53 + log\ensuremath{_{2}}\ensuremath{\epsilon} (up to 52). So closeness of the declared measurement cannot be relied on to carry a key-rate claim to a nearby declaration.}
\claimbookmark{meat-core/prop:optol-nobound}{Proposition}{prop:optol-nobound}{Instability of the rate under measurement perturbations}
\claimleadin{%
One might admit a declaration whose measurement operators are only close
to exact ones, with an error allowance depending on the distance. The
following proposition shows that no allowance computed from the declared
dimensions, alphabet sizes and $\alpha$ can be useful. It compares values
$\hlow$ of \eqref{eq:hlow-def}, the infimum of the objective
\eqref{eq:objective} over the acceptance set and the feasible states. That
objective adds to its entropy terms the acceptance-test penalty
$\frac{\alpha}{\alpha-1}\sfD(\mathbf q\|\Phimap[\omega])$, where
$\mathbf q$ is a frequency distribution from the acceptance set and
$\Phimap[\omega]$ the distribution of the test statistic under $\omega$;
the penalty grows without bound as $\mathbf q$ demands an outcome that
$\omega$ makes improbable. The hypothesis $2^{N}\le\epsilon\,2^{53}$ is
what the binary64 number model allows
(Remark~\ref{rem:representability-ceiling}). Part~(iii) is arithmetic on
the formula $\ell_{\max}$ of Definition~\ref{def:max-key-length}, which
Theorem~\ref{thm:keylength} reads as a key length on $(1,2)$ under its own
assumptions.
}
\begin{proposition}[Instability of the rate under measurement perturbations]\label{prop:optol-nobound}
Let $\alpha\in(1,\infty)$ and $B>0$, write $N\coloneqq\lceil B\rceil$, and
let $\epsilon\in(0,\tfrac12]$ satisfy $2^{N}\le\epsilon\,2^{53}$.
Write $\Decl^{(i)}$ for the declaration obtained from a given one by
replacing its measurement $\{F_v\}$ on $\Qsys$ by $\{F^{(i)}_v\}$, and
$\hlow[\,\cdot\,]$ for the value \eqref{eq:hlow-def} of a declaration.
\begin{enumerate}[label=(\roman*),leftmargin=2.4em]
\item There is a declaration $\Decl$ satisfying Definition~\ref{def:admissible-declaration}
in full, with $\{F_v\}$ itself \exttermref{../foundations/contract.pdf}{term-exactly-admissible}{exactly admissible}, and there are two
exactly admissible families $\{F^{(1)}_v\}$, $\{F^{(2)}_v\}$ with
$\max_{v}\|F^{(i)}_v-F_v\|\le\epsilon$ for $i=1,2$, such that
\begin{equation}\label{eq:optol-gap}
  \hlow\big[\Decl^{(1)}\big]-\hlow\big[\Decl^{(2)}\big]\;\ge\;B .
\end{equation}
The dimensions and alphabets of $\Decl$ are $d_A=d_{\Qsys}=1$, one
sector, $|\Xset|=1$, $|\Vset|=|\Chatset|=2$ and $d_Z=1$, the same for
every $B$ and every $\epsilon$.
\item Let $\mathrm{err}$ be a function of the declared dimensions,
alphabet sizes and $\alpha$ whose value on the data of (i) is less than
$\frac{\alpha}{\alpha-1}\,N$. Then the bound
$\big|\hlow[\Decl^{(1)}]-\hlow[\Decl^{(2)}]\big|\le\mathrm{err}$, with
$\mathrm{err}$ evaluated on the data of the declaration, fails for some
declaration and some pair of exactly admissible families within
$\epsilon$ of each other.
\item For $\alpha\in(1,2)$ and every number $n$ of rounds, the
\termref{term-maximal-key-length}{maximal admissible key lengths} $\ell_{\max}$ of
$\Decl^{(1)}$ and $\Decl^{(2)}$
(Definition~\ref{def:max-key-length}), computed with the same $\lEC$, $\eev$ and
$\epa$, differ by at least $nB-1$ bits whenever the smaller of the two is
positive.
\end{enumerate}
\end{proposition}
\claimmeta{\metalabel{Uses}\mref{Def.}{def:declaration}, \mref{Def.}{def:admissible-declaration}, \mref{Def.}{def:feasible}, \mref{Prop.}{prop:trivial-alice-system}, \mref{Prop.}{prop:value-range}, \mref{Thm.}{thm:keylength}, \mref{Def.}{def:max-key-length}, \extref{../foundations/contract.pdf}{foundations/def:kl}
\metasep \metalabel{Used by}\mref{Rem.}{rem:declared-measurement}, \mref{Req.}{req:exact-admission}, \mref{Rem.}{rem:fo-arithmetic}
\metasep newly derived
\metasep accepted
\metasep \metalabel{Oracles}\coderef{https://github.com/fraunhofersg/certiqs-key/blob/v0.1.0/tests/oracles/registry.py}{registry.py}}

\begin{proof}
\emph{Idea.} The instance has one-dimensional systems, so a single state;
the measurement on $\Qsys$ has two outcomes, $0$ with probability $g_0$
and $1$ with probability $g_1$; the acceptance test demands outcome~$0$ in
every round; and no round contributes key, so no acceptance value is
key-live. Then $\hlow$ is the acceptance-test penalty alone,
$\frac{\alpha}{\alpha-1}\log_2\frac1{g_0}$. Dividing $g_0$, the largest
power of two at most $\epsilon$, by $2^{N}$ raises $\hlow$ by
$\frac{\alpha}{\alpha-1}N$ and moves each entry of the family by at most
$\epsilon$, so the declared family and its perturbation lie within
$\epsilon$ of each other.

\emph{Step 1 (the instance).} Take the instance of
Proposition~\ref{prop:value-range}\,(iii): $d_A=d_{\Qsys}=1$ (so both source modes
coincide, Proposition~\ref{prop:trivial-alice-system}, and $\Feas=\{[1]\}$), one
sector, $\Xset=\{0\}$ with $M_0=[1]$, $\Vset=\{0,1\}$, $a_A$ constant and
$a_B=\idmap$, $\phi(\kappa)=c_B$ so that $\Chatset=\Vset$, all-fail
sifting ($\varsigma\equiv\mathrm{fail}$, $r\equiv\perp$, $\hold=\mathsf B$
with $r_{\mathsf B}\equiv\perp$, hence $\Chkey=\varnothing$),
$\CList=\varnothing$, the default \exttermref{../protocol-compiler/contract.pdf}{term-party-partition}{party partition}, and the singleton
acceptance set $\acSet=\{(1,0)\}$. For scalars $g_0,g_1\ge0$ write
$\Decl_{(g_0,g_1)}$ for the declaration whose measurement on $\Qsys$ is
$F_0=[g_0]$, $F_1=[g_1]$. Every clause of Definition~\ref{def:admissible-declaration} other
than the $\{F_v\}$-clause of (iii) holds for it: all systems and
alphabets are finite, sifting is announcement-determined, the key map is
party-computable with $r_{\mathsf B}\equiv\perp$,
$\Feas=\{[1]\}\ne\varnothing$, $\acSet$ is closed and convex, and
$\dim\Hs{\Etil}\ge1$. That clause holds exactly when $g_0+g_1=1$, and the
family is exactly admissible in the sense of
Definition~\ref{def:declaration} exactly when in addition both entries are
\exttermref{../foundations/contract.pdf}{term-binary64-numbers}{binary64 numbers} and the one addition of \eqref{eq:declorder} is exact.
Since $\Chkey=\varnothing$, the objective \eqref{eq:objective} keeps only
its first sum, the acceptance-test penalty
$\frac{\alpha}{\alpha-1}\sum_{\hat c}q(\hat c)\log_2\frac{q(\hat c)}{P_{\hat c}(\omega)}$,
and $\hlow$ of \eqref{eq:hlow-def} is its infimum over
$\mathbf q\in\acSet=\{(1,0)\}$ and $\omega\in\Feas$. $\Feas$ is the single
state $[1]$, at which $P_{\hat c_0}=g_0$ and $P_{\hat c_1}=g_1$; with the
convention $0\log_2\frac0v=0$ of \xref{foundations/def:kl} this gives
\begin{equation}\label{eq:optol-instance}
  \hlow\big[\Decl_{(g_0,g_1)}\big]
  =\frac{\alpha}{\alpha-1}\Big(1\cdot\log_2\frac1{g_0}
  +0\cdot\log_2\frac{0}{g_1}\Big)
  =\frac{\alpha}{\alpha-1}\,\log_2\frac1{g_0}
  \qquad (g_0>0,\ g_1\ge0),
\end{equation}
a computation that uses no property of $(g_0,g_1)$ beyond those two
inequalities, in particular not $g_0+g_1=1$.

\emph{Step 2 (part (i)).} Let $\epsilon'$ be the largest power of two
with $\epsilon'\le\epsilon$, so $\epsilon'=2^{-k}$ for an integer
$k\ge1$. The hypothesis says $2^{N-53}\le\epsilon$, and $2^{N-53}$ is a
power of two, so $2^{N-53}\le\epsilon'$ as well. Take
$\Decl\coloneqq\Decl_{(\epsilon',\,1-\epsilon')}$,
$\{F^{(2)}_v\}\coloneqq\{F_v\}$ (distance $0$) and $\{F^{(1)}_v\}$ the
POVM of $\Decl_{(\epsilon'2^{-N},\,1-\epsilon'2^{-N})}$, whose distance to
$\{F_v\}$ is $\epsilon'(1-2^{-N})\le\epsilon'\le\epsilon$ in each of its
two entries. All three families have non-negative entries summing to $1$;
each has $|\Vset|=2$, so \eqref{eq:declorder} performs one addition. By
$2^{N-53}\le\epsilon'$, $\epsilon'2^{-N}=2^{-k-N}$ is at least $2^{-53}$,
so $\epsilon'2^{-N}$ and $\epsilon'$ are powers of two in
$[2^{-53},\tfrac12]$, hence binary64 numbers and multiples of $2^{-53}$;
$\epsilon=\tfrac12$ with $N\le52$ is one such choice. Then every entry is
exactly representable ($1-\epsilon'$ and $1-\epsilon'2^{-N}$ are
multiples of $2^{-53}$ in $[\tfrac12,1)$, where consecutive binary64
numbers are $2^{-53}$ apart), and the single addition returns the
representable $1$ and is therefore exact, which is the clause of
Definition~\ref{def:declaration} that no addition of \eqref{eq:declorder}
rounds. So all three families are exactly admissible and $\Decl$
satisfies Definition~\ref{def:admissible-declaration} in full. By \eqref{eq:optol-instance},
\begin{equation*}
  \hlow\big[\Decl^{(1)}\big]-\hlow\big[\Decl^{(2)}\big]
  =\frac{\alpha}{\alpha-1}
  \Big(\log_2\frac{1}{\epsilon'2^{-N}}-\log_2\frac1{\epsilon'}\Big)
  =\frac{\alpha}{\alpha-1}\,N\;\ge\;N\;\ge\;B ,
\end{equation*}
using $\frac{\alpha}{\alpha-1}>1$ for $\alpha>1$.

\emph{Step 3 (part (ii)).} Let $\mathrm{err}$ be as in (ii). The
declaration $\Decl$ of Step~2 has the dimensions and alphabets listed in
(i), which are the same for every $B$ and $\epsilon$, so $\mathrm{err}$
takes one value on its data, and that value is less than
$\frac{\alpha}{\alpha-1}\,N$. The families $\{F^{(1)}_v\}$ and
$\{F^{(2)}_v\}=\{F_v\}$ of Step~2 are exactly admissible and lie within
$\epsilon$ of each other, and their values differ by
$\frac{\alpha}{\alpha-1}\,N$, which exceeds the prescribed $B$ for every
$B$ the representability condition admits and exceeds that value of
$\mathrm{err}$. So the bound fails for this pair. Since the value of
$\mathrm{err}$ does not depend on $B$, the argument may be run at the
largest $N$ the representability condition admits.

\emph{Step 4 (part (iii)).} Let $\alpha\in(1,2)$. By
Definition~\ref{def:max-key-length}, for $n$ rounds
\begin{equation*}
  \ell_{\max}=\max\Big\{0,\Big\lfloor n\hlow-\lEC
  -\Big\lceil\log_2\frac1{\eev}\Big\rceil
  -\frac{\alpha}{\alpha-1}\log_2\frac1{\epa}+1\Big\rfloor\Big\};
\end{equation*}
the terms after $n\hlow$ do not depend on $\hlow$ and are the same for
both declarations. Insert the gap \eqref{eq:optol-gap}: if the smaller of the
two lengths is positive, the two lengths differ by at least $nB-1$ bits.
\end{proof}

% claim meat-core/rem:representability-ceiling
\claimsummary{meat-core/rem:representability-ceiling}{The condition N \ensuremath{\le} 53 + log\ensuremath{_{2}}\ensuremath{\epsilon} of \hyperlink{claim-meat-core/prop:optol-nobound}{Proposition~\ref*{prop:optol-nobound}} is what exact admissibility forces on a two-outcome binary64 family. It caps N at 52, attained at \ensuremath{\epsilon} = 1/2, and belongs to the number model, not to the map from the measurement to the rate.}
\claimbookmark{meat-core/rem:representability-ceiling}{Remark}{rem:representability-ceiling}{The representability ceiling}
\begin{remark}[The representability ceiling]\label{rem:representability-ceiling}
Proposition~\ref{prop:optol-nobound} needs a declared family $(g_0,g_1)$
with $g_0+g_1=1$ and both entries binary64: exact admissibility is an
identity of the declared numbers (Definition~\ref{def:declaration}, the
exact-admissibility paragraph), not of the reals they approximate, and it
additionally requires the one addition of \eqref{eq:declorder} to be
exact. The smaller entry $g_0$ satisfies $g_0\le\tfrac12$, so
$g_1=1-g_0$ lies in the binade $[\tfrac12,1)$, the binary64 numbers with
exponent $-1$, where consecutive numbers are $2^{-53}$ apart (the unit in
the last place); hence $g_1$ is representable only if $g_0$ is a multiple
of $2^{-53}$, and $g_0>0$ makes that $g_0\ge2^{-53}$. With
$g_0=\epsilon2^{-N}$ this reads $\epsilon2^{-N}\ge2^{-53}$, that is
$2^{N}\le\epsilon\,2^{53}$, the inequality the proposition assumes of
$N$ and $\epsilon$; since $\epsilon\le\tfrac12$, it gives $N\le52$. The
value is attained: $\epsilon=\tfrac12$ and $N=52$ give
$(2^{-53},1-2^{-53})$, with the addition exact. The proof builds its
families at the largest power of two $\epsilon'\le\epsilon$, with
$g_0=\epsilon'2^{-N}$, and loses no $N$ by it: $2^{N-53}$ is itself a
power of two, so the inequality holds at $\epsilon'$ exactly when it
holds at $\epsilon$.

The ceiling is a property of the number model at these fixed dimensions,
not of the map $\{F_v\}\mapsto\hlow$: in exact real arithmetic the same
construction runs at every $B>0$, and it is the number model that stops
it. What the number model does is cap how much of the gap a family of
finite value, $g_0>0$, can express, at $\frac{\alpha}{\alpha-1}N$ with
$2^{N}\le\epsilon\,2^{53}$; the exactly admissible family $(0,1)$ has the
value $\hlow=+\infty$ (Proposition~\ref{prop:value-range}\,(iv)) and carries
no key-length claim (Requirement~\ref{req:infinite-value}), so it is set
aside. At $\epsilon=2^{-40}$ the instance still
delivers $N=13$, thirteen bits per round or more, and it runs down to
$\epsilon=2^{-52}$, where $N=1$ still gives $\frac{\alpha}{\alpha-1}>1$
bit per round. At the instance's $d_Z=1$ both exceed
$\log_2|\Zset|=1$ bit per round, the entropy-part ceiling of
Proposition~\ref{prop:value-range}\,(ii): more than the entropy terms of $\hlow$ can
contribute at any declaration with $d_Z=1$, the rest of $\hlow$ being the
acceptance-test penalty. The comparison is the instance's: at a larger
$d_Z$ the entropy-part ceiling $\log_2(d_Z+1)$ can exceed a gap of that
size (at $\alpha=1.9$, $d_Z=7$ and $N=1$ it is $3$, against
$\frac{\alpha}{\alpha-1}N\approx2.1$). The representability ceiling
shapes part~(ii) of the proposition, which rules out allowances below the
exhibited scale rather than every allowance. It shapes as well
Remark~\ref{rem:classes},
which cites part~(i) to explain why a tolerance may refuse the deletion of
a sector but never license one, and which states the gap there as
$\frac{\alpha}{\alpha-1}N$ with $N$ up to $52$. The conclusion that
matters is what the bounded statement still delivers: on the data of
part~(i), part~(ii) with $N=1$ refutes, at every
$\epsilon\in[2^{-52},\tfrac12]$, every $\mathrm{err}$ below
$\frac{\alpha}{\alpha-1}$, so every $\mathrm{err}$ at or below that
entropy-part ceiling of $1$ bit per round. The cap moves the refuted range
from every allowance to the allowances below
$\frac{\alpha}{\alpha-1}\cdot52$, still far above that ceiling, and no
rule of this module turns on the difference.
\end{remark}
\claimmeta{\metalabel{Uses}\mref{Def.}{def:declaration}, \mref{Prop.}{prop:value-range}
\metasep newly derived
\metasep accepted}

% claim meat-core/rem:optol-direction
\claimsummary{meat-core/rem:optol-direction}{The gap of \hyperlink{claim-meat-core/prop:optol-nobound}{Proposition~\ref*{prop:optol-nobound}} involves no defect, lives in the acceptance-test penalty and can point either way. Among families of finite value the instance leaves an allowance of 53\ensuremath{\cdot}\ensuremath{\alpha}/(\ensuremath{\alpha}\ensuremath{-}1) bits per round or more unrefuted; at a key rate below one bit no key-length claim can rest on it.}
\claimbookmark{meat-core/rem:optol-direction}{Remark}{rem:optol-direction}{Scope of the perturbation gap}
\claimleadin{%
Proposition~\ref{prop:optol-nobound} perturbs only the measurement
$\{F_v\}$ on $\Qsys$, not $\tau_A$, $\{M_x\}$ or the constraint data, and
it compares values of $\hlow$, the number Theorem~\ref{thm:keylength}
turns into a key length and Theorem~\ref{thm:reduction} places below
$\hqkd$, not values of $\hqkd$ itself. Its pair stops at
$\frac{\alpha}{\alpha-1}\cdot52$, while the figure $53$ in the remark is
the largest value of $\log_2\frac1{g_0}$ at an exactly admissible family
of the instance with $g_0>0$, since such a family has $g_0\ge2^{-53}$
(Remark~\ref{rem:representability-ceiling}); between the two scales the
proposition decides nothing. Its part~(iii) applies when the smaller of
the two key lengths is positive, which depends on $n$ and on the other
terms of $\ell_{\max}$; the proposition does not say for which $n$ that
happens on the instance.
}
\begin{remark}[Scope of the perturbation gap]\label{rem:optol-direction}
Proposition~\ref{prop:optol-nobound} does not say that a small
perturbation always matters: it says that no bound in the declared
metadata alone can say that it does not. It does not locate the
fragility in any defect (a failure of Hermiticity, positivity or
completeness in the declared numbers): both families are exact POVMs, so the phenomenon belongs to the map
$\{F_v\}\mapsto\hlow$ and would be there for a \exttermref{../foundations/contract.pdf}{term-producer}{producer} that emitted
perfect numbers for a slightly different device. It does not depend on
the entropy term: the gap lives entirely in the acceptance-test penalty
$\frac{\alpha}{\alpha-1}\sfD(\mathbf q\|\Phimap[\omega])$ of
\eqref{eq:objective}, with $\mathbf q$ ranging over the acceptance set, a
term every declaration carries, and it is the same unboundedness that
Proposition~\ref{prop:value-range}\,(iii) records for $\hlow$ itself, there
as a range statement and here as a stability statement.

And it does not refute every allowance: an $\mathrm{err}$ whose value on
the data of part~(i) is at least $\frac{\alpha}{\alpha-1}\cdot53$ is not
refuted at any admissible $\epsilon$ by a pair of families of the instance
of finite value, that is with $g_0>0$, and this module neither exhibits nor
rules out one. The family $(0,1)$ is exactly admissible as well and lies
within $2^{-53}\le\epsilon$ of $(2^{-53},1-2^{-53})$, but its value is
$\hlow=+\infty$ (Proposition~\ref{prop:value-range}\,(iv)), which carries
no key-length claim (Requirement~\ref{req:infinite-value}); it is set aside
here, since a gap to an infinite value is not a gap between two rates. On those data such an $\mathrm{err}$
would license a discrepancy of more than $53$ bits per round
($\frac{\alpha}{\alpha-1}>1$ at every $\alpha>1$, and it diverges as
$\alpha\downarrow1$) against a key rate whose entropy part is at most
$\log_2|\Zset|=1$ bit per round there ($d_Z=1$), by the entropy-part
ceiling of Proposition~\ref{prop:value-range}\,(ii) (the bound
$\log_2|\Zset|=\log_2(d_Z+1)$ on each entropy term, the rest of $\hlow$
being the acceptance-test penalty). For the protocols this module rates, a
key rate lies below one bit per round in practice, and against such a rate
a robustness allowance of tens of bits per round is not a statement any
key-length claim can rest on. At
a larger $d_Z$ the entropy-part ceiling grows, and the proposition
constrains $\mathrm{err}$ on the data of part~(i) only.

What the proposition does say is, first, that the direction is not
controlled: the gap \eqref{eq:optol-gap} has a sign, and exchanging the
roles of the two families reverses it, so a perturbation of a declared
family may move $\hlow$ up (for a key-length claim, the unsafe direction)
as easily as down. Second, that no threshold escapes the difficulty by
being small: $\epsilon$ may be taken as small as the number model
permits, the gap being bought by the second parameter of the instance,
$N$, rather than by the size of the perturbation. Third, for
$\alpha\in(1,2)$, that the gap reaches the output: by part~(iii) the key
lengths themselves differ, by at least $nB-1$ bits whenever the smaller is
positive, so the gap is not confined to an unused intermediate quantity. At $\alpha\ge2$ the same computation compares
two values of the expression of Definition~\ref{def:max-key-length}, to which this
module gives no key-length reading (Remark~\ref{rem:alpha-domain}).
\end{remark}
\claimmeta{\metalabel{Uses}\mref{Def.}{def:feasible}, \mref{Prop.}{prop:value-range}
\metasep \metalabel{Used by}\mref{Rem.}{rem:declared-measurement}, \mref{Rem.}{rem:fo-arithmetic}
\metasep newly derived
\metasep accepted
\metasep \metalabel{Oracles}\coderef{https://github.com/fraunhofersg/certiqs-key/blob/v0.1.0/tests/oracles/registry.py}{registry.py}}

% claim meat-core/rem:declared-measurement
\claimsummary{meat-core/rem:declared-measurement}{A tiny change of the declared measurement can move the key rate by many bits per round, so the model is a premise of the rate. Where a handed reference's companion counts and the key length is formed on it, the snap only steers the search; nothing bounds model to device.}
\claimbookmark{meat-core/rem:declared-measurement}{Remark}{rem:declared-measurement}{The declared measurement as a premise of the rate}
\claimleadin{%
A key rate of this contract is computed from a declaration whose
measurement on $\Qsys$ is a model of a device, and
Proposition~\ref{prop:optol-nobound} shows how much the number rests on that
model. A domination weight declared for the device, outright or derived from
ranges of its efficiencies and dark-count probabilities, would carry part of
the gap between model and device into the flags of a \exttermref{../reductions/contract.pdf}{term-domination-link}{domination link}; it is
future scope (\href{../../../spec/future-scope.md}{\nolinkurl{spec/future-scope.md}}~\S13).
}
\begin{remark}[The declared measurement as a premise of the rate]\label{rem:declared-measurement}
The value $\hlow$ is not controlled by a tolerance on the declared
measurement $\{F_v\}$ on $\Qsys$ (Definition~\ref{def:declaration}) at any
tolerance down to $2^{-52}$. For every $\epsilon\in(0,\tfrac12]$ and
every positive integer $N$ with $2^{N}\le\epsilon\,2^{53}$ there is a
declaration on which two \exttermref{../foundations/contract.pdf}{term-exactly-admissible}{exactly admissible} families within
$\epsilon$ of each other in every element give values
$\frac{\alpha}{\alpha-1}N$ bits per round apart, and no allowance computed
from the declared dimensions, alphabet sizes and $\alpha$ below that scale
bounds the change (Proposition~\ref{prop:optol-nobound}\,(i)--(ii)); below
$\epsilon=2^{-52}$ the proposition exhibits no such pair. For
$\alpha\in(1,2)$ the gap reaches the key length of
Theorem~\ref{thm:keylength}: over $n$ rounds the two lengths differ by at
least $nN-1$ bits whenever the smaller is positive
(Proposition~\ref{prop:optol-nobound}\,(iii)). The gap is bought by a ratio,
not by a size: on the instance of that proposition
$\hlow=\frac{\alpha}{\alpha-1}\log_2\frac1{g_0}$, where $g_0$ is the
probability of the one outcome the acceptance test demands, and moving
$g_0$ from $\epsilon'$, the largest power of two at most $\epsilon$, to
$\epsilon'2^{-N}$ changes the element by at most $\epsilon$ and the
probability by the factor $2^{N}$. Its limit is a change of support: at
$g_0=0$ the measurement treats the demanded outcome as impossible, no
feasible state gives it positive probability, and $\hlow=+\infty$
(Proposition~\ref{prop:value-range}\,(iv)), a value that carries no
key-length claim (Requirement~\ref{req:infinite-value}). A change of which
outcomes the measurement treats as impossible, or nearly so, can thus move
the rate by far more than the size of the change, in either direction; it
need not, since the proposition exhibits one declaration on which it does
and states no property of every declaration
(Remark~\ref{rem:optol-direction}). The declared measurement model is
therefore itself a premise of every number computed from it.

Between the device and the declared measurement lie two moves. The first is
the \exttermref{../foundations/contract.pdf}{term-snap}{snap} of a family read as a measurement onto its lattice, whose
displacement, however small, bounds the change of the rate by nothing.
Where the declarer holds a \exttermref{../foundations/contract.pdf}{term-reference}{reference} of the family and hands it
beside the snap, a claim establishes for the instance that the construction
may run on it (\xref{reductions/req:reference-companion}), the companion is
compiled and admitted (Requirement~\ref{req:companion-admission}), and the
run's bound on the companion counts
(Requirement~\ref{req:bracket-on-reference}\,(iii)) and its key length is
formed from a number on the companion, not from one certified on the snapped
declaration in the fallback of
Requirement~\ref{req:certified-evaluation}\,(i), the run need
not bound that change at all. Its first-order bound is then evaluated on the compiled reference companion
(Proposition~\ref{prop:reference-companion}), at the points the search on the
snapped declaration produced and at those of a short search continued on the
companion from them (Corollary~\ref{cor:fo-reference},
Requirement~\ref{req:bracket-on-reference}); the snap only steers the choice
of those points, and a bound that holds at any admissible point holds at
them. Every
block on which the companion carries the reference is
\exttermref{../foundations/contract.pdf}{term-evaluated-off-its-snap}{evaluated off its snap} (\xref{foundations/def:evaluated-off-snap}),
nothing the key length rests on reads its snapped data, and it rests on no
premise about the snap. In that fallback the number the key length is formed
from reads the snapped data of those blocks, which are then among the blocks
the run does not evaluate off its snap, below
(Requirement~\ref{req:bracket-on-reference}\,(v)). Outside that fallback,
the key length is one of the model the reference describes, and what lies between it and
the device is the second move below. The evaluation of an enclosed
reference at its centres is itself a displacement of the measured family, of
the size of the enclosure's radius, to which
Proposition~\ref{prop:optol-nobound} applies as it does to the snap's: the
premise about the snap is gone, the exposure to the radius is not, and the
recorded largest radius is the only handle on it
(Remark~\ref{rem:fo-arithmetic}). The rest of this paragraph concerns the
blocks a run does not evaluate off its snap. A block kept because its
reference is exact and already on the lattice rests on nothing, the snap
being the identity there. Every other block is a \exttermref{../foundations/contract.pdf}{term-moved}{moved} block of the
snap (\xref{foundations/def:snap}), and no claim relates it to the
reference: domination, which would bound the snapped family by the
reference, is not established for a snap, and an idealised model may
dominate no lattice family at a small weight at all
(\xref{foundations/rem:noise-floor}). On such a block the run proceeds,
unless even the last resort yields no snap, resting on
\xref{foundations/ass:snapped-measurement}: that on the blocks its record
names the device measures the snapped family, or, where the snapped family is
a producer's emission, that what the run concludes from the emission holds
of the snapped family. No claim then relates its key length to the
reference, and a run that declares the strict option refuses instead
(\xref{foundations/req:snap-reference}\,(iv)). The second move is
the gap between the model as declared, or its reference, and the physical
device. No claim bounds it: the run assumes, without checking, that the
model describes the device, for a receiver that the receiver module models
through \xref{receiver/ass:trust}, and by
Proposition~\ref{prop:optol-nobound} a small distance between model and
device bounds the change of the rate no more than a small distance between
two declared families does.
\end{remark}
\claimreading{The key rate is computed from the measurement model you declare for your receiver, and that model is a premise of the number: nothing checks that your device measures what it says (for a receiver the engine models, \hyperref{../receiver/contract.pdf}{}{claim-receiver/ass:trust}{A-TRUST}). A change of the model too small to see in its entries, notably one that makes an outcome nearly or exactly impossible, can move the rate by many bits per round in either direction; no result bounds that gap. The engine also moves your model onto its exact number grid (the snap) for its search. Where your model reaches the engine as a reference, the bound is evaluated on your model itself and needs no premise about the snap, unless the rate is certified only on the moved model, as the report then says; any other block rests on \hyperref{../foundations/contract.pdf}{}{claim-foundations/ass:snapped-measurement}{A-SNAP}, or the run refuses if it declares the strict option. Validate the model, above all which outcomes it treats as impossible.}
\claimmeta{\metalabel{Uses}\mref{Def.}{def:declaration}, \mref{Prop.}{prop:value-range}, \mref{Prop.}{prop:optol-nobound}, \mref{Rem.}{rem:optol-direction}, \mref{Prop.}{prop:reference-companion}, \mref{Req.}{req:infinite-value}, \extref{../foundations/contract.pdf}{foundations/def:evaluated-off-snap}, \extref{../foundations/contract.pdf}{foundations/def:snap}, \extref{../foundations/contract.pdf}{foundations/rem:noise-floor}, \extref{../foundations/contract.pdf}{foundations/req:snap-reference}
\metasep newly derived
\metasep accepted}

\section{The reference companion}\label{group-meat-core/prop:reference-companion}
\sectionintro{Where the receiver segment of one or more measured factors carries the reference companion, the same reduction run on the receiver's reference in place of its snap, the compiler forms beside the declaration a second tuple from the same specification, with every such factor replaced. It keeps every component of the declaration but the measured families and the entries of those segments, and so its joint sectors and the points of its program; it is admissible save possibly for a non-empty feasible set, and its entries, rounded with their constants relaxed, only enlarge that set (\hyperlink{claim-meat-core/prop:reference-companion}{Proposition~\ref*{prop:reference-companion}}). Where the reference carries factor enclosures and the reduction combines blocks with non-negative weights, the companion's blocks are members of the factor enclosures the reduction and the compiler carry to them (\hyperlink{claim-meat-core/cor:reference-companion-factors}{Corollary~\ref*{cor:reference-companion-factors}}). The Frank-Wolfe submodule evaluates the first-order bound on it, at points of the snapped search and of a short search continued on it, and certifies it.}
\sectionimports{\extref{../foundations/contract.pdf}{foundations/def:factor-enclosure}, \extref{../protocol-compiler/contract.pdf}{protocol-compiler/def:compile}, \extref{../protocol-compiler/contract.pdf}{protocol-compiler/def:weighted}, \extref{../protocol-compiler/contract.pdf}{protocol-compiler/lem:weighted}, \extref{../protocol-compiler/contract.pdf}{protocol-compiler/lem:weighted-factor}, \extref{../reductions/contract.pdf}{reductions/cor:companion-factor-enclosure}, \extref{../reductions/contract.pdf}{reductions/cor:reference-companion}, \extref{../reductions/contract.pdf}{reductions/def:reduction-instance}, \extref{../reductions/contract.pdf}{reductions/lem:row-relaxation}}

% claim meat-core/prop:reference-companion
\claimsummary{meat-core/prop:reference-companion}{Compiled with reference companions in place of the snapped targets of one or more measured factors, the declaration keeps every other component and is admissible save possibly for a non-empty feasible set. Its relaxed entries only enlarge that set.}
\claimbookmark{meat-core/prop:reference-companion}{Proposition}{prop:reference-companion}{The compiled reference companion is an admissible declaration}
\claimleadin{%
A run searches on the declaration $\Decl$ compiled from the snapped
measurement. Where the \exttermref{../reductions/contract.pdf}{term-receiver-segment}{receiver segment} of a \exttermref{../protocol-compiler/contract.pdf}{term-measured-factor}{measured factor} also
carries the reference companion, the compiler forms beside $\Decl$ a second
tuple from the same specification with the companion in place of the snapped
target, for every measured factor whose receiver segment carries one, both
stations of an entanglement-based declaration among them
(\xref{protocol-compiler/req:reference-companion-compiled}). The proposition
says that this tuple is a declaration of the kind this contract's results are
stated for, so that a lower bound on its value is a lower bound of the
ordinary kind, and that it shares with $\Decl$ everything a point of the
program is made of. It is stated for all the replaced factors at once: the
tuple with one factor replaced is not shown to have a non-empty
\exttermref{../protocol-compiler/contract.pdf}{term-feasible-set}{feasible set}, so the case of one factor cannot be applied again to it.
Its entries are the companions', each of which carries its own certificate at
the reference; an entry requested of $\Decl$ whose certificate is not
established at the reference is left out, which only enlarges the feasible
set, so that a bound on it stays a bound on the reference model, if a lower
one. Which of its conditions the engine decides, and which it imports, is
Requirement~\ref{req:companion-admission}.
}
\begin{proposition}[The compiled reference companion is an admissible declaration]\label{prop:reference-companion}
Let $\Decl$ be an \termref{term-admissible-declaration}{admissible declaration}
(Definition~\ref{def:admissible-declaration}) that the \exttermref{../protocol-compiler/contract.pdf}{term-compilation-map}{compilation map}
(\xref{protocol-compiler/def:compile}) emits from a specification, and let
$\mathbb{F}^{\mathrm{ref}}$ be a non-empty set of its
\exttermref{../protocol-compiler/contract.pdf}{term-measured-factor}{measured factors} (the one \exttermref{../reductions/contract.pdf}{term-measured-system}{measured system} in
the \exttermref{../protocol-compiler/contract.pdf}{term-basic-format}{basic format}) such that, for each $f\in\mathbb{F}^{\mathrm{ref}}$,
the last link of the receiver segment of $f$ is the emission of a
\exttermref{../reductions/contract.pdf}{term-reduction-instance}{reduction instance} $\mathcal{J}_f$ satisfying (H-REF) of
\xref{reductions/cor:reference-companion}. Let $\Decl^{\mathrm{ref}}$ be the
tuple the compilation map gives from the same specification with, for every
$f\in\mathbb{F}^{\mathrm{ref}}$, that link's \exttermref{../reductions/contract.pdf}{term-target-family}{target family} replaced by
the target $\{F_v[A^{\mathrm{ref}}_f]\}$ of the reference companion of
$\mathcal{J}_f$ of that corollary, $A^{\mathrm{ref}}_f$ its datum, and that
link's entries replaced by the relaxed entries
(\xref{reductions/lem:row-relaxation}) of the companion's entries, each
rounded operator Hermitian and with the sector blocks of the entry it rounds.
The companion's entries of such a link are the instance's own entries at the
companion's datum and the companion forms
(\xref{reductions/cor:reference-companion}\,(vi)) of those entries requested
of the link whose companion forms carry their certificates at that datum; a
requested entry whose companion form is not established to carry one has no
entry in $\Decl^{\mathrm{ref}}$. Then:
\begin{enumerate}[label=(\roman*),leftmargin=2.4em]
\item \emph{(what is shared)} $\Decl^{\mathrm{ref}}$ has the systems,
  sectors and dimensions, the alphabets, Alice's \exttermref{../protocol-compiler/contract.pdf}{term-marginal-slot}{marginal slot} and
  measurement, the \exttermref{../protocol-compiler/contract.pdf}{term-party-partition}{party partition}, the announcement and sifting maps, the
  raw-key map and the \exttermref{../protocol-compiler/contract.pdf}{term-index-map}{index map} of $\Decl$, and in the general format its
  factor list and Alice's \exttermref{../foundations/contract.pdf}{term-block-list}{block lists}; it differs from
  $\Decl$ at most in the measured families of the factors of
  $\mathbb{F}^{\mathrm{ref}}$ and in the entries of their last links. In
  particular the two have the same \exttermref{../protocol-compiler/contract.pdf}{term-joint-sector}{joint sectors}
  (Definition~\ref{def:joint-sectors}) with the same block dimensions and the
  same coupling condition;
\item \emph{(admissible)} $\Decl^{\mathrm{ref}}$ satisfies conditions
  (i)--(vi) and (viii) of Definition~\ref{def:admissible-declaration};
\item \emph{(the relaxed entries enlarge the \exttermref{../protocol-compiler/contract.pdf}{term-feasible-set}{feasible set})} let
  $\Decl^{\mathrm{ref}}_{\mathrm{ex}}$ be formed as $\Decl^{\mathrm{ref}}$
  but with the companions' entries themselves in place of their relaxed
  entries; then the feasible set of $\Decl^{\mathrm{ref}}_{\mathrm{ex}}$
  (Definition~\ref{def:feasible}) is contained in that of
  $\Decl^{\mathrm{ref}}$, and for every acceptance set
  $\hlow[\Decl^{\mathrm{ref}}]\le\hlow[\Decl^{\mathrm{ref}}_{\mathrm{ex}}]$.
\end{enumerate}
\end{proposition}
\claimmeta{\metalabel{Uses}\mref{Def.}{def:declaration}, \mref{Def.}{def:admissible-declaration}, \mref{Def.}{def:homog}, \mref{Def.}{def:joint-sectors}, \mref{Def.}{def:feasible}, \extref{../protocol-compiler/contract.pdf}{protocol-compiler/def:compile}, \extref{../protocol-compiler/contract.pdf}{protocol-compiler/def:weighted}, \extref{../protocol-compiler/contract.pdf}{protocol-compiler/lem:weighted}, \extref{../reductions/contract.pdf}{reductions/cor:reference-companion}, \extref{../reductions/contract.pdf}{reductions/lem:row-relaxation}
\metasep \metalabel{Used by}\mref{Rem.}{rem:declared-measurement}, \mref{Cor.}{cor:reference-companion-factors}, \mref{Req.}{req:companion-admission}, \mref{Cor.}{cor:fo-reference}, \mref{Lem.}{lem:certified-bracket-enclosed}, \mref{Cor.}{cor:certified-keylength-reference}, \mref{Req.}{req:bracket-on-reference}
\metasep newly derived
\metasep accepted}

\begin{proof}
(i) The compilation map passes through the \exttermref{../foundations/contract.pdf}{term-declared-data}{declared data} it reads and
constructs only the raw-key map, the index map and the weighted families
(\xref{protocol-compiler/def:compile}), each measured factor's family from
the last link of that factor's receiver segment, weighted, where it is, by
that factor's own weighting and applied to nothing else. The raw-key map is
formed from the public maps and the alphabets; the index map from the
alphabets, the \exttermref{../reductions/contract.pdf}{term-sector-set}{sector set} and the \exttermref{../foundations/contract.pdf}{term-declared-order}{declared
orders} that arrive with them; the weighted families from the delivered
families and the declared weights $P$. By
\xref{reductions/cor:reference-companion}\,(i), for each
$f\in\mathbb{F}^{\mathrm{ref}}$ the target of the reference companion of
$\mathcal{J}_f$ has the alphabet, with its order, and the \exttermref{../foundations/contract.pdf}{term-sector-map}{sector map}
of the target of $f$'s last link, so every constructed object other than the
weighted families of the factors of $\mathbb{F}^{\mathrm{ref}}$ is formed
from the same data as for $\Decl$ and is the same, and so is the family of
every other measured factor. The constraint list is the concatenation of the
specification's entries, the \exttermref{../transmitter/contract.pdf}{term-source-segment}{source segment}'s and the receiver
segments', each with its datum and tag; only the entries of the replaced
links differ. The joint sectors are the pairs of an Alice sector and a sector
of the measured system, with the dimensions of both, and the coupling
condition reads Alice's slot and these dimensions; all of them are shared.

(ii) Conditions (i), (ii), (iv), (v) and (viii) read only components that
(i) shows to be those of $\Decl$, or, for (i), the shape of the measured
families, finite families of finite matrices over the shared sectors;
$\Decl$ satisfies them. Condition (iii): $\{M_x\}$ is $\Decl$'s. The family
of each $f\in\mathbb{F}^{\mathrm{ref}}$ is the target of the reference
companion of $\mathcal{J}_f$, a POVM block diagonal for the sector map by
\xref{reductions/cor:reference-companion}\,(i), or, where that factor's
family is weighted after its chain by that factor's weighting $P_f$
(\xref{protocol-compiler/def:compile}), the family
$P_f(\varphi\,|\,t(v))\,F_v$ over $(v,\varphi)\in\Vset_{P_f}$, a POVM
by \xref{protocol-compiler/lem:weighted}\,(i), each element a non-negative
multiple of a block-diagonal operator and so block diagonal. In the basic
format that family is Bob's. In the general format every factor's family
outside $\mathbb{F}^{\mathrm{ref}}$ is $\Decl$'s, and the product family
\eqref{eq:productfamily} is then a POVM on $\Qsys$, a tensor product of
positive semidefinite operators being positive semidefinite and
$\sum_{v}\bigotimes_{f'}F^{f'}_{v_{f'}}=\bigotimes_{f'}\sum_{v_{f'}}F^{f'}_{v_{f'}}
=\idon{\Qsys}$. Every entry operator is Hermitian and block diagonal for its
scope: the shared entries as in $\Decl$, and every relaxed entry by
hypothesis. Condition (vi): a relaxed entry has a relation in
$\{\preceq,\succeq\}$, or is the entry itself where its bound is zero, and a
real constant, and it carries the scope of the entry it relaxes and, where
that entry is party-local, its party-local tag. The compilation map tags it
as it tags every entry (\xref{protocol-compiler/def:compile}): an entry
without the party-local tag by the type of its datum, a real constant,
which is the homogeneous form of Definition~\ref{def:homog}, and a
party-local entry by the tag field it arrives with beside that constant.

(iii) The two tuples differ only in the entries of the replaced links, each
exact entry $(\hat O_k,\bowtie_k,c_k)$ of $\Decl^{\mathrm{ref}}_{\mathrm{ex}}$
standing against its relaxed entries in $\Decl^{\mathrm{ref}}$. Let
$\omega$ be feasible for $\Decl^{\mathrm{ref}}_{\mathrm{ex}}$, with $E_k$ and
$\delta_k$ as in \xref{reductions/lem:row-relaxation}, whose part (i) gives
$-\delta_k\idon{}\preceq E_k\preceq\delta_k\idon{}$ and hence the same for
$\idon{}\otimes E_k$ on any larger system. In
\exttermref{../transmitter/contract.pdf}{term-trusted-source-mode}{trusted-source mode} the entry reads
$\tr_{\Qsys}[(\idon{A}\otimes\hat O_k)\omega]\bowtie_kc_k\tau_A$ with
$\tr_{\Qsys}\omega=\tau_A$, and part (iv) of that lemma with $X=A$ gives
the relaxed entries at their constants. In \exttermref{../transmitter/contract.pdf}{term-untrusted-source-mode}{untrusted-source mode} it
reads $\tr[(\idon{A}\otimes\hat O_k)\omega]\bowtie_kc_k$ with
$\tr\omega=1$, and
$\tr[(\idon{A}\otimes E_k)\omega]\in[-\delta_k,\delta_k]$ gives them. A
party-local entry reads \eqref{eq:partylocal-entry}, the case of part (iv)
with $X$ the other systems, and an untagged entry of a factor is read for
its \exttermref{../protocol-compiler/contract.pdf}{term-lifted-operator}{lifted operator} $\idon{}\otimes\hat O_k$, whose rounding error is
$\idon{}\otimes E_k$. An equality entry gives both relaxed entries. So
$\omega$ satisfies every relaxed entry, of whichever replaced link, and is
feasible for $\Decl^{\mathrm{ref}}$, the other conditions of
Definition~\ref{def:feasible} being the same for both. The objective and the
acceptance set are the same for both tuples, so the infimum
\eqref{eq:hlow-def} over the larger set is at most that over the smaller.
\end{proof}

% claim meat-core/cor:reference-companion-factors
\claimsummary{meat-core/cor:reference-companion-factors}{Where a handed reference carries factor enclosures and the reduction forms its datum as a non-negative combination of blocks, each block of the compiled companion's measured families read with a factor enclosure is a member of it, which the certified bracket over the companion assumes.}
\claimbookmark{meat-core/cor:reference-companion-factors}{Corollary}{cor:reference-companion-factors}{Factor enclosures of the compiled reference companion}
\claimleadin{%
Lemma~\ref{lem:certified-bracket-enclosed} reads some blocks of a compiled
reference companion through \exttermref{../foundations/contract.pdf}{term-factor-enclosure}{factor enclosures} and assumes, as its (H-FAC),
that the companion's own block is a member of each. The factor enclosures
are formed by the module that constructs the reference, which establishes
that its blocks are members of them, as it does for its
\exttermref{../foundations/contract.pdf}{term-entrywise-enclosure}{entrywise enclosures}
(\xref{foundations/req:snap-reference}\,(i)); the reduction and the compiler
then carry them, each step keeping the block a member. This corollary
composes those two steps for the compiled companion of
Proposition~\ref{prop:reference-companion}, one replaced factor at a time,
each replaced factor's family being formed from its own link alone.
}
\begin{corollary}[Factor enclosures of the compiled reference companion]\label{cor:reference-companion-factors}
Let $\Decl$, $\mathbb{F}^{\mathrm{ref}}$, the instances $\mathcal{J}_f$ and
$\Decl^{\mathrm{ref}}$ be as in Proposition~\ref{prop:reference-companion},
and for $f\in\mathbb{F}^{\mathrm{ref}}$ let $G_f$ be the family of
\xref{reductions/cor:reference-companion} from which the datum
$A^{\mathrm{ref}}_f$ of the reference companion of $\mathcal{J}_f$ is
formed. Let $f\in\mathbb{F}^{\mathrm{ref}}$, and read the block on a
\exttermref{../reductions/contract.pdf}{term-data-sector}{data sector} $s$ of $\mathcal{J}_f$'s target of an element of the
measured family of factor $f$ of $\Decl^{\mathrm{ref}}$ formed from the
target element $F_v[A^{\mathrm{ref}}_f]$ with the triple that part (i) of
\xref{reductions/cor:companion-factor-enclosure} forms for the element
$A^{\mathrm{ref},(s)}_{f,v}$ of $A^{\mathrm{ref}}_f$ on $s$, where the rule
of (I1) of $\mathcal{J}_f$ is a non-negative combination (that corollary)
and that part forms one from \exttermref{../foundations/contract.pdf}{term-factor-enclosure}{factor enclosures}
(\xref{foundations/def:factor-enclosure}) that blocks of $G_f$ carry, its
core multiplied by $P_f(\varphi\,|\,t(v))$ where the family is weighted after
the chain by factor $f$'s weighting $P_f$ and the element is
$P_f(\varphi\,|\,t(v))\,F_v[A^{\mathrm{ref}}_f]$ with
$(v,\varphi)\in\Vset_{P_f}$ (\xref{protocol-compiler/def:weighted}, with
$P\coloneqq P_f$); and read every other block of the
measured families of $\Decl^{\mathrm{ref}}$ with none. Then every block so
read with a triple is a member of it, the triple being a factor enclosure of
that block.
\end{corollary}
\claimmeta{\metalabel{Uses}\mref{Prop.}{prop:reference-companion}, \extref{../foundations/contract.pdf}{foundations/def:factor-enclosure}, \extref{../protocol-compiler/contract.pdf}{protocol-compiler/def:compile}, \extref{../protocol-compiler/contract.pdf}{protocol-compiler/def:weighted}, \extref{../protocol-compiler/contract.pdf}{protocol-compiler/lem:weighted}, \extref{../protocol-compiler/contract.pdf}{protocol-compiler/lem:weighted-factor}, \extref{../reductions/contract.pdf}{reductions/cor:companion-factor-enclosure}, \extref{../reductions/contract.pdf}{reductions/cor:reference-companion}, \extref{../reductions/contract.pdf}{reductions/def:reduction-instance}
\metasep \metalabel{Used by}\mref{Req.}{req:companion-admission}, \mref{Cor.}{cor:certified-keylength-reference}
\metasep newly derived
\metasep accepted}

\begin{proof}
By the definition of $\Decl^{\mathrm{ref}}$ in
Proposition~\ref{prop:reference-companion}, the compilation map
(\xref{protocol-compiler/def:compile}) forms the measured family of each
factor $f\in\mathbb{F}^{\mathrm{ref}}$ from the companion's target
$\{F_v[A^{\mathrm{ref}}_f]\}$ as it forms that of $\Decl$ from the snapped
target of $f$'s last link: the target itself, or, where the family is
weighted after the chain, the family
$\{P_f(\varphi\,|\,t(v))\,F_v[A^{\mathrm{ref}}_f]\}_{(v,\varphi)\in\Vset_{P_f}}$ of
\xref{protocol-compiler/def:weighted}, over the sector map of the target in
either case; the families of other factors take no part in it. On a data
sector $s$ the block of $F_v[A^{\mathrm{ref}}_f]$ is
$A^{\mathrm{ref},(s)}_{f,v}$, by the target formula of
\xref{reductions/def:reduction-instance}, and parts (i) and (ii) of
\xref{reductions/cor:companion-factor-enclosure}, applied to $\mathcal{J}_f$,
$G_f$ and $A^{\mathrm{ref}}_f$, make the triple formed there a factor
enclosure of it. In the weighted case the block of
$P_f(\varphi\,|\,t(v))\,F_v[A^{\mathrm{ref}}_f]$ on $s$ is
$P_f(\varphi\,|\,t(v))\,A^{\mathrm{ref},(s)}_{f,v}$
(\xref{protocol-compiler/lem:weighted}\,(ii)), and
\xref{protocol-compiler/lem:weighted-factor} makes the same box with the core
multiplied by the weight a factor enclosure of it. A factor enclosure of a
matrix has that matrix among its members
(\xref{foundations/def:factor-enclosure}), so each block read with a triple
is a member of it.
\end{proof}

\section{Assumptions}\label{group-meat-core/ass:keyreg}
\sectionintro{The assumptions fix what the post-processing must satisfy for the results of this contract to hold, and what a run must attest about its laboratories, its seeds, its fixed parameters and its authentication, and, for several executions whose parameters are to add, their separation. What the declaration and the acceptance set must satisfy is decided by the engine and is stated as definitions (\hyperlink{claim-meat-core/def:admissible-declaration}{Definition~\ref*{def:admissible-declaration}}, \hyperlink{claim-meat-core/def:admissible-acceptance-set}{Definition~\ref*{def:admissible-acceptance-set}}); the product marginal of Alice's systems in trusted-source mode is a lemma under the transmitter's premises and this module's A-SEED (v) and A-LAB at the party holding A (\hyperlink{claim-meat-core/lem:product-marginal}{Lemma~\ref*{lem:product-marginal}}). No assumption restricts the adversary's memory: \hyperlink{claim-meat-core/ass:eve-dim}{A-EVE-DIM}, which would make it finite-dimensional, is assumed by no claim (\hyperlink{claim-meat-core/cor:composed-security}{Corollary~\ref*{cor:composed-security}}). The attack-set licence and round homogeneity are the compiler's assumptions, cited where they are used. Each assumption states its conditions only, with a plain reading for the user.}
\sectionimports{\extref{../foundations/contract.pdf}{foundations/def:round-model}}

% claim meat-core/ass:keyreg
\claimsummary{meat-core/ass:keyreg}{The post-processing acts on the raw-key string the rounds write, the declared holder computes that string, and reconciliation runs one way from the holder through a padded fixed-length transcript. The leakage term counts the register's full length.}
\claimbookmark{meat-core/ass:keyreg}{Assumption}{ass:keyreg}{Dedicated key register and one-way reconciliation}
\claimleadin{%
The chain-rule penalty counts the dimension of the transcript register, not
the average length of a transcript, which is why the register has a fixed
length. Whether the raw-key map is one the holder can compute from its own
readout is the party-computability condition of
Definition~\ref{def:declaration}, decided on the declared tables; what is
assumed here is that the run's post-processing acts on that string, one way.
}
\SetShortId{A-KEYREG}
\begin{assumption}[Dedicated key register and one-way reconciliation]\label{ass:keyreg}
(i) Error correction, error verification and privacy amplification are
applied to $Z_1^n$, the string of raw-key values $r(x,v)$ that the rounds
write to the registers $Z_1,\dots,Z_n$
(Definition~\ref{def:postprocessing-model});
(ii) the declared key holder $\hold$ (Definition~\ref{def:declaration})
computes $Z_1^n$ from its own outcome record and the public announcements;
(iii) reconciliation is one-way from $\hold$: its transcript
$L_{\mathrm{EC}}$ is a function of $Z_1^n$ and public seeds only;
(iv) $L_{\mathrm{EC}}$ occupies a fixed-length register of $\lEC$ bits,
shorter transcripts being padded, so that $\dim(L_{\mathrm{EC}})=2^{\lEC}$.
\end{assumption}
\claimverification{The engine does not verify this assumption; a run that rests on it records it as unverified.}
\claimreading{Error correction is one-way. The key holder computes its raw key from its own record (its letters, the detection results of its measuring stations, or both, whichever it holds) and the public announcements, and its messages depend only on that raw key and public random seeds. The transcript always has exactly \ensuremath{\lambda}\_EC bits, shorter messages being padded. Interactive error correction (for example Cascade) or a variable-length transcript is not covered by this key length.}
\claimmeta{\metalabel{Uses}\mref{Def.}{def:declaration}, \mref{Def.}{def:postprocessing-model}
\metasep \metalabel{Category}post processing
\metasep \metalabel{Affects}secrecy
\metasep \metalabel{Assumed by}\mref{Thm.}{thm:keylength}
\metasep \metalabel{Machine key}\iface{key_register}
\metasep newly derived
\metasep accepted}

% claim meat-core/ass:hash
\claimsummary{meat-core/ass:hash}{Privacy amplification uses an ideal two-universal family, collision probability exactly 2\textasciicircum{}\ensuremath{-}\ensuremath{\ell} for distinct inputs, and error verification a two-universal family. The first is what leftover hashing reads, the second what the verification collision bound reads.}
\claimbookmark{meat-core/ass:hash}{Assumption}{ass:hash}{Ideal two-universal privacy amplification}
\claimleadin{%
The R\'enyi leftover-hashing lemma \cite[Thm.~4]{kamin2025renyi} is applied
under the reading of \cite[Lemma~6.5 and Rem.~6.1]{tupkary2026decoy}, which
asks for an ideal family; the correctness argument uses the collision
probability of the verification family only as an upper bound. The two
notions are those of \cite[Def.~3.5]{tupkary2026decoy}.
}
\SetShortId{A-HASH}
\begin{assumption}[Ideal two-universal privacy amplification]\label{ass:hash}
(i) The privacy-amplification family of
Definition~\ref{def:postprocessing-model}, applied to $Z_1^n$, is ideal
two-universal: for any two distinct inputs, the probability under the
family's law that their $\ell$-bit values coincide is exactly $2^{-\ell}$;
(ii) the error-verification family of
Definition~\ref{def:postprocessing-model} is two-universal: for any two
distinct inputs, the probability under the family's law that their
$\lceil\log_2\frac1{\eev}\rceil$-bit values coincide is at most
$2^{-\lceil\log_2(1/\eev)\rceil}$.
\end{assumption}
\claimverification{The engine does not verify this assumption; a run that rests on it records it as unverified.}
\claimreading{Privacy amplification uses a hash family in which any two different inputs collide with probability exactly 2\textasciicircum{}\ensuremath{-}\ensuremath{\ell} (an ideal two-universal family, such as Toeplitz hashing), and error verification uses a two-universal hash family. If your privacy-amplification family is only two-universal (collision probability at most 2\textasciicircum{}\ensuremath{-}\ensuremath{\ell}), shorten the key by one further bit.}
\claimmeta{\metalabel{Uses}\mref{Def.}{def:postprocessing-model}
\metasep \metalabel{Category}post processing
\metasep \metalabel{Affects}secrecy, correctness
\metasep \metalabel{Assumed by}\mref{Lem.}{lem:ev-correctness}, \mref{Thm.}{thm:keylength}
\metasep \metalabel{Machine key}\iface{pa_hash_family}
\metasep newly derived
\metasep accepted}

% claim meat-core/ass:seed
\claimsummary{meat-core/ass:seed}{The post-processing seeds are drawn with their declared laws, privately, independently and afresh; in trusted-source mode the source's letters and prepared systems are independent of the adversary, who never gets a purification of them. Leftover hashing and source replacement rest on it.}
\claimbookmark{meat-core/ass:seed}{Assumption}{ass:seed}{Private, independent and fresh local randomness}
\claimleadin{%
The parties' local randomness is taken to be ideal and never leaked beyond
what the protocol announces \cite[Sec.~4, opening paragraphs]{tupkary2026decoy},
\cite[Sec.~3]{tupkary2025survey}; a generator that is only close to ideal
in the composable sense adds its own error parameter
\cite[Sec.~5.6.1]{tupkary2025survey}. The published premise covers the
randomness used to choose the signal states as well as the hash seeds: the
protocol begins with local randomness in classical registers and ``the
adversary is uncorrelated to this randomness''
\cite[Sec.~4, opening paragraphs]{tupkary2026decoy}, and in particular the
adversary holds no purification of a prepared signal state
\cite[Protocol~1, step 1(a)(i)]{tupkary2026decoy}. Clause (v) states that
part for the source's draws, in \exttermref{../transmitter/contract.pdf}{term-trusted-source-mode}{trusted-source mode}, where the analysis
describes the source by them.
}
\SetShortId{A-SEED}
\begin{assumption}[Private, independent and fresh local randomness]\label{ass:seed}
(i) The error-correction seed $S_{\mathrm{pub}}$, the error-verification seed
$S_{\mathrm{EV}}$ and the privacy-amplification hash $G$ of
Definition~\ref{def:postprocessing-model} are drawn with the laws that
definition gives them, $S_{\mathrm{EV}}$ and $G$ with the laws of their hash
families;
(ii) each is drawn privately: its value is available to the adversary only
once the protocol announces it;
(iii) each is independent of every other register of the protocol and of the
adversary's information at the time it is drawn;
(iv) each is drawn afresh in every run;
(v) in \exttermref{../transmitter/contract.pdf}{term-trusted-source-mode}{trusted-source mode}, the letters Alice's source draws in the
$n$ rounds, and the systems it prepares from them, are jointly independent of
the adversary's registers at the start of the protocol: the joint state of
the letters, the prepared systems and those registers is the product of its
marginal on the letters and the prepared systems together and its marginal on
those registers; in particular the adversary holds no purification of a
prepared system, and no purification of a prepared system is made available
to the adversary during or after the protocol.
\end{assumption}
\claimverification{The engine does not verify this assumption; a run that rests on it records it as unverified.}
\claimreading{Your random choices come from a good private random number generator. The error-correction seed, the error-verification hash and the privacy-amplification hash are each drawn with their intended distribution, kept secret until the protocol announces them, independent of each other, of every other record of the protocol (the raw key included) and of everything the attacker knows, and drawn anew in every run. With a trusted source, the randomness with which it picks each round's letter and prepares its state is uncorrelated with the attacker at the start, and no purification of a prepared state is ever available to the attacker. A seed reused across runs, fixed in the device or disclosed early, or a predictable or pre-shared random source, violates this.}
\claimmeta{\metalabel{Uses}\mref{Def.}{def:postprocessing-model}
\metasep \metalabel{Category}local randomness
\metasep \metalabel{Affects}secrecy, correctness
\metasep \metalabel{Assumed by}\mref{Lem.}{lem:product-marginal}, \mref{Thm.}{thm:accumulated-entropy}, \mref{Lem.}{lem:ev-correctness}, \mref{Thm.}{thm:keylength}
\metasep \metalabel{Machine key}\iface{post_processing_seeds}
\metasep newly derived
\metasep accepted}

% claim meat-core/ass:fixed
\claimsummary{meat-core/ass:fixed}{The acceptance set, the round count, the key length, the transcript length, the error budgets, the authentication parameters and the R\'enyi order are fixed before the run and applied unchanged. The key length is that of a fixed-length protocol, not one chosen from the run's data.}
\claimbookmark{meat-core/ass:fixed}{Assumption}{ass:fixed}{Protocol parameters fixed before the run}
\claimleadin{%
A fixed-length protocol outputs a key of a length determined before the
protocol begins, or aborts \cite[Sec.~2.2]{tupkary2025survey}; choosing the
length or the test from the observed data makes it a variable-length
protocol, which needs a different analysis. The R\'enyi order $\alpha$ is a
parameter of the proof, not of the protocol: the run never reads it, and it
enters only the computation of $\ell$. Clause (iii) fixes it before the run
all the same, so that the order, and with it the length computed at it, is
never a function of the run's data; an order chosen by comparing the
lengths that the \exttermref{../reductions/contract.pdf}{term-declared-model}{declared model} gives at several orders satisfies it.
No proof consumes clause (iii): the key-length theorem holds at each
$\alpha\in(1,2)$, and the clause is stated so that the attestation names
every input of $\ell$.
}
\SetShortId{A-FIXED}
\begin{assumption}[Protocol parameters fixed before the run]\label{ass:fixed}
(i) The acceptance set $\acSet$, the round count $n$, the output length
$\ell$, the transcript length $\lEC$, the error budgets $\epa$ and $\eev$ of
Definition~\ref{def:postprocessing-model}, and the authentication parameters
$\varepsilon_{\mathrm{auth}}$ and $k_{\mathrm{auth}}$, the declared
authentication error bound and key consumption
(Definition~\ref{def:authenticated-channel}\,(b)), are fixed before the
first round of the run;
(ii) the run applies them unchanged: none of them is chosen, adjusted or
recomputed from data observed in the run;
(iii) the R\'enyi order $\alpha\in(1,2)$ at which $\ell$ is computed is
fixed before the first round of the run, from data available before it, and
is not chosen, adjusted or recomputed from data observed in the run.
\end{assumption}
\claimverification{The engine does not verify this assumption; a run that rests on it records it as unverified.}
\claimreading{Every parameter the key length rests on is fixed before the run starts: the acceptance test, the number of rounds, the key length, the error-correction budget \ensuremath{\lambda}\_EC, the error budgets \ensuremath{\varepsilon}\_PA and \ensuremath{\varepsilon}\_EV, and the authentication parameters \ensuremath{\varepsilon}\_auth and k\_auth. The R\'enyi order \ensuremath{\alpha} at which the key length is computed is fixed before the run too; choosing it by comparing the key lengths your declared model gives at several orders, as the run tier's outer search over \ensuremath{\alpha} does, reads no data of the run and is such a choice. None of them may be chosen or adjusted after looking at the run's data, for example by relaxing the acceptance threshold or picking the key length from the observed error rate. That would be a variable-length protocol, which this analysis does not cover.}
\claimmeta{\metalabel{Uses}\mref{Def.}{def:postprocessing-model}, \mref{Def.}{def:authenticated-channel}
\metasep \metalabel{Category}protocol structure
\metasep \metalabel{Affects}secrecy, correctness
\metasep \metalabel{Assumed by}\mref{Thm.}{thm:keylength}, \mref{Cor.}{cor:composed-security}
\metasep \metalabel{Machine key}\iface{parameters_fixed_before_run}
\metasep newly derived
\metasep accepted}

% claim meat-core/ass:conform
\claimsummary{meat-core/ass:conform}{The implementation applies the acceptance test and the verification abort, sends no message beyond the model's, outputs the model's key or nothing, uses no key before the execution ends and aborts on an authentication failure. It ties the analysed protocol to the one that runs.}
\claimbookmark{meat-core/ass:conform}{Assumption}{ass:conform}{The implementation runs the analysed post-processing}
\claimleadin{%
That the devices execute the protocol as prescribed is one of the standard
premises of QKD \cite[Sec.~IV.A]{portmann2022security}; the post-processing
prescribed here is the fixed-length instance of
\cite[Sec.~4.1, Protocol~1]{tupkary2026decoy}.
}
\SetShortId{A-CONFORM}
\begin{assumption}[The implementation runs the analysed post-processing]\label{ass:conform}
(i) The implementation evaluates the acceptance test of
Definition~\ref{def:postprocessing-model}\,(i) and aborts unless $\OmAT$
holds;
(ii) after the rounds, the only classical messages either party sends are
those of Definition~\ref{def:postprocessing-model} and of the
\termref{term-authentication-post-processing}{authentication post-processing}
(Definition~\ref{def:authenticated-channel}\,(c));
(iii) both parties abort when the verification flag is $\mathrm{rej}$;
(iv) each party outputs the $\ell$-bit key that
Definition~\ref{def:postprocessing-model}\,(iv) assigns to it, or $\perp$;
(v) neither party uses any part of its key before the
\termref{term-end-of-an-execution}{end of an execution} (Definition~\ref{def:execution-end});
(vi) a party that receives an authentication failure acts as the
authentication post-processing prescribes.
\end{assumption}
\claimverification{The engine does not verify this assumption; a run that rests on it records it as unverified.}
\claimreading{Your implementation runs exactly the analysed post-processing. It aborts unless the acceptance statistics lie in the declared acceptance set; after the rounds it sends no public message other than the error-correction transcript of at most \ensuremath{\lambda}\_EC bits and its seed, the verification seed, hash and result, the privacy-amplification hash choice and the closing accept or abort decisions; it aborts when verification fails or a message fails authentication (an auth-abort); it outputs the analysed \ensuremath{\ell}-bit key, or nothing. No part of the key is used before the execution has ended: the key holder, having sent its final decision, waits until the partner has received it or an auth-abort in its place, and the partner until one of these has arrived. Interactive error correction, a resent syndrome or early key use violate this. Tightening 2\ensuremath{\varepsilon}\_auth to \ensuremath{\varepsilon}\_auth under this timing is an open question.}
\claimmeta{\metalabel{Uses}\mref{Def.}{def:postprocessing-model}, \mref{Def.}{def:authenticated-channel}, \mref{Def.}{def:execution-end}
\metasep \metalabel{Category}post processing
\metasep \metalabel{Affects}secrecy, correctness
\metasep \metalabel{Assumed by}\mref{Lem.}{lem:ev-correctness}, \mref{Lem.}{lem:honest-attacks}, \mref{Thm.}{thm:keylength}, \mref{Cor.}{cor:composed-security}
\metasep \metalabel{Machine key}\iface{post_processing_as_analysed}
\metasep newly derived
\metasep accepted}

% claim meat-core/ass:auth
\claimsummary{meat-core/ass:auth}{Every classical message travels through an authentication scheme using at most k\_auth bits of pre-shared key and realising the authenticated classical channel within \ensuremath{\varepsilon}\_auth, a reused key's distance from an ideal one included. It is the premise by which authentication enters the key's security.}
\claimbookmark{meat-core/ass:auth}{Assumption}{ass:auth}{Authentication of the classical channel by a keyed scheme}
\claimleadin{%
A real scheme authenticates with a short pre-shared key and an
unauthenticated channel, for instance by Wegman--Carter hashing; its
composable security is the distance of
\cite[Sec.~VII.A, Eq.~(36)]{portmann2022security} from the authenticated
channel, and an authentic classical channel is one of the standard premises
of QKD \cite[Sec.~IV.A]{portmann2022security},
\cite[Sec.~3]{tupkary2025survey}.
}
\SetShortId{A-AUTH}
\begin{assumption}[Authentication of the classical channel by a keyed scheme]\label{ass:auth}
(i) All classical messages of the protocol, the round announcements
included, are sent through one authentication scheme using a pre-shared key
and an unauthenticated classical channel, or all over a classical channel
that is itself the \termref{term-authenticated-classical-channel}{authenticated classical channel} of
Definition~\ref{def:authenticated-channel}\,(a), in which case
$\varepsilon_{\mathrm{auth}}=k_{\mathrm{auth}}=0$;
(ii) in the first case, the scheme's key consumption is at most
$k_{\mathrm{auth}}$ bits;
(iii) run on two identical copies of a uniformly distributed key independent
of every other register of the protocol and of the adversary's information,
the scheme $\varepsilon_0$-realises the authenticated classical channel
(Definition~\ref{def:authenticated-channel}\,(b)), for some
$\varepsilon_0\ge0$;
(iv) before the run starts, the joint state of the two parties' copies of
the pre-shared key the run consumes with every other register of the
protocol and with the adversary's information is within trace distance
$\delta\ge0$ of the state in which those copies are replaced by two
identical copies of a uniformly distributed key independent of them;
(v) $\varepsilon_0+\delta\le\varepsilon_{\mathrm{auth}}$.
\end{assumption}
\claimverification{The engine does not verify this assumption; a run that rests on it records it as unverified.}
\claimreading{Every classical message of your protocol, per-round announcements included, is protected by one authentication scheme: the attacker may read, delay or block any message, but one she alters or injects, or delivers before it was sent, is rejected except within an error \ensuremath{\varepsilon}\_auth. It uses at most k\_auth bits of pre-shared secret key per run. That error is its composable error on two identical copies of a uniform, independent key plus the distance of the parties' actual copies from such a pair (zero for a fresh key; for a key from an earlier run, a distance you must bound, which that run's security parameter alone does not bound); declare that sum as \ensuremath{\varepsilon}\_auth and the bits used as k\_auth. No composable proof for this timing-aware channel is published yet; Wegman\textendash{}Carter authentication with separate keys per direction approximates it. Both declared zero assert a keyless authenticated channel.}
\claimmeta{\metalabel{Uses}\mref{Def.}{def:authenticated-channel}
\metasep \metalabel{Category}channel and adversary
\metasep \metalabel{Affects}secrecy, correctness
\metasep \metalabel{Assumed by}\mref{Cor.}{cor:composed-security}
\metasep \metalabel{Machine key}\iface{classical_channel_authenticated}
\metasep newly derived
\metasep accepted}

% claim meat-core/ass:separate-executions
\claimsummary{meat-core/ass:separate-executions}{Several executions, fixed in advance and run against one adversary, touch none of each other's registers, leave their keys unused until all have ended, and keep their own premises with the others counted as the adversary's. It is what lets their security parameters add.}
\claimbookmark{meat-core/ass:separate-executions}{Assumption}{ass:separate-executions}{Executions that share no key and no device state}
\claimleadin{%
Composable frameworks add the errors of protocols composed in sequence or
in parallel \cite[Sec.~II.E]{portmann2022security}. For the trace-distance
parameter of this contract, the corresponding sum needs the executions to
be separate in the sense below, so that each can be analysed with every
other one on the adversary's side.
}
\SetShortId{A-SEPARATE}
\begin{assumption}[Executions that share no key and no device state]\label{ass:separate-executions}
For executions $1,\dots,N$ of the protocol, run against one adversary:
(i) $N$ and, for every execution, the parameters fixed before its run (its
acceptance set, round count, output length, transcript length, error
budgets and authentication parameters) are fixed before the first of the
$N$ executions starts;
(ii) no operation of the honest parties of one execution acts on a register
that the honest parties of another execution use, hold or output: no
execution uses any part of another execution's key, local randomness,
seeds or pre-shared key, no device, memory or random source carries a
state from one execution into another, and the executions interact only
through systems and classical channels the adversary has access to;
(iii) from the \termref{term-end-of-an-execution}{end of an execution}
(Definition~\ref{def:execution-end}) until the end of the last of the $N$
executions, no part of that execution's key is used, read or otherwise
acted on, and no information about any other register of its parties
leaves their laboratories;
(iv) for every execution, every assumption on which its key and its
security parameter rest holds when every register of the honest parties of
every other execution is counted as part of the adversary's information,
and every operation of those parties as part of the adversary's actions.
\end{assumption}
\claimverification{The engine does not verify this assumption; a run that rests on it records it as unverified.}
\claimreading{You run several executions, one after another or at the same time, and want their security parameters to add. Their number and parameters are fixed before the first starts. Each execution has its own key material, randomness, seeds and pre-shared authentication key, such as a separate part of one pool. No device, memory or random source carries information between them. An execution has ended once both parties have fixed their output, key or abort, after the closing messages; no key is used or read, and nothing else leaks from an ended one, until all have ended. The assumptions each execution's run record lists must still hold when an attacker holds and runs every other execution, her memory, which then holds their systems, of any dimension. A key that authenticates another execution, a key used early, or parameters chosen from an earlier execution's data violate this.}
\claimmeta{\metalabel{Uses}\mref{Def.}{def:execution-end}
\metasep \metalabel{Category}protocol structure
\metasep \metalabel{Affects}secrecy, correctness
\metasep \metalabel{Assumed by}\mref{Cor.}{cor:separate-executions}
\metasep \metalabel{Machine key}\iface{separate_executions}
\metasep newly derived
\metasep accepted}

% claim meat-core/ass:lab
\claimsummary{meat-core/ass:lab}{Nothing about a party's records leaves its laboratory but its key, the public messages and, in trusted-source mode at the party of A, the emitted systems; the adversary acts on it only through what its declared measurements receive and the messages. It excludes side channels at both parties.}
\claimbookmark{meat-core/ass:lab}{Assumption}{ass:lab}{Each party's laboratory is sealed}
\claimleadin{%
An information carrier of an implementation that the model omits is a side
channel that no security proof accounts for, and sources and photon detectors
are known places for one \cite[Sec.~I.B]{portmann2022security}; that the
devices do exactly what they are instructed to do is one of the standard
premises of QKD \cite[Sec.~IV.A]{portmann2022security}. The premise is stated
for each party of the declaration, whatever that party does in the protocol:
what crosses the boundary of its laboratory is read off the
\exttermref{../foundations/contract.pdf}{term-round-model}{round model} and the \exttermref{../protocol-compiler/contract.pdf}{term-party-partition}{party partition}, in both directions.
}
\SetShortId{A-LAB}
\begin{assumption}[Each party's laboratory is sealed]\label{ass:lab}
Let $\mathsf P\in\{\mathsf A,\mathsf B\}$ be a party of the
\exttermref{../protocol-compiler/contract.pdf}{term-protocol-declaration}{protocol declaration} (Definition~\ref{def:declaration}). Its
\emph{records} are its letters, where the \exttermref{../protocol-compiler/contract.pdf}{term-party-partition}{party partition} assigns
Alice's system $A$ to $\mathsf P$; the private settings and the outcomes of
the measurements made in its laboratory, the public choices it draws before
they are announced, and its readout $\rd{\mathsf P}$; the inputs of its
announcements; its raw or final key; and its local randomness. The \emph{systems handed out of its laboratory} are, in
\exttermref{../transmitter/contract.pdf}{term-trusted-source-mode}{trusted-source mode} and where the party partition assigns $A$ to
$\mathsf P$, the systems its source emits into the quantum channel, which the
\exttermref{../foundations/contract.pdf}{term-round-model}{round model} (Definition~\ref{def:roundmodel}) gives to the adversary,
$A$ being no output of the attack in that mode \eqref{eq:attackchannel} and
staying with $\mathsf P$; in every other case there are none. In particular,
in \exttermref{../transmitter/contract.pdf}{term-untrusted-source-mode}{untrusted-source mode}, where the attack supplies $A$ and $\Qsys$
\eqref{eq:attackchannel}, no laboratory hands out any system, and a source the declaration treats as untrusted lies
outside every laboratory wherever it physically stands. The \emph{systems
delivered into its laboratory} are, in untrusted-source mode and where the
party partition assigns $A$ to $\mathsf P$, the system $A$; and, in either
mode and in each sector $s$, the tensor factors $\Qsys_{s,i}$ of the declared
factorisation of $\Hs{\Qsys_s}$ that the party partition labels $\mathsf P$
(in the general format, every \exttermref{../protocol-compiler/contract.pdf}{term-measured-factor}{measured factor} labelled $\mathsf P$),
physically the input modes those factors model.
\begin{enumerate}[label=(\roman*),leftmargin=2.2em]
\item During the protocol, no information about $\mathsf P$'s records leaves
its laboratory, other than its key as the protocol's output, except through
the systems handed out of its laboratory and through the public messages of
the protocol (the announcements of the
\termref{term-round-channel}{round channels} of Definition~\ref{def:roundmodel}, the
messages of Definition~\ref{def:postprocessing-model} and of the
\termref{term-authentication-post-processing}{authentication post-processing}
(Definition~\ref{def:authenticated-channel}\,(c)), and, where an
authentication scheme carries them, the scheme's own transmissions over the
unauthenticated classical channel): no light emitted by its devices other
than those systems, no electromagnetic, acoustic, power or timing emanation
and no side channel of its classical processing carries such information, and
neither the timing of its messages nor any auxiliary or synchronisation
signal does;
\item the adversary acts on $\mathsf P$'s laboratory only through the systems
delivered into its laboratory, as the declared measurements receive them, and
through the public messages the protocol delivers to $\mathsf P$ and, where a
scheme carries them, the transmissions it receives over the unauthenticated
channel: no light or signal sent into the laboratory, through the systems
handed out of it or otherwise, reads out any of its records, and nothing else
influences the state, configuration or calibration of its devices (its source
included, where it has one), its measurements or its classical processing.
\end{enumerate}
\end{assumption}
\claimverification{The engine does not verify this assumption; a run that rests on it records it as unverified.}
\claimreading{Whatever a party's role, nothing about its letters, settings, choices, detection results, keys or random numbers leaves its laboratory but the final key, the public messages and, where the party sends with a trusted source, the signals that source emits; timing and auxiliary signals reveal no more. An untrusted source counts as outside every laboratory, even one it is in. The attacker reaches a laboratory only via the systems its declared measurements receive and the public messages. Light sent in to read out a letter, a setting, a choice or a detector state (a Trojan-horse attack, at a receiver as at a modulator), emanations from any device, detector back-flash and anything but the declared input modes influencing a measurement violate this, a clock or trigger line from an untrusted source included; where it sends with a trusted source, so does light seeding or recalibrating its laser.}
\claimmeta{\metalabel{Uses}\mref{Def.}{def:declaration}, \mref{Def.}{def:roundmodel}, \mref{Def.}{def:postprocessing-model}, \mref{Def.}{def:authenticated-channel}
\metasep \metalabel{Category}laboratory
\metasep \metalabel{Affects}secrecy
\metasep \metalabel{Assumed by}\mref{Lem.}{lem:product-marginal}, \mref{Thm.}{thm:accumulated-entropy}, \mref{Lem.}{lem:honest-attacks}, \mref{Thm.}{thm:keylength}
\metasep \metalabel{Machine key}\iface{laboratory_sealed}
\metasep newly derived
\metasep accepted}

% claim meat-core/ass:eve-dim
\claimsummary{meat-core/ass:eve-dim}{Every register of the adversary, her initial register and her memories, is finite-dimensional, with no bound on its dimension. No claim rests on it: the analysed models have such registers by definition, and \hyperlink{claim-meat-core/cor:composed-security}{Corollary~\ref*{cor:composed-security}} reaches an adversary whose memory has any dimension.}
\claimbookmark{meat-core/ass:eve-dim}{Assumption}{ass:eve-dim}{The adversary's registers are finite-dimensional}
\claimleadin{%
The source of the imported accumulation theorem states its convention as
``we will assume that all systems are finite-dimensional, but we will not
impose any bounds on the system dimensions unless otherwise specified''
\cite[Sec.~2, the convention stated below its Table~1]{arqand2025meat}. Read for the adversary's registers, that convention would
be a premise about her, and this assumption states it. The analysed models
of these contracts keep it instead as part of their definitions, at every
fixed finite dimension, and the security statement of a run reaches an
adversary whose memory has any dimension by a limit over those dimensions;
so no claim assumes it.
}
\SetShortId{A-EVE-DIM}
\begin{assumption}[The adversary's registers are finite-dimensional]\label{ass:eve-dim}
(i) Every register of the adversary in the \exttermref{../foundations/contract.pdf}{term-round-model}{round model}
(\xref{foundations/def:round-model}), her initial register $E_0$, her
outgoing memories $E'_j$ and her registers $E_j$ for $j=1,\dots,n$, is
finite-dimensional, of arbitrary finite dimension.
\end{assumption}
\claimverification{The engine does not verify this assumption; a run that rests on it records it as unverified.}
\claimreading{No run rests on this assumption: the security statement of a run holds against an attacker whose quantum memory has any size, finite or infinite (\hyperlink{claim-meat-core/cor:composed-security}{Corollary~\ref*{cor:composed-security}}). The assumption states that the attacker's quantum memory is finite-dimensional, of any size she likes.}
\claimmeta{\metalabel{Uses}\extref{../foundations/contract.pdf}{foundations/def:round-model}
\metasep \metalabel{Category}channel and adversary
\metasep \metalabel{Affects}secrecy
\metasep \metalabel{Machine key}\iface{adversary_registers_finite}
\metasep newly derived
\metasep superseded}

\clearpage
\part{Engine requirements}
\sectionintro{What an implementation of this module must do, record or refuse; each requirement names the result that motivates it.}

\section{Requirements: admission and the declared structure}\label{group-meat-core/req:exact-admission}
\sectionintro{The engine admits declared data only when they are exactly admissible, decides every condition of an \termref{term-admissible-declaration}{admissible declaration} but the non-empty feasible set itself, admits a reference companion only beside its declaration, deciding exactly what the two share and nothing of its family, and admits an acceptance set only when it decides that the set is admissible, and an implemented \termref{term-branch-map}{branch map} drops only components, of registers other than the raw-key register, that are constant across its Kraus terms. It decides the sector classes by exact tests, flags a sector removable only when an exhibited witness meets the whole removal checklist, and reports every class and every deletion. It changes the declared sectors in no other way than the removal theorem and the composition rule allow, and deletes at most one sector where it does not check the composition hypothesis.}
\sectionimports{\extref{../foundations/contract.pdf}{foundations/def:exact-reference}, \extref{../protocol-compiler/contract.pdf}{protocol-compiler/req:reference-companion-compiled}, \extref{../reductions/contract.pdf}{reductions/cor:reference-companion}, \extref{../reductions/contract.pdf}{reductions/req:reference-companion}}

% claim meat-core/req:exact-admission
\claimsummary{meat-core/req:exact-admission}{The engine admits a declaration only if it is exactly admissible and refuses every other, with no tolerance on declared data, save a reference companion, admitted on its own terms. \hyperlink{claim-meat-core/prop:optol-nobound}{Proposition~\ref*{prop:optol-nobound}} shows, for the measurement on Q, that closeness to admissible data could license no key-length claim.}
\claimbookmark{meat-core/req:exact-admission}{Requirement}{req:exact-admission}{Exact admission of declared data}
\claimleadin{%
Exact admissibility asks the conditions of
Definition~\ref{def:declaration} of the declared
\exttermref{../foundations/contract.pdf}{term-binary64-numbers}{binary64 numbers} themselves, not of reals they approximate. For a
declaration with one Alice sector, one \exttermref{../protocol-compiler/contract.pdf}{term-measured-factor}{measured factor} and Alice's
measurement in \exttermref{../protocol-compiler/contract.pdf}{term-povm-form}{POVM form}, this means that, block by block, each
declared operator is exactly Hermitian, each measurement operator and the
declared $\tau_A$ are exactly positive semidefinite, and each completeness
sum, formed in the \exttermref{../foundations/contract.pdf}{term-declared-order}{declared order} \eqref{eq:declorder}, is exactly
the identity with no addition rounding; the paragraph of that definition on
exact admissibility in the general format says how the conditions read
otherwise. How the rule meets an instance that moves its own data within
its \exttermref{../foundations/contract.pdf}{term-declared-tolerance}{declared tolerance} $\tnum$ is
\xref{foundations/req:declared-tolerance}, read with
\xref{foundations/req:consumer-decides-exactly}.
}
\begin{requirement}[Exact admission of declared data]\label{req:exact-admission}
The engine takes no tolerance on \exttermref{../foundations/contract.pdf}{term-declared-data}{declared data} and decides no
condition on it more loosely than exactly
(\xref{foundations/req:consumer-decides-exactly}): it admits a declaration
only if the declaration is \exttermref{../foundations/contract.pdf}{term-exactly-admissible}{exactly admissible} in the sense of
Definition~\ref{def:declaration}, and refuses it otherwise. No
perturbation of a declared family is licensed to carry the key length of
the family it approximates. Proposition~\ref{prop:optol-nobound} is the
reason: it shows, for the measurement on $\Qsys$, that no tolerance could
license a key-length claim. A compiled reference companion is not admitted
under this requirement but under Requirement~\ref{req:companion-admission},
which decides exactly what it shares with the declaration and decides none
of its family; the binary64 numbers it is evaluated at are no declared
family whose key length is claimed, since the bound evaluated there is that
of the reference, at the unbudgeted arithmetic of
Remark~\ref{rem:fo-arithmetic}.
\end{requirement}
\claimmeta{\metalabel{Motivated by}\mref{Def.}{def:declaration}, \mref{Prop.}{prop:optol-nobound}, \mref{Req.}{req:companion-admission}, \mref{Rem.}{rem:fo-arithmetic}, \extref{../foundations/contract.pdf}{foundations/req:consumer-decides-exactly}, \extref{../foundations/contract.pdf}{foundations/req:declared-tolerance}
\metasep implemented (effective: accepted)
\metasep \metalabel{Code}\coderef{https://github.com/fraunhofersg/certiqs-key/blob/v0.1.0/src/certiqs_key/meat_core/declaration/general_validation.py}{certiqs_key.meat_core.declaration.general_validation:validate_general_declaration}}

% claim meat-core/req:declaration-admission
\claimsummary{meat-core/req:declaration-admission}{The engine computes nothing from a declaration that fails a condition of admissibility other than the non-empty feasible set, each decided exactly on the declared data. The feasible set it leaves to the compiler where constraints are declared.}
\claimbookmark{meat-core/req:declaration-admission}{Requirement}{req:declaration-admission}{Admission of the declaration}
\claimleadin{%
Seven of the eight conditions of admissibility are conditions on the
declared numbers and tables, and the engine decides them itself rather than
trusting the compiler's decision
(Requirement~\ref{req:exact-admission} for the operator conditions). The
eighth, a non-empty \exttermref{../protocol-compiler/contract.pdf}{term-feasible-set}{feasible set}, holds of every declaration with an empty
constraint list, which has an explicit feasible state
(\xref{protocol-compiler/thm:feasibility}\,(i)); with declared constraints it
is a semidefinite feasibility question that the compiler decides on the
declaration it emits, and Requirement~\ref{req:compiled-declaration-only}
keeps every other declaration away from a key length.
}
\begin{requirement}[Admission of the declaration]\label{req:declaration-admission}
The engine decides conditions (i)--(vi) and (viii) of
Definition~\ref{def:admissible-declaration}, in the general format as that
definition reads them there, exactly on the \exttermref{../foundations/contract.pdf}{term-declared-data}{declared data}, and
computes nothing from a declaration that fails one of them, which it
refuses. It does not decide condition~(vii) on a declaration whose
constraint list is non-empty. A compiled reference companion is admitted
instead under Requirement~\ref{req:companion-admission}, which decides
exactly the components it shares with the declaration it accompanies and
none of its measured family.
\end{requirement}
\claimmeta{\metalabel{Motivated by}\mref{Def.}{def:admissible-declaration}, \mref{Req.}{req:exact-admission}, \mref{Req.}{req:companion-admission}, \mref{Req.}{req:compiled-declaration-only}, \extref{../protocol-compiler/contract.pdf}{protocol-compiler/thm:feasibility}
\metasep implemented (effective: accepted)
\metasep \metalabel{Code}\coderef{https://github.com/fraunhofersg/certiqs-key/blob/v0.1.0/src/certiqs_key/meat_core/declaration/general_validation.py}{certiqs_key.meat_core.declaration.general_validation:validate_general_declaration}}

% claim meat-core/req:companion-admission
\claimsummary{meat-core/req:companion-admission}{A compiled reference companion is admitted beside its declaration: shared components are decided equal, the acceptance set is read from it, and the maps are generated from the companion's own family, imported to be a POVM, its binary64 defects measured; only a short polish searches on it.}
\claimbookmark{meat-core/req:companion-admission}{Requirement}{req:companion-admission}{Admission of the reference companion beside its declaration}
\claimleadin{%
Exact admission exists so that the binary64 numbers of a declaration denote
exactly what was declared, which a run needs whose estimate and key length
are read off a solve over them. A reference companion has no run of its own:
its value is bounded below at points another declaration's search produced
and at the points of a short search continued on it from them, which serve
only as points, the hypotheses of the bound being decided at each; and the
numbers it is evaluated at are an evaluation of a family that is not on the
binary64 grid. Its admission
therefore decides exactly what can be decided exactly, the components it
shares with the declaration, and imports the rest from the claims that
establish it; the binary64 defects of its evaluation are measured and shown,
and decide nothing. What is evaluated on it, and when its number counts, is
Requirement~\ref{req:bracket-on-reference}, and how its number is certified
is Requirement~\ref{req:certified-evaluation}.
}
\begin{requirement}[Admission of the reference companion beside its declaration]\label{req:companion-admission}
Where the protocol compiler emits a compiled reference companion
$\Decl^{\mathrm{ref}}$ beside a declaration $\Decl$
(\xref{protocol-compiler/req:reference-companion-compiled}), the engine
admits it only beside $\Decl$, after admitting $\Decl$
(Requirements~\ref{req:exact-admission} and~\ref{req:declaration-admission}),
and only as follows.
\begin{enumerate}[label=(\roman*),leftmargin=2.4em]
\item \emph{Decided equal.} It decides exactly, on the declared data, that
  every component Proposition~\ref{prop:reference-companion}\,(i) names as
  shared is equal to $\Decl$'s, entry for entry and label for label: the
  shapes, the \exttermref{../foundations/contract.pdf}{term-sector-map}{sector maps} and the \exttermref{../protocol-compiler/contract.pdf}{term-joint-sector}{joint sectors}, the
  alphabets and their orders, Alice's slots, the \exttermref{../protocol-compiler/contract.pdf}{term-party-partition}{party partition}, the
  announcement, sifting and raw-key maps; that the companion's entries are
  those of $\Decl$ save the entries of the links the companion replaced, one
  per measured factor whose family it replaces, each of which carries the relaxed mark of
  \xref{reductions/req:reference-companion}\,(iii) and a Hermitian binary64
  operator with the blocks of its scope; and it reads the acceptance set of
  $\Decl$, the acceptance set being one per run. A companion that fails a
  decision of this item is not admitted, the reason naming the component,
  and the run proceeds without it. The \termref{term-statistics-map}{statistics map}, the
  \termref{term-single-round-protocol-map}{single-round protocol map}, the \termref{term-branch-map}{branch maps}, the
  objective and its gradient of the companion are generated from the
  companion's own slots, its measured family among them
  (Definition~\ref{def:maps}), as for any declaration; no map generated from
  $\Decl$'s measured family is used for the companion. The announcement,
  sifting and raw-key maps decided equal above are declared data, not
  generated maps.
\item \emph{Imported, never decided.} It decides none of the conditions of
  Definition~\ref{def:admissible-declaration} on the companion's measured
  family. That the family it denotes is a POVM, block diagonal for its
  sector map, is \xref{reductions/cor:reference-companion}\,(i), which rests
  on the claims that establish the \exttermref{../foundations/contract.pdf}{term-handed-reference}{handed reference}
  (\xref{foundations/def:exact-reference}) on the blocks where one is
  handed, on the exact decisions of the declarer's snap on the blocks where
  the companion carries the snapped block, and on (H-REF) for each instance
  whose link it replaces;
  and with (i) above it gives the companion conditions (i)--(vi) and (viii)
  by Proposition~\ref{prop:reference-companion}\,(ii). Condition (vii) is not
  decided on it. Nor does it decide that a block of the companion's measured
  family is a member of a \exttermref{../foundations/contract.pdf}{term-factor-enclosure}{factor enclosure} the companion carries with
  it, which is Corollary~\ref{cor:reference-companion-factors}, on the claims
  that establish the factor enclosures of the handed reference.
\item \emph{Evaluated and measured.} It reads the companion's family at the
  binary64 numbers nearest to the centres of its enclosure, rounded once,
  entrywise to nearest, when the family is docked, and records, per block of
  the measured family, the completeness defect and the least eigenvalue of
  those numbers, measured and not decided, and the largest radius of the
  enclosure. A defect so measured refuses nothing and marks nothing
  inadmissible; what these numbers mean for the bound evaluated on them is
  Remark~\ref{rem:fo-arithmetic}.
\item \emph{Nothing else.} It computes from the companion only what
  Requirements~\ref{req:bracket-on-reference}
  and~\ref{req:certified-evaluation} ask. The one program solved over it is
  the polish of Requirement~\ref{req:bracket-on-reference}\,(i), a search
  continued from a point of the search on $\Decl$, whose points serve only
  as points at which those requirements form brackets; no estimate is
  reported from it, and no key length is reported from it save through
  those requirements.
\end{enumerate}
Requirements~\ref{req:exact-admission} and~\ref{req:declaration-admission}
bind every declaration but a companion admitted under this requirement.
\end{requirement}
\claimmeta{\metalabel{Motivated by}\mref{Def.}{def:declaration}, \mref{Def.}{def:admissible-declaration}, \mref{Def.}{def:maps}, \mref{Def.}{def:branch}, \mref{Def.}{def:joint-sectors}, \mref{Prop.}{prop:reference-companion}, \mref{Cor.}{cor:reference-companion-factors}, \mref{Req.}{req:exact-admission}, \mref{Req.}{req:declaration-admission}, \mref{Rem.}{rem:fo-arithmetic}, \mref{Req.}{req:bracket-on-reference}, \mref{Req.}{req:certified-evaluation}, \extref{../foundations/contract.pdf}{foundations/def:exact-reference}, \extref{../protocol-compiler/contract.pdf}{protocol-compiler/req:reference-companion-compiled}, \extref{../reductions/contract.pdf}{reductions/cor:reference-companion}, \extref{../reductions/contract.pdf}{reductions/req:reference-companion}
\metasep \metalabel{Used by}\mref{Req.}{req:bracket-on-reference}
\metasep implemented (effective: accepted)
\metasep \metalabel{Code}\coderef{https://github.com/fraunhofersg/certiqs-key/blob/v0.1.0/src/certiqs_key/meat_core/declaration/companion.py}{certiqs_key.meat_core.declaration.companion:admit_reference_companion}}

% claim meat-core/req:acceptance-set-admission
\claimsummary{meat-core/req:acceptance-set-admission}{The engine admits an acceptance set only when it decides on the declared data that the set is admissible, through a geometry tag whose every well-formed datum denotes a closed convex subset of the simplex and data well formed for that tag.}
\claimbookmark{meat-core/req:acceptance-set-admission}{Requirement}{req:acceptance-set-admission}{Admission of the acceptance set}
\claimleadin{%
Whether a declared set is closed and convex is not read off its points: it
is carried by the form in which the set is declared. The two forms of the
statistics module's taxonomy, a single distribution and a polytope cut out
of the simplex by linear inequalities, denote an \termref{term-admissible-acceptance-set}{admissible acceptance set} whenever their data are well
formed, so deciding admissibility is deciding the tag and the data.
}
\begin{requirement}[Admission of the acceptance set]\label{req:acceptance-set-admission}
The engine admits an acceptance set only if it decides, exactly and on the
declared data, that the set is an admissible acceptance set
(Definition~\ref{def:admissible-acceptance-set}): the set carries a geometry
tag every well-formed datum of which denotes such a set, and its data are
well formed for that tag. It refuses every other declared acceptance set
(\iface{acceptance_set_ill_formed} for data ill formed for the geometry they
claim).
\end{requirement}
\claimmeta{\metalabel{Motivated by}\mref{Def.}{def:admissible-acceptance-set}
\metasep implemented (effective: accepted)
\metasep \metalabel{Code}\coderef{https://github.com/fraunhofersg/certiqs-key/blob/v0.1.0/src/certiqs_key/meat_core/declaration/validation.py}{certiqs_key.meat_core.declaration.validation:validate_acceptance_set}}

% claim meat-core/req:branch-register-components
\claimsummary{meat-core/req:branch-register-components}{An implementation of a branch map may drop a component of a register other than the raw-key register that is constant across its Kraus terms, and keeps the raw-key register and every varying component. By \hyperlink{claim-meat-core/rem:inert}{Remark~\ref*{rem:inert}} the first changes no quantity; the others carry what the derivations use.}
\claimbookmark{meat-core/req:branch-register-components}{Requirement}{req:branch-register-components}{Register components of an implemented branch map}
\claimleadin{%
By Remark~\ref{rem:inert}, a component of a register other than the
raw-key register $R$ that is constant across the Kraus terms of a
\termref{term-branch-map}{branch map} can be removed from it without changing any quantity of
this module, while a component that varies across them, and $R$ itself,
carry what the derivations use.
}
\begin{technicalclaim}
\begin{requirement}[Register components of an implemented branch map]\label{req:branch-register-components}
In implementing a branch map $\cG_{\hat c}$ or $\Theta_{\hat c}$ of
Definition~\ref{def:branch}, or a sector version $\cG_{\hat c,s}$ or
$\Theta_{\hat c,s}$ of either, the engine may drop a component of a register
other than $R$ whose value is constant across all Kraus terms of that map,
for example the $C_AC_B$ factor when the branch has a single announcement
class, and it retains the register $R$ and every component that varies
across the Kraus terms: the class label $\kappa$ across classes and the
outcome labels within a class.
\end{requirement}
\end{technicalclaim}
\claimmeta{\metalabel{Motivated by}\mref{Def.}{def:branch}, \mref{Rem.}{rem:inert}
\metasep accepted}

% claim meat-core/req:sector-report
\claimsummary{meat-core/req:sector-report}{The engine reports the four classes of every sector, and justifies every sector it deletes by the checked witness, the order in which the witness was evaluated and the retention of the witness's targets.}
\claimbookmark{meat-core/req:sector-report}{Requirement}{req:sector-report}{The sector report}
\claimleadin{%
A deletion of a set is licensed by the hypothesis of
Proposition~\ref{prop:removal-compose}, which is read on the recorded
witnesses, and the evaluation order of a witness is not part of the
declaration (Remark~\ref{rem:witness-exactness}); so the record has to carry
both for a deletion to be checked.
}
\begin{requirement}[The sector report]\label{req:sector-report}
The engine emits, for every sector, the four predicates of
Definition~\ref{def:classes} in \iface{sector_report}, and justifies there
every sector it deletes by the six conditions (W1)--(W6) checked on the
exhibited witness, by that witness with its evaluation order of (W6),
declared with the witness and recorded with it as
\iface{witness_evaluation_order}, and by the retention of the witness's
targets.
\end{requirement}
\claimmeta{\metalabel{Motivated by}\mref{Def.}{def:classes}, \mref{Prop.}{prop:removal-compose}, \mref{Rem.}{rem:witness-exactness}
\metasep implemented (effective: accepted)
\metasep \metalabel{Code}\coderef{https://github.com/fraunhofersg/certiqs-key/blob/v0.1.0/src/certiqs_key/meat_core/declaration/general_sectors.py}{certiqs_key.meat_core.declaration.general_sectors:general_sector_report}}

% claim meat-core/req:state-blind-test
\claimsummary{meat-core/req:state-blind-test}{The engine decides state-blindness by an exact identity of declared data: every off-diagonal entry of each block is the datum 0 and every diagonal entry the same datum. A block that is only nearly scalar is not state-blind.}
\claimbookmark{meat-core/req:state-blind-test}{Requirement}{req:state-blind-test}{Bit-level test of state-blindness}
\claimleadin{%
Classification runs only on a declaration already admitted, and admission
is exact (Requirement~\ref{req:exact-admission}): every declared block is
exactly Hermitian (Definition~\ref{def:declaration}, the
exact-admissibility paragraph), so a diagonal entry is real and the sign
condition of Definition~\ref{def:classes}\,(c) compares real numbers.
}
\begin{requirement}[Bit-level test of state-blindness]\label{req:state-blind-test}
The engine decides whether a sector $s$ is \termref{term-state-blind}{state-blind}
(Definition~\ref{def:classes}\,(c)) by a bit-level identity of declared
data. It reads $c_{v,s}$ off the declared block as its first diagonal
entry, $c_{v,s}\coloneqq(F^{(s)}_v)_{11}$, and requires
$F^{(s)}_v-c_{v,s}\idon{\Qsys_s}=0$ entrywise and exactly: every
off-diagonal entry of the block is the datum $0$, and every diagonal entry
is the same datum. Likewise, whenever $\CList\neq\varnothing$, it requires
$\hat O^{(s)}_k-o_{k,s}\idon{\Qsys_s}=0$ with
$o_{k,s}\coloneqq(\hat O^{(s)}_k)_{11}$. No subtraction, no norm and no
window enters: two binary64 data are equal or they are not, equality being
equality of the values they denote and not of their bit patterns
(\xref{foundations/def:binary64}). The engine decides the sign condition
$c_{v,s}\ge0$ by comparing that same datum with zero as a real number. A
block that is only \emph{nearly} scalar fails, and $s$ is then not
state-blind.
\end{requirement}
\claimmeta{\metalabel{Motivated by}\mref{Def.}{def:declaration}, \mref{Def.}{def:classes}, \mref{Req.}{req:exact-admission}, \extref{../foundations/contract.pdf}{foundations/def:binary64}
\metasep implemented (effective: accepted)
\metasep \metalabel{Code}\coderef{https://github.com/fraunhofersg/certiqs-key/blob/v0.1.0/src/certiqs_key/meat_core/declaration/general_sectors.py}{certiqs_key.meat_core.declaration.general_sectors:is_state_blind}}

% claim meat-core/req:removal-checklist
\claimsummary{meat-core/req:removal-checklist}{The engine flags a sector removable only when a witness it exhibits meets all six conditions of the removal checklist, each decided exactly. Every other outcome retains the sector, and no tolerance licenses a removal.}
\claimbookmark{meat-core/req:removal-checklist}{Requirement}{req:removal-checklist}{Setting the removable flag}
\claimleadin{%
The flag rests on the verification of the checklist and on nothing about
where the witness came from, so the witness may come from any search. A
deletion on an approximate match would overstate the value
(Remark~\ref{rem:classes}), so every doubt retains the sector; and with a
\exttermref{../protocol-compiler/contract.pdf}{term-party-local}{party-local} entry or several
\exttermref{../protocol-compiler/contract.pdf}{term-measured-factor}{measured factors} the checklist has nothing to
certify or nothing to read (Remark~\ref{rem:removal-general-format}).
}
\begin{requirement}[Setting the removable flag]\label{req:removal-checklist}
The engine sets the flag \iface{removable} for a sector $s^{*}$ of a
declaration with one \exttermref{../protocol-compiler/contract.pdf}{term-measured-factor}{measured factor} when, and only when, a candidate
witness exhibited by the implementation is a witness for $s^{*}$ in the
sense of Definition~\ref{def:classes}\,(d): the list $\CList$ carries no
entry declared \exttermref{../protocol-compiler/contract.pdf}{term-party-local}{party-local} and the candidate meets all six of
(W1)--(W6); it sets the flag on no other ground. The
witness may come from any search, and it is never read as declared input: a
declaration supplies no witness. Each condition is decided exactly, at no
window: (W1) by the bit-level identity of
Requirement~\ref{req:state-blind-test}, (W2) as a comparison of sector
indices, (W5) as the finite combinatorial predicate of key-inertness, (W3)
and (W4) on the exhibited data as those conditions state, and (W6) on the
residual formed in the exhibited evaluation order; a witness whose
evaluation rounds anywhere, or whose residual is not the datum $0$, sets no
flag. Every other outcome --- a nonzero residual, a rounding operation, an
absent witness, a witness that names $s^{*}$, a target that is not
key-inert, a block that is not exactly scalar --- retains the sector. The
flag is set for no sector of a declaration whose list carries an entry
declared party-local, nor of a declaration with more than one measured
factor. No tolerance licenses a removal:
\iface{sector_criterion} is decided at exactly zero and not at the
\exttermref{../foundations/contract.pdf}{term-declared-tolerance}{declared tolerance} $\tnum$.
\end{requirement}
\claimmeta{\metalabel{Motivated by}\mref{Def.}{def:classes}, \mref{Thm.}{thm:removal}, \mref{Rem.}{rem:removal-general-format}, \mref{Rem.}{rem:classes}, \mref{Req.}{req:state-blind-test}
\metasep implemented (effective: accepted)
\metasep \metalabel{Code}\coderef{https://github.com/fraunhofersg/certiqs-key/blob/v0.1.0/src/certiqs_key/meat_core/declaration/general_sectors.py}{certiqs_key.meat_core.declaration.general_sectors:general_sector_report}}

% claim meat-core/req:structure-operations
\claimsummary{meat-core/req:structure-operations}{The engine never edits the declared sector list: it evaluates the objective over the key-active sectors and deletes from the variable list only a set of flagged sectors that meets the composition hypothesis. It never merges or averages sectors.}
\claimbookmark{meat-core/req:structure-operations}{Requirement}{req:structure-operations}{Operations on the declared sectors}
\claimleadin{%
Dropping a constraint-active sector (Proposition~\ref{prop:retention}) and
averaging sectors (Proposition~\ref{prop:no-averaging}) can both overstate
the key, while the two operations below are identities.
}
\begin{requirement}[Operations on the declared sectors]\label{req:structure-operations}
The engine treats the declared sector list as input it never edits; sectors
leave the variable list only by the deletion of~(ii). The operations
it may perform on the declared sectors are exactly: (i) evaluate the
objective over the key-active sectors only (Definition~\ref{def:keyactive};
an identity, by Proposition~\ref{prop:exact-cancellation}); (ii) delete
sectors carrying the flag \iface{removable} of
Definition~\ref{def:classes}\,(d), set as
Requirement~\ref{req:removal-checklist} states --- and only those, and among those only
a set meeting the hypothesis of Proposition~\ref{prop:removal-compose},
every deleted sector's witness naming targets outside the deleted set (an
identity, by Theorem~\ref{thm:removal} and that proposition); (iii) nothing
else. In particular no simplification step merges sectors, and the only
structure-changing simplification is the deletion of~(ii).
\end{requirement}
\claimmeta{\metalabel{Motivated by}\mref{Prop.}{prop:exact-cancellation}, \mref{Def.}{def:keyactive}, \mref{Def.}{def:classes}, \mref{Thm.}{thm:removal}, \mref{Prop.}{prop:retention}, \mref{Prop.}{prop:no-averaging}, \mref{Prop.}{prop:removal-compose}, \mref{Req.}{req:removal-checklist}
\metasep implemented (effective: accepted)
\metasep \metalabel{Code}\coderef{https://github.com/fraunhofersg/certiqs-key/blob/v0.1.0/src/certiqs_key/meat_core/declaration/sectors.py}{certiqs_key.meat_core.declaration.sectors:delete_sectors}}

% claim meat-core/req:single-sector-deletion
\claimsummary{meat-core/req:single-sector-deletion}{An implementation that does not check the set hypothesis of \hyperlink{claim-meat-core/prop:removal-compose}{Proposition~\ref*{prop:removal-compose}} deletes at most one sector per run. A single deleted sector meets that hypothesis automatically.}
\claimbookmark{meat-core/req:single-sector-deletion}{Requirement}{req:single-sector-deletion}{At most one deleted sector without the composition check}
\begin{technicalclaim}
\begin{requirement}[At most one deleted sector without the composition check]\label{req:single-sector-deletion}
An implementation that will not check the hypothesis of
Proposition~\ref{prop:removal-compose} deletes at most one sector per run,
which is the safe reading of that hypothesis: a singleton
$\SecSet_{\mathrm{del}}$ satisfies it automatically, by (W2) of
Definition~\ref{def:classes}.
\end{requirement}
\end{technicalclaim}
\claimmeta{\metalabel{Motivated by}\mref{Def.}{def:classes}, \mref{Prop.}{prop:removal-compose}
\metasep accepted}

\section{Requirements: the Choi program and gradients}\label{group-meat-core/req:frame-constructed}
\sectionintro{The engine constructs every \termref{term-support-frame}{support frame} it uses rather than accepting one, checks it exactly, uses the \termref{term-diagonal-frame}{diagonal frame} only for an exactly diagonal marginal and records the frame; it forms every Choi block in one fixed convention. It drops the conjugation of a constraint row only where the frame's construction exhibits tightness, and records the form used. It evaluates gradients in split form and only at maximal support, never reports the order-one limit as a bound, and records a \termref{term-dropped-set}{dropped set} with the value computed from it.}

% claim meat-core/req:frame-constructed
\claimsummary{meat-core/req:frame-constructed}{The engine constructs the declared isometry of every support frame it uses and never takes it as declared input. It decides the two identities of the frame on that isometry exactly, at zero, and refuses a frame that fails them.}
\claimbookmark{meat-core/req:frame-constructed}{Requirement}{req:frame-constructed}{Construction and exact check of the support frame}
\claimleadin{%
Lemma~\ref{lem:choi-marginal} and Theorem~\ref{thm:choi} hold on any
isometry satisfying \eqref{eq:iota}, and both constructions of
Definition~\ref{def:support-frame} satisfy it by arithmetic on $0$ and $1$,
the \termref{term-full-frame}{full frame} for every admissible $\tau_A$. A frame is therefore
constructed to satisfy \eqref{eq:iota} rather than fitted to it.
}
\begin{requirement}[Construction and exact check of the support frame]\label{req:frame-constructed}
In \exttermref{../transmitter/contract.pdf}{term-trusted-source-mode}{trusted-source mode} the engine constructs the
\termref{term-declared-isometry}{declared isometry} $\iota$ of every \termref{term-support-frame}{support frame} it uses and
never takes $\iota$ as declared input:
by the full frame, which is always available, by the
\termref{term-diagonal-frame}{diagonal frame}, or by another construction for which it exhibits
\eqref{eq:iota} exactly. It decides both identities of \eqref{eq:iota} on
the isometry it constructed, entrywise as binary64 data at exactly zero
(\iface{choi_isometry}), and refuses a frame that fails either of them
with \iface{choi_isometry_inexact}.
\end{requirement}
\claimmeta{\metalabel{Motivated by}\mref{Def.}{def:support-frame}, \mref{Lem.}{lem:choi-marginal}, \mref{Thm.}{thm:choi}
\metasep implemented (effective: accepted)
\metasep \metalabel{Code}\coderef{https://github.com/fraunhofersg/certiqs-key/blob/v0.1.0/src/certiqs_key/meat_core/objective/joint.py}{certiqs_key.meat_core.objective.joint:marginal_frame}}

% claim meat-core/req:diagonal-frame-test
\claimsummary{meat-core/req:diagonal-frame-test}{The engine uses the diagonal frame only when every off-diagonal entry of the declared \ensuremath{\tau}\_A is the datum 0, compared at no window. For a \ensuremath{\tau}\_A that is only nearly diagonal the diagonal frame is unavailable, and the full frame always is.}
\claimbookmark{meat-core/req:diagonal-frame-test}{Requirement}{req:diagonal-frame-test}{Exact diagonality test for the diagonal frame}
\claimleadin{%
The diagonal frame is tight by construction only because $\tau_A$ is
diagonal. Read at a positive window, the test would pass a singular
$\tau_A$ with small off-diagonal entries and return a frame that is not
tight (Remark~\ref{rem:tightness-by-construction}).
}
\begin{requirement}[Exact diagonality test for the diagonal frame]\label{req:diagonal-frame-test}
The engine uses the \termref{term-diagonal-frame}{diagonal frame} of
Definition~\ref{def:support-frame}\,(b) only when every off-diagonal entry
of the declared $\tau_A$ is the datum $0$. It decides this by comparing
each off-diagonal entry with $0$ as a binary64 value, at no window, as
Requirement~\ref{req:state-blind-test} compares the entries of a sector
block, and the comparison is a clause of \iface{alice_marginal_exact}. A
$\tau_A$ that is only \emph{nearly} diagonal is not diagonal: the diagonal
frame is unavailable for it, and the \termref{term-full-frame}{full frame} always is.
\end{requirement}
\claimmeta{\metalabel{Motivated by}\mref{Def.}{def:admissible-declaration}, \mref{Def.}{def:support-frame}, \mref{Rem.}{rem:tightness-by-construction}, \mref{Req.}{req:state-blind-test}
\metasep implemented (effective: accepted)
\metasep \metalabel{Code}\coderef{https://github.com/fraunhofersg/certiqs-key/blob/v0.1.0/src/certiqs_key/meat_core/objective/joint.py}{certiqs_key.meat_core.objective.joint:marginal_frame}}

% claim meat-core/req:frame-recorded
\claimsummary{meat-core/req:frame-recorded}{In trusted-source mode every run records the support frame it used, with the declared isometry, beside any Choi block it reports. In untrusted-source mode there is no frame, and none is recorded.}
\claimbookmark{meat-core/req:frame-recorded}{Requirement}{req:frame-recorded}{Record of the support frame}
\claimleadin{%
The reported Choi blocks and the channel reading of
Remark~\ref{rem:ordering} are relative to the declared isometry, while
$\hlow$ and every quantity the module outputs within one frame are not
(Remark~\ref{rem:basis-invariance}); so the frame is recorded rather than
made canonical.
}
\begin{requirement}[Record of the support frame]\label{req:frame-recorded}
In trusted-source mode every run records, beside any Choi block it
reports, the \termref{term-support-frame}{support frame} it used: the construction that produced
it, the dimension $d_{A_0}$ and the \termref{term-declared-isometry}{declared isometry} $\iota$, per
Alice sector in the general format, as \iface{choi_input_basis} in
\iface{optimizer_certificate}. In \exttermref{../transmitter/contract.pdf}{term-untrusted-source-mode}{untrusted-source mode} no $A_0$ and no
$\iota$ exist (Definition~\ref{def:untrusted-choi}\,(a)), and the field is
absent.
\end{requirement}
\claimmeta{\metalabel{Motivated by}\mref{Def.}{def:support-frame}, \mref{Rem.}{rem:ordering}, \mref{Rem.}{rem:basis-invariance}, \mref{Def.}{def:untrusted-choi}
\metasep accepted
\metasep \metalabel{Code}\coderef{https://github.com/fraunhofersg/certiqs-key/blob/v0.1.0/src/certiqs_key/solvers/frank_wolfe/joint_driver.py}{certiqs_key.solvers.frank_wolfe.joint_driver:joint_run_metadata}}

% claim meat-core/req:choi-convention
\claimsummary{meat-core/req:choi-convention}{The engine forms and reads every Choi block with the input factor A\ensuremath{_{0}} first and takes the transpose of the channel reading in the declared basis of A\ensuremath{_{0}}.}
\claimbookmark{meat-core/req:choi-convention}{Requirement}{req:choi-convention}{Factor order and transpose basis of the Choi blocks}
\claimleadin{%
The channel a Choi block represents depends on the factor order and on the
basis the transpose is taken in (Remark~\ref{rem:ordering}).
}
\begin{technicalclaim}
\begin{requirement}[Factor order and transpose basis of the Choi blocks]\label{req:choi-convention}
The engine forms and reads every Choi block $J_s$ of Theorem~\ref{thm:choi}
with the input factor $A_0$ first, and takes the transpose of the channel
reading of Remark~\ref{rem:ordering} in the
\termref{term-declared-basis-of-a-0}{declared basis of $A_0$}, as that remark declares them.
\end{requirement}
\end{technicalclaim}
\claimmeta{\metalabel{Motivated by}\mref{Thm.}{thm:choi}, \mref{Rem.}{rem:ordering}
\metasep accepted}

% claim meat-core/req:row-form-by-construction
\claimsummary{meat-core/req:row-form-by-construction}{The engine imposes a constraint row without its conjugation only where the construction of the frame exhibits a trivial kernel of W\_\ensuremath{\tau} without computing a rank, and the conjugated row everywhere else. No numerical test on \ensuremath{\tau}\_A licenses the frame-free form.}
\claimbookmark{meat-core/req:row-form-by-construction}{Requirement}{req:row-form-by-construction}{Frame-free rows only where the construction exhibits tightness}
\claimleadin{%
By Proposition~\ref{prop:frame-free-rows} the frame-free form of a row is
exact on a tight frame and only a restriction on any other, while the
conjugated form is exact on every frame (Theorem~\ref{thm:choi}). Deciding
tightness of a given $\tau_A$ would be a rank decision on a computed
spectrum; the construction of the frame can exhibit it instead
(Remark~\ref{rem:tightness-by-construction}).
}
\begin{requirement}[Frame-free rows only where the construction exhibits tightness]\label{req:row-form-by-construction}
In trusted-source mode the engine imposes a homogeneous row of
\eqref{eq:choi-feas} in its frame-free form \eqref{eq:frame-free-row}, and
in the general format the row of an Alice sector $t$ of a homogeneous or a
party-local entry in its unconjugated form, only where the construction
that produced the frame (of sector $t$) exhibits $\ker W_{\tau}=0$
($\ker W_t=0$) exactly, with no rank computed and no window declared.
Of the two constructions of Definition~\ref{def:support-frame} that is the
\termref{term-diagonal-frame}{diagonal frame}; on a frame produced by the \termref{term-full-frame}{full frame}
construction the conjugated form is imposed, whatever $\tau_A$ is, and so
it is on every frame whose construction does not exhibit
$\ker W_{\tau}=0$. The engine licenses the frame-free form by no numerical
test on $\tau_A$.
\end{requirement}
\claimmeta{\metalabel{Motivated by}\mref{Def.}{def:support-frame}, \mref{Thm.}{thm:choi}, \mref{Prop.}{prop:frame-free-rows}, \mref{Rem.}{rem:tightness-by-construction}
\metasep implemented (effective: accepted)
\metasep \metalabel{Code}\coderef{https://github.com/fraunhofersg/certiqs-key/blob/v0.1.0/src/certiqs_key/solvers/oracles/joint_rows.py}{certiqs_key.solvers.oracles.joint_rows:prebuilt_rows}}

% claim meat-core/req:row-form-record
\claimsummary{meat-core/req:row-form-record}{In trusted-source mode a run records, for every declared constraint row, whether it imposed the conjugated or the frame-free form, per Alice sector in the general format.}
\claimbookmark{meat-core/req:row-form-record}{Requirement}{req:row-form-record}{Record of the form each constraint row was imposed in}
\claimleadin{%
The two forms of a row define the same set on a tight frame and in general
not on another (Proposition~\ref{prop:frame-free-rows}), so which one a run imposed is part
of what its reported iterates mean.
}
\begin{requirement}[Record of the form each constraint row was imposed in]\label{req:row-form-record}
In trusted-source mode a run records, for every declared constraint row,
which of the two forms of Proposition~\ref{prop:frame-free-rows} it
imposed, per Alice sector in the general format, as
\iface{constraint_row_form} in \iface{optimizer_certificate}, beside the
frame it records under Requirement~\ref{req:frame-recorded}.
\end{requirement}
\claimmeta{\metalabel{Motivated by}\mref{Prop.}{prop:frame-free-rows}, \mref{Req.}{req:frame-recorded}
\metasep implemented (effective: accepted)
\metasep \metalabel{Code}\coderef{https://github.com/fraunhofersg/certiqs-key/blob/v0.1.0/src/certiqs_key/solvers/frank_wolfe/joint_driver.py}{certiqs_key.solvers.frank_wolfe.joint_driver:joint_run_metadata}}

% claim meat-core/req:split-form-gradient
\claimsummary{meat-core/req:split-form-gradient}{The engine evaluates the gradient of the pinched trace power in the split form of \hyperlink{claim-meat-core/thm:gradient}{Theorem~\ref*{thm:gradient}} (i): divided differences over the positive spectrum, pinched operators on the ambient space. It compresses the whole computation only to a pinching-invariant space containing the branch support.}
\claimbookmark{meat-core/req:split-form-gradient}{Requirement}{req:split-form-gradient}{Split-form evaluation of the trace-power gradient}
\claimleadin{%
A branch operator is in general singular, and its
\termref{term-maximal-branch-support}{maximal branch support} is in general not invariant under the
\termref{term-key-pinching}{key pinching} (Lemma~\ref{lem:branch-support}\,(iv)); compressing the whole
computation to its support then computes a different quantity, as the
rank-one witness in the proof of Theorem~\ref{thm:gradient} shows.
}
\begin{requirement}[Split-form evaluation of the trace-power gradient]\label{req:split-form-gradient}
The engine evaluates the gradient of the pinched trace power at a branch
operator $X$ in the split form of Theorem~\ref{thm:gradient}\,(i): the
divided-difference map $DX^{p}[\cdot]$ compressed to
$\operatorname{supp}X$, its kernel summed over the strictly positive
eigenvalues of $X$ only, and $\cZ(X^{p})$, $Y^{\alpha-1}$ and
$\cZ(Y^{\alpha-1})$ formed on the ambient branch space. It compresses the
whole computation to a subspace $S'$ only where $S'$ contains the maximal
branch support $W_{\hat c,s}$ of Definition~\ref{def:branch-support} and
$Z_jS'\subseteq S'$ for every $j$.
\end{requirement}
\claimmeta{\metalabel{Motivated by}\mref{Def.}{def:branch-support}, \mref{Lem.}{lem:branch-support}, \mref{Thm.}{thm:gradient}
\metasep implemented (effective: accepted)
\metasep \metalabel{Code}\coderef{https://github.com/fraunhofersg/certiqs-key/blob/v0.1.0/src/certiqs_key/meat_core/objective/gradient.py}{certiqs_key.meat_core.objective.gradient:pullback_trace_power_gradient}}

% claim meat-core/req:gradient-at-maximal-support
\claimsummary{meat-core/req:gradient-at-maximal-support}{The engine evaluates a gradient of the pinched trace power only at an iterate whose support is kept maximal, for instance after an explicit perturbation. At lower support a direction that enlarges the support can have an infinite one-sided derivative (\hyperlink{claim-meat-core/thm:gradient}{Theorem~\ref*{thm:gradient}} (ii)).}
\claimbookmark{meat-core/req:gradient-at-maximal-support}{Requirement}{req:gradient-at-maximal-support}{Gradient evaluation at maximal support}
\claimleadin{%
Theorem~\ref{thm:gradient}\,(ii) places every infinite one-sided
derivative at a direction that strictly enlarges the support, and at
maximal support no feasible direction does (Lemma~\ref{lem:fo-face}). A
gradient taken after an explicit perturbation is the gradient of another
function (Remark~\ref{rem:fo-nopert}), and on the first-order path it
removes that iterate's bracket from the bound
(Requirement~\ref{req:first-order-withdrawal}).
}
\begin{requirement}[Gradient evaluation at maximal support]\label{req:gradient-at-maximal-support}
The engine evaluates a gradient of the pinched trace power only at an
iterate whose support is (kept) maximal, maximality being that of
\eqref{eq:Wcs} for the branch operators, for instance after an explicit
perturbation.
\end{requirement}
\claimmeta{\metalabel{Motivated by}\mref{Def.}{def:branch-support}, \mref{Lem.}{lem:branch-support}, \mref{Thm.}{thm:gradient}, \mref{Lem.}{lem:fo-face}, \mref{Rem.}{rem:fo-nopert}, \mref{Req.}{req:first-order-withdrawal}
\metasep implemented (effective: accepted)
\metasep \metalabel{Code}\coderef{https://github.com/fraunhofersg/certiqs-key/blob/v0.1.0/src/certiqs_key/meat_core/objective/gradient.py}{certiqs_key.meat_core.objective.gradient:deficit_gradient}}

% claim meat-core/req:limit-not-a-bound
\claimsummary{meat-core/req:limit-not-a-bound}{The engine reports the order-one limit of \hyperlink{claim-meat-core/thm:oracle}{Theorem~\ref*{thm:oracle}}, however computed, as a bound of no kind, and never as the first-order bound. The limit validates an implementation as \ensuremath{\alpha} \ensuremath{\to} 1 and says nothing about h\_low at a finite order.}
\claimbookmark{meat-core/req:limit-not-a-bound}{Requirement}{req:limit-not-a-bound}{The order-one limit as a validation reference only}
\claimleadin{%
Remark~\ref{rem:oracle} says what the limit $h_1$ of
Theorem~\ref{thm:oracle} validates: the $\alpha\to1^{+}$ behaviour of an
implementation, and no finite-$\alpha$ value.
}
\begin{requirement}[The order-one limit as a validation reference only]\label{req:limit-not-a-bound}
The engine never reports the limit $h_1$ of \eqref{eq:oracle}, however
computed, as a bound of any kind; in particular it never reports it as the
first-order bound $\hlb$ of Theorem~\ref{thm:firstorder}.
\end{requirement}
\claimmeta{\metalabel{Motivated by}\mref{Thm.}{thm:oracle}, \mref{Rem.}{rem:oracle}, \mref{Thm.}{thm:firstorder}
\metasep accepted}

% claim meat-core/req:dropped-set-record
\claimsummary{meat-core/req:dropped-set-record}{A run records its dropped set and reports the dropped value under its own name together with that set; for a non-empty set, never under the name of the undropped value.}
\claimbookmark{meat-core/req:dropped-set-record}{Requirement}{req:dropped-set-record}{Recording the dropped set}
\claimleadin{%
A \termref{term-dropped-set}{dropped set} is a choice of the run and not \exttermref{../foundations/contract.pdf}{term-declared-data}{declared data},
and the value it gives lies at or below $\hlow$
(Proposition~\ref{prop:drop-bound}); the number means something only
beside the set it was computed for.
}
\begin{requirement}[Recording the dropped set]\label{req:dropped-set-record}
A run records its dropped set $\mathcal{D}$ (Definition~\ref{def:dropped})
with the run, and reports the number it computes from \eqref{eq:hlowD} as
$\underline{h}^{\mathcal{D}}_{\alpha}$ together with $\mathcal{D}$, never,
when $\mathcal{D}\neq\varnothing$, under the name $\hlow$.
\end{requirement}
\claimmeta{\metalabel{Motivated by}\mref{Def.}{def:dropped}, \mref{Prop.}{prop:drop-bound}
\metasep implemented (effective: accepted)
\metasep \metalabel{Code}\coderef{https://github.com/fraunhofersg/certiqs-key/blob/v0.1.0/src/certiqs_key/solvers/frank_wolfe/general_keyrate.py}{certiqs_key.solvers.frank_wolfe.general_keyrate:GeneralKeyRateResult}}

\section{Requirements: values and key lengths}\label{group-meat-core/req:infinite-value}
\sectionintro{The engine carries an infinite value through as the infeasibility of the acceptance set and emits no key length for it, and on a run whose honest statistics are supplied or whose declaration carries a constraint the caller declared it bounds the value with every pair dropped once, at an order fixed before any order is evaluated, records the cap it implies on the probability that devices the declaration admits pass the acceptance test, and refuses the run where that cap is at most a threshold, one half by default. It reports the \termref{term-maximal-key-length}{maximal key length} with its floor and its clamp at 0. It evaluates the objective only at a R\'enyi order in (1, \ensuremath{\infty}) and reports a key length only at one in (1, 2), and only on a declaration the compiler emitted. It refuses a run whose security budgets, the authentication error included, sum to at least 1, where the security statement is empty. An instance that adopts the weaker secrecy convention declares it. A key length under declared constraints, or with a declared marginal, is emitted only behind the gates that license it, and a degenerate licensed-class check is reported as no evidence.}
\sectionimports{\extref{../foundations/contract.pdf}{foundations/def:tiers}}

% claim meat-core/req:infinite-value
\claimsummary{meat-core/req:infinite-value}{The engine carries h\_low = +\ensuremath{\infty} as a value of its extended-real output and clamps it to no finite range. At +\ensuremath{\infty} it emits no key length and reports the mismatch between the declaration and the acceptance set.}
\claimbookmark{meat-core/req:infinite-value}{Requirement}{req:infinite-value}{Propagation of an infinite value}
\claimleadin{%
The value has no finite a-priori ceiling, and it is $+\infty$ when the
declared acceptance set is infeasible against the declaration
(Proposition~\ref{prop:value-range}\,(iii), (iv)); on such a declaration no
attack the declared list admits is accepted and no run produces a key
(Remark~\ref{rem:hlow-range}). It is $+\infty$ as well when the
\exttermref{../protocol-compiler/contract.pdf}{term-feasible-set}{feasible set} itself is empty, \eqref{eq:hlow-def} being then an infimum
over nothing; on an \termref{term-admissible-declaration}{admissible declaration} that cannot happen
(Definition~\ref{def:admissible-declaration}\,(vii)), and the bounds of
Theorem~\ref{thm:firstorder}\,(v) cover it where condition (vii)
is not decided.
}
\begin{requirement}[Propagation of an infinite value]\label{req:infinite-value}
The engine propagates $\hlow=+\infty$, the infeasibility of the declared
acceptance set, rather than clamping it to any finite range; its output
\iface{h_low} has the extended-real domain $[0,+\infty]$. If $\hlow=+\infty$
reaches Theorem~\ref{thm:keylength}, the engine emits no key-length claim and
no numeric key length, and reports the mismatch between the declaration and
the acceptance set (\iface{infeasible_acceptance_set}).
\end{requirement}
\claimmeta{\metalabel{Motivated by}\mref{Def.}{def:admissible-declaration}, \mref{Def.}{def:feasible}, \mref{Prop.}{prop:value-range}, \mref{Rem.}{rem:hlow-range}, \mref{Thm.}{thm:keylength}, \mref{Thm.}{thm:firstorder}
\metasep \metalabel{Used by}\mref{Rem.}{rem:declared-measurement}
\metasep accepted}

% claim meat-core/req:unreachable-acceptance-set
\claimsummary{meat-core/req:unreachable-acceptance-set}{On a checked run the engine bounds the penalty-only value once, before any order is evaluated, and records the cap on the probability that devices as declared pass the acceptance test, proven where b is certified. A run whose cap is at most a threshold, one half by default, is refused.}
\claimbookmark{meat-core/req:unreachable-acceptance-set}{Requirement}{req:unreachable-acceptance-set}{The check of the acceptance set against the feasible statistics}
\claimleadin{%
The engine's own honest models are built to match the declared Alice
marginal, but no claim states that the law they produce is the statistics of
a state of the \exttermref{../protocol-compiler/contract.pdf}{term-feasible-set}{feasible set}, and a constraint entry the caller declares
can exclude every state that produces it; a law or envelope the caller
supplies carries no such construction at all. Where no accepted distribution
is the statistics $\Phimap[\omega]$ of a feasible state, the value has the
positive floor $\frac{\alpha}{\alpha-1}\sfD_{\min}$, which grows without limit
as $\alpha$ tends to one (Proposition~\ref{prop:minimal-divergence}\,(ii)),
while the entropy terms stay bounded
(Proposition~\ref{prop:value-range}\,(ii)): a key length formed from it is a
bound for the declaration, but part of it is a penalty for the mismatch
between the declared devices and the accepted statistics rather than
entropy. What a positive bound rules out is a per-round law: no feasible
state has its statistics in the acceptance set. It does not rule out that
the frequencies of a finite run of such a state fall in the set, but
Proposition~\ref{prop:model-pass-bound} bounds the probability of that
event: a run whose rounds follow statistics of feasible states passes the
test with probability at most $2^{-n\sfD_{\min}}$. That cap is the
completeness statement of the model the declaration states, proven where
the bound it is formed from is certified and of the first-order tier
otherwise. Where it is at most one half, devices behaving as the
declaration says would fail the acceptance test at least as often as pass
it: a protocol whose own model
expects to abort is not the protocol the caller means to run, and a key
length formed under it is withheld. The infinite value is
the extreme case of the same mismatch
(Requirement~\ref{req:infinite-value}). The penalty-only program is the
\termref{term-dropped-program}{dropped program} with every pair dropped, which the first-order bound
(Corollary~\ref{cor:fo-dropped}) and its certified evaluation
(Lemma~\ref{lem:certified-bracket}\,(iv),
Requirement~\ref{req:certified-evaluation}) already cover, so no other
solver is needed. By Proposition~\ref{prop:minimal-divergence}\,(iv) one
solve at an order $\alpha_0$ bounds $\sfD_{\min}$ from below, and the cap
reads only that bound; so the solve is made once, at
an order fixed before any is evaluated and below none the run evaluates, the
lowest the run evaluates where its order search fixes that order in advance
and otherwise the lower end of the search's declared domain, with a bounded
share of the wall budget of one evaluation of the run: a first-order number
bounds the value from below at any effort, so the share bears on the size of
the bound only, never on its soundness. Nor does the
sign of the solve's estimate bear on the bound: the first-order bound holds
at points of the \termref{term-relaxed-feasible-set}{relaxed feasible set}, where the objective may fall below
zero, and on statistics that a feasible state produces the value is zero, so
the estimate's sign is set by rounding and by the relaxation window, at a
relaxed point as low as
$-\frac{\alpha_0}{\alpha_0-1}\log_2(1+d_{A_0}\varepsilon)$
(Proposition~\ref{prop:relaxed-set-properties}\,(c$''$), through
\eqref{eq:penaltyshift-fw}, where $d_{A_0}\varepsilon<1$). Where the acceptance set is built from
an honest envelope, a positive bound shows that the envelope's premise
\xref{statistics/ass:reach} fails.
}
\begin{requirement}[The check of the acceptance set against the feasible statistics]\label{req:unreachable-acceptance-set}
A run is \emph{checked} when its acceptance set is built from an honest law
or an honest envelope that the caller gives as numbers, such as one measured
on a device or read from a table, rather than one the engine computes from
the declared devices, or when its declaration carries a constraint entry
that the caller declares rather than one the compiler or a reduction
derives, unless the run sets the option \iface{statistics_check} to
\iface{off} (default \iface{on}). On every run, checked or not, the engine
records under \iface{statistics_consistency} whether the run is checked and,
where it is not, why (out of scope, or opted out), and the report says so
beside every key length; a run that is not checked names the premise
\xref{statistics/ass:reach} as unchecked, the caller's to vouch for. On
every checked run the engine does the following once, before it evaluates
the run at any R\'enyi order and before it reports a key length.
\begin{enumerate}[label=(\roman*),leftmargin=2.4em]
\item \emph{The order and the solve.} It takes $\alpha_0$ to be fixed before
  any order is evaluated: the run's fixed order, or the lowest order its
  order search evaluates, where the search fixes that order before it
  evaluates any and evaluates none below it, and otherwise the lower end of
  the domain the search declares, which every order the search evaluates is
  at least. The engine evaluates the program only at an order in
  $(1,\infty)$ (Requirement~\ref{req:objective-alpha-domain}), so the solve
  runs only where $\alpha_0>1$, and at any other $\alpha_0$ it is refused,
  as (ii) records; a key length is reported only at an order in $(1,2)$
  (Requirement~\ref{req:key-length-alpha-domain}), so a fixed order at
  which one is reported lies there, and that the orders a search evaluates
  and the lower end of its declared domain lie in $(1,2)$ is the bound the
  run tier places on the domain it accepts, which no claim of this contract
  states. On the
  declaration the run searches and with the run's
  acceptance set, it runs the first-order search of the
  \termref{term-dropped-program}{dropped program} for the \termref{term-dropped-set}{dropped set}
  $\mathcal{D}_{\mathrm{all}}$ of every pair
  (Proposition~\ref{prop:minimal-divergence}) at $\alpha_0$, with the
  method, the iteration cap and the gap target the run applies to one
  evaluation of its own, and with the share
  \iface{statistics_check_budget_share} of the wall budget of one
  evaluation of the run, where it has one, as its wall budget. A smaller
  share can leave a mismatch unproven: $b$ smaller, or not positive, so
  that (iv) may leave unrefused a run that a longer solve would have
  refused, and the
  cap of (iv) looser; no bound is made unsound by it.
\item \emph{The bound.} From that solve it forms the first-order number of
  the dropped program and, unless the run sets \iface{certified_evaluation}
  to \iface{off}, the certified evaluation of that number, on the
  declaration the run searches and at that solve's points, by parts
  (ii)--(v) of Requirement~\ref{req:certified-evaluation}, whatever
  declaration the run's key length is formed on. The number $b$ is the
  certified number where a point is certified and the first-order number
  otherwise, with its \exttermref{../foundations/contract.pdf}{term-tier}{tier}. The solve is not refused on the sign of
  its estimate $O_0$ of (v), which is no \exttermref{../foundations/contract.pdf}{term-estimate}{estimate} of the run and
  which Requirement~\ref{req:negative-estimate-refused} does not read: $b$
  is formed from the solve's points whatever that sign, and a negative
  $O_0$ is recorded as a diagnostic and refuses nothing. Where no bound is
  formed (the first-order number withdrawn, the solve not completed, or the
  solve refused at an order outside $(1,\infty)$), the run proceeds with
  the reason recorded.
\item \emph{The proven part.} Where $b>0$, for every order $\alpha'$ the run
  evaluates, the proven floor
  $f(\alpha')\coloneqq\frac{\alpha'}{\alpha'-1}\cdot
  \frac{\alpha_0-1}{\alpha_0}\,b$ satisfies
  $f(\alpha')\le\frac{\alpha'}{\alpha'-1}\sfD_{\min}\le\underline{h}_{\alpha'}$
  (Proposition~\ref{prop:minimal-divergence}\,(ii), (iv)), with $\sfD_{\min}$
  and $\underline{h}_{\alpha'}$ those of the declaration the run searches: a
  part of that declaration's value at $\alpha'$ that the mismatch alone
  accounts for. Where $b\le0$, $f(\alpha')\coloneqq0$. Where no bound was
  formed, $f(\alpha')$ is not defined.
\item \emph{The cap and the refusal.} Let $\tau\in(0,1]$ be the run's option
  \iface{min_model_pass_probability}, an exact rational number, default
  $\frac12$. Where $b$ was formed and $b>0$, the engine forms the exponent
  $x\coloneqq n\,\frac{\alpha_0-1}{\alpha_0}\,b$ exactly from the run's
  round count $n$, $\alpha_0$ and $b$, with the tier of $b$. By
  Proposition~\ref{prop:model-pass-bound}\,(ii), devices the declaration the
  run searches admits, run so that each round's acceptance value, given the
  earlier ones, follows the statistics $\Phimap[\omega]$ of some $\omega$ in
  that declaration's $\Feas$, pass the acceptance test with probability at
  most $2^{-x}$, the \emph{cap}. Where $2^{-x}\le\tau$, that is
  $x\ge\log_2\frac1\tau$, the run is refused as a whole with
  \iface{acceptance_set_unreachable}: it is evaluated at no order and
  reports no key length. The engine decides $x\ge\log_2\frac1\tau$ exactly
  where $\tau=2^{-k}$ for an integer $k\ge0$, by comparing $x$ with $k$;
  for any other $\tau$ it forms a rigorous enclosure
  $[\lambda_-,\lambda_+]\ni\log_2\frac1\tau$ and refuses where
  $x\ge\lambda_+$, refuses nothing where $x<\lambda_-$, and refuses where
  $\lambda_-\le x<\lambda_+$, the comparison being undecided there. The
  decision carries the tier of $b$: where $b$ is a first-order number not
  certified, the refusal or its absence is of the first-order tier. Where
  $b\le0$ or no bound was formed there is no cap, and this item refuses
  nothing. With $\tau=1$ every run with $b>0$ is refused. Where a reference
  companion carries the key length, the cap is still that of the
  declaration the run searches, which no claim relates to the companion's
  feasible set; the refusal only withholds a key.
\item \emph{The record.} On a checked run the engine records under
  \iface{statistics_consistency}: the declaration the solve ran on;
  $\alpha_0$ and which rule of (i) fixed it; $b$ with its tier and the lower bound
  $\frac{\alpha_0-1}{\alpha_0}\,b$ on $\sfD_{\min}$, or the reason none was
  formed; $\tau$; where $b>0$, the exponent $x$ of (iv) with the tier of
  $b$, rounded downward where the record cannot carry it exactly, so that
  $2^{-x}$ at the recorded $x$ stays an upper bound on the pass
  probability, and whether the comparison of (iv) was decided exactly, by
  the enclosure, or left undecided and refused; where $b\le0$ or no bound
  was formed, that there is no cap; whether (iv) refused the run, as (iv)
  decided it on the exact $x$ and never as a comparison of the recorded or
  displayed values with $\tau$; where the
  solve ran, the share of the wall budget and the wall
  budget it was given, the wall time it took and the cause it stopped on, as the solver
  reports it (such as the gap target, the iteration cap or the budget);
  where the solve returned a point, the
  estimate $O_0\coloneqq O^{\mathcal{D}_{\mathrm{all}}}_{\alpha_0}$ at that
  point, whether it is negative, and whether its membership in the relaxed set was established;
  where $b$ was formed, the gap $O_0-b$; and, where a bound was formed,
  $n\,f(\alpha')$ in bits at each order $\alpha'$ the run evaluates, a
  diagnostic of the part of the value the mismatch accounts for, which
  satisfies $\frac{\alpha'-1}{\alpha'}\,n\,f(\alpha')
  =n\,\frac{\alpha_0-1}{\alpha_0}\,b$
  (Proposition~\ref{prop:model-pass-bound}\,(ii)). The report shows, beside
  every key length it reports and, on a run refused under (iv), in place of
  one, the cap as a probability: $2^{-x}$ at the recorded $x$ together with
  a decimal approximation of it rounded upward, with the tier of $b$, as
  the probability with which devices the declaration the run searches
  admits, run as (iv) says, pass the acceptance test, both labelled as
  rounded upward; and whether (iv) refused the run, with $\tau$, as the
  record carries that decision from the exact comparison, so that a
  displayed cap a little above $\tau$ beside a refusal reads as the
  upward rounding it is. Where there is no cap it says so and why
  ($b\le0$, or no bound formed and the reason). It states the cap as a
  probability and not as a number of bits, and it does not show
  $n\,f(\alpha')$. No key length and no other number of the run reads
  these entries.
\end{enumerate}
\end{requirement}
\claimreading{If you supply the honest statistics or declare a constraint, the engine checks, unless switched off, whether a state the declaration admits has statistics the test accepts. Once per run, before any key length is computed and with a bounded share of one evaluation's time, it bounds the mismatch from below; a short solve can leave a mismatch unproven. From a positive bound it forms a cap on the probability that devices behaving as the declaration says pass your acceptance test, proven where the bound is certified, and the record and report show that cap as a probability. Where the cap is at most one half (adjustable), the model's own devices would fail the test at least as often as pass it, and the run is refused with no key length. The record says whether the run was checked.}
\claimmeta{\metalabel{Motivated by}\mref{Prop.}{prop:value-range}, \mref{Def.}{def:dropped}, \mref{Prop.}{prop:minimal-divergence}, \mref{Prop.}{prop:model-pass-bound}, \mref{Req.}{req:infinite-value}, \mref{Req.}{req:objective-alpha-domain}, \mref{Req.}{req:key-length-alpha-domain}, \mref{Prop.}{prop:relaxed-set-properties}, \mref{Cor.}{cor:fo-dropped}, \mref{Lem.}{lem:certified-bracket}, \mref{Req.}{req:negative-estimate-refused}, \mref{Req.}{req:certified-evaluation}, \extref{../foundations/contract.pdf}{foundations/def:tiers}, \extref{../statistics/contract.pdf}{statistics/ass:reach}
\metasep implemented (effective: accepted)
\metasep \metalabel{Code}\coderef{https://github.com/fraunhofersg/certiqs-key/blob/v0.1.0/src/certiqs_key/meat_core/finitekey/consistency.py}{certiqs_key.meat_core.finitekey.consistency:decide_pass_cap}}

% claim meat-core/req:objective-alpha-domain
\claimsummary{meat-core/req:objective-alpha-domain}{The engine evaluates the objective, its gradient and the program only at a R\'enyi order \ensuremath{\alpha} in (1, \ensuremath{\infty}), where the objective is defined, and refuses any other order.}
\claimbookmark{meat-core/req:objective-alpha-domain}{Requirement}{req:objective-alpha-domain}{Rényi order of an evaluated objective}
\claimleadin{%
The \termref{term-objective}{objective} \eqref{eq:objective} is defined for
$\alpha\in(1,\infty)$ (Definition~\ref{def:feasible}), the range on which
\xref{foundations/def:renyi} defines the sandwiched entropies and on which
the entropy, program and gradient results are stated
(Remark~\ref{rem:alpha-domain}).
}
\begin{technicalclaim}
\begin{requirement}[Rényi order of an evaluated objective]\label{req:objective-alpha-domain}
The engine evaluates the objective, its gradient and the program only at a
R\'enyi order $\alpha\in(1,\infty)$, and refuses any other order
(\iface{alpha_in_the_profile_domain}).
\end{requirement}
\end{technicalclaim}
\claimmeta{\metalabel{Motivated by}\mref{Def.}{def:feasible}, \mref{Rem.}{rem:alpha-domain}, \extref{../foundations/contract.pdf}{foundations/def:renyi}
\metasep \metalabel{Used by}\mref{Req.}{req:unreachable-acceptance-set}
\metasep implemented (effective: accepted)
\metasep \metalabel{Code}\coderef{https://github.com/fraunhofersg/certiqs-key/blob/v0.1.0/src/certiqs_key/meat_core/objective/objective.py}{certiqs_key.meat_core.objective.objective:check_alpha_structural}}

% claim meat-core/req:reported-key-length
\claimsummary{meat-core/req:reported-key-length}{The engine reports the maximal key length built from the lower bound h on h\_low it uses, the right-hand side of the key-length bound at max\{0, h\} floored and clamped at 0, as its conservative integer output.}
\claimbookmark{meat-core/req:reported-key-length}{Requirement}{req:reported-key-length}{Reported key length}
\claimleadin{%
The engine never has the exact $\hlow$, only a lower bound $h$ on it that a
solver establishes, and Corollary~\ref{cor:keylength-bound} is where it
obtains its key length from one. The key length is a non-negative integer, both steps that
make it one, the floor and the clamp at $0$, go in the safe direction
(Definition~\ref{def:max-key-length}), and, when $h$ depends on none of
$(\lEC,\eev,\epa)$, Corollary~\ref{cor:keylength-bound}\,(a) carries the
monotonicity of Proposition~\ref{prop:monotone} over to it.
}
\begin{requirement}[Reported key length]\label{req:reported-key-length}
The engine reports as \iface{key_length} the \termref{term-maximal-key-length}{maximal key length}
built from the lower bound $h\le\hlow$ it uses, that is $\ell_{\max}$ of
Definition~\ref{def:max-key-length} with $\max\{0,h\}$ in place of $\hlow$
(Corollary~\ref{cor:keylength-bound}\,(a)), with its floor and its clamp at
$0$: that is its conservative integer-output rule.
\end{requirement}
\claimmeta{\metalabel{Motivated by}\mref{Thm.}{thm:keylength}, \mref{Def.}{def:max-key-length}, \mref{Prop.}{prop:monotone}, \mref{Cor.}{cor:keylength-bound}
\metasep \metalabel{Used by}\mref{Req.}{req:certified-evaluation}
\metasep accepted
\metasep \metalabel{Code}\coderef{https://github.com/fraunhofersg/certiqs-key/blob/v0.1.0/src/certiqs_key/meat_core/finitekey/keylength.py}{certiqs_key.meat_core.finitekey.keylength:key_length}}

% claim meat-core/req:hash-penalty
\claimsummary{meat-core/req:hash-penalty}{For a privacy-amplification family that is only two-universal, collision probability at most 2\textasciicircum{}\ensuremath{-}\ensuremath{\ell}, the key-length bound is tightened by one further bit, \ensuremath{\ell} \ensuremath{\to} \ensuremath{\ell} \ensuremath{-} 1.}
\claimbookmark{meat-core/req:hash-penalty}{Requirement}{req:hash-penalty}{Key-length penalty for a merely two-universal family}
\claimleadin{%
The leftover-hashing step of Theorem~\ref{thm:keylength} is taken under a
reading that requires an ideal two-universal family; for a family that is only two-universal that
requirement can be relaxed at the cost of one bit of key length
\cite[Rem.~6.1]{tupkary2026decoy} (Remark~\ref{rem:hash-penalty}).
}
\begin{requirement}[Key-length penalty for a merely two-universal family]\label{req:hash-penalty}
For a privacy-amplification family that is only two-universal, collision
probability at most $2^{-\ell}$, \eqref{eq:keylength} must be tightened by
one further bit, $\ell\to\ell-1$.
\end{requirement}
\claimmeta{\metalabel{Motivated by}\mref{Thm.}{thm:keylength}, \mref{Rem.}{rem:hash-penalty}
\metasep accepted
\metasep \metalabel{Code}\coderef{https://github.com/fraunhofersg/certiqs-key/blob/v0.1.0/src/certiqs_key/meat_core/finitekey/keylength.py}{certiqs_key.meat_core.finitekey.keylength:key_length}}

% claim meat-core/req:key-length-alpha-domain
\claimsummary{meat-core/req:key-length-alpha-domain}{The engine reports a key length only at a R\'enyi order \ensuremath{\alpha} in (1, 2), the order at which every key-length statement holds, and reports none at any other order.}
\claimbookmark{meat-core/req:key-length-alpha-domain}{Requirement}{req:key-length-alpha-domain}{Rényi order of a reported key length}
\claimleadin{%
Theorem~\ref{thm:keylength}, Corollary~\ref{cor:keylength-bound} and
Corollary~\ref{cor:fo-keylength} give a key length only for
$\alpha\in(1,2)$, the domain of the R\'enyi leftover-hashing lemma
(Remark~\ref{rem:alpha-domain}).
}
\begin{requirement}[Rényi order of a reported key length]\label{req:key-length-alpha-domain}
The engine reports a key length only at a R\'enyi order $\alpha\in(1,2)$;
at any other order it reports no key length
(\iface{alpha_in_the_profile_domain}).
\end{requirement}
\claimmeta{\metalabel{Motivated by}\mref{Thm.}{thm:keylength}, \mref{Rem.}{rem:alpha-domain}, \mref{Cor.}{cor:keylength-bound}, \mref{Cor.}{cor:fo-keylength}
\metasep implemented (effective: accepted)
\metasep \metalabel{Code}\coderef{https://github.com/fraunhofersg/certiqs-key/blob/v0.1.0/src/certiqs_key/meat_core/finitekey/keylength.py}{certiqs_key.meat_core.finitekey.keylength:check_alpha_key_length}}

% claim meat-core/req:compiled-declaration-only
\claimsummary{meat-core/req:compiled-declaration-only}{The engine reports a key length only on a declaration that the protocol compiler emitted, and none on a declaration from anywhere else. Every condition of admissibility is then decided on every route that reports a key length.}
\claimbookmark{meat-core/req:compiled-declaration-only}{Requirement}{req:compiled-declaration-only}{Key lengths only on declarations the compiler emitted}
\claimleadin{%
The key-length results hold of an \termref{term-admissible-declaration}{admissible declaration}
(Definition~\ref{def:admissible-declaration}) that describes the protocol
being run. On a declaration that the compiler did not emit and whose
constraint list is non-empty, no code path of this module decides
condition~(vii), the non-empty \exttermref{../protocol-compiler/contract.pdf}{term-feasible-set}{feasible set}, exactly, and nothing in this
module relates such a declaration to a specification of the protocol. On a
declaration the compiler emitted, the compiler has decided condition~(vii)
(\xref{protocol-compiler/thm:feasibility}), and it has filled the declaration
from a specification of the protocol (\xref{protocol-compiler/def:compile}) whose
agreement with the apparatus is what the compiler's own assumptions state.
Refusing a key length elsewhere therefore leaves no condition of
admissibility undecided, and needs no premise that a declaration from
elsewhere describes the protocol. Numerical use of any declaration outside
the key length stays open.
}
\begin{requirement}[Key lengths only on declarations the compiler emitted]\label{req:compiled-declaration-only}
The engine reports a key length, whether as a key-length claim of
Theorem~\ref{thm:keylength}, Corollary~\ref{cor:keylength-bound} or
Corollary~\ref{cor:fo-keylength} or as a numeric key length, only on a
declaration that the protocol compiler emitted; on any other declaration it
reports none.
\end{requirement}
\claimmeta{\metalabel{Motivated by}\mref{Def.}{def:admissible-declaration}, \mref{Thm.}{thm:keylength}, \mref{Cor.}{cor:keylength-bound}, \mref{Cor.}{cor:fo-keylength}, \extref{../protocol-compiler/contract.pdf}{protocol-compiler/def:compile}, \extref{../protocol-compiler/contract.pdf}{protocol-compiler/thm:feasibility}
\metasep implemented (effective: accepted)
\metasep \metalabel{Code}\coderef{https://github.com/fraunhofersg/certiqs-key/blob/v0.1.0/src/certiqs_key/meat_core/finitekey/keylength.py}{certiqs_key.meat_core.finitekey.keylength:declaration_provenance_failure}}

% claim meat-core/req:security-statement-not-vacuous
\claimsummary{meat-core/req:security-statement-not-vacuous}{The engine refuses a run whose budgets satisfy \ensuremath{\varepsilon}\_PA + \ensuremath{\varepsilon}\_EV + 2\ensuremath{\varepsilon}\_auth \ensuremath{\ge} 1, deciding the sum exactly on the three declared numbers before any program is built. The security statement of \hyperlink{claim-meat-core/cor:composed-security}{Corollary~\ref*{cor:composed-security}} is then empty, and \hyperlink{claim-meat-core/cor:keylength-bound}{Corollary~\ref*{cor:keylength-bound}} (b) does not claim the corner.}
\claimbookmark{meat-core/req:security-statement-not-vacuous}{Requirement}{req:security-statement-not-vacuous}{Refusal of budgets with an empty security statement}
\claimleadin{%
A security statement with parameter at least $1$ excludes nothing: every
protocol meets it. The proof of Corollary~\ref{cor:keylength-bound}\,(b)
exhibits such a corner inside the hypotheses of
Theorem~\ref{thm:keylength}, $\ell=1$ at $\epa=\eev=1$ and $\lEC=0$. The
parameter a run reports is the composed one of
Corollary~\ref{cor:composed-security}\,(iii), which at the default
$\varepsilon_{\mathrm{auth}}=0$ of Assumption~\ref{ass:auth} is equal to
$\epa+\eev$, the parameter of (ii) (and numerically of (i)).
}
\begin{requirement}[Refusal of budgets with an empty security statement]\label{req:security-statement-not-vacuous}
The engine refuses a run whose budgets satisfy
$\epa+\eev+2\varepsilon_{\mathrm{auth}}\ge1$, deciding the sum exactly on the
three declared numbers before any program is built, and reports
\iface{security_statement_vacuous}: the
$(\epa+\eev+2\varepsilon_{\mathrm{auth}})$-security statement of
Corollary~\ref{cor:composed-security}\,(iii), at
$\varepsilon_{\mathrm{auth}}=0$ equal to $\epa+\eev$, the parameter of (ii)
(and numerically of (i)), is then empty, and Corollary~\ref{cor:keylength-bound}\,(b) does
not claim the corner.
\end{requirement}
\claimmeta{\metalabel{Motivated by}\mref{Thm.}{thm:keylength}, \mref{Cor.}{cor:composed-security}, \mref{Cor.}{cor:keylength-bound}, \mref{Ass.}{ass:auth}
\metasep implemented (effective: accepted)
\metasep \metalabel{Code}\coderef{https://github.com/fraunhofersg/certiqs-key/blob/v0.1.0/src/certiqs_key/meat_core/finitekey/keylength.py}{certiqs_key.meat_core.finitekey.keylength:check_security_statement}}

% claim meat-core/req:secrecy-convention
\claimsummary{meat-core/req:secrecy-convention}{An instance that adds back the one bit the key-length bound spends adopts the secrecy definition conditioned on the acceptance test alone, and declares explicitly that it does so. Its key length carries no claim of this module.}
\claimbookmark{meat-core/req:secrecy-convention}{Requirement}{req:secrecy-convention}{Declared choice of the source's secrecy convention}
\claimleadin{%
The constant of \eqref{eq:keylength} spends one bit to reach the accepted
state; the secrecy convention of the derivation it generalises does not
(Remark~\ref{rem:plus2}).
}
\begin{requirement}[Declared choice of the source's secrecy convention]\label{req:secrecy-convention}
An \exttermref{../foundations/contract.pdf}{term-instance}{instance} that adds the one bit of \eqref{eq:keylength} back adopts
the secrecy definition of the $\wedge\OmAT$ state, the convention of the
derivation Remark~\ref{rem:plus2} compares, and declares explicitly that it
does so. A key length so obtained is not one of
Theorem~\ref{thm:keylength} and carries no claim of this module.
\end{requirement}
\claimmeta{\metalabel{Motivated by}\mref{Thm.}{thm:keylength}, \mref{Rem.}{rem:plus2}
\metasep accepted}

% claim meat-core/req:constrained-key-gate
\claimsummary{meat-core/req:constrained-key-gate}{A key length is emitted with a non-empty constraint list only if the list is licensed for the declared source mode and the attack-set licence is present. Otherwise the request is refused, and the numbers serve only for exploration.}
\claimbookmark{meat-core/req:constrained-key-gate}{Requirement}{req:constrained-key-gate}{Two-part gate on a key length with declared constraints}
\begin{requirement}[Two-part gate on a key length with declared constraints]\label{req:constrained-key-gate}
A key-length output may be combined with $\CList\neq\varnothing$
\emph{only if} (i) the list lies in the licensed class of the declared
source mode (Definition~\ref{def:homog}), a condition that is mechanically
checkable and degenerate in the cases of Remark~\ref{rem:cgate}, and
(ii) the declaration carries the attack-set licence
\iface{attack_set_license} (Assumption~A-ATK,
\xref{protocol-compiler/ass:atk}), an attestation no program can verify.
The engine refuses any request that couples $\CList\neq\varnothing$ with a
key-length output unless \emph{both} parts are present; without them the
numerical machinery may be used on $\CList$ for exploration, but its output
does not reach \eqref{eq:keylength}.
\end{requirement}
\claimmeta{\metalabel{Motivated by}\mref{Def.}{def:homog}, \mref{Thm.}{thm:keylength}, \mref{Rem.}{rem:cgate}, \extref{../protocol-compiler/contract.pdf}{protocol-compiler/ass:atk}
\metasep \metalabel{Used by}\mref{Req.}{req:degenerate-gate-report}
\metasep implemented (effective: accepted)
\metasep \metalabel{Code}\coderef{https://github.com/fraunhofersg/certiqs-key/blob/v0.1.0/src/certiqs_key/meat_core/finitekey/keylength.py}{certiqs_key.meat_core.finitekey.keylength:check_key_length_permitted}}

% claim meat-core/req:degenerate-gate-report
\claimsummary{meat-core/req:degenerate-gate-report}{Where the licensed-class half of the constraint gate is degenerate, at d\_A = 1 and in untrusted-source mode, the engine never reports its pass as evidence of anything, nor as evidence about a party-local entry. The key length there rests on the attack-set licence alone.}
\claimbookmark{meat-core/req:degenerate-gate-report}{Requirement}{req:degenerate-gate-report}{A degenerate licensed-class check reported as no evidence}
\claimleadin{%
In the degenerate cases of Remark~\ref{rem:cgate} the constrained key length
rests entirely on the attestation, so a passing part~(i) there carries no
information.
}
\begin{requirement}[A degenerate licensed-class check reported as no evidence]\label{req:degenerate-gate-report}
Where part~(i) of the two-part gate of
Requirement~\ref{req:constrained-key-gate} is degenerate, at $d_A=1$ and in
\exttermref{../transmitter/contract.pdf}{term-untrusted-source-mode}{untrusted-source mode} (Remark~\ref{rem:cgate}), the engine never
reports a passing part~(i) as evidence of anything, and it never reports a
passing part~(i) as evidence of anything about a party-local entry either.
\end{requirement}
\claimmeta{\metalabel{Motivated by}\mref{Rem.}{rem:cgate}, \mref{Req.}{req:constrained-key-gate}
\metasep implemented (effective: accepted)
\metasep \metalabel{Code}\coderef{https://github.com/fraunhofersg/certiqs-key/blob/v0.1.0/src/certiqs_key/meat_core/finitekey/keylength.py}{certiqs_key.meat_core.finitekey.keylength:constraint_gate_report}}

% claim meat-core/req:marginal-licence-gate
\claimsummary{meat-core/req:marginal-licence-gate}{A declared trusted marginal without the source-replacement licence is refused before any key-length output. At a one-dimensional Alice system the licence is vacuous and is not demanded.}
\claimbookmark{meat-core/req:marginal-licence-gate}{Requirement}{req:marginal-licence-gate}{Source licence for a declared marginal}
\begin{requirement}[Source licence for a declared marginal]\label{req:marginal-licence-gate}
The refusal discipline of the two-part gate of
Requirement~\ref{req:constrained-key-gate} governs the marginal as well: the engine refuses a declared $\tau_A$
without \iface{source_replacement_license} (Assumption~A-SRC,
\xref{transmitter/ass:src}) before any key-length output, and at $d_A=1$,
where the licence is vacuous, it does not demand it
(Proposition~\ref{prop:trivial-alice-system}\,(iv)).
\end{requirement}
\claimmeta{\metalabel{Motivated by}\mref{Prop.}{prop:trivial-alice-system}, \mref{Req.}{req:constrained-key-gate}, \extref{../transmitter/contract.pdf}{transmitter/ass:src}
\metasep implemented (effective: accepted)
\metasep \metalabel{Code}\coderef{https://github.com/fraunhofersg/certiqs-key/blob/v0.1.0/src/certiqs_key/meat_core/finitekey/keylength.py}{certiqs_key.meat_core.finitekey.keylength:check_source_license}}

\clearpage
\part{Frank-Wolfe first-order bound}
\sectionintro{The Frank-Wolfe submodule supplies the first-order lower bound on the core's value h\_low, and its main results are that bound, its forms for the dropped program and for the reference companion, and the key length from its certified form. At a point of the Choi program where the relevant probabilities are positive and every key-active branch operator has maximal support, the objective is differentiable along every feasible direction, and convexity turns its linearisation into a lower bound on h\_low once lower bounds on the linear minima over the acceptance set and over the feasible set are supplied (\hyperlink{claim-meat-core/frank-wolfe/thm:firstorder}{Theorem~\ref*{thm:firstorder}}); the same holds for the dropped program, with maximal support asked only of the kept pairs (\hyperlink{claim-meat-core/frank-wolfe/cor:fo-dropped}{Corollary~\ref*{cor:fo-dropped}}). The acceptance-side minimum is computed exactly on the exact island, and the feasible-side minimum is bounded from below in exact arithmetic from any multipliers, the solver's dual among them, with the rounding of the frame the rows are built with charged, so no bound trusts a solver's optimality (\hyperlink{claim-meat-core/frank-wolfe/lem:feasible-side-dual}{Lemma~\ref*{lem:feasible-side-dual}}). Iterates need only lie in a relaxed feasible set read at declared windows, both minima being taken over the exact sets, and reading the set at a window can only lower the value. The theorem holds at any admissible points, wherever they came from, so a run that holds a reference of its snapped measurement forms the bound on the reference companion at the points its search on the snapped declaration produced and at those of a short search continued on the companion from them (\hyperlink{claim-meat-core/frank-wolfe/cor:fo-reference}{Corollary~\ref*{cor:fo-reference}}). The resulting bound feeds the core's key-length corollary; the objective at the returned point remains an estimate, and for the first-order number the objective, the gradient and the data the engine computes from the declaration are evaluated in binary64 with no rounding-error certificate. At the run's final points the bound is formed again from enclosures of the gradient alone, the objective's homogeneity making a bracket the sum of its two lower bounds, every step from the declaration's data exact or rigorously enclosed; the best certified bracket gives the key length at the certified tier, a run whose companion certifies no point can be certified on its searched declaration under the snap premise instead, unless it declares strict\_snaps, and a run that certifies no point keeps its first-order number with the reasons (\hyperlink{claim-meat-core/frank-wolfe/lem:certified-bracket}{Lemma~\ref*{lem:certified-bracket}}, \hyperlink{claim-meat-core/frank-wolfe/cor:certified-keylength}{Corollary~\ref*{cor:certified-keylength}}).}

\section{The relaxed feasible set}\label{group-meat-core/frank-wolfe/def:relaxed-feasible-set}
\sectionintro{A run establishes membership of its iterates only up to declared windows on computed residuals, so the feasible set is read at a window: the marginal or coupling condition and the declared rows are relaxed, the normalisation with them, and positivity and the block structure not at all (\hyperlink{claim-meat-core/frank-wolfe/def:relaxed-feasible-set}{Definition~\ref*{def:relaxed-feasible-set}}). Every hypothesis of this module that names the feasible set survives on the relaxed set, and reading the set at a window can only lower the value (\hyperlink{claim-meat-core/frank-wolfe/prop:relaxed-set-properties}{Proposition~\ref*{prop:relaxed-set-properties}}). What does not survive is the non-negativity of the objective, which falls below zero there by at most an explicit floor.}
\sectionimports{\extref{../foundations/contract.pdf}{foundations/lem:pinchmono}}

% claim meat-core/frank-wolfe/def:relaxed-feasible-set
\claimsummary{meat-core/frank-wolfe/def:relaxed-feasible-set}{At a window \ensuremath{\varepsilon} \ensuremath{\ge} 0 the relaxed feasible set keeps positivity and the block structure exact and relaxes the marginal or coupling condition, the declared rows and, with them, the normalisation. At \ensuremath{\varepsilon} = 0 it is the feasible set of the program, in the state form and in the Choi form.}
\claimbookmark{meat-core/frank-wolfe/def:relaxed-feasible-set}{Definition}{def:relaxed-feasible-set}{The relaxed feasible set}
\claimleadin{%
The conditions defining $\Feas$ in \eqref{eq:feas}--\eqref{eq:feas-modes}
are of two kinds. Some compare the variable against \exttermref{../foundations/contract.pdf}{term-declared-data}{declared data}:
the marginal condition against $\tau_A$, and each relation $\bowtie_k$
against $\Theta_k$ or $c_k$. The others, positive semidefiniteness, the
sector-block structure and the normalisation $\tr\omega=1$, name no
declared datum. A run establishes membership of its iterates only up to
windows on computed residuals, so the set is read at a tolerance; positive
semidefiniteness and the block structure stay exact. The normalisation is
relaxed with the condition that couples the sector blocks because it is an
arithmetic consequence of that condition, and in one mode it is that
condition: in the state form in \exttermref{../transmitter/contract.pdf}{term-trusted-source-mode}{trusted-source mode}, $\tr\omega=1$
follows from $\omega_A=\tau_A$; in \exttermref{../transmitter/contract.pdf}{term-untrusted-source-mode}{untrusted-source mode},
$\sum_s\tr\omega_s=1$ is the coupling condition itself
(Definition~\ref{def:untrusted-choi}\,(b)); in the Choi form,
$\tr[\chi_{\Qsys}(J)]=1$ follows from
$\sum_s\tr_{\Qsys_s}J_s=\idon{A_0}$ (proof of Theorem~\ref{thm:choi}).
Holding the normalisation exact beside a relaxed coupling condition would
exclude every point whose coupling residual $R$ has
$\tr[\iota^{\dagger}\tau_A\iota\,R]\ne0$
(Proposition~\ref{prop:relaxed-set-properties}\,(c$'$)): the generic point
with a nonzero residual, the residual being the drift over the convex
combinations that built the iterate, and a nonzero residual is what the
coupling check of a run accepts in the normal case. The dimension factors
$d_A$ and $d_{A_0}$ of the normalisation clauses of
Definition~\ref{def:relaxed-feasible-set} are written into them as \xref{foundations/req:norm-conversion} requires.
}
\begin{definition}[The relaxed feasible set]\label{def:relaxed-feasible-set}
For $\varepsilon\ge0$ the
\termdef{term-relaxed-feasible-set}{relaxed feasible set}
$\Feas(\varepsilon)$ is the set of sector-block-diagonal
$\omega\in\Pos{A\Qsys}$ satisfying, in trusted-source mode,
\begin{equation}\label{eq:feas-eps}
  \big\|\omega_A-\tau_A\big\|_{\max}\le\varepsilon,
  \quad
  \big|\tr\omega-1\big|\le d_A\varepsilon,
  \quad
  \begin{cases}
  \lmin\big(\Theta_k-A_k(\omega)\big)\ge-\varepsilon, & \bowtie_k\ =\ \preceq,\\
  \lmin\big(A_k(\omega)-\Theta_k\big)\ge-\varepsilon, & \bowtie_k\ =\ \succeq,\\
  \big\|A_k(\omega)-\Theta_k\big\|_{\max}\le\varepsilon, & \bowtie_k\ =\ =,
  \end{cases}
\end{equation}
where $A_k(\omega)\coloneqq\tr_{\Qsys}[(\idon{A}\otimes\hat O_k)\,\omega]$
is the left-hand side of the $k$-th declared relation of
\eqref{eq:feas-modes}, and, in untrusted-source mode, the normalisation
$\big|\sum_s\tr\omega_s-1\big|\le\varepsilon$ together with the scalar
forms
$a_k(\omega)\le c_k+\varepsilon$, $a_k(\omega)\ge c_k-\varepsilon$,
$|a_k(\omega)-c_k|\le\varepsilon$, with
$a_k(\omega)\coloneqq\tr[(\idon{A}\otimes\hat O_k)\,\omega]$ and no
marginal condition. In the Choi form $\Feas^{J}(\varepsilon)$ is the set of
$(J_s)_s$ with $J_s\in\Pos{A_0\Qsys_s}$,
$\big|\tr[\chi_{\Qsys}(J)]-1\big|\le d_{A_0}\varepsilon$,
$\big\|\sum_s\tr_{\Qsys_s}J_s-\idon{A_0}\big\|_{\max}\le\varepsilon$ and
the same relaxation of the $\bowtie_k$ rows.

The normalisation clause of each form is implied by the clause written
before it, by the marginal in trusted-source mode and by the coupling
residual in the Choi form, and is written out only so that the set can be
read without re-deriving it; Proposition~\ref{prop:relaxed-set-properties}\,(c$'$)
proves both implications and derives the two dimension factors. In
untrusted-source mode there is no clause before it: there the
normalisation \emph{is} the relaxed coupling condition
(Definition~\ref{def:untrusted-choi}\,(b)), which is why it carries no
dimension factor.

Which predicate is which is part of the definition and is never
converted: an inequality between operators is relaxed by a \emph{spectral
floor} $\lmin(\cdot)\ge-\varepsilon$, an operator equality by an
\emph{entrywise} ball $\|\cdot\|_{\max}\le\varepsilon$, a scalar equality
by an interval, and these are different quantities against the same number
(\xref{foundations/lem:floor-stricter}: the entrywise reading and the
spectral floor). The factors $d_A$ and $d_{A_0}$ are not such a
conversion: each is the dimension factor of an explicit trace bound
derived in Proposition~\ref{prop:relaxed-set-properties}\,(c$'$), written
into the clause, and neither reads an entrywise defect as a spectral one.

At $\varepsilon=0$ both sets reduce to \eqref{eq:feas} and
\eqref{eq:choi-feas} respectively: each normalisation clause is then the
exact $\tr\omega=1$, respectively $\tr[\chi_{\Qsys}(J)]=1$, which the
coupling condition already forces, so re-imposing it changes nothing. We
write $\Feas$ for $\Feas(0)$ and $\Feas^{J}$ for $\Feas^{J}(0)$
throughout.

\emph{In the general format} of Definition~\ref{def:declaration} both sets
are read on lists over the \exttermref{../protocol-compiler/contract.pdf}{term-joint-sector}{joint sectors}, with every
line of \eqref{eq:feas-general}, respectively of the general-format Choi
set of Theorem~\ref{thm:choi}, relaxed by the predicate its relation calls
for, as in \eqref{eq:feas-eps}, the \exttermref{../protocol-compiler/contract.pdf}{term-party-local}{party-local} rows included. The
marginal condition and the coupling condition are then read per Alice
sector:
$\|\sum_s\tr_{\Qsys_s}\omega^{(\mathsf t,s)}-\tau^{(\mathsf t)}_A\|_{\max}
\le\varepsilon$, respectively
$\|R_{\mathsf t}\|_{\max}\le\varepsilon$ with
$R_{\mathsf t}\coloneqq\sum_s\tr_{\Qsys_s}J_{(\mathsf t,s)}
-\idon{A_{0,\mathsf t}}$, for every $\mathsf t\in\mathbb{T}$. The residual
of the single marginal clause of \eqref{eq:feas-eps}, respectively of the
single coupling clause of the Choi form, is then the block-diagonal operator
$\omega_A-\tau_A=\bigoplus_{\mathsf t}(\sum_s\tr_{\Qsys_s}
\omega^{(\mathsf t,s)}-\tau^{(\mathsf t)}_A)$, respectively
$\bigoplus_{\mathsf t}R_{\mathsf t}$ with
$J_s\coloneqq\bigoplus_{\mathsf t}J_{(\mathsf t,s)}$, whose entrywise
maximum is the largest over $\mathsf t$ of the blocks' ones; so the clauses
per Alice sector together are that single clause. The normalisation clauses
are those of the first paragraph of this definition. At $\varepsilon=0$ the
two sets are $\Feas$ of \eqref{eq:feas-general} and the general-format
$\Feas^{J}$ of Theorem~\ref{thm:choi}.
\end{definition}
\claimmeta{\metalabel{Uses}\mref{Def.}{def:feasible}, \mref{Thm.}{thm:choi}
\metasep \metalabel{Used by}\mref{Prop.}{prop:relaxed-set-properties}, \mref{Rem.}{rem:feas-tol}, \mref{Lem.}{lem:fo-khyp}, \mref{Lem.}{lem:fo-face}, \mref{Thm.}{thm:firstorder}, \mref{Cor.}{cor:fo-dropped}, \mref{Cor.}{cor:fo-reference}, \mref{Lem.}{lem:certified-bracket}, \mref{Cor.}{cor:certified-keylength}, \mref{Rem.}{rem:fo-interface}, \mref{Rem.}{rem:fo-iterate-checks}, \mref{Req.}{req:iterate-checks}
\metasep newly derived
\metasep implemented (effective: accepted)
\metasep \metalabel{Code}\coderef{https://github.com/fraunhofersg/certiqs-key/blob/v0.1.0/src/certiqs_key/solvers/frank_wolfe/joint_driver.py}{certiqs_key.solvers.frank_wolfe.joint_driver:_check_point}}

% claim meat-core/frank-wolfe/prop:relaxed-set-properties
\claimsummary{meat-core/frank-wolfe/prop:relaxed-set-properties}{The relaxed feasible sets grow with the window, are convex and compact, and keep every hypothesis of this module that names the feasible set, positivity carrying them; the objective falls below zero there by at most an explicit floor. Reading the set at a window can only lower the value.}
\claimbookmark{meat-core/frank-wolfe/prop:relaxed-set-properties}{Proposition}{prop:relaxed-set-properties}{Properties of the relaxed feasible set}
\claimleadin{%
The relaxed feasible sets keep every hypothesis of this module that names
the feasible set, and positivity of the variable is what carries them:
(a) and (b) are the monotonicity, convexity and compactness of the
sets, (c) the facts the objective needs on them, (c$'$) the implication
behind each normalisation clause, (c$''$) what relaxing the normalisation
costs, (d), (d$'$) that reading the set at a tolerance is conservative,
and (e) that all of this holds in the general format.
The window at which a run checks the positivity of a point it consumes is
a number this module declares and records
(\xref{foundations/req:tolerance-scope}\,(ii)); item (c) \emph{Fourth}
gives the reason it reads $0$ (Requirements~\ref{req:iterate-checks}
and~\ref{req:iterate-repair}).
}
\begin{proposition}[Properties of the relaxed feasible set]\label{prop:relaxed-set-properties}
Let $\Decl$ satisfy Definition~\ref{def:admissible-declaration}, let $\varepsilon\ge0$, and
let $\Feas(\varepsilon)$ and $\Feas^{J}(\varepsilon)$ be as in
Definition~\ref{def:relaxed-feasible-set}.
\begin{enumerate}[label=(\alph*),leftmargin=2.4em]
\item \emph{(Monotone.)} For $0\le\varepsilon\le\varepsilon'$,
  $\Feas(\varepsilon)\subseteq\Feas(\varepsilon')$ and
  $\Feas^{J}(\varepsilon)\subseteq\Feas^{J}(\varepsilon')$; in particular
  $\Feas\subseteq\Feas(\varepsilon)$ and $\Feas^{J}\subseteq\Feas^{J}(\varepsilon)$
  for every $\varepsilon\ge0$.
\item \emph{(Convex and compact.)} $\Feas(\varepsilon)$ and
  $\Feas^{J}(\varepsilon)$ are convex and compact. Every
  $\omega\in\Feas(\varepsilon)$ has $\tr\omega\le1+d_A\varepsilon$ in
  trusted-source mode and $\tr\omega\le1+\varepsilon$ in untrusted-source
  mode, and every $J\in\Feas^{J}(\varepsilon)$ has
  $\sum_s\tr J_s\le d_{A_0}(1+\varepsilon)$, a bound that reads the coupling
  clause and never the normalisation clause.
\item \emph{(What positivity carries.)} $\Feas(\varepsilon)\subseteq\Pos{A\Qsys}$
  and $\chi_{\Qsys}\big(\Feas^{J}(\varepsilon)\big)\subseteq\Pos{A\Qsys}$.
  \emph{First}, the bounds $0\le K_{\hat c}(\omega)\le P_{\hat c}(\omega)$
  hold on the whole cone. \emph{Second}, $K_{\hat c}$ is concave on the
  sector-block-diagonal part of $\Pos{A\Qsys}$. \emph{Third},
  Theorem~\ref{thm:convexity} holds verbatim with $\Feas(\varepsilon)$ in
  place of $\Feas$, and $O^{J}_{\alpha}$ is convex and lower semicontinuous
  on $\acSet\times\Feas^{J}(\varepsilon)$. \emph{Fourth}, none of this
  survives a relaxation of positivity: the objective factors through
  $X\mapsto X^{p}$, which is undefined off the cone. One fact of the exact
  set does not survive: for $\varepsilon>0$ and $\omega\in\Feas(\varepsilon)$,
  $\sum_{\hat c}P_{\hat c}(\omega)=\tr\omega$ for an exactly admissible
  $\{F_v\}$, which the relaxed
  normalisation does not pin to $1$, so $\Phimap[\omega]$ need not be a
  probability vector and the range fact $\hlow\ge0$ of
  Proposition~\ref{prop:value-range}\,(i) is not available on
  $\Feas(\varepsilon)$ or on $\Feas^{J}(\varepsilon)$; it holds on the exact
  sets $\Feas$ and $\Feas^{J}$.
\item[(c$'$)] \emph{(Each normalisation clause is implied by the clause
  before it.)} Let $J=(J_s)_s$ be per-block positive semidefinite with
  $\|R\|_{\max}\le\varepsilon$, where
  $R\coloneqq\sum_s\tr_{\Qsys_s}J_s-\idon{A_0}$, and put
  $G\coloneqq\iota^{\dagger}\tau_A\iota=W_{\tau}^{\dagger}W_{\tau}$. Then
  $G$ is a state,
  \begin{equation}\label{eq:choi-trace-residual}
    \tr[\chi_{\Qsys}(J)]
    =\tr\Big[W_{\tau}^{\dagger}W_{\tau}\,
      \textstyle\sum_s\tr_{\Qsys_s}J_s\Big]
    =\tr\big[G(\idon{A_0}+R)\big]
    =1+\tr[GR],
  \end{equation}
  and $\big|\tr[\chi_{\Qsys}(J)]-1\big|\le d_{A_0}\varepsilon$; in
  particular $\tr[\chi_{\Qsys}(J)]=1$ holds at a point of nonzero residual
  $R$ exactly where $\tr[GR]=0$. In the state form in trusted-source mode,
  $\|\omega_A-\tau_A\|_{\max}\le\varepsilon$ implies
  $|\tr\omega-1|\le d_A\varepsilon$. In untrusted-source mode the
  normalisation clause is the relaxed coupling condition itself.
\item[(c$''$)] \emph{(What relaxing the normalisation costs.)} At a point
  of $\Feas^{J}(\varepsilon)$ put
  $\varsigma\coloneqq\tr[\chi_{\Qsys}(J)]-1$, so that
  $|\varsigma|\le d_{A_0}\varepsilon$ and
  $\sum_{\hat c}P^{J}_{\hat c}(J)=1+\varsigma$, and, where $1+\varsigma>0$,
  $\bar{\mathbf P}\coloneqq\mathbf P/(1+\varsigma)$, a probability vector.
  Then for every probability vector $\mathbf q$
  \begin{equation}\label{eq:penaltyshift-fw}
    \sum_{\hat c}q(\hat c)\log_2\frac{q(\hat c)}{P^{J}_{\hat c}}
    =\sum_{\hat c}q(\hat c)\log_2\frac{q(\hat c)}{(1+\varsigma)\bar P_{\hat c}}
    =\sfD(\mathbf q\|\bar{\mathbf P})-\log_2(1+\varsigma),
  \end{equation}
  and, provided $d_{A_0}\varepsilon<1$,
  \begin{equation}\label{eq:relaxed-objective-floor}
    O^{J}_{\alpha}(\mathbf q,J)\;\ge\;
    -\frac{\alpha}{\alpha-1}\,\log_2\!\big(1+d_{A_0}\varepsilon\big)
    \qquad\text{on }\acSet\times\Feas^{J}(\varepsilon),
  \end{equation}
  whose right-hand side is $0$ at $\varepsilon=0$ and negative otherwise.
  The state form carries the same statement with $d_A\varepsilon$ in
  trusted-source mode and $\varepsilon$ in untrusted-source mode in place of
  $d_{A_0}\varepsilon$. At $d_{A_0}\varepsilon\ge1$ the identity
  \eqref{eq:penaltyshift-fw} still holds wherever $1+\varsigma>0$ but
  carries no bound.
\item[(d)] \emph{(Reading the set at a tolerance is conservative.)} With
  $\hlow[\Feas(\varepsilon)]\coloneqq\inf_{\mathbf q\in\acSet}
  \inf_{\omega\in\Feas(\varepsilon)}O_{\alpha}(\mathbf q,\omega)$,
  \begin{equation}\label{eq:feas-tol-mono}
    \varepsilon\le\varepsilon'\ \Longrightarrow\
    \hlow\big[\Feas(\varepsilon')\big]\;\le\;\hlow\big[\Feas(\varepsilon)\big]
    \;\le\;\hlow\big[\Feas(0)\big]=\hlow\;\le\;\hqkd .
  \end{equation}
\item[(d$'$)] \emph{(The Choi-form quantity.)} With
  \begin{equation}\label{eq:hlow-feasJ}
    \hlow\big[\Feas^{J}(\varepsilon)\big]\coloneqq
    \inf_{\mathbf q\in\acSet}\ \inf_{J\in\Feas^{J}(\varepsilon)}
    O^{J}_{\alpha}(\mathbf q,J),
  \end{equation}
  the two relaxed quantities agree at $\varepsilon=0$, and
  \begin{equation}\label{eq:feas-tol-mono-J}
    \varepsilon\le\varepsilon'\ \Longrightarrow\
    \hlow\big[\Feas^{J}(\varepsilon')\big]\;\le\;
    \hlow\big[\Feas^{J}(\varepsilon)\big]\;\le\;
    \hlow\big[\Feas^{J}(0)\big]\;=\;\hlow\;\le\;\hqkd .
  \end{equation}
  At $\varepsilon>0$ no relation between $\hlow[\Feas^{J}(\varepsilon)]$
  and $\hlow[\Feas(\varepsilon)]$ is asserted.
\item[(e)] \emph{(The general format.)} For a declaration in the general
  format, with both sets read as Definition~\ref{def:relaxed-feasible-set}
  reads them there and $\chi_{\Qsys}$ the assignment
  $(J_{(\mathsf t,s)})\mapsto(\omega^{(\mathsf t,s)})$ of
  Theorem~\ref{thm:choi}, (a) to (d$'$) hold as stated. In (c$'$) the
  residual is $R=\bigoplus_{\mathsf t}R_{\mathsf t}$, $G=\bigoplus_{\mathsf
  t}G_{\mathsf t}$ with $G_{\mathsf t}\coloneqq W_{\mathsf t}^{\dagger}
  W_{\mathsf t}$, and
  $\tr[\chi_{\Qsys}(J)]=1+\sum_{\mathsf t}\tr[G_{\mathsf t}R_{\mathsf t}]$.
\end{enumerate}
\end{proposition}
\claimmeta{\metalabel{Uses}\mref{Def.}{def:admissible-declaration}, \mref{Lem.}{lem:branch-maps}, \mref{Thm.}{thm:branch-formula}, \mref{Lem.}{lem:cone-concavity}, \mref{Lem.}{lem:joint-sectors}, \mref{Def.}{def:feasible}, \mref{Thm.}{thm:reduction}, \mref{Cor.}{cor:reduction-constrained}, \mref{Prop.}{prop:value-range}, \mref{Thm.}{thm:convexity}, \mref{Def.}{def:support-frame}, \mref{Lem.}{lem:choi-factor}, \mref{Thm.}{thm:choi}, \mref{Def.}{def:relaxed-feasible-set}, \extref{../foundations/contract.pdf}{foundations/lem:pinchmono}
\metasep \metalabel{Used by}\mref{Rem.}{rem:feas-tol}, \mref{Thm.}{thm:firstorder}, \mref{Rem.}{rem:iterate-relaxed-oracle-exact}, \mref{Cor.}{cor:fo-dropped}, \mref{Rem.}{rem:fo-iterate-checks}, \mref{Rem.}{rem:fo-negative-values}, \mref{Req.}{req:negative-estimate-refused}, \mref{Req.}{req:iterate-checks}, \mref{Req.}{req:per-sector-coupling}
\metasep newly derived
\metasep implemented (effective: accepted)
\metasep \metalabel{Code}\coderef{https://github.com/fraunhofersg/certiqs-key/blob/v0.1.0/src/certiqs_key/solvers/frank_wolfe/joint_driver.py}{certiqs_key.solvers.frank_wolfe.joint_driver:_check_point}}

\begin{proof}
\emph{Idea.} Every relaxed condition weakens as the window grows, which
gives (a), and is an affine preimage of a closed convex set, which with
the trace bounds gives (b). The facts of (c) read only positivity
of the variable, which the relaxation keeps, and transfer to the Choi form
along the linear map $\chi_{\Qsys}$. Item (c$'$) is a trace estimate on the
coupling residual, (c$''$) a rescaling of the statistics vector, and (d),
(d$'$) compare infima over nested sets.

\emph{(a).} For $0\le\varepsilon\le\varepsilon'$ each condition of
\eqref{eq:feas-eps} at $\varepsilon$ implies the same condition at
$\varepsilon'$, and so does each condition of the untrusted-source and the
Choi forms, so $\Feas(\varepsilon)\subseteq\Feas(\varepsilon')$ and
$\Feas^{J}(\varepsilon)\subseteq\Feas^{J}(\varepsilon')$. At
$\varepsilon=0$ the sets are $\Feas$ and $\Feas^{J}$
(Definition~\ref{def:relaxed-feasible-set}), which gives the inclusions
for every $\varepsilon\ge0$.

\emph{(b).} Each relaxed condition is convex:
$\lmin(B-A(\omega))\ge-\varepsilon$ reads
$B-A(\omega)+\varepsilon\idon{}\succeq0$, an affine preimage of the
positive semidefinite cone; $\|\cdot\|_{\max}\le\varepsilon$ is an affine
preimage of a norm ball; each normalisation clause is the preimage of a
closed interval under the affine map $\omega\mapsto\tr\omega$,
respectively $J\mapsto\tr\chi_{\Qsys}(J)$, $\chi_{\Qsys}$ being linear;
and $\omega\succeq0$ and the block structure are as for $\Feas$. Each is
closed, so $\Feas(\varepsilon)$ is closed, and it is bounded because
$\omega\succeq0$ with $\tr\omega\le1+d_A\varepsilon$ in trusted-source
mode and $\tr\omega\le1+\varepsilon$ in untrusted-source mode: the relaxed
normalisation enters here as a bound on the trace in place of its exact
value $1$. $\Feas^{J}(\varepsilon)$ is closed for the same reasons and
bounded because $J_s\succeq0$ with
$\sum_s\tr J_s=\tr\big[\sum_s\tr_{\Qsys_s}J_s\big]
\le d_{A_0}\,(1+\varepsilon)$: the Choi input register $A_0$ has dimension
$d_{A_0}$ (Definition~\ref{def:support-frame}), and
$|\tr X|\le d\,\|X\|_{\max}$ for a $d\times d$ matrix, applied to
$X=\sum_s\tr_{\Qsys_s}J_s-\idon{A_0}$, gives
$|\tr[\sum_s\tr_{\Qsys_s}J_s]-d_{A_0}|\le d_{A_0}\,\varepsilon$. That
bound reads the coupling clause and never the normalisation clause.

\emph{(c), state form.} $\Feas(\varepsilon)\subseteq\Pos{A\Qsys}$ by
definition, and that inclusion is all that \emph{First} to \emph{Third}
need. \emph{First}: this is the cone clause of
Theorem~\ref{thm:branch-formula}, and its derivation is short: by
\eqref{eq:fK} and \eqref{eq:traces},
$K_{\hat c}=T_{\hat c}+\tr\Theta_{\hat c}$, both summands non-negative for
$\omega\succeq0$; and $K_{\hat c}\le P_{\hat c}$ is
\xref{foundations/lem:pinchmono} applied to $\tau=\cG_{\hat c}(\omega)$,
that lemma being stated on all of $\Pos{}$ and consuming no trace
hypothesis. \emph{Second}: Lemma~\ref{lem:cone-concavity}\,(b),\,(c),
whose hypothesis on $\{F_v\}$, sector-block-diagonal positive semidefinite
operators, is implied by Definition~\ref{def:admissible-declaration}. That lemma, and not
Corollary~\ref{cor:concavity}, is the source of the concavity here: the
corollary's domain is $\Sleq{}$, which a point of $\Feas(\varepsilon)$
need not meet once the normalisation is relaxed, and
Lemma~\ref{lem:cone-concavity} is stated for exactly this situation, its
part (c) being that substitution. \emph{Third}: with those two facts the
proof of Theorem~\ref{thm:convexity} runs verbatim with
$\Feas(\varepsilon)$ in place of $\Feas$, since it consumes affinity of
$P_{\hat c}$, non-negativity, concavity and continuity of $K_{\hat c}$, and
convexity and compactness of the feasible set, which are (b), and nothing
else. \emph{Fourth}: the objective factors through $X\mapsto X^{p}$, by
way of the functional $T$ of \eqref{eq:Tcone}, and $X\mapsto X^{p}$ is
undefined off the cone, so none of \emph{First} to \emph{Third} survives a
relaxation of positivity.

\emph{(c), Choi form.} Every $J\in\Feas^{J}(\varepsilon)$ has
$J_s\succeq0$ per block, the one condition of the set that is not relaxed,
and $\chi_{\Qsys}(J)=\bigoplus_s
(W_{\tau}\otimes\idon{})J_s(W_{\tau}^{\dagger}\otimes\idon{})$ is
sector-block-diagonal and positive semidefinite; hence
$\chi_{\Qsys}\big(\Feas^{J}(\varepsilon)\big)\subseteq\Pos{A\Qsys}$, the
one property of $\Feas(\varepsilon)$ the state-form argument uses. So the
bounds $0\le K_{\hat c}\le P_{\hat c}$, the cone concavity of
Lemma~\ref{lem:cone-concavity} and Theorem~\ref{thm:convexity} all apply
to $O_{\alpha}$ on $\acSet\times\chi_{\Qsys}(\Feas^{J}(\varepsilon))$; and
$O^{J}_{\alpha}=O_{\alpha}\circ(\idmap\times\chi_{\Qsys})$ with
$\chi_{\Qsys}$ \emph{linear}, so convexity and lower semicontinuity pull
back to the convex preimage $\acSet\times\Feas^{J}(\varepsilon)$, a
pullback along a linear map that needs no injectivity, as in Step~2 of the
proof of Theorem~\ref{thm:firstorder}. Convexity and compactness of
$\Feas^{J}(\varepsilon)$ itself are (b). The transfer uses no bijection.

\emph{(c), the range fact.} On $\Feas(\varepsilon)$ with $\varepsilon>0$,
$\sum_{\hat c}P_{\hat c}(\omega)=\tr\omega$ for an exactly admissible
$\{F_v\}$ (Definition~\ref{def:declaration}): both POVMs being complete,
the sum over $\hat c$ of \eqref{eq:Phimap} is
$\tr[(\sum_xM_x\otimes\sum_vF_v)\,\omega]=\tr\omega$, and the relaxed
normalisation no longer pins that to $1$. The proof of
Proposition~\ref{prop:value-range}\,(i) reads a KL divergence between two
probability vectors, so it is not available on the relaxed set; read at
$\varepsilon=0$, on $\Feas$ itself, it gives $\hlow\ge0$ there. In the
Choi form the range fact does not transfer either: it is read on the exact
set $\Feas^{J}(0)$, where $\Phimap$ is a probability vector.

\emph{(c$'$).} \emph{Choi form.} $R$ is Hermitian because each $J_s$ is
positive semidefinite, and $G=\iota^{\dagger}\tau_A\iota
=W_{\tau}^{\dagger}W_{\tau}$ by \eqref{eq:Wtau}. Then $G\succeq0$ and
$\tr G=\tr[\tau_A\,\iota\iota^{\dagger}]=\tr[\Pi_{\iota}\tau_A]
=\tr\tau_A=1$ by \eqref{eq:iota}, so $G$ is a state. In
\eqref{eq:choi-trace-residual} the first equality is the computation in
the proof of Theorem~\ref{thm:choi}. Hence
$\big|\tr[\chi_{\Qsys}(J)]-1\big|=|\tr[GR]|\le\|R\|\,\tr G=\|R\|
\le d_{A_0}\|R\|_{\max}\le d_{A_0}\varepsilon$, using
$|\tr[Y\varrho]|\le\|Y\|\tr\varrho$ for $\varrho\succeq0$ and the
inequality $\|X\|\le\|X\|_F\le d\,\|X\|_{\max}$ for a $d\times d$
matrix $X$, here on the $d_{A_0}\times d_{A_0}$ matrix $R$. This is the
explicit operator-norm conversion that
\xref{foundations/req:norm-conversion} asks for, with its dimension factor
written into the clause, and it is not
a conversion between the spectral and the entrywise relaxation shapes of
Definition~\ref{def:relaxed-feasible-set}: $R$ is measured in
$\|\cdot\|_{\max}$ on both sides of the derivation and no spectral floor
enters it. By \eqref{eq:choi-trace-residual}, $\tr[\chi_{\Qsys}(J)]=1$
exactly when $\tr[GR]=0$. \emph{State form, trusted-source mode.}
$|\tr\omega-1|=|\tr(\omega_A-\tau_A)|\le d_A\|\omega_A-\tau_A\|_{\max}
\le d_A\varepsilon$, by $\tr\tau_A=1$ and $|\tr X|\le d\|X\|_{\max}$ on
the $d_A\times d_A$ matrix $\omega_A-\tau_A$. \emph{Untrusted-source
mode.} There is no clause before it:
$|\sum_s\tr\omega_s-1|\le\varepsilon$ \emph{is} the relaxed coupling
condition (Definition~\ref{def:untrusted-choi}\,(b)).

\emph{(c$''$).} $|\varsigma|\le d_{A_0}\varepsilon$ is (c$'$), and
$\sum_{\hat c}P^{J}_{\hat c}(J)=1+\varsigma$ is the sum computed in the
proof of the range fact above, read at $\chi_{\Qsys}(J)$. The identity
\eqref{eq:penaltyshift-fw} reads only that
$\sum_{\hat c}P_{\hat c}=1+\varsigma$ with every $P_{\hat c}\ge0$, never
where the displacement came from; so the full-alphabet penalty sum is
$\sfD(\mathbf q\|\bar{\mathbf P})-\log_2(1+\varsigma)$ with
$\bar{\mathbf P}$ a probability vector. Writing the objective as in the
proof of Proposition~\ref{prop:value-range}\,(i),
$O^{J}_{\alpha}=\frac{\alpha}{\alpha-1}
\big[\sfD(\mathbf q\|\bar{\mathbf P})-\log_2(1+\varsigma)\big]
+\sum_{\hat c\in\Chkey}q(\hat c)e_{\hat c}$ with every $e_{\hat c}\ge0$ by
\emph{First} of (c); where some $P^{J}_{\hat c}=0$ with $q(\hat c)>0$ the
objective is $+\infty$ and the bound is trivial. So, provided
$d_{A_0}\varepsilon<1$, \eqref{eq:relaxed-objective-floor} holds. The
state form is the same computation with the factor of (c$'$) for its mode.

\emph{(d).} An infimum over a larger set is no larger, so by (a) the first
two inequalities of \eqref{eq:feas-tol-mono} hold;
$\hlow[\Feas(0)]=\hlow$ because $\Feas(0)=\Feas$; and the last step is
Theorem~\ref{thm:reduction} (Corollary~\ref{cor:reduction-constrained} for
a nonempty $\CList$).

\emph{(d$'$).} At $\varepsilon=0$ the two quantities agree by
Theorem~\ref{thm:choi}. The inclusion
$\Feas^{J}=\Feas^{J}(0)\subseteq\Feas^{J}(\varepsilon)$ of (a) is seen
directly as well: a point of $\Feas^{J}$ has
$\sum_s\tr_{\Qsys_s}J_s=\idon{A_0}$, hence coupling residual
$0\le\varepsilon$; has $\tr\chi_{\Qsys}(J)=1$ by
Theorem~\ref{thm:choi}; is per-block positive semidefinite; and satisfies
every $\bowtie_k$ row exactly, hence in its relaxed form. An infimum over a
larger set is no larger, which gives \eqref{eq:feas-tol-mono-J}, the
equality being Theorem~\ref{thm:choi} and the last step
\eqref{eq:feas-tol-mono}.

\emph{(e).} In the general format the frame is
$\iota=\bigoplus_{\mathsf t}\iota_{\mathsf t}$, which satisfies
\eqref{eq:iota}, and $W_{\tau}=\bigoplus_{\mathsf t}W_{\mathsf t}$
(Lemma~\ref{lem:choi-factor}\,(ii)); so
$G=W_{\tau}^{\dagger}W_{\tau}=\bigoplus_{\mathsf t}G_{\mathsf t}
=\iota^{\dagger}\tau_A\iota$ is a state, the coupling residual of
$J_s\coloneqq\bigoplus_{\mathsf t}J_{(\mathsf t,s)}$ is
$R=\bigoplus_{\mathsf t}R_{\mathsf t}$ on
$A_0=\bigoplus_{\mathsf t}A_{0,\mathsf t}$, and
$\|R\|_{\max}=\max_{\mathsf t}\|R_{\mathsf t}\|_{\max}$. With these
identifications each proof above applies as written. (a) and (b) read each
relaxed line as an affine preimage of a closed convex set, the
\exttermref{../protocol-compiler/contract.pdf}{term-party-local}{party-local} rows among them, and bound the trace from the coupling
clause, $\sum_{\mathsf t,s}\tr J_{(\mathsf t,s)}\le d_{A_0}(1+\varepsilon)$
with $d_{A_0}=\sum_{\mathsf t}d_{A_{0,\mathsf t}}$. (c) reads per-block
positive semidefiniteness, the linearity of $\chi_{\Qsys}$ and
Theorem~\ref{thm:convexity}, whose inputs in the general format are the
convexity of the feasible set and the concavity of $K_{\hat c}$ on the
joint-sector lists (Definition~\ref{def:feasible},
Lemma~\ref{lem:joint-sectors}). (c$'$) reads $G$, $R$ and \eqref{eq:iota}:
$\tr[\chi_{\Qsys}(J)]=\tr[G(\idon{A_0}+R)]=1+\tr[GR]
=1+\sum_{\mathsf t}\tr[G_{\mathsf t}R_{\mathsf t}]$ and
$|\tr[GR]|\le\|R\|\le d_{A_0}\|R\|_{\max}\le d_{A_0}\varepsilon$. (c$''$)
reads (c$'$) and the statistics, and (d), (d$'$) read (a),
Theorem~\ref{thm:reduction} (Corollary~\ref{cor:reduction-constrained})
and Theorem~\ref{thm:choi}, all stated for the general format.
\end{proof}

% claim meat-core/frank-wolfe/rem:feas-tol
\claimsummary{meat-core/frank-wolfe/rem:feas-tol}{A run reads the relaxed feasible set at the larger of its two residual windows; every result holds at every window, and each reading is conservative. The relaxation always has a subject, and a supplier's window, the declared tolerance and a run's windows relax the set independently.}
\claimbookmark{meat-core/frank-wolfe/rem:feas-tol}{Remark}{rem:feas-tol}{Scope of reading the feasible set at a tolerance}
\begin{remark}[Scope of reading the feasible set at a tolerance]\label{rem:feas-tol}
\emph{The window a run reads.} A run reads the relaxed feasible set at the
larger of its two declared windows on computed residuals, the window on
the coupling residual and the window on the declared-row residuals
(Requirement~\ref{req:iterate-checks}). A per-block positive semidefinite
$J$ whose coupling residual and whose $\bowtie_k$ residuals each pass
their own window lies in $\Feas^{J}(\varepsilon)$ at that larger
$\varepsilon$, by Proposition~\ref{prop:relaxed-set-properties}\,(a) from
the smaller window to the larger and by its (c$'$) for the normalisation
clause, which the coupling clause implies; so the normalisation clause
adds no check to any run. In untrusted-source mode the coupling residual
a run measures is $|\sum_s\tr\omega_s-1|$, the relaxed normalisation
clause itself (Definition~\ref{def:untrusted-choi}\,(b)). Every statement
of Proposition~\ref{prop:relaxed-set-properties} holds for every
$\varepsilon\ge0$, so no claim of this module depends on which window a
run reads, and each candidate reading is conservative by its (a) and (d).

\emph{What the conservative reading buys.} By
Proposition~\ref{prop:relaxed-set-properties}\,(d), every lower-bound
claim of this module holds verbatim with $\Feas$ read at any
$\varepsilon\ge0$: the number can only come out lower, never higher, and a
run that reads the set at any such $\varepsilon$, in particular at the
windows it declares, reports a bound on the same $\hqkd$.
Theorem~\ref{thm:keylength} and Corollary~\ref{cor:fo-keylength} inherit
this, their right-hand sides being non-decreasing in the bound they
consume. The Choi-form quantity of (d$'$) is conservative by its own
inclusion $\Feas^{J}\subseteq\Feas^{J}(\varepsilon)$, not by a bijection
with the state form.

\emph{The two relaxed quantities are not compared.} At $\varepsilon>0$ the
quantities $\hlow[\Feas^{J}(\varepsilon)]$ and $\hlow[\Feas(\varepsilon)]$
need not agree, and no relation between them is asserted:
$\chi_{\Qsys}$ carries $\Feas^{J}(\varepsilon)$ neither into nor onto
$\Feas(\varepsilon)$ in general. Forward, the coupling residual is
transported \emph{conjugated},
$\omega_A-\tau_A=W_{\tau}\big(\sum_s\tr_{\Qsys_s}J_s-\idon{A_0}\big)
W_{\tau}^{\dagger}$ by \eqref{eq:WWtau}, and
$\|W_{\tau}YW_{\tau}^{\dagger}\|_{\max}$ is not $\|Y\|_{\max}$; backward,
for a singular $\tau_A$ the set $\Feas(\varepsilon)$ contains states not
supported on $\operatorname{supp}\tau_A\otimes\Hs{\Qsys}$, which are the
image of no $J$ whatever.

\emph{The floor of (c$''$) and the window of the penalty sum.} The floor
\eqref{eq:relaxed-objective-floor} is not the two-sided window
$\log_2\frac{1}{1-d_{A_0}\varepsilon}$ that \eqref{eq:penaltyshift-fw}
yields on the penalty sum alone
($|\log_2(1+\varsigma)|\le\log_2\frac{1}{1-d_{A_0}\varepsilon}$ for
$|\varsigma|\le d_{A_0}\varepsilon<1$): that window bounds the penalty
\emph{sum} on both sides, whereas the floor bounds the objective from below
and carries the prefactor $\frac{\alpha}{\alpha-1}>1$ that the objective
attaches to that sum, applied to the one-sided shift
$\log_2(1+x)$, $x\coloneqq d_{A_0}\varepsilon$, which is smaller than
$\log_2\frac{1}{1-x}$ for $x\in(0,1)$. The floor is therefore less than
$\frac{\alpha}{\alpha-1}$ times the window, their ratio tending to that
factor as $x\to0$, and it can lie on either side of the window itself: at
$\alpha=1.1$ and $x=0.01$ the two are $0.158$ and $0.0145$, at $\alpha=10$
and $x=0.5$ they are $0.650$ and $1$. The objective at a point of the
relaxed set may be negative, by at most the floor; what a run does with a
negative estimate is Requirement~\ref{req:negative-estimate-refused}.

\emph{(e) What it is worth, and when it is worth nothing.} With
$\CList=\varnothing$ and in untrusted-source mode the one condition left
to relax is the normalisation, and
$\Feas(\varepsilon)=\{\omega\in\Pos{A\Qsys}\ \text{sector-block-diagonal}:
|\sum_s\tr\omega_s-1|\le\varepsilon\}$, which is strictly larger than
$\Feas$ for $\varepsilon>0$. The relaxation is never without a subject,
and that is why: in untrusted-source mode the normalisation is the single
unconditional constraint the feasible set carries
(Definition~\ref{def:feasible}), a run measures its residual like any
other, and a relaxation that exempted it would exempt the only condition
there is. With $\CList=\varnothing$ in \emph{trusted}-source mode the
marginal row is relaxed and the normalisation follows it by
Proposition~\ref{prop:relaxed-set-properties}\,(c$'$). Where the
relaxation has a declared-row subject its content is that a declaration
whose constraint rows or whose marginal are satisfied only to within a
positive window by any reachable state is not thereby made infeasible.

\emph{(f) Composition, and one obligation this places on a producer.} A
\exttermref{../foundations/contract.pdf}{term-supplier}{supplier} that certifies the non-emptiness of the feasible set from
a supplied witness certifies it of $\Feas(\varepsilon)$ at \emph{that
supplier's own} declared window, not of $\Feas$: under an exact reading a
binary64 witness certifies nothing, and a set read at some positive number
is what such a witness can be a witness for. That window is a declared
number of the supplier's contract, and it is neither the
\exttermref{../foundations/contract.pdf}{term-declared-tolerance}{declared tolerance} $\tnum$ nor either of the two windows at which
this module reads its own feasible set; a consumer that identifies any two
of the three has performed a conversion this module does not license, and
\xref{foundations/req:norm-conversion} says what such a conversion must
carry. What \emph{is} claimed of all of them is (a) and (d) of
Proposition~\ref{prop:relaxed-set-properties}: each is a relaxation of
$\Feas$, so each reading is conservative, and no relation between the three
sets is asserted or needed. They can disagree, in a known direction: this
module decides every predicate on declared data \emph{exactly}
(Requirement~\ref{req:exact-admission};
\xref{foundations/req:consumer-decides-exactly}), so a declaration a
supplier would admit at some tolerance of its own is refused here. That is
the safe direction and the licensed one; what it asks of a
\exttermref{../foundations/contract.pdf}{term-producer}{producer} is that it emit exactly admissible data
(\xref{foundations/req:exact-emission}), which the compiler decides on the
numbers it emits (\xref{protocol-compiler/def:exact}). The normalisation is
on neither side of that comparison: it is relaxed with the coupling
residual, at that residual's own window and with the dimension factor
Proposition~\ref{prop:relaxed-set-properties}\,(c$'$) derives, and it is
compared with $\tnum$ nowhere.
\end{remark}
\claimmeta{\metalabel{Uses}\mref{Def.}{def:admissible-declaration}, \mref{Def.}{def:relaxed-feasible-set}, \mref{Prop.}{prop:relaxed-set-properties}
\metasep newly derived
\metasep accepted}

\section{Admissible directions}\label{group-meat-core/frank-wolfe/def:first-order-setting}
\sectionintro{The first-order argument runs on the Choi program, pairs acceptance vectors over the \termref{term-acceptance-support}{acceptance support}, the outcomes that some accepted distribution charges, and names the two linear-minimisation inequalities it reads (\hyperlink{claim-meat-core/frank-wolfe/def:first-order-setting}{Definition~\ref*{def:first-order-setting}}). Positivity of the probabilities is needed only there (\hyperlink{claim-meat-core/frank-wolfe/lem:fo-khyp}{Lemma~\ref*{lem:fo-khyp}}), and at a branch operator of maximal support no Hermitian direction leaves its support, so the face derivative covers every feasible direction (\hyperlink{claim-meat-core/frank-wolfe/lem:fo-face}{Lemma~\ref*{lem:fo-face}}).}
\sectionimports{\extref{../foundations/contract.pdf}{foundations/def:kl}}

% claim meat-core/frank-wolfe/def:first-order-setting
\claimsummary{meat-core/frank-wolfe/def:first-order-setting}{The first-order argument runs on the Choi program, the state form being primary in untrusted-source mode. Acceptance vectors are paired over the acceptance support and operators by the trace, and the acceptance-side and feasible-side linear-minimisation inequalities are named.}
\claimbookmark{meat-core/frank-wolfe/def:first-order-setting}{Definition}{def:first-order-setting}{Ambient program, pairings, acceptance support and linear minimisation}
\claimleadin{%
The lemmas, the theorem and the corollaries of the first-order bound share
one ambient program, one pair of inner products and the two linear
minimisations whose minima the bound reads at each iterate, through a pair
satisfying the inequalities below or through lower bounds on the minima.
The pairing on the
acceptance side runs over the outcomes that some accepted distribution
charges; Remark~\ref{rem:pairing-on-acceptance-support} explains why that
restriction loses nothing, and Lemma~\ref{lem:fo-khyp} and
Remark~\ref{rem:acc-side} show that the acceptance-side gradient lives
on those outcomes.
}
\begin{definition}[Ambient program, pairings, acceptance support and linear minimisation]\label{def:first-order-setting}
From this definition on, the ambient program of this submodule is the
Choi program \eqref{eq:choi-program} of Theorem~\ref{thm:choi}: the
variable is $x=(\mathbf q,J)$ with $\mathbf q\in\acSet$ and
$J=(J_s)_{s\in\SecSet}\in\Feas^{J}$ of \eqref{eq:choi-feas}, and the
objective is the function $O^{J}_{\alpha}$ displayed in
\eqref{eq:choi-program}, so that
$\hlow=\inf\{O^{J}_{\alpha}(x):x\in\acSet\times\Feas^{J}\}$ by that
theorem. In \exttermref{../transmitter/contract.pdf}{term-untrusted-source-mode}{untrusted-source mode} the state form is primary: there
$J_s$, $\Feas^{J}$, $\chi_{\Qsys_s}$ and $O^{J}_{\alpha}$ denote $\omega_s$,
$\Feas$, the identity and $O_{\alpha}$
(Definition~\ref{def:untrusted-choi}), and every later statement of this
submodule is read with these substitutions, as
Corollary~\ref{cor:full-gradient} reads the gradients.

The
\termdef{term-acceptance-support}{acceptance support}
of the acceptance set $\acSet$ (Definition~\ref{def:admissible-acceptance-set}) is
\begin{equation}\label{eq:chacc-def}
  \Chacc\coloneqq\big\{\hat c\in\Chatset:\ q(\hat c)>0
  \ \text{for some }\mathbf q\in\acSet\big\}.
\end{equation}
Pairings are the real inner products
$\langle\mathbf a,\mathbf b\rangle=
\sum_{\hat c\in\Chacc}a(\hat c)b(\hat c)$ on $\mathbb{R}^{\Chacc}$, the
acceptance support being where the acceptance-side gradient lives
(Lemma~\ref{lem:fo-khyp}), and
$\langle A,B\rangle=\tr[A B]$ on Hermitian operators, the latter real and,
on a tuple $(A_s)_s$, summed over $s$.

For linear forms $\lmoq_t\in\mathbb{R}^{\Chacc}$ and
$\lmoJ_t=(\lmoJ_{t,s})_{s\in\SecSet}$, a tuple of Hermitian operators of the
block shape of $\Feas^{J}$, a pair $(v_t,S_t)$ of a vector in
$\mathbb{R}^{\Chacc}$ and such a tuple satisfies the
\termdef{term-acceptance-side-inequality}{acceptance-side inequality}
of \eqref{eq:lmo} when the first inequality below holds, and the
\termdef{term-feasible-side-inequality}{feasible-side inequality}
when the second holds:
\begin{equation}\label{eq:lmo}
  \langle\lmoq_t,v_t\rangle\le
  \min_{\mathbf q'\in\acSet}\langle\lmoq_t,\mathbf q'\rangle
  \qquad\text{and}\qquad
  \langle\lmoJ_t,S_t\rangle\le
  \min_{J'\in\Feas^{J}}\langle\lmoJ_t,J'\rangle ,
\end{equation}
the second minimum being taken over the exact $\Feas^{J}$ of
\eqref{eq:choi-feas}.
\end{definition}
\claimmeta{\metalabel{Uses}\mref{Thm.}{thm:choi}, \mref{Def.}{def:untrusted-choi}
\metasep \metalabel{Used by}\mref{Lem.}{lem:fo-khyp}, \mref{Lem.}{lem:fo-face}, \mref{Thm.}{thm:firstorder}, \mref{Rem.}{rem:iterate-relaxed-oracle-exact}, \mref{Cor.}{cor:fo-dropped}, \mref{Cor.}{cor:fo-reference}, \mref{Rem.}{rem:fo-general-format}, \mref{Lem.}{lem:feasible-side-dual}, \mref{Lem.}{lem:certified-bracket}, \mref{Ass.}{ass:lmo}, \mref{Rem.}{rem:fo-interface}, \mref{Rem.}{rem:fo-iterate-checks}, \mref{Req.}{req:certified-evaluation}
\metasep newly derived
\metasep implemented (effective: accepted)
\metasep \metalabel{Code}\coderef{https://github.com/fraunhofersg/certiqs-key/blob/v0.1.0/src/certiqs_key/meat_core/objective/support.py}{certiqs_key.meat_core.objective.support:acceptance_support}}

% claim meat-core/frank-wolfe/lem:fo-khyp
\claimsummary{meat-core/frank-wolfe/lem:fo-khyp}{Every accepted distribution vanishes outside the acceptance support, so the terms of the objective indexed there vanish at every per-block positive point. The gradient's positivity hypotheses are therefore needed only on the acceptance support, and for the branch quantity only on its key-live part.}
\claimbookmark{meat-core/frank-wolfe/lem:fo-khyp}{Lemma}{lem:fo-khyp}{Positivity needed only on the acceptance support}
\claimleadin{%
Hypothesis (H-POS) of Theorem~\ref{thm:firstorder} asks for positive
probabilities only on the \termref{term-acceptance-support}{acceptance support} $\Chacc$, because that
is where the lemma below shows it is consumed; on all of $\Chatset$ it
would be unsatisfiable whenever $\Chacc\subsetneq\Chatset$, a case
Definition~\ref{def:admissible-acceptance-set} admits and Remark~\ref{rem:acc-side} handles, so
the \exttermref{../foundations/contract.pdf}{term-first-order}{first-order} tier would be unreachable for a whole class of
declarations for no mathematical reason. Part~(i) names the non-emptiness
of $\acSet$ because Definition~\ref{def:admissible-acceptance-set} admits $\varnothing$, where
the first clause has no witness to name. That case is degenerate
throughout and nothing rests on it: $\Chacc=\varnothing$ by
\eqref{eq:chacc-def}, no $x_t\in\acSet\times\Feas^{J}$ exists, so $\Itset$
is empty and no $\hlb$ is reported
(Requirement~\ref{req:first-order-withdrawal}),
$\hlow=+\infty$ as an infimum over an empty set, and the hypothesis
$\Pr[\OmAT]>0$ of Theorem~\ref{thm:accumulated-entropy} fails, so no key
length is claimed (Remark~\ref{rem:hlow-range}).
}
\begin{lemma}[Positivity needed only on the acceptance support]\label{lem:fo-khyp}
Let $\Chacc$ be the acceptance support \eqref{eq:chacc-def} of the
acceptance set $\acSet$ of Definition~\ref{def:admissible-acceptance-set}.
\begin{enumerate}[label=(\roman*),leftmargin=2.2em]
\item If $\acSet\neq\varnothing$, $\Chacc$ is the support of a single
element of $\acSet$; and, whether or not $\acSet$ is empty, every
$\mathbf q\in\acSet$ vanishes outside $\Chacc$.
\item On $\acSet\times\Feas^{J}(\varepsilon)$, for every
$\varepsilon\ge0$ and in particular on the exact
$\acSet\times\Feas^{J}$, every term of $O^{J}_{\alpha}$ indexed by
a $\hat c\notin\Chacc$ vanishes identically, so $O^{J}_{\alpha}$ restricted
to that set is a function of $(\mathbf q|_{\Chacc},J)$ alone; and for
$\mathbf q,\mathbf q'\in\acSet$ every feasible direction has
$q'(\hat c)-q(\hat c)=0$ at every $\hat c\notin\Chacc$. The only property
of the $J$ factor the first clause reads is per-block positive
semidefiniteness, so the first clause holds at every per-block positive
semidefinite point whatever, of which $\Feas^{J}(\varepsilon)$ is one
family.
\item Corollary~\ref{cor:full-gradient} holds on
$\acSet\times\Feas^{J}(\varepsilon)$, for every $\varepsilon\ge0$,
with its hypothesis ``$q(\hat c)>0$, $P^{J}_{\hat c}>0$ and
$K^{J}_{\hat c}>0$ for all $\hat c$'' weakened to: $q(\hat c)>0$ and
$P^{J}_{\hat c}>0$ for every $\hat c\in\Chacc$, and $K^{J}_{\hat c}>0$ for
every $\hat c\in\Chkey\cap\Chacc$. Then \eqref{eq:grad-q} is read as a
vector in $\mathbb{R}^{\Chacc}$ and the sums of \eqref{eq:grad-J} are
restricted to $\Chacc$, respectively to $\Chkey\cap\Chacc$; both
restrictions are identities on that set, not truncations.
\end{enumerate}
\end{lemma}
\claimmeta{\metalabel{Uses}\mref{Def.}{def:admissible-acceptance-set}, \mref{Thm.}{thm:accumulated-entropy}, \mref{Thm.}{thm:branch-formula}, \mref{Prop.}{prop:value-range}, \mref{Thm.}{thm:choi}, \mref{Cor.}{cor:full-gradient}, \mref{Def.}{def:relaxed-feasible-set}, \mref{Def.}{def:first-order-setting}, \extref{../foundations/contract.pdf}{foundations/def:kl}
\metasep \metalabel{Used by}\mref{Thm.}{thm:firstorder}, \mref{Cor.}{cor:fo-dropped}, \mref{Lem.}{lem:certified-bracket}, \mref{Rem.}{rem:acc-side}, \mref{Rem.}{rem:pairing-on-acceptance-support}, \mref{Rem.}{rem:fo-reachability}
\metasep newly derived
\metasep implemented (effective: accepted)
\metasep \metalabel{Code}\coderef{https://github.com/fraunhofersg/certiqs-key/blob/v0.1.0/src/certiqs_key/meat_core/objective/gradient.py}{certiqs_key.meat_core.objective.gradient:check_gradient_hypotheses}
\metasep \metalabel{Oracles}\coderef{https://github.com/fraunhofersg/certiqs-key/blob/v0.1.0/tests/oracles/registry.py}{registry.py}}

\begin{proof}
(i) Suppose $\acSet\neq\varnothing$. Choose one
$\mathbf q_{\hat c}\in\acSet$ witnessing each
$\hat c\in\Chacc$ and average the finitely many of them with equal
positive weights: the result lies in $\acSet$, which is convex by
Definition~\ref{def:admissible-acceptance-set}, and the support of a positively weighted convex
combination of non-negative vectors is the union of their supports, which
is $\Chacc$. (If $\Chacc=\varnothing$ the average is over the empty index
set and any element of the non-empty $\acSet$ serves, its support being
empty by \eqref{eq:chacc-def}.) No element of $\acSet$ has support outside
$\Chacc$, by \eqref{eq:chacc-def}, and that clause quantifies over a set
that may be empty, so it needs no witness.

(ii) Let $\hat c\notin\Chacc$, so $q(\hat c)=0$ for every
$\mathbf q\in\acSet$ by (i), whence also $q'(\hat c)-q(\hat c)=0$. By
the cone clause of Theorem~\ref{thm:branch-formula},
$0\le K^{J}_{\hat c}\le P^{J}_{\hat c}$ at every per-block positive
semidefinite point, hence at every point of $\Feas^{J}(\varepsilon)$ for
every $\varepsilon\ge0$, per-block positive semidefiniteness being the
one condition that set does not relax
(Definition~\ref{def:relaxed-feasible-set}); and by the convention
$0\log_2\frac0v=0$ of \xref{foundations/def:kl}, which is stated for every
$v\ge0$, the case $v=0$ included, the term
$\frac{\alpha}{\alpha-1}q(\hat c)\log_2\frac{q(\hat c)}{P^{J}_{\hat c}}$
and, when $\hat c\in\Chkey$, the pair
$q(\hat c)\log_2\frac{q(\hat c)}{P^{J}_{\hat c}}
+\frac{q(\hat c)}{\alpha-1}\log_2\frac{q(\hat c)}{K^{J}_{\hat c}}$
are $0$ at \emph{every} point of $\acSet\times\Feas^{J}(\varepsilon)$,
whatever $P^{J}_{\hat c}$ and $K^{J}_{\hat c}$ are there and in particular
when one of them vanishes. Deleting those terms changes no value of
$O^{J}_{\alpha}$ on that set. No normalisation, coupling condition or
declared row is read in this argument, which is why the statement is on
the relaxed sets and not on the exact set alone.

(iii) By (ii), on $\acSet\times\Feas^{J}(\varepsilon)$ the corollary
differentiates a sum from which the $\hat c\notin\Chacc$ terms have been
deleted, so neither its statement nor its proof evaluates $q(\hat c)$,
$P^{J}_{\hat c}$ or $K^{J}_{\hat c}$ at such a $\hat c$; and
$K^{J}_{\hat c}$ occurs in \eqref{eq:choi-program} only inside the sum
over $\hat c\in\Chkey$, so it is evaluated at no $\hat c\notin\Chkey$
either: in \eqref{eq:grad-q} the $K$-bearing branch is the case
$\hat c\in\Chkey$, and in \eqref{eq:grad-J} the $K$-bearing terms are the
sum over $\hat c\in\Chkey$. A hypothesis with no consumer in either the
statement or the proof may be dropped from both. That the surviving
gradient is the right object for every use this submodule makes of it is
the second half of (ii): a component of \eqref{eq:grad-q} indexed outside
$\Chacc$ is, in every pairing of this submodule, multiplied by a component
of a feasible direction that is $0$, so it contributes nothing and its
value, finite, undefined or $-\infty$, is never read.
\end{proof}

% claim meat-core/frank-wolfe/lem:fo-face
\claimsummary{meat-core/frank-wolfe/lem:fo-face}{The image of every Hermitian input direction under a branch map lies in the maximal branch support, so at a branch operator with maximal support no direction leaves its support. The face derivative then covers every feasible direction, and no one-sided derivative is infinite along them.}
\claimbookmark{meat-core/frank-wolfe/lem:fo-face}{Lemma}{lem:fo-face}{Every direction is admissible at maximal support}
\claimleadin{%
Theorem~\ref{thm:firstorder} pairs the gradient of the objective against
every feasible direction. The lemma below is what makes maximal support
of the branch operators enough for the linearisation to exist along every
direction the theorem pairs against: a \termref{term-branch-map}{branch map} sends no
Hermitian direction outside the \termref{term-maximal-branch-support}{maximal branch support}.
}
\begin{lemma}[Every direction is admissible at maximal support]\label{lem:fo-face}
In the setting of Definition~\ref{def:first-order-setting}, let
$\hat c\in\Chkey$ and $s\in\SecSet_K(\hat c)$, let
$J_s\in\Pos{A_0\Qsys_s}$, and put $X\coloneqq\cG^{J}_{\hat c,s}(J_s)$.
Then for \emph{every} Hermitian $\Delta$ on $\Hs{A_0\Qsys_s}$ the
operator $H\coloneqq\cG^{J}_{\hat c,s}(\Delta)$ satisfies
\begin{equation}\label{eq:fo-face}
  \operatorname{supp}H\;\subseteq\;W_{\hat c,s}
\end{equation}
with $W_{\hat c,s}$ the maximal branch support \eqref{eq:Wcs}.
Consequently, if $X$ attains maximal support,
$\operatorname{supp}X=W_{\hat c,s}$, for which $J_s\succ0$ suffices,
then $\operatorname{supp}H\subseteq\operatorname{supp}X$ for every such
$\Delta$, so that:
\begin{enumerate}[label=(\alph*),leftmargin=2.2em]
\item Theorem~\ref{thm:gradient}\,(i) applies to the pair $(X,H)$, and
  $t\mapsto T(X+tH)$ is differentiable at $t=0$, two-sidedly, the
  hypothesis being symmetric in $H\mapsto-H$;
\item the $+\infty$ case of Theorem~\ref{thm:gradient}\,(ii) is
  \emph{unreachable} from $X$ along the image of any Hermitian direction
  $\Delta$ for which $X+t H\succeq0$ for small $t>0$. That is the class
  that theorem is stated on and the class along which $T$ is defined at
  all, and it contains every direction $y-x$ that
  Theorem~\ref{thm:firstorder} pairs against, since
  $X+t\cG^{J}_{\hat c,s}(J'_s-J_s)=\cG^{J}_{\hat c,s}\big((1-t)J_s+tJ'_s\big)
  \succeq0$ for $t\in[0,1]$ whenever $J_s$ and $J'_s$ are positive
  semidefinite. No such direction strictly enlarges the support of $X$,
  because by \eqref{eq:fo-face} none leaves
  $W_{\hat c,s}=\operatorname{supp}X$.
\end{enumerate}
\end{lemma}
\claimmeta{\metalabel{Uses}\mref{Lem.}{lem:branch-maps}, \mref{Thm.}{thm:choi}, \mref{Def.}{def:branch-support}, \mref{Lem.}{lem:branch-support}, \mref{Thm.}{thm:gradient}, \mref{Cor.}{cor:full-gradient}, \mref{Def.}{def:relaxed-feasible-set}, \mref{Def.}{def:first-order-setting}
\metasep \metalabel{Used by}\mref{Thm.}{thm:firstorder}, \mref{Cor.}{cor:fo-dropped}
\metasep newly derived
\metasep implemented (effective: accepted)
\metasep \metalabel{Code}\coderef{https://github.com/fraunhofersg/certiqs-key/blob/v0.1.0/src/certiqs_key/meat_core/objective/gradient.py}{certiqs_key.meat_core.objective.gradient:z_invariant_face}
\metasep \metalabel{Oracles}\coderef{https://github.com/fraunhofersg/certiqs-key/blob/v0.1.0/tests/oracles/registry.py}{registry.py}}

\begin{proof}
Write $\Delta=\Delta_{+}-\Delta_{-}$ for the Jordan decomposition, with
$\Delta_{\pm}\succeq0$ supported on orthogonal subspaces. The map
$\cG^{J}_{\hat c,s}=\cG_{\hat c,s}\circ\chi_{\Qsys_s}$ is a composition of
two completely positive maps, $\chi_{\Qsys_s}$ being the conjugation
$\Delta\mapsto(W_{\tau}\otimes\idon{})\Delta(W_{\tau}^{\dagger}\otimes
\idon{})$ of Theorem~\ref{thm:choi} and $\cG_{\hat c,s}$ being completely
positive by Lemma~\ref{lem:branch-maps}\,(i); hence it is completely
positive and in particular positive: $\cG^{J}_{\hat c,s}(\Delta_{\pm})\succeq0$.
Each of $\chi_{\Qsys_s}(\Delta_{\pm})$ lies in $\Pos{A\Qsys_s}$ and is of
the form Lemma~\ref{lem:branch-support}\,(i) considers, so
\eqref{eq:Wcs-bound} gives
$\operatorname{supp}\cG^{J}_{\hat c,s}(\Delta_{\pm})\subseteq W_{\hat c,s}$.
For positive semidefinite $A,B$ one has
$\operatorname{supp}(A-B)\subseteq\operatorname{supp}A
+\operatorname{supp}B$, since a vector orthogonal to both supports is
annihilated by $A$ and by $B$ and hence by $A-B$, and the support of a
Hermitian operator is the orthogonal complement of its kernel. Applying
this to $A=\cG^{J}_{\hat c,s}(\Delta_{+})$,
$B=\cG^{J}_{\hat c,s}(\Delta_{-})$ gives \eqref{eq:fo-face}, the
right-hand side being a subspace and therefore closed under the sum. In
untrusted-source mode $\chi_{\Qsys_s}=\idmap$, $\Delta$ is Hermitian on
$\Hs{A\Qsys_s}$, and the same two steps apply with
$\Sadm=\Hs{A}$.

For (a): at maximal support \eqref{eq:fo-face} reads
$\operatorname{supp}H\subseteq\operatorname{supp}X$, which is the
hypothesis of Theorem~\ref{thm:gradient}\,(i) with $S=\operatorname{supp}X$;
that part concludes differentiability of $t\mapsto T(X+tH)$ at $t=0$, and
its hypothesis holds for $-H$ as well, so the derivative is two-sided.
For (b): Theorem~\ref{thm:gradient}\,(ii), whose hypothesis
$X+tH\succeq0$ for small $t>0$ is the scope clause of (b), locates the
$+\infty$ case at directions that \emph{strictly enlarge}
$\operatorname{supp}X$, and by \eqref{eq:fo-face} no image of a Hermitian
direction does. That $J_s\succ0$ suffices for maximal support is the
equality clause of Lemma~\ref{lem:branch-support}\,(i).
\end{proof}

\section{The first-order lower bound}\label{group-meat-core/frank-wolfe/thm:firstorder}
\sectionintro{The section's main results are \hyperlink{claim-meat-core/frank-wolfe/thm:firstorder}{Theorem~\ref*{thm:firstorder}}, which in exact real arithmetic bounds h\_low from below by the best linearisation over the iterates when positivity, maximal support and an unperturbed gradient hold and the two linear minima over the exact sets are bounded from below, by the pairing with a valid minimiser or by any lower bound, with iterates in the relaxed set, and \hyperlink{claim-meat-core/frank-wolfe/cor:fo-dropped}{Corollary~\ref*{cor:fo-dropped}}, the same bound for the \termref{term-dropped-program}{dropped program} with maximal support asked only of the kept key-active joint pairs. Read for the reference companion at points a search on the snapped declaration found, the theorem bounds the reference model's value (\hyperlink{claim-meat-core/frank-wolfe/cor:fo-reference}{Corollary~\ref*{cor:fo-reference}}). The key length follows from the core's corollary (\hyperlink{claim-meat-core/frank-wolfe/cor:fo-keylength}{Corollary~\ref*{cor:fo-keylength}}). The remarks show that a failure of maximal support falls on the unsafe side, that a perturbed gradient bounds another function, and that the theorem holds as written in the general format.}
\sectionimports{\extref{../foundations/contract.pdf}{foundations/def:kl}}

% claim meat-core/frank-wolfe/thm:firstorder
\claimsummary{meat-core/frank-wolfe/thm:firstorder}{In exact real arithmetic the best linearisation over the iterates, even ones feasible only up to windows, bounds h\_low from below given positive probabilities, maximal branch support, unperturbed gradients and bounds on both linear minima. An exactly feasible point's objective bounds it from above.}
\claimbookmark{meat-core/frank-wolfe/thm:firstorder}{Theorem}{thm:firstorder}{The estimate and the first-order lower bound}
\claimleadin{%
The \emph{first-order bound} $\hlb$ is the largest, over the iterates the
optimiser passed through, of the value of the linearisation of the
objective at that iterate, minimised over the feasible set. Part~(ii)
proves $\hlb\le\hlow$ in exact real arithmetic, at maximal support, with
no perturbation of the iterate and given the two inequalities of
\eqref{eq:lmo}; part~(iii) admits an iterate that lies only in
the \termref{term-relaxed-feasible-set}{relaxed feasible set}, which is where a run's iterates lie
(Remark~\ref{rem:iterate-relaxed-oracle-exact}). Part~(i) restates
Lemma~\ref{lem:estimate}\,(ii) of the core, because it is the whole
warrant of the number a run reports as its estimate, and the membership
$x\in\acSet\times\Feas^{J}$ is exactly what must be checked for that
warrant to hold (Requirement~\ref{req:iterate-checks}). That $\acSet$ is an
\termref{term-admissible-acceptance-set}{admissible acceptance set} (Definition~\ref{def:admissible-acceptance-set}) is
a hypothesis of the theorem itself and not only of the claims it consumes:
convexity of $\acSet$ is used in Lemma~\ref{lem:fo-khyp}\,(i) and
closedness in Step~2 of the proof. Part~(iv) lets a lower bound on either
minimum stand in for the pairing on its side: on the acceptance side a
certified bound (Remark~\ref{rem:acc-side}), on the feasible side the bound
of Lemma~\ref{lem:feasible-side-dual}, formed as
Requirement~\ref{req:verified-feasible-side} says. Where both stand in, the
bracket reads no point an oracle returned, and the theorem asks nothing of
an oracle. Part~(v) makes the bound independent of
condition~(vii) of an \termref{term-admissible-declaration}{admissible declaration}, the non-empty feasible set,
which the engine does not decide on every declaration: on an empty feasible
set $\hlow$ is $+\infty$, and every real number bounds it.
}
\begin{mainclaim}
\begin{theorem}[The estimate and the first-order lower bound]
\label{thm:firstorder}
Let $\alpha\in(1,\infty)$, let $\Decl$ satisfy
Definition~\ref{def:admissible-declaration}, in the \exttermref{../protocol-compiler/contract.pdf}{term-basic-format}{basic format} or in the general
format of Definition~\ref{def:declaration}, and let $\acSet$ satisfy
Definition~\ref{def:admissible-acceptance-set}. In the general format $\Feas^{J}$ is the set of
lists $(J_{(\mathsf t,s)})$ over the \exttermref{../protocol-compiler/contract.pdf}{term-joint-sector}{joint sectors} of
Theorem~\ref{thm:choi}, with $J_s\coloneqq\bigoplus_{\mathsf t}
J_{(\mathsf t,s)}$ for every sector $s$ of the measured system, and
$\Feas^{J}(\varepsilon)$ is read as
Definition~\ref{def:relaxed-feasible-set} reads it there.
\begin{enumerate}[label=(\roman*),leftmargin=2.2em]
\item \emph{(The estimate over-estimates; Lemma~\ref{lem:estimate}\,(ii)
of the core, restated.)} For every $x\in\acSet\times\Feas^{J}$,
\begin{equation}\label{eq:fo-estimate}
  O^{J}_{\alpha}(x)\;\ge\;\hlow .
\end{equation}
\item \emph{(The first-order lower bound.)} Let $\Itset$ be a finite
non-empty index set, and let
$x_t=(\mathbf q_t,J_t)\in\acSet\times\Feas^{J}$, $t\in\Itset$, satisfy
for each $t$:
\begin{itemize}[leftmargin=2.6em]
\item[(H-POS)] $q_t(\hat c)>0$ and $P^{J}_{\hat c}(J_t)>0$ for every
  $\hat c\in\Chacc$, and $K^{J}_{\hat c}(J_t)>0$ for every
  $\hat c\in\Chkey\cap\Chacc$, with $\Chacc$ the \termref{term-acceptance-support}{acceptance support}
  \eqref{eq:chacc-def};
\item[(H-SUPP)] every key-active branch operator
  $X_{\hat c,s}\coloneqq\cG^{J}_{\hat c,s}(J_{t,s})$,
  $\hat c\in\Chkey$, $s\in\SecSet_K(\hat c)$, attains maximal support:
  $\operatorname{supp}X_{\hat c,s}=W_{\hat c,s}$, the
  \termref{term-maximal-branch-support}{maximal branch support} \eqref{eq:Wcs}, equivalently
  $\operatorname{rank}X_{\hat c,s}=\dim W_{\hat c,s}
  =\operatorname{rank}\mathbf K_{\hat c,s}$ by \eqref{eq:Kstack-rank},
  the inclusion $\subseteq$ holding always by \eqref{eq:Wcs-bound};
  $J_{t,s}\succ0$ suffices for it (Lemma~\ref{lem:branch-support}\,(i));
\item[(H-NOPERT)] the pair
  $(\lmoq_t,\lmoJ_t)\coloneqq\nabla O^{J}_{\alpha}(x_t)$ is the
  split-form gradient of Theorem~\ref{thm:gradient}\,(i)(a)--(b) assembled
  by Corollary~\ref{cor:full-gradient} \emph{at $x_t$ itself}: no
  perturbation of $x_t$, of any $J_{t,s}$, or of any branch operator
  built from them is applied before the gradient is taken.
\end{itemize}
Let $(v_t,S_t)$ satisfy both inequalities of \eqref{eq:lmo} at
$(\lmoq_t,\lmoJ_t)$ (Definition~\ref{def:first-order-setting}): the
\termref{term-acceptance-side-inequality}{acceptance-side inequality} and the
\termref{term-feasible-side-inequality}{feasible-side inequality}.
Then, \emph{in exact real arithmetic},
\begin{equation}\label{eq:hlb}
  \hlb\;\coloneqq\;\max_{t\in\Itset}\Big[
  O^{J}_{\alpha}(x_t)
  +\langle\lmoq_t,v_t-\mathbf q_t\rangle
  +\langle\lmoJ_t,S_t-J_t\rangle\Big]
  \;\le\;\hlow ,
\end{equation}
and $\max\{0,\hlb\}\le\hlow$ as well.
\item \emph{(The iterate read at a tolerance, the minima over the exact
set.)} Let $\varepsilon\ge0$ and let $\Feas^{J}(\varepsilon)$ be the
relaxed feasible set of Definition~\ref{def:relaxed-feasible-set}.
\begin{enumerate}[label=(iii\alph*),leftmargin=2.4em]
\item \emph{The estimate.} For every
$x\in\acSet\times\Feas^{J}(\varepsilon)$,
$O^{J}_{\alpha}(x)\ge\hlow[\Feas^{J}(\varepsilon)]$, the Choi-form
relaxed infimum \eqref{eq:hlow-feasJ}; and
$\hlow[\Feas^{J}(\varepsilon)]\le\hlow\le\hqkd$ by
\eqref{eq:feas-tol-mono-J}.
\item \emph{The bound.} Let $x_t\in\acSet\times\Feas^{J}(\varepsilon)$,
$t\in\Itset$, satisfy (H-POS), (H-SUPP) and (H-NOPERT) as stated in (ii),
and let $(v_t,S_t)$ satisfy \eqref{eq:lmo} \emph{as written}, its
acceptance-side inequality and its feasible-side inequality, with the two
minima over $\acSet$ and over the \emph{exact} $\Feas^{J}$. Then, in exact
real arithmetic,
$\hlb\le\hlow$ and $\max\{0,\hlb\}\le\hlow$, with $\hlb$ of
\eqref{eq:hlb}.
\end{enumerate}
\item \emph{(Lower bounds in place of the pairings.)} The conclusions of
(ii) and of (iiib) hold as well when, at some or all $t\in\Itset$, the
pairing $\langle\lmoq_t,v_t\rangle$ in \eqref{eq:hlb} is replaced by a real
number $b^{\mathrm A}_t$, or the pairing $\langle\lmoJ_t,S_t\rangle$ by a
real number $b^{\mathrm F}_t$, or both, with
\begin{equation}\label{eq:lmo-bounds}
  b^{\mathrm A}_t\;\le\;\min_{\mathbf q'\in\acSet}\langle\lmoq_t,\mathbf q'\rangle
  \qquad\text{and}\qquad
  b^{\mathrm F}_t\;\le\;\min_{J'\in\Feas^{J}}\langle\lmoJ_t,J'\rangle ,
\end{equation}
each inequality asked only where its number is used, the second minimum
taken over the \emph{exact} $\Feas^{J}$, and the pairings
$\langle\lmoq_t,\mathbf q_t\rangle$ and $\langle\lmoJ_t,J_t\rangle$ being
kept. At such a $t$ the inequality of \eqref{eq:lmo} on a replaced side is
not asked of $(v_t,S_t)$, and where both sides are replaced no pair
$(v_t,S_t)$ is asked at all: the bracket of \eqref{eq:hlb} at that $t$ is
then
\begin{equation*}
  O^{J}_{\alpha}(x_t)+b^{\mathrm A}_t-\langle\lmoq_t,\mathbf q_t\rangle
  +b^{\mathrm F}_t-\langle\lmoJ_t,J_t\rangle .
\end{equation*}
\item \emph{(An empty feasible set.)} In place of the hypothesis on $\Decl$
above, let $\Decl$ be a tuple of the form of
Definition~\ref{def:declaration} that satisfies conditions (i)--(vi) and
(viii) of Definition~\ref{def:admissible-declaration}, and form $\Feas$, the
objective $O_{\alpha}$ and $\hlow$ from it by the formulas of
Definition~\ref{def:feasible}. Let $\Feas=\varnothing$. Then $\hlow=+\infty$,
so every real number is a lower bound on $\hlow$, the value of
\eqref{eq:hlb} that a run forms at its iterates included. Consequently the
conclusions $\hlb\le\hlow$ and $\max\{0,\hlb\}\le\hlow$ of (ii), (iiib) and
(iv), for a real $\hlb$, hold of every tuple of the form of
Definition~\ref{def:declaration} that satisfies conditions (i)--(vi) and
(viii) of Definition~\ref{def:admissible-declaration}, whether or not it
satisfies condition~(vii): under the hypotheses of those parts where it does,
and because $\hlow=+\infty$ where it does not.
\end{enumerate}
\end{theorem}
\end{mainclaim}
\claimreading{The headline key rate rests on this bound, which needs no exact optimum. The per-round bound h\_low is the minimum of a convex optimisation. At each step whose positivity and full-support conditions hold, the tangent plane of the objective lies below it over the feasible set, so its lowest point bounds h\_low from below; the engine reports the best such bound, h\_lb. That point is bounded without trusting a solver: one part is computed exactly, the other verified in exact arithmetic from the solver's dual against the exact constraint rows, the floating-point frame W\_t's error in those rows deducted. The engine evaluates the objective, its gradient and the maps and frame they are built from in floating point, with no error budget, so a computed h\_lb can exceed h\_low by an amount no claim bounds. Its points are feasible only within a window, so its reported estimate need not lie above h\_low.}
\claimmeta{\metalabel{Uses}\mref{Def.}{def:declaration}, \mref{Def.}{def:admissible-declaration}, \mref{Def.}{def:admissible-acceptance-set}, \mref{Thm.}{thm:branch-formula}, \mref{Lem.}{lem:joint-sectors}, \mref{Def.}{def:feasible}, \mref{Prop.}{prop:value-range}, \mref{Thm.}{thm:convexity}, \mref{Thm.}{thm:choi}, \mref{Lem.}{lem:estimate}, \mref{Def.}{def:branch-support}, \mref{Lem.}{lem:branch-support}, \mref{Thm.}{thm:gradient}, \mref{Cor.}{cor:full-gradient}, \mref{Def.}{def:relaxed-feasible-set}, \mref{Prop.}{prop:relaxed-set-properties}, \mref{Def.}{def:first-order-setting}, \mref{Lem.}{lem:fo-khyp}, \mref{Lem.}{lem:fo-face}
\metasep \metalabel{Used by}\mref{Rem.}{rem:iterate-relaxed-oracle-exact}, \mref{Rem.}{rem:fo-danger}, \mref{Rem.}{rem:fo-nopert}, \mref{Cor.}{cor:fo-keylength}, \mref{Cor.}{cor:fo-dropped}, \mref{Cor.}{cor:fo-reference}, \mref{Rem.}{rem:fo-general-format}, \mref{Lem.}{lem:certified-bracket}, \mref{Lem.}{lem:certified-bracket-enclosed}, \mref{Cor.}{cor:certified-keylength}, \mref{Cor.}{cor:certified-keylength-reference}, \mref{Rem.}{rem:acc-side}, \mref{Rem.}{rem:fo-arithmetic}, \mref{Rem.}{rem:certified-scope}, \mref{Rem.}{rem:fo-interface}, \mref{Rem.}{rem:fo-iterate-checks}, \mref{Rem.}{rem:fo-returned-point}, \mref{Rem.}{rem:fo-one-ruler}, \mref{Rem.}{rem:fo-reachability}, \mref{Rem.}{rem:fo-check-exposures}, \mref{Rem.}{rem:fo-oracle-repair}, \mref{Rem.}{rem:fo-iterate-repair}, \mref{Rem.}{rem:fo-negative-values}, \mref{Req.}{req:first-order-withdrawal}, \mref{Req.}{req:returned-point-measured}, \mref{Req.}{req:first-order-clamp}, \mref{Req.}{req:no-imported-penalty}, \mref{Req.}{req:bracket-on-reference}, \mref{Req.}{req:support-rank-check}, \mref{Req.}{req:iterate-repair}, \mref{Req.}{req:verified-feasible-side}, \mref{Req.}{req:primal-pairing}, \mref{Req.}{req:certified-evaluation}
\metasep newly derived
\metasep implemented (effective: accepted)
\metasep \metalabel{Code}\coderef{https://github.com/fraunhofersg/certiqs-key/blob/v0.1.0/src/certiqs_key/solvers/frank_wolfe/joint_driver.py}{certiqs_key.solvers.frank_wolfe.joint_driver:run_joint_frank_wolfe}
\metasep \metalabel{Oracles}\coderef{https://github.com/fraunhofersg/certiqs-key/blob/v0.1.0/tests/oracles/registry.py}{registry.py}}

\begin{proof}
\emph{Idea.} At each iterate the objective has a gradient, supplied by
Corollary~\ref{cor:full-gradient} with its positivity weakened to the
acceptance support by Lemma~\ref{lem:fo-khyp}, and at maximal support that
gradient is a genuine directional derivative in every feasible direction
(Lemma~\ref{lem:fo-face}). Convexity of the objective turns the
derivative into a linear minorant; its minimum over the product
$\acSet\times\Feas^{J}$ separates into the two linear minimisations of
\eqref{eq:lmo}; and a maximum of lower bounds is a lower bound.
In (iiib) the point lies in the relaxed set while the comparison point
lies in the exact one, and the chord argument runs in the relaxed set,
which is convex.

\emph{(i).} This is Lemma~\ref{lem:estimate}\,(ii) of the core.

\emph{(ii).} Fix $t$ and write $x\coloneqq x_t$,
$\nabla O\coloneqq(\lmoq_t,\lmoJ_t)$.

\emph{Step 1 (the gradient exists and is a genuine directional derivative
in every feasible direction).} By (H-POS) and (H-SUPP), and by
Lemma~\ref{lem:fo-khyp}\,(iii), which is what lets (H-POS) be assumed
on $\Chacc$ and its $K$-part only on $\Chkey\cap\Chacc$,
Corollary~\ref{cor:full-gradient} applies at $x$ and supplies
$\nabla O$ through \eqref{eq:grad-q}--\eqref{eq:grad-pieces}, with
$\lmoq_t\in\mathbb{R}^{\Chacc}$; by
(H-NOPERT) that is the object the bracket uses. Let
$y=(\mathbf q',J')\in\acSet\times\Feas^{J}$ and put $h\coloneqq y-x$. Its
$J$-component $\Delta_s\coloneqq J'_s-J_{t,s}$ is a difference of positive
semidefinite operators and hence Hermitian, so by
Lemma~\ref{lem:fo-face} every branch operator's direction
$\cG^{J}_{\hat c,s}(\Delta_s)$ is supported inside
$\operatorname{supp}X_{\hat c,s}$, and $t'\mapsto T(X_{\hat c,s}
+t'\cG^{J}_{\hat c,s}(\Delta_s))$ is differentiable at $t'=0$. The
remaining dependence of $O^{J}_{\alpha}$ on $x$ is through the affine maps
$\mathbf q\mapsto q(\hat c)$ and $J\mapsto P^{J}_{\hat c}(J)$ and through
finitely many compositions of $\log_2$ with quantities that (H-POS) keeps
strictly positive, all differentiable there, the sums running over
$\Chacc$, respectively $\Chkey\cap\Chacc$, by
Lemma~\ref{lem:fo-khyp}\,(ii), which deletes the remaining terms
identically on $\acSet\times\Feas^{J}$ rather than bounding them. Hence
$t'\mapsto O^{J}_{\alpha}(x+t'h)$ is differentiable at $t'=0$ with
derivative $\langle\nabla O,h\rangle$, by the chain and sum rules, exactly
as in the proof of Corollary~\ref{cor:full-gradient}, whose
maximal-support step is Lemma~\ref{lem:fo-face} applied to that same
$\Delta_s$. Theorem~\ref{thm:gradient}\,(ii) is what makes this a
\emph{complete} statement rather than an assumption: the only way a
one-sided derivative of $T$ can fail to be finite is a support-enlarging
direction, and Lemma~\ref{lem:fo-face}\,(b) says there is none.

\emph{Step 2 (convexity turns the derivative into an inequality).} By
Theorem~\ref{thm:convexity}, $O^{J}_{\alpha}$ is convex with values in
$(-\infty,+\infty]$ on the convex set $\acSet\times\Feas^{J}$:
Theorem~\ref{thm:convexity} states this for $O_{\alpha}$ on
$\Delta(\Chatset)\times\Feas$, and $O^{J}_{\alpha}=O_{\alpha}\circ
(\idmap\times\chi_{\Qsys})$ with $\chi_{\Qsys}$ \emph{linear}, so the
composition is convex on the convex preimage, a pullback along a linear
map that needs no injectivity, which is Theorem~\ref{thm:choi}'s own
convexity clause. For a convex extended-real function the difference
quotient $t'\mapsto\big(O^{J}_{\alpha}(x+t'h)-O^{J}_{\alpha}(x)\big)/t'$
is non-decreasing on $(0,1]$: this is the three-chord inequality, and it
holds with the value $+\infty$ admitted. $O^{J}_{\alpha}(x)$ is finite by
(H-POS), and $x+t'h=(1-t')x+t'y$ stays in the convex set
$\acSet\times\Feas^{J}$ for $t'\in[0,1]$, so the quotient is defined
there. Its limit as $t'\downarrow0$ is $\langle\nabla O,h\rangle$ by
Step~1, and its value at $t'=1$ is
$O^{J}_{\alpha}(y)-O^{J}_{\alpha}(x)$; monotonicity gives
$\langle\nabla O,h\rangle\le O^{J}_{\alpha}(y)-O^{J}_{\alpha}(x)$, that
is
\begin{equation}\label{eq:fo-subgrad}
  O^{J}_{\alpha}(y)\;\ge\;O^{J}_{\alpha}(x)
  +\langle\nabla O,\,y-x\rangle
  \qquad\text{for every }y\in\acSet\times\Feas^{J},
\end{equation}
including every $y$ at which $O^{J}_{\alpha}=+\infty$, where the
inequality is trivial. Taking the infimum over $y$ and using
Theorem~\ref{thm:choi},
\begin{equation}\label{eq:fo-inf}
  \hlow\;\ge\;O^{J}_{\alpha}(x)
  +\min_{y\in\acSet\times\Feas^{J}}\langle\nabla O,\,y-x\rangle ,
\end{equation}
the minimum being attained because $\acSet\times\Feas^{J}$ is compact and
the map is linear. That $\Feas^{J}$ is compact is read off
\eqref{eq:choi-feas} directly: it is closed, being an intersection of the
positive semidefinite cones with affine equalities and inequalities; and
it is bounded, because $J_s\succeq0$ with
$\sum_s\tr_{\Qsys_s}J_s=\idon{A_0}$ forces
$\sum_s\tr J_s=\tr\idon{A_0}=d_{A_0}$ and hence
$\|J_s\|\le\tr J_s\le d_{A_0}$ for every $s$. $\acSet$ is
closed in the simplex by Definition~\ref{def:admissible-acceptance-set}. In untrusted-source
mode the same reading applies to $\Feas$ with $\tr\omega=1$ in place of
the coupling condition.

\emph{Step 3 (the minimum separates over the product).} The map
$y\mapsto\langle\nabla O,y\rangle
=\langle\lmoq_t,\mathbf q'\rangle+\langle\lmoJ_t,J'\rangle$ is a sum of
one term in each factor, and $\acSet\times\Feas^{J}$ is a product set, so
\begin{equation*}
  \min_{y\in\acSet\times\Feas^{J}}\langle\nabla O,y\rangle
  =\min_{\mathbf q'\in\acSet}\langle\lmoq_t,\mathbf q'\rangle
  +\min_{J'\in\Feas^{J}}\langle\lmoJ_t,J'\rangle
  \;\ge\;\langle\lmoq_t,v_t\rangle+\langle\lmoJ_t,S_t\rangle ,
\end{equation*}
the last step being \eqref{eq:lmo}. The pair $(v_t,S_t)$ is \emph{not}
required to be feasible, only to satisfy the two inequalities.
Substituting into \eqref{eq:fo-inf},
\begin{equation*}
  \hlow\;\ge\;O^{J}_{\alpha}(x_t)
  +\langle\lmoq_t,v_t-\mathbf q_t\rangle
  +\langle\lmoJ_t,S_t-J_t\rangle .
\end{equation*}

\emph{Step 4 (a maximum of lower bounds is a lower bound).} The display
holds for each $t\in\Itset$, which is finite and non-empty, so it holds
for the largest of the right-hand sides, which is \eqref{eq:hlb}. Finally
$\hlow\ge0$ by Proposition~\ref{prop:value-range}\,(i), whose proof
consumes only Definition~\ref{def:admissible-declaration} and
Theorem~\ref{thm:branch-formula}; so $\max\{0,\hlb\}$ is the maximum of
two lower bounds on $\hlow$ and is therefore one.

\emph{(iiia).} $\hlow[\Feas^{J}(\varepsilon)]$ is by \eqref{eq:hlow-feasJ}
the infimum of $O^{J}_{\alpha}$ over $\acSet\times\Feas^{J}(\varepsilon)$
and $x$ is a point of that set, which is the whole of (i)'s argument with
the set substituted; nothing else of (i) is used, and in particular no
property of the set is. That
$\hlow[\Feas^{J}(\varepsilon)]\le\hlow\le\hqkd$ is
\eqref{eq:feas-tol-mono-J}, whose one substantive step is the inclusion
$\Feas^{J}\subseteq\Feas^{J}(\varepsilon)$ proved there.

\emph{(iiib).} The proof is (ii)'s four steps with one change, made in
Step~2: the point $x_t$ ranges over the \emph{relaxed} set while the
comparison point $y$ ranges over the \emph{exact} one. Fix
$t$ and write $x\coloneqq x_t$, $\nabla O\coloneqq(\lmoq_t,\lmoJ_t)$.

\emph{Step 1 is unchanged, and reads only per-block positivity.} By
(H-POS) and (H-SUPP) and Lemma~\ref{lem:fo-khyp}\,(iii),
Corollary~\ref{cor:full-gradient} applies at $x$ and supplies $\nabla O$.
That the corollary applies here, at a point of the relaxed set, is a
property of its hypothesis and not of this step: it asks $\mathbf q\in\acSet$ and
$J_s\in\Pos{A_0\Qsys_s}$ per block, names no normalisation, no coupling
condition and no declared row, and every point of $\Feas^{J}(\varepsilon)$
meets it, per-block positive semidefiniteness being the one condition the
relaxed set does not relax. And Lemma~\ref{lem:fo-khyp}\,(iii), the
statement that weakens the corollary's positivity clause to $\Chacc$, is
itself stated on $\acSet\times\Feas^{J}(\varepsilon)$ for every
$\varepsilon\ge0$, so the weakening reaches this point too.
By (H-NOPERT) the object so supplied is the one the bracket uses. Let
$y=(\mathbf q',J')\in\acSet\times\Feas^{J}$, the \emph{exact} set,
and put $h\coloneqq y-x$. Its $J$-component
$\Delta_s\coloneqq J'_s-J_{t,s}$ is a difference of two positive
semidefinite operators, $J'_s$ from $\Feas^{J}$ and $J_{t,s}$ from
$\Feas^{J}(\varepsilon)$ (per-block positive semidefiniteness being the
one condition the relaxed set does not relax,
Definition~\ref{def:relaxed-feasible-set}), hence Hermitian. That is all
Lemma~\ref{lem:fo-face} asks: its hypotheses are $J_s\in\Pos{A_0\Qsys_s}$
and $\Delta$ Hermitian, and it names no feasible set, no normalisation,
no coupling condition and no declared row. So every branch operator's
direction $\cG^{J}_{\hat c,s}(\Delta_s)$ is supported inside
$\operatorname{supp}X_{\hat c,s}$ and the map
$t'\mapsto T(X_{\hat c,s}+t'\cG^{J}_{\hat c,s}(\Delta_s))$ is
differentiable at $t'=0$. The remaining dependence is through the affine
$\mathbf q\mapsto q(\hat c)$ and $J\mapsto P^{J}_{\hat c}(J)$ and through
compositions of $\log_2$ with quantities (H-POS) keeps strictly positive.
Lemma~\ref{lem:fo-khyp}\,(ii) deletes the terms indexed outside $\Chacc$
identically here as well: its input is $0\le K^{J}_{\hat c}\le
P^{J}_{\hat c}$ on the set the points lie in, which
Proposition~\ref{prop:relaxed-set-properties}\,(c) establishes on the
whole cone (\emph{First}), from positivity alone and consuming no trace
hypothesis, and in the Choi form on $\Feas^{J}(\varepsilon)$. Hence
$t'\mapsto O^{J}_{\alpha}(x+t'h)$ is differentiable at $t'=0$ with
derivative $\langle\nabla O,h\rangle$.

\emph{Step 2 (the segment lies in the relaxed set, and the chord argument
runs there).} $\acSet\times\Feas^{J}(\varepsilon)$ is convex
(Proposition~\ref{prop:relaxed-set-properties}\,(b); $\acSet$ by
Definition~\ref{def:admissible-acceptance-set}), and $O^{J}_{\alpha}$ is convex on it with
values in $(-\infty,+\infty]$ by
Proposition~\ref{prop:relaxed-set-properties}\,(c), whose Choi form pulls
Theorem~\ref{thm:convexity} back along the \emph{linear} $\chi_{\Qsys}$,
needing no injectivity. Both endpoints lie in that convex set:
$x\in\acSet\times\Feas^{J}(\varepsilon)$ by hypothesis, and
$y\in\acSet\times\Feas^{J}\subseteq\acSet\times\Feas^{J}(\varepsilon)$ by
Proposition~\ref{prop:relaxed-set-properties}\,(a). So $x+t'h=(1-t')x+t'y$
stays in it for $t'\in[0,1]$, $O^{J}_{\alpha}(x)$ is finite by (H-POS),
and the three-chord inequality gives, exactly as in (ii),
$\langle\nabla O,h\rangle\le O^{J}_{\alpha}(y)-O^{J}_{\alpha}(x)$, that
is, \eqref{eq:fo-subgrad} \emph{for every $y\in\acSet\times\Feas^{J}$}.
Taking the infimum over that set, and not over the relaxed one, and using
Theorem~\ref{thm:choi}, which says the infimum of $O^{J}_{\alpha}$ over
$\acSet\times\Feas^{J}$ is $\hlow$,
\begin{equation}\label{eq:fo-inf-hybrid}
  \hlow\;\ge\;O^{J}_{\alpha}(x)
  +\min_{y\in\acSet\times\Feas^{J}}\langle\nabla O,\,y-x\rangle ,
\end{equation}
the minimum attained because $\acSet\times\Feas^{J}$ is compact and the
map linear, compactness being read off \eqref{eq:choi-feas} exactly as in
(ii), and non-emptiness following from Definition~\ref{def:admissible-declaration}\,(vii)
with Theorem~\ref{thm:choi} mapping $\Feas^{J}$ onto $\Feas$.

\emph{Steps 3 and 4 are those of (ii).} The map
$y\mapsto\langle\nabla O,y\rangle$ separates over the product
$\acSet\times\Feas^{J}$, and the two minima that appear are exactly the
two of \eqref{eq:lmo} as the hypothesis of (iiib) states them; substituting
and maximising over the finite non-empty $\Itset$ gives $\hlb\le\hlow$.
Finally $\hlow\ge0$ by Proposition~\ref{prop:value-range}\,(i), a
statement about the \emph{exact} set that is not claimed at
$\varepsilon>0$ (Proposition~\ref{prop:relaxed-set-properties}\,(c),
(c$''$)) and is not needed there, so $\max\{0,\hlb\}$ is the maximum of
two lower bounds on $\hlow$ and is one. Part~(ii) is the case
$\varepsilon=0$.

\emph{(iv).} Steps~1, 2 and~4, of (ii) and of (iiib), do not read the pair
$(v_t,S_t)$, and Step~3 reads of it only the two inequalities of
\eqref{eq:lmo}, the first of its acceptance half and the second of its
feasible half, each to bound one of the two minima into which the minimum
over the product separates. At a $t$ where a replacement is made, put
$a_t\coloneqq b^{\mathrm A}_t$ if the acceptance side is replaced and
$a_t\coloneqq\langle\lmoq_t,v_t\rangle$ otherwise, and
$f_t\coloneqq b^{\mathrm F}_t$ if the feasible side is replaced and
$f_t\coloneqq\langle\lmoJ_t,S_t\rangle$ otherwise. The separation of
Step~3, with \eqref{eq:lmo-bounds} on each replaced side and \eqref{eq:lmo}
on each side not replaced, gives
\begin{equation*}
  \min_{y\in\acSet\times\Feas^{J}}\langle\nabla O,y\rangle
  =\min_{\mathbf q'\in\acSet}\langle\lmoq_t,\mathbf q'\rangle
  +\min_{J'\in\Feas^{J}}\langle\lmoJ_t,J'\rangle
  \;\ge\;a_t+f_t ,
\end{equation*}
and substituting into \eqref{eq:fo-inf}, respectively
\eqref{eq:fo-inf-hybrid}, gives
$\hlow\ge O^{J}_{\alpha}(x_t)+a_t-\langle\lmoq_t,\mathbf q_t\rangle
+f_t-\langle\lmoJ_t,J_t\rangle$, which is the bracket of \eqref{eq:hlb}
with the replacements made. Where both sides are replaced, neither $v_t$
nor $S_t$ enters any step. Step~4 then runs unchanged.

\emph{(v).} No claim is used for $\hlow=+\infty$. By \eqref{eq:hlow-def},
$\hlow$ is the infimum of $O_{\alpha}$ over $\mathbf q\in\acSet$ and
$\omega\in\Feas$; with $\Feas=\varnothing$ the inner infimum is over the empty
set and is $+\infty$ for every $\mathbf q$, and so is the outer one, whether
$\acSet$ is empty or not. Every real number is below $+\infty$. For the
consequence: where the tuple satisfies condition~(vii) as well, it
satisfies Definition~\ref{def:admissible-declaration} in full, and with it
the two conditions Definition~\ref{def:declaration} places on its raw-key
map, which are conditions (v) and (viii); it is then a protocol declaration
satisfying the hypothesis of (ii), (iiib) and (iv); where it does not, $\Feas=\varnothing$, so
$\hlow=+\infty$ and the real number $\hlb$ is below it.

\emph{What (iiib) uses of the relaxed set.} Only these: that
$\Feas^{J}\subseteq\Feas^{J}(\varepsilon)$
(Proposition~\ref{prop:relaxed-set-properties}\,(a)); that
$\acSet\times\Feas^{J}(\varepsilon)$ is convex (\emph{ibid.}\,(b)); that
$O^{J}_{\alpha}$ is convex on it (\emph{ibid.}\,(c)); that every member is
per-block positive semidefinite (Definition~\ref{def:relaxed-feasible-set}:
the one condition not relaxed, what Lemma~\ref{lem:fo-face} consumes, and
what the hypothesis of Corollary~\ref{cor:full-gradient} asks of $x$, the
corollary being the statement that supplies $(\lmoq_t,\lmoJ_t)$ and
therefore the one whose scope decides whether \eqref{eq:hlb} names an
object at all); and that $0\le K^{J}_{\hat c}\le P^{J}_{\hat c}$ on it
(\emph{ibid.}\,(c), \emph{First}, which is the cone clause of
Theorem~\ref{thm:branch-formula} and what Lemma~\ref{lem:fo-khyp}\,(ii)
consumes). The list does not say that compactness of the relaxed set is
unused: no step of this proof reads it, the minimum being taken over the
exact set, but the convexity of $O^{J}_{\alpha}$ that Step~2 imports is
Proposition~\ref{prop:relaxed-set-properties}\,(c) \emph{Third}, whose own
derivation consumes the convexity \emph{and} the compactness of (b), the
latter through the continuity of $K_{\hat c}$; so compactness is used one
level down and not in the minimisation. Nor does it say that lower
semicontinuity is unavailable: the chord argument needs convexity alone,
but the fact imported carries both. Of the \emph{exact} set the proof uses
compactness, non-emptiness, and that its infimum is $\hlow$
(Theorem~\ref{thm:choi}). No statement is made about
$\hlow[\Feas(\varepsilon)]$ here, and none is used: the state-form relaxed
set enters this part nowhere.

\emph{The general format.} No step above reads a constraint row, the
number of \exttermref{../protocol-compiler/contract.pdf}{term-measured-factor}{measured factors} or the number of Alice
sectors, so the proof applies verbatim with $\Feas^{J}$ the joint-sector
Choi set. Step~1 reads the hypothesis of
Corollary~\ref{cor:full-gradient}, per-block positive semidefiniteness of
the point, which names no feasible set, no normalisation and no declared
row; Lemma~\ref{lem:fo-face}, whose hypotheses are
$J_s\in\Pos{A_0\Qsys_s}$ and $\Delta$ Hermitian and which reads only
\eqref{eq:Wcs-bound}, stated for the admissible Alice subspace in either
format; Theorem~\ref{thm:gradient}\,(ii), a statement about the functional
on the positive cone of a branch space; and
Lemma~\ref{lem:fo-khyp}\,(ii) and~(iii), which read only the acceptance
support, per-block positive semidefiniteness and the cone clause. Step~2
reads convexity of $\acSet\times\Feas^{J}$, convexity of the objective on
it, compactness of $\Feas^{J}$ and the identification of the infimum over
it with $\hlow$, all four supplied by the general-format paragraph of
Theorem~\ref{thm:choi}, including for the \exttermref{../protocol-compiler/contract.pdf}{term-party-local}{party-local} rows, which are
affine semidefinite conditions there like every other row; that paragraph
cites the convexity of the objective from Theorem~\ref{thm:convexity},
whose inputs, convexity of the feasible set and concavity of $K_{\hat c}$,
Definition~\ref{def:feasible} and Lemma~\ref{lem:joint-sectors} supply in
the general format. Steps~3 and~4 read that the variable is a product and
that the gradient pairs term by term over its factors, that $\Itset$ is
finite and non-empty, and, for the clamp at zero,
Proposition~\ref{prop:value-range}\,(i), which consumes only the
conventions of the kernel, the objective and the cone clause. Part~(iiib)
reads in addition (a), (b) and (c) of
Proposition~\ref{prop:relaxed-set-properties} and per-block positive
semidefiniteness of the relaxed set, which hold in the general format by
that proposition's (e); and (iv) reads nothing more.
\end{proof}

% claim meat-core/frank-wolfe/rem:iterate-relaxed-oracle-exact
\claimsummary{meat-core/frank-wolfe/rem:iterate-relaxed-oracle-exact}{A run establishes its iterate only in the relaxed feasible set, while the linear minima its bracket bounds are over the exact one, and part (iii) of \hyperlink{claim-meat-core/frank-wolfe/thm:firstorder}{Theorem~\ref*{thm:firstorder}} asks each half of the object that supplies it. Its bound compares h\_lb with the exact h\_low directly, through no relaxed quantity.}
\claimbookmark{meat-core/frank-wolfe/rem:iterate-relaxed-oracle-exact}{Remark}{rem:iterate-relaxed-oracle-exact}{Iterate in the relaxed set, minima over the exact set}
\begin{remark}[Iterate in the relaxed set, minima over the exact set]\label{rem:iterate-relaxed-oracle-exact}
Part~(iii) of Theorem~\ref{thm:firstorder} substitutes the
\termref{term-relaxed-feasible-set}{relaxed feasible set} $\Feas^{J}(\varepsilon)$ in different places
in its two halves, and the asymmetry is deliberate. A run of this module
is always in that part: it verifies the $J$ factor of its iterate's
membership at positive declared windows
(Requirement~\ref{req:iterate-checks}) and holds the $\acSet$ factor
exactly, by forming $\mathbf q_t$ on the \exttermref{../foundations/contract.pdf}{term-exact-island}{exact island}
(\xref{foundations/def:exact-island}; Remark~\ref{rem:acc-side}), so what
it establishes is $x_t\in\acSet\times\Feas^{J}(\varepsilon)$ and not exact
membership in the product; while the minima of \eqref{eq:lmo}, and the
lower bounds that part~(iv) of the theorem puts in their place, are taken
over the declared conic program \eqref{eq:choi-feas}, which is the
\emph{exact} $\Feas^{J}$. \emph{Iterate in the relaxed set, minima over
the exact set}: (iiib) asks each half of the object that supplies it. Its
hypotheses (H-POS), (H-SUPP) and (H-NOPERT) are conditions on the iterate
alone and name no set. Substituting $\Feas^{J}(\varepsilon)$ into
\eqref{eq:lmo} as well would ask for a bound against a \emph{larger} set,
hence below a \emph{smaller} minimum: since
$\Feas^{J}\subseteq\Feas^{J}(\varepsilon)$, one has
$\min_{\Feas^{J}(\varepsilon)}\le\min_{\Feas^{J}}$, so the inequality over
the exact set is the weaker of the two. It is also the one
Lemma~\ref{lem:feasible-side-dual} supplies: the trace identity that lemma
rests on is the exact coupling condition, which a point of the relaxed set
meets only within its window.

Part~(iiib) concludes against the exact $\hlow$ directly, with no
intermediate quantity: $\hlow[\Feas^{J}(\varepsilon)]$ appears in (iiia)
and nowhere in (iiib), and no non-negativity of it is claimed; it may be
negative, by at most $\frac{\alpha}{\alpha-1}\log_2(1+d_{A_0}\varepsilon)$
(Proposition~\ref{prop:relaxed-set-properties}\,(c$''$)). The clamp at
zero reads the range fact $\hlow\ge0$ of
Proposition~\ref{prop:value-range}\,(i) on the exact set, which is the only
reading it needs; on the relaxed set that fact is not available, and
nothing uses it there. Corollary~\ref{cor:fo-keylength} consumes the same
inequality on this reading as on that of~(ii), (ii) being the case
$\varepsilon=0$ of (iiib). Which $\varepsilon$ a run uses, the larger of
its two declared windows on computed residuals and never the declared
tolerance, is fixed by Requirement~\ref{req:iterate-checks}, and part~(iii)
is insensitive to it, being proved for every $\varepsilon\ge0$.

The relaxed quantity in (iiia) is the Choi-form one and not the state-form
$\hlow[\Feas(\varepsilon)]$ of
Proposition~\ref{prop:relaxed-set-properties}\,(d): at $\varepsilon>0$ the
two relaxed sets are not in bijection under $\chi_{\Qsys}$, for the two
reasons Remark~\ref{rem:feas-tol} gives, and no relation between the two
numbers is claimed or needed. Part~(iiib) uses neither relaxed value: it
compares $\hlb$ against the exact $\hlow$ directly, so the chain
\eqref{eq:feas-tol-mono-J} is consumed by the estimate half of part~(iii)
and not by the bound half.
\end{remark}
\claimmeta{\metalabel{Uses}\mref{Prop.}{prop:relaxed-set-properties}, \mref{Def.}{def:first-order-setting}, \mref{Thm.}{thm:firstorder}
\metasep newly derived
\metasep accepted}

% claim meat-core/frank-wolfe/rem:fo-danger
\claimsummary{meat-core/frank-wolfe/rem:fo-danger}{Where maximal support fails, a support-enlarging direction can give the objective a \ensuremath{-}\ensuremath{\infty} one-sided derivative, so the linearisation over-estimates the objective and the bracket can exceed h\_low. That is why (H-SUPP) is a hypothesis of \hyperlink{claim-meat-core/frank-wolfe/thm:firstorder}{Theorem~\ref*{thm:firstorder}} and a failed check drops the iterate.}
\claimbookmark{meat-core/frank-wolfe/rem:fo-danger}{Remark}{rem:fo-danger}{Failure of maximal support lies on the unsafe side}
\claimleadin{%
Which way a hypothesis of Theorem~\ref{thm:firstorder} fails decides what
the runtime checks of it are worth (Remark~\ref{rem:fo-check-exposures}). For
(H-SUPP) the failure is on the unsafe side.
}
\begin{remark}[Failure of maximal support lies on the unsafe side]\label{rem:fo-danger}
The functional $T$ is \emph{concave} on the positive semidefinite cone
(Lemma~\ref{lem:cone-concavity}), hence so is each deficit
$f_{\hat c,s}$ of \eqref{eq:f-def} and so is
$K_{\hat c}=P_{\hat c}+\sum_sf_{\hat c,s}$ of \eqref{eq:fK}. The objective
carries $K_{\hat c}$ only through the term
$\frac{q(\hat c)}{\alpha-1}\log_2\frac{q(\hat c)}{K_{\hat c}}$, whose
partial derivative in $K_{\hat c}$ is
$-\frac{q(\hat c)\log_2\mathrm e}{(\alpha-1)K_{\hat c}}<0$ for
$\alpha>1$, $q(\hat c)>0$, $K_{\hat c}>0$: the objective is
\emph{decreasing} in $K_{\hat c}$, through a negated logarithm.
Consequently, at a point where some branch operator does \emph{not} attain
maximal support and a feasible direction $h$ enlarges it, the
$+\infty$ one-sided derivative that
Theorem~\ref{thm:gradient}\,(ii) permits for $T$, and hence for
$K_{\hat c}$, becomes a $-\infty$ one-sided derivative for
$O^{J}_{\alpha}$ along $h$. Two things follow. First, no subgradient of
$O^{J}_{\alpha}$ exists at such a point in such a direction, so
\eqref{eq:fo-subgrad} is not merely unproved there, it is false for the
split-form gradient: the objective falls away faster than any linear
model. Second, the split-form gradient then makes the linearisation
\emph{over}-estimate $O^{J}_{\alpha}$ away from the face, so the bracket
\eqref{eq:hlb} can exceed $\hlow$. The failure is on the \emph{unsafe}
side. That is why (H-SUPP) is a hypothesis of the theorem rather than a
convenience of the implementation, and why a failed check of it
(Requirement~\ref{req:support-rank-check}) removes that iterate's bracket
from the bound (Requirement~\ref{req:first-order-withdrawal}).
\end{remark}
\claimmeta{\metalabel{Uses}\mref{Thm.}{thm:branch-formula}, \mref{Thm.}{thm:gradient}, \mref{Thm.}{thm:firstorder}
\metasep \metalabel{Used by}\mref{Rem.}{rem:fo-check-exposures}
\metasep newly derived
\metasep accepted}

% claim meat-core/frank-wolfe/rem:fo-nopert
\claimsummary{meat-core/frank-wolfe/rem:fo-nopert}{A gradient taken at a perturbed operator is the gradient of a different function, and no statement here relates the two, so the explicit perturbation route has no accounted instance. Every path here differentiates on a face at maximal support; an iterate with a perturbed gradient adds no bracket.}
\claimbookmark{meat-core/frank-wolfe/rem:fo-nopert}{Remark}{rem:fo-nopert}{A perturbed gradient bounds another function}
\begin{remark}[A perturbed gradient bounds another function]\label{rem:fo-nopert}
The rules drawn from Theorem~\ref{thm:gradient}\,(i) and~(ii), the face
rule of Requirement~\ref{req:split-form-gradient} and the one-sided
rule of Requirement~\ref{req:gradient-at-maximal-support}, are the rule of
every path this module specifies, without exception, and there is no
second path: no numerical lower-bound certificate is specified here, so no
route perturbs into the strictly positive cone instead of differentiating
on a face, and neither rule is another path's concern. Requirement~\ref{req:gradient-at-maximal-support} names
two routes to a usable gradient: keep the support maximal, ``for instance
after an explicit perturbation''. Only the first is available to
Theorem~\ref{thm:firstorder}, and (H-NOPERT) says so; it is the only route
this module supplies an accounting for.

A perturbation is not unsound in principle; a gradient taken at a
perturbed operator is the gradient of a \emph{different function}, and
this module states no relation between the two. Concretely: replacing $X$
by $X^{(\delta)}=(1-\delta)X+\delta\,\sigma$ for some $\sigma\succ0$ and
evaluating $\nabla T$ there yields the linearisation of an objective
$O^{(\delta)}_{\alpha}$ built from the perturbed branch operators. To
conclude anything about $\hlow=\inf O^{J}_{\alpha}$ from a bound on
$\inf O^{(\delta)}_{\alpha}$ one needs a \emph{domination} statement,
$O^{(\delta)}_{\alpha}\le O^{J}_{\alpha}$ pointwise or an explicit
two-sided error, and this module proves none. The phrase ``after an
explicit perturbation'' in
Requirement~\ref{req:gradient-at-maximal-support} therefore has no
accounted instance anywhere in this module. It stands there as a statement
about what a solver may do to keep a gradient defined, which is true and
is all it says.

The consequence for a run is a \emph{withdrawal}, not a weakening: an
iterate whose gradient was taken after a perturbation contributes no
bracket to $\hlb$ (Requirement~\ref{req:first-order-withdrawal}), so a run
that perturbs the gradient at every iterate reports no $\hlb$ at all. Reporting a perturbed
bracket as a \exttermref{../foundations/contract.pdf}{term-first-order}{first-order} bound would be reporting a bound on a
function nobody asked about.
\end{remark}
\claimmeta{\metalabel{Uses}\mref{Thm.}{thm:gradient}, \mref{Thm.}{thm:firstorder}
\metasep \metalabel{Used by}\mref{Rem.}{rem:fo-iterate-repair}
\metasep newly derived
\metasep accepted}

% claim meat-core/frank-wolfe/cor:fo-keylength
\claimsummary{meat-core/frank-wolfe/cor:fo-keylength}{With max\{0, h\_lb\} in place of the rate, the fixed-length key-length bound gives an \ensuremath{\varepsilon}\_pa-secret and \ensuremath{\varepsilon}\_ev-correct key for a R\'enyi order in (1, 2), whether the iterates are exact or relaxed. If h\_lb \ensuremath{\le} 0 and \ensuremath{\varepsilon}\_ev \textless{} 1, no key is claimed.}
\claimbookmark{meat-core/frank-wolfe/cor:fo-keylength}{Corollary}{cor:fo-keylength}{Key length from the first-order bound}
\claimleadin{%
The first-order bound enters the core's key-length corollary as the lower
bound that corollary asks for. Parts~(ii) and~(iiib) of
Theorem~\ref{thm:firstorder} conclude the same inequality
$\max\{0,\hlb\}\le\hlow$ against the same exact $\hlow$, and (iiib), which
admits an iterate in the \termref{term-relaxed-feasible-set}{relaxed feasible set}, is the part a run of
this module is in (Requirement~\ref{req:iterate-checks}).
}
\begin{corollary}[Key length from the first-order bound]\label{cor:fo-keylength}
Let $\alpha\in(1,2)$, let the hypotheses of Theorem~\ref{thm:keylength}
hold in its unconstrained or constrained branch, the run as in
Lemma~\ref{lem:memory-lift}\,(a) and hypothesis (H-CHAIN) included, at every
all-at-once attack with finite registers, with $\Feas^{J}$ carrying
the declared $\CList$ in the constrained branch, let those of
Lemma~\ref{lem:ev-correctness} hold, A-SEQ with its clause (v) among the
former, and let $\hlb$ be as in
Theorem~\ref{thm:firstorder}\,(ii) or~(iiib), at some or all iterates
with the lower bounds of its part~(iv) in place of the pairings. If the
output length
satisfies
\begin{equation}\label{eq:keylength-fo}
  \ell\;\le\;n\,\max\{0,\hlb\}-\lEC-\Big\lceil\log_2\frac1{\eev}\Big\rceil
  -\frac{\alpha}{\alpha-1}\log_2\frac1{\epa}+1,
\end{equation}
then the protocol is $\epa$-secret and $\eev$-correct on the accepted
state, hence, for $\ell\ge1$, $(\epa+\eev)$-secure in the
\termref{term-honest-authentication-setting}{honest-authentication setting}, and under the further hypotheses of
Corollary~\ref{cor:composed-security}\,(iii)
$(\epa+\eev+2\varepsilon_{\mathrm{auth}})$-secure over the classical channel
of that part, with net key gain $\ell-k_{\mathrm{auth}}$ when both parties
output a key, the secrecy and the security statements holding at every
output state, whatever the dimension of the adversary's registers and the
order in which the source releases its signals, as
Corollary~\ref{cor:composed-security} states them, and the correctness
statement at every attack; the four conclusions of
Proposition~\ref{prop:monotone} hold verbatim for the
\termref{term-maximal-key-length}{maximal key length} built from \eqref{eq:keylength-fo}; and if
$\hlb\le0$ and $\eev<1$ no key is claimed.
\end{corollary}
\claimmeta{\metalabel{Uses}\mref{Lem.}{lem:memory-lift}, \mref{Lem.}{lem:honest-attacks}, \mref{Cor.}{cor:keylength-bound}, \mref{Thm.}{thm:firstorder}
\metasep \metalabel{Used by}\mref{Cor.}{cor:fo-reference}, \mref{Cor.}{cor:certified-keylength}, \mref{Cor.}{cor:certified-keylength-reference}, \mref{Rem.}{rem:fo-arithmetic}, \mref{Rem.}{rem:fo-returned-point}, \mref{Rem.}{rem:fo-negative-values}, \mref{Req.}{req:first-order-clamp}, \mref{Req.}{req:bracket-on-reference}
\metasep newly derived
\metasep implemented (effective: accepted)
\metasep \metalabel{Code}\coderef{https://github.com/fraunhofersg/certiqs-key/blob/v0.1.0/src/certiqs_key/solvers/frank_wolfe/general_keyrate.py}{certiqs_key.solvers.frank_wolfe.general_keyrate:key_length_from_first_order}
\metasep \metalabel{Oracles}\coderef{https://github.com/fraunhofersg/certiqs-key/blob/v0.1.0/tests/oracles/registry.py}{registry.py}}

\begin{proof}
This is Corollary~\ref{cor:keylength-bound} of the core with
$h\coloneqq\hlb$. Theorem~\ref{thm:firstorder}\,(ii), respectively~(iiib),
with its part~(iv) where lower bounds stand in for the pairings, gives
$\hlb\le\hlow$ in exact real arithmetic, which is that corollary's
hypothesis \eqref{eq:lower-bound-h}, and \eqref{eq:keylength-fo} is its
\eqref{eq:keylength-h}. Its clause~(a) applies because $\hlb$ depends on
the iterates, the pairs or lower bounds its brackets use and the
declaration, and on none of $(\lEC,\eev,\epa)$; its clause~(b) is the last
sentence of the statement.
Nothing is added to the core's argument: no additional subtrahend, no
floor beyond the one $\max\{0,\cdot\}$ already supplies, and no arithmetic
requirement of any kind.
\end{proof}

% claim meat-core/frank-wolfe/cor:fo-dropped
\claimsummary{meat-core/frank-wolfe/cor:fo-dropped}{In either source mode the first-order bound holds for the dropped program (\hyperlink{claim-meat-core/def:dropped}{Definition~\ref*{def:dropped}}), with maximal support asked only of key-active pairs (acceptance value, joint sector) whose entropy is kept. It bounds h\_low from below, also at iterates feasible only up to windows.}
\claimbookmark{meat-core/frank-wolfe/cor:fo-dropped}{Corollary}{cor:fo-dropped}{First-order bound for the dropped program}
\claimleadin{%
A run that gives up the deficits of a \termref{term-dropped-set}{dropped set} takes its
first-order bound from the \termref{term-dropped-program}{dropped program}. The corollary is to that
program what Theorem~\ref{thm:firstorder}\,(ii) and~(iiib) are to the
undropped one, with the support hypothesis asked of each kept key-active
joint pair and positivity asked of the dropped branch quantity, which is
never smaller than the undropped one. A \termref{term-fully-dropped}{fully dropped} joint sector
needs no maximal support, no rank decision and no matrix function. In the
relaxed reading the three hypotheses are conditions on the iterate alone
and name no set, and no quantity relaxed at $\varepsilon>0$ appears in the
conclusion: the chain
\eqref{eq:hlb-dropped} compares
$\underline{h}^{\mathrm{FO},\mathcal{D}}_{\alpha}$ against the exact
$\underline{h}^{\mathcal{D}}_{\alpha}$ directly, and no non-negativity of
a relaxed infimum is claimed. On an empty feasible set both values are
$+\infty$, so the bound needs the non-empty feasible set of an admissible
declaration only to be informative. The corollary holds in either source
mode: in \exttermref{../transmitter/contract.pdf}{term-untrusted-source-mode}{untrusted-source mode} the Choi step is the identity, so the
frame and the operator coupling of the trusted-source proof become the
state blocks and the normalisation.
}
\begin{mainclaim}
\begin{corollary}[First-order bound for the dropped program]
\label{cor:fo-dropped}
Let $\alpha\in(1,\infty)$, let $\Decl$ satisfy Definition~\ref{def:admissible-declaration} in
the general format of Definition~\ref{def:declaration}, in either source
mode, in \exttermref{../transmitter/contract.pdf}{term-trusted-source-mode}{trusted-source mode} with a frame per Alice sector as in
Theorem~\ref{thm:choi}, let $\acSet$ satisfy
Definition~\ref{def:admissible-acceptance-set}, and let $\mathcal{D}$ be a dropped set
(Definition~\ref{def:dropped}). Write an Alice sector as
$\mathsf t\in\mathbb{T}$ in this claim, the letter $t$ indexing iterates as
in Theorem~\ref{thm:firstorder}. Put
$O^{J,\mathcal{D}}_{\alpha}(\mathbf q,J)\coloneqq
O^{\mathcal{D}}_{\alpha}(\mathbf q,\chi(J))$ and
$K^{J,\mathcal{D}}_{\hat c}\coloneqq K^{\mathcal{D}}_{\hat c}\circ\chi$, with
$\chi$ the assignment of that theorem in trusted-source mode. In
untrusted-source mode the program is read as
Definition~\ref{def:untrusted-choi} reads it, with no frame: $\chi$ is the
identity, $J=\omega$ is the list of the blocks
$J_{(\mathsf t,s)}=\omega^{(\mathsf t,s)}$ over the joint sectors,
$A_{0,\mathsf t}$ and $A_0$ stand for $A_{\mathsf t}$ and $A$, and
$\Feas^{J}$ is $\Feas$ of
Definition~\ref{def:feasible} in the general format, whose condition
coupling the blocks is the normalisation $\sum_{(\mathsf t,s)}
\tr\omega^{(\mathsf t,s)}=1$; so $O^{J,\mathcal{D}}_{\alpha}$ and
$K^{J,\mathcal{D}}_{\hat c}$ are $O^{\mathcal{D}}_{\alpha}$ and
$K^{\mathcal{D}}_{\hat c}$. Let
$x_t=(\mathbf q_t,J_t)\in\acSet\times\Feas^{J}$, $t\in\Itset$, satisfy
\begin{itemize}[leftmargin=3.4em]
\item[(H-POS$^{\mathcal{D}}$)] (H-POS) of Theorem~\ref{thm:firstorder} with
  $K^{J,\mathcal{D}}_{\hat c}$ in place of $K^{J}_{\hat c}$;
\item[(H-SUPP$^{\mathcal{D}}$)] for every $\hat c\in\Chkey$ and every joint
  pair $(\mathsf t,s)\in\mathbb{J}_K(\hat c)$ with
  $(\hat c,(\mathsf t,s))\notin\mathcal{D}$, the branch operator
  $\cG^{J}_{\hat c,s}(J_{t,(\mathsf t,s)})$ of the block
  $J_{t,(\mathsf t,s)}\in\Pos{A_{0,\mathsf t}\Qsys_s}$ attains the maximal
  support of that joint pair, the support it has at
  $J_{t,(\mathsf t,s)}\succ0$; nothing is asked of a pair that is key-inert
  or dropped;
\item[(H-NOPERT$^{\mathcal{D}}$)] $(\lmoq_t,\lmoJ_t)$ is the gradient of
  $O^{J,\mathcal{D}}_{\alpha}$ at $x_t$ itself, assembled as in (H-NOPERT)
  with the kept summands only and with $K^{J,\mathcal{D}}_{\hat c}$ in the
  scalar prefactors, its blocks on fully dropped joint sectors being
  \eqref{eq:drop-gradient} in Choi coordinates;
\end{itemize}
and let $(v_t,S_t)$ satisfy \eqref{eq:lmo} for these linear forms, or let,
at some or all $t$, lower bounds $b^{\mathrm A}_t$, $b^{\mathrm F}_t$ on
one or both of its minima, as in \eqref{eq:lmo-bounds} for these linear
forms, stand in for the pairings as Theorem~\ref{thm:firstorder}\,(iv)
lets them. Then, in
exact real arithmetic, with $\underline{h}^{\mathrm{FO},\mathcal{D}}_{\alpha}$
the expression \eqref{eq:hlb} for $O^{J,\mathcal{D}}_{\alpha}$, with those
replacements where they are made,
\begin{equation}\label{eq:hlb-dropped}
  \max\{0,\underline{h}^{\mathrm{FO},\mathcal{D}}_{\alpha}\}\;\le\;\underline{h}^{\mathcal{D}}_{\alpha}
  \;\le\;\hlow\;\le\;\hqkd(\CList).
\end{equation}
(H-POS$^{\mathcal{D}}$) is implied by (H-POS), and
(H-SUPP$^{\mathcal{D}}$) by (H-SUPP) read on joint blocks. For
$\mathcal{D}=\varnothing$ and $|\mathbb{T}|=1$ the statement is
Theorem~\ref{thm:firstorder}\,(ii), with its part~(iv) where lower bounds
stand in.

\emph{The relaxed reading.} Let $\varepsilon\ge0$ and let
$\Feas^{J}(\varepsilon)$ be the \termref{term-relaxed-feasible-set}{relaxed feasible set} of
Definition~\ref{def:relaxed-feasible-set}, read in the general format, with
one coupling clause per Alice sector in trusted-source mode and its
normalisation clause in untrusted-source mode. Let
$x_t\in\acSet\times\Feas^{J}(\varepsilon)$, $t\in\Itset$, satisfy
(H-POS$^{\mathcal{D}}$), (H-SUPP$^{\mathcal{D}}$) and
(H-NOPERT$^{\mathcal{D}}$) as stated, and let $(v_t,S_t)$ satisfy
\eqref{eq:lmo} \emph{as written}, or the lower bounds standing in for its
pairings satisfy \eqref{eq:lmo-bounds}, with the two minima over $\acSet$
and over the \emph{exact} $\Feas^{J}$. Then \eqref{eq:hlb-dropped} holds in
full, and the case $\varepsilon=0$ is the statement above.

\emph{An empty feasible set.} In place of the hypothesis on $\Decl$ above,
let $\Decl$ be a tuple of the form of Definition~\ref{def:declaration}, in
the general format and in either source mode, that satisfies conditions
(i)--(vi) and (viii) of Definition~\ref{def:admissible-declaration}, and form
$\Feas$, $O_{\alpha}$, $\hlow$ and $\underline{h}^{\mathcal{D}}_{\alpha}$ from
it by the formulas of Definitions~\ref{def:feasible} and~\ref{def:dropped}.
Let $\Feas=\varnothing$. Then
$\underline{h}^{\mathcal{D}}_{\alpha}=\hlow=+\infty$, so every real number is
a lower bound on both, the value of \eqref{eq:hlb} for
$O^{J,\mathcal{D}}_{\alpha}$ that a run forms at its iterates included.
Consequently the first two inequalities of \eqref{eq:hlb-dropped}, for a real
$\underline{h}^{\mathrm{FO},\mathcal{D}}_{\alpha}$, hold of every such
tuple, whether or not it satisfies condition~(vii): under the hypotheses above where it does,
and by the preceding sentence where it does not.
\end{corollary}
\end{mainclaim}
\claimmeta{\metalabel{Uses}\mref{Def.}{def:admissible-declaration}, \mref{Def.}{def:admissible-acceptance-set}, \mref{Def.}{def:joint-sectors}, \mref{Lem.}{lem:joint-sectors}, \mref{Def.}{def:feasible}, \mref{Lem.}{lem:choi-factor}, \mref{Thm.}{thm:choi}, \mref{Def.}{def:untrusted-choi}, \mref{Def.}{def:branch-support}, \mref{Lem.}{lem:branch-support}, \mref{Thm.}{thm:gradient}, \mref{Cor.}{cor:full-gradient}, \mref{Def.}{def:dropped}, \mref{Prop.}{prop:drop-bound}, \mref{Cor.}{cor:drop-program}, \mref{Def.}{def:relaxed-feasible-set}, \mref{Prop.}{prop:relaxed-set-properties}, \mref{Def.}{def:first-order-setting}, \mref{Lem.}{lem:fo-khyp}, \mref{Lem.}{lem:fo-face}, \mref{Thm.}{thm:firstorder}, \extref{../foundations/contract.pdf}{foundations/def:kl}
\metasep \metalabel{Used by}\mref{Lem.}{lem:certified-bracket}, \mref{Lem.}{lem:certified-bracket-enclosed}, \mref{Cor.}{cor:certified-keylength}, \mref{Cor.}{cor:certified-keylength-reference}, \mref{Rem.}{rem:fo-returned-point}
\metasep newly derived
\metasep implemented (effective: accepted)
\metasep \metalabel{Code}\coderef{https://github.com/fraunhofersg/certiqs-key/blob/v0.1.0/src/certiqs_key/solvers/frank_wolfe/joint_driver.py}{certiqs_key.solvers.frank_wolfe.joint_driver:run_joint_frank_wolfe}}

\begin{proof}
\emph{Idea.} The first inequality of \eqref{eq:hlb-dropped} follows the
four steps of the proof of Theorem~\ref{thm:firstorder}\,(ii) with the
changes below, preceded by the two things the dropped program needs that
that theorem's proof does not state: the support geometry of a joint pair
(Step~0) and the acceptance-support bookkeeping for $K^{J,\mathcal{D}}$
(Step~0$'$). The relaxed reading changes Step~2 only. The proof is
written once for both source modes, with the reading of untrusted-source
mode fixed in the paragraph on notation.

\emph{The last two inequalities} of \eqref{eq:hlb-dropped} are
Proposition~\ref{prop:drop-bound}\,(iii).

\emph{Notation.} Write $p=(\mathsf t,s)$ for a joint pair. In
trusted-source mode put
$W_{\mathsf t}\coloneqq(\tau^{(\mathsf t)}_A)^{1/2}\iota_{\mathsf t}$ and
$\chi_p(\cdot)\coloneqq(W_{\mathsf t}\otimes\idon{\Qsys_s})(\cdot)
(W_{\mathsf t}^{\dagger}\otimes\idon{\Qsys_s})$ for the assignment of
Theorem~\ref{thm:choi} in the general format, so that
$\chi=\bigoplus_p\chi_p$ and $W\coloneqq\bigoplus_{\mathsf t}W_{\mathsf t}$
satisfies $WW^{\dagger}=\tau_A$
(Lemma~\ref{lem:choi-factor}\,(ii)), and
$S_{\mathsf t}\coloneqq\operatorname{supp}\tau^{(\mathsf t)}_A
=\operatorname{ran}W_{\mathsf t}$
(Lemma~\ref{lem:choi-factor}\,(i)). In untrusted-source mode
(Definition~\ref{def:untrusted-choi}\,(a)) let $W_{\mathsf t}$ be the
inclusion of $\Hs{A_{\mathsf t}}$ into $\Hs{A}$ and $\chi_p$ the conjugation
by $W_{\mathsf t}\otimes\idon{\Qsys_s}$, that is, the inclusion of
$\Pos{A_{\mathsf t}\Qsys_s}$ into $\Pos{A\Qsys_s}$ as the subcone of
Definition~\ref{def:joint-sectors}; then $\chi=\bigoplus_p\chi_p$ is the
identity on lists over the joint sectors, and
$S_{\mathsf t}\coloneqq\Hs{A_{\mathsf t}}=\operatorname{ran}W_{\mathsf t}$.
In both modes $\chi_p$ is linear and preserves positive semidefiniteness, and
$\operatorname{ran}W_{\mathsf t}=S_{\mathsf t}\subseteq\Hs{A_{\mathsf t}}$.
Write
$\cG^{J}_{\hat c,p}\coloneqq\cG_{\hat c,p}\circ\chi_p$, which is the
restriction of $\cG^{J}_{\hat c,s}$ to the blocks supported on
$\Hs{A_{0,\mathsf t}\Qsys_s}$, the reading (H-SUPP$^{\mathcal{D}}$) writes,
and $X_{\hat c,p}\coloneqq\cG^{J}_{\hat c,p}(J_{t,p})$. Let
$\widetilde\Pi_{\mathsf t}$ be $\Pi^{A}_{\mathsf t}$ acting on the $A$ factor
of the branch space $\Hs{R\,XY\,C_AC_B\,A\Qsys_s}$, and put
\begin{equation}\label{eq:Wcs-joint}
  W_{\hat c,p}\coloneqq\widetilde\Pi_{\mathsf t}\,W_{\hat c,s}
  \qquad\text{with }W_{\hat c,s}\text{ of }\eqref{eq:Wcs}.
\end{equation}

\emph{Step 0 (the support geometry of a joint pair).} \emph{(a)} By
\eqref{eq:Wcs}, $W_{\hat c,s}=\sum_{\kappa\in\Kann_{\mathrm{pass}}(\hat c)}
K^{(s)}_{\kappa}(\Sadm\otimes\Hs{\Qsys_s})$, and
$\Sadm=\bigoplus_{\mathsf t}S_{\mathsf t}$. In trusted-source mode
$\Sadm=\operatorname{supp}\tau_A$, and
$\tau_A=\bigoplus_{\mathsf t}\tau^{(\mathsf t)}_A$ is block
diagonal along $\mathbb{T}$ by Definition~\ref{def:admissible-declaration}\,(iii) in the general
format, so $\Sadm=\bigoplus_{\mathsf t}\operatorname{supp}
\tau^{(\mathsf t)}_A$; in untrusted-source mode $\Sadm=\Hs{A}$
(Definition~\ref{def:branch-support}), which is
$\bigoplus_{\mathsf t}\Hs{A_{\mathsf t}}$ by the general format of
Definition~\ref{def:declaration}. By the computation in the proof of
Lemma~\ref{lem:joint-sectors}\,(i),
$K^{(s)}_{\kappa}(\Pi^{A}_{\mathsf t}\otimes\idon{})
=\widetilde\Pi_{\mathsf t}K^{(s)}_{\kappa}$, so each
$K^{(s)}_{\kappa}(S_{\mathsf t}\otimes
\Hs{\Qsys_s})$ lies in the range of $\widetilde\Pi_{\mathsf t}$; summing over
$\mathsf t$ and $\kappa$ exhibits $W_{\hat c,s}$ as a sum of subspaces of
\emph{mutually orthogonal} summands, hence
\begin{equation}\label{eq:Wcs-joint-split}
  W_{\hat c,s}=\bigoplus_{\mathsf t\in\mathbb{T}}W_{\hat c,(\mathsf t,s)},
  \qquad
  W_{\hat c,(\mathsf t,s)}=\sum_{\kappa\in\Kann_{\mathrm{pass}}(\hat c)}
  K^{(s)}_{\kappa}\big(S_{\mathsf t}\otimes
  \Hs{\Qsys_s}\big).
\end{equation}
\emph{(b)} For every Hermitian $\Delta$ on $\Hs{A_{0,\mathsf t}\Qsys_s}$,
$\operatorname{supp}\cG^{J}_{\hat c,p}(\Delta)\subseteq W_{\hat c,p}$. Indeed
$\Delta\oplus0$ is Hermitian on $\Hs{A_0\Qsys_s}$, so
Lemma~\ref{lem:fo-face}, \eqref{eq:fo-face} gives
$\operatorname{supp}\cG^{J}_{\hat c,s}(\Delta\oplus0)\subseteq W_{\hat c,s}$;
and $\chi_p(\Delta_{\pm})$ is supported on
$\Hs{A_{\mathsf t}\Qsys_s}$, $\operatorname{ran}W_{\mathsf t}$ lying in
$\Hs{A_{\mathsf t}}$, so by
Lemma~\ref{lem:joint-sectors}\,(i) the image of each lies in the
range of $\widetilde\Pi_{\mathsf t}$ and so does their difference. Intersecting
the two inclusions gives the claim by \eqref{eq:Wcs-joint}. \emph{(c)} For
$\varrho\succeq0$ the support of $\sum_{\kappa}K_{\kappa}\varrho
K_{\kappa}^{\dagger}$ is $\sum_{\kappa}K_{\kappa}(\operatorname{supp}\varrho)$
(the support of a sum, in the proof of Lemma~\ref{lem:branch-support}). At
$J_{t,p}\succ0$, $\operatorname{supp}\chi_p(J_{t,p})
=\operatorname{ran}(W_{\mathsf t}\otimes\idon{})
=S_{\mathsf t}\otimes\Hs{\Qsys_s}$, so
$\operatorname{supp}X_{\hat c,p}=W_{\hat c,p}$ by
\eqref{eq:Wcs-joint-split}, while for any $J_{t,p}\succeq0$ the inclusion
$\subseteq$ holds. So the ``maximal support of that joint pair'' of
(H-SUPP$^{\mathcal{D}}$) is $W_{\hat c,p}$, and that hypothesis reads
$\operatorname{supp}X_{\hat c,p}=W_{\hat c,p}$ for every kept key-active pair.

\emph{Step 0$'$ (the bookkeeping, which is Lemma~\ref{lem:fo-khyp} for
$K^{J,\mathcal{D}}$).} Part~(i) of that lemma is a statement about $\acSet$
alone and is unchanged: every $\mathbf q\in\acSet$ vanishes outside the
\termref{term-acceptance-support}{acceptance support} $\Chacc$ of \eqref{eq:chacc-def}, and for
$\mathbf q,\mathbf q'\in\acSet$ every feasible direction has
$q'(\hat c)-q(\hat c)=0$ at every $\hat c\notin\Chacc$. Part~(ii) holds for
$O^{J,\mathcal{D}}_{\alpha}$ on $\acSet\times\Feas^{J}(\varepsilon)$, for
every $\varepsilon\ge0$, with the cone clause of
Theorem~\ref{thm:branch-formula} replaced by
Proposition~\ref{prop:drop-bound}\,(i), which states
$0\le K^{\mathcal{D}}_{\hat c}\le P_{\hat c}$ for every $\omega$ block diagonal
along the joint sectors with positive semidefinite blocks; every point of
$\Feas^{J}(\varepsilon)$ is per-block positive semidefinite, the one
condition that set does not relax (Definition~\ref{def:relaxed-feasible-set}),
while $\chi$ preserves the joint-sector block structure and positivity
(Lemma~\ref{lem:choi-factor}\,(ii); in untrusted-source mode $\chi$ is the
identity). So for
$\hat c\notin\Chacc$, where $q(\hat c)=0$, the convention
$0\log_2\frac0v=0$ of \xref{foundations/def:kl}, stated for every
$v\ge0$ and the case $v=0$ included, makes the term
$\frac{\alpha}{\alpha-1}q(\hat c)\log_2\frac{q(\hat c)}{P^{J}_{\hat c}}$ and,
for $\hat c\in\Chkey$, the pair
$q(\hat c)\log_2\frac{q(\hat c)}{P^{J}_{\hat c}}
+\frac{q(\hat c)}{\alpha-1}\log_2\frac{q(\hat c)}{K^{J,\mathcal{D}}_{\hat c}}$
vanish at \emph{every} point of that set, whatever $P^{J}_{\hat c}$ and
$K^{J,\mathcal{D}}_{\hat c}$ are there and in particular when one of them
vanishes. Deleting those terms changes no value, so
$O^{J,\mathcal{D}}_{\alpha}$ is on that set a function of
$(\mathbf q|_{\Chacc},J)$ alone, (H-POS$^{\mathcal{D}}$) is stated on
$\Chacc$, and on $\Chkey\cap\Chacc$ for the $K$-part, exactly where it is
consumed, and a component of the acceptance-side gradient indexed outside
$\Chacc$ is paired against a direction component that is $0$ in every pairing
below, so its value is never read. No normalisation, coupling condition
or declared row is read in this step, which is why it holds on the
relaxed sets and not on the exact set alone.

\emph{Step 1 (the gradient exists and is a directional derivative in every
feasible direction).} Fix $t$, write $x\coloneqq x_t=(\mathbf q,J)$ and
$\nabla O\coloneqq(\lmoq_t,\lmoJ_t)$, let
$y=(\mathbf q',J')\in\acSet\times\Feas^{J}$ and $h\coloneqq y-x$, and put
$\Delta_p\coloneqq J'_p-J_{t,p}$, Hermitian as a difference of positive
semidefinite operators. By Step~0$'$ the objective along the segment is the
finite sum over $\hat c\in\Chacc$ of
$\frac{\alpha}{\alpha-1}q(\hat c)\log_2\frac{q(\hat c)}{P^{J}_{\hat c}}$ for
$\hat c\notin\Chkey$ and of the pair of Step~0$'$ for
$\hat c\in\Chkey\cap\Chacc$. In it $q(\hat c)$ and $P^{J}_{\hat c}$ are affine
and strictly positive at $x$ by (H-POS$^{\mathcal{D}}$), and
\begin{equation}\label{eq:KD-choi}
  K^{J,\mathcal{D}}_{\hat c}(J)=P^{J}_{\hat c}(J)
  +\sum_{p:\,(\hat c,p)\notin\mathcal{D}}
  f_{\hat c,p}\big(\chi_p(J_p)\big)
\end{equation}
by \eqref{eq:KD}, with $f_{\hat c,p}=T\circ\cG_{\hat c,p}
-\tr\circ\,\cG_{\hat c,p}$ and the second summand linear in $J_p$. Each kept
summand of \eqref{eq:KD-choi} falls into one of two cases. If $p$ is key-inert
for $\hat c$ then $f_{\hat c,p}\equiv0$ on the whole subcone
(Lemma~\ref{lem:joint-sectors}\,(iii)), hence identically along
the segment: the summand is constant, contributes $0$ to the derivative, and
no support hypothesis on it is consumed, which is why
(H-SUPP$^{\mathcal{D}}$) asks nothing of it. If $p\in\mathbb{J}_K(\hat c)$ and
$(\hat c,p)\notin\mathcal{D}$ then, by (H-SUPP$^{\mathcal{D}}$) with Step~0(c),
$\operatorname{supp}X_{\hat c,p}=W_{\hat c,p}$, and by Step~0(b)
$\operatorname{supp}\cG^{J}_{\hat c,p}(\Delta_p)\subseteq W_{\hat c,p}
=\operatorname{supp}X_{\hat c,p}$; so
Theorem~\ref{thm:gradient}\,(i) applies to the pair
$(X_{\hat c,p},\cG^{J}_{\hat c,p}(\Delta_p))$ and
$t'\mapsto T(X_{\hat c,p}+t'\cG^{J}_{\hat c,p}(\Delta_p))$ is differentiable at
$t'=0$, two-sidedly, the hypothesis being symmetric under
$\Delta_p\mapsto-\Delta_p$. And the $+\infty$ case of
Theorem~\ref{thm:gradient}\,(ii) is unreachable: that case sits
at directions strictly enlarging the support, Step~0(b) says no image of a
Hermitian direction leaves $W_{\hat c,p}$, and the directions this proof pairs
against are admissible for it because
$X_{\hat c,p}+t'\cG^{J}_{\hat c,p}(\Delta_p)
=\cG^{J}_{\hat c,p}\big((1-t')J_{t,p}+t'J'_p\big)\succeq0$ for $t'\in[0,1]$.
This is Lemma~\ref{lem:fo-face}\,(a), (b) read on the joint pair, with
Step~0(b) in place of \eqref{eq:fo-face}. A pair with
$(\hat c,p)\in\mathcal{D}$ is absent from \eqref{eq:KD-choi} altogether, so no
matrix power, no eigendecomposition and no support condition of a fully dropped
joint sector enters
(Corollary~\ref{cor:drop-program}\,(iii)). The remaining
dependence is through finitely many compositions of $\log_2$ with quantities
(H-POS$^{\mathcal{D}}$) keeps strictly positive, all differentiable there.
Hence $t'\mapsto O^{J,\mathcal{D}}_{\alpha}(x+t'h)$ is differentiable at
$t'=0$ with derivative $\langle\nabla O,h\rangle$, where, by the chain and sum
rules exactly as in the proof of
Corollary~\ref{cor:full-gradient}, whose maximal-support step is
Step~0(b) here, the acceptance-side component is \eqref{eq:grad-q} with
$K^{J,\mathcal{D}}_{\hat c}$ in place of $K^{J}_{\hat c}$, read as a vector in
$\mathbb{R}^{\Chacc}$, and the block of the feasible-side component at the
joint pair $p$ is \eqref{eq:grad-J} with the sums restricted to $\Chacc$,
respectively $\Chkey\cap\Chacc$, with $K^{J,\mathcal{D}}_{\hat c}$ in the
scalar prefactors, with $\Gamma^{(p)}_{\hat c}$ of
\eqref{eq:grad-pieces} read on the joint pair (the operators
$W_{\mathsf t}^{\dagger}M^{(\mathsf t)}_xW_{\mathsf t}\otimes F^{(s)}_v$), and
with the deficit piece $\Delta^{(p)}_{\hat c}$ carrying, in place of
$\mathbb 1[s\in\SecSet_K(\hat c)]$, the indicator
$\mathbb 1[p\in\mathbb{J}_K(\hat c)\text{ and }(\hat c,p)\notin\mathcal{D}]$.
That is the object (H-NOPERT$^{\mathcal{D}}$) names, and on a fully
dropped joint sector every deficit piece is absent, so the block is
$\chi_p^{\dagger}$ applied to \eqref{eq:drop-gradient}, which is what that
hypothesis says of such a sector.

\emph{Step 2 (convexity turns the derivative into an inequality).}
$O^{\mathcal{D}}_{\alpha}$ is jointly convex on $\Delta(\Chatset)\times\Feas$
with values in $[0,+\infty]$
(Corollary~\ref{cor:drop-program}\,(ii)), and
$O^{J,\mathcal{D}}_{\alpha}=O^{\mathcal{D}}_{\alpha}\circ(\idmap\times\chi)$
with $\chi$ \emph{linear}, so it is convex on the convex preimage
$\acSet\times\Feas^{J}$, a pullback along a linear map that needs no
injectivity, and $\Feas^{J}$ is convex and compact by
Theorem~\ref{thm:choi} in its general format, in untrusted-source mode
because it is $\Feas$ (Definition~\ref{def:feasible}). $x$ and $y$ lie in
that convex set and $O^{J,\mathcal{D}}_{\alpha}(x)$ is finite by
(H-POS$^{\mathcal{D}}$), so the three-chord inequality for a convex
extended-real function gives, exactly as in Step~2 of the proof of
Theorem~\ref{thm:firstorder}\,(ii),
\begin{equation}\label{eq:fo-subgrad-dropped}
  O^{J,\mathcal{D}}_{\alpha}(y)\;\ge\;O^{J,\mathcal{D}}_{\alpha}(x)
  +\langle\nabla O,\,y-x\rangle
  \qquad\text{for every }y\in\acSet\times\Feas^{J},
\end{equation}
the case $O^{J,\mathcal{D}}_{\alpha}(y)=+\infty$ included. Taking the infimum
over $y$: by Theorem~\ref{thm:choi} in its general format
$\chi(\Feas^{J})=\Feas$, in untrusted-source mode because $\chi$ is the
identity and $\Feas^{J}=\Feas$, and surjectivity is all a value needs, so
$\inf_{\acSet\times\Feas^{J}}O^{J,\mathcal{D}}_{\alpha}
=\inf_{\acSet\times\Feas}O^{\mathcal{D}}_{\alpha}
=\underline{h}^{\mathcal{D}}_{\alpha}$ by \eqref{eq:hlowD}; hence
$\underline{h}^{\mathcal{D}}_{\alpha}\ge O^{J,\mathcal{D}}_{\alpha}(x)
+\min_{y\in\acSet\times\Feas^{J}}\langle\nabla O,y-x\rangle$, the minimum
attained because $\acSet\times\Feas^{J}$ is compact and the map linear.

\emph{Steps 3 and 4.} The map $y\mapsto\langle\nabla O,y\rangle$ is a sum of
one term in each factor and $\acSet\times\Feas^{J}$ is a product set, so the
minimum separates over the product and \eqref{eq:lmo}, at the linear forms
$(\lmoq_t,\lmoJ_t)$ of this corollary, bounds each part below by the oracle
return; substituting gives
$\underline{h}^{\mathcal{D}}_{\alpha}\ge O^{J,\mathcal{D}}_{\alpha}(x_t)
+\langle\lmoq_t,v_t-\mathbf q_t\rangle+\langle\lmoJ_t,S_t-J_t\rangle$ for each
$t\in\Itset$. Where $b^{\mathrm A}_t$ or $b^{\mathrm F}_t$ stands in for a
pairing, \eqref{eq:lmo-bounds} bounds that part below by it instead, and
the same substitution gives the bracket with the replacement made, exactly
as in the proof of Theorem~\ref{thm:firstorder}\,(iv); where both stand in,
no oracle return enters. $\Itset$ is finite and non-empty, so the largest of those
right-hand sides, which is \eqref{eq:hlb} for
$O^{J,\mathcal{D}}_{\alpha}$, that is
$\underline{h}^{\mathrm{FO},\mathcal{D}}_{\alpha}$, is a lower bound too.
Finally $\underline{h}^{\mathcal{D}}_{\alpha}\ge0$ by
Proposition~\ref{prop:drop-bound}\,(iii), a statement about
the exact set, so $\max\{0,
\underline{h}^{\mathrm{FO},\mathcal{D}}_{\alpha}\}$ is the maximum of two lower
bounds on $\underline{h}^{\mathcal{D}}_{\alpha}$ and is one. This is the first
inequality of \eqref{eq:hlb-dropped}.

\emph{The implications, and the degenerate case.} (H-POS)
$\Rightarrow$ (H-POS$^{\mathcal{D}}$): $K^{J,\mathcal{D}}_{\hat c}\ge
K^{J}_{\hat c}$ by
Proposition~\ref{prop:drop-bound}\,(i), so
$K^{J}_{\hat c}>0$ gives $K^{J,\mathcal{D}}_{\hat c}>0$, and the other clauses
are identical. (H-SUPP) on joint blocks $\Rightarrow$
(H-SUPP$^{\mathcal{D}}$): with $J_s=\bigoplus_{\mathsf t}J_{t,(\mathsf t,s)}$,
Lemma~\ref{lem:joint-sectors}\,(i) gives
$X_{\hat c,s}=\bigoplus_{\mathsf t}X_{\hat c,(\mathsf t,s)}$ along mutually
orthogonal subspaces, so $\operatorname{supp}X_{\hat c,s}
=\bigoplus_{\mathsf t}\operatorname{supp}X_{\hat c,(\mathsf t,s)}$, while
$W_{\hat c,s}=\bigoplus_{\mathsf t}W_{\hat c,(\mathsf t,s)}$ by
\eqref{eq:Wcs-joint-split} with
$\operatorname{supp}X_{\hat c,(\mathsf t,s)}\subseteq
W_{\hat c,(\mathsf t,s)}$ by Step~0(c); two direct sums over the same
orthogonal summands that are equal with term-wise inclusions are equal term by
term, so maximal support of the sector gives maximal support of each of its
joint pairs. At $\mathcal{D}=\varnothing$ and $|\mathbb{T}|=1$:
$K^{\mathcal{D}}_{\hat c}=K_{\hat c}$ and
$O^{\mathcal{D}}_{\alpha}=O_{\alpha}$ by
Definition~\ref{def:dropped},
$\mathbb{J}_K(\hat c)=\{t_0\}\times\SecSet_K(\hat c)$ and
$W_{\hat c,(t_0,s)}=W_{\hat c,s}$ by
Definition~\ref{def:joint-sectors} and
\eqref{eq:Wcs-joint-split}, so the three hypotheses are (H-POS), (H-SUPP),
(H-NOPERT) and \eqref{eq:hlb-dropped} is \eqref{eq:hlb} together with
$\hlow\le\hqkd(\CList)$: Theorem~\ref{thm:firstorder}\,(ii).

\emph{The relaxed reading.} The proof is the one above with a single change,
made in Step~2: $x$ ranges over $\acSet\times\Feas^{J}(\varepsilon)$ while the
comparison point $y$ ranges over the \emph{exact} $\acSet\times\Feas^{J}$.
Steps~0, 0$'$ and~1 are unchanged and read only per-block positive
semidefiniteness of $x$: Step~0(b), (c) ask that of $J_{t,p}$, Step~0$'$
asks it through
Proposition~\ref{prop:drop-bound}\,(i), and $\Delta_p$ is
Hermitian as a difference of a positive semidefinite $J'_p$ from $\Feas^{J}$
and a positive semidefinite $J_{t,p}$ from $\Feas^{J}(\varepsilon)$, per-block
positive semidefiniteness being the one condition that set does not relax
(Definition~\ref{def:relaxed-feasible-set}). In Step~2 both endpoints lie in
$\acSet\times\Feas^{J}(\varepsilon)$, which is convex
(Proposition~\ref{prop:relaxed-set-properties}\,(b); $\acSet$ by
Definition~\ref{def:admissible-acceptance-set}), and
$O^{J,\mathcal{D}}_{\alpha}$ is convex on it. The transfer is the one
Proposition~\ref{prop:relaxed-set-properties}\,(c) performs for
Theorem~\ref{thm:convexity}, and it goes through for the dropped
objective for the same reason: the convexity argument of
Corollary~\ref{cor:drop-program}\,(ii) reads of its domain only
that it is a convex set on which $P_{\hat c}$ is affine and
$K^{\mathcal{D}}_{\hat c}$ is non-negative, concave and continuous;
Corollary~\ref{cor:drop-program}\,(i) states that concavity and
continuity on the whole \emph{cone} of operators block diagonal along the joint
sectors, while
Proposition~\ref{prop:drop-bound}\,(i) states
$0\le K^{\mathcal{D}}_{\hat c}\le P_{\hat c}$ there; and
$\chi(\Feas^{J}(\varepsilon))$ is a convex subset of that cone, so
$O^{\mathcal{D}}_{\alpha}$ is convex on
$\acSet\times\chi(\Feas^{J}(\varepsilon))$ and pulls back along the linear
$\chi$. So the three-chord inequality gives
\eqref{eq:fo-subgrad-dropped} for every $y\in\acSet\times\Feas^{J}$, and the
infimum is taken over \emph{that} set, not over the relaxed one; Steps~3
and~4 are then as above, and the clamp reads
Proposition~\ref{prop:drop-bound}\,(iii) at $\varepsilon=0$,
where it holds. What this uses of the relaxed set, in full: that
$\Feas^{J}\subseteq\Feas^{J}(\varepsilon)$
(Proposition~\ref{prop:relaxed-set-properties}\,(a)); that
$\acSet\times\Feas^{J}(\varepsilon)$ is convex (\emph{ibid.}\,(b)); that every
member is per-block positive semidefinite; and the convexity transfer just
made, whose own derivation consumes the compactness of (b) through the
continuity of $K^{\mathcal{D}}_{\hat c}$, so compactness is used one level down
and never in the minimisation. Of the \emph{exact} set it uses compactness,
non-emptiness, and that the infimum of $O^{J,\mathcal{D}}_{\alpha}$ over it is
$\underline{h}^{\mathcal{D}}_{\alpha}$. No statement is made about the
infimum of $O^{J,\mathcal{D}}_{\alpha}$ over the relaxed set, and none is
used.

\emph{An empty feasible set.} No claim is used for the two values being
$+\infty$. By
\eqref{eq:hlowD} and \eqref{eq:hlow-def}, $\underline{h}^{\mathcal{D}}_{\alpha}$
and $\hlow$ are infima over $\mathbf q\in\acSet$ and $\omega\in\Feas$; with
$\Feas=\varnothing$ each inner infimum is over the empty set and is $+\infty$
for every $\mathbf q$, and so is the outer one, whether $\acSet$ is empty or
not. Every real number is below $+\infty$. For the consequence: where the
tuple satisfies condition~(vii) as well, it satisfies
Definition~\ref{def:admissible-declaration} in full, and with it the two
conditions Definition~\ref{def:declaration} places on its raw-key map, which
are conditions (v) and (viii); it is then a protocol declaration satisfying
the hypothesis of the statement above; where it does not, $\Feas=\varnothing$, both values
are $+\infty$, and a real number is below them.
\end{proof}

% claim meat-core/frank-wolfe/cor:fo-reference
\claimsummary{meat-core/frank-wolfe/cor:fo-reference}{At any points, such as those of a search on the snapped declaration or of one continued on the companion, the bracket formed with the reference companion's objective, gradient, entries and linear minima bounds the companion's value from below; its key length is the reference model's.}
\claimbookmark{meat-core/frank-wolfe/cor:fo-reference}{Corollary}{cor:fo-reference}{The first-order bound on the reference companion}
\claimleadin{%
Theorem~\ref{thm:firstorder} asks nothing about where its points came from:
they may be any finite set of points of the relaxed feasible set at which its
hypotheses hold. A run may therefore search on the declaration $\Decl$
compiled from the snapped measurement and form the bracket \eqref{eq:hlb} on
the compiled reference companion $\Decl^{\mathrm{ref}}$ at the points the
search produced. The snap then only steers the search, and its displacement
plays no part in the bound
(Remark~\ref{rem:declared-measurement}). The corollary is the theorem read
for $\Decl^{\mathrm{ref}}$, with the hypotheses that carry over from $\Decl$
named; what a run evaluates, and how, is
Requirement~\ref{req:bracket-on-reference}, and the arithmetic of that
evaluation is Remark~\ref{rem:fo-arithmetic}.
}
\begin{mainclaim}
\begin{corollary}[The first-order bound on the reference companion]\label{cor:fo-reference}
Let $\Decl$ and $\Decl^{\mathrm{ref}}$ be as in
Proposition~\ref{prop:reference-companion}, let $\alpha\in(1,\infty)$, and let
$\acSet$ be an \termref{term-admissible-acceptance-set}{admissible acceptance set}
(Definition~\ref{def:admissible-acceptance-set}), the one of the run, read
for both. Let $\varepsilon\ge0$ and let $x_t=(\mathbf q_t,J_t)$,
$t\in\Itset$, be a finite non-empty set of points, produced by any
computation, a search on $\Decl$ among them, such that every $x_t$ lies in
$\acSet\times\Feas^{J}(\varepsilon)$ of $\Decl^{\mathrm{ref}}$
(Definition~\ref{def:relaxed-feasible-set}) and satisfies (H-POS), (H-SUPP)
and (H-NOPERT) of Theorem~\ref{thm:firstorder} for the objective, the \termref{term-branch-map}{branch
maps} and the gradient of $\Decl^{\mathrm{ref}}$, and let $(v_t,S_t)$
satisfy \eqref{eq:lmo} at the gradient of $\Decl^{\mathrm{ref}}$ at $x_t$,
its feasible-side minimum over the exact $\Feas^{J}$ of $\Decl^{\mathrm{ref}}$
(Definition~\ref{def:first-order-setting}), or let, at some or all $t$,
lower bounds on one or both of those minima stand in for the pairings as
Theorem~\ref{thm:firstorder}\,(iv) lets them, \eqref{eq:lmo-bounds} read
at that gradient and over those sets. Write
$\hlb[\Decl^{\mathrm{ref}}]$ for \eqref{eq:hlb} so formed. Then:
\begin{enumerate}[label=(\roman*),leftmargin=2.4em]
\item in exact real arithmetic,
  $\max\{0,\hlb[\Decl^{\mathrm{ref}}]\}\le\hlow[\Decl^{\mathrm{ref}}]$, and
  $\hlow[\Decl^{\mathrm{ref}}]\le\hlow[\Decl^{\mathrm{ref}}_{\mathrm{ex}}]$
  for the companion with its unrelaxed entries;
\item the Choi variables of $\Decl$ and $\Decl^{\mathrm{ref}}$ are lists over
  the same \exttermref{../protocol-compiler/contract.pdf}{term-joint-sector}{joint sectors} with the same block dimensions,
  so every point of the one program is a point of the other; where no link
  the companion replaces carries entries, as on the encoded route,
  $\acSet\times\Feas^{J}(\varepsilon)$ is the same set for both; and a point with $J_{t,s}\succ0$ for every sector
  $s$ satisfies (H-SUPP) for $\Decl^{\mathrm{ref}}$, whatever its measured
  family;
\item where the hypotheses of Corollary~\ref{cor:fo-keylength} hold for
  $\Decl^{\mathrm{ref}}$, its conclusions hold with
  $\hlb[\Decl^{\mathrm{ref}}]$ in \eqref{eq:keylength-fo}: the key length so
  formed is secret and correct for the protocol whose declaration is
  $\Decl^{\mathrm{ref}}$.
\end{enumerate}
\end{corollary}
\end{mainclaim}
\claimmeta{\metalabel{Uses}\mref{Def.}{def:admissible-declaration}, \mref{Def.}{def:admissible-acceptance-set}, \mref{Def.}{def:joint-sectors}, \mref{Lem.}{lem:branch-support}, \mref{Prop.}{prop:reference-companion}, \mref{Def.}{def:relaxed-feasible-set}, \mref{Def.}{def:first-order-setting}, \mref{Thm.}{thm:firstorder}, \mref{Cor.}{cor:fo-keylength}
\metasep \metalabel{Used by}\mref{Cor.}{cor:certified-keylength-reference}, \mref{Rem.}{rem:fo-arithmetic}, \mref{Req.}{req:bracket-on-reference}
\metasep newly derived
\metasep accepted}

\begin{proof}
(i) By Proposition~\ref{prop:reference-companion}\,(ii),
$\Decl^{\mathrm{ref}}$ is a tuple of the form of
Definition~\ref{def:admissible-declaration} that satisfies its conditions
(i)--(vi) and (viii). Where it satisfies (vii) as well it is an admissible
declaration, and Theorem~\ref{thm:firstorder}\,(iiib) applies to it with
the given points and pairs, with its part~(iv) where lower bounds stand in,
whose hypotheses are exactly those listed: membership in the relaxed set of
$\Decl^{\mathrm{ref}}$, (H-POS), (H-SUPP) and (H-NOPERT) for its objects,
and \eqref{eq:lmo}, or \eqref{eq:lmo-bounds} where a lower bound stands in,
over the exact sets of $\Decl^{\mathrm{ref}}$. The theorem reads of the
points nothing but these hypotheses, so it holds whatever computation
produced them. Where (vii) fails, Theorem~\ref{thm:firstorder}\,(v)
gives the same conclusion, $\hlow[\Decl^{\mathrm{ref}}]=+\infty$. The second
inequality is Proposition~\ref{prop:reference-companion}\,(iii).

(ii) The Choi variable of a declaration in the general format is a list of
positive semidefinite blocks over its joint sectors, of dimension the
product of the Alice sector's and the measured sector's
(Definition~\ref{def:joint-sectors}); these are shared by
Proposition~\ref{prop:reference-companion}\,(i). The relaxed set of
Definition~\ref{def:relaxed-feasible-set} is cut out by per-block positivity,
the coupling condition and the declared entries; the first two are shared,
and the entries differ only in those of the replaced links
(Proposition~\ref{prop:reference-companion}\,(i)), so where none of those
links has any the two sets coincide. By
Lemma~\ref{lem:branch-support}\,(i), $J_{t,s}\succ0$ gives every
key-active branch operator its maximal support, for any declaration, which
is (H-SUPP).

(iii) Corollary~\ref{cor:fo-keylength} holds for any $\hlb$ of
Theorem~\ref{thm:firstorder}\,(ii) or (iiib), with the lower bounds of its
part~(iv) where they stand in, under its hypotheses, and (i) gives one for
$\Decl^{\mathrm{ref}}$.
\end{proof}

% claim meat-core/frank-wolfe/rem:fo-general-format
\claimsummary{meat-core/frank-wolfe/rem:fo-general-format}{The first-order theorem holds as written in the general format, since no step of its proof reads a constraint row, the number of measured factors or the number of Alice sectors. In trusted-source mode the linear minimisation splits into one conic program per Alice sector.}
\claimbookmark{meat-core/frank-wolfe/rem:fo-general-format}{Remark}{rem:fo-general-format}{The first-order bound in the general format}
\begin{remark}[The first-order bound in the general format]\label{rem:fo-general-format}
\emph{(1) Theorem~\ref{thm:firstorder} is stated for a declaration in the
general format} of Definition~\ref{def:declaration} as well, with
$\Feas^{J}$ the set of lists over the joint sectors of
Theorem~\ref{thm:choi} and $J_s\coloneqq\bigoplus_{\mathsf t}
J_{(\mathsf t,s)}$ for every sector $s$ of the measured system, and the
last paragraph of its proof checks each step there. That check is made on
the general format itself, since the \exttermref{../protocol-compiler/contract.pdf}{term-basic-format}{basic format} cannot express
the general case: a \exttermref{../protocol-compiler/contract.pdf}{term-party-local}{party-local} entry, a
\exttermref{../protocol-compiler/contract.pdf}{term-transcript-protocol}{transcript protocol} of more than two messages and an
\exttermref{../protocol-compiler/contract.pdf}{term-ensemble-form}{ensemble form} measurement slot are each outside it, so no appeal to
``the same declaration with its factor tagged $A$ forgotten'' would carry
the general case, and none is made.
\emph{What Proposition~\ref{prop:a-neutral} does and does not contribute
here.} It says that restricting the variable to states block diagonal along
Alice's sectors costs nothing, which is why the joint-sector program is the
right program to run and not a weaker one; it is not what carries the
theorem, and the inclusion of the joint-sector Choi set into a Choi set of
the $A$-forgotten declaration is not used, nor is it available in general,
that declaration's Choi feasible set carrying only
\exttermref{../protocol-compiler/contract.pdf}{term-lifted-operator}{lifted-operator} rows and no party-local row. In the
pairing $\langle\lmoJ_t,S_t-J_t\rangle$ only the blocks of the gradient on
the joint sectors enter, $S_t$ and $J_t$ being lists over them. (H-SUPP) is
asked of $\cG^{J}_{\hat c,s}(J_s)$ on the whole sector $s$, and
$J_{(\mathsf t,s)}\succ0$ for every $\mathsf t$ suffices for it, since then
$J_s\succ0$. The weaker hypothesis per joint pair, which does not ask
anything of key-inert joint pairs inside a key-active sector, is
Corollary~\ref{cor:fo-dropped} at $\mathcal{D}=\varnothing$.

\emph{(2) The linear minimisation splits.} In trusted-source mode
$\Feas^{J}=\prod_{\mathsf t}\Feas^{J}_{\mathsf t}$
(Theorem~\ref{thm:choi}, general format), so
$\min_{J'\in\Feas^{J}}\langle\lmoJ_t,J'\rangle
=\sum_{\mathsf t}\min_{J'_{\mathsf t}\in\Feas^{J}_{\mathsf t}}
\langle\lmoJ_{t,\mathsf t},J'_{\mathsf t}\rangle$: one independent conic
program per Alice sector, each with its own coupling constraint of size
$d_{A_{0,\mathsf t}}$. A lower bound on each sector's minimum, summed over
the sectors, is therefore a lower bound on the minimum over the product,
which is how Lemma~\ref{lem:feasible-side-dual}\,(ii) forms the feasible
side of a bracket, one bound per Alice sector; and a pairing at most each
sector's minimum, summed, is at most the minimum over the product, which is
the second inequality of \eqref{eq:lmo}.
In untrusted-source mode the single normalisation and every scalar entry
couple all joint sectors and nothing splits. What does not split in either
mode is the iteration: the gradient couples all blocks through the scalars
$q(\hat c)/P_{\hat c}$ and $q(\hat c)/K_{\hat c}$.

\emph{(3) The relaxed feasible set per Alice sector.} In the general format
the \termref{term-relaxed-feasible-set}{relaxed feasible set} of Definition~\ref{def:relaxed-feasible-set}
and the recorded fields of Remarks~\ref{rem:fo-interface},
\ref{rem:fo-iterate-checks}, \ref{rem:fo-returned-point},
\ref{rem:fo-one-ruler}, \ref{rem:fo-reachability},
\ref{rem:fo-check-exposures}, \ref{rem:fo-oracle-repair},
\ref{rem:fo-iterate-repair} and~\ref{rem:fo-negative-values} are read with the
coupling constraint per Alice sector, and no statement of
Proposition~\ref{prop:relaxed-set-properties}, of
Remark~\ref{rem:feas-tol} or of those remarks changes;
what a run decides and records per sector is
Requirement~\ref{req:per-sector-coupling}.
\end{remark}
\claimmeta{\metalabel{Uses}\mref{Prop.}{prop:a-neutral}, \mref{Thm.}{thm:choi}, \mref{Lem.}{lem:branch-support}, \mref{Def.}{def:first-order-setting}, \mref{Thm.}{thm:firstorder}
\metasep \metalabel{Used by}\mref{Req.}{req:per-sector-coupling}
\metasep newly derived
\metasep accepted}

\section{The verified feasible side}\label{group-meat-core/frank-wolfe/lem:feasible-side-dual}
\sectionintro{The feasible-side minimum is bounded from below by weak duality: the coupling condition fixes the total trace of every feasible point, so any multipliers, feasible for the dual or not, give a valid bound once the least eigenvalue of the residual operator is bounded from below (\hyperlink{claim-meat-core/frank-wolfe/lem:feasible-side-dual}{Lemma~\ref*{lem:feasible-side-dual}}). The bound reads no point a solver returned, and with a one-dimensional coupling and no rows it is the minimum itself. Rows built with a computed frame in place of the exact one cost a deduction linear in the row multipliers and in the frame's error, after which the bound holds over the exact set again.}

% claim meat-core/frank-wolfe/lem:feasible-side-dual
\claimsummary{meat-core/frank-wolfe/lem:feasible-side-dual}{Any multipliers for the coupling condition and the declared rows give a lower bound on a linear form's minimum over the exact Choi feasible set, dual-feasible or not, since the coupling fixes every feasible point's total trace. Rows built from an inexact frame cost an explicit deduction.}
\claimbookmark{meat-core/frank-wolfe/lem:feasible-side-dual}{Lemma}{lem:feasible-side-dual}{A lower bound on the feasible-side minimum from any multipliers}
\claimleadin{%
Part~(iv) of Theorem~\ref{thm:firstorder} admits any real number below the
feasible-side minimum $\min_{J'\in\Feas^{J}}\langle\lmoJ_t,J'\rangle$ in
place of the pairing with a point an oracle returned. The lemma supplies one
by weak duality, from multipliers for the coupling condition and for the
declared rows. What makes it hold for \emph{any} multipliers, and not only at
a feasible point of a dual program, is that the coupling condition fixes the
total trace of every point of the \exttermref{../protocol-compiler/contract.pdf}{term-feasible-set}{feasible set} exactly: the least
eigenvalue of the residual operator, however negative, then enters with a
known coefficient. Part~(i) states this for an abstract set cut out by a
coupling condition and affine operator rows, part~(ii) reads the Choi
feasible set of Theorem~\ref{thm:choi} in that form, and part~(iii) shows
that the bound is attained on a set with a one-dimensional coupling and no
rows. A run can form the rows only with a computed frame, which in general
differs from the exact $W_{\mathsf t}$ of Theorem~\ref{thm:choi}; part~(iv)
says what that costs: a deduction linear in the row multipliers and in a
bound on the frame's error, after which the bound holds over the exact set
again.
}
\begin{lemma}[A lower bound on the feasible-side minimum from any multipliers]\label{lem:feasible-side-dual}
\begin{enumerate}[label=(\roman*),leftmargin=2.2em]
\item \emph{(Weak duality with a trace identity.)} Let
  $\mathcal{H}_{\mathrm{in}}$ be a Hilbert space of finite dimension
  $d_{\mathrm{in}}\ge1$, let $P$ be a finite non-empty set, and for every
  $p\in P$ let $\mathcal{H}_p$ be a finite-dimensional Hilbert space. Let
  $\mathcal{R}$ be a finite, possibly empty, set of \emph{rows}; a row $r$
  consists of a finite-dimensional Hilbert space $\mathcal{K}_r$, a
  Hermitian operator $\Theta_r$ on it, a relation
  $\bowtie_r\in\{\preceq,\succeq,=\}$ and, for every $p\in P$, a linear map
  $L_{r,p}$ from operators on $\mathcal{H}_{\mathrm{in}}\otimes\mathcal{H}_p$
  to operators on $\mathcal{K}_r$ that maps Hermitian operators to Hermitian
  ones. Write $L^{*}_{r,p}$ for its adjoint, the real-linear map from
  Hermitian operators on $\mathcal{K}_r$ to Hermitian operators on
  $\mathcal{H}_{\mathrm{in}}\otimes\mathcal{H}_p$ with
  $\tr[\Lambda\,L_{r,p}(X)]=\tr[L^{*}_{r,p}(\Lambda)\,X]$ for all Hermitian
  $\Lambda$ and $X$. Let $\mathcal{J}$ be the set of lists
  $J=(J_p)_{p\in P}$ of positive semidefinite operators $J_p$ on
  $\mathcal{H}_{\mathrm{in}}\otimes\mathcal{H}_p$ with
  \begin{equation}\label{eq:abstract-feasible}
    \sum_{p\in P}\tr_{\mathcal{H}_p}J_p=\id_{\mathcal{H}_{\mathrm{in}}}
    \qquad\text{and}\qquad
    L_r(J)\coloneqq\sum_{p\in P}L_{r,p}(J_p)\ \bowtie_r\ \Theta_r
    \quad\text{for every }r\in\mathcal{R}.
  \end{equation}
  Let $G=(G_p)_{p\in P}$ be a list of Hermitian operators $G_p$ on
  $\mathcal{H}_{\mathrm{in}}\otimes\mathcal{H}_p$, let $Y$ be a Hermitian
  operator on $\mathcal{H}_{\mathrm{in}}$, and for every row $r$ let
  $\Lambda_r$ be a Hermitian operator on $\mathcal{K}_r$, with
  $\Lambda_r\succeq0$ when $\bowtie_r$ is $\preceq$ or $\succeq$. Put
  $\sigma_r\coloneqq-1$ when $\bowtie_r$ is $\preceq$ and
  $\sigma_r\coloneqq+1$ otherwise,
  \begin{equation}\label{eq:dual-residual}
    R_p\;\coloneqq\;G_p-Y\otimes\id_{\mathcal{H}_p}
    -\sum_{r\in\mathcal{R}}\sigma_r\,L^{*}_{r,p}(\Lambda_r)
    \qquad(p\in P),
  \end{equation}
  and let $\lambda$ be a real number with $\lambda\le\lmin(R_p)$ for every
  $p\in P$. Then for every $J\in\mathcal{J}$
  \begin{equation}\label{eq:feasible-side-dual}
    \langle G,J\rangle\coloneqq\sum_{p\in P}\tr[G_pJ_p]
    \;\ge\;b(G;Y,\Lambda,\lambda)\coloneqq
    d_{\mathrm{in}}\,\lambda+\tr Y
    +\sum_{r\in\mathcal{R}}\sigma_r\tr[\Lambda_r\Theta_r],
  \end{equation}
  so that $b(G;Y,\Lambda,\lambda)\le\inf_{J\in\mathcal{J}}\langle G,J\rangle$.
  No other condition is placed on $Y$, on the $\Lambda_r$ or on $\lambda$:
  $Y$ and the $\Lambda_r$ need not be feasible for any dual program,
  $\lambda$ may be negative, and $Y=0$ with every $\Lambda_r=0$ is admitted.
\item \emph{(The Choi feasible set is of that form.)} Let $\Decl$ be a
  declaration in the general format of Definition~\ref{def:declaration}, the
  \exttermref{../protocol-compiler/contract.pdf}{term-basic-format}{basic format} included, with its Alice sectors
  $\mathsf t\in\mathbb{T}$.
  \begin{enumerate}[label=(\alph*),leftmargin=2.2em]
  \item In \exttermref{../transmitter/contract.pdf}{term-trusted-source-mode}{trusted-source mode}, with a frame and the factors
    $W_{\mathsf t}$ as in Theorem~\ref{thm:choi}, the factor
    $\Feas^{J}_{\mathsf t}$ of $\Feas^{J}=\prod_{\mathsf t}\Feas^{J}_{\mathsf t}$
    (Theorem~\ref{thm:choi}, the general format) is the set $\mathcal{J}$ of
    (i) with $\mathcal{H}_{\mathrm{in}}=\Hs{A_{0,\mathsf t}}$, $P=\SecSet$,
    $\mathcal{H}_s=\Hs{\Qsys_s}$ and $J_s=J_{(\mathsf t,s)}$, and with the
    following rows. For every entry $k$ read through its
    \exttermref{../protocol-compiler/contract.pdf}{term-lifted-operator}{lifted operator}, one row with $\mathcal{K}=\Hs{A_{\mathsf t}}$,
    $\Theta=\Theta^{(\mathsf t)}_k$, the relation $\bowtie_k$ and
    \begin{equation*}
      L_{k,s}(X)=W_{\mathsf t}\,\tr_{\Qsys_s}\!\big[(\idon{A_{0,\mathsf t}}
      \otimes\hat O^{\uparrow(s)}_k)X\big]\,W_{\mathsf t}^{\dagger},
      \qquad
      L^{*}_{k,s}(\Lambda)=(W_{\mathsf t}^{\dagger}\Lambda W_{\mathsf t})
      \otimes\hat O^{\uparrow(s)}_k .
    \end{equation*}
    For every \exttermref{../protocol-compiler/contract.pdf}{term-party-local}{party-local} entry $k$, scoped to the measured factor $f$,
    and every $s_{f^c}$, one row with
    $\mathcal{K}=\Hs{A_{\mathsf t}}\otimes\Hs{\Qsys_{f^c,s_{f^c}}}$,
    $\Theta=0$, the relation $\bowtie_k$ and, at $s=(s_{f^c},s_f)$,
    \begin{gather*}
      L_{(k,s_{f^c}),s}(X)=(W_{\mathsf t}\otimes\id)\,
      \tr_{\Qsys_{f,s_f}}\!\big[(\id\otimes D^{(s_f)}_k)X\big]\,
      (W_{\mathsf t}^{\dagger}\otimes\id),\\
      L^{*}_{(k,s_{f^c}),s}(\Lambda)=\big((W_{\mathsf t}^{\dagger}\otimes\id)
      \Lambda(W_{\mathsf t}\otimes\id)\big)\otimes D^{(s_f)}_k ,
    \end{gather*}
    where $D^{(s_f)}_k\coloneqq\hat O^{(s_f)}_k-c_k\idon{}$ acts on
    $\Hs{\Qsys_{f,s_f}}$ at the place of that factor in $\Hs{\Qsys_s}$, while
    $L_{(k,s_{f^c}),s'}=0$ and $L^{*}_{(k,s_{f^c}),s'}=0$ at every sector
    $s'$ whose component on the other factors is not $s_{f^c}$. The two
    kinds of map, and the formulas for their adjoints, are defined, and the
    formulas hold, for any linear map $W$ from $\Hs{A_{0,\mathsf t}}$ to
    $\Hs{A_{\mathsf t}}$ in place of $W_{\mathsf t}$; we write
    $L^{W}_{r,s}$ for a row's map so built. A
    homogeneous or party-local row in the frame-free form of
    Proposition~\ref{prop:frame-free-rows}, on a \termref{term-tight}{tight} frame, is a row
    of the same kind with $\idon{A_{0,\mathsf t}}$ in place of
    $W_{\mathsf t}$, hence with $\Hs{A_{0,\mathsf t}}$ in place of
    $\Hs{A_{\mathsf t}}$ in its $\mathcal{K}$, and with
    $\Theta=c_k\idon{A_{0,\mathsf t}}$, respectively $\Theta=0$; and
    $\Feas^{J}_{\mathsf t}$ is unchanged when it replaces the conjugated
    row. Hence, for $Y_{\mathsf t}$, $\Lambda_{\mathsf t}=(\Lambda_r)_r$ and
    $\lambda_{\mathsf t}$ as in (i) for each $\mathsf t$, and with
    $G_{\mathsf t}\coloneqq(G_{(\mathsf t,s)})_{s\in\SecSet}$,
    \begin{equation}\label{eq:feasible-side-sum}
      b_{\mathrm F}\;\coloneqq\;\sum_{\mathsf t\in\mathbb{T}}
      b(G_{\mathsf t};Y_{\mathsf t},\Lambda_{\mathsf t},\lambda_{\mathsf t})
      \;\le\;\langle G,J'\rangle
      \qquad\text{for every }J'\in\Feas^{J}.
    \end{equation}
  \item In \exttermref{../transmitter/contract.pdf}{term-untrusted-source-mode}{untrusted-source mode}, $\Feas^{J}$, which is $\Feas$
    (Definition~\ref{def:first-order-setting}), is the set $\mathcal{J}$ of
    (i) with $\mathcal{H}_{\mathrm{in}}$ one-dimensional, so that $Y=y$ is a
    real number, $Y\otimes\id=y\,\id$, $\tr Y=y$ and $d_{\mathrm{in}}=1$;
    with $P$ the set of \exttermref{../protocol-compiler/contract.pdf}{term-joint-sector}{joint sectors},
    $\mathcal{H}_{(\mathsf t,s)}=\Hs{A_{\mathsf t}\Qsys_s}$ and
    $J_{(\mathsf t,s)}=\omega^{(\mathsf t,s)}$; and with these rows. For every
    scalar entry $k$, one row with $\mathcal{K}=\mathbb{C}$, $\Theta=c_k$,
    the relation $\bowtie_k$,
    $L_{k,(\mathsf t,s)}(X)=\tr[(\idon{A_{\mathsf t}}\otimes
    \hat O^{\uparrow(s)}_k)X]$ and
    $L^{*}_{k,(\mathsf t,s)}(\mu)=\mu\,\idon{A_{\mathsf t}}\otimes
    \hat O^{\uparrow(s)}_k$; for every party-local entry $k$ and every pair
    $(\mathsf t,s_{f^c})$, one row as in (a) with $\idon{A_{\mathsf t}}$ in
    place of $W_{\mathsf t}$, its maps vanishing outside the joint sectors
    $(\mathsf t,(s_{f^c},s_f))$. Hence
    $b(G;y,\Lambda,\lambda)\le\langle G,J'\rangle$ for every
    $J'\in\Feas^{J}$, the whole list being one set of the form of (i).
  \end{enumerate}
\item \emph{(A scalar coupling and no rows.)} If $d_{\mathrm{in}}=1$ and
  $\mathcal{R}=\varnothing$ in (i), then, with $Y=0$ and
  $\lambda=\min_{p\in P}\lmin(G_p)$, the infimum in (i) is attained and
  $b(G;0,\Lambda,\lambda)=\min_{J\in\mathcal{J}}\langle G,J\rangle$.
\item \emph{(Rows built from an inexact frame.)} In trusted-source mode fix
  an Alice sector $\mathsf t$, a linear map $\widetilde W_{\mathsf t}$ from
  $\Hs{A_{0,\mathsf t}}$ to $\Hs{A_{\mathsf t}}$ and a real number
  $\varepsilon_{\mathsf t}\ge\opn{\widetilde W_{\mathsf t}-W_{\mathsf t}}$.
  For every row $r$ of $\Feas^{J}_{\mathsf t}$ in (ii)(a) let
  $\widetilde L_{r,s}\coloneqq L^{\widetilde W_{\mathsf t}}_{r,s}$, a row in
  frame-free form keeping its map, every other datum of the row
  ($\Theta_r$, $\hat O^{\uparrow(s)}_k$, $D^{(s_f)}_k$, $c_k$) unchanged. Let
  $G_{\mathsf t}$, $Y$ and the $\Lambda_r$ be as in (i), form
  $\widetilde R_s$ by \eqref{eq:dual-residual} with
  $\widetilde L^{*}_{r,s}$ in place of $L^{*}_{r,s}$, and let
  $\widetilde\lambda\le\lmin(\widetilde R_s)$ for every $s\in\SecSet$. Then
  for every $J\in\Feas^{J}_{\mathsf t}$, the factor cut out by the rows
  built from the exact $W_{\mathsf t}$,
  \begin{equation}\label{eq:frame-deficit}
    \begin{gathered}
    \langle G_{\mathsf t},J\rangle\;\ge\;
    b(G_{\mathsf t};Y,\Lambda,\widetilde\lambda)-\sum_{r}\kappa_r,\\
    \kappa_r\coloneqq d_{A_{0,\mathsf t}}\,\varepsilon_{\mathsf t}
    \big(2\opn{\widetilde W_{\mathsf t}}+\varepsilon_{\mathsf t}\big)\,
    \opn{\Lambda_r}\,\max_{s\in\SecSet}\opn{O_{r,s}},
    \end{gathered}
  \end{equation}
  where $O_{r,s}$ is $\hat O^{\uparrow(s)}_k$ for a lifted-operator row,
  $D^{(s_f)}_k$ for a party-local row at $s=(s_{f^c},s_f)$ and $0$ at the
  other sectors, and $\kappa_r\coloneqq0$ for a row in frame-free form.
  Summed over $\mathsf t$ as in (ii)(a), the right-hand sides bound
  $\langle G,J'\rangle$ from below for every $J'\in\Feas^{J}$.
\end{enumerate}
\end{lemma}
\claimmeta{\metalabel{Uses}\mref{Def.}{def:declaration}, \mref{Def.}{def:joint-sectors}, \mref{Def.}{def:feasible}, \mref{Thm.}{thm:choi}, \mref{Prop.}{prop:frame-free-rows}, \mref{Def.}{def:untrusted-choi}, \mref{Def.}{def:first-order-setting}
\metasep \metalabel{Used by}\mref{Lem.}{lem:certified-bracket}, \mref{Rem.}{rem:fo-arithmetic}, \mref{Req.}{req:verified-feasible-side}, \mref{Req.}{req:certified-evaluation}
\metasep newly derived
\metasep implemented (effective: accepted)
\metasep \metalabel{Code}\coderef{https://github.com/fraunhofersg/certiqs-key/blob/v0.1.0/src/certiqs_key/solvers/frank_wolfe/verified.py}{certiqs_key.solvers.frank_wolfe.verified:verified_feasible_side}}

\begin{proof}
\emph{Idea.} The linear form is split into the residual paired with $J$ and
two terms the constraints fix or bound: the coupling condition turns the
$Y$-term into $\tr Y$, and each row turns its term into at least
$\sigma_r\tr[\Lambda_r\Theta_r]$. The residual term is at least $\lambda$
times the total trace, which the coupling condition fixes at
$d_{\mathrm{in}}$, so a $\lambda$ of either sign enters with a known
coefficient.

\emph{(i), Step 1 (two identities).} The adjoint exists and is unique: the
Hermitian operators on a finite-dimensional space form a finite-dimensional
real inner product space under $(A,B)\mapsto\tr[AB]$, which is real for
Hermitian $A$ and $B$ since $\overline{\tr[AB]}=\tr[(AB)^{\dagger}]=\tr[BA]$,
and $L_{r,p}$ restricts to a real-linear map between two such spaces. Let
$J\in\mathcal{J}$. For every operator $X$ on
$\mathcal{H}_{\mathrm{in}}\otimes\mathcal{H}_p$,
$\tr[(Y\otimes\id_{\mathcal{H}_p})X]=\tr[Y\,\tr_{\mathcal{H}_p}X]$, the
defining property of the partial trace; summing over $p$ and using the
coupling condition of \eqref{eq:abstract-feasible},
\begin{equation*}
  \sum_{p\in P}\tr[(Y\otimes\id_{\mathcal{H}_p})J_p]
  =\tr\Big[Y\sum_{p\in P}\tr_{\mathcal{H}_p}J_p\Big]=\tr Y .
\end{equation*}
Every $J_p$ is Hermitian, being positive semidefinite, so by the defining
identity of the adjoint, for every row $r$,
$\sum_p\tr[L^{*}_{r,p}(\Lambda_r)J_p]=\sum_p\tr[\Lambda_r\,L_{r,p}(J_p)]
=\tr[\Lambda_r\,L_r(J)]$.

\emph{Step 2 (the rearrangement).} By \eqref{eq:dual-residual},
$G_p=R_p+Y\otimes\id_{\mathcal{H}_p}+\sum_r\sigma_rL^{*}_{r,p}(\Lambda_r)$,
so by Step~1
\begin{equation*}
  \langle G,J\rangle=\sum_{p\in P}\tr[R_pJ_p]+\tr Y
  +\sum_{r\in\mathcal{R}}\sigma_r\tr[\Lambda_r\,L_r(J)] .
\end{equation*}

\emph{Step 3 (the residual term).} For every $p$, $R_p$ is Hermitian, $G_p$,
$Y\otimes\id$ and every $L^{*}_{r,p}(\Lambda_r)$ being so, and
$R_p-\lambda\id$ is positive semidefinite, its eigenvalues being those of
$R_p$ less $\lambda$, which are non-negative by the choice of $\lambda$. For
positive semidefinite $A$ and $B$, $\tr[AB]=\tr[B^{1/2}AB^{1/2}]\ge0$. Hence
$\tr[R_pJ_p]\ge\lambda\tr J_p$. The trace of $J_p$ is that of
$\tr_{\mathcal{H}_p}J_p$, so by the coupling condition
$\sum_p\tr J_p=\tr\id_{\mathcal{H}_{\mathrm{in}}}=d_{\mathrm{in}}$ exactly,
and
\begin{equation*}
  \sum_{p\in P}\tr[R_pJ_p]\;\ge\;\lambda\sum_{p\in P}\tr J_p
  \;=\;d_{\mathrm{in}}\,\lambda ,
\end{equation*}
for either sign of $\lambda$, the total trace being an equality and not a
bound.

\emph{Step 4 (the rows).} If $\bowtie_r$ is $\succeq$, then
$L_r(J)-\Theta_r\succeq0$ and $\Lambda_r\succeq0$, so
$\tr[\Lambda_r(L_r(J)-\Theta_r)]\ge0$ by the trace fact of Step~3, which is
$\sigma_r\tr[\Lambda_rL_r(J)]\ge\sigma_r\tr[\Lambda_r\Theta_r]$ with
$\sigma_r=+1$. If $\bowtie_r$ is $\preceq$, then $\Theta_r-L_r(J)\succeq0$
and $\Lambda_r\succeq0$ give $\tr[\Lambda_rL_r(J)]\le\tr[\Lambda_r\Theta_r]$,
which is the same inequality with $\sigma_r=-1$. If $\bowtie_r$ is $=$, then
$L_r(J)=\Theta_r$ and the two sides are equal for every Hermitian
$\Lambda_r$.

Substituting Steps~3 and~4 into Step~2 gives \eqref{eq:feasible-side-dual}.
As $J\in\mathcal{J}$ was arbitrary, $b(G;Y,\Lambda,\lambda)$ is below every
value of $\langle G,\cdot\rangle$ on $\mathcal{J}$ and hence below its
infimum, which is $+\infty$ when $\mathcal{J}$ is empty. No step used a
property of $Y$, of the $\Lambda_r$ or of $\lambda$ beyond those stated in
(i).

\emph{(ii)(a).} By Theorem~\ref{thm:choi} in the general format, $\Feas^{J}$
is the set of lists $(J_{(\mathsf t,s)})$ over the joint sectors with
$J_{(\mathsf t,s)}\in\Pos{A_{0,\mathsf t}\Qsys_s}$, the coupling
constraints $\sum_s\tr_{\Qsys_s}J_{(\mathsf t,s)}=\idon{A_{0,\mathsf t}}$,
for every entry read through its lifted operator the rows
$W_{\mathsf t}\big(\sum_s\tr_{\Qsys_s}[(\idon{A_{0,\mathsf t}}\otimes
\hat O^{\uparrow(s)}_k)J_{(\mathsf t,s)}]\big)W_{\mathsf t}^{\dagger}
\bowtie_k\Theta^{(\mathsf t)}_k$, and for every party-local entry the last
line of \eqref{eq:feas-general} at $J$, conjugated by
$W_{\mathsf t}\otimes\id$, which is one operator condition on
$\Hs{A_{\mathsf t}}\otimes\Hs{\Qsys_{f^c,s_{f^c}}}$ for every
$(\mathsf t,s_{f^c})$; every constraint involves the blocks of one Alice
sector only, and $\Feas^{J}=\prod_{\mathsf t}\Feas^{J}_{\mathsf t}$ with
$\Feas^{J}_{\mathsf t}$ cut out by the constraints of $\mathsf t$. In the
basic format the same holds with one Alice sector, the lifted operator of an
entry being its own operator, and no party-local entry
(\eqref{eq:choi-feas}). These constraints of $\mathsf t$ are
\eqref{eq:abstract-feasible} with the data listed: the row of a
lifted-operator entry is $\sum_sL_{k,s}(J_s)\bowtie_k\Theta^{(\mathsf t)}_k$,
and the row of a party-local entry at $s_{f^c}$ is the sum over $s_f$ of
$L_{(k,s_{f^c}),(s_{f^c},s_f)}$, which is the sum over all $s$, the other
terms being $0$. Each listed map takes Hermitian operators to Hermitian ones:
for a Hermitian $X$ and a Hermitian operator $O$ acting on one tensor
factor $\Qsys'$ alone,
$\big(\tr_{\Qsys'}[(\id\otimes O)X]\big)^{\dagger}=\tr_{\Qsys'}[X(\id\otimes
O)]=\tr_{\Qsys'}[(\id\otimes O)X]$, an operator on the traced factor
alone moving cyclically under the partial trace over that factor, and a
congruence preserves Hermiticity. Here $\hat O^{\uparrow(s)}_k$ is
Hermitian, every $\hat O_k$ being exactly Hermitian per sector block
(Definition~\ref{def:declaration}, exact admissibility), and so is
$D^{(s_f)}_k$, $c_k$ being real. The adjoints: for Hermitian $\Lambda$ on
$\Hs{A_{\mathsf t}}$ and $X$ on $\Hs{A_{0,\mathsf t}\Qsys_s}$, by
cyclicity of the trace and the partial-trace identity of Step~1,
\begin{equation*}
  \tr\big[\Lambda\,W_{\mathsf t}\tr_{\Qsys_s}[(\id\otimes
  \hat O^{\uparrow(s)}_k)X]W_{\mathsf t}^{\dagger}\big]
  =\tr\big[\big((W_{\mathsf t}^{\dagger}\Lambda W_{\mathsf t})\otimes\id\big)
  (\id\otimes\hat O^{\uparrow(s)}_k)X\big]
  =\tr\big[\big((W_{\mathsf t}^{\dagger}\Lambda W_{\mathsf t})\otimes
  \hat O^{\uparrow(s)}_k\big)X\big],
\end{equation*}
and the operator in the last trace is Hermitian; for a party-local row the
same computation holds with $(W_{\mathsf t}^{\dagger}\otimes\id)\Lambda
(W_{\mathsf t}\otimes\id)$ on $\Hs{A_{0,\mathsf t}}\otimes
\Hs{\Qsys_{f^c,s_{f^c}}}$ in place of $W_{\mathsf t}^{\dagger}\Lambda
W_{\mathsf t}$ and $D^{(s_f)}_k$ on $\Hs{\Qsys_{f,s_f}}$ in place of
$\hat O^{\uparrow(s)}_k$. Neither computation, nor the Hermiticity of the
maps, reads any property of $W_{\mathsf t}$ beyond its being linear, so
both hold for any linear $W$ in its place. For the frame-free forms: on a tight frame a
homogeneous row, read per Alice sector, and a party-local row hold for a
list in their frame-free form exactly when they hold in their conjugated
form (Proposition~\ref{prop:frame-free-rows}\,(i) and (iii)), so
$\Feas^{J}_{\mathsf t}$ is unchanged, and the frame-free row is of the
stated kind by the same two computations with the identity in place of
$W_{\mathsf t}$. Finally, for $J'\in\Feas^{J}$ its part
$J'_{\mathsf t}\coloneqq(J'_{(\mathsf t,s)})_s$ lies in
$\Feas^{J}_{\mathsf t}$, and
$\langle G,J'\rangle=\sum_{\mathsf t}\langle G_{\mathsf t},J'_{\mathsf t}\rangle
\ge\sum_{\mathsf t}b(G_{\mathsf t};Y_{\mathsf t},\Lambda_{\mathsf t},
\lambda_{\mathsf t})$ by (i) applied to each factor, which is
\eqref{eq:feasible-side-sum}.

\emph{(ii)(b).} In untrusted-source mode the Choi step is the identity and
$\Feas^{J}$ denotes $\Feas$ (Definitions~\ref{def:untrusted-choi}
and~\ref{def:first-order-setting}). By \eqref{eq:choi-feas-un} in the basic
format and by the lines of \eqref{eq:feas-general} that hold in
untrusted-source mode in the general format (Definition~\ref{def:feasible}),
$\Feas$ is the set of lists of positive semidefinite blocks
$\omega^{(\mathsf t,s)}$ over the joint sectors with
$\sum_{(\mathsf t,s)}\tr\omega^{(\mathsf t,s)}=1$, the scalar rows
$\sum_{(\mathsf t,s)}\tr[(\idon{A_{\mathsf t}}\otimes
\hat O^{\uparrow(s)}_k)\omega^{(\mathsf t,s)}]\bowtie_kc_k$, and the
party-local rows, unconjugated, one for every $(\mathsf t,s_{f^c})$. With
$\mathcal{H}_{\mathrm{in}}$ one-dimensional, $\tr_{\mathcal{H}_p}J_p$ is the
number $\tr J_p$ and $\id_{\mathcal{H}_{\mathrm{in}}}$ is $1$, so the
coupling condition of \eqref{eq:abstract-feasible} is the normalisation.
The scalar row's map is real on Hermitian $X$, a trace of a product of two
Hermitian operators, and
$\mu\tr[(\id\otimes\hat O^{\uparrow(s)}_k)X]
=\tr[(\mu\,\id\otimes\hat O^{\uparrow(s)}_k)X]$ gives its adjoint; the
party-local rows are those of (a) with the identity in place of
$W_{\mathsf t}$. Then (i) applies to the whole list.

\emph{(iii).} With $\mathcal{R}=\varnothing$ and $Y=0$,
\eqref{eq:dual-residual} is $R_p=G_p$, so
$\lambda=\min_p\lmin(G_p)$ is admissible in (i) and
$b(G;0,\Lambda,\lambda)=d_{\mathrm{in}}\lambda=\lambda$ is at most the
infimum. Let $p^{*}$ attain the minimum and let $u$ be a unit eigenvector of
$G_{p^{*}}$ for $\lmin(G_{p^{*}})$. The list with $J_{p^{*}}=uu^{\dagger}$
and $J_p=0$ for $p\ne p^{*}$ is positive semidefinite, and, as
$\mathcal{H}_{\mathrm{in}}$ is one-dimensional, its coupling condition reads
$\sum_p\tr J_p=1$, which holds; there are no rows. So it lies in
$\mathcal{J}$, and $\langle G,J\rangle=u^{\dagger}G_{p^{*}}u=\lambda$. The
infimum is therefore attained and equals $b(G;0,\Lambda,\lambda)$.

\emph{(iv).} Write $W\coloneqq W_{\mathsf t}$,
$\widetilde W\coloneqq\widetilde W_{\mathsf t}$,
$\varepsilon\coloneqq\varepsilon_{\mathsf t}$ and
$d\coloneqq d_{A_{0,\mathsf t}}$, and let $J\in\Feas^{J}_{\mathsf t}$: its
blocks are positive semidefinite, it satisfies the exact coupling condition,
and it satisfies every row built from $W$. \emph{Step 1.} Steps~1 to~3 of
(i), run with the maps $\widetilde L_{r,s}$ in place of $L_{r,s}$, read of
the rows only the adjoint identity, which $\widetilde L_{r,s}$ satisfies by
(ii)(a) for any linear $W$, and of $J$ only positivity and the coupling
condition; they give
$\langle G_{\mathsf t},J\rangle\ge d\,\widetilde\lambda+\tr Y
+\sum_r\sigma_r\tr[\Lambda_r\widetilde L_r(J)]$.
\emph{Step 2 (the exact rows).} For each row put
$\Delta_{r,s}\coloneqq\widetilde L^{*}_{r,s}(\Lambda_r)-L^{*}_{r,s}(\Lambda_r)$,
a Hermitian operator. By the adjoint identities of both maps,
$\tr[\Lambda_r\widetilde L_r(J)]=\tr[\Lambda_rL_r(J)]
+\sum_s\tr[\Delta_{r,s}J_s]$. $J$ satisfies the row built from $W$, so
Step~4 of (i) gives
$\sigma_r\tr[\Lambda_rL_r(J)]\ge\sigma_r\tr[\Lambda_r\Theta_r]$. For a
Hermitian $\Delta$ and a positive semidefinite $J_s$,
$-\opn{\Delta}\id\preceq\Delta\preceq\opn{\Delta}\id$ and the trace fact of
Step~3 give $|\tr[\Delta J_s]|\le\opn{\Delta}\tr J_s$, so, the total trace
being $d$ by the coupling condition,
$\big|\sum_s\tr[\Delta_{r,s}J_s]\big|\le d\max_s\opn{\Delta_{r,s}}$. Hence
$\langle G_{\mathsf t},J\rangle\ge b(G_{\mathsf t};Y,\Lambda,\widetilde\lambda)
-d\sum_r\max_s\opn{\Delta_{r,s}}$.
\emph{Step 3 (the size of $\Delta$).} For a lifted-operator row,
$\Delta_{r,s}=(\widetilde W^{\dagger}\Lambda_r\widetilde W
-W^{\dagger}\Lambda_rW)\otimes\hat O^{\uparrow(s)}_k$ and
$\widetilde W^{\dagger}\Lambda_r\widetilde W-W^{\dagger}\Lambda_rW
=(\widetilde W-W)^{\dagger}\Lambda_r\widetilde W
+W^{\dagger}\Lambda_r(\widetilde W-W)$. The operator norm is
submultiplicative and invariant under the adjoint,
$\opn{\widetilde W-W}\le\varepsilon$ and
$\opn{W}\le\opn{\widetilde W}+\varepsilon$, so the bracket has norm at most
$\varepsilon\opn{\Lambda_r}(2\opn{\widetilde W}+\varepsilon)$; and the
norm of a tensor product is the product of the norms, the singular values
of $A\otimes B$ being the products of those of $A$ and of $B$. So
$\opn{\Delta_{r,s}}\le\varepsilon(2\opn{\widetilde W}+\varepsilon)
\opn{\Lambda_r}\opn{O_{r,s}}$. For a party-local row the same computation
holds with $\widetilde W\otimes\id$ and $W\otimes\id$, whose difference has
the norm of $\widetilde W-W$ and whose norms are those of $\widetilde W$ and
$W$, and with $D^{(s_f)}_k$ in place of $\hat O^{\uparrow(s)}_k$; at a
sector whose component on the other factors is not $s_{f^c}$ both maps
vanish and $\Delta_{r,s}=0$. For a row in frame-free form
$\widetilde L_{r,s}=L_{r,s}$ and $\Delta_{r,s}=0$. Substituting into Step~2
gives \eqref{eq:frame-deficit}. The sum over $\mathsf t$ follows as in
(ii)(a), $J'_{\mathsf t}$ lying in $\Feas^{J}_{\mathsf t}$ for every
$J'\in\Feas^{J}$.
\end{proof}

\section{The certified bracket}\label{group-meat-core/frank-wolfe/lem:certified-bracket}
\sectionintro{The objective is positively homogeneous of degree one, so at a point with positive definite blocks the gradient paired with the point is the objective there, and a bracket is the sum of its two lower bounds alone (\hyperlink{claim-meat-core/frank-wolfe/lem:certified-bracket}{Lemma~\ref*{lem:certified-bracket}}). Those bounds need only enclosures of the gradient: an entrywise one on the acceptance side, and on the feasible side a ball around an exact residual, the frame of the rows enclosed with it. Over a reference companion known through enclosures, entrywise or, for an element of rank below its dimension, through an enclosed factor, one enclosure serves the companion itself where every element the gradient reads has its zero pattern and its rank fixed throughout (\hyperlink{claim-meat-core/frank-wolfe/lem:certified-bracket-enclosed}{Lemma~\ref*{lem:certified-bracket-enclosed}}). The best certified bracket gives the key length as any first-order bound does (\hyperlink{claim-meat-core/frank-wolfe/cor:certified-keylength}{Corollary~\ref*{cor:certified-keylength}}), on a companion for the protocol it declares (\hyperlink{claim-meat-core/frank-wolfe/cor:certified-keylength-reference}{Corollary~\ref*{cor:certified-keylength-reference}}).}
\sectionimports{\extref{../foundations/contract.pdf}{foundations/cor:least-eigenvalue-over-enclosure}, \extref{../foundations/contract.pdf}{foundations/cor:rank-over-factor-enclosure}, \extref{../foundations/contract.pdf}{foundations/def:exact-reference}, \extref{../foundations/contract.pdf}{foundations/def:factor-enclosure}, \extref{../foundations/contract.pdf}{foundations/def:kl}}

% claim meat-core/frank-wolfe/lem:certified-bracket
\claimsummary{meat-core/frank-wolfe/lem:certified-bracket}{At a point with positive definite blocks and positive probabilities the gradient paired with the point is the objective, so the bracket is the sum of its two lower bounds and needs only enclosures of the gradient.}
\claimbookmark{meat-core/frank-wolfe/lem:certified-bracket}{Lemma}{lem:certified-bracket}{The certified bracket from enclosures of the gradient}
\claimleadin{%
Part~(iv) of Theorem~\ref{thm:firstorder} lets lower bounds on the two linear
minima stand in for the pairings with an oracle's return. The objective is
positively homogeneous of degree one
(Lemma~\ref{lem:objective-homogeneous}), so its derivative along the point
itself is its value there, and the bracket loses the objective and the two
pairings with the point: it is the sum of the two lower bounds, part~(i).
What remains to be bounded is the gradient alone, and parts~(ii) and~(iii)
say how an enclosure of it yields the two lower bounds, the feasible one from
a residual known up to a ball around an exact matrix, the frame of the rows
enclosed with it, so that no deduction for the frame arises.
Corollary~\ref{cor:certified-keylength} takes the key length from this
lemma, and Requirement~\ref{req:certified-evaluation} is how a run forms
it. Lemma~\ref{lem:certified-bracket-enclosed} says when one enclosure of
the gradient serves a reference companion known through enclosures.
}
\begin{lemma}[The certified bracket from enclosures of the gradient]\label{lem:certified-bracket}
Let $\alpha\in(1,\infty)$, let $\Decl$ be a tuple of the form of
Definition~\ref{def:declaration}, in the \exttermref{../protocol-compiler/contract.pdf}{term-basic-format}{basic format} or the general
format, that satisfies conditions (i)--(vi) and (viii) of
Definition~\ref{def:admissible-declaration}, and let $\acSet$ satisfy
Definition~\ref{def:admissible-acceptance-set}. Let $x=(\mathbf q,J)$ with
$\mathbf q\in\acSet$ and every block of $J$ positive definite, on every
\exttermref{../protocol-compiler/contract.pdf}{term-joint-sector}{joint sector} in the general format, satisfy (H-POS) of
Theorem~\ref{thm:firstorder}, and let $(\lmoq,\lmoJ)$ be the gradient of
$O^{J}_{\alpha}$ at $x$ assembled by Corollary~\ref{cor:full-gradient},
read with Lemma~\ref{lem:fo-khyp}\,(iii) and, in the general format, on the
joint sectors as Theorem~\ref{thm:firstorder} reads it.
\begin{enumerate}[label=(\roman*),leftmargin=2.4em]
\item \emph{(The Euler form.)}
  \begin{equation}\label{eq:euler-form}
    \langle\lmoq,\mathbf q\rangle+\langle\lmoJ,J\rangle=O^{J}_{\alpha}(x).
  \end{equation}
  Consequently, for real numbers $b^{\mathrm A}$ and $b^{\mathrm F}$ with
  $b^{\mathrm A}\le\min_{\mathbf q'\in\acSet}\langle\lmoq,\mathbf q'\rangle$
  and $b^{\mathrm F}\le\min_{J'\in\Feas^{J}}\langle\lmoJ,J'\rangle$, the
  minimum over the exact $\Feas^{J}$, the bracket of
  Theorem~\ref{thm:firstorder}\,(iv) at $x$ with both pairings replaced is
  $b^{\mathrm A}+b^{\mathrm F}$, and
  \begin{equation}\label{eq:certified-bracket}
    \max\{0,\,b^{\mathrm A}+b^{\mathrm F}\}\;\le\;\hlow ,
  \end{equation}
  in exact real arithmetic.
\item \emph{(The acceptance side from an enclosure.)} If
  $\hat g\in\mathbb{R}^{\Chacc}$ and $\delta\ge0$ satisfy
  $|\lmoq(\hat c)-\hat g(\hat c)|\le\delta$ for every $\hat c\in\Chacc$, then
  $\min_{\mathbf q'\in\acSet}\langle\lmoq,\mathbf q'\rangle\ge
  \min_{\mathbf q'\in\acSet}\langle\hat g,\mathbf q'\rangle-\delta$.
\item \emph{(The feasible side from enclosed residuals.)} For each group of
  Lemma~\ref{lem:feasible-side-dual}\,(ii), an Alice sector $\mathsf t$ in
  \exttermref{../transmitter/contract.pdf}{term-trusted-source-mode}{trusted-source mode}, its rows built from the exact frame
  $W_{\mathsf t}$ of Theorem~\ref{thm:choi}, or the one group in
  \exttermref{../transmitter/contract.pdf}{term-untrusted-source-mode}{untrusted-source mode}, let $Y$ and the $\Lambda_r$ be as in
  Lemma~\ref{lem:feasible-side-dual}\,(i), let $R_p$ be the residual
  \eqref{eq:dual-residual} formed with $\lmoJ$, and for every block $p$ of the
  group let $M_p$ be a Hermitian matrix, $\rho_p\ge0$ with
  $\|R_p-M_p\|_{\mathrm F}\le\rho_p$, the Frobenius norm of
  \xref{foundations/cor:least-eigenvalue-over-enclosure}, and $\ell_p$ a
  lower bound on $\lmin(M_p)$. Then the group's
  $\lambda\coloneqq\min_p(\ell_p-\rho_p)$ satisfies $\lambda\le\lmin(R_p)$
  for every $p$, and the sum over the groups of $b(\lmoJ;Y,\Lambda,\lambda)$
  of \eqref{eq:feasible-side-dual} is a number
  $b^{\mathrm F}\le\min_{J'\in\Feas^{J}}\langle\lmoJ,J'\rangle$.
\item \emph{(The \termref{term-dropped-program}{dropped program}.)} In the setting of
  Corollary~\ref{cor:fo-dropped}, with a \termref{term-dropped-set}{dropped set} $\mathcal{D}$,
  (i)--(iii) hold with $O^{J,\mathcal{D}}_{\alpha}$ in place of
  $O^{J}_{\alpha}$, (H-POS$^{\mathcal{D}}$) in place of (H-POS), its gradient
  of (H-NOPERT$^{\mathcal{D}}$) in place of $(\lmoq,\lmoJ)$, the kept pairs in
  place of the pairs that are not key-inert, and
  $\underline{h}^{\mathcal{D}}_{\alpha}$ in place of $\hlow$; and
  $\underline{h}^{\mathcal{D}}_{\alpha}\le\hlow$.
\end{enumerate}
\end{lemma}
\claimmeta{\metalabel{Uses}\mref{Def.}{def:declaration}, \mref{Def.}{def:admissible-declaration}, \mref{Def.}{def:admissible-acceptance-set}, \mref{Thm.}{thm:branch-formula}, \mref{Def.}{def:keyinert}, \mref{Def.}{def:keyactive}, \mref{Def.}{def:joint-sectors}, \mref{Lem.}{lem:joint-sectors}, \mref{Thm.}{thm:choi}, \mref{Lem.}{lem:branch-support}, \mref{Thm.}{thm:gradient}, \mref{Cor.}{cor:full-gradient}, \mref{Def.}{def:dropped}, \mref{Prop.}{prop:drop-bound}, \mref{Lem.}{lem:objective-homogeneous}, \mref{Def.}{def:relaxed-feasible-set}, \mref{Def.}{def:first-order-setting}, \mref{Lem.}{lem:fo-khyp}, \mref{Thm.}{thm:firstorder}, \mref{Cor.}{cor:fo-dropped}, \mref{Lem.}{lem:feasible-side-dual}, \extref{../foundations/contract.pdf}{foundations/cor:least-eigenvalue-over-enclosure}, \extref{../foundations/contract.pdf}{foundations/def:kl}
\metasep \metalabel{Used by}\mref{Lem.}{lem:certified-bracket-enclosed}, \mref{Cor.}{cor:certified-keylength}, \mref{Cor.}{cor:certified-keylength-reference}, \mref{Rem.}{rem:certified-scope}, \mref{Req.}{req:certified-evaluation}
\metasep newly derived
\metasep implemented (effective: accepted)
\metasep \metalabel{Code}\coderef{https://github.com/fraunhofersg/certiqs-key/blob/v0.1.0/src/certiqs_key/solvers/frank_wolfe/certified.py}{certiqs_key.solvers.frank_wolfe.certified:certified_bracket}}

\begin{proof}
\emph{Idea.} Homogeneity makes the derivative along the point equal to the
value, and the gradient's formulas give the same derivative as the pairing
with the point; the two lower bounds then carry the whole bracket. An
entrywise enclosure of the acceptance-side gradient moves its minimum over a
set of distributions by at most the radius, and a ball around the residual
moves its least eigenvalue by at most the ball's radius.

\emph{(i), Step 1 (the point is admissible to the theorem).} Every block of
$J$ is positive semidefinite, and each residual of
Definition~\ref{def:relaxed-feasible-set} at $J$, of the coupling condition,
of the normalisation and of every declared row, is a finite number, so $J$
lies in $\Feas^{J}(\varepsilon)$ for every $\varepsilon$ at least each of
them, the normalisation's divided by the dimension factor its clause
carries; with
$\mathbf q\in\acSet$, $x\in\acSet\times\Feas^{J}(\varepsilon)$. Every block
being positive definite, every key-active branch operator attains its
\termref{term-maximal-branch-support}{maximal branch support} (Lemma~\ref{lem:branch-support}\,(i), per joint
sector in the general format), which is (H-SUPP); and $(\lmoq,\lmoJ)$ is the
gradient at $x$ itself, which is (H-NOPERT).

\emph{Step 2 (the value along the ray).} By (H-POS) and
Lemma~\ref{lem:fo-khyp}\,(i), $q(\hat c)>0$ exactly on $\Chacc$, and
$P^{J}_{\hat c}(J)>0$ there, $K^{J}_{\hat c}(J)>0$ on $\Chkey\cap\Chacc$, so
every term of $O^{J}_{\alpha}(x)$ indexed in $\Chacc$ is real, and every
term indexed outside it has $q(\hat c)=0$ and is $0$ by
\xref{foundations/def:kl}, the second arguments being non-negative by the
cone clause of Theorem~\ref{thm:branch-formula}. So $O^{J}_{\alpha}(x)$ is
real. By Lemma~\ref{lem:objective-homogeneous}\,(iii),
$O^{J}_{\alpha}((1+\tau)x)=(1+\tau)\,O^{J}_{\alpha}(x)$ for every
$\tau>-1$, so $\tau\mapsto O^{J}_{\alpha}(x+\tau x)$ is differentiable at
$\tau=0$ with derivative $O^{J}_{\alpha}(x)$.

\emph{Step 3 (the same derivative from the gradient).} The terms indexed
outside $\Chacc$ vanish at every $(1+\tau)x$, $\tau>-1$, by Step~2 applied to
the scaled point, $(1+\tau)q(\hat c)=0$ there; they contribute $0$ to the
derivative, and the pairing $\langle\lmoq,\mathbf q\rangle$ of
Definition~\ref{def:first-order-setting} and the sums of \eqref{eq:grad-J}
read with Lemma~\ref{lem:fo-khyp}\,(iii) omit them. For $\hat c\in\Chacc$
the term is a smooth function of $(q(\hat c),P^{J}_{\hat c},K^{J}_{\hat c})$
near the point, these being positive at $x$ and scaling with $1+\tau$
(Lemma~\ref{lem:objective-homogeneous}\,(i)). Along $\tau$,
$q(\hat c)$ and $P^{J}_{\hat c}$ move linearly, with derivatives $q(\hat c)$
and $P^{J}_{\hat c}(J)=\sum_{s'}\tr[\Gamma^{(s')}_{\hat c}J_{s'}]$
(\eqref{eq:grad-pieces}). Each deficit of a key-active pair is
$f^{J}_{\hat c,s}(J_s+\tau J_s)=T(X+\tau X)-\tr(X+\tau X)$ with
$X\coloneqq\cG^{J}_{\hat c,s}(J_s)$, the \termref{term-branch-map}{branch map} being linear; the
direction $X$ is Hermitian with support $\operatorname{supp}X$, so
Theorem~\ref{thm:gradient}\,(i) gives the derivative
$\tr[\nabla T(X)\,X]-\tr X=\tr\big[(\nabla T(X)-\id)\,\cG^{J}_{\hat c,s}(J_s)\big]
=\tr\big[\chi^{\dagger}_{\Qsys_s}\big(\cG^{\dagger}_{\hat c,s}[\nabla T(X)-\id]\big)
J_s\big]$, which is $\tr[\Delta^{(s)}_{\hat c}J_s]$ of \eqref{eq:grad-pieces};
the deficits of key-inert pairs are identically $0$
(Definition~\ref{def:keyactive}, Lemma~\ref{lem:joint-sectors}\,(iii)).
The chain rule through $u\log_2\frac uv$ at $u,v>0$, exactly as in the proof
of Corollary~\ref{cor:full-gradient}, then gives the derivative of the term as
its coefficients of \eqref{eq:grad-q} and \eqref{eq:grad-J} paired with
$q(\hat c)$ and with the blocks of $J$, and summing over $\Chacc$, the
derivative of $\tau\mapsto O^{J}_{\alpha}(x+\tau x)$ at $0$ is
$\langle\lmoq,\mathbf q\rangle+\langle\lmoJ,J\rangle$. With Step~2 this is
\eqref{eq:euler-form}.

\emph{Step 4 (the bracket).} By Step~1, Theorem~\ref{thm:firstorder}\,(iiib)
with its part~(iv) applies to the one point $x$ with both pairings replaced,
where $\Decl$ satisfies condition~(vii) as well; its bracket is
$O^{J}_{\alpha}(x)+b^{\mathrm A}-\langle\lmoq,\mathbf q\rangle
+b^{\mathrm F}-\langle\lmoJ,J\rangle=b^{\mathrm A}+b^{\mathrm F}$ by
\eqref{eq:euler-form}, and the theorem gives \eqref{eq:certified-bracket}.
Where condition~(vii) fails, Theorem~\ref{thm:firstorder}\,(v) gives the same,
$\hlow=+\infty$.

\emph{(ii).} Let $\mathbf q'\in\acSet$, a non-empty set as it contains
$\mathbf q$. It is a probability vector vanishing outside $\Chacc$
(Lemma~\ref{lem:fo-khyp}\,(i)), so $q'(\hat c)\ge0$ and
$\sum_{\hat c\in\Chacc}q'(\hat c)=1$. Hence
$\langle\lmoq,\mathbf q'\rangle=\langle\hat g,\mathbf q'\rangle
+\sum_{\hat c\in\Chacc}(\lmoq(\hat c)-\hat g(\hat c))q'(\hat c)
\ge\langle\hat g,\mathbf q'\rangle-\delta
\ge\min_{\mathbf q''\in\acSet}\langle\hat g,\mathbf q''\rangle-\delta$, the
minimum attained, $\acSet$ being compact
(Definition~\ref{def:admissible-acceptance-set}) and the map linear. The
minimum over $\mathbf q'$ of the left-hand side is therefore at least the
right-hand side.

\emph{(iii).} $R_p$ is Hermitian, $\lmoJ$, $Y\otimes\id$ and every
$L^{*}_{r,p}(\Lambda_r)$ being so
(Lemma~\ref{lem:feasible-side-dual}\,(i), (ii)). By
\xref{foundations/cor:least-eigenvalue-over-enclosure}\,(i),
$\lmin(R_p)\ge\lmin(M_p)-\rho_p\ge\ell_p-\rho_p\ge\lambda$. So $\lambda$ is
admissible in Lemma~\ref{lem:feasible-side-dual}\,(i) for the group, whose
rows, built from the exact frame, cut out the factor of the exact
$\Feas^{J}$ (part~(ii)(a) of that lemma; part~(ii)(b) in untrusted-source
mode), and the sum over the groups bounds $\langle\lmoJ,J'\rangle$ from
below for every $J'\in\Feas^{J}$, which is
\eqref{eq:feasible-side-sum}.

\emph{(iv).} $O^{J,\mathcal{D}}_{\alpha}$ is positively homogeneous of degree
one (Lemma~\ref{lem:objective-homogeneous}\,(iii)), and Steps~2 and~3 run
unchanged with $K^{J,\mathcal{D}}_{\hat c}$ in place of $K^{J}_{\hat c}$
and the deficits of the kept pairs only, the dropped ones being absent from
\eqref{eq:KD}; (H-SUPP$^{\mathcal{D}}$) holds at a positive definite block as
in Step~1. Step~4 runs with Corollary~\ref{cor:fo-dropped}, whose
conclusion \eqref{eq:hlb-dropped} admits the lower bounds of
Theorem~\ref{thm:firstorder}\,(iv), in place of the theorem. Parts (ii) and
(iii) do not read the objective. Finally
$\underline{h}^{\mathcal{D}}_{\alpha}\le\hlow$ is
Proposition~\ref{prop:drop-bound}.
\end{proof}

% claim meat-core/frank-wolfe/lem:certified-bracket-enclosed
\claimsummary{meat-core/frank-wolfe/lem:certified-bracket-enclosed}{Over a reference companion whose measured families, of one or more replaced factors, are known through entrywise enclosures and on some blocks enclosed factors, one enclosure of the gradient serves the companion where every element the gradient reads has its zero pattern and rank fixed.}
\claimbookmark{meat-core/frank-wolfe/lem:certified-bracket-enclosed}{Lemma}{lem:certified-bracket-enclosed}{The certified bracket over a measured family known through enclosures}
\claimleadin{%
Lemma~\ref{lem:certified-bracket} needs, at a point, only enclosures of the
gradient. A reference companion known through enclosures leaves a whole set of
families consistent with its data, the companion's own among them, and this
lemma says when one enclosure of the gradient serves every one of them: when
the zero pattern and the rank of every element that the gradient's formula
reads are the same throughout. The companion may replace the measured
families of several factors, as it replaces both stations of an
entanglement-based declaration, and each of them may be known through
enclosures, so that an element of the product family may be a tensor product
of several blocks that are only enclosed. An element of rank below its
dimension whose support no coordinate vectors span has its rank fixed by no
entrywise enclosure with positive radii, and is read instead through an
enclosed factor of fixed rank. A block may so be
known twice, and the statistics may read it through the one enclosure and the
deficits through the other; part~(c) says that an enclosure of the gradient
formed in that way, over every pair of the two readings, serves the companion,
whose own family is read twice among those pairs. Part~(b) gives the factor
structure that keeps the factors of a \termref{term-pass-class}{pass class} of known rank where
some of them are only enclosed: a tensor product of factors of full row rank
has full row rank, the tensor product of their Gram matrices being positive
definite, so blocks of fixed rank multiply into elements of fixed rank.
Elements whose blocks known only through boxes are the same blocks have
their rows grouped by those blocks, named as blocks and never compared by
their boxes, so that an exact dependence between their exact factors is
reduced exactly.
Corollary~\ref{cor:certified-keylength-reference} takes the key length from
the lemma.
}
\begin{lemma}[The certified bracket over a measured family known through enclosures]\label{lem:certified-bracket-enclosed}
Let $\Decl^{\mathrm{ref}}$ be a compiled reference companion of
Proposition~\ref{prop:reference-companion}, with $\mathbb{F}^{\mathrm{ref}}$
the set of the \exttermref{../protocol-compiler/contract.pdf}{term-measured-factor}{measured factors} it replaces (the one
measured system in the basic format), let $\alpha$, $\acSet$ and
$x=(\mathbf q,J)$ be as in Lemma~\ref{lem:certified-bracket} with
$\Decl\coloneqq\Decl^{\mathrm{ref}}$, save that (H-POS) of
Theorem~\ref{thm:firstorder} is not assumed of $x$, and, where (H-POS) holds
for $\Decl^{\mathrm{ref}}$ at $x$, let $(\lmoq,\lmoJ)$ be the gradient that
lemma reads there, the blocks of the measured families of the factors of
$\mathbb{F}^{\mathrm{ref}}$ each known through an \exttermref{../foundations/contract.pdf}{term-entrywise-enclosure}{entrywise enclosure}
(\xref{foundations/def:exact-reference}), some of them read, besides, with a
triple $(L_c,t,W)$ of a \exttermref{../foundations/contract.pdf}{term-matrix-box}{matrix box} and a core as in
\xref{foundations/def:factor-enclosure}, the block's \emph{factor
enclosure}, and every other datum exactly. Call a tuple
$F=(F^{f})_{f\in\mathbb{F}^{\mathrm{ref}}}$ of families, each $F^{f}$
indexed as the measured family of factor $f$ is, a \emph{family of the
replaced factors}, its elements $F^{(s)}_v$ being the products
\eqref{eq:productfamily} of its blocks and of the blocks of the other factors
of $\Decl^{\mathrm{ref}}$; call it an \emph{entrywise member} when each of
its blocks is a Hermitian matrix that the entrywise enclosure of that block
contains, and a \emph{member} when, moreover, each of its blocks that is read
with a factor enclosure is a member of it. Suppose
\begin{itemize}[leftmargin=4.6em]
\item[(H-FAC)] each block of the companion's own family of the replaced
  factors that is read with a factor enclosure is a member of it, the triple
  being a factor enclosure of that block, so that the companion's family of
  the replaced factors is a member.
\end{itemize}
Call a block \emph{decided} when every entry of its entrywise enclosure has
centre and radius $0$; or some entry's centre $C_{il}$ and radius $r_{il}$
satisfy $|C_{il}|>r_{il}$; or its factor enclosure has no column; or its
factor enclosure has $k\ge1$ columns and there are $G$, $\rho$ and $\ell$ as
in \xref{foundations/cor:rank-over-factor-enclosure}\,(ii) with
$\ell-\rho>0$. Call a block \emph{of fixed rank} when it is of one of three
kinds: of the first kind when it is zero for every member; of the second kind
when there are a matrix $E$ of full row rank and, for every member, a
positive definite $\widetilde X$ with the member's block equal to
$E^{\dagger}\widetilde XE$, $E$ the same for every member
(\xref{foundations/cor:least-eigenvalue-over-enclosure}\,(iv) gives such an
$E$); and of the third kind when there are an integer $k\ge1$, a matrix box
of complex $k\times n$ matrices, $n$ the dimension of the block, and a
positive definite $k\times k$ matrix $\widetilde X$ and, for every member, a
matrix $E_{*}$ of rank $k$ that the box contains, with the member's block
equal to $E_{*}^{\dagger}\widetilde XE_{*}$
(\xref{foundations/cor:rank-over-factor-enclosure}\,(ii) gives such a box,
that of the adjoints of the factors, with $\widetilde X$ the core). Suppose
that every block of a factor of $\mathbb{F}^{\mathrm{ref}}$ that is a factor
of an element $F^{(s)}_v$ with $(x',v)$ an outcome pair of a \termref{term-pass-class}{pass class}
of some $\hat c\in\Chkey$ and $M^{(\mathsf t)}_{x'}\ne0$, no other block of
which is zero for every member, is decided, and that every such block is of
fixed rank at the pairs $(\hat c,(\mathsf t,s))$ that are not
\termref{term-key-inert}{key-inert}. Then:
  \begin{enumerate}[label=(\alph*),leftmargin=2.2em]
  \item every member gives the same
    \termref{term-effective-outcome-set}{effective outcome sets} on every joint
    sector, the same key-inert pairs and the same
    \termref{term-joint-key-active-index-set}{joint key-active index sets}, those
    read off the enclosures;
  \item fix, for every block that the hypothesis asks to be of fixed rank,
    one kind it is of, with the data of that kind, and, where that kind is
    the third, at every member one matrix $E_{*}$ as that kind asks, one for
    the block and not one for each element it is a factor of; below, the
    kind of a block is the one so fixed for it. At every pair
    $(\hat c,(\mathsf t,s))$ that is not key-inert, every
    element $F^{(s)}_v$ of an outcome pair $(x',v)$ of a pass class of
    $\hat c$ with $M^{(\mathsf t)}_{x'}\ne0$ is, for every member,
    $E^{(s)\dagger}_v\widetilde F^{(s)}_vE^{(s)}_v$ with
    $\widetilde F^{(s)}_v$ positive definite and $E^{(s)}_v$ of full row
    rank, the case $F^{(s)}_v=0$ included as $E^{(s)}_v$ with no rows, where
    $E^{(s)}_v$ is the same for every member if no block of $F^{(s)}_v$ is of
    the third kind or some block of it is zero for every member, and
    otherwise its rank is the same for every member and it is the tensor
    product, in the order of $\mathbb{F}$, of one matrix for each block of
    $F^{(s)}_v$: the matrix $E_{*}$ fixed for the block at that member where
    the block is of the third kind, and a matrix the same for every member
    otherwise. For an element of the latter case, call its \emph{tuple}
    $\beta(v)$ the list, in the order of $\mathbb{F}$, of its blocks of the
    third kind, each block $F^{f,(s_f)}_{v_f}$ named by the triple
    $(f,s_f,v_f)$, so that two entries of tuples agree exactly when they
    name the same block, whatever their boxes; and for a tuple $\beta$ let
    $E^{\mathrm{box}}_{\beta}$ be the tensor product, in that order, of the
    matrices $E_{*}$ fixed for its blocks, which at each member is one
    matrix for every element whose tuple is $\beta$.
    Moreover, for every pass class of $\hat c$ and any factor forms
    $M^{(\mathsf t)}_{x}=E^{(\mathsf t)\dagger}_{x}\widetilde M_{x}
    E^{(\mathsf t)}_{x}$ of Alice's elements with $E^{(\mathsf t)}_{x}$ of
    full row rank and $\widetilde M_{x}$ positive definite, the matrix
    $\widetilde E$ that stacks the matrices
    $E^{(\mathsf t)}_{x}\otimes E^{(s)}_v$ over the class's outcome pairs
    $(x,v)$ with $M^{(\mathsf t)}_{x}\ne0$, in a fixed order, is, for every
    member, $\widetilde E=\widetilde B\,C$ with $\widetilde B$ of full column
    rank and the same for every member, and $C$ a stack of a matrix $C_0$
    and, one for each tuple $\beta$ of an element of the latter case in the
    class, of a matrix $(C_{\beta}\otimes E^{\mathrm{box}}_{\beta})U_{\beta}$,
    with $C_0$, every $C_{\beta}$ and every permutation matrix $U_{\beta}$
    the same for every member; where the $E^{(\mathsf t)}_{x}$, the
    $E^{(s)}_v$ that are the same for every member and the matrices the same
    for every member of which the other $E^{(s)}_v$ are tensor products have
    entries in $\mathbb{Q}(i)$, $\widetilde B$, $C_0$ and every $C_{\beta}$
    can be taken with entries in $\mathbb{Q}(i)$;
  \item call a pair $(F',F'')$ of families of the replaced factors, the
    first read by the statistics and the second by the deficits,
    \emph{admissible} when every block of $F''$ that is a factor of an
    element $F^{(s)}_v$ of an outcome pair $(x',v)$ of a pass class of
    $\hat c$ with $M^{(\mathsf t)}_{x'}\ne0$, at a pair
    $(\hat c,(\mathsf t,s))$ that is not key-inert, no other block of which
    is zero for every member, is a positive semidefinite matrix. For such a pair form, at $x$, $P^{J}_{\hat c}$ and
    every $\Gamma^{(s')}_{\hat c}$ with $F'$, every deficit
    $f^{J}_{\hat c,s}$ and every $\Delta^{(s')}_{\hat c}$ with $F''$, an
    outcome pair whose element of Alice's is zero, or whose element
    $F^{(s)}_v$ has a block that is zero for every member, contributing
    nothing to a branch map or to the statistics, over the key-inert pairs and the joint key-active index sets
    of (a), and
    $K^{J}_{\hat c}$ as $P^{J}_{\hat c}$ plus those deficits; where
    $q(\hat c)$, $P^{J}_{\hat c}$ and $K^{J}_{\hat c}$ so formed are positive
    as (H-POS) of Theorem~\ref{thm:firstorder} asks, let
    $(\lmoq,\lmoJ)[F',F'']$ be the pair given by \eqref{eq:grad-q} and
    \eqref{eq:grad-J} with these pieces, read as
    Lemma~\ref{lem:certified-bracket} reads them. Let $\mathcal P$ be a set
    of admissible pairs that contains $(F,F)$ for every member $F$, such as
    the set of the admissible pairs of an entrywise member $F'$ and a family
    $F''$ of the replaced factors each of whose blocks lies in its entrywise
    enclosure or is a member of its factor enclosure, the one or the other
    chosen block by block. Then
    a set that contains $P^{J}_{\hat c}$, or $K^{J}_{\hat c}$, formed with
    every pair of $\mathcal P$ contains that of $\Decl^{\mathrm{ref}}$ at
    $x$; and where (H-POS) holds for $\Decl^{\mathrm{ref}}$, enclosures
    $(\hat g,\delta)$ and $(M_p,\rho_p)$ that hold, as in parts (ii) and
    (iii) of Lemma~\ref{lem:certified-bracket}, at
    $(\lmoq,\lmoJ)[F',F'']$ for every pair of $\mathcal P$ at which it is
    formed hold at the gradient of $\Decl^{\mathrm{ref}}$, and parts
    (i)--(iii) of that lemma then give
    $\max\{0,b^{\mathrm A}+b^{\mathrm F}\}\le\hlow[\Decl^{\mathrm{ref}}]$;
  \item in the setting of
    Corollary~\ref{cor:fo-dropped}, with a \termref{term-dropped-set}{dropped set}
    $\mathcal{D}$, (a)--(c) hold, the hypotheses read the same way, with
    $O^{J,\mathcal{D}}_{\alpha}$ in place of $O^{J}_{\alpha}$,
    (H-POS$^{\mathcal{D}}$) in place of (H-POS), its gradient of
    (H-NOPERT$^{\mathcal{D}}$) in place of $(\lmoq,\lmoJ)$ and
    $K^{J,\mathcal{D}}_{\hat c}$ in place of $K^{J}_{\hat c}$, the kept
    pairs in place of the pairs that are not key-inert, part (iv) of
    Lemma~\ref{lem:certified-bracket} in place of its parts (i)--(iii),
    and $\underline{h}^{\mathcal{D}}_{\alpha}$ in place of $\hlow$.
  \end{enumerate}
\end{lemma}
\claimmeta{\metalabel{Uses}\mref{Def.}{def:declaration}, \mref{Def.}{def:admissible-declaration}, \mref{Def.}{def:branch}, \mref{Def.}{def:keyinert}, \mref{Def.}{def:joint-sectors}, \mref{Thm.}{thm:choi}, \mref{Cor.}{cor:full-gradient}, \mref{Def.}{def:dropped}, \mref{Prop.}{prop:reference-companion}, \mref{Thm.}{thm:firstorder}, \mref{Cor.}{cor:fo-dropped}, \mref{Lem.}{lem:certified-bracket}, \extref{../foundations/contract.pdf}{foundations/cor:least-eigenvalue-over-enclosure}, \extref{../foundations/contract.pdf}{foundations/cor:rank-over-factor-enclosure}, \extref{../foundations/contract.pdf}{foundations/def:exact-reference}, \extref{../foundations/contract.pdf}{foundations/def:factor-enclosure}
\metasep \metalabel{Used by}\mref{Cor.}{cor:certified-keylength-reference}, \mref{Rem.}{rem:certified-scope}, \mref{Req.}{req:certified-evaluation}
\metasep newly derived
\metasep implemented (effective: accepted)
\metasep \metalabel{Code}\coderef{https://github.com/fraunhofersg/certiqs-key/blob/v0.1.0/src/certiqs_key/solvers/frank_wolfe/certified.py}{certiqs_key.solvers.frank_wolfe.certified:_companion_element}}

\begin{proof}
\emph{(a).} By Definition~\ref{def:joint-sectors}, an element
$F^{(s)}_v$ of a factorised family is non-zero exactly when each of its
factor blocks is, and the effective outcome set of a pass class on a joint
sector collects the outcome pairs with $M^{(\mathsf t)}_{x'}\ne0$ and
$F^{(s)}_v\ne0$. Alice's blocks and the blocks of the factors outside
$\mathbb{F}^{\mathrm{ref}}$ are the same for every member. A decided block of
a factor of $\mathbb{F}^{\mathrm{ref}}$ is zero for every member when every
entry has centre and radius $0$, since then every entry of a contained matrix
is $0$, and non-zero for every member when some entry has
$|C_{il}|>r_{il}$, since then $|X_{il}|\ge|C_{il}|-r_{il}>0$. When its
factor enclosure has no column, every member of that enclosure is the zero
matrix (\xref{foundations/cor:rank-over-factor-enclosure}\,(i)), and so is
the block of every member, which is one of them; when it has $k\ge1$ columns
with $\ell-\rho>0$, every member of the enclosure has rank $k$ and is not
zero (part (ii) of that corollary), and so the block of every member is not
zero. An element one of whose blocks is zero for every member is zero for
every member, whatever its other blocks, of which the hypothesis therefore
asks nothing at that element. So each effective outcome set is the same for every member, and
key-inertness, which reads only these sets and the raw-key map, and the joint
key-active index sets with it, are the same.

\emph{(b), the factor forms.} Let $F$ be a member. A block of a factor of
$\mathbb{F}^{\mathrm{ref}}$ of fixed rank is, at $F$, by the kind fixed for
it: zero, of the first kind; $E^{\dagger}\widetilde XE$ with $E$ the same
for every member, of the second; or $E_{*}^{\dagger}\widetilde XE_{*}$ with
$E_{*}$ the matrix fixed for the block at $F$, in the box, of rank $k$, and
$\widetilde X$ the same for every member, of the third, so that the block
enters every element it is a factor of through that one $E_{*}$. Every
other factor block $B$, a block of a factor outside
$\mathbb{F}^{\mathrm{ref}}$, is the same for every member and positive
semidefinite, and with $B=\sum_{k'\le m}b_{k'}u_{k'}u_{k'}^{\dagger}$,
$b_{k'}>0$ and $u_{k'}$ orthonormal, it is
$E_B^{\dagger}\operatorname{diag}(b_1,\dots,b_m)E_B$ with $E_B$ the matrix
of rows $u_{k'}^{\dagger}$, of full row rank. A tensor product of such forms
is $(\bigotimes E)^{\dagger}(\bigotimes\widetilde X)(\bigotimes E)$, where
$\bigotimes E$ has full row rank, the tensor product of injective adjoints
being injective: its Gram matrix
$(\bigotimes E)(\bigotimes E)^{\dagger}=\bigotimes(EE^{\dagger})$ is a tensor
product of positive definite matrices, each $EE^{\dagger}$ being positive
definite for a matrix $E$ of full row rank, and so positive definite, and its
rank, its number of rows, is the product of the ranks of the factors; and
$\bigotimes\widetilde X$ is positive definite. A positive multiple
$wE^{\dagger}\widetilde XE=E^{\dagger}(w\widetilde X)E$ keeps the form, and a
zero weight or a zero factor gives the zero element; an element one of whose
blocks is zero for every member is so zero for every member, its
$E^{(s)}_v$ with no rows the same for every member, whatever its other
blocks. The weight and every
factor of $\bigotimes E$ but those of the blocks of the third kind are the
same for every member, $E$ being the same for every member at a block of the
second kind and $E_B$ at a block outside $\mathbb{F}^{\mathrm{ref}}$. So,
where no block of the element is zero for every member, $\bigotimes E$, taken
in the order of $\mathbb{F}$, is the same for every member where no block of
the element is of the third kind, and otherwise it is the tensor product of
the matrices $E_{*}$ fixed for its blocks of the third kind and of matrices
the same for every member, one per block, as the statement says, and its
rank is the same for every member: the product of the ranks of its factors,
each the same for every member, a block of the third kind contributing its
$k$ at every member, whichever other blocks of the third kind stand beside
it.

\emph{(b), the reordering.} Let $v$ be an element of the latter case, with
tuple $\beta=\beta(v)$, let $f_1,\dots,f_m$ be the factors of $\mathbb{F}$ in
their order, and let $E^{(s)}_v=E_{v,1}\otimes\dots\otimes E_{v,m}$, with
$E_{v,j}$ the matrix of its block of factor $f_j$, of $p_j$ rows and $n_j$
columns; let $\operatorname{pos}\beta\subseteq\{1,\dots,m\}$, non-empty, be
the positions of the blocks of $\beta$. Each $p_j$ is the same for every
member: the $k$ of the block where it is of the third kind, and otherwise
the number of rows of a matrix the same for every member. Each $n_j$ is the
block dimension $d_{f_j,s_{f_j}}$ of the joint sector $s$, the same for
every element of the class. For positive integers $q_1,\dots,q_m$ let
$\Sigma_{\beta}(q)$ be the linear map with
$\Sigma_{\beta}(q)(u_1\otimes\dots\otimes u_m)=u_{\sigma(1)}\otimes\dots
\otimes u_{\sigma(m)}$ for $u_j\in\mathbb{C}^{q_j}$, where $\sigma$ lists the
positions outside $\operatorname{pos}\beta$ in increasing order and then
those in $\operatorname{pos}\beta$ in increasing order. Defined on the
tensor products of standard basis vectors, which it maps bijectively onto
those of $\mathbb{C}^{q_{\sigma(1)}}\otimes\dots\otimes\mathbb{C}^{q_{\sigma(m)}}$,
and extended linearly, it is a permutation matrix, so that
$\Sigma_{\beta}(q)^{-1}=\Sigma_{\beta}(q)^{\top}$, and the identity holds for
all vectors $u_j$, both sides being multilinear in them. For matrices $A_j$
of size $p_j\times n_j$, both sides applied to $u_1\otimes\dots\otimes u_m$
give $\Sigma_{\beta}(p)(A_1\otimes\dots\otimes A_m)=(A_{\sigma(1)}\otimes
\dots\otimes A_{\sigma(m)})\Sigma_{\beta}(n)$, and the Kronecker product
being associative,
\begin{equation*}
  E^{(s)}_v=R_v\big(E^{\mathrm{ex}}_v\otimes E^{\mathrm{box}}_{\beta}\big)
  \Sigma_{\beta}(n),\qquad R_v\coloneqq\Sigma_{\beta}(p)^{\top},
\end{equation*}
with $E^{\mathrm{ex}}_v$ the tensor product, in order, of the $E_{v,j}$ with
$j\notin\operatorname{pos}\beta$, the $1\times1$ matrix $1$ where there is
none, the same for every member, and $E^{\mathrm{box}}_{\beta}$ that of the
$E_{v,j}$ with $j\in\operatorname{pos}\beta$, the matrices $E_{*}$ fixed for
the blocks of $\beta$, as the statement defines it. $R_v$ reads only
$\operatorname{pos}\beta$ and the $p_j$, and $\Sigma_{\beta}(n)$ only
$\operatorname{pos}\beta$ and the $n_j$; so $R_v$ is the same for every
member, and $\Sigma_{\beta}(n)$ is the same for every member and for every
element of the class whose tuple is $\beta$. At a member, every element whose
tuple is $\beta$ has, at each block of $\beta$, the one $E_{*}$ fixed for that
block, so $E^{\mathrm{box}}_{\beta}$ is one matrix for all of them. This is
where $E_{*}$ is fixed per block and not per element, and where a tuple names
blocks and not boxes: two distinct blocks with equal boxes are distinct
entries, and their matrices $E_{*}$ may differ at a member.

\emph{(b), the stack.} Fix a pass class, and let $V_1$ be the elements $v$ of
its outcome pairs $(x,v)$ with $M^{(\mathsf t)}_{x}\ne0$ that are of the
latter case. Which elements these are, and their tuples, depend only on the
kinds fixed for the blocks and on the zero pattern, the same for every member.
The rows of $\widetilde E$ that belong to the pairs $(x,v)$ with
$v\notin V_1$ form a matrix $S_0$, the same for every member. Let $a$ be the
number of columns of the $E^{(\mathsf t)}_{x}$, the dimension of Alice's
elements on the sector $\mathsf t$, and $e_x$ the number of rows of
$E^{(\mathsf t)}_{x}$. For a tuple $\beta$ of an element of $V_1$, the pairs
$(x,v)$ with $v\in V_1$ and $\beta(v)=\beta$ have, by the reordering, the
associativity of the Kronecker product and the mixed-product rule
$(\idon{e_x}\otimes R)(A\otimes M)(\idon{a}\otimes\Sigma)=A\otimes RM\Sigma$,
the rows
\begin{equation*}
  E^{(\mathsf t)}_{x}\otimes E^{(s)}_v=(\idon{e_x}\otimes R_v)
  \big((E^{(\mathsf t)}_{x}\otimes E^{\mathrm{ex}}_v)\otimes
  E^{\mathrm{box}}_{\beta}\big)U_{\beta},\qquad
  U_{\beta}\coloneqq\idon{a}\otimes\Sigma_{\beta}(n).
\end{equation*}
The matrices $E^{(\mathsf t)}_{x}\otimes E^{\mathrm{ex}}_v$ of these pairs
all have the same number of columns, $a$ times the product of the $n_j$
with $j\notin\operatorname{pos}\beta$; let
$S_{\beta}$ stack them in the fixed order, the same for every member.
Stacking on the left commutes with the tensor product by a fixed right
factor, $\big[\begin{smallmatrix}A\\ A'\end{smallmatrix}\big]\otimes
E=\big[\begin{smallmatrix}A\otimes E\\ A'\otimes E\end{smallmatrix}\big]$,
and the right factor $E^{\mathrm{box}}_{\beta}$ is, at a member, one matrix
for all these pairs; so their rows, in the fixed order, form
$D_{\beta}(S_{\beta}\otimes E^{\mathrm{box}}_{\beta})U_{\beta}$, with
$D_{\beta}$ the block-diagonal matrix of the $\idon{e_x}\otimes R_v$, a
permutation matrix the same for every member. Every matrix $S$ has a rank
factorisation $S=BC'$ with $B$ of full column rank: over a field that contains
its entries, $B$ its pivot columns and $C'$ the non-zero rows of its reduced
row echelon form, an empty factorisation where $S$ is zero or has no row. Take
$S_0=B_0C_0$ and $S_{\beta}=B_{\beta}C_{\beta}$ so. By the mixed-product
rule, $S_{\beta}\otimes E^{\mathrm{box}}_{\beta}=(B_{\beta}\otimes
\idon{k_{\beta}})(C_{\beta}\otimes E^{\mathrm{box}}_{\beta})$, with
$k_{\beta}$ the number of rows of $E^{\mathrm{box}}_{\beta}$, the product of
the $k$ of the blocks of $\beta$, the same for every member; nothing else of
$E^{\mathrm{box}}_{\beta}$ enters. Let $\Pi$ be the permutation matrix that
returns the rows from this grouping, that of $S_0$ and one group per tuple,
to the fixed order. Then $\widetilde E=\Pi\operatorname{diag}\big(B_0,
(D_{\beta}(B_{\beta}\otimes\idon{k_{\beta}}))_{\beta}\big)\,C$ with $C$ the
stack of $C_0$ and the $(C_{\beta}\otimes E^{\mathrm{box}}_{\beta})U_{\beta}$;
the first factor, $\widetilde B$, is the same for every member and has full
column rank, being an invertible matrix times a block-diagonal matrix whose
blocks have full column rank, $D_{\beta}(B_{\beta}\otimes\idon{k_{\beta}})$
among them as an invertible matrix times a tensor product of injective maps.
Where the $E^{(\mathsf t)}_{x}$, the $E^{(s)}_v$ for $v\notin V_1$ and, for
$v\in V_1$, the $E_{v,j}$ with $j\notin\operatorname{pos}\beta(v)$ have
entries in $\mathbb{Q}(i)$, so do $S_0$ and every $S_{\beta}$, and the
factorisations over the field $\mathbb{Q}(i)$ give $B_0$, $C_0$, $B_{\beta}$
and $C_{\beta}$ there; $\Pi$, every $D_{\beta}$ and every $U_{\beta}$ have
entries $0$ and $1$.

\emph{(c).} Let $F^{\mathrm{ref}}$ be the companion's family of the replaced
factors. The family of each factor of $\Decl^{\mathrm{ref}}$ is a positive
operator-valued measure (Proposition~\ref{prop:reference-companion}\,(ii),
with condition (iii) of Definition~\ref{def:admissible-declaration}, which in
the general format asks it of every factor's family), so the blocks of
$F^{\mathrm{ref}}$ are Hermitian and positive semidefinite; they lie in their
entrywise enclosures, which are enclosures of them, and, by (H-FAC), in their
factor enclosures. So $F^{\mathrm{ref}}$ is a member,
$(F^{\mathrm{ref}},F^{\mathrm{ref}})$ is an admissible pair, and it lies in
$\mathcal P$. By (a), the key-inert pairs and joint key-active index sets of
$\Decl^{\mathrm{ref}}$ are those read off the enclosures, so the
$P^{J}_{\hat c}$, $\Gamma^{(s')}_{\hat c}$, deficits,
$\Delta^{(s')}_{\hat c}$ and $K^{J}_{\hat c}=P^{J}_{\hat c}+\sum
f^{J}_{\hat c,s}$ (Theorem~\ref{thm:choi}) formed with this pair are those of
$\Decl^{\mathrm{ref}}$ at $x$, every piece being formed from its one family:
an outcome pair whose element of Alice's is zero, which (c) leaves out of the
branch maps, contributes nothing to those of $\Decl^{\mathrm{ref}}$ either,
its Kraus factor $\sqrt{M^{(\mathsf t)}_{x'}}\otimes\sqrt{F^{(s)}_v}$ being
zero (Definition~\ref{def:branch}); and so does one whose element
$F^{(s)}_v$ has a block zero for every member, which (c) leaves out of the
branch maps and the statistics, that element being zero at
$F^{\mathrm{ref}}$, a member, so that its Kraus factor and its contribution
to the statistics are zero there. A set containing such a quantity formed
with every pair of
$\mathcal P$ contains that of $\Decl^{\mathrm{ref}}$. Where (H-POS) holds
for $\Decl^{\mathrm{ref}}$, these $q(\hat c)$, $P^{J}_{\hat c}$ and
$K^{J}_{\hat c}$ are positive as it asks, so
$(\lmoq,\lmoJ)[F^{\mathrm{ref}},F^{\mathrm{ref}}]$ is formed, and by the
formulas of Corollary~\ref{cor:full-gradient} it is the gradient
$(\lmoq,\lmoJ)$ of $\Decl^{\mathrm{ref}}$ at $x$ as
Lemma~\ref{lem:certified-bracket} reads it; the enclosures hold there.
$\Decl^{\mathrm{ref}}$ satisfies conditions (i)--(vi) and (viii) of
Definition~\ref{def:admissible-declaration}
(Proposition~\ref{prop:reference-companion}\,(ii)), so parts (ii), (iii)
and then (i) of Lemma~\ref{lem:certified-bracket} apply to it at $x$, where
(H-POS) holds. Finally, a
member $F$ is an entrywise member each of whose blocks lies in its entrywise
enclosure and in its factor enclosure where one is read, so $(F,F)$ belongs
to the set of pairs named in (c), whatever enclosure is chosen for each
block, and it is admissible: every block that admissibility reads is a factor
of an element no other block of which is zero for every member, at a pair
that is not key-inert, so the hypothesis asks it to be of fixed rank, and it
is zero or $E^{\dagger}\widetilde XE$ at $F$, and so positive semidefinite.
An element one of whose blocks is zero for every member is zero at $F$ and
contributes nothing to a piece formed with $(F,F)$, whatever its other
blocks.

\emph{(d).} Parts (a) and (b) read only the zero pattern and the ranks, the
hypothesis that the blocks are decided being unchanged and the fixed-rank
hypothesis now only at the kept pairs, and (c) runs with
$K^{J,\mathcal{D}}_{\hat c}$, formed with the deficits of the kept pairs
only, and with part (iv) of Lemma~\ref{lem:certified-bracket}, whose parts
(ii) and (iii) do not read the objective.
\end{proof}

% claim meat-core/frank-wolfe/cor:certified-keylength
\claimsummary{meat-core/frank-wolfe/cor:certified-keylength}{The largest bracket formed by \hyperlink{claim-meat-core/frank-wolfe/lem:certified-bracket}{Lemma~\ref*{lem:certified-bracket}} over finitely many points is a first-order bound of \hyperlink{claim-meat-core/frank-wolfe/thm:firstorder}{Theorem~\ref*{thm:firstorder}}, so it is at most the value of the declaration it is formed on and gives that declaration's fixed-length key length exactly as h\_lb does.}
\claimbookmark{meat-core/frank-wolfe/cor:certified-keylength}{Corollary}{cor:certified-keylength}{Key length from the certified bracket}
\claimleadin{%
By Lemma~\ref{lem:certified-bracket}\,(i) the sum of the two lower bounds is
the bracket of Theorem~\ref{thm:firstorder}\,(iv) at its point, so the best
of finitely many such sums is a first-order bound like any other, and the
key-length corollaries take it unchanged. What makes it certified is how the
two lower bounds are formed (Requirement~\ref{req:certified-evaluation}),
and what it still rests on is Remark~\ref{rem:certified-scope}. On a
reference companion, Corollary~\ref{cor:certified-keylength-reference}
carries the conclusion over to the protocol the companion declares.
}
\begin{mainclaim}
\begin{corollary}[Key length from the certified bracket]\label{cor:certified-keylength}
Let $\alpha$, $\Decl$ and $\acSet$ be as in
Lemma~\ref{lem:certified-bracket}, let $\Itset$ be a finite non-empty set of
points $x_t=(\mathbf q_t,J_t)$ each satisfying the hypotheses of that lemma,
and let $b^{\mathrm A}_t$ and $b^{\mathrm F}_t$ be real numbers with
\eqref{eq:lmo-bounds} at the gradient at $x_t$, as parts (ii) and (iii) of
that lemma form them. Put
\begin{equation}\label{eq:hcert}
  \underline{h}^{\mathrm{cert}}_{\alpha}\coloneqq
  \max_{t\in\Itset}\big(b^{\mathrm A}_t+b^{\mathrm F}_t\big).
\end{equation}
\begin{enumerate}[label=(\roman*),leftmargin=2.4em]
\item $\underline{h}^{\mathrm{cert}}_{\alpha}$ is the value of \eqref{eq:hlb}
  at the points of $\Itset$ with both pairings replaced as
  Theorem~\ref{thm:firstorder}\,(iv) replaces them, and
  $\max\{0,\underline{h}^{\mathrm{cert}}_{\alpha}\}\le\hlow$ in exact real
  arithmetic.
\item Where $\Decl$ satisfies Definition~\ref{def:admissible-declaration} and
  the hypotheses of Corollary~\ref{cor:fo-keylength} hold, its conclusions
  hold with $\underline{h}^{\mathrm{cert}}_{\alpha}$ in place of $\hlb$ in
  \eqref{eq:keylength-fo}: the key length so formed is $\epa$-secret and
  $\eev$-correct, with the composed-security and monotonicity conclusions of
  that corollary, the secrecy and security statements at every output
  state, whatever the dimension of the adversary's registers and the order
  in which the source releases its signals.
\item Under a dropped set, with the lemma's part (iv),
  $\max\{0,\underline{h}^{\mathrm{cert}}_{\alpha}\}\le
  \underline{h}^{\mathcal{D}}_{\alpha}\le\hlow$; and where $\Decl$
  satisfies Definition~\ref{def:admissible-declaration} and the hypotheses
  of Corollary~\ref{cor:keylength-bound} other than
  \eqref{eq:lower-bound-h} hold, the conclusions of that corollary hold at
  $h=\underline{h}^{\mathrm{cert}}_{\alpha}$.
\end{enumerate}
\end{corollary}
\end{mainclaim}
\claimreading{The certified key rate is the best of a few brackets, each formed from enclosures of the gradient at a point the search found, with every arithmetic step from the declaration's data exact or rigorously enclosed. It is a lower bound on the minimum h\_low of the declaration it is formed on, and the key length built from it is secret and correct under the same premises as the first-order key length: the declaration's modelling premises remain, but no premise about how a number was computed is added, the key length itself being computed exactly, its one logarithm enclosed until the integer part is decided. It is not a bound on any other declaration, so a run formed on the snapped declaration still rests where its record says on the premise that the device measures the snapped family.}
\claimmeta{\metalabel{Uses}\mref{Def.}{def:admissible-declaration}, \mref{Prop.}{prop:drop-bound}, \mref{Cor.}{cor:keylength-bound}, \mref{Def.}{def:relaxed-feasible-set}, \mref{Thm.}{thm:firstorder}, \mref{Cor.}{cor:fo-keylength}, \mref{Cor.}{cor:fo-dropped}, \mref{Lem.}{lem:certified-bracket}
\metasep \metalabel{Used by}\mref{Cor.}{cor:certified-keylength-reference}, \mref{Rem.}{rem:certified-scope}, \mref{Req.}{req:certified-evaluation}
\metasep newly derived
\metasep implemented (effective: accepted)
\metasep \metalabel{Code}\coderef{https://github.com/fraunhofersg/certiqs-key/blob/v0.1.0/src/certiqs_key/meat_core/finitekey/keylength.py}{certiqs_key.meat_core.finitekey.keylength:certified_key_length}}

\begin{proof}
(i) By Step~1 of the proof of Lemma~\ref{lem:certified-bracket}, each $x_t$
lies in $\acSet\times\Feas^{J}(\varepsilon_t)$ for some $\varepsilon_t\ge0$
and satisfies (H-POS), (H-SUPP) and (H-NOPERT); every point then lies in
$\acSet\times\Feas^{J}(\varepsilon)$ with $\varepsilon\coloneqq
\max_t\varepsilon_t$, the relaxed sets growing with $\varepsilon$
(Definition~\ref{def:relaxed-feasible-set}, each clause bounding a quantity
by a multiple of $\varepsilon$). By Lemma~\ref{lem:certified-bracket}\,(i)
the bracket of Theorem~\ref{thm:firstorder}\,(iv) at $x_t$ with both
pairings replaced is $b^{\mathrm A}_t+b^{\mathrm F}_t$, so
$\underline{h}^{\mathrm{cert}}_{\alpha}$ is \eqref{eq:hlb} so formed, and
Theorem~\ref{thm:firstorder}\,(iiib) and (iv), or (v) where condition (vii)
of Definition~\ref{def:admissible-declaration} fails, give the inequality.

(ii) Corollary~\ref{cor:fo-keylength} holds for any $\hlb$ of
Theorem~\ref{thm:firstorder}\,(iiib) with the lower bounds of its part (iv)
in place of the pairings, under its hypotheses, and (i) shows that
$\underline{h}^{\mathrm{cert}}_{\alpha}$ is one; it depends on the points,
the lower bounds and the declaration, and on none of
$(\lEC,\eev,\epa)$.

(iii) Lemma~\ref{lem:certified-bracket}\,(iv) gives
$\max\{0,b^{\mathrm A}_t+b^{\mathrm F}_t\}\le
\underline{h}^{\mathcal{D}}_{\alpha}$ at each point and
$\underline{h}^{\mathcal{D}}_{\alpha}\le\hlow$
(Proposition~\ref{prop:drop-bound}), so
$\underline{h}^{\mathrm{cert}}_{\alpha}$ is a lower bound on $\hlow$, the
hypothesis \eqref{eq:lower-bound-h} of
Corollary~\ref{cor:keylength-bound}, and it depends on none of
$(\lEC,\eev,\epa)$, the condition of that corollary's clause (a). Its
remaining hypotheses are those the second clause of (iii) assumes, so its
conclusions hold there at $h=\underline{h}^{\mathrm{cert}}_{\alpha}$.
\end{proof}

% claim meat-core/frank-wolfe/cor:certified-keylength-reference
\claimsummary{meat-core/frank-wolfe/cor:certified-keylength-reference}{On a reference companion, the best certified bracket, formed over its enclosed measured family as \hyperlink{claim-meat-core/frank-wolfe/lem:certified-bracket-enclosed}{Lemma~\ref*{lem:certified-bracket-enclosed}} allows, gives the key length of the protocol whose declaration is the companion, as the first-order bound on the companion does.}
\claimbookmark{meat-core/frank-wolfe/cor:certified-keylength-reference}{Corollary}{cor:certified-keylength-reference}{Key length on a reference companion from the certified bracket}
\claimleadin{%
Corollary~\ref{cor:certified-keylength} applies to a reference companion as
to any declaration; this corollary carries its conclusion over to the
protocol whose declaration the companion is, as
Corollary~\ref{cor:fo-reference} does for the first-order bound. Where the
companion's measured family is known only through enclosures, the two lower
bounds at a point are formed from enclosures that hold over every family the
enclosures leave possible, and part~(c) of
Lemma~\ref{lem:certified-bracket-enclosed}, or part~(d) under a dropped set,
is what makes them bounds at the companion's own gradient.
}
\begin{corollary}[Key length on a reference companion from the certified bracket]\label{cor:certified-keylength-reference}
Let $\Decl^{\mathrm{ref}}$ be a compiled reference companion of
Proposition~\ref{prop:reference-companion}, and let $\alpha$, $\acSet$ and
$\Itset$, with its points $x_t$, be as in
Corollary~\ref{cor:certified-keylength} with
$\Decl\coloneqq\Decl^{\mathrm{ref}}$. At each $t\in\Itset$ let
$b^{\mathrm A}_t$ and $b^{\mathrm F}_t$ be real numbers formed in one of two
ways: as in that corollary, with \eqref{eq:lmo-bounds} at the gradient of
$\Decl^{\mathrm{ref}}$ at $x_t$; or, where blocks of the measured family of
$\Decl^{\mathrm{ref}}$ are known through enclosures, entrywise or by factor
enclosures, as parts (ii) and (iii) of Lemma~\ref{lem:certified-bracket} form
them from enclosures $(\hat g,\delta)$ and $(M_p,\rho_p)$ that hold at every
pair of a set $\mathcal P$ as in
Lemma~\ref{lem:certified-bracket-enclosed}\,(c), at a point $x_t$ at which
the hypotheses of that lemma hold and (H-POS) holds for
$\Decl^{\mathrm{ref}}$, the blocks it reads with factor enclosures being read
with those of Corollary~\ref{cor:reference-companion-factors}, under that
corollary's hypotheses (for each replaced factor whose blocks are so read,
the rule of (I1) of the reduction instance of its last link a non-negative
combination, and the blocks of its handed reference carrying those factor
enclosures), so that its conclusion is (H-FAC). Put $\underline{h}^{\mathrm{cert}}_{\alpha}\coloneqq
\max_{t\in\Itset}(b^{\mathrm A}_t+b^{\mathrm F}_t)$, as in \eqref{eq:hcert}.
\begin{enumerate}[label=(\roman*),leftmargin=2.4em]
\item Where $\Decl^{\mathrm{ref}}$ satisfies
  Definition~\ref{def:admissible-declaration} and the hypotheses of
  Corollary~\ref{cor:fo-keylength} hold, the conclusions of
  Corollary~\ref{cor:fo-reference}\,(iii) hold with
  $\underline{h}^{\mathrm{cert}}_{\alpha}$ in place of
  $\hlb[\Decl^{\mathrm{ref}}]$, for the protocol whose declaration is
  $\Decl^{\mathrm{ref}}$.
\item In the setting of Corollary~\ref{cor:fo-dropped}, with a dropped set
  $\mathcal{D}$ and with the numbers formed at the gradient of the dropped
  objective and, in the second way, with part (d) of
  Lemma~\ref{lem:certified-bracket-enclosed} in place of its part (c):
  $\max\{0,\underline{h}^{\mathrm{cert}}_{\alpha}\}\le
  \underline{h}^{\mathcal{D}}_{\alpha}[\Decl^{\mathrm{ref}}]\le
  \hlow[\Decl^{\mathrm{ref}}]$, and where $\Decl^{\mathrm{ref}}$ satisfies
  Definition~\ref{def:admissible-declaration} and the hypotheses of
  Corollary~\ref{cor:keylength-bound} other than \eqref{eq:lower-bound-h}
  hold, the conclusions of
  Corollary~\ref{cor:keylength-bound} hold at
  $h=\underline{h}^{\mathrm{cert}}_{\alpha}$, for the protocol whose
  declaration is $\Decl^{\mathrm{ref}}$.
\end{enumerate}
\end{corollary}
\claimmeta{\metalabel{Uses}\mref{Def.}{def:admissible-declaration}, \mref{Cor.}{cor:keylength-bound}, \mref{Prop.}{prop:reference-companion}, \mref{Cor.}{cor:reference-companion-factors}, \mref{Thm.}{thm:firstorder}, \mref{Cor.}{cor:fo-keylength}, \mref{Cor.}{cor:fo-dropped}, \mref{Cor.}{cor:fo-reference}, \mref{Lem.}{lem:certified-bracket}, \mref{Lem.}{lem:certified-bracket-enclosed}, \mref{Cor.}{cor:certified-keylength}
\metasep \metalabel{Used by}\mref{Rem.}{rem:certified-scope}, \mref{Req.}{req:certified-evaluation}
\metasep newly derived
\metasep accepted}

\begin{proof}
\emph{The numbers are bounds at the companion's gradient.} At a point where
$b^{\mathrm A}_t$ and $b^{\mathrm F}_t$ are formed in the second way,
hypothesis (H-FAC) of Lemma~\ref{lem:certified-bracket-enclosed} holds by
Corollary~\ref{cor:reference-companion-factors}, under the hypotheses on
which the reduction and the compiler carry the factor enclosures, so part
(c) of that lemma shows that the enclosures
hold at the gradient of $\Decl^{\mathrm{ref}}$ at $x_t$, (H-POS) holding for
$\Decl^{\mathrm{ref}}$ there, and parts (ii) and (iii) of
Lemma~\ref{lem:certified-bracket}, which $\Decl^{\mathrm{ref}}$ is admitted
to by Proposition~\ref{prop:reference-companion}\,(ii), turn them into
$b^{\mathrm A}_t\le\min_{\mathbf q'\in\acSet}\langle\lmoq,\mathbf q'\rangle$
and $b^{\mathrm F}_t\le\min_{J'\in\Feas^{J}}\langle\lmoJ,J'\rangle$, the
gradient and the exact $\Feas^{J}$ being those of $\Decl^{\mathrm{ref}}$:
that is \eqref{eq:lmo-bounds} at its gradient, as in the first way.

\emph{(i).} Corollary~\ref{cor:fo-reference}\,(iii) holds for any bound
formed as its statement admits, lower bounds on the minima over the
companion's exact sets standing in for the pairings at points of its relaxed
set at which (H-POS), (H-SUPP) and (H-NOPERT) hold for its objects. Each
$x_t$ satisfies the hypotheses of Lemma~\ref{lem:certified-bracket} for
$\Decl^{\mathrm{ref}}$, which gives that by Step~1 of its proof, and the
numbers satisfy \eqref{eq:lmo-bounds} at the gradient of
$\Decl^{\mathrm{ref}}$, so Corollary~\ref{cor:certified-keylength}\,(i)
applied to $\Decl^{\mathrm{ref}}$ shows that
$\underline{h}^{\mathrm{cert}}_{\alpha}$ is such a bound: the value of
\eqref{eq:hlb} at the points of $\Itset$ with both pairings replaced.

\emph{(ii).} The same argument with part (d) of
Lemma~\ref{lem:certified-bracket-enclosed}, which runs part (c) for the
dropped objective, and part (iv) of Lemma~\ref{lem:certified-bracket} gives
the numbers as lower bounds on the two minima at the gradient of the dropped
objective of $\Decl^{\mathrm{ref}}$, and
Corollary~\ref{cor:certified-keylength}\,(iii) applied to
$\Decl^{\mathrm{ref}}$ gives the two inequalities and, under the hypotheses
that corollary's second clause and (ii) here assume, the conclusions of
Corollary~\ref{cor:keylength-bound}.
\end{proof}

\section{The two minimisations and the arithmetic}\label{group-meat-core/frank-wolfe/ass:lmo}
\sectionintro{The assumption that a solver's feasible-side return reached the minimum (\hyperlink{claim-meat-core/frank-wolfe/ass:lmo}{A-LMO}) is superseded, and no claim rests on it: the feasible side is bounded by \hyperlink{claim-meat-core/frank-wolfe/lem:feasible-side-dual}{Lemma~\ref*{lem:feasible-side-dual}}. The acceptance-side minimisation is a hypothesis of the bound that is discharged exactly on the exact island, for a singleton or a polytope acceptance set (\hyperlink{claim-meat-core/frank-wolfe/rem:acc-side}{Remark~\ref*{rem:acc-side}}). The remarks show that pairing over the \termref{term-acceptance-support}{acceptance support} loses nothing and say what exact real arithmetic means for a printed bound, and which of its terms a run still computes in binary64. A certified number rests on the claims, the declaration's modelling premises and the correctness of the implementation, and on no premise about arithmetic (\hyperlink{claim-meat-core/frank-wolfe/rem:certified-scope}{Remark~\ref*{rem:certified-scope}}).}
\sectionimports{\extref{../foundations/contract.pdf}{foundations/cor:least-eigenvalue-over-enclosure}, \extref{../foundations/contract.pdf}{foundations/cor:rank-over-factor-enclosure}, \extref{../foundations/contract.pdf}{foundations/def:tiers}}

% claim meat-core/frank-wolfe/ass:lmo
\claimsummary{meat-core/frank-wolfe/ass:lmo}{At every iterate a first-order bound uses, the block list consumed from the oracle, any positivity repair included, puts the linear form of the gradient its bracket consumes at or below its minimum over the exact Choi feasible set. No claim rests on it: \hyperlink{claim-meat-core/frank-wolfe/lem:feasible-side-dual}{Lemma~\ref*{lem:feasible-side-dual}} bounds that minimum instead.}
\claimbookmark{meat-core/frank-wolfe/ass:lmo}{Assumption}{ass:lmo}{Exact linear minimisation over the feasible set}
\claimleadin{%
The first-order bound consumes at each iterate two linear minimisations: one
over the acceptance set, which is exact rational data and is minimised
exactly, and one over the Choi feasible set, a semidefinite program a backend
solves in binary64. The second is what this assumption trusts. A positivity
repair the adapter applies to the backend's return lies inside the trusted
object, and checking that the return is feasible cannot replace the
assumption, since feasibility gives the opposite inequality.
}
\SetShortId{A-LMO}
\begin{assumption}[Exact linear minimisation over the feasible set]\label{ass:lmo}
(i) At each iterate $x_t$ at which the engine accumulates a bracket of the
first-order bound, the block list $S_t$ it consumes, which is the backend's
return after any positivity repair the adapter applies to it, satisfies the
\termref{term-feasible-side-inequality}{feasible-side inequality} of \eqref{eq:lmo}
(Definition~\ref{def:first-order-setting}), its minimum taken over the
\emph{exact} $\Feas^{J}$ and not over the \termref{term-relaxed-feasible-set}{relaxed feasible set}, at
$\lmoJ_t$, the block part of the gradient $(\lmoq_t,\lmoJ_t)$ the bracket
consumes at $x_t$, of the objective, or of that of a \termref{term-dropped-program}{dropped program}.
\end{assumption}
\claimverification{The engine does not verify this assumption; a run that rests on it records it as unverified.}
\claimreading{No run rests on this assumption: the engine verifies a lower bound on the solver's linear minimisation in exact arithmetic instead of trusting it. The assumption states that at every step whose bound the engine uses, the semidefinite-programming solver returned a block list at which the linear objective it was asked to minimise, the gradient of the objective the bound is taken for, is no larger than its true minimum over the exact feasible set, after the engine's small positivity repair of the return.}
\claimmeta{\metalabel{Uses}\mref{Def.}{def:first-order-setting}
\metasep \metalabel{Category}numerical trust
\metasep \metalabel{Affects}secrecy
\metasep \metalabel{Machine key}\iface{exact_linear_oracle}
\metasep newly derived
\metasep superseded}

% claim meat-core/frank-wolfe/rem:acc-side
\claimsummary{meat-core/frank-wolfe/rem:acc-side}{For a singleton or a polytope acceptance set the acceptance-side half of the linear minimisation is discharged by an exact minimisation, and a certified lower bound may replace the minimum, the bound staying valid. An acceptance iterate formed exactly stays in the set by induction.}
\claimbookmark{meat-core/frank-wolfe/rem:acc-side}{Remark}{rem:acc-side}{Exact minimisation on the acceptance side}
\begin{remark}[Exact minimisation on the acceptance side]\label{rem:acc-side}
The acceptance half of \eqref{eq:lmo}, the
\termref{term-acceptance-side-inequality}{acceptance-side inequality}, is discharged exactly, because its
minimisation over $\acSet$ lies on the \exttermref{../foundations/contract.pdf}{term-exact-island}{exact island}
(\xref{foundations/def:exact-island}\,(ii)): $\acSet$ is declared as exact rational
data and the minimisation is a linear program over it. Two geometries are
admitted, the singleton~(a) and the polytope~(b), and each supplies what
the bracket needs.

The \termref{term-acceptance-support}{acceptance support} $\Chacc$ is \eqref{eq:chacc-def}, and by
Lemma~\ref{lem:fo-khyp}\,(i) it is the support of a single element of
$\acSet$ while no element has support outside it. So minimising a linear
functional over $\acSet$ may be done in the coordinates $\Chacc$ alone,
which is also where the acceptance-side gradient $\lmoq_t$ of
Theorem~\ref{thm:firstorder} lives and where the first pairing of
Definition~\ref{def:first-order-setting} is taken --- the three statements
agree, and the case $\Chacc\subsetneq\Chatset$ is admitted throughout rather
than excluded by a hypothesis.

Given an exact vector $\lmoq\in\mathbb{R}^{\Chacc}$, a number $\bacc$ with
\begin{equation}\label{eq:bacc-def}
  \bacc\;\le\;\min_{\mathbf q'\in\acSet}\
  \sum_{\hat c\in\Chacc}\lmoq(\hat c)\,q'(\hat c)
\end{equation}
is a certified lower bound on the acceptance-side minimum, and a point
attaining the minimum discharges the first inequality of \eqref{eq:lmo}
outright. \emph{The bracket may consume $\bacc$ in place of the pairing
with such a point, and the substitution is exact rather than
approximate}: at $\lmoq=\lmoq_t$, \eqref{eq:bacc-def} is the hypothesis on
$b^{\mathrm A}_t$ of Theorem~\ref{thm:firstorder}\,(iv), so replacing
$\langle\lmoq_t,v_t\rangle$ by $\bacc$ in \eqref{eq:hlb} --- and only
there, the subtrahend $\langle\lmoq_t,\mathbf q_t\rangle$ being evaluated
as before --- keeps the conclusion $\hlb\le\hlow$ (a run that does so
records which of the two it used, Requirement~\ref{req:bacc-recorded}). Two instantiations are part of this submodule; an $\acSet$
that fits neither is refused at admission
(Requirement~\ref{req:acceptance-geometries}).
\emph{(a) Singleton} $\acSet=\{\mathbf q^{*}\}$: the minimum is the single
evaluation $\sum_{\hat c}\lmoq(\hat c)q^{*}(\hat c)$, attained at
$v_t=\mathbf q^{*}$, and $\bacc$ is any downward-rounded evaluation of it.
\emph{(b) Polytope}
$\acSet=\{\mathbf q:\ \mathbf q\ge0,\ \textstyle\sum_{\hat c}q(\hat c)=1,\
B\mathbf q\le\mathbf b\}$ for an exact matrix $B$ and vector
$\mathbf b$, read in the coordinates $\Chacc$: the minimum is attained at
a vertex and is computed exactly over the rationals, which discharges
\eqref{eq:lmo}'s first inequality with equality. A dual witness certifies
it independently: for any $\mathbf u\ge0$ and $v\in\mathbb{R}$ with
$\lmoq(\hat c)+(B^{\transpose}\mathbf u)(\hat c)+v\ge0$ for every $\hat
c\in\Chacc$,
\begin{equation*}
  \min_{\mathbf q'\in\acSet}\langle\lmoq,\mathbf q'\rangle
  \;\ge\;-\langle\mathbf b,\mathbf u\rangle-v :
\end{equation*}
indeed $\langle\lmoq,\mathbf q'\rangle
=\langle\lmoq+B^{\transpose}\mathbf u+v\mathbf 1,\mathbf q'\rangle
-\langle\mathbf u,B\mathbf q'\rangle-v
\ge0-\langle\mathbf u,\mathbf b\rangle-v$, using $\mathbf q'\ge0$
componentwise, $\mathbf u\ge0$ with $B\mathbf q'\le\mathbf b$, and
$\sum_{\hat c}q'(\hat c)=1$.

\emph{The acceptance-side iterate is on the island too, and that is what
makes $\mathbf q_t\in\acSet$ hold by construction.} The membership
hypothesis of Theorem~\ref{thm:firstorder}\,(ii) and~(iiib) is on a
\emph{product}, $x_t\in\acSet\times\Feas^{J}(\varepsilon)$, and the three
checks of Requirement~\ref{req:iterate-checks} are all conditions on the $J$
factor: per-block positive semidefiniteness, the coupling residual and the
declared rows. \emph{Nothing checks the other factor, and nothing needs
to}, provided the acceptance-side iterate is formed on the exact island
under two conditions (Requirement~\ref{req:acceptance-iterate-exact}):
\begin{itemize}[leftmargin=2.4em]
\item[($\alpha$)] the acceptance-side iterate is initialised at an exact
rational point $\mathbf q_{0}\in\acSet$ --- one exists whenever
$\acSet\neq\varnothing$, the acceptance data being exact rationals by
Definition~\ref{def:admissible-acceptance-set} and \xref{foundations/def:exact-island}, so that
the single point of geometry~(a) and the vertices of geometry~(b) are
rational; and Lemma~\ref{lem:fo-khyp}\,(i)'s equally weighted average of
finitely many of them is rational too, and is supported on all of
$\Chacc$, which Remark~\ref{rem:fo-reachability} also needs;
\item[($\beta$)] every acceptance-side update is a convex combination,
with \emph{exact rational} weights in $[0,1]$, of points already
established to lie in $\acSet$ --- in the intended use
$\mathbf q_{t+1}=(1-\gamma_t)\mathbf q_t+\gamma_t\,u_t$ with
$u_t\in\acSet$ exact and rational, and in geometries (a) and (b) the
acceptance-side minimiser $v_t$ is such a point: the single
$\mathbf q^{*}$ in (a), a vertex computed exactly over the rationals in
(b).
\end{itemize}
Under ($\alpha$) and ($\beta$), $\mathbf q_t\in\acSet$ for every $t$ by
induction, $\acSet$ being convex (Definition~\ref{def:admissible-acceptance-set}) and the
arithmetic exact --- so the property is \emph{true} rather than tested,
there is no check for a run to fail and no residual for it to expose, and
\eqref{eq:hlb} is evaluated at a point whose acceptance factor is where
the theorem says it is. \emph{Why it must be exact and not merely
checked}: on the polytope geometry (b) a binary64 convex combination of
exact points of $\acSet$ can violate $\sum_{\hat c}q(\hat c)=1$ or a row
of $B\mathbf q\le\mathbf b$, and three consumers break if it does ---
Lemma~\ref{lem:fo-khyp}\,(ii)'s identical deletion of the
$\hat c\notin\Chacc$ terms is unlicensed and $O^{J}_{\alpha}(x_t)$ may be
$+\infty$; the pairing on $\mathbb{R}^{\Chacc}$ loses its justification,
since $v_t-\mathbf q_t$ need no longer vanish outside $\Chacc$ and
$\langle\lmoq_t,v_t-\mathbf q_t\rangle$ would silently drop terms of
unknown sign from \eqref{eq:hlb}; and Step~2's chord argument needs $x_t$
in the convex set it runs in. The singleton geometry (a) is immune ---
there $\acSet$ has one point and every convex combination of it is it ---
but the polytope geometry is a first-class instantiation of this submodule
and is not.

\emph{What this leaves.} With the acceptance half discharged, the other
half of \eqref{eq:lmo}, the \termref{term-feasible-side-inequality}{feasible-side inequality}, concerns the
minimisation over $\Feas^{J}$ --- a semidefinite program in the block
variables, solved in binary64 by a backend. It is not asked of the backend's
return at all: part~(iv) of Theorem~\ref{thm:firstorder} takes a lower
bound on that minimum in place of the pairing, and
Lemma~\ref{lem:feasible-side-dual} supplies one from the backend's
multipliers, formed in exact arithmetic over the exact $\Feas^{J}$, the
rounding of the frame the rows are built with charged
(Requirement~\ref{req:verified-feasible-side}). Both sides of a bracket are
then lower bounds over the exact sets, certified exactly for the linear
forms the run computed; what the binary64 computation of those forms
leaves is Remark~\ref{rem:fo-arithmetic}'s.
\end{remark}
\claimmeta{\metalabel{Uses}\mref{Def.}{def:admissible-acceptance-set}, \mref{Lem.}{lem:fo-khyp}, \mref{Thm.}{thm:firstorder}
\metasep \metalabel{Used by}\mref{Rem.}{rem:pairing-on-acceptance-support}, \mref{Rem.}{rem:fo-interface}, \mref{Rem.}{rem:fo-iterate-checks}, \mref{Rem.}{rem:fo-returned-point}, \mref{Rem.}{rem:fo-check-exposures}, \mref{Req.}{req:acceptance-geometries}, \mref{Req.}{req:acceptance-iterate-exact}, \mref{Req.}{req:bacc-recorded}, \mref{Req.}{req:certified-evaluation}
\metasep newly derived
\metasep accepted}

% claim meat-core/frank-wolfe/rem:pairing-on-acceptance-support
\claimsummary{meat-core/frank-wolfe/rem:pairing-on-acceptance-support}{Restricting the acceptance-side pairing to the acceptance support loses nothing: a feasible direction vanishes outside it, and the oracle term is defined by that pairing. The oracle return therefore need not be feasible, though in both admitted acceptance geometries it is.}
\claimbookmark{meat-core/frank-wolfe/rem:pairing-on-acceptance-support}{Remark}{rem:pairing-on-acceptance-support}{Pairing over the acceptance support}
\begin{remark}[Pairing over the acceptance support]\label{rem:pairing-on-acceptance-support}
Restricting the first pairing of Definition~\ref{def:first-order-setting}
to the \termref{term-acceptance-support}{acceptance support} $\Chacc$ costs nothing and is not a
choice, and the reason splits in two according to what the pairing is
applied to. \emph{Where a directional derivative is taken}, the vector on
the right is $\mathbf q'-\mathbf q_t$ with \emph{both} entries in
$\acSet$: $\mathbf q'$ by the quantifier of \eqref{eq:fo-subgrad}, and
$\mathbf q_t$ because the acceptance-side iterate is formed on the
\exttermref{../foundations/contract.pdf}{term-exact-island}{exact island} and lies in $\acSet$ by construction
(Remark~\ref{rem:acc-side}, ($\alpha$) and ($\beta$)). So it vanishes
outside $\Chacc$ by Lemma~\ref{lem:fo-khyp}\,(i), and the sum over
$\Chacc$ and the sum over $\Chatset$ agree, the second being the one that
may fail to be defined. \emph{Where the oracle return $v_t$ enters}, the
$\Chacc$ sum is not an abbreviation of the $\Chatset$ sum but the
definition: \eqref{eq:lmo} is stated in it, \eqref{eq:hlb} evaluates it,
and the two occurrences are the same functional, so the algebra closes
whether or not $v_t$ lies in $\acSet$; that is why \eqref{eq:lmo} can
decline to require feasibility of the pair. In both geometries Remark~\ref{rem:acc-side} admits, $v_t$
\emph{does} lie in $\acSet$ (the singleton's one point; a vertex computed
exactly over the rationals), or the acceptance half is discharged through
$\bacc$ and no $v_t$ enters at all.
\end{remark}
\claimmeta{\metalabel{Uses}\mref{Lem.}{lem:fo-khyp}, \mref{Rem.}{rem:acc-side}
\metasep newly derived
\metasep accepted}

% claim meat-core/frank-wolfe/rem:fo-arithmetic
\claimsummary{meat-core/frank-wolfe/rem:fo-arithmetic}{Both linear minima are bounded exactly over the exact sets, but for the gradient as computed: the objective, the gradient and the maps and frame behind them are binary64, and a reference's centres move the measured family. So a computed h\_lb may exceed h\_low, and the zero clamp repairs neither.}
\claimbookmark{meat-core/frank-wolfe/rem:fo-arithmetic}{Remark}{rem:fo-arithmetic}{Exact real arithmetic in the first-order bound}
\begin{remark}[Exact real arithmetic in the first-order bound]\label{rem:fo-arithmetic}
Theorem~\ref{thm:firstorder}\,(ii), respectively~(iiib) --- the part a
run is in, its iterates being established only in the relaxed set
(Requirement~\ref{req:iterate-checks}) --- with the lower bounds of its
part~(iv) in place of both pairings, which is how a run forms every bracket
(Remark~\ref{rem:acc-side}, Requirement~\ref{req:verified-feasible-side}),
concludes $\hlb\le\hlow$ in exact real arithmetic, for the exact objective
and the exact gradient. No assumption about a solver enters it. What the
theorem still asks, and a run cannot verify, is exact real arithmetic in
the terms a run computes in binary64. What that means for a number a run
prints is narrower than the theorem, and a reader who takes \eqref{eq:hlb}
for a runtime invariant is wrong in the unsafe direction.

\emph{What is exact, and of what.} The two lower bounds that stand in for
the minima are certified exactly
(\xref{foundations/def:exact-island}\,(ii), (iii)), and both are bounds over
the exact sets, but of the linear form at the gradient as the run computed
it. The acceptance-side one is the exact minimum over $\acSet$ of that
linear form. The feasible-side one is at most the minimum over the exact
$\Feas^{J}$ of that linear form: the run builds its constraint rows with a
binary64 frame $\widetilde W_{\mathsf t}$ in place of the frame
$W_{\mathsf t}$ of Theorem~\ref{thm:choi}, which carries a square root of
the declared marginal, and Lemma~\ref{lem:feasible-side-dual}\,(iv) charges
the difference, a deduction linear in the row multipliers and in a bound on
$\opn{\widetilde W_{\mathsf t}-W_{\mathsf t}}$; where no such bound is
formed the row multipliers are zero and the rows are not read. So no status
of the backend, no dual it returns and no rounding of the frame can raise
the feasible-side bound above that minimum. Without the deduction it could
be raised: Lemma~\ref{lem:feasible-side-dual}\,(i) applied to rows built
from $\widetilde W_{\mathsf t}$ bounds the minimum over the set those rows
cut out, which need not contain $\Feas^{J}$ and may be empty, and over an
empty set every number is a lower bound.

\emph{The conclusion is not a floating-point invariant.} The objective
$O^{J}_{\alpha}(x_t)$, the gradient $(\lmoq_t,\lmoJ_t)$ and the subtracted
pairings $\langle\lmoq_t,\mathbf q_t\rangle$ and $\langle\lmoJ_t,J_t\rangle$
are computed in binary64 (\xref{foundations/req:exact-island}), and so are
the data the engine derives from the declaration for them: the
\termref{term-statistics-map}{statistics map}, the \termref{term-branch-map}{branch maps} and their supports, and the frame
$\widetilde W_{\mathsf t}$
through which they read Alice's side. The gap between the gradient as
computed and the exact gradient is part of that unbudgeted evaluation
arithmetic, and no claim of this module bounds it. A computed $\hlb$ may
therefore exceed the computed estimate, and may exceed the true $\hlow$, by
an amount no claim of this module bounds. Two consequences follow, and
neither is a defect of the theorem.
\begin{enumerate}[label=(\alph*),leftmargin=2.2em]
\item \emph{$\max\{0,\hlb\}$ does not repair it.} The clamp is sound and
  it is what Corollary~\ref{cor:fo-keylength} consumes, but it corrects
  only the \emph{low} side: clamping at zero cannot bring down a bracket
  that came out too high. The clamp is no general safety net, and a run is
  not to be written, or read, as if it were.
\item \emph{A comparison of $\hlb$ against the reported estimate is a
  diagnostic, never an assertion.} The ordering of $\hlb$ below the
  objective at the returned point is a consequence of exact arithmetic and
  of the returned point's exact membership in $\acSet\times\Feas^{J}$
  (Theorem~\ref{thm:firstorder}\,(i) and~(ii)), neither of which this module
  establishes for a run: its returned point is established at most in the
  relaxed set, where the objective is bounded below only by
  $\hlow[\Feas^{J}(\varepsilon)]\le\hlow$
  (Theorem~\ref{thm:firstorder}\,(iiia)), and near convergence, where
  $\hlb$ is close to $\hlow$, that alone can put the estimate below $\hlb$.
  The difference $\hlb$ minus the estimate is recorded with its sign
  (Requirement~\ref{req:bound-estimate-difference}), and a positive
  difference, a bound above the estimate, is evidence about the arithmetic
  or about the relaxation, and not by itself evidence that the first-order
  tier is unsound.
\end{enumerate}

\emph{Evaluation on a reference companion.} Where the bound is formed on a
compiled reference companion $\Decl^{\mathrm{ref}}$
(Corollary~\ref{cor:fo-reference},
Requirement~\ref{req:bracket-on-reference}), the theorem is applied to the
companion, a declaration of the reference model, but what is evaluated is
not that declaration. Its measured family is known through an entrywise
enclosure, recentred so that every entry whose enclosure contains zero has
centre zero (\xref{foundations/lem:recentred-enclosure}), and it is
evaluated at the binary64 numbers nearest to the centres. The evaluated
family differs from the companion's, entry by entry, by at most the radius
and one rounding, and it need not be a POVM: its completeness defect and its
least eigenvalue are measured and recorded, never decided
(Requirement~\ref{req:companion-admission}\,(iii)). The objective, the
gradient, the \termref{term-branch-map}{branch maps} and their supports, the key-inert pairs and the
rank comparison behind (H-SUPP) are all formed from those numbers, and so is
the linear form whose minimum the feasible-side bound bounds. Recentring is
applied to the handed
blocks and again to the companion's datum after the construction has
combined their entries (\xref{reductions/req:reference-companion}\,(i)), so
that an entry of the datum whose exact value is a sum that cancels to zero
is not left with a centre that is not zero. After it, a centre is not zero
only where its enclosure excludes zero; the compiler's weighting multiplies
a centre by a weight and its radius by the weight's modulus, and a product of
factor families multiplies single entries, so the compiled family keeps that
property, and rounding a centre to binary64 can turn it to zero but never
the reverse. The evaluated family therefore has a non-zero entry only where
the companion's family has one; it may lack one where the companion's entry
is no larger than the recentred radius there, and a support,
or a key-active pair, that the companion has through such entries alone is
not seen by the evaluation. What the theorem asks of the point itself,
(H-SUPP), holds of the companion whenever every block of $J_t$ is positive
definite (Corollary~\ref{cor:fo-reference}\,(ii)). This exposure is
larger in kind than the binary64 evaluation of the objective and the
gradient on any declaration. There the inputs are exact and only the
formulas are evaluated in binary64. Here the measured family itself is
displaced, by up to the radius and one rounding per entry, with a relative
error up to order one on an entry not much larger than its radius. The
displacement from the companion to its evaluation is a perturbation of the
measured family, and Proposition~\ref{prop:optol-nobound} applies to it:
closeness of two measured families bounds the change of the value by
nothing, and the change may point either way
(Remark~\ref{rem:optol-direction}); the evaluated family need not even be a
POVM. Its size is set by the radius of the enclosure, and not by binary64
rounding alone, save on an exact reference, whose radius is zero. No claim
of this module bounds its effect, and the largest radius, which the record
carries, is the only handle on it. Evaluating on the companion removes the
premise that the device measures the snapped family
(\xref{foundations/ass:snapped-measurement}); it does not remove this
exposure, which narrows only with the enclosure. The companion's entries are
not part of this exposure: they are emitted in binary64 as the oracle reads
them, their rounding absorbed into their constants
(\xref{reductions/lem:row-relaxation}), so they are the companion's own
exact rows, and the feasible-side bound is over the companion's exact
$\Feas^{J}$, the frame's rounding charged as on any declaration. The
first-order number is not evaluated over the enclosure; the certified number
of Requirement~\ref{req:certified-evaluation} is, where every element the
gradient reads has a rank fixed throughout the enclosure
(Lemma~\ref{lem:certified-bracket-enclosed}), and it then holds of the
companion itself.

\emph{Where the exposure is.} Of the ingredients of \eqref{eq:hlb}, the
objective and the gradient are evaluated by formulas this module states
in closed form, at points where (H-POS) keeps every denominator bounded
away from zero; the subtracted pairings are inner products; and the two
lower bounds are certifications of the \exttermref{../foundations/contract.pdf}{term-exact-island}{exact island} over the exact
sets, for the gradient as computed. No ingredient is a solver's return: the
oracle's point and its status are diagnostics
(Requirement~\ref{req:verified-feasible-side}), and a solve that stopped
short of its optimum, or failed, loosens a bracket and cannot raise it
above the value it has at the exact minima of the computed linear forms.
What a run leaves unverified is named: the binary64 evaluation of the
objective, of the gradient and of the maps and frame they are built from,
of the subtracted pairings, and of the checks whose one-sided exposures
Remark~\ref{rem:fo-check-exposures} names; and on a reference companion
the displacement to its centres. It is \emph{unbudgeted}: no claim of this
module bounds it, and no assumption of this module covers it, the engine
computing outside the exact island without an error budget, as
\xref{foundations/req:exact-island} allows. That is what the
\exttermref{../foundations/contract.pdf}{term-first-order}{first-order} tier costs: it is weaker than a verified numerical
certificate of $\hlow$, a first-order number being a lower bound
established under the hypotheses that the claim stating it names
(\xref{foundations/def:tiers}), evaluated in arithmetic that no claim of this
module budgets. The \exttermref{../foundations/contract.pdf}{term-certified}{certified} tier of the same definition is the bound
formed again with that arithmetic exact or enclosed, and what it still rests
on is Remark~\ref{rem:certified-scope}.
\end{remark}
\claimmeta{\metalabel{Uses}\mref{Prop.}{prop:optol-nobound}, \mref{Rem.}{rem:optol-direction}, \mref{Thm.}{thm:firstorder}, \mref{Cor.}{cor:fo-keylength}, \mref{Cor.}{cor:fo-reference}, \mref{Lem.}{lem:feasible-side-dual}
\metasep \metalabel{Used by}\mref{Rem.}{rem:certified-scope}, \mref{Rem.}{rem:fo-iterate-repair}, \mref{Rem.}{rem:fo-negative-values}, \mref{Req.}{req:bound-estimate-difference}
\metasep newly derived
\metasep accepted}

% claim meat-core/frank-wolfe/rem:certified-scope
\claimsummary{meat-core/frank-wolfe/rem:certified-scope}{A certified number rests on the claims, the modelling premises of its declaration and the correctness of the code, its ball-arithmetic library included, and on no premise about arithmetic. It does not reach data computed before the declaration; a receiver weight in floating point withholds it.}
\claimbookmark{meat-core/frank-wolfe/rem:certified-scope}{Remark}{rem:certified-scope}{What a certified number rests on}
\begin{remark}[What a certified number rests on]\label{rem:certified-scope}
A certified number, Requirement~\ref{req:certified-evaluation}, is by
Corollary~\ref{cor:certified-keylength} a lower bound on the value of the
declaration it is formed on, and it rests on three things.

\emph{The claims.} Lemma~\ref{lem:certified-bracket} and
Corollary~\ref{cor:certified-keylength}, on a reference companion with
Lemma~\ref{lem:certified-bracket-enclosed} and
Corollary~\ref{cor:certified-keylength-reference}, with what they use, the hypotheses of
Theorem~\ref{thm:firstorder} among them, each of which the evaluation
establishes at the point itself: positive definite blocks, which give
(H-SUPP) by Lemma~\ref{lem:branch-support}\,(i); (H-POS) with certified
signs; and the exact gradient, which is (H-NOPERT). The binary64 checks whose
exposures Remark~\ref{rem:fo-check-exposures} names, the oracle's return and
the gradient the search computed enter no certified number, and the
exposures of Remark~\ref{rem:fo-arithmetic} are those of the first-order
number, not of this one.

\emph{The modelling premises of that declaration.} They are the premises of
the first-order number on the same declaration, and certification decides
none of them and adds none. Where the number is formed on the declaration
the run searched, the blocks that rest there on
\xref{foundations/ass:snapped-measurement} rest on it for the certified number
as well. That includes the fallback of
Requirement~\ref{req:certified-evaluation}\,(i): where a companion counts but
certifies no point and the number is certified on the searched declaration,
the blocks the companion carried rest on that assumption for the reported
key length, though they would not for the companion's first-order number
(Requirement~\ref{req:bracket-on-reference}\,(v)). On a reference companion
the number holds of the companion itself, not of the centres of its
enclosure; and the premises of the source and the receiver remain, on every
route.

\emph{The implementation.} Every arithmetic step from the declaration's data
is exact or \exttermref{../foundations/contract.pdf}{term-rigorously-enclosed}{rigorously enclosed}, so no premise is about how a number
was computed and none of the category numerical-trust enters through the
evaluation (\xref{foundations/def:tiers}). What remains is that the code
computes what the claims state, the library of the rigorous arithmetic
included. That is the correctness of the engine, held to review and to tests
like every other part of it; no assumption is stated for it, and a defect
there is a defect of the engine and not the failure of a premise.

The chain begins at the declaration's data. A supplier's floating-point
evaluation that produced those data, under an allowance or an assumption of
its own, lies before it. Such an evaluation of the receiver's blocks reaches
a number formed on the declaration the run searched only through the snap,
where the run rests on \xref{foundations/ass:snapped-measurement} as the
paragraph on modelling premises above says, and the tier is not withheld
there. The tier is withheld in the two cases
Requirement~\ref{req:certified-evaluation}\,(vii) names: where an
assumption of the category numerical-trust is invoked on the run, and where
the declaration carries a weight routed from a receiver's high-photon weights
into which a per-setting weight enters that the receiver evaluated in
floating point, the uniform weight counting every setting, a number that
holds of the device only as \xref{reductions/ass:recv}\,(ii) attests it and
that the receiver records so (\xref{receiver/req:evaluation-label}). A receiver weight
evaluated rigorously is a bound by the receiver's own theorem
(\xref{receiver/thm:enclosed-weight}), and rests on that assumption only for
the artifact's provenance and for the device premises of the receiver, where
the ledger carries its row; for an artifact the run built itself, exact and
with every weight evaluated rigorously, the row is omitted
(\xref{reductions/req:recv-row-provenance}), and the receiver's own
ledger carries those device premises.

Over a companion known through enclosures, the number exists only where
every element the gradient reads has a fixed rank
(Lemma~\ref{lem:certified-bracket-enclosed}). An element whose principal
block on the rows its entrywise enclosure does not fix at zero is singular
has no positive lower bound on that block's least eigenvalue, the element
being among the matrices its enclosure contains, so the rule of
\xref{foundations/cor:least-eigenvalue-over-enclosure}\,(iv) fixes no rank
for it, at any radius. Among the elements whose enclosure carries a
positive radius, that is the case of every rank-deficient element whose
support is spanned by no set of coordinate vectors, by part (iii) of that
corollary, such as the projector of an ideal analyser at an angle whose
projector has irrational entries, and of every rank-deficient element whose
zero rows the enclosure does not fix at zero. Such an element has a fixed
rank all the same where it is carried with a factor enclosure that has no
column, or whose columns have a Gram matrix with a positive floor over its
box, its rank then being the number of columns at every member
(\xref{foundations/cor:rank-over-factor-enclosure}); so the analyser's
projector certifies on the companion where the receiver emits its elements
in that form. A factor enclosure fixes no rank where the columns of a member
are linearly dependent, its Gram matrix being singular there, and where the
member's rank is also below its dimension, neither rule fixes one: a
multi-photon element of a receiver with loss and no dark counts can be of
this kind. Nor is a rank fixed for a rank-deficient element carried with no
factor enclosure, as where the receiver forms its blocks without an exact
dilation of its instrument, for a generic component with none, or where a squashing map acts on a data sector after
the factor is formed. A run whose companion reads such an element at a pair
that is not \termref{term-key-inert}{key-inert} certifies no point on the companion, and it
reaches the certified tier only through the fallback of
Requirement~\ref{req:certified-evaluation}\,(i), on the declaration it
searched, which does not run under the strict option \emph{strict\_snaps}
(\xref{foundations/req:snap-reference}\,(iv)), resting there on
\xref{foundations/ass:snapped-measurement} for the blocks the companion
carried. An element the companion carries exactly,
with radius zero, is factorised exactly and is not subject to these rules.
\end{remark}
\claimreading{Certified means that every number between the declaration and the key rate was computed exactly or with a rigorous bound on its rounding. It does not mean that the declaration describes your devices: the source, receiver and laboratory premises in the report still apply, and a number formed on the snapped measurement rests where the record says on the device measuring it. Nor is the software assumed correct by a premise: its correctness, its interval-arithmetic library included, is checked by review and tests. Ideal analysers with irrational projector entries (at pi/8, say) certify on your model where the receiver gives their elements in factor form; otherwise, as for some lossy multi-photon receivers, the rate can be certified on the snapped measurement, resting there on that premise, unless the run sets strict\_snaps. A floating-point multiclick weight keeps the first-order number.}
\claimmeta{\metalabel{Uses}\mref{Lem.}{lem:branch-support}, \mref{Thm.}{thm:firstorder}, \mref{Lem.}{lem:certified-bracket}, \mref{Lem.}{lem:certified-bracket-enclosed}, \mref{Cor.}{cor:certified-keylength}, \mref{Cor.}{cor:certified-keylength-reference}, \mref{Rem.}{rem:fo-arithmetic}, \mref{Rem.}{rem:fo-check-exposures}, \mref{Req.}{req:certified-evaluation}, \extref{../foundations/contract.pdf}{foundations/cor:least-eigenvalue-over-enclosure}, \extref{../foundations/contract.pdf}{foundations/cor:rank-over-factor-enclosure}, \extref{../foundations/contract.pdf}{foundations/def:tiers}
\metasep newly derived
\metasep accepted}

\section{The reported numbers and the runtime checks}\label{group-meat-core/frank-wolfe/rem:fo-interface}
\sectionintro{These remarks give the mathematics behind the reporting rules. A run reports its estimate and, where the hypotheses held, the first-order bound, and withdrawal never weakens what the bound means (\hyperlink{claim-meat-core/frank-wolfe/rem:fo-interface}{Remark~\ref*{rem:fo-interface}}). The iterate checks decide the \termref{term-relaxed-feasible-set}{relaxed feasible set} on the Choi factor only and serve the membership hypothesis and the estimate's premise (\hyperlink{claim-meat-core/frank-wolfe/rem:fo-iterate-checks}{Remark~\ref*{rem:fo-iterate-checks}}); maximal support is a rank equality read with one ruler (\hyperlink{claim-meat-core/frank-wolfe/rem:fo-one-ruler}{Remark~\ref*{rem:fo-one-ruler}}); and the checks, computed in binary64, leave one-sided exposures that no claim bounds. A repaired oracle return is a step of the search and enters no bracket, a repaired iterate is a construction step, and a negative bound is clamped while a negative estimate is refused.}

% claim meat-core/frank-wolfe/rem:fo-interface
\claimsummary{meat-core/frank-wolfe/rem:fo-interface}{A run reports its estimate always, and h\_lb only where the hypotheses of \hyperlink{claim-meat-core/frank-wolfe/thm:firstorder}{Theorem~\ref*{thm:firstorder}} held, membership of the consumed point in the acceptance set times the relaxed set among them. No hypothesis is read of the solver, and withdrawal never weakens what h\_lb means.}
\claimbookmark{meat-core/frank-wolfe/rem:fo-interface}{Remark}{rem:fo-interface}{Two reported numbers and one first-order claim}
\claimleadin{%
Every run reports its \exttermref{../foundations/contract.pdf}{term-estimate}{estimate}, and where the hypotheses of
Theorem~\ref{thm:firstorder} held the \exttermref{../foundations/contract.pdf}{term-first-order}{first-order} bound $\hlb$ as
well, naming the \exttermref{../foundations/contract.pdf}{term-tier}{tier} that carries its claim
(\xref{foundations/req:tiers}). The requirements of this submodule say when
$\hlb$ is withdrawn and what a run checks, repairs and records; this remark and
Remarks~\ref{rem:fo-iterate-checks}, \ref{rem:fo-returned-point},
\ref{rem:fo-one-ruler}, \ref{rem:fo-reachability},
\ref{rem:fo-check-exposures}, \ref{rem:fo-oracle-repair},
\ref{rem:fo-iterate-repair} and~\ref{rem:fo-negative-values} give the
mathematics behind them, one topic each.
}
\begin{remark}[Two reported numbers and one first-order claim]\label{rem:fo-interface}
The objective at the point a run returned is an upper estimate by
Theorem~\ref{thm:firstorder}\,(i), respectively~(iiia), which is the part a
run is in: its point is established, where it is established at all, only
in the \termref{term-relaxed-feasible-set}{relaxed feasible set} (Remark~\ref{rem:fo-iterate-checks}), and
Remark~\ref{rem:fo-returned-point} says what a run records of it, which is
less than the whole premise. Where the hypotheses of that theorem held, the
run reports $\hlb$ as well; Requirement~\ref{req:first-order-withdrawal} says
when it does not.

Among those hypotheses is the membership of the point consumed at an iterate
in the \emph{product} $\acSet\times\Feas^{J}(\varepsilon)$: it is a
hypothesis of Theorem~\ref{thm:firstorder}\,(iiib) as much as (H-POS),
(H-SUPP) and (H-NOPERT). Its two factors are established by different
means: the $J$ factor by the three checks of
Remark~\ref{rem:fo-iterate-checks}, after any repair of
Remark~\ref{rem:fo-iterate-repair}, and the acceptance factor by forming
$\mathbf q_t$ on the \exttermref{../foundations/contract.pdf}{term-exact-island}{exact island} under ($\alpha$) and ($\beta$) of
Remark~\ref{rem:acc-side}. A failure of either leaves the membership not
established.

No hypothesis is read of the solver. A return that is merely near the
minimum satisfies no inequality in the safe direction, which is why no
bracket pairs the gradient with it: both sides of a bracket are lower
bounds on the two minima over the exact sets, the feasible one verified
from the solver's multipliers in exact arithmetic, the rounding of the
frame its rows are built with deducted
(Requirement~\ref{req:verified-feasible-side}). So a solve that stopped
short of its optimum, or failed, can loosen a bracket but cannot raise
either side above the exact minimum of the linear form at the gradient as
computed, and neither its return nor its status withdraws one; what the
gradient's own binary64 evaluation leaves is
Remark~\ref{rem:fo-arithmetic}'s. Withdrawal is never a weakening of what
$\hlb$ means; there is no reduced form of the claim.
\end{remark}
\claimmeta{\metalabel{Uses}\mref{Def.}{def:relaxed-feasible-set}, \mref{Def.}{def:first-order-setting}, \mref{Thm.}{thm:firstorder}, \mref{Rem.}{rem:acc-side}
\metasep \metalabel{Used by}\mref{Req.}{req:first-order-withdrawal}
\metasep newly derived
\metasep accepted}

% claim meat-core/frank-wolfe/rem:fo-iterate-checks
\claimsummary{meat-core/frank-wolfe/rem:fo-iterate-checks}{The three iterate checks decide every condition of the relaxed feasible set, all of them on the J factor; the acceptance factor holds by construction. They serve the membership hypothesis and the estimate's premise; no check of a solver's return would serve the bound.}
\claimbookmark{meat-core/frank-wolfe/rem:fo-iterate-checks}{Remark}{rem:fo-iterate-checks}{What the iterate checks establish}
\begin{remark}[What the iterate checks establish]\label{rem:fo-iterate-checks}
Requirement~\ref{req:iterate-checks} decides, on the point consumed at each
iterate, every condition defining $\Feas^{J}$ in \eqref{eq:choi-feas}:
per-block positive semidefiniteness, the coupling condition
$\sum_s\tr_{\Qsys_s}J_s=\idon{A_0}$, \emph{and the declared constraint rows}
$\bowtie_k$. Checking two of the three is not checking the condition.
\emph{No such check is made on a point a solver returned, and none would
serve the bound}: verifying $S_t\in\Feas^{J}$ gives
$\langle\lmoJ_t,S_t\rangle\ge\min_{J'\in\Feas^{J}}\langle\lmoJ_t,J'\rangle$,
the \emph{opposite} inequality to \eqref{eq:lmo}. Feasibility bounds the
linear form from below, which is exactly the side the first-order claim
cannot use; one can verify \emph{which set} a solver minimised over, never
that it minimised. The feasible side of a bracket is therefore a lower bound
formed without any point of the solver
(Requirement~\ref{req:verified-feasible-side}). What the checks do buy is
the membership hypothesis of Theorem~\ref{thm:firstorder}\,(iiib) on its
$J$ factor, and Theorem~\ref{thm:firstorder}\,(i),
respectively~(iiia): the point the run reports its estimate at is a convex
combination of oracle returns; \eqref{eq:fo-estimate} holds only for points
of $\acSet\times\Feas^{J}$, which is what a run does \emph{not} establish,
and (iiia) is the part it is in, giving
$O^{J}_{\alpha}(x)\ge\hlow[\Feas^{J}(\varepsilon)]$ at the relaxed
membership the checks do establish --- on the point the run returned, the
membership the end-of-run measurement of
Remark~\ref{rem:fo-returned-point} establishes, which is of the $J$ factor
only. So the checks serve the bound's membership hypothesis and the premise
of the estimate tier, which is the tier reported on every run. Their windows are numbers of their own and not
the declared tolerance $\tnum$, the object checked being one the optimiser
computed (\xref{foundations/req:tolerance-scope}\,(ii)).

Between them the three checks are what membership in
$\Feas^{J}(\varepsilon)$ at the larger of the two declared windows asks
for, and they cover all of it: per-block positive semidefiniteness, which
the family does not relax and which is decided on the point consumed at
exactly $0$ and at no window (Requirement~\ref{req:iterate-repair} says
what may be done to a block that fails it, and the record it must carry);
the coupling residual and each $\bowtie_k$ residual, each inside its own
window, hence inside the larger of the two by monotonicity
(Proposition~\ref{prop:relaxed-set-properties}\,(a)); and the family's
normalisation clause with them, that clause being implied by the coupling
clause and adding no check
(Proposition~\ref{prop:relaxed-set-properties}\,(c$'$)). No fourth check is
needed.

\emph{All three are conditions on the $J$ factor.} The hypothesis they
serve is on a \emph{product}, $x_t\in\acSet\times\Feas^{J}(\varepsilon)$,
and nothing in these checks touches the other factor: \emph{the three
checks establish the $J$ half and no more, and the $\acSet$ half is not
checked at all}. It is instead true \emph{by construction}, the
acceptance-side iterate being formed in exact rational arithmetic on the
\exttermref{../foundations/contract.pdf}{term-exact-island}{exact island} under ($\alpha$) and ($\beta$) of
Remark~\ref{rem:acc-side}: a convex combination, with exact rational
weights, of exact points of a convex set, so $\mathbf q_t\in\acSet$ by
induction. Remark~\ref{rem:acc-side} names the three consumers of
Theorem~\ref{thm:firstorder} that break if it does not hold. A run that
forms $\mathbf q_t$ in binary64 instead has not met the membership
hypothesis (Requirement~\ref{req:acceptance-iterate-exact}); there is no
check for it to pass on that factor, and none is offered. The coverage
claim of the previous paragraph is therefore a claim about one factor of a
product.

\emph{The window on the positivity check decides whether the coverage claim
is true.} At a positive window $w$ the positivity check would accept a
point with $\lmin\in(-w,0)$, which lies in \emph{no}
$\Feas^{J}(\varepsilon)$ for any $\varepsilon$ --- the family relaxes the
coupling and the $\bowtie_k$ rows and relaxes positivity not at all
(Proposition~\ref{prop:relaxed-set-properties}\,(c), \emph{Fourth}) --- so
at such a point $O^{J}_{\alpha}$ is not even defined and neither (i) nor
(ii) nor (iii) of Theorem~\ref{thm:firstorder} applies. That is why
positivity is decided at $0$, and the drift a declared window would let
pass silently is repaired and recorded under
Requirement~\ref{req:iterate-repair}. By
Theorem~\ref{thm:firstorder}\,(iiia) that reading is the conservative one:
membership in the wider $\Feas^{J}(\varepsilon)$ gives
$O^{J}_{\alpha}(x)\ge\hlow[\Feas^{J}(\varepsilon)]$, the \emph{Choi-form}
relaxed infimum \eqref{eq:hlow-feasJ}, and
$\hlow[\Feas^{J}(\varepsilon)]\le\hlow$ by \eqref{eq:feas-tol-mono-J}, so
the estimate over-estimates a quantity that is itself no larger than
$\hlow$. (The state-form $\hlow[\Feas(\varepsilon)]$ is a different number
and is not what the checks buy;
Proposition~\ref{prop:relaxed-set-properties}\,(d$'$) says why the two are
not interchangeable at $\varepsilon>0$.)
\end{remark}
\claimmeta{\metalabel{Uses}\mref{Thm.}{thm:choi}, \mref{Def.}{def:relaxed-feasible-set}, \mref{Prop.}{prop:relaxed-set-properties}, \mref{Def.}{def:first-order-setting}, \mref{Thm.}{thm:firstorder}, \mref{Rem.}{rem:acc-side}
\metasep newly derived
\metasep accepted}

% claim meat-core/frank-wolfe/rem:fo-returned-point
\claimsummary{meat-core/frank-wolfe/rem:fo-returned-point}{The point a run returns is checked once more, on its J factor only, so a pass establishes that factor of the estimate's premise and a failure leaves the premise unestablished. Neither h\_lb nor the key length rests on that point, which bears on a run only through the refusal of a negative estimate.}
\claimbookmark{meat-core/frank-wolfe/rem:fo-returned-point}{Remark}{rem:fo-returned-point}{What the measurement of the returned point establishes}
\begin{remark}[What the measurement of the returned point establishes]\label{rem:fo-returned-point}
The checks of Remark~\ref{rem:fo-iterate-checks} are made on the point
consumed at each iterate, and the point a run returns need not be one of
those: a run that stops at its iteration limit returns a point formed after
the last point it consumed, and
Requirement~\ref{req:returned-point-measured} applies the three checks to it
once more. The measurement changes nothing about the point, and it is
binary64 arithmetic like every application of those checks (the third
exposure of Remark~\ref{rem:fo-check-exposures}).

\emph{What the record says is about the $J$ factor only.} The three checks
are the $J$-side checks of Remark~\ref{rem:fo-iterate-checks}; the premise is
membership in the \emph{product} $\acSet\times\Feas^{J}(\varepsilon)$, and
its $\acSet$ factor is not checked but holds by construction, and only where
the acceptance-side point was formed on the \exttermref{../foundations/contract.pdf}{term-exact-island}{exact island} under
($\alpha$) and ($\beta$) of Remark~\ref{rem:acc-side}. Whether the
acceptance-side point the run returned was so formed is \emph{not} part of
this measurement: the run's record states that property of the points
consumed at the iterates (Requirement~\ref{req:acceptance-iterate-exact})
and does not state it of the returned point as such; where that point is the
last one consumed, the record's entry for that iterate carries it.

\emph{When the point passes all three}, its membership in
$\Feas^{J}(\varepsilon)$ is established as far as those checks establish
any membership; the premise of Theorem~\ref{thm:firstorder}\,(iiia) is then
established at that point exactly when its acceptance-side component was
formed on the exact island, a property the record does not state of the
returned point as such. \emph{When it fails any of the three}, the $J$ half
of that premise, and so the premise of Theorem~\ref{thm:firstorder}\,(i),
respectively~(iiia), is \emph{not established at that point}, so the
reported objective is not established to be an upper estimate even of
$\hlow[\Feas^{J}(\varepsilon)]$, and the record says so. The estimate is
still reported, as on every run: it is never a bound, so there is no claim
to withdraw from it, and the record names the premise it lacks rather than
hiding the number.

\emph{Neither the first-order bound nor the key length rests on any
property of that point.} $\hlb$ is the maximum in \eqref{eq:hlb} over the
iterates of $\Itset$, each of which met every condition of
Requirement~\ref{req:first-order-withdrawal}; under a non-empty
\termref{term-dropped-set}{dropped set} $\mathcal{D}$ the bound reported instead is
$\underline{h}^{\mathrm{FO},\mathcal{D}}_{\alpha}$ of
\eqref{eq:hlb-dropped}, the same expression read on the dropped objective
over the iterates at which the hypotheses of the \termref{term-dropped-program}{dropped program}
held (Corollary~\ref{cor:fo-dropped}). The key length reads this submodule
only through $\hlb$, respectively
$\underline{h}^{\mathrm{FO},\mathcal{D}}_{\alpha}$
(Corollary~\ref{cor:fo-keylength}, respectively
Corollary~\ref{cor:keylength-bound} at
$h=\underline{h}^{\mathrm{FO},\mathcal{D}}_{\alpha}$ by
Corollary~\ref{cor:fo-dropped}); the measurement of the returned point
enters neither the bound, nor the tier, nor the key length. The one way that
point bears on a run is Requirement~\ref{req:negative-estimate-refused}: a
negative objective there refuses the run, and no key length is emitted on a
refused run --- a refusal, never a changed value.
\end{remark}
\claimmeta{\metalabel{Uses}\mref{Cor.}{cor:keylength-bound}, \mref{Thm.}{thm:firstorder}, \mref{Cor.}{cor:fo-keylength}, \mref{Cor.}{cor:fo-dropped}, \mref{Rem.}{rem:acc-side}
\metasep newly derived
\metasep accepted}

% claim meat-core/frank-wolfe/rem:fo-one-ruler
\claimsummary{meat-core/frank-wolfe/rem:fo-one-ruler}{(H-SUPP) is equivalent to rank X = dim W, never to positivity of X on its ambient space, whose spectrum always contains 0. Both sides are counted relative to a spectral radius at one cutoff, the reference rank at its square root, since a direction of singular value \ensuremath{\sigma} enters X with weight \ensuremath{\sigma}\ensuremath{^{2}}.}
\claimbookmark{meat-core/frank-wolfe/rem:fo-one-ruler}{Remark}{rem:fo-one-ruler}{Maximal support as a rank equality read with one ruler}
\begin{remark}[Maximal support as a rank equality read with one ruler]\label{rem:fo-one-ruler}
(H-SUPP) is decided by a rank comparison at one declared cutoff, not by a
proof. (H-SUPP) says $\operatorname{supp}X_{\hat c,s}=W_{\hat c,s}$, and
since $\operatorname{supp}X_{\hat c,s}\subseteq W_{\hat c,s}$ always
(\eqref{eq:Wcs-bound}), it is equivalent to the \emph{rank equality}
\begin{equation}\label{eq:hsupp-check}
  \operatorname{rank}X_{\hat c,s}=\dim W_{\hat c,s},
\end{equation}
equivalently to strict positivity of the compression of $X_{\hat c,s}$ to
$W_{\hat c,s}$. \emph{It is not, and must not be read as, strict positivity
of $X_{\hat c,s}$ on its ambient space}: $X_{\hat c,s}$ lives on
$\Hs{R\,XY\,C_AC_B\,A\Qsys_s}$ while $\dim W_{\hat c,s}\le
|\Kann_{\mathrm{pass}}(\hat c)|\cdot\dim\Sadm\cdot d_s$ by
\eqref{eq:Kstack-rank}, so the ambient spectrum always contains $0$ and an
ambient positivity test is never satisfied.
Requirement~\ref{req:support-rank-check} computes both sides of
\eqref{eq:hsupp-check} and reads both with one declared cutoff: the left side
as the count of eigenvalues of $X_{\hat c,s}$ above the cutoff,
\emph{relative to that block's own spectral radius}; the right side as
$\operatorname{rank}\mathbf K_{\hat c,s}$ of \eqref{eq:Kstack}, read at the
square root of the cutoff relative to that matrix's largest singular value,
equivalently the count of eigenvalues of
$\mathbf K_{\hat c,s}\mathbf K_{\hat c,s}^{\dagger}$ above the cutoff
relative to its spectral radius, which is the left side's own count taken at
the identity input ($J_s=\idon{A_0\Qsys_s}$, respectively
$\omega_s=\idon{A\Qsys_s}$ in \exttermref{../transmitter/contract.pdf}{term-untrusted-source-mode}{untrusted-source mode}). The relative
reading is right rather than convenient, because \eqref{eq:Wcs-bound}'s
condition is a rank condition, invariant under scaling the block, and
$\nabla T$ is positively homogeneous of degree zero, so a uniformly small
block is not a degenerate one.

\emph{Why one ruler.} The branch operator is
$X_{\hat c,s}=\mathbf K_{\hat c,s}(\id\otimes J_s)
\mathbf K_{\hat c,s}^{\dagger}$, the identity acting on the index of the
\termref{term-pass-class}{pass classes} (with $\omega_s$ for $J_s$ in
untrusted-source mode), so
$\lambda_{\min}(J_s)\,\mathbf K_{\hat c,s}\mathbf K_{\hat c,s}^{\dagger}
\preceq X_{\hat c,s}\preceq
\lambda_{\max}(J_s)\,\mathbf K_{\hat c,s}\mathbf K_{\hat c,s}^{\dagger}$:
a direction along a singular value $\sigma$ of $\mathbf K_{\hat c,s}$
carries weight of order $\sigma^2$ in $X_{\hat c,s}$, never of order
$\sigma$, and $\operatorname{supp}X_{\hat c,s}\subseteq
\operatorname{supp}\mathbf K_{\hat c,s}\mathbf K_{\hat c,s}^{\dagger}
=W_{\hat c,s}$, which is \eqref{eq:Wcs-bound}. A window at the cutoff's
own value on $\sigma$ would count in $\dim W_{\hat c,s}$ directions whose
relative weight in $X_{\hat c,s}$ is below that value at every iterate with
$J_s$ not far from a multiple of the identity --- for a direction at
relative singular value $\sigma_{\mathrm{rel}}$, every iterate at which
$\lambda_{\max}(J_s)/\lambda_{\min}(J_s)$ is at most the cutoff divided by
$\sigma_{\mathrm{rel}}^2$ --- and the comparison would then fail at every
such iterate however exactly (H-SUPP) holds; read at the square root, both
sides count eigenvalues of operators of the one form
$\mathbf K_{\hat c,s}(\id\otimes J)\mathbf K_{\hat c,s}^{\dagger}$ at one
relative window.
\end{remark}
\claimmeta{\metalabel{Uses}\mref{Def.}{def:branch-support}, \mref{Lem.}{lem:branch-support}, \mref{Thm.}{thm:firstorder}
\metasep newly derived
\metasep accepted}

% claim meat-core/frank-wolfe/rem:fo-reachability
\claimsummary{meat-core/frank-wolfe/rem:fo-reachability}{An iterate with positive weight on a point of positive definite Choi blocks and an acceptance vector positive on the acceptance support meets (H-SUPP) and (H-POS)'s first clause; such a point exists if the declared rows admit one. No optimiser is fixed: a run without such a component may lack h\_lb.}
\claimbookmark{meat-core/frank-wolfe/rem:fo-reachability}{Remark}{rem:fo-reachability}{Reachability of positivity and maximal support}
\begin{remark}[Reachability of positivity and maximal support]\label{rem:fo-reachability}
What makes (H-SUPP) and (H-POS) reachable at all is a separate question from
what discharges them. $J_{t,s}\succ0$ suffices for (H-SUPP) by
\eqref{eq:Wcs-bound}, and $q_t$ strictly positive on $\Chacc$ gives the
first clause of (H-POS); neither is automatic, and an oracle return is
typically an extreme point of its feasible set and so neither strictly
positive nor of full support on $\Chacc$. What supplies both is the
\emph{formation} of the iterate: a point that carries a strictly positive
weight on some $x_{0}\in\acSet\times\Feas^{J}$ with $J_{0,s}\succ0$ for
every $s$ and $q_{0}(\hat c)>0$ for every $\hat c\in\Chacc$ inherits both
properties, a sum of a positive definite operator and positive semidefinite
ones being positive definite. Such an $x_0$ exists whenever the declared
rows admit one --- with $\CList=\varnothing$,
$J_{0,s}=|\SecSet|^{-1}\,\idon{A_0}\otimes d_s^{-1}\idon{\Qsys_s}$ is
feasible and strictly positive, and Lemma~\ref{lem:fo-khyp}\,(i) supplies a
$\mathbf q_{0}\in\acSet$ supported on all of $\Chacc$. This module
prescribes no optimiser and does not require one: the statement is what the
first-order tier costs, not an algorithm. A run whose iterates carry no such
component may legitimately fail (H-POS) or (H-SUPP) at every iterate, in
which case $\Itset$ is empty, $\hlb$ is absent, and the run reports its
estimate, correctly: the tier is conditional on how the iterates are formed,
not unreachable.
\end{remark}
\claimmeta{\metalabel{Uses}\mref{Lem.}{lem:branch-support}, \mref{Lem.}{lem:fo-khyp}, \mref{Thm.}{thm:firstorder}
\metasep newly derived
\metasep accepted}

% claim meat-core/frank-wolfe/rem:fo-check-exposures
\claimsummary{meat-core/frank-wolfe/rem:fo-check-exposures}{The rank comparison for (H-SUPP) and the three iterate checks are computed in binary64 and leave three one-sided exposures, two on the unsafe side, that no claim of this module bounds. A reference rank computed too large is harmless, and the acceptance factor, exact by construction, adds none.}
\claimbookmark{meat-core/frank-wolfe/rem:fo-check-exposures}{Remark}{rem:fo-check-exposures}{One-sided exposures of the runtime checks}
\begin{remark}[One-sided exposures of the runtime checks]\label{rem:fo-check-exposures}
The rank comparison of Remark~\ref{rem:fo-one-ruler} buys a refusal on the
safe side in the common case, a rank-deficient block being refused rather
than differentiated. It does not buy the hypothesis, and there are
\emph{three} one-sided residual exposures, the first two on the unsafe side
of Remark~\ref{rem:fo-danger} and the third on the membership hypothesis.
\emph{First}, a block whose true support is a proper subspace of
$W_{\hat c,s}$ but whose smallest relevant eigenvalue is reported above the
cutoff by the arithmetic of the spectral decomposition is differentiated on
a non-maximal face silently. \emph{Second}, a $\dim W_{\hat c,s}$ computed
\emph{too small} --- a direction of $W_{\hat c,s}$ that is genuinely present
but whose relative weight in
$\mathbf K_{\hat c,s}\mathbf K_{\hat c,s}^{\dagger}$, $\sigma^2$ over the
largest $\sigma^2$, is at most the cutoff, and which is therefore not
counted --- lowers the bar in \eqref{eq:hsupp-check} and passes a
rank-deficient block for the same reason. A $\dim W_{\hat c,s}$ computed too
\emph{large} is harmless for the claim: it can only fail the comparison,
which removes the iterate's bracket from $\hlb$ --- at every iterate, so
that $\hlb$ is withdrawn, if it counts a direction no
$\operatorname{rank}X_{\hat c,s}$ can count, which is why the two sides are
read with one ruler. \emph{Third}, and belonging to the checks of
Remark~\ref{rem:fo-iterate-checks} rather than to the rank comparison: each
of the three quantities those checks compare is itself computed in binary64.
$\lmin$ of an iterate block is computed, so a block whose true minimum
eigenvalue is negative within the eigensolver's own backward error passes at
exactly $0$ and is not repaired, and the reassembly of a block that
\emph{is} repaired is itself rounded; and the coupling residual
$\|\sum_s\tr_{\Qsys_s}J_s-\idon{A_0}\|_{\max}$ and each $\bowtie_k$ residual
are formed by partial traces and sums that round, so a point whose
\emph{true} residual exceeds its window can be measured inside it. So the
consumed point lies in $\Feas^{J}(\varepsilon)$ by construction only up to
that arithmetic.

\emph{The acceptance factor adds no fourth exposure to this list}: it is not
measured at all, being exact by construction on the \exttermref{../foundations/contract.pdf}{term-exact-island}{exact island}
(Remark~\ref{rem:acc-side}), so the count is three because the three are the
$J$-side checks, and the product's other factor is established in a way
that has no residual to expose. Nor does the feasible side of a bracket:
its value is a certification computed exactly from the numbers it is given
(Requirement~\ref{req:verified-feasible-side}) and decides nothing at a
window, the rounding of the frame its rows are built with being charged
within it, and how far the gradient it is given lies from the exact one is
the binary64 evaluation of Remark~\ref{rem:fo-arithmetic}, not a check. No
claim of this module bounds any of the
three errors; what the declared thresholds and the exact-$0$ check buy is
that each comparison is well posed, that its reference is fixed by this
module rather than chosen by an implementation, and that no window
\emph{this module declares} admits a point at which the theorem does not
apply.
\end{remark}
\claimmeta{\metalabel{Uses}\mref{Thm.}{thm:firstorder}, \mref{Rem.}{rem:fo-danger}, \mref{Rem.}{rem:acc-side}
\metasep \metalabel{Used by}\mref{Rem.}{rem:certified-scope}
\metasep newly derived
\metasep accepted}

% claim meat-core/frank-wolfe/rem:fo-oracle-repair
\claimsummary{meat-core/frank-wolfe/rem:fo-oracle-repair}{Repairs are permitted on computed objects and never on declared data. A repaired oracle return feeds only the search and a recorded diagnostic, and its record bounds nothing: a window on the block bounds no error on a pairing whose gradient can grow without limit.}
\claimbookmark{meat-core/frank-wolfe/rem:fo-oracle-repair}{Remark}{rem:fo-oracle-repair}{A repaired oracle return as a step of the search}
\begin{remark}[A repaired oracle return as a step of the search]\label{rem:fo-oracle-repair}
This remark and Remark~\ref{rem:fo-iterate-repair} concern two positivity
repairs, on two different objects. Both are on objects this module
computed, so both are permitted where no repair of \exttermref{../foundations/contract.pdf}{term-declared-data}{declared data}
ever is (contrast Requirement~\ref{req:exact-admission}: declared data is
exactly admissible or refused, and this module repairs none of it). The
first, on an \emph{oracle return}, touches only the search; the second, on
a \emph{formed iterate}, is a construction step that makes a hypothesis of
Theorem~\ref{thm:firstorder} hold by construction.

\emph{On an oracle return: a step of the search.} A backend may return a
block whose spectrum is slightly negative, and
Requirement~\ref{req:oracle-repair} lets the engine reject the return or
project the block within a window of its own. The repaired return is a
direction for the search and the point of the recorded diagnostic pairing;
no bracket of $\hlb$ reads it, the feasible side of every bracket being a
lower bound formed from the backend's multipliers
(Requirement~\ref{req:verified-feasible-side}), so the repair needs no
account in any claim of this module. Nor could its record supply one: a
declared window bounds $\|S_t^{+}-S_t\|$, but converting that into an error
on $\langle\lmoJ_t,S_t\rangle$ needs a bound on $\|\lmoJ_t\|$, and
\eqref{eq:grad-q}--\eqref{eq:grad-J} carry the ratios $q/P^{J}_{\hat c}$
and $q/K^{J}_{\hat c}$, which Proposition~\ref{prop:value-range}\,(iii)
exhibits declarations driving up without limit.
\end{remark}
\claimmeta{\metalabel{Uses}\mref{Prop.}{prop:value-range}, \mref{Cor.}{cor:full-gradient}, \mref{Thm.}{thm:firstorder}
\metasep \metalabel{Used by}\mref{Req.}{req:oracle-repair}
\metasep newly derived
\metasep accepted}

% claim meat-core/frank-wolfe/rem:fo-iterate-repair
\claimsummary{meat-core/frank-wolfe/rem:fo-iterate-repair}{A block repaired while the iterate is formed makes the repaired point the iterate, at which every term of the bracket is evaluated, so no perturbation that (H-NOPERT) forbids occurs. The repair moves the coupling residual by at most the dimension times its size and says nothing about the backend.}
\claimbookmark{meat-core/frank-wolfe/rem:fo-iterate-repair}{Remark}{rem:fo-iterate-repair}{A repaired iterate as a construction step}
\begin{remark}[A repaired iterate as a construction step]\label{rem:fo-iterate-repair}
\emph{On a formed iterate: a construction step, and the only repair this
module licenses on an iterate.} The point consumed at iterate $t$ is formed
by the engine --- in the intended use a convex combination of oracle
returns --- and Requirement~\ref{req:iterate-checks} requires it to be
per-block positive semidefinite at exactly $0$; a block failing that check
is rejected or projected (Requirement~\ref{req:iterate-repair}).
\emph{The repaired point is $x_t$.} The repair belongs to the formation of
the iterate, upstream of the point the theorem is applied at, and it is a
different act from the oracle repair of Remark~\ref{rem:fo-oracle-repair},
which touches only the search; everything in \eqref{eq:hlb} --- the
objective $O^{J}_{\alpha}(x_t)$, the gradient $(\lmoq_t,\lmoJ_t)$ and the
pairings with $x_t$ --- is then evaluated at that one point. \emph{This is not the perturbation route
Remark~\ref{rem:fo-nopert} closes, and (H-NOPERT) is not violated}: that
route takes the gradient at a perturbed operator and concludes something
about the objective at the unperturbed one, which needs a domination
statement this module does not prove, whereas here no second function
enters the argument at all --- Theorem~\ref{thm:firstorder} concludes about
the infimum over the whole set and asks of $x_t$ only that it lie in
$\acSet\times\Feas^{J}(\varepsilon)$. What (H-NOPERT) forbids is a
perturbation applied \emph{to} $x_t$, between the point being fixed and the
gradient being taken; a repair applied \emph{before} the point is fixed is a
different act, and it is the one Requirement~\ref{req:iterate-repair}
licenses.

\emph{The repair moves the residuals by a bounded amount.} A projection
moves the coupling and $\bowtie_k$ residuals of
Remark~\ref{rem:fo-iterate-checks}: if the repair moves each block by at
most $\delta$ in the operator norm, the coupling residual moves by at most
$d_{\Qsys}\,\delta$ in $\|\cdot\|_{\max}$, since
$\|\tr_{\Qsys_s}\Delta_s\|\le d_s\|\Delta_s\|$, $\sum_sd_s=d_{\Qsys}$ and
$\|\cdot\|_{\max}\le\|\cdot\|$, a bound that carries its dimension factor
as \xref{foundations/req:norm-conversion} asks. The requirement nevertheless
measures the residuals again rather than bounding them in advance.
What the repair establishes is a hypothesis of
Theorem~\ref{thm:firstorder}, on an object the engine itself chose, and it
says nothing whatever about the backend.

\emph{What it does not establish.} $\lmin$ is itself computed, so a block
whose true minimum eigenvalue is negative within the eigensolver's own
backward error passes the check unrepaired, and the reassembly of a
projected block is itself rounded. That residue is a third one-sided
exposure of the kind Remark~\ref{rem:fo-check-exposures} names, and it is
named there with the other two. Its \emph{nature} is the unbudgeted
arithmetic this module declares everywhere (Remark~\ref{rem:fo-arithmetic}),
and not a window this module chose while knowing it admitted points at which
its own theorem does not apply.
\end{remark}
\claimmeta{\metalabel{Uses}\mref{Thm.}{thm:firstorder}, \mref{Rem.}{rem:fo-nopert}, \mref{Rem.}{rem:fo-arithmetic}
\metasep \metalabel{Used by}\mref{Req.}{req:iterate-repair}
\metasep newly derived
\metasep accepted}

% claim meat-core/frank-wolfe/rem:fo-negative-values
\claimsummary{meat-core/frank-wolfe/rem:fo-negative-values}{A negative computed h\_lb gives way to max\{0, h\_lb\}, itself a valid lower bound on h\_low. The objective is non-negative on the exact feasible set but not on the relaxed one, so a negative estimate is evidence that a hypothesis fails, and it is refused rather than replaced by zero.}
\claimbookmark{meat-core/frank-wolfe/rem:fo-negative-values}{Remark}{rem:fo-negative-values}{Opposite treatment of a negative bound and a negative estimate}
\begin{remark}[Opposite treatment of a negative bound and a negative estimate]\label{rem:fo-negative-values}
A negative number can arise on each path, both cases in floating point, and
the two are handled oppositely.

\emph{On the first-order path.} A computed $\hlb$ may be negative --- the
linearisation is weak far from the optimum, or the arithmetic crosses
zero --- and what reaches Corollary~\ref{cor:fo-keylength} is
$\max\{0,\hlb\}$ (Requirement~\ref{req:first-order-clamp}), which is a
valid lower bound on $\hlow$ by Theorem~\ref{thm:firstorder}\,(ii),
respectively~(iiib) --- the second being the part a run is in, since (ii)
needs exact membership and Remark~\ref{rem:fo-iterate-checks} says no run
establishes it; the two conclude the same inequality against the same exact
$\hlow$, (ii) being the case $\varepsilon=0$ of (iiib). This is not a repair
of anybody's data; it is the substitution of one valid lower bound for
another, and Remark~\ref{rem:fo-arithmetic}\,(a) says what it does not fix.

\emph{On the estimate path.} A computed objective at a point of the
\emph{exact} feasible set $\Feas^{J}$ cannot be negative: writing
$O^{J}_{\alpha}$ as in the proof of
Proposition~\ref{prop:value-range}\,(i), it is a non-negative multiple of a
KL divergence between two probability vectors plus a non-negative
combination of the branch terms
$e_{\hat c}=\frac{1}{1-\alpha}\log_2\frac{K_{\hat c}}{P_{\hat c}}$, each
non-negative because $K_{\hat c}\le P_{\hat c}$
(Theorem~\ref{thm:branch-formula}, \emph{cone clause}) and
$\frac{1}{1-\alpha}<0$; that argument is pointwise and needs no infimum.
\emph{At a point of the relaxed set $\Feas^{J}(\varepsilon)$ it can be
negative}, and the statement is written on the exact set for that reason:
there the \emph{second} argument of the divergence term is not a
probability vector, and by
Proposition~\ref{prop:relaxed-set-properties}\,(c$''$) the objective is
bounded below by $-\frac{\alpha}{\alpha-1}\log_2(1+d_{A_0}\varepsilon)$ and
by nothing larger. A run reads its feasible set at the two declared windows
of Requirement~\ref{req:iterate-checks}, so it has a floor below zero, and
the refusal of Requirement~\ref{req:negative-estimate-refused} can fire
against a point that is otherwise in order. That is intended: the refusal is
on the safe side, no key length is emitted on a refused run, and what is
refused is an estimate and not a bound. So a negative computed estimate is
\emph{evidence that a hypothesis of this module does not hold} of the point
it was computed at --- most likely that the point is not in the exact set
(Remark~\ref{rem:fo-iterate-checks}), that its distance from it exceeds the
scale just named, or that the declaration was not exactly admissible, which
its admission excludes (Definition~\ref{def:declaration}). Unlike $\hlb$ the
estimate is not a lower bound, and zero cannot stand in for it: replacing it
by zero would report a plausible number in place of a detected
inconsistency.
\end{remark}
\claimmeta{\metalabel{Uses}\mref{Thm.}{thm:branch-formula}, \mref{Prop.}{prop:value-range}, \mref{Prop.}{prop:relaxed-set-properties}, \mref{Thm.}{thm:firstorder}, \mref{Cor.}{cor:fo-keylength}, \mref{Rem.}{rem:fo-arithmetic}
\metasep newly derived
\metasep accepted}

\section{Requirements: reported numbers and withdrawal}\label{group-meat-core/frank-wolfe/req:first-order-withdrawal}
\sectionintro{A run reports its estimate and, where the hypotheses held, the first-order bound; an iterate at which a hypothesis fails or is not established adds no bracket to the bound, with a reason, and with no iterate left the engine withdraws the bound. The backend's return and status withdraw nothing. Where a reference companion is admitted, the run's number is the bracket on the companion at the best iterates of the search and its last and at the points of a short search continued on the companion, the bound on the searched declaration a diagnostic beside it. It clamps the bound at zero with a flag, refuses a negative estimate, and imports no penalty from another formulation. It measures the returned point once more and records the difference between the bound and the estimate, resting nothing on either.}
\sectionimports{\extref{../foundations/contract.pdf}{foundations/def:evaluated-off-snap}, \extref{../foundations/contract.pdf}{foundations/req:snap-reference}}

% claim meat-core/frank-wolfe/req:first-order-withdrawal
\claimsummary{meat-core/frank-wolfe/req:first-order-withdrawal}{An iterate at which a first-order hypothesis fails or is not established contributes no bracket and its reason is recorded; with none left h\_lb is withdrawn and the run falls to the estimate tier. The oracle's return and its status withdraw nothing.}
\claimbookmark{meat-core/frank-wolfe/req:first-order-withdrawal}{Requirement}{req:first-order-withdrawal}{Withdrawal of the first-order bound}
\claimleadin{%
Every run reports its \exttermref{../foundations/contract.pdf}{term-estimate}{estimate} and names the \exttermref{../foundations/contract.pdf}{term-tier}{tier} that carries
its claim (\xref{foundations/req:tiers}). The \exttermref{../foundations/contract.pdf}{term-first-order}{first-order} bound $\hlb$
of Theorem~\ref{thm:firstorder} holds only where the theorem's hypotheses
hold at the iterates its maximum runs over, and a gradient taken after a
perturbation bounds another function (Remark~\ref{rem:fo-nopert}).
No hypothesis of that theorem is read of the backend: the feasible side of
a bracket is a lower bound formed without its return or its status
(Requirement~\ref{req:verified-feasible-side}), so neither withdraws a
bracket. Withdrawal is never a weakening of what $\hlb$ means
(Remark~\ref{rem:fo-interface}).
}
\begin{requirement}[Withdrawal of the first-order bound]\label{req:first-order-withdrawal}
An iterate at which any of the following holds is not a member of
$\Itset$: it contributes no bracket to the maximum in \eqref{eq:hlb}, and the
engine records the machine-readable reason.
\begin{enumerate}[label=(\roman*),leftmargin=2.4em]
\item (H-POS) fails (\iface{gradient_hypothesis_positivity});
\item (H-SUPP) fails (\iface{gradient_hypothesis_support});
\item a gradient perturbation was applied, that is, (H-NOPERT) fails
  (\iface{gradient_perturbation_applied});
\item the point consumed at that iterate fails one of the three checks of
  Requirement~\ref{req:iterate-checks} and is not repairable under
  Requirement~\ref{req:iterate-repair}, or its acceptance-side component
  $\mathbf q_t$ was not formed on the \exttermref{../foundations/contract.pdf}{term-exact-island}{exact island} under ($\alpha$)
  and ($\beta$) of Remark~\ref{rem:acc-side}, so that its membership in the
  product $\acSet\times\Feas^{J}(\varepsilon)$ is not established
  (\iface{iterate_not_feasible} in either case).
\end{enumerate}
Neither the point the oracle returned at that iterate nor its status
removes the iterate from $\Itset$.
If no iterate is left, $\Itset$ is empty: the engine withdraws $\hlb$ and
the tier falls to the estimate, with a reason from the same
machine-readable vocabulary.
\end{requirement}
\claimmeta{\metalabel{Motivated by}\mref{Thm.}{thm:firstorder}, \mref{Rem.}{rem:fo-nopert}, \mref{Rem.}{rem:acc-side}, \mref{Rem.}{rem:fo-interface}, \mref{Req.}{req:iterate-checks}, \mref{Req.}{req:iterate-repair}, \mref{Req.}{req:verified-feasible-side}, \extref{../foundations/contract.pdf}{foundations/req:tiers}
\metasep \metalabel{Used by}\mref{Req.}{req:bracket-on-reference}, \mref{Req.}{req:certified-evaluation}
\metasep implemented (effective: accepted)
\metasep \metalabel{Code}\coderef{https://github.com/fraunhofersg/certiqs-key/blob/v0.1.0/src/certiqs_key/meat_core/reasons.py}{certiqs_key.meat_core.reasons:AssuranceReason}}

% claim meat-core/frank-wolfe/req:returned-point-measured
\claimsummary{meat-core/frank-wolfe/req:returned-point-measured}{Every run applies the three iterate checks once more to the point it returned and records whether that point passes, with the three quantities measured. Neither the bound, nor the tier, nor the key length rests on that measurement.}
\claimbookmark{meat-core/frank-wolfe/req:returned-point-measured}{Requirement}{req:returned-point-measured}{Measurement of the returned point}
\claimleadin{%
The checks of Requirement~\ref{req:iterate-checks} are made on the point
consumed at each iterate, and the point a run returns, at which its
\exttermref{../foundations/contract.pdf}{term-estimate}{estimate} is taken, need not be one of those: a run that stops at its
iteration limit returns a point formed after the last point it consumed.
What the measurement establishes of that point, the $J$ factor of the
estimate's premise and no more, is Remark~\ref{rem:fo-returned-point}.
}
\begin{requirement}[Measurement of the returned point]\label{req:returned-point-measured}
Every run applies the three checks of Requirement~\ref{req:iterate-checks}
once more, at the end of the run, to the point it returned: per-block
positivity at exactly $0$, the coupling residual at \iface{iterate_coupling}
and each $\bowtie_k$ residual at \iface{constraint_rows}. It records whether
that point passes all three, together with the three quantities measured,
in \iface{optimizer_certificate}. The measurement enters neither the
first-order bound, nor the tier, nor the key length.
\end{requirement}
\claimmeta{\metalabel{Motivated by}\mref{Thm.}{thm:firstorder}, \mref{Rem.}{rem:fo-returned-point}, \mref{Req.}{req:iterate-checks}
\metasep implemented (effective: accepted)
\metasep \metalabel{Code}\coderef{https://github.com/fraunhofersg/certiqs-key/blob/v0.1.0/src/certiqs_key/solvers/frank_wolfe/joint_driver.py}{certiqs_key.solvers.frank_wolfe.joint_driver:run_joint_frank_wolfe}}

% claim meat-core/frank-wolfe/req:first-order-clamp
\claimsummary{meat-core/frank-wolfe/req:first-order-clamp}{What reaches the key length is max\{0, h\_lb\}, and the engine reports the unclamped h\_lb beside a flag recording the clamp. The clamp substitutes one valid lower bound for another and repairs nobody's data.}
\claimbookmark{meat-core/frank-wolfe/req:first-order-clamp}{Requirement}{req:first-order-clamp}{Clamping of the first-order bound at zero}
\claimleadin{%
A computed $\hlb$ may be negative, the linearisation being weak far from the
optimum or the arithmetic crossing zero, and $\max\{0,\hlb\}$ is still a
lower bound on $\hlow$ by Theorem~\ref{thm:firstorder}\,(ii),
respectively~(iiib) (Remark~\ref{rem:fo-negative-values}); the clamp corrects
only the low side (Remark~\ref{rem:fo-arithmetic}\,(a)). The unclamped value
lets a reader see what was computed and not only what was used.
}
\begin{requirement}[Clamping of the first-order bound at zero]\label{req:first-order-clamp}
The engine passes $\max\{0,\hlb\}$ to the key length of
Corollary~\ref{cor:fo-keylength}, and records the clamp: it emits the
boolean \iface{first_order_clamped} beside \iface{h_lb}, and \iface{h_lb}
carries the unclamped value.
\end{requirement}
\claimmeta{\metalabel{Motivated by}\mref{Thm.}{thm:firstorder}, \mref{Cor.}{cor:fo-keylength}, \mref{Rem.}{rem:fo-arithmetic}, \mref{Rem.}{rem:fo-negative-values}
\metasep \metalabel{Used by}\mref{Req.}{req:bracket-on-reference}, \mref{Req.}{req:certified-evaluation}
\metasep implemented (effective: accepted)
\metasep \metalabel{Code}\coderef{https://github.com/fraunhofersg/certiqs-key/blob/v0.1.0/src/certiqs_key/solvers/frank_wolfe/general_keyrate.py}{certiqs_key.solvers.frank_wolfe.general_keyrate:key_length_from_first_order}}

% claim meat-core/frank-wolfe/req:negative-estimate-refused
\claimsummary{meat-core/frank-wolfe/req:negative-estimate-refused}{A run whose search returns a negative estimate is refused with a typed reason naming the value and the point, and the estimate is never clamped. Unlike h\_lb it is no lower bound, so zero cannot stand in for it.}
\claimbookmark{meat-core/frank-wolfe/req:negative-estimate-refused}{Requirement}{req:negative-estimate-refused}{Refusal of a negative estimate}
\claimleadin{%
The objective is non-negative at every point of the exact feasible set, by a
pointwise argument (Proposition~\ref{prop:value-range}\,(i)), while at a
point of the \termref{term-relaxed-feasible-set}{relaxed feasible set} it is bounded below only by
$-\frac{\alpha}{\alpha-1}\log_2(1+d_{A_0}\varepsilon)$, where
$d_{A_0}\varepsilon<1$, and in the state form with $d_A\varepsilon$ in
trusted-source mode and $\varepsilon$ in untrusted-source mode in place of
$d_{A_0}\varepsilon$
(Proposition~\ref{prop:relaxed-set-properties}\,(c$''$)). A negative estimate
is therefore evidence that a hypothesis does not hold of the point it was
computed at, and unlike $\hlb$ it is no lower bound
(Remark~\ref{rem:fo-negative-values}). No key length is emitted on a refused
run.
}
\begin{requirement}[Refusal of a negative estimate]\label{req:negative-estimate-refused}
The engine refuses a run whose computed \exttermref{../foundations/contract.pdf}{term-estimate}{estimate}, the objective at
the point the run's search on the declaration it searches returned at an
order the run evaluates, is negative, with the
machine-readable reason \iface{negative_objective} naming the value and the
point, and never clamps the estimate to zero. A bare exception is not a
conforming implementation of this rule; the reason must be readable by the
caller. The rule binds that search alone: the polish of
Requirement~\ref{req:bracket-on-reference}\,(i) and the solve of
Requirement~\ref{req:unreachable-acceptance-set}\,(i) form no estimate of
the run, and a negative objective value of theirs refuses nothing.
\end{requirement}
\claimmeta{\metalabel{Motivated by}\mref{Prop.}{prop:value-range}, \mref{Req.}{req:unreachable-acceptance-set}, \mref{Prop.}{prop:relaxed-set-properties}, \mref{Rem.}{rem:fo-negative-values}, \mref{Req.}{req:bracket-on-reference}
\metasep implemented (effective: accepted)
\metasep \metalabel{Code}\coderef{https://github.com/fraunhofersg/certiqs-key/blob/v0.1.0/src/certiqs_key/solvers/frank_wolfe/general_keyrate.py}{certiqs_key.solvers.frank_wolfe.general_keyrate:NegativeObjectiveError}}

% claim meat-core/frank-wolfe/req:bound-estimate-difference
\claimsummary{meat-core/frank-wolfe/req:bound-estimate-difference}{The engine records h\_lb minus the reported estimate, with its sign, as a diagnostic, and asserts nothing of it. Their ordering holds only in exact arithmetic and at an exactly feasible point, neither of which a run establishes.}
\claimbookmark{meat-core/frank-wolfe/req:bound-estimate-difference}{Requirement}{req:bound-estimate-difference}{The recorded difference between the bound and the estimate}
\claimleadin{%
The ordering of $\hlb$ below the objective at the returned point is a
consequence of exact arithmetic and of the returned point's exact
membership in $\acSet\times\Feas^{J}$, neither of which a run establishes;
a positive difference, a bound above the estimate, is evidence about the
arithmetic or the relaxation, and not by itself evidence that the
first-order tier is unsound (Remark~\ref{rem:fo-arithmetic}\,(b)).
}
\begin{requirement}[The recorded difference between the bound and the estimate]\label{req:bound-estimate-difference}
The engine records the difference $\hlb$ minus the reported
\exttermref{../foundations/contract.pdf}{term-estimate}{estimate}, with its sign, as a diagnostic, and asserts nothing of it.
\end{requirement}
\claimmeta{\metalabel{Motivated by}\mref{Rem.}{rem:fo-arithmetic}
\metasep \metalabel{Used by}\mref{Req.}{req:bracket-on-reference}
\metasep accepted
\metasep \metalabel{Code}\coderef{https://github.com/fraunhofersg/certiqs-key/blob/v0.1.0/src/certiqs_key/solvers/frank_wolfe/general_keyrate.py}{certiqs_key.solvers.frank_wolfe.general_keyrate:GeneralKeyRateResult}}

% claim meat-core/frank-wolfe/req:no-imported-penalty
\claimsummary{meat-core/frank-wolfe/req:no-imported-penalty}{The engine imports no relative-entropy penalty onto the objective: it subtracts no penalty term taken from a von Neumann entropy key-rate formulation. The objective is no relative entropy, and no claim here relates the two at finite \ensuremath{\alpha}.}
\claimbookmark{meat-core/frank-wolfe/req:no-imported-penalty}{Requirement}{req:no-imported-penalty}{The program's own objective, with no imported penalty}
\claimleadin{%
The objective of \eqref{eq:choi-program} is not a relative entropy: its
$K$-terms are R\'enyi-order-$\alpha$ trace powers, and no statement of this
module relates the two at finite $\alpha$, so a penalty term taken from a
von-Neumann-entropy key-rate formulation, however standard there, would
claim a theorem this objective does not have. Face handling here is the compression of
Theorem~\ref{thm:gradient}\,(i), unbudgeted, which is why this tier is weaker
than a certificate.
}
\begin{requirement}[The program's own objective, with no imported penalty]\label{req:no-imported-penalty}
The engine imports no relative-entropy penalty onto $O^{J}_{\alpha}$ from any
source: it subtracts from it no penalty term taken from a
von-Neumann-entropy key-rate formulation.
\end{requirement}
\claimmeta{\metalabel{Motivated by}\mref{Thm.}{thm:choi}, \mref{Thm.}{thm:gradient}, \mref{Thm.}{thm:firstorder}
\metasep implemented (effective: accepted)
\metasep \metalabel{Code}\coderef{https://github.com/fraunhofersg/certiqs-key/blob/v0.1.0/src/certiqs_key/meat_core/objective/joint.py}{certiqs_key.meat_core.objective.joint:joint_objective_value}}

% claim meat-core/frank-wolfe/req:bracket-on-reference
\claimsummary{meat-core/frank-wolfe/req:bracket-on-reference}{Where a reference companion is admitted, the run's first-order number is the best bracket on the companion at the best points of the snapped search and of a short search continued on the companion from them, each assessed by the usual gate; the snapped bound is recorded as a diagnostic only.}
\claimbookmark{meat-core/frank-wolfe/req:bracket-on-reference}{Requirement}{req:bracket-on-reference}{The first-order bound evaluated on the reference companion}
\claimleadin{%
The search runs on the snapped declaration; the bound is formed on the
companion. Any points at which Corollary~\ref{cor:fo-reference} holds give a
valid bound, so the choice of points bears on its size only. The companion
and the snapped declaration differ only in the measured families of the
factors the companion replaces and in the entries of their last links, so the
points the snapped bracket ranks highest are the natural candidates, and a
few of them keep the extra oracle calls few; the last point is added because
it is the point the search ended at. Those points can still sit badly on the
companion. The search fits the acceptance-side component to the snapped
statistics, and where the snap moved an entry that is small beside its
outcome's probability, the companion's statistics at the same point differ
from that component by a relative amount of the move's size, which the
factor $\alpha/(\alpha-1)$ of the objective magnifies: the bracket on the
companion can then fall far below the snapped one, though the two values are
close. The polish, a short search continued on the companion's own program
from the best of those points, moves the point to where the companion's
bracket is large. Its points are points like any other, since
Corollary~\ref{cor:fo-reference} reads of a point nothing but its
hypotheses, and the gate decides them at each point whatever produced it:
in exact arithmetic the points it adds can only raise a maximum of valid
brackets. That is not the whole of it at the first-order tier. Where the
companion carries a positive radius, its first-order bracket is formed from
the binary64 numbers nearest the centres of its enclosure, which are
displaced from the companion itself by an amount no claim bounds
(Remark~\ref{rem:fo-arithmetic}, Proposition~\ref{prop:optol-nobound}); a
search that maximises that evaluation makes it more likely that such an
unbudgeted excess reaches the first-order number than points found on
another declaration do. The polish's points count at the first-order tier
all the same, as the search's do; where every radius is zero only the
rounding of the evaluation is exposed, as at any point; and the certified
tier, formed from the exact centres and radii
(Requirement~\ref{req:certified-evaluation}), is unaffected. Which blocks then rest
on no premise about the snap follows from what the companion carries, block
by block (\xref{foundations/def:evaluated-off-snap}). It follows also from
which number the key length is formed from: where the companion's bound
counts but its certified evaluation certifies no point, a number certified
on the snapped declaration may carry the key length instead
(Requirement~\ref{req:certified-evaluation}), and the blocks the companion
carried then rest on the snap premise after all.
}
\begin{requirement}[The first-order bound evaluated on the reference companion]\label{req:bracket-on-reference}
Where the engine has admitted a compiled reference companion
$\Decl^{\mathrm{ref}}$ beside the declaration $\Decl$ it searched on
(Requirement~\ref{req:companion-admission}), it forms the run's first-order
number as follows, after the search on $\Decl$ has ended.
\begin{enumerate}[label=(\roman*),leftmargin=2.4em]
\item \emph{The points.} It evaluates the points of two sets. The first
  holds at most five points of the search on $\Decl$: the at most four
  members of the index set of $\Decl$'s bound
  (Requirement~\ref{req:first-order-withdrawal}) with the largest brackets
  on $\Decl$, ties broken by the order of the run, and the last point the
  search assessed; where $\Decl$'s index set is empty, the last point alone.
  The second holds the points kept by the \emph{polish}, a search on the
  Choi program of $\Decl^{\mathrm{ref}}$ that continues the search on
  $\Decl$ after it has ended. The polish reads the objective, the gradient,
  the maps and the entries of $\Decl^{\mathrm{ref}}$
  (Requirement~\ref{req:companion-admission}\,(i)) and no map generated from
  $\Decl$'s measured family; it starts at the point of the first set with
  the largest bracket on $\Decl$, at the last point where $\Decl$'s index
  set is empty, and forms its acceptance-side iterate as
  Requirement~\ref{req:acceptance-iterate-exact} requires, from that point's
  exact acceptance-side component. Where that component was not formed on
  the \exttermref{../foundations/contract.pdf}{term-exact-island}{exact island} (Requirement~\ref{req:acceptance-iterate-exact}),
  the polish does not run, the second set is empty and the record names
  that reason. It assesses each point it consumes by the
  gate of (ii) and keeps the at most four of them that contribute with the
  largest brackets on $\Decl^{\mathrm{ref}}$, ties broken by its own order,
  and the last point it assessed; a polish that assesses no point keeps
  none. It runs within the share \iface{reference_polish_share} of the
  wall budget of one evaluation of the run, where the run has one, set
  aside for it before the search on $\Decl$ begins, and ends where that
  share is spent, where an iteration cap it records is reached or where its
  own stopping rule holds, whichever comes first; a polish that fails or is
  cut short keeps, as above, the points it assessed until then. The share
  so set aside shortens the search on $\Decl$, which bears on the run's
  estimate, on $\hlb[\Decl]$, on a number certified on $\Decl$ in the
  fallback of Requirement~\ref{req:certified-evaluation}\,(i), and on the
  soundness of none of them. Its objective values are not the run's
  \exttermref{../foundations/contract.pdf}{term-estimate}{estimate}: none is reported, and none refuses the run
  (Requirement~\ref{req:negative-estimate-refused} reads the estimate of the
  search on $\Decl$ alone).
  Each point $x_t=(\mathbf q_t,J_t)$ of either set is taken as the search
  or the polish consumed it, $\mathbf q_t$ included, and is a point of the
  Choi program of $\Decl^{\mathrm{ref}}$
  (Corollary~\ref{cor:fo-reference}\,(ii)).
\item \emph{The gate.} At each such point it forms the bracket of
  \eqref{eq:hlb} for $\Decl^{\mathrm{ref}}$: the objective of
  $\Decl^{\mathrm{ref}}$ at $x_t$, its gradient at $x_t$ with no
  perturbation, the acceptance-side minimisation of
  Requirements~\ref{req:acceptance-geometries}
  and~\ref{req:acceptance-iterate-exact}, and the feasible side of
  Requirement~\ref{req:verified-feasible-side}, from a fresh feasible-side
  oracle call on the linear form of that gradient over
  $\Decl^{\mathrm{ref}}$'s entries. It decides (H-POS) and (H-SUPP) on
  $\Decl^{\mathrm{ref}}$'s maps (Requirement~\ref{req:support-rank-check})
  and makes every check of Requirement~\ref{req:iterate-checks} against
  $\Decl^{\mathrm{ref}}$'s entries, on the point; a point at which
  Requirement~\ref{req:first-order-withdrawal} withdraws it on these terms
  contributes no bracket, with that requirement's reason.
\item \emph{The number.} Where some evaluated point contributes and the
  bracket is not withdrawn by (iv), the run's
  first-order number is $\hlb[\Decl^{\mathrm{ref}}]$, the largest bracket so
  formed, its clamp and key length are those of
  Requirement~\ref{req:first-order-clamp} and
  Corollary~\ref{cor:fo-keylength} at it, with the parameters of $\Decl$'s
  run, and its tier is \emph{first-order}. The bound on $\Decl$,
  $\hlb[\Decl]$, is recorded beside it with its brackets as a diagnostic,
  and no key length is formed from it. The difference
  $\hlb[\Decl^{\mathrm{ref}}]-\hlb[\Decl]$ is recorded with its sign, as a
  diagnostic, beside the difference of
  Requirement~\ref{req:bound-estimate-difference}, which is taken against
  the estimate of $\Decl$. A point of the second set that contributes
  counts as any evaluated point does: where no point of the first set
  contributes, it makes this item apply in place of the fallback of (iv),
  so that the run's first-order number moves from $\hlb[\Decl]$ to the
  companion.
\item \emph{The bracket withdrawn.} The bracket on the companion is
  withdrawn for one of two causes: no evaluated point contributes; or the
  certified evaluation finds the two enclosures of a block the companion
  carries inconsistent, the cause \iface{factor_enclosure_inconsistent}
  (Requirement~\ref{req:certified-evaluation}\,(iv)). An inconsistency shows
  that a module carrying the enclosures is defective, and the first-order
  number on the companion reads its entrywise enclosure, which may be the
  defective one; the second cause is decided after that number was formed,
  and the run then forms its first-order number again as this item says,
  before any key length is reported. Where the bracket is withdrawn, the
  run's first-order number is $\hlb[\Decl]$, formed as without a
  companion, with the reason \iface{reference_bracket_withdrawn}, its cause
  and the reasons of every evaluated point; no block is then evaluated off its snap.
  The record names in the group \emph{reference bracket withdrawn}, with the
  cause (\xref{foundations/req:snap-reference}\,(iii)(b)), every moved block
  on which a declaring instance whose link the companion replaces handed the
  reference and every moved class of each snapped emission of such a link
  whose normalizer the companion bypasses; it names as it does without a
  companion every other moved block or class that rests on
  \xref{foundations/ass:snapped-measurement}, so that a moved block on which a declaring
  instance whose link the companion does not replace handed a reference
  rests in the group \emph{reference bracket not evaluated}
  (\xref{foundations/req:snap-reference}\,(iii)(c)). A run
  that declares the strict option \iface{strict_snaps} refuses there instead
  and reports no key length (\xref{foundations/req:snap-reference}\,(iv)).
\item \emph{The blocks.} Where (iii) applies and the key length the run
  reports is not formed in the fallback of
  Requirement~\ref{req:certified-evaluation}\,(i), the record of the run names as
  \exttermref{../foundations/contract.pdf}{term-evaluated-off-its-snap}{evaluated off its snap} (\xref{foundations/def:evaluated-off-snap})
  every \exttermref{../foundations/contract.pdf}{term-moved}{moved} block of the snap of each declaring instance whose
  link the companion replaces, on which the companion carries the handed
  reference, and every data sector of each snapped emission of such a link
  that holds a moved class and whose normalizer the companion bypasses; a kept block or sector is named as kept and is not
  counted among them. It names every moved block on which the companion
  carries the snapped block as not so evaluated, so that it rests on
  \xref{foundations/ass:snapped-measurement}
  (\xref{foundations/req:snap-reference}\,(iii)), and a run that declares
  \iface{strict_snaps} refuses there (\xref{foundations/req:snap-reference}\,(iv)).
  Where (iii) applies but the key length the run reports is formed from a
  number certified on $\Decl$ in the fallback of
  Requirement~\ref{req:certified-evaluation}\,(i), no number formed on the
  companion reaches that key length, and the number it is formed from reads
  the snapped blocks: the moved blocks on which the companion carries the
  handed reference and the data sectors whose normalizer it bypasses are
  then not evaluated off their snap, condition (ii) of
  \xref{foundations/def:evaluated-off-snap} failing for them. The record
  names them as resting on \xref{foundations/ass:snapped-measurement} in
  the group \emph{reference bracket withdrawn} with the cause
  \iface{companion_not_certified}
  (\xref{foundations/req:snap-reference}\,(iii)(b)), beside the moved blocks
  on which the companion carries the snapped block, which rest as above. The
  run decides this after the certified evaluation and before any key length
  is reported, and no block is then named as evaluated off its
  snap. A run that declares \iface{strict_snaps} forms no such
  number (Requirement~\ref{req:certified-evaluation}\,(i)).
\item \emph{The record.} The record carries, per evaluated point, the set
  of (i) it belongs to, its bracket on $\Decl^{\mathrm{ref}}$ and the
  withdrawal reasons there, and for a point of the first set its bracket on
  $\Decl$; for the polish, whether it ran and, where it did not, why, the
  point it started from, the share of the wall budget set aside for it and
  the wall time it took, its iteration cap, the number of points it
  assessed and the cause it stopped on; the measurements of
  Requirement~\ref{req:companion-admission}\,(iii), the largest radius of the
  companion's enclosure among them, which bear on the arithmetic of
  Remark~\ref{rem:fo-arithmetic} and on nothing else; and, per block,
  whether the companion
  carries the handed reference or the snapped block.
\end{enumerate}
A run with no admitted companion forms its first-order number on $\Decl$
alone, as the other requirements of this submodule state; a moved block on
which a reference was handed, and a moved class of an emission beside which a
companion that bypasses its normalizer was formed, then rest in the group
\emph{reference bracket not evaluated}
(\xref{foundations/req:snap-reference}\,(iii)).
\end{requirement}
\claimmeta{\metalabel{Motivated by}\mref{Prop.}{prop:optol-nobound}, \mref{Prop.}{prop:reference-companion}, \mref{Req.}{req:companion-admission}, \mref{Thm.}{thm:firstorder}, \mref{Cor.}{cor:fo-keylength}, \mref{Cor.}{cor:fo-reference}, \mref{Rem.}{rem:fo-arithmetic}, \mref{Req.}{req:first-order-withdrawal}, \mref{Req.}{req:first-order-clamp}, \mref{Req.}{req:negative-estimate-refused}, \mref{Req.}{req:bound-estimate-difference}, \mref{Req.}{req:iterate-checks}, \mref{Req.}{req:support-rank-check}, \mref{Req.}{req:acceptance-geometries}, \mref{Req.}{req:acceptance-iterate-exact}, \mref{Req.}{req:verified-feasible-side}, \mref{Req.}{req:certified-evaluation}, \extref{../foundations/contract.pdf}{foundations/ass:snapped-measurement}, \extref{../foundations/contract.pdf}{foundations/def:evaluated-off-snap}, \extref{../foundations/contract.pdf}{foundations/req:snap-reference}
\metasep \metalabel{Used by}\mref{Req.}{req:certified-evaluation}
\metasep implemented (effective: accepted)
\metasep \metalabel{Code}\coderef{https://github.com/fraunhofersg/certiqs-key/blob/v0.1.0/src/certiqs_key/solvers/frank_wolfe/reference.py}{certiqs_key.solvers.frank_wolfe.reference:bracket_on_reference}}

\section{Requirements: checks and repairs of the iterates}\label{group-meat-core/frank-wolfe/req:iterate-checks}
\sectionintro{The engine checks every condition of the \termref{term-relaxed-feasible-set}{relaxed feasible set} on each consumed point, decides maximal support by one rank comparison, repairs a block only within a declared window and records every repair, an oracle return's repair serving the search alone. In the general format it decides the coupling condition per Alice sector.}

% claim meat-core/frank-wolfe/req:iterate-checks
\claimsummary{meat-core/frank-wolfe/req:iterate-checks}{At each iterate the engine decides block positivity of the consumed point at exactly 0, and the coupling and constraint-row residuals of that point at their own windows. The larger window is the \ensuremath{\varepsilon} at which a run reads the relaxed set; no declared tolerance is that number.}
\claimbookmark{meat-core/frank-wolfe/req:iterate-checks}{Requirement}{req:iterate-checks}{Feasibility checks at each iterate}
\claimleadin{%
Every statement about the \termref{term-relaxed-feasible-set}{relaxed feasible set} holds for every
$\varepsilon\ge0$ (Proposition~\ref{prop:relaxed-set-properties}), so no
claim depends on which $\varepsilon$ a run reads, and each candidate reading
is conservative; an implementation nevertheless needs one number. The set
relaxes positivity not at all, which is why positivity is decided at
exactly $0$ (Proposition~\ref{prop:relaxed-set-properties}\,(c),
\emph{Fourth}); a point that passes the three checks lies in
$\Feas^{J}(\varepsilon)$ at the larger window, by monotonicity from the
smaller one (\emph{ibid.}\,(a)) and because the coupling clause implies the
normalisation clause (\emph{ibid.}\,(c$'$)). The windows are numbers of
their own on a \exttermref{../foundations/contract.pdf}{term-computed-object}{computed object}
(\xref{foundations/req:tolerance-scope}\,(ii)), and this module takes no
declared tolerance (Requirement~\ref{req:exact-admission}). What the checks
buy, and why no check of a point the oracle returned would serve the bound,
is Remark~\ref{rem:fo-iterate-checks}.
}
\begin{requirement}[Feasibility checks at each iterate]\label{req:iterate-checks}
On the point it consumes at each iterate, after any repair, the engine
decides every condition defining $\Feas^{J}$ in \eqref{eq:choi-feas}, as the
relaxed feasible set relaxes it, all three and none of them optional:
\begin{enumerate}[label=(\roman*),leftmargin=2.4em]
\item per-block positive semidefiniteness, at exactly $0$ and at no window
  (\iface{iterate_positivity});
\item the coupling condition, its residual
  $\|\sum_s\tr_{\Qsys_s}J_s-\idon{A_0}\|_{\max}$, or, in
  \exttermref{../transmitter/contract.pdf}{term-untrusted-source-mode}{untrusted-source mode}, where the coupling condition is the
  normalisation $\sum_s\tr\omega_s=1$
  (Definition~\ref{def:untrusted-choi}\,(b)), its residual
  $|\sum_s\tr\omega_s-1|$, at the declared window \iface{iterate_coupling};
\item the declared constraint rows $\bowtie_k$, each residual at the declared
  window \iface{constraint_rows}.
\end{enumerate}
No fourth check is required, and none is required of the point the oracle
returned at that iterate. The
$\varepsilon$ at which a run reads the relaxed feasible set is the larger of
the two declared windows of (ii) and~(iii). No \exttermref{../foundations/contract.pdf}{term-declared-tolerance}{declared tolerance} (\xref{foundations/def:tolerance}) is
that number: the windows govern conditions checked on an iterate and never on
a declaration.
\end{requirement}
\claimmeta{\metalabel{Motivated by}\mref{Thm.}{thm:choi}, \mref{Def.}{def:untrusted-choi}, \mref{Req.}{req:exact-admission}, \mref{Def.}{def:relaxed-feasible-set}, \mref{Prop.}{prop:relaxed-set-properties}, \mref{Rem.}{rem:fo-iterate-checks}, \extref{../foundations/contract.pdf}{foundations/def:tolerance}, \extref{../foundations/contract.pdf}{foundations/req:tolerance-scope}
\metasep \metalabel{Used by}\mref{Req.}{req:returned-point-measured}, \mref{Req.}{req:bracket-on-reference}
\metasep implemented (effective: accepted)
\metasep \metalabel{Code}\coderef{https://github.com/fraunhofersg/certiqs-key/blob/v0.1.0/src/certiqs_key/solvers/frank_wolfe/joint_driver.py}{certiqs_key.solvers.frank_wolfe.joint_driver:_check_point}}

% claim meat-core/frank-wolfe/req:support-rank-check
\claimsummary{meat-core/frank-wolfe/req:support-rank-check}{The engine decides (H-SUPP) as a rank equality, both sides computed and read with one declared cutoff relative to a spectral radius. The reference rank is computed once per acceptance value and sector from the declaration and the support frame, never from an iterate.}
\claimbookmark{meat-core/frank-wolfe/req:support-rank-check}{Requirement}{req:support-rank-check}{Rank test for maximal support}
\claimleadin{%
Hypothesis (H-SUPP) of Theorem~\ref{thm:firstorder} is equivalent to the rank
equality \eqref{eq:hsupp-check}, and no positivity test on the ambient space
of the branch operator can decide it (Remark~\ref{rem:fo-one-ruler}). A
direction along a singular value $\sigma$ of the stacked Kraus matrix carries
weight of order $\sigma^2$ in the branch operator, which is why the right
side is read at the square root of the cutoff: both sides then count with
one ruler. The comparison discharges nothing;
Remark~\ref{rem:fo-check-exposures} names the residual exposures it leaves.
}
\begin{requirement}[Rank test for maximal support]\label{req:support-rank-check}
The engine decides (H-SUPP) as the rank equality \eqref{eq:hsupp-check},
both sides computed and both read with the one declared threshold
\iface{support_cutoff}:
\begin{enumerate}[label=(\roman*),leftmargin=2.4em]
\item the left side is the count of eigenvalues of $X_{\hat c,s}$ above
  \iface{support_cutoff}, read relative to that block's own spectral radius;
\item the right side is $\operatorname{rank}\mathbf K_{\hat c,s}$ of
  \eqref{eq:Kstack}, read at the square root of \iface{support_cutoff}
  relative to that matrix's largest singular value, a window derived from
  the same threshold and not declared beside it; equivalently the count of
  eigenvalues of $\mathbf K_{\hat c,s}\mathbf K_{\hat c,s}^{\dagger}$ above
  \iface{support_cutoff} relative to its spectral radius. It is computed once
  per $(\hat c,s)$ from the declaration and $W_{\tau}$, and never from an
  iterate.
\end{enumerate}
\end{requirement}
\claimmeta{\metalabel{Motivated by}\mref{Def.}{def:branch-support}, \mref{Lem.}{lem:branch-support}, \mref{Thm.}{thm:firstorder}, \mref{Rem.}{rem:fo-one-ruler}, \mref{Rem.}{rem:fo-check-exposures}
\metasep \metalabel{Used by}\mref{Req.}{req:bracket-on-reference}
\metasep implemented (effective: accepted)
\metasep \metalabel{Code}\coderef{https://github.com/fraunhofersg/certiqs-key/blob/v0.1.0/src/certiqs_key/meat_core/objective/gradient.py}{certiqs_key.meat_core.objective.gradient:attains_the_reference_rank}}

% claim meat-core/frank-wolfe/req:oracle-repair
\claimsummary{meat-core/frank-wolfe/req:oracle-repair}{An oracle return with a slightly negative spectrum may be rejected, or projected onto the positive cone within a declared window of its own. Every repair is recorded with its block, iterate and depth, apart from iterate repairs, and the repaired return serves the search alone.}
\claimbookmark{meat-core/frank-wolfe/req:oracle-repair}{Requirement}{req:oracle-repair}{Repair of an oracle return}
\claimleadin{%
No bracket reads a point the backend returned
(Requirement~\ref{req:verified-feasible-side}), so a repair of a return
changes only the search's next step and a recorded diagnostic. A window
bounds the displacement of the block but gives no bound on a pairing with
it (Remark~\ref{rem:fo-oracle-repair}), and the object repaired is a
\exttermref{../foundations/contract.pdf}{term-computed-object}{computed object}, whose window is a number of its own
(\xref{foundations/req:tolerance-scope}\,(ii)).
}
\begin{requirement}[Repair of an oracle return]\label{req:oracle-repair}
The engine may reject outright an oracle return a block of which has a
negative spectrum, or project that block onto the positive semidefinite cone
within the declared window \iface{oracle_repair_window}, a number of its own,
declared beside $\tnum$ and never as an instance of it. It records every
such repair, with the block, the iterate and the measured depth, in
\iface{optimizer_certificate}, kept apart there from the iterate repairs of
Requirement~\ref{req:iterate-repair}, and uses the repaired return only as
a step of the search and in the diagnostic pairing of
Requirement~\ref{req:verified-feasible-side}\,(vi), never in a bracket and
never as a bound on anything.
\end{requirement}
\claimmeta{\metalabel{Motivated by}\mref{Rem.}{rem:fo-oracle-repair}, \mref{Req.}{req:iterate-repair}, \mref{Req.}{req:verified-feasible-side}, \extref{../foundations/contract.pdf}{foundations/req:tolerance-scope}
\metasep implemented (effective: accepted)
\metasep \metalabel{Code}\coderef{https://github.com/fraunhofersg/certiqs-key/blob/v0.1.0/src/certiqs_key/solvers/oracles/qics/adapter.py}{certiqs_key.solvers.oracles.qics.adapter:QicsSectorChoiOracle}}

% claim meat-core/frank-wolfe/req:iterate-repair
\claimsummary{meat-core/frank-wolfe/req:iterate-repair}{A block of a formed iterate that fails the positivity check is rejected, or projected within a declared window of its own. The repaired point is the iterate, its residuals are measured again, and every repair is recorded apart from oracle repairs.}
\claimbookmark{meat-core/frank-wolfe/req:iterate-repair}{Requirement}{req:iterate-repair}{Repair of a formed iterate}
\claimleadin{%
The repair belongs to the formation of the iterate, upstream of the point
Theorem~\ref{thm:firstorder} is applied at, so it is not the perturbation
that (H-NOPERT) excludes; it establishes a hypothesis of that theorem on an
object the engine chose and says nothing about the backend
(Remark~\ref{rem:fo-iterate-repair}). Re-measuring the residuals is the one
place this repair can cost a run.
}
\begin{requirement}[Repair of a formed iterate]\label{req:iterate-repair}
When a block of the point consumed at iterate $t$ fails the positivity check
of Requirement~\ref{req:iterate-checks}, the engine either rejects it, in
which case $t$ contributes no bracket and is not a member of $\Itset$, or
projects it onto the positive semidefinite cone, provided its measured depth
is within the declared window \iface{iterate_repair_window}, a number of its
own, declared beside $\tnum$ and never as an instance of it. No other repair
is made on an iterate. A projection is governed by three rules:
\begin{enumerate}[label=(\roman*),leftmargin=2.4em]
\item the repaired point is $x_t$, the point at which the objective, the
  gradient and the pairings with $x_t$ of \eqref{eq:hlb} are evaluated;
\item the coupling and $\bowtie_k$ residuals of
  Requirement~\ref{req:iterate-checks} are measured again on the repaired
  point and never carried over from the point before it, and a residual that
  then fails its window refuses the iterate;
\item every repair is recorded, with the block, the iterate and the measured
  depth, in \iface{optimizer_certificate}, beside the oracle-return repairs
  of Requirement~\ref{req:oracle-repair} and distinguished from them.
\end{enumerate}
\end{requirement}
\claimmeta{\metalabel{Motivated by}\mref{Thm.}{thm:firstorder}, \mref{Rem.}{rem:fo-iterate-repair}, \mref{Req.}{req:iterate-checks}, \mref{Req.}{req:oracle-repair}
\metasep implemented (effective: accepted)
\metasep \metalabel{Code}\coderef{https://github.com/fraunhofersg/certiqs-key/blob/v0.1.0/src/certiqs_key/solvers/frank_wolfe/joint_driver.py}{certiqs_key.solvers.frank_wolfe.joint_driver:_check_point}}

% claim meat-core/frank-wolfe/req:per-sector-coupling
\claimsummary{meat-core/frank-wolfe/req:per-sector-coupling}{In the general format a run decides the coupling condition of the relaxed feasible set per Alice sector, one residual per sector at the one declared coupling window, and records the frame and the row form per sector.}
\claimbookmark{meat-core/frank-wolfe/req:per-sector-coupling}{Requirement}{req:per-sector-coupling}{Coupling residual and record per Alice sector}
\claimleadin{%
In the general format the Choi program carries one coupling constraint per
Alice sector (Theorem~\ref{thm:choi}), and the relaxed feasible set is read
with the coupling constraint per Alice sector
(Remark~\ref{rem:fo-general-format}\,(3)).
}
\begin{requirement}[Coupling residual and record per Alice sector]\label{req:per-sector-coupling}
On a declaration in the general format the engine decides the coupling
condition of the relaxed feasible set of
Definition~\ref{def:relaxed-feasible-set} per Alice sector, one coupling
residual per Alice sector, each at the declared window
\iface{iterate_coupling}; and the record names the frame and the row form
per Alice sector (Requirements~\ref{req:frame-recorded}
and~\ref{req:row-form-record}).
\end{requirement}
\claimmeta{\metalabel{Motivated by}\mref{Thm.}{thm:choi}, \mref{Req.}{req:frame-recorded}, \mref{Req.}{req:row-form-record}, \mref{Def.}{def:relaxed-feasible-set}, \mref{Prop.}{prop:relaxed-set-properties}, \mref{Rem.}{rem:fo-general-format}
\metasep implemented (effective: accepted)
\metasep \metalabel{Code}\coderef{https://github.com/fraunhofersg/certiqs-key/blob/v0.1.0/src/certiqs_key/solvers/frank_wolfe/joint_driver.py}{certiqs_key.solvers.frank_wolfe.joint_driver:coupling_residuals}}

\section{Requirements: the acceptance side}\label{group-meat-core/frank-wolfe/req:acceptance-geometries}
\sectionintro{The engine solves the acceptance-side minimisation exactly for a singleton or a polytope acceptance set and admits no other acceptance set. It forms the acceptance iterate exactly on the exact island and records, per iterate, whether it did; a run that replaces the acceptance-side pairing by a certified bound records which of the two the bracket used.}

% claim meat-core/frank-wolfe/req:acceptance-geometries
\claimsummary{meat-core/frank-wolfe/req:acceptance-geometries}{The engine solves the acceptance-side minimisation exactly for a singleton or a polytope acceptance set and admits no other geometry. A declaration with any other acceptance set is refused before any solve.}
\claimbookmark{meat-core/frank-wolfe/req:acceptance-geometries}{Requirement}{req:acceptance-geometries}{The admitted acceptance geometries}
\claimleadin{%
Definition~\ref{def:admissible-acceptance-set} admits any closed convex
subset of the probability simplex, while Remark~\ref{rem:acc-side} solves the
acceptance-side minimisation exactly for two geometries only. For any other
set the bracket would need a certified $\bacc$ of \eqref{eq:bacc-def} from
elsewhere, which no requirement decides and no assumption carries, so the
engine admits no other set. The refusal is this requirement's own: the
general admission rule, Requirement~\ref{req:acceptance-set-admission}, would
admit a third tag whose every datum denotes a closed convex set.
}
\begin{requirement}[The admitted acceptance geometries]\label{req:acceptance-geometries}
The engine solves the acceptance-side minimisation exactly for the two
geometries of Remark~\ref{rem:acc-side}, a singleton~(a) and a polytope~(b),
and admits an acceptance set in no other geometry: a declaration whose
$\acSet$ fits neither is refused at admission, before any solve, and no
\exttermref{../foundations/contract.pdf}{term-first-order}{first-order} bound is formed for it.
\end{requirement}
\claimmeta{\metalabel{Motivated by}\mref{Def.}{def:admissible-acceptance-set}, \mref{Req.}{req:acceptance-set-admission}, \mref{Rem.}{rem:acc-side}
\metasep \metalabel{Used by}\mref{Req.}{req:bracket-on-reference}, \mref{Req.}{req:certified-evaluation}
\metasep accepted
\metasep \metalabel{Code}\coderef{https://github.com/fraunhofersg/certiqs-key/blob/v0.1.0/src/certiqs_key/solvers/oracles/acceptance.py}{certiqs_key.solvers.oracles.acceptance:acceptance_oracle}}

% claim meat-core/frank-wolfe/req:acceptance-iterate-exact
\claimsummary{meat-core/frank-wolfe/req:acceptance-iterate-exact}{The engine forms the acceptance-side iterate on the exact island, from an exact rational start in the acceptance set by exact rational convex combinations of points established to lie in it, and records per iterate whether it did. An iterate where it did not adds nothing to the first-order bound.}
\claimbookmark{meat-core/frank-wolfe/req:acceptance-iterate-exact}{Requirement}{req:acceptance-iterate-exact}{Exact formation of the acceptance-side iterate}
\claimleadin{%
Under conditions ($\alpha$) and ($\beta$) of Remark~\ref{rem:acc-side} the
acceptance-side iterate lies in $\acSet$ at every step by induction, so the
acceptance factor of the membership hypothesis of
Theorem~\ref{thm:firstorder}\,(ii) and~(iiib) is true by construction and
needs no check; a binary64 convex combination of exact points of a polytope
can leave the set, and three consumers of that theorem then break (the same
remark). Every computation of the \exttermref{../foundations/contract.pdf}{term-exact-island}{exact island} is exact rational
arithmetic (\xref{foundations/req:exact-island}).
}
\begin{requirement}[Exact formation of the acceptance-side iterate]\label{req:acceptance-iterate-exact}
The engine forms the acceptance-side iterate on the \exttermref{../foundations/contract.pdf}{term-exact-island}{exact island},
under conditions ($\alpha$) and ($\beta$) of Remark~\ref{rem:acc-side}: it
initialises it at an exact rational point of $\acSet$, and updates it only by
convex combinations, with exact rational weights in $[0,1]$, of points
already established to lie in $\acSet$. An iterate whose acceptance-side
component was not so formed has not met the membership hypothesis and
contributes no bracket to $\hlb$, with the reason of
Requirement~\ref{req:first-order-withdrawal}\,(iv). The run's record states in
\iface{optimizer_certificate}, of the point consumed at each iterate, whether
its acceptance-side component was so formed.
\end{requirement}
\claimmeta{\metalabel{Motivated by}\mref{Thm.}{thm:firstorder}, \mref{Rem.}{rem:acc-side}, \mref{Req.}{req:first-order-withdrawal}, \extref{../foundations/contract.pdf}{foundations/req:exact-island}
\metasep \metalabel{Used by}\mref{Req.}{req:bracket-on-reference}, \mref{Req.}{req:certified-evaluation}
\metasep implemented (effective: accepted)
\metasep \metalabel{Code}\coderef{https://github.com/fraunhofersg/certiqs-key/blob/v0.1.0/src/certiqs_key/solvers/frank_wolfe/acceptance_iterate.py}{certiqs_key.solvers.frank_wolfe.acceptance_iterate:AcceptanceIterate}}

% claim meat-core/frank-wolfe/req:bacc-recorded
\claimsummary{meat-core/frank-wolfe/req:bacc-recorded}{A run that replaces the acceptance-side pairing in the bracket of the first-order bound by a certified lower bound on the minimum records which of the two it used. The substitution is exact and leaves the first-order bound valid.}
\claimbookmark{meat-core/frank-wolfe/req:bacc-recorded}{Requirement}{req:bacc-recorded}{Record of the acceptance-side quantity in the bracket}
\claimleadin{%
A certified $\bacc$ of \eqref{eq:bacc-def} may replace $\langle\lmoq_t,v_t\rangle$ in
\eqref{eq:hlb}, and the first-order bound stays valid
(Theorem~\ref{thm:firstorder}\,(iv); Remark~\ref{rem:acc-side}).
}
\begin{technicalclaim}
\begin{requirement}[Record of the acceptance-side quantity in the bracket]\label{req:bacc-recorded}
A run that replaces $\langle\lmoq_t,v_t\rangle$ in the bracket of
\eqref{eq:hlb} by a certified $\bacc$ of \eqref{eq:bacc-def} records which of
the two it used.
\end{requirement}
\end{technicalclaim}
\claimmeta{\metalabel{Motivated by}\mref{Thm.}{thm:firstorder}, \mref{Rem.}{rem:acc-side}
\metasep accepted}

\section{Requirements: the feasible side}\label{group-meat-core/frank-wolfe/req:verified-feasible-side}
\sectionintro{The feasible side of every bracket is the verified lower bound of \hyperlink{claim-meat-core/frank-wolfe/lem:feasible-side-dual}{Lemma~\ref*{lem:feasible-side-dual}} over the exact feasible set, formed in exact rational arithmetic from the gradient as computed, the declared rows and the backend's dual, less a deduction for the rounding of the frame, with zero multipliers where the backend returned none and zero row multipliers where the frame's error is not bounded. The oracle's primal pairing, the dual gap and the deduction are recorded beside it. Beside it stands the superseded requirement that paired the gradient with the oracle's return, which no run follows.}
\sectionimports{\extref{../foundations/contract.pdf}{foundations/def:exact-island}, \extref{../foundations/contract.pdf}{foundations/lem:least-eigenvalue-bound}, \extref{../foundations/contract.pdf}{foundations/req:exact-island}}

% claim meat-core/frank-wolfe/req:verified-feasible-side
\claimsummary{meat-core/frank-wolfe/req:verified-feasible-side}{The feasible side of every bracket is the lower bound of \hyperlink{claim-meat-core/frank-wolfe/lem:feasible-side-dual}{Lemma~\ref*{lem:feasible-side-dual}} over the exact feasible set, formed exactly from the computed gradient and the backend's dual, less a deduction for the computed frame, or with zero row multipliers. The oracle's return, status and objective enter no number.}
\claimbookmark{meat-core/frank-wolfe/req:verified-feasible-side}{Requirement}{req:verified-feasible-side}{The verified feasible side of every bracket}
\claimleadin{%
Part~(iv) of Theorem~\ref{thm:firstorder} accepts any real number below the
feasible-side minimum over the exact $\Feas^{J}$ in place of the pairing
with the oracle's return. Lemma~\ref{lem:feasible-side-dual} gives one from
any multipliers, and its part~(iv) charges the rounding of the frame a run
builds its rows with. Formed as below, $b^{\mathrm F}_t$ is at most the
minimum, over the exact $\Feas^{J}$, of the linear form at the gradient as
the run computed it, whatever the backend's status, its dual or the
rounding of the frame: it rests on no property of the backend. A dual
objective the backend reports is a different number, which an
approximately dual-feasible return can put above the minimum, and it is
never used. What this requirement does not reach is the gradient itself,
computed in binary64 from maps and a frame computed in binary64; that is
named in Remark~\ref{rem:fo-arithmetic} and bounded by no claim of this
module. The same formation serves the bracket on a reference companion
(Corollary~\ref{cor:fo-reference},
Requirement~\ref{req:bracket-on-reference}).
}
\begin{requirement}[The verified feasible side of every bracket]\label{req:verified-feasible-side}
At every iterate $t$ at which the engine forms a bracket of \eqref{eq:hlb},
on the declaration it searched and on an admitted reference companion alike,
it uses in place of the pairing $\langle\lmoJ_t,S_t\rangle$ the number
$b^{\mathrm F}_t$ of Theorem~\ref{thm:firstorder}\,(iv), formed as follows,
and never the pairing with a point the backend returned, a dual objective
the backend reported, or the backend's status.
\begin{enumerate}[label=(\roman*),leftmargin=2.4em]
\item \emph{The groups and the data.} $b^{\mathrm F}_t$ is a sum over the
  groups of the program, one per Alice sector in \exttermref{../transmitter/contract.pdf}{term-trusted-source-mode}{trusted-source mode}
  and one in \exttermref{../transmitter/contract.pdf}{term-untrusted-source-mode}{untrusted-source mode}, read as
  Lemma~\ref{lem:feasible-side-dual}\,(ii) reads them. Its $G$ is the exact
  Hermitian part of the gradient blocks as the engine computed them at
  $x_t$, the same gradient whose acceptance-side part the bracket uses; its
  coupling is the exact coupling condition of the group; its rows are the
  maps of Lemma~\ref{lem:feasible-side-dual}\,(ii), in the conjugated or
  frame-free form the engine imposes, built from the declared data taken
  exactly: $\Theta^{(\mathsf t)}_k$ as the declared block or the exact
  product $c_k\tau^{(\mathsf t)}_A$, the lifted operators, and
  $D^{(s_f)}_k=\hat O^{(s_f)}_k-c_k\idon{}$ formed exactly, all Hermitian by
  exact admissibility; and in trusted-source mode the frame in those maps is
  the binary64 $\widetilde W_{\mathsf t}$ the engine constructed for
  $W_{\mathsf t}$, read as the dyadic matrix it is. Every binary64 number is
  read as the dyadic rational it is.
\item \emph{The multipliers.} $Y$ and every $\Lambda_r$ are read off the
  dual the backend returned for that same solve, as the operators that the
  multipliers of the coupling condition and of row $r$ represent in the
  encoding the oracle used, each replaced by its exact Hermitian part. On a
  group with a one-dimensional coupling and no rows, and on any group whose
  solve returned no finite dual, whatever its status, $Y=0$ and every
  $\Lambda_r=0$; on an Alice sector for which (iv) forms no
  $\varepsilon_{\mathsf t}$, every $\Lambda_r$ of its rows is $0$ and $Y$ is
  kept. The multiplier of an inequality row is replaced by
  $\Lambda_r+\mu_r\id$, with $\mu_r$ a dyadic rational at least
  $\max\{0,-\ell_r\}$ and $\ell_r$ a lower bound on $\lmin(\Lambda_r)$ given
  by \xref{foundations/lem:least-eigenvalue-bound}, so that
  $\Lambda_r\succeq0$ holds exactly.
\item \emph{The eigenvalue.} The group's $\lambda$ is the least, over the
  blocks of the group, of the lower bounds on $\lmin(R_p)$ that
  \xref{foundations/lem:least-eigenvalue-bound} gives at a matrix $V$ of the
  engine's choosing, with $R_p$ formed from the rows of (i). The bound holds
  at every $V$ that lemma admits, the identity ($\theta=0$) always among
  them, and the choice bears only on how close $\lambda$ lies to
  $\min_p\lmin(R_p)$.
\item \emph{The frame.} For each Alice sector the engine forms a dyadic
  $\varepsilon_{\mathsf t}\ge\opn{\widetilde W_{\mathsf t}-W_{\mathsf t}}$,
  established in exact rational arithmetic, where it can: $0$ where
  $\widetilde W_{\mathsf t}=W_{\mathsf t}$ is decided exactly; and where
  $\tau^{(\mathsf t)}_A$ is diagonal, the frame is the \termref{term-full-frame}{full frame} or
  the \termref{term-diagonal-frame}{diagonal frame} (Definition~\ref{def:support-frame}), and
  $\widetilde W_{\mathsf t}$ vanishes exactly wherever $W_{\mathsf t}$
  does, so that both have at most one non-zero entry in each row and each
  column and $\opn{\widetilde W_{\mathsf t}-W_{\mathsf t}}$ is the largest
  modulus of their entrywise difference: a dyadic at least
  $|\widetilde w-\sqrt{\tau_{ii}}|$ at every entry where $W_{\mathsf t}$
  carries $\sqrt{\tau_{ii}}$, $\widetilde w$ being the entry of
  $\widetilde W_{\mathsf t}$ there (for a real $\widetilde w>0$,
  $|\widetilde w^{2}-\tau_{ii}|/\widetilde w$ is one, since
  $|\widetilde w-\sqrt{\tau_{ii}}|=|\widetilde w^{2}-\tau_{ii}|/
  (\widetilde w+\sqrt{\tau_{ii}})$). Where every row of the sector is
  imposed in frame-free form, no deduction is formed: those rows read no
  frame, every $\kappa_r$ is $0$, and the sector's bound is that of
  Lemma~\ref{lem:feasible-side-dual}\,(i). In every other case no
  $\varepsilon_{\mathsf t}$ is formed. Where it is formed, the sector's bound
  is that of Lemma~\ref{lem:feasible-side-dual}\,(iv),
  $b(G_{\mathsf t};Y,\Lambda,\widetilde\lambda)-\sum_r\kappa_r$, with each
  $\kappa_r$ formed from dyadic upper bounds on
  $\opn{\widetilde W_{\mathsf t}}$, $\opn{\Lambda_r}$ and $\opn{O_{r,s}}$
  (for instance on their Frobenius norms, by an integer square root rounded
  up). Where it is not, the row multipliers are zero by (ii) and the
  sector's bound is that of Lemma~\ref{lem:feasible-side-dual}\,(i), which
  then reads no row. In untrusted-source mode there is no frame and no
  deduction.
\item \emph{The arithmetic.} The residuals, the shifted multipliers, every
  eigenvalue and norm bound, every deduction and $b^{\mathrm F}_t$ are
  computed in exact rational arithmetic on those dyadic data, as a
  certification of the \exttermref{../foundations/contract.pdf}{term-exact-island}{exact island}
  (\xref{foundations/def:exact-island}\,(iii)), the backend's dual entering
  it only as a candidate (\xref{foundations/req:exact-island}).
  $b^{\mathrm F}_t$ enters the binary64 bracket rounded toward $-\infty$, the
  direction that weakens a lower bound, as
  \xref{foundations/req:rounding-direction} rounds an emitted one.
\item \emph{The record.} Per bracket the engine records $b^{\mathrm F}_t$
  (\iface{feasible_side_bound}); the sum of the deductions $\kappa_r$
  (\iface{frame_deduction}) and the sum of the bounds on
  $\opn{\Lambda_r}$ it used (\iface{row_multiplier_norm}); where the backend
  returned a point, the pairing $\langle\lmoJ_t,S_t\rangle$ with it
  (\iface{oracle_pairing}) and the difference
  $\langle\lmoJ_t,S_t\rangle-b^{\mathrm F}_t$ with its sign
  (\iface{dual_gap}); and, per group, whether its multipliers were the
  backend's, zero, or the backend's with the row multipliers set to zero
  (\iface{feasible_side_multipliers}). None of these but
  $b^{\mathrm F}_t$ enters any number.
\end{enumerate}
\end{requirement}
\claimmeta{\metalabel{Motivated by}\mref{Def.}{def:support-frame}, \mref{Thm.}{thm:firstorder}, \mref{Cor.}{cor:fo-reference}, \mref{Lem.}{lem:feasible-side-dual}, \mref{Rem.}{rem:fo-arithmetic}, \mref{Req.}{req:bracket-on-reference}, \extref{../foundations/contract.pdf}{foundations/def:exact-island}, \extref{../foundations/contract.pdf}{foundations/lem:least-eigenvalue-bound}, \extref{../foundations/contract.pdf}{foundations/req:exact-island}, \extref{../foundations/contract.pdf}{foundations/req:rounding-direction}
\metasep \metalabel{Used by}\mref{Req.}{req:bracket-on-reference}, \mref{Req.}{req:certified-evaluation}
\metasep implemented (effective: accepted)
\metasep \metalabel{Code}\coderef{https://github.com/fraunhofersg/certiqs-key/blob/v0.1.0/src/certiqs_key/solvers/frank_wolfe/verified.py}{certiqs_key.solvers.frank_wolfe.verified:verified_feasible_side}}

% claim meat-core/frank-wolfe/req:primal-pairing
\claimsummary{meat-core/frank-wolfe/req:primal-pairing}{The bracket of the first-order bound pairs the gradient with the oracle's primal return, never with a dual objective the backend reports. No bracket uses either: \hyperlink{claim-meat-core/frank-wolfe/req:verified-feasible-side}{Requirement~\ref*{req:verified-feasible-side}} puts a verified lower bound in place of the pairing.}
\claimbookmark{meat-core/frank-wolfe/req:primal-pairing}{Requirement}{req:primal-pairing}{The primal pairing in the first-order bracket}
\claimleadin{%
A backend that solves the linear subproblem by an interior-point method
reports a \emph{dual} objective, and it is tempting to read that number as a
lower bound on $\min_{J'\in\Feas^{J}}\langle\lmoJ_t,J'\rangle$ and to use it
in place of $\langle\lmoJ_t,S_t\rangle$ in \eqref{eq:hlb}. Weak duality
bounds the primal minimum from below by the dual objective at a
\emph{feasible} dual point; an interior-point return is only approximately
dual-feasible, and for an infeasibility residual $R$ the corrected bound
carries a term of size at most
$\|R\|\cdot\sup_{J'\in\Feas^{J}}\|J'\|_1$, of \emph{unknown sign} without
$R$ itself. A scalar convergence indicator does not supply $R$, and this
module requires no oracle to report it. Consequently a reported dual
objective can lie above the primal value at the returned point, and
substituting it would put \eqref{eq:hlb} on the unsafe side of the very
quantity it is supposed to bound.
}
\begin{requirement}[The primal pairing in the first-order bracket]\label{req:primal-pairing}
In the bracket of \eqref{eq:hlb} the engine uses the primal pairing
$\langle\lmoJ_t,S_t\rangle$ at the returned point $S_t$ and nothing else; in
particular it never uses a dual objective reported by the backend in its
place.
\end{requirement}
\claimmeta{\metalabel{Motivated by}\mref{Thm.}{thm:firstorder}
\metasep superseded}

\section{Requirements: the certified evaluation}\label{group-meat-core/frank-wolfe/req:certified-evaluation}
\sectionintro{By default the engine forms the bracket again at the points of the run's first-order number from the declaration's exact data, the maps in factor form and every step exact or rigorously enclosed, and reports the best certified bracket at the certified tier (\hyperlink{claim-meat-core/frank-wolfe/req:certified-evaluation}{Requirement~\ref*{req:certified-evaluation}}). An element of a companion's measured family whose rank neither its entrywise enclosure nor an enclosed factor fixes, or a block not shown positive definite, leaves a point uncertified. Where a companion certifies no point, the evaluation runs once more on the searched declaration, and a number certified there is reported at the certified tier, the blocks the companion carried resting on the snap premise; a run with no certified point keeps its first-order number with the reasons, or refuses under the strict option.}
\sectionimports{\extref{../foundations/contract.pdf}{foundations/cor:least-eigenvalue-over-enclosure}, \extref{../foundations/contract.pdf}{foundations/cor:rank-over-factor-enclosure}, \extref{../foundations/contract.pdf}{foundations/def:factor-enclosure}, \extref{../foundations/contract.pdf}{foundations/def:tiers}, \extref{../foundations/contract.pdf}{foundations/lem:least-eigenvalue-bound}, \extref{../foundations/contract.pdf}{foundations/lem:matrix-function-perturbation}, \extref{../foundations/contract.pdf}{foundations/lem:recentred-enclosure}, \extref{../foundations/contract.pdf}{foundations/req:exact-island}, \extref{../foundations/contract.pdf}{foundations/req:rounding-direction}, \extref{../foundations/contract.pdf}{foundations/req:snap-reference}, \extref{../foundations/contract.pdf}{foundations/req:tiers}}

% claim meat-core/frank-wolfe/req:certified-evaluation
\claimsummary{meat-core/frank-wolfe/req:certified-evaluation}{The engine forms the bracket again at the final points, every step exact or rigorously enclosed, and reports the best at the certified tier; if a companion certifies none, on the searched declaration under \hyperref{../foundations/contract.pdf}{}{claim-foundations/ass:snapped-measurement}{A-SNAP} unless strict\_snaps is set. Else the first-order number stays, or a strict run refuses.}
\claimbookmark{meat-core/frank-wolfe/req:certified-evaluation}{Requirement}{req:certified-evaluation}{The certified evaluation of the first-order bound}
\claimleadin{%
The first-order number is formed in binary64, from maps, a frame and a
gradient that no claim bounds (Remark~\ref{rem:fo-arithmetic}). Lemma~\ref{lem:certified-bracket}
needs, at a point, only enclosures of the gradient, and
Corollary~\ref{cor:certified-keylength} turns the best certified bracket into
a key length. The gradient's maps can be built from factors whose rank is
fixed (Lemma~\ref{lem:kraus-freedom}), every matrix function then acts on a
positive definite matrix whose spectrum an enclosure bounds
(\xref{foundations/lem:matrix-function-perturbation}), and an enclosure of an
element of the measured family fixes its rank where its non-zero part is
positive definite throughout
(\xref{foundations/cor:least-eigenvalue-over-enclosure}) or, for an element
carried with an enclosed factor, where the factor's columns are independent
throughout (\xref{foundations/cor:rank-over-factor-enclosure}). A computation
in ball arithmetic is \exttermref{../foundations/contract.pdf}{term-rigorously-enclosed}{rigorously enclosed}; it lies outside the
\exttermref{../foundations/contract.pdf}{term-exact-island}{exact island} and is no error budget
(\xref{foundations/req:exact-island}). Of the search the evaluation reads
the points, and the multipliers of its solves as candidates of which nothing
is assumed; any point at which the lemma's hypotheses are established serves,
and a point that fails anywhere is dropped, never repaired. Certifying only
a few of the points the first-order number reads keeps the evaluation a
small fraction of the run's cost; a maximum over fewer certified brackets is
still a lower bound. Where a
companion certifies no point, the evaluation is made once more on the
declaration the run searched: a number certified there bounds that
declaration's value, so the blocks the companion carried rest on the snap
premise (\xref{foundations/ass:snapped-measurement}), and the number is
never combined with the companion's. Remark~\ref{rem:certified-scope} says
what the resulting number still rests on.
}
\begin{requirement}[The certified evaluation of the first-order bound]\label{req:certified-evaluation}
On every run that forms a first-order number, unless it sets the option
\iface{certified_evaluation} to \iface{off}, its default being \iface{on},
the engine forms a certified number as follows, after the first-order number
is formed and before a key length is reported.
\begin{enumerate}[label=(\roman*),leftmargin=2.4em]
\item \emph{The declaration and the points.} It works on the declaration the
  run's first-order number is formed on, at the points that number is formed
  from: on an admitted reference companion where
  Requirement~\ref{req:bracket-on-reference}\,(iii) applies, at the at most
  five of the points part~(i) of that requirement evaluates, of the search
  and of the polish together, that contribute there with the largest
  brackets on the companion, a point that is in both sets counted once and
  ties broken by the sequence in which the search and then the polish
  consumed them, each with the solve part~(ii) of it
  makes there; otherwise on the declaration the run searched, at the at most
  four members of the index set of its bound
  (Requirement~\ref{req:first-order-withdrawal}) with the largest brackets,
  ties broken by the order of the run, and the last point the search
  assessed, each with the search's solve there. \emph{The fallback}: where it
  has worked on a companion by (ii)--(v) and certified no point there, it
  works once more, by (ii)--(v), on the declaration the run searched, at the
  points and with the solves of the second case, unless the run declares the
  strict option \iface{strict_snaps} of
  \xref{foundations/req:snap-reference}\,(iv); the companion's results are
  kept for the record (ix). The fallback runs nowhere else: not where a point
  of the companion is certified, not where no companion's bracket counts, not
  where the companion's bracket is withdrawn by (iv), and not where the
  arithmetic the evaluation needs is not available. Where the companion's
  bracket is withdrawn by (iv), the run's first-order number is formed on the
  declaration the run searched (Requirement~\ref{req:bracket-on-reference}\,(iv)),
  and the evaluation works once more there, by the second case above, as on
  any run whose first-order number is formed there. Each point
  is taken exactly: its acceptance-side component as the exact rational
  iterate formed on the exact island
  (Requirement~\ref{req:acceptance-iterate-exact}), a point whose component
  was not so formed being not certified
  (\iface{acceptance_iterate_not_exact}), and each Choi block as the exact
  Hermitian part of the dyadic matrix its binary64 entries are.
\item \emph{The data.} Every quantity is formed from the data of the
  declaration it works on, taken exactly, every binary64 number as the dyadic
  rational it is and the acceptance set as its exact rational data, and, for
  the blocks a reference companion replaces, from the exact centres and radii
  of the entrywise enclosures its compilation carries, recentred as
  \xref{foundations/lem:recentred-enclosure} recentres them, and from the
  exact centres and radii of the \exttermref{../foundations/contract.pdf}{term-matrix-box}{matrix box} and the exact core of
  every \exttermref{../foundations/contract.pdf}{term-factor-enclosure}{factor enclosure} its compilation carries
  (\xref{foundations/def:factor-enclosure}). No map, frame,
  gradient, support, rank, eigenvalue or check value the engine computed in
  binary64 enters it. In \exttermref{../transmitter/contract.pdf}{term-trusted-source-mode}{trusted-source mode} the frame of each Alice
  sector is $W_{\mathsf t}$ of Theorem~\ref{thm:choi} for the
  \termref{term-support-frame}{support frame} the engine constructed
  (Definition~\ref{def:support-frame}), held as an enclosure formed by a
  rigorously enclosed computation; where the engine forms none, every point
  is not certified (\iface{frame_not_enclosed}).
\item \emph{The maps in factor form.} At every joint pair that is not
  \termref{term-key-inert}{key-inert}, every kept one under a \termref{term-dropped-set}{dropped set}, the branch
  operators and their pinched blocks are formed through factors
  $S'_{xv}$ with the Gram matrices of
  Lemma~\ref{lem:kraus-freedom}\,(v), built from factor forms
  $E^{\dagger}\widetilde XE$ of the elements of Alice's measurement or
  ensemble and of the measured family that enter them, $\widetilde X$
  positive definite and $E$ either an exact matrix of full row rank
  (\xref{foundations/cor:least-eigenvalue-over-enclosure}\,(iii)) or, on a
  block whose rank the factor rule of (iv) fixes, the adjoint $L^{\dagger}$
  of a factor $L$ of its factor enclosure, held through the box of the
  adjoints, of full row rank at every factor
  (\xref{foundations/cor:rank-over-factor-enclosure}\,(ii)), with
  $\widetilde X$ the exact core: on exact data from an exact factorisation,
  on a block a companion carries with a positive radius by the rule of (iv).
  In the general format an element of the measured family is the tensor
  product of its blocks, one per measured factor, and its factor form is
  the tensor product of theirs (Lemma~\ref{lem:certified-bracket-enclosed}\,(b)),
  $E$ and $\widetilde X$ the tensor products of the blocks' $E$ and
  $\widetilde X$, however many of its blocks a companion replaces and
  however many of them are held through boxes: $E$ is then held through an
  enclosure of the Kronecker product of their enclosures, which contains it
  at every member, and has full row rank at every member, $EE^{\dagger}$
  being the tensor product of the blocks' $EE^{\dagger}$, each positive
  definite. No rank is read off a computed spectrum, and a matrix function applied to
  an operator $A$, a branch operator $X$ or a pinched block
  $(Z_j\otimes\id)X^{p}(Z_j\otimes\id)$ of one, acts on its compression
  $Q^{\dagger}AQ$ by an isometry $Q$, held through an enclosure, as follows.
  \begin{enumerate}[label=(\alph*),leftmargin=2.2em]
  \item \emph{Containment.} The range of $Q$ is established to contain
    $\operatorname{supp}A$, exactly or by an argument on the factors and the
    enclosures that holds at every value they enclose. The range of the
    side-by-side listing of the pair's pass-class factors contains the support
    of every branch operator they form, and the pinching projector applied to
    a subspace containing $\operatorname{supp}X$ contains the support of the
    pinched block (Lemma~\ref{lem:kraus-freedom}\,(iv)(a)). Where the
    listing factors as $T_{\mathrm{out}}\widetilde B\,T_{\mathrm{in}}$, with
    $\widetilde B$ an exact matrix of full column rank, formed exactly,
    $T_{\mathrm{out}}$ the maps applied after it, certified injective by a
    Gram matrix certified positive definite, and $T_{\mathrm{in}}$ the maps
    applied before it, an orthonormal basis of
    $\operatorname{ran}(T_{\mathrm{out}}\widetilde B)$ is established from
    $\widetilde B$. On exact data $\widetilde B$ comes from an exact rank
    factorisation of the exact matrix of the factors of the
    \termref{term-pass-class}{pass classes}, and $T_{\mathrm{in}}$ holds the rest of
    that factorisation and the frame $W_{\mathsf t}\otimes\id$ in the Choi
    form. Where an element is held through the box of the adjoints of the
    factors of one or more of its blocks, $\widetilde B$ is, for each pass class, the exact factor of
    Lemma~\ref{lem:certified-bracket-enclosed}\,(b), formed by exact rank
    factorisations of the rows whose factors are exact and, per tuple of the
    blocks held through boxes, of the stack of Alice's factors tensored with
    the exact factors of the elements whose blocks held through boxes are
    that tuple, the kind that lemma fixes for a block being the third for a
    block that the factor rule of (iv)(b) decides not zero, fixing its rank,
    the first for a block decided zero, and the second otherwise, an exact
    block included. The
    rows of such elements are grouped by their tuples, each block named by
    its factor, its sector and its letter, and never by the equality of
    their boxes: elements whose tuples name the same blocks share a group,
    whatever their boxes, and two distinct blocks with equal boxes are
    distinct entries of the tuples. Where the form the engine holds of such
    an element is not of this shape (its parts and the blocks naming them
    differing in number, a split into exact and boxed parts held where no
    block of the element is held through a box or missing where one is, or
    one tuple of blocks at two positions or of two dimensions), no claim
    admits it, and every point is not certified
    (\iface{boxed_form_inconsistent}), a defect of the engine and not of the
    declaration. $T_{\mathrm{in}}$ holds that lemma's
    matrix $C$, the stack of $C_0$ and of each tuple's
    $(C_{\beta}\otimes E^{\mathrm{box}}_{\beta})U_{\beta}$, the columns of
    each such block permuted by its own $U_{\beta}$ and by no other
    permutation, which lies in an enclosure containing it at every member,
    and the frame. That range contains the range of the listing at every member,
    and equals it at every member where the product of the maps in
    $T_{\mathrm{in}}$ is established to have full row rank at every member;
    otherwise the basis spans a subspace containing the range, which (a)
    admits and the floor of (b) decides. $T_{\mathrm{out}}$ is block
    diagonal over the outcome pairs, so the pinching projectors commute with
    it, and $(Z_j\otimes\id)\operatorname{ran}(T_{\mathrm{out}}\widetilde B)
    =\operatorname{ran}\big(T_{\mathrm{out}}(Z_j\otimes\id)\widetilde B\big)$,
    whose basis comes in the same way from an exact rank factorisation of
    $(Z_j\otimes\id)\widetilde B$. A point at which containment is not
    established is not certified (\iface{support_not_established}).
  \item \emph{The floor.} A positive lower bound on $\lmin(Q^{\dagger}AQ)$
    is then established, over the enclosures at every value they contain,
    which with (a) proves $\operatorname{ran}Q=\operatorname{supp}A$
    (Lemma~\ref{lem:kraus-freedom}\,(iv)(b)) at every member and makes the
    compression that of part (iv)(c) of that lemma. A point at which it is
    not is not certified (\iface{floor_not_certified}).
  \end{enumerate}
\item \emph{The rank rule over an enclosure.} A block a companion carries
  with a positive radius is decided, and its rank fixed, by the first of
  two rules that does so.
  \begin{enumerate}[label=(\alph*),leftmargin=2.2em]
  \item \emph{The entrywise rule.} The block is decided when every entry of
    its entrywise enclosure has centre and radius $0$, then zero, or some
    entry's centre exceeds its radius in modulus, then not zero. Its factor
    form is that of
    \xref{foundations/cor:least-eigenvalue-over-enclosure}\,(iv): the rows
    and columns whose every entry has centre and radius $0$ are dropped, and
    the rest must have a positive lower bound on its least eigenvalue over
    the enclosure.
  \item \emph{The factor rule}, where the compilation carries a factor
    enclosure of the block and the entrywise rule leaves the block
    undecided or fixes no rank for it. With no column the block is decided
    zero (\xref{foundations/cor:rank-over-factor-enclosure}\,(i)). With
    $k\ge1$ columns and a floor of part (ii) of that corollary established,
    $G$ and $\rho$ formed as its part (iii) forms them or by a rigorously
    enclosed computation of $L^{\dagger}L$ over the box, the block is
    decided not zero and its factor form is $L^{\dagger}$, held through the
    box of the adjoints, with the core as $\widetilde X$.
  \item \emph{The cross-check}, wherever the compilation carries a factor
    enclosure of the block, whichever rule fixes its rank. The engine
    encloses $LWL^{\dagger}$ entry by entry over the box of the factor
    enclosure, in exact rational arithmetic or by a rigorously enclosed
    computation, so that the enclosure contains $LWL^{\dagger}$ for every
    factor $L$, every member of the factor enclosure; and it decides, as (b)
    decides them, whether the factor enclosure has no column and, where it
    has columns, whether a floor of part (ii) of
    \xref{foundations/cor:rank-over-factor-enclosure} is established.
  \end{enumerate}
  A block that neither rule decides is undecided, and a block whose rank
  neither fixes has no certified rank. The two enclosures of a block are
  \emph{inconsistent} where the enclosure of $LWL^{\dagger}$ of (c) and the
  block's entrywise enclosure are disjoint at some entry, or where one of the
  two decides the block zero and the other decides it not zero, as (a) and
  (c) decide them. The companion's own block lies in its entrywise enclosure
  and, by (H-FAC) of Lemma~\ref{lem:certified-bracket-enclosed}, is a member
  of its factor enclosure, so it lies in both enclosures and both decisions
  hold of it: an inconsistency shows that one of the two was not formed or
  carried as the claims that establish it say. Where the two enclosures of a
  block are inconsistent, the factor rule is set aside for the block, every
  point of the companion is not certified
  (\iface{factor_enclosure_inconsistent}), whatever joint pairs the block is a
  factor at, key-inert or not, and whichever rule would fix its rank, and the
  companion's bracket is withdrawn with that cause
  (Requirement~\ref{req:bracket-on-reference}\,(iv)), its first-order number
  reading the entrywise enclosure, which may be the defective one. On a recentred
  enclosure, the entrywise rule leaves a block undecided exactly when its
  centre vanishes and some radius does not, its enclosure then deciding
  neither that it vanishes nor that it does not. Where an undecided block is
  a factor of an element of an outcome pair of a \termref{term-pass-class}{pass class} of a
  \termref{term-key-live-acceptance-value}{key-live acceptance value} whose element of Alice's is not zero, or
  a block with no certified rank is such a factor at a joint pair that is
  not key-inert, every point is not certified (\iface{rank_not_certified}),
  save where another block of that element is decided zero, an exact block
  that is zero or a block a rule decides zero: the element is then zero for
  every member, whatever its other blocks.
  What the rule establishes is the hypothesis of
  Lemma~\ref{lem:certified-bracket-enclosed} that such blocks are decided
  and of fixed rank, in the sense of that lemma; its hypothesis (H-FAC), that
  the companion's block is a member of the factor enclosure carried with it,
  is established by the claims of the modules that form and carry that
  enclosure, and decided by no rule here.
\item \emph{What is established at a point.} At each point the engine
  establishes, each step in exact rational arithmetic or in a rigorously
  enclosed computation at a working precision of at least $128$ bits:
  \begin{enumerate}[label=(\alph*),leftmargin=2.2em]
  \item every Choi block positive definite, by a positive lower bound of
    \xref{foundations/lem:least-eigenvalue-bound} on its exact Hermitian
    part (\iface{block_not_positive_definite}), which gives (H-SUPP) by
    Lemma~\ref{lem:branch-support}\,(i);
  \item (H-POS), with $q_t(\hat c)>0$ decided exactly and the lower ends of
    the enclosures of $P^{J}_{\hat c}(J_t)$ and $K^{J}_{\hat c}(J_t)$
    positive (\iface{positivity_not_certified}), on the exact
    \termref{term-acceptance-support}{acceptance support} $\Chacc$ of \eqref{eq:chacc-def}, which is
    also where $\hat g$ of (c) lives: for a singleton the support of its
    one point, and for a polytope the union of the supports of its exact
    vertices (Remark~\ref{rem:acc-side});
  \item an enclosure $(\hat g,\delta)$ of the acceptance-side gradient and,
    for every block $p$ of every group, an exact Hermitian $M_p$ and a
    rational $\rho_p$ with $\|R_p-M_p\|_{\mathrm F}\le\rho_p$ for the
    residual of Lemma~\ref{lem:certified-bracket}\,(iii), both formed by a
    rigorously enclosed computation of the gradient of
    Corollary~\ref{cor:full-gradient} at the point (of the dropped objective
    under a dropped set) and of the frame: the enclosure of $R_p$ is made
    Hermitian-centred with symmetric radii by
    \xref{foundations/lem:recentred-enclosure}\,(iv), $M_p$ is its centre
    rounded entrywise to dyadic rationals and kept Hermitian, and $\rho_p$
    is a rational at least the radius that
    \xref{foundations/cor:least-eigenvalue-over-enclosure}\,(ii) gives for
    that centre plus a rational bound on the Frobenius norm of the
    rounding, so that the triangle inequality gives
    $\|R_p-M_p\|_{\mathrm F}\le\rho_p$. On a companion, the computation
    reads each block of the measured family through its entrywise enclosure
    where it forms $P^{J}_{\hat c}$, the term $P^{J}_{\hat c}$ of
    $K^{J}_{\hat c}$ and the $\Gamma^{(s')}_{\hat c}$, and, where it forms
    the deficits and the $\Delta^{(s')}_{\hat c}$, through the enclosure
    that holds the block's factor form of (iii), its entrywise enclosure
    where the entrywise rule of (iv) fixed the rank and its factor
    enclosure where the factor rule did, so that the enclosures of (b) and
    (c) hold at every pair of a set as in
    Lemma~\ref{lem:certified-bracket-enclosed}\,(c). The
    multipliers $Y$ and $\Lambda_r$ are read off the dual of that point's
    solve as Requirement~\ref{req:verified-feasible-side}\,(ii) reads them in
    these respects only: the operators the dual's multipliers represent,
    each replaced by its exact Hermitian part; zero where the solve returned
    no finite dual; and the exact shift that makes the multiplier of an
    inequality row positive semidefinite. The rule of that item that zeroes
    the row multipliers of an Alice sector without a bound on its frame's
    error is not read, the rows here being built from the enclosed exact
    frame;
  \item $b^{\mathrm A}_t$ as the exact minimum of
    $\langle\hat g,\cdot\rangle$ over $\acSet$ for a geometry of
    Requirement~\ref{req:acceptance-geometries} that it solves exactly, less
    $\delta$, a point of any other geometry being not certified
    (\iface{acceptance_geometry_not_certified}); and $b^{\mathrm F}_t$ by
    Lemma~\ref{lem:certified-bracket}\,(iii), each $\ell_p$ from
    \xref{foundations/lem:least-eigenvalue-bound}, in exact rational
    arithmetic;
  \end{enumerate}
  and the point's certified bracket $b^{\mathrm A}_t+b^{\mathrm F}_t$. A
  point at which any of these is not established is not certified, with the
  reason named, and contributes nothing.
\item \emph{The number, its key length and its tier.} Where some point is
  certified, the run's certified number
  $\underline{h}^{\mathrm{cert}}_{\alpha}$ (\iface{certified_bound}) is the
  largest certified bracket, \eqref{eq:hcert}, an exact rational, entering
  every binary64 use rounded toward $-\infty$ as
  \xref{foundations/req:rounding-direction} rounds an emitted bound; the run
  reports it at the \exttermref{../foundations/contract.pdf}{term-certified}{certified} tier (\xref{foundations/req:tiers}).
  The key length it reports is that of
  Corollary~\ref{cor:certified-keylength} (on a reference companion, for the
  protocol it declares, Corollary~\ref{cor:certified-keylength-reference}) at
  $h\coloneqq\max\{0,\underline{h}^{\mathrm{cert}}_{\alpha}\}$, with the
  parameters of the run, and it is the \termref{term-maximal-key-length}{maximal key length} of
  Definition~\ref{def:max-key-length} with $h$ in place of $\hlow$, computed
  exactly: $c\coloneqq n\,h-\lEC-\lceil\log_2(1/\eev)\rceil+1$ is formed
  exactly, a rational, $\lceil\log_2(1/\eev)\rceil$ being the least integer
  $k$ with $2^{k}\ge1/\eev$, decided by comparing the rational $1/\eev$ with
  powers of two; the order $\alpha$ is taken exactly, as the dyadic rational
  its binary64 value is, so that $\frac{\alpha}{\alpha-1}$ is a non-zero
  rational; the term
  $\tau\coloneqq\frac{\alpha}{\alpha-1}\log_2(1/\epa)$ exactly where
  $1/\epa$ is a power of two, $\tau$ being then rational, and otherwise
  through an enclosure $[L,U]\ni\tau$ formed in exact arithmetic or by a
  rigorously enclosed computation and refined until
  $\lfloor c-U\rfloor=\lfloor c-L\rfloor$, which is then
  $\lfloor c-\tau\rfloor$; the refinement ends, $\tau$ being irrational
  there, as a non-zero rational multiple of $\log_2(1/\epa)$, which is
  irrational for a dyadic rational $\epa$ whose reciprocal is not a power
  of two, so that $c-\tau$ is not an integer. The key length is
  $\max\{0,\lfloor c-\tau\rfloor\}$, and the rule of
  Requirement~\ref{req:reported-key-length} is met as written. The clamp is
  recorded as
  Requirement~\ref{req:first-order-clamp} records it. The
  first-order number and the key length formed from it, by
  Requirement~\ref{req:bracket-on-reference}\,(iii) or without a companion,
  are recorded beside it (\iface{first_order_key_length}). In the fallback
  of (i), the certified number is the largest certified bracket on the
  declaration the run searched, a lower bound on that declaration's value
  and not on the companion's, and its key length is that of
  Corollary~\ref{cor:certified-keylength} on that declaration, never that of
  Corollary~\ref{cor:certified-keylength-reference}; it is reported at the
  certified tier, and the blocks that
  Requirement~\ref{req:bracket-on-reference}\,(v) names for the fallback
  rest on \xref{foundations/ass:snapped-measurement}. The companion's
  first-order number stays the run's first-order number, recorded beside it
  with its key length, and no number is formed from the two together: the
  key length is that of the fallback's certified number alone, whether the
  companion's first-order number gives a larger key length or a smaller one.
  Where the companion's bracket is withdrawn by (iv), the run's first-order
  number is that of Requirement~\ref{req:bracket-on-reference}\,(iv) on the
  declaration the run searched, a number certified there is a lower bound on
  that declaration's value with the key length of
  Corollary~\ref{cor:certified-keylength}, and the blocks the companion
  carried rest on \xref{foundations/ass:snapped-measurement} with the cause
  \iface{factor_enclosure_inconsistent}; a run that declares
  \iface{strict_snaps} refuses there (\xref{foundations/req:snap-reference}\,(iv)).
  Where no point is certified, or the arithmetic the evaluation needs is not
  available to the engine (\iface{certifier_unavailable}), the run's number,
  key length and tier are those it has without this requirement, recorded
  with the reason \iface{not_certified} and the reasons of every point.
\item \emph{Premises and evaluations upstream of the declaration.} A run is
  not reported at the certified tier where (a) an assumption of the category
  numerical-trust is invoked on it, reason
  \iface{numerical_trust_premise}: an assumption of its assumption closure
  counts as invoked unless the module that owns it records that it was
  invoked nowhere on the run, as the receiver records its eigenvalue-slack
  assumption beside the settings at which it was invoked
  (\xref{receiver/req:eig-ledger-row}); or (b) the declaration it is formed
  on carries a weight routed from a receiver's high-photon weights
  (\xref{reductions/wpfss/def:weight-route}) into which a per-setting weight
  enters that the receiver's artifact records as evaluated in floating point
  (\xref{receiver/req:evaluation-label}), the uniform weight of a slot entry
  counting the weight of every setting as entering it, reason
  \iface{receiver_weight_floating_point}. Its certified number is then
  recorded with the reason, and the run's number, key length and tier are
  those it has without this requirement.
\item \emph{The strict option.} A run that sets
  \iface{certified_evaluation} to \iface{strict} and is not reported at the
  certified tier by (vi) and (vii) refuses
  (\iface{certification_refused}), naming the reasons, and reports no key
  length. A run reported at the certified tier through the fallback of (i)
  is so reported under this option too and does not refuse: the premise it
  rests on there, \xref{foundations/ass:snapped-measurement}, is named in its
  record, and the option that refuses a run resting on it is
  \iface{strict_snaps}, under which the fallback does not run.
\item \emph{The record.} The run records the option; per point, whether it
  was certified, $b^{\mathrm A}_t$, $b^{\mathrm F}_t$, its certified bracket,
  the largest radius of its enclosure of the gradient and its reasons; and
  per run $\underline{h}^{\mathrm{cert}}_{\alpha}$, the largest such radius
  over the certified points (\iface{max_gradient_radius}), the working
  precision (\iface{working_precision_bits}), the library of the
  rigorous arithmetic with its version (\iface{arithmetic_library}), and,
  as a diagnostic, $\underline{h}^{\mathrm{cert}}_{\alpha}$ less the
  first-order number on the same declaration, with its sign
  (\iface{certified_minus_first_order}), which in the fallback is the bound
  on the searched declaration that
  Requirement~\ref{req:bracket-on-reference}\,(iii) records and is absent
  where that declaration has none. It records further the declaration the
  per-point entries were evaluated on (\iface{declaration}:
  \iface{reference_companion} or \iface{searched}); the premise at the
  snap the certified number rests on (\iface{premises}): the short id of
  \xref{foundations/ass:snapped-measurement} wherever a moved block or class
  of the declaration the number is formed on rests on that assumption
  (\xref{foundations/req:snap-reference}\,(iii)), among them, on the
  declaration the run searched, every moved block or class it has, in the
  fallback, without a companion and where the companion's bracket is
  withdrawn alike, and, on a companion, every moved block on which it carries
  the snapped block and every moved block of another snap it shares with the
  declaration the run searched, such as that of another measured factor;
  and none otherwise; with it, the blocks and classes that
  rest on it (\iface{resting_blocks}), with their groups and causes as the
  record of the run names them where the key length is formed from a number
  on that declaration, those of
  Requirement~\ref{req:bracket-on-reference}\,(v) among them in the
  fallback, with the cause \iface{companion_not_certified}; and, wherever
  the fallback ran or the companion's bracket was withdrawn by (iv), the
  companion's points, each with whether it was certified and its reasons
  (\iface{companion_points}), which say why. The fields \iface{premises}
  and \iface{resting_blocks} describe the declaration that
  \iface{declaration} names, whatever the status: they name what a number
  certified there rests on, also where none is reported at the certified
  tier. In a fallback not reported at the certified tier, the status being
  \iface{withheld} or \iface{not_certified}, they therefore name the blocks
  of Requirement~\ref{req:bracket-on-reference}\,(v) with the cause
  \iface{companion_not_certified}, under which that item names them where
  the fallback's number forms the key length, while the run's key length is
  formed from the companion's first-order number and the record of the run
  names those blocks as evaluated off their snap. None of these but
  $\underline{h}^{\mathrm{cert}}_{\alpha}$ enters a number.
\end{enumerate}
\end{requirement}
\claimreading{After the search, the engine recomputes the bound at the few points the reported number came from, in an arithmetic that encloses every rounding, from the declaration's data. Where that succeeds at some point, the key rate is certified: no step between the declaration and the number rests on floating-point arithmetic being accurate enough. Where it fails on the model the engine holds of your device, it is tried again on the moved measurement the search ran on, unless the run sets strict\_snaps; a rate certified there is reported as certified, and the report names the premise that your device performs the moved measurement on the blocks concerned (\hyperref{../foundations/contract.pdf}{}{claim-foundations/ass:snapped-measurement}{A-SNAP}). Where both fail, for instance because a detector element has a rank no enclosure fixes or a multiclick weight was computed in floating point, the report keeps the first-order number and says why; a strict run is refused instead.}
\claimmeta{\metalabel{Motivated by}\mref{Def.}{def:support-frame}, \mref{Thm.}{thm:choi}, \mref{Lem.}{lem:branch-support}, \mref{Cor.}{cor:full-gradient}, \mref{Lem.}{lem:kraus-freedom}, \mref{Def.}{def:max-key-length}, \mref{Req.}{req:reported-key-length}, \mref{Def.}{def:first-order-setting}, \mref{Thm.}{thm:firstorder}, \mref{Lem.}{lem:feasible-side-dual}, \mref{Lem.}{lem:certified-bracket}, \mref{Lem.}{lem:certified-bracket-enclosed}, \mref{Cor.}{cor:certified-keylength}, \mref{Cor.}{cor:certified-keylength-reference}, \mref{Rem.}{rem:acc-side}, \mref{Rem.}{rem:fo-arithmetic}, \mref{Rem.}{rem:certified-scope}, \mref{Req.}{req:first-order-withdrawal}, \mref{Req.}{req:first-order-clamp}, \mref{Req.}{req:bracket-on-reference}, \mref{Req.}{req:acceptance-geometries}, \mref{Req.}{req:acceptance-iterate-exact}, \mref{Req.}{req:verified-feasible-side}, \extref{../foundations/contract.pdf}{foundations/ass:snapped-measurement}, \extref{../foundations/contract.pdf}{foundations/cor:least-eigenvalue-over-enclosure}, \extref{../foundations/contract.pdf}{foundations/cor:rank-over-factor-enclosure}, \extref{../foundations/contract.pdf}{foundations/def:factor-enclosure}, \extref{../foundations/contract.pdf}{foundations/def:tiers}, \extref{../foundations/contract.pdf}{foundations/lem:least-eigenvalue-bound}, \extref{../foundations/contract.pdf}{foundations/lem:matrix-function-perturbation}, \extref{../foundations/contract.pdf}{foundations/lem:recentred-enclosure}, \extref{../foundations/contract.pdf}{foundations/req:exact-island}, \extref{../foundations/contract.pdf}{foundations/req:rounding-direction}, \extref{../foundations/contract.pdf}{foundations/req:snap-reference}, \extref{../foundations/contract.pdf}{foundations/req:tiers}, \extref{../receiver/contract.pdf}{receiver/req:eig-ledger-row}, \extref{../receiver/contract.pdf}{receiver/req:evaluation-label}, \extref{../reductions/contract.pdf}{reductions/wpfss/def:weight-route}
\metasep \metalabel{Used by}\mref{Rem.}{rem:certified-scope}
\metasep implemented (effective: accepted)
\metasep \metalabel{Code}\coderef{https://github.com/fraunhofersg/certiqs-key/blob/v0.1.0/src/certiqs_key/solvers/frank_wolfe/certified.py}{certiqs_key.solvers.frank_wolfe.certified:evaluate}}
\clearpage
\bookmarksetupnext{level=part}
\section*{Interface}\addcontentsline{toc}{section}{Interface}
\sectionintro{The declared inputs and outputs of the module and its submodules: each key with its symbol, when it is present and the claims that fix it. The domains, types and refusal rules are in the interface files themselves.}
\subsection*{Inputs: MEAT key-rate engine (meat-core)}
\noindent{\small From \href{interface.yaml}{\nolinkurl{interface.yaml}}.}\par\smallskip
{\footnotesize\begin{xltabular}{\linewidth}{@{}>{\raggedright\arraybackslash}p{4.4cm}>{\raggedright\arraybackslash}X>{\raggedright\arraybackslash}p{2.9cm}@{}}
\toprule
Key and symbol & Presence & Fixed by\\
\midrule
\endhead
\iface{alice_dim}\newline d\_A & always & \mref{Def.}{def:declaration}, \mref{Def.}{def:admissible-declaration}, \mref{Rem.}{rem:modes-da1}\\
\addlinespace
\iface{alice_marginal}\newline tau\_A, or null & always & \extref{../protocol-compiler/contract.pdf}{protocol-compiler/def:declaration}, \mref{Def.}{def:declaration}, \mref{Def.}{def:admissible-declaration}, \mref{Rem.}{rem:modes-da1}, \mref{Def.}{def:support-frame}, \mref{Lem.}{lem:choi-marginal}, \mref{Rem.}{rem:basis-invariance}\\
\addlinespace
\iface{sector_index_set}\newline S & always & \mref{Def.}{def:declaration}, \mref{Def.}{def:admissible-declaration}, \mref{Def.}{def:feasible}\\
\addlinespace
\iface{sector_dims}\newline (d\_s)\_s & always & \mref{Def.}{def:declaration}, \mref{Def.}{def:admissible-declaration}\\
\addlinespace
\iface{sector_block_kind} & always & \extref{../protocol-compiler/contract.pdf}{protocol-compiler/def:declaration}, \mref{Def.}{def:declaration}\\
\addlinespace
\iface{alice_povm}\newline \{M\_x\}\_\{x in X\} & always & \mref{Def.}{def:declaration}, \mref{Def.}{def:admissible-declaration}, \mref{Thm.}{thm:protocol-channel}\\
\addlinespace
\iface{bob_povm}\newline \{F\_v\}\_\{v in V\} & always & \mref{Def.}{def:declaration}, \mref{Req.}{req:exact-admission}, \mref{Def.}{def:admissible-declaration}, \extref{../protocol-compiler/contract.pdf}{protocol-compiler/def:exact}\\
\addlinespace
\iface{constraint_list}\newline CList = \{(hat O\_k, bowtie\_k, c\_k)\}\_\{k=1..m\} & always; m \textgreater{}= 0, and m = 0 is the default route & \mref{Def.}{def:declaration}, \mref{Def.}{def:admissible-declaration}, \mref{Def.}{def:homog}, \mref{Rem.}{rem:cgate}, \mref{Rem.}{rem:homog}\\
\addlinespace
\iface{attack_set_license} & optional; required by part (ii) of the two-part gate whenever m \textgreater{}= 1 and a key-length statement is wanted. Absent there, the engine refuses the key-length output (req:constrained-key-gate): the key-length entry points raise SecurityGateError naming the failing part, and a run emits no key length, its reason naming the failing part unless h\_low is infinite, which is reported first. The program is still solved and its numbers are reported for exploration; they do not reach eq:keylength. & \extref{../protocol-compiler/contract.pdf}{protocol-compiler/ass:atk}, \mref{Thm.}{thm:constrained-accumulation}, \mref{Cor.}{cor:constrained-key-length}, \mref{Rem.}{rem:cgate}, \mref{Rem.}{rem:not-licensed}, \mref{Req.}{req:constrained-key-gate}\\
\addlinespace
\iface{source_replacement_license}\newline the transmitter's source-replacement licence record, or absent & read in trusted-source mode at d\_A \textgreater{}= 2; at d\_A = 1 it is vacuous and must not be demanded (prop:trivial-alice-system (iv); req:marginal-licence-gate; transmitter/req:licence-scope) & \mref{Req.}{req:marginal-licence-gate}, \mref{Rem.}{rem:modes-da1}, \mref{Prop.}{prop:trivial-alice-system}, \extref{../transmitter/contract.pdf}{transmitter/ass:src}, \extref{../protocol-compiler/contract.pdf}{protocol-compiler/def:eligible}\\
\addlinespace
\iface{party_partition}\newline Party & always & \mref{Def.}{def:declaration}, \mref{Def.}{def:admissible-declaration}, \mref{Ass.}{ass:lab}, \extref{../transmitter/contract.pdf}{transmitter/ass:src}\\
\addlinespace
\iface{public_structure}\newline (a\_A, a\_B, phi, r, h) & always & \mref{Def.}{def:declaration}, \mref{Def.}{def:admissible-declaration}\\
\addlinespace
\iface{label_index_map}\newline idxm & always & \mref{Def.}{def:declaration}, \extref{../foundations/contract.pdf}{foundations/req:declared-order}\\
\addlinespace
\iface{acceptance_set}\newline Qacc & always & \mref{Def.}{def:admissible-acceptance-set}, \mref{Def.}{def:feasible}, \mref{Rem.}{rem:hlow-range}, \mref{Lem.}{lem:fo-khyp}\\
\addlinespace
\iface{rounds}\newline n & always & \mref{Thm.}{thm:keylength}, \mref{Cor.}{cor:keylength-bound}, \mref{Def.}{def:admissible-declaration}\\
\addlinespace
\iface{sector_factorisation}\newline the ordered factor list of (DX1) & general format only; absent is one measured factor and no A factor & \extref{../protocol-compiler/contract.pdf}{protocol-compiler/def:declaration}, \mref{Def.}{def:declaration}, \mref{Def.}{def:feasible}, \mref{Thm.}{thm:choi}, \mref{Def.}{def:joint-sectors}\\
\addlinespace
\iface{alice_form}\newline `povm` or `ensemble` & general format only; absent is `povm` & \extref{../protocol-compiler/contract.pdf}{protocol-compiler/def:declaration}, \mref{Lem.}{lem:ensemble-form}, \mref{Lem.}{lem:choi-factor}\\
\addlinespace
\iface{alice_ensemble}\newline \{N\_x\}\_\{x in X\} & general format, ensemble form only; in place of `alice\_povm` & \extref{../protocol-compiler/contract.pdf}{protocol-compiler/def:declaration}, \mref{Def.}{def:declaration}, \mref{Lem.}{lem:ensemble-form}, \mref{Lem.}{lem:choi-factor}, \mref{Thm.}{thm:choi}\\
\addlinespace
\iface{bob_povm.per_factor}\newline (\{F\textasciicircum{}f\_v\_f\})\_\{f in F\} & general format only; one factor is the `bob\_povm` row & \extref{../protocol-compiler/contract.pdf}{protocol-compiler/def:declaration}, \mref{Def.}{def:declaration}, \mref{Def.}{def:joint-sectors}, \extref{../reductions/contract.pdf}{reductions/thm:emitted-product}\\
\addlinespace
\iface{transcript}\newline ((pi\_m, C\_m, a\_m))\_\{m = 1..L\} & general format only; L = 2 is the pair `announce\_alice`, `announce\_bob` & \extref{../protocol-compiler/contract.pdf}{protocol-compiler/def:declaration}, \mref{Def.}{def:declaration}, \mref{Def.}{def:maps}, \mref{Def.}{def:admissible-declaration}\\
\addlinespace
\iface{constraint_scope_and_party_local_tag}\newline (O\_k, rel\_k, c\_k; f\_k) & general format only & \extref{../protocol-compiler/contract.pdf}{protocol-compiler/def:declaration}, \extref{../protocol-compiler/contract.pdf}{protocol-compiler/def:homog}, \mref{Def.}{def:homog}, \mref{Def.}{def:feasible}, \mref{Def.}{def:classes}, \mref{Rem.}{rem:removal-general-format}, \mref{Req.}{req:removal-checklist}, \mref{Thm.}{thm:removal}\\
\addlinespace
\iface{reference_companion}\newline D\textasciicircum{}ref, beside D & where the protocol compiler emits one beside the declaration (protocol-compiler/req:reference-companion-compiled) & \mref{Prop.}{prop:reference-companion}, \mref{Cor.}{cor:reference-companion-factors}, \mref{Req.}{req:companion-admission}, \mref{Req.}{req:exact-admission}, \mref{Req.}{req:declaration-admission}, \mref{Cor.}{cor:fo-reference}\\
\bottomrule
\end{xltabular}}
\subsection*{Outputs: MEAT key-rate engine (meat-core)}
\noindent{\small From \href{interface.yaml}{\nolinkurl{interface.yaml}}.}\par\smallskip
{\footnotesize\begin{xltabular}{\linewidth}{@{}>{\raggedright\arraybackslash}p{4.4cm}>{\raggedright\arraybackslash}X>{\raggedright\arraybackslash}p{2.9cm}@{}}
\toprule
Key and symbol & Presence & Fixed by\\
\midrule
\endhead
\iface{h_low}\newline h\_low & always & \mref{Def.}{def:feasible}, \mref{Lem.}{lem:estimate}, \mref{Rem.}{rem:hlow-range}, \mref{Thm.}{thm:reduction}, \mref{Cor.}{cor:reduction-constrained}\\
\addlinespace
\iface{h_lb}\newline h\_lb & The submodule's row, restated here only as far as a consumer of this module needs it: present exactly when at least one iterate contributes a bracket, every hypothesis of frank-wolfe/thm:firstorder having held at it; absent otherwise, with a machine-readable reason naming the hypothesis that failed (frank-wolfe/req:first-order-withdrawal). The oracle's return and its status remove no iterate, the feasible side of every bracket being a lower bound verified without them (frank-wolfe/req:verified-feasible-side). The domain, the withdrawal vocabulary, the clamp and the recorded repairs are spec/modules/meat-core/frank-wolfe/interface.yaml's. & \mref{Thm.}{thm:firstorder}, \mref{Req.}{req:verified-feasible-side}, \mref{Req.}{req:first-order-withdrawal}, \mref{Req.}{req:first-order-clamp}, \mref{Rem.}{rem:fo-interface}, \mref{Rem.}{rem:fo-negative-values}, \mref{Rem.}{rem:fo-arithmetic}\\
\addlinespace
\iface{reported_tier} & always & \extref{../foundations/contract.pdf}{foundations/def:tiers}, \extref{../foundations/contract.pdf}{foundations/req:tiers}, \mref{Rem.}{rem:fo-interface}, \mref{Cor.}{cor:fo-keylength}, \mref{Req.}{req:certified-evaluation}, \mref{Cor.}{cor:certified-keylength}\\
\addlinespace
\iface{key_length}\newline ell & Emitted only on a run whose `reported\_tier` is `first-order` or `certified`, and formed from the number of that tier. Absent -- not zero, not clamped -- when `h\_lb` is absent, when h\_low = +infinity, and on a declaration the run tier did not pass as the compiler's emission (req:compiled-declaration-only; `key\_length\_reason` then opens with `declaration\_not\_compiled`); in every case the reason is machine-readable and no key-length claim is made. & \mref{Cor.}{cor:keylength-bound}, \mref{Thm.}{thm:keylength}, \mref{Cor.}{cor:fo-keylength}, \mref{Cor.}{cor:certified-keylength}, \mref{Cor.}{cor:certified-keylength-reference}, \mref{Prop.}{prop:monotone}, \mref{Def.}{def:max-key-length}, \mref{Req.}{req:reported-key-length}, \mref{Req.}{req:key-length-alpha-domain}, \mref{Rem.}{rem:alpha-domain}\\
\addlinespace
\iface{security_parameter}\newline eps\_PA + eps\_EV + 2 eps\_auth & emitted exactly when `key\_length` is & \mref{Cor.}{cor:composed-security}, \mref{Ass.}{ass:auth}, \mref{Cor.}{cor:separate-executions}, \mref{Rem.}{rem:repeated-executions}\\
\addlinespace
\iface{net_key_length}\newline ell - k\_auth & emitted exactly when `key\_length` is & \mref{Cor.}{cor:composed-security}, \mref{Ass.}{ass:auth}\\
\addlinespace
\iface{sector_report} & always & \mref{Def.}{def:classes}, \mref{Def.}{def:keyinert}, \mref{Thm.}{thm:removal}, \mref{Prop.}{prop:retention}, \mref{Prop.}{prop:no-averaging}, \mref{Prop.}{prop:removal-compose}, \mref{Rem.}{rem:classes}, \mref{Req.}{req:sector-report}, \mref{Req.}{req:removal-checklist}\\
\addlinespace
\iface{sector_weight_report}\newline (w\_s)\_s & always & \mref{Thm.}{thm:choi}, \mref{Def.}{def:relaxed-feasible-set}, \mref{Prop.}{prop:relaxed-set-properties}, \mref{Def.}{def:classes}\\
\addlinespace
\iface{map_validation_report} & always & \mref{Def.}{def:maps}, \mref{Thm.}{thm:protocol-channel}, \mref{Lem.}{lem:branch-maps}\\
\addlinespace
\iface{optimizer_certificate} & always & \mref{Rem.}{rem:fo-returned-point}, \mref{Rem.}{rem:fo-iterate-checks}, \mref{Rem.}{rem:fo-oracle-repair}, \mref{Rem.}{rem:fo-iterate-repair}, \mref{Req.}{req:verified-feasible-side}, \mref{Req.}{req:frame-recorded}, \mref{Req.}{req:row-form-record}, \mref{Def.}{def:support-frame}, \mref{Rem.}{rem:ordering}, \mref{Rem.}{rem:basis-invariance}, \mref{Def.}{def:untrusted-choi}, \mref{Prop.}{prop:frame-free-rows}\\
\addlinespace
\iface{refusal} & present exactly when this module refuses & \mref{Req.}{req:exact-admission}, \mref{Def.}{def:declaration}, \mref{Req.}{req:first-order-withdrawal}, \mref{Req.}{req:negative-estimate-refused}\\
\addlinespace
\iface{assumption_ledger}\newline A-... & always; empty on a run that reports neither a first-order bound nor a key length & ---\\
\addlinespace
\iface{assumption_bindings} & always; one entry per `assumption\_ledger` row, in its order & \mref{Ass.}{ass:keyreg}, \mref{Ass.}{ass:hash}, \mref{Ass.}{ass:seed}, \mref{Ass.}{ass:fixed}, \mref{Ass.}{ass:conform}, \mref{Ass.}{ass:auth}, \extref{../protocol-compiler/contract.pdf}{protocol-compiler/ass:atk}, \extref{../transmitter/contract.pdf}{transmitter/ass:src}, \extref{../protocol-compiler/contract.pdf}{protocol-compiler/ass:rhom}, \extref{../transmitter/contract.pdf}{transmitter/ass:seq}\\
\addlinespace
\iface{mathematical_spec_version} & always & ---\\
\addlinespace
\iface{sector_report.joint_pairs}\newline J\_K(c-hat) and the key-inert joint pairs & every run (one Frank-Wolfe driver: an accepted-format declaration is solved as its degenerate general embedding and writes the general report) & \mref{Def.}{def:joint-sectors}, \mref{Lem.}{lem:joint-sectors}, \mref{Def.}{def:classes}\\
\addlinespace
\texttt{sector\_report.records[].removable}\newline the `removable` flag & always & \mref{Thm.}{thm:removal}, \mref{Req.}{req:removal-checklist}, \mref{Rem.}{rem:removal-general-format}\\
\addlinespace
\iface{dropped_set}\newline D, and h\_low\textasciicircum{}D & only when the run chooses a non-empty dropped set & \mref{Def.}{def:dropped}, \mref{Prop.}{prop:drop-bound}, \mref{Cor.}{cor:drop-program}, \mref{Req.}{req:dropped-set-record}, \mref{Cor.}{cor:fo-dropped}\\
\addlinespace
\iface{statistics_consistency}\newline b, alpha\_0, x, tau, n f(alpha') & always & \mref{Req.}{req:unreachable-acceptance-set}, \mref{Prop.}{prop:minimal-divergence}, \mref{Prop.}{prop:model-pass-bound}, \extref{../statistics/contract.pdf}{statistics/ass:reach}\\
\bottomrule
\end{xltabular}}
\subsection*{Inputs: Frank-Wolfe first-order bound}
\noindent{\small From \href{frank-wolfe/interface.yaml}{\nolinkurl{frank-wolfe/interface.yaml}}.}\par\smallskip
{\footnotesize\begin{xltabular}{\linewidth}{@{}>{\raggedright\arraybackslash}p{4.4cm}>{\raggedright\arraybackslash}X>{\raggedright\arraybackslash}p{2.9cm}@{}}
\toprule
Key and symbol & Presence & Fixed by\\
\midrule
\endhead
\iface{choi_program}\newline the feasible set Feas\textasciicircum{}J and its objective & always & \mref{Thm.}{thm:choi}, \mref{Def.}{def:support-frame}, \mref{Lem.}{lem:choi-marginal}, \mref{Def.}{def:untrusted-choi}, \mref{Def.}{def:feasible}\\
\addlinespace
\iface{gradient}\newline the split-form gradient & at each iterate at which a bracket is accumulated & \mref{Thm.}{thm:gradient}, \mref{Cor.}{cor:full-gradient}, \mref{Lem.}{lem:dk}, \mref{Lem.}{lem:fo-khyp}\\
\addlinespace
\iface{acceptance_set_exact_data}\newline Qacc & always & \mref{Def.}{def:admissible-acceptance-set}, \mref{Def.}{def:first-order-setting}, \mref{Lem.}{lem:fo-khyp}, \mref{Rem.}{rem:acc-side}, \mref{Req.}{req:acceptance-geometries}\\
\addlinespace
\iface{docking_claims} & always & \mref{Lem.}{lem:estimate}, \mref{Cor.}{cor:keylength-bound}\\
\addlinespace
\iface{linear_oracle_return}\newline the primal block list and the dual the backend returns at an iterate & at each iterate at which a bracket is accumulated, where the backend returns them & \mref{Req.}{req:verified-feasible-side}, \mref{Lem.}{lem:feasible-side-dual}, \mref{Thm.}{thm:firstorder}, \mref{Def.}{def:first-order-setting}, \mref{Rem.}{rem:fo-oracle-repair}, \mref{Rem.}{rem:fo-iterate-checks}\\
\addlinespace
\iface{dropped_set}\newline D & optional; empty is the run without a dropped set & \mref{Def.}{def:dropped}, \mref{Prop.}{prop:drop-bound}, \mref{Cor.}{cor:drop-program}, \mref{Cor.}{cor:fo-dropped}\\
\bottomrule
\end{xltabular}}
\subsection*{Outputs: Frank-Wolfe first-order bound}
\noindent{\small From \href{frank-wolfe/interface.yaml}{\nolinkurl{frank-wolfe/interface.yaml}}.}\par\smallskip
{\footnotesize\begin{xltabular}{\linewidth}{@{}>{\raggedright\arraybackslash}p{4.4cm}>{\raggedright\arraybackslash}X>{\raggedright\arraybackslash}p{2.9cm}@{}}
\toprule
Key and symbol & Presence & Fixed by\\
\midrule
\endhead
\iface{h_lb}\newline h\_lb & Present exactly when the index set of contributing iterates is non-empty -- when at least one iterate satisfied (H-POS), (H-SUPP) and (H-NOPERT) and lay in the relaxed feasible set with its acceptance factor formed on the exact island. Absent otherwise, with a machine-readable reason from `withdrawal\_reasons` naming what failed; if no iterate survives, the index set is empty and the number is absent for that reason, with the same vocabulary (req:first-order-withdrawal). The oracle's return and status decide nothing here. & \mref{Thm.}{thm:firstorder}, \mref{Req.}{req:verified-feasible-side}, \mref{Req.}{req:first-order-withdrawal}, \mref{Req.}{req:first-order-clamp}, \mref{Rem.}{rem:fo-interface}, \mref{Rem.}{rem:fo-negative-values}, \mref{Rem.}{rem:fo-arithmetic}, \mref{Cor.}{cor:fo-reference}, \mref{Req.}{req:bracket-on-reference}\\
\addlinespace
\iface{reference_bracket}\newline h\_lb[D\textasciicircum{}ref], beside h\_lb[D] & wherever a compiled reference companion was handed to the run (protocol-compiler/req:reference-companion-compiled), whether or not it was admitted; null otherwise. A companion not admitted (meat-core/req:companion-admission (i)) is recorded here with status not\_evaluated and the admission's refusal. & \mref{Req.}{req:bracket-on-reference}, \mref{Req.}{req:certified-evaluation}, \mref{Cor.}{cor:fo-reference}, \mref{Rem.}{rem:fo-arithmetic}, \mref{Req.}{req:companion-admission}, \extref{../foundations/contract.pdf}{foundations/def:evaluated-off-snap}, \extref{../foundations/contract.pdf}{foundations/req:snap-reference}\\
\addlinespace
\iface{first_order_clamped} & always, on a run that emits the bound & \mref{Req.}{req:first-order-clamp}, \mref{Rem.}{rem:fo-negative-values}, \mref{Thm.}{thm:firstorder}, \mref{Cor.}{cor:fo-keylength}\\
\addlinespace
\iface{key_length}\newline ell & the core's row; emitted only on a run whose tier is the first-order or the certified tier, from the number of that tier & \mref{Cor.}{cor:fo-keylength}, \mref{Cor.}{cor:certified-keylength}, \mref{Cor.}{cor:certified-keylength-reference}, \mref{Req.}{req:certified-evaluation}, \mref{Cor.}{cor:keylength-bound}, \mref{Thm.}{thm:keylength}, \mref{Prop.}{prop:monotone}\\
\addlinespace
\iface{optimizer_certificate} & always & \mref{Rem.}{rem:fo-returned-point}, \mref{Rem.}{rem:fo-iterate-checks}, \mref{Rem.}{rem:fo-oracle-repair}, \mref{Rem.}{rem:fo-iterate-repair}, \mref{Req.}{req:iterate-repair}, \mref{Req.}{req:oracle-repair}, \mref{Req.}{req:returned-point-measured}, \mref{Req.}{req:bacc-recorded}, \mref{Req.}{req:verified-feasible-side}, \mref{Rem.}{rem:acc-side}, \mref{Req.}{req:acceptance-iterate-exact}, \mref{Req.}{req:frame-recorded}, \mref{Req.}{req:row-form-record}, \mref{Rem.}{rem:ordering}, \mref{Def.}{def:untrusted-choi}\\
\addlinespace
\iface{h_lb_dropped}\newline h\textasciicircum{}FO,D\_alpha & present exactly when a non-empty dropped set was used and the index set of contributing iterates is non-empty & \mref{Cor.}{cor:fo-dropped}, \mref{Def.}{def:dropped}, \mref{Req.}{req:dropped-set-record}, \mref{Cor.}{cor:keylength-bound}\\
\addlinespace
\iface{certified_bound}\newline h\textasciicircum{}cert\_alpha & Present exactly when the certified evaluation ran and certified at least one point (req:certified-evaluation (vi)); absent where the option is `off`, where the run formed no first-order number, and where no point was certified, with the reasons under `certified\_evaluation` below. & \mref{Req.}{req:certified-evaluation}, \mref{Cor.}{cor:certified-keylength}, \mref{Cor.}{cor:certified-keylength-reference}, \mref{Lem.}{lem:certified-bracket}, \mref{Lem.}{lem:certified-bracket-enclosed}, \mref{Rem.}{rem:certified-scope}, \extref{../foundations/contract.pdf}{foundations/def:tiers}, \extref{../foundations/contract.pdf}{foundations/ass:snapped-measurement}\\
\bottomrule
\end{xltabular}}
\clearpage
\bookmarksetupnext{level=part}
\section*{Claim index}\addcontentsline{toc}{section}{Claim index}
{\small\begin{xltabular}{\linewidth}{@{}>{\raggedright\arraybackslash}p{2.7cm}>{\raggedright\arraybackslash}Xr@{}}
\toprule
Claim & Title and identifier & Page\\
\midrule
\endhead
\mref{Definition}{def:declaration} & The protocol declaration as the MEAT core reads it\newline{\footnotesize\color{metaGrey}\claimid{meat-core/def:declaration}} & \pageref{def:declaration}\\
\addlinespace[2pt]
\mref{Definition}{def:admissible-declaration} & Admissible declaration\newline{\footnotesize\color{metaGrey}\claimid{meat-core/def:admissible-declaration}} & \pageref{def:admissible-declaration}\\
\addlinespace[2pt]
\mref{Lemma}{lem:ensemble-form} & The measurement an ensemble-form declaration denotes\newline{\footnotesize\color{metaGrey}\claimid{meat-core/lem:ensemble-form}} & \pageref{lem:ensemble-form}\\
\addlinespace[2pt]
\mref{Definition}{def:maps} & Generated registers, isometries and maps\newline{\footnotesize\color{metaGrey}\claimid{meat-core/def:maps}} & \pageref{def:maps}\\
\addlinespace[2pt]
\mref{Theorem}{thm:protocol-channel} & The declaration generates a protocol channel\newline{\footnotesize\color{metaGrey}\claimid{meat-core/thm:protocol-channel}} & \pageref{thm:protocol-channel}\\
\addlinespace[2pt]
\mref{Definition}{def:branch} & Branch maps\newline{\footnotesize\color{metaGrey}\claimid{meat-core/def:branch}} & \pageref{def:branch}\\
\addlinespace[2pt]
\mref{Lemma}{lem:branch-maps} & Properties of the pass and fail branch maps\newline{\footnotesize\color{metaGrey}\claimid{meat-core/lem:branch-maps}} & \pageref{lem:branch-maps}\\
\addlinespace[2pt]
\mref{Remark}{rem:inert} & Inert register components of a branch map\newline{\footnotesize\color{metaGrey}\claimid{meat-core/rem:inert}} & \pageref{rem:inert}\\
\addlinespace[2pt]
\mref{Fact}{fact:meat} & Marginal-constrained entropy accumulation, sequential form\newline{\footnotesize\color{metaGrey}\claimid{meat-core/fact:meat}} & \pageref{fact:meat}\\
\addlinespace[2pt]
\mref{Definition}{def:homog} & Licensed constraint list and its two classes\newline{\footnotesize\color{metaGrey}\claimid{meat-core/def:homog}} & \pageref{def:homog}\\
\addlinespace[2pt]
\mref{Definition}{def:roundmodel} & Round model and licensed attack sets\newline{\footnotesize\color{metaGrey}\claimid{meat-core/def:roundmodel}} & \pageref{def:roundmodel}\\
\addlinespace[2pt]
\mref{Definition}{def:hqkd-list} & Constrained target set and single-round rate\newline{\footnotesize\color{metaGrey}\claimid{meat-core/def:hqkd-list}} & \pageref{def:hqkd-list}\\
\addlinespace[2pt]
\mref{Definition}{def:admissible-acceptance-set} & Admissible acceptance set\newline{\footnotesize\color{metaGrey}\claimid{meat-core/def:admissible-acceptance-set}} & \pageref{def:admissible-acceptance-set}\\
\addlinespace[2pt]
\mref{Lemma}{lem:cond-more} & Conditioning on more side information\newline{\footnotesize\color{metaGrey}\claimid{meat-core/lem:cond-more}} & \pageref{lem:cond-more}\\
\addlinespace[2pt]
\mref{Lemma}{lem:purifier-indep} & Independence of the purifier size\newline{\footnotesize\color{metaGrey}\claimid{meat-core/lem:purifier-indep}} & \pageref{lem:purifier-indep}\\
\addlinespace[2pt]
\mref{Lemma}{lem:range} & Constraint transport and range inclusion\newline{\footnotesize\color{metaGrey}\claimid{meat-core/lem:range}} & \pageref{lem:range}\\
\addlinespace[2pt]
\mref{Lemma}{lem:product-marginal} & Product marginal of Alice's retained systems\newline{\footnotesize\color{metaGrey}\claimid{meat-core/lem:product-marginal}} & \pageref{lem:product-marginal}\\
\addlinespace[2pt]
\mref{Theorem}{thm:accumulated-entropy} & Accumulated entropy bound under adaptive attacks\newline{\footnotesize\color{metaGrey}\claimid{meat-core/thm:accumulated-entropy}} & \pageref{thm:accumulated-entropy}\\
\addlinespace[2pt]
\mref{Remark}{rem:accumulation-scope} & Scope of the accumulated entropy bound\newline{\footnotesize\color{metaGrey}\claimid{meat-core/rem:accumulation-scope}} & \pageref{rem:accumulation-scope}\\
\addlinespace[2pt]
\mref{Remark}{rem:no-constraints} & Where declared constraints may enter the rate\newline{\footnotesize\color{metaGrey}\claimid{meat-core/rem:no-constraints}} & \pageref{rem:no-constraints}\\
\addlinespace[2pt]
\mref{Theorem}{thm:constrained-accumulation} & Constrained accumulation under the attack-set licence\newline{\footnotesize\color{metaGrey}\claimid{meat-core/thm:constrained-accumulation}} & \pageref{thm:constrained-accumulation}\\
\addlinespace[2pt]
\mref{Remark}{rem:tagged-entries-and-accumulation} & Party-local entries and constrained accumulation\newline{\footnotesize\color{metaGrey}\claimid{meat-core/rem:tagged-entries-and-accumulation}} & \pageref{rem:tagged-entries-and-accumulation}\\
\addlinespace[2pt]
\mref{Example}{ex:fss} & The flag-state squasher as a licensed constraint list\newline{\footnotesize\color{metaGrey}\claimid{meat-core/ex:fss}} & \pageref{ex:fss}\\
\addlinespace[2pt]
\mref{Remark}{rem:homog} & Canonicity of the licensed class\newline{\footnotesize\color{metaGrey}\claimid{meat-core/rem:homog}} & \pageref{rem:homog}\\
\addlinespace[2pt]
\mref{Remark}{rem:scalar-declaration} & Scalar declaration of a homogeneous right-hand side\newline{\footnotesize\color{metaGrey}\claimid{meat-core/rem:scalar-declaration}} & \pageref{rem:scalar-declaration}\\
\addlinespace[2pt]
\mref{Remark}{rem:not-licensed} & Limits of the attack-set licence\newline{\footnotesize\color{metaGrey}\claimid{meat-core/rem:not-licensed}} & \pageref{rem:not-licensed}\\
\addlinespace[2pt]
\mref{Theorem}{thm:branch-formula} & Branch entropy from the protocol marginal\newline{\footnotesize\color{metaGrey}\claimid{meat-core/thm:branch-formula}} & \pageref{thm:branch-formula}\\
\addlinespace[2pt]
\mref{Corollary}{cor:concavity} & Concavity of the branch quantities on subnormalised states\newline{\footnotesize\color{metaGrey}\claimid{meat-core/cor:concavity}} & \pageref{cor:concavity}\\
\addlinespace[2pt]
\mref{Lemma}{lem:cone-concavity} & Concavity on the positive cone\newline{\footnotesize\color{metaGrey}\claimid{meat-core/lem:cone-concavity}} & \pageref{lem:cone-concavity}\\
\addlinespace[2pt]
\mref{Corollary}{cor:f-nonpos} & Non-positivity of the sector deficits\newline{\footnotesize\color{metaGrey}\claimid{meat-core/cor:f-nonpos}} & \pageref{cor:f-nonpos}\\
\addlinespace[2pt]
\mref{Definition}{def:keyinert} & Key-inert pair\newline{\footnotesize\color{metaGrey}\claimid{meat-core/def:keyinert}} & \pageref{def:keyinert}\\
\addlinespace[2pt]
\mref{Proposition}{prop:exact-cancellation} & Exact cancellation of key-inert sectors\newline{\footnotesize\color{metaGrey}\claimid{meat-core/prop:exact-cancellation}} & \pageref{prop:exact-cancellation}\\
\addlinespace[2pt]
\mref{Definition}{def:keyactive} & Key-active index set\newline{\footnotesize\color{metaGrey}\claimid{meat-core/def:keyactive}} & \pageref{def:keyactive}\\
\addlinespace[2pt]
\mref{Definition}{def:joint-sectors} & Joint-sector deficits and the joint key-active index set\newline{\footnotesize\color{metaGrey}\claimid{meat-core/def:joint-sectors}} & \pageref{def:joint-sectors}\\
\addlinespace[2pt]
\mref{Lemma}{lem:joint-sectors} & The objective split over joint sectors\newline{\footnotesize\color{metaGrey}\claimid{meat-core/lem:joint-sectors}} & \pageref{lem:joint-sectors}\\
\addlinespace[2pt]
\mref{Definition}{def:feasible} & Feasible set and computable objective\newline{\footnotesize\color{metaGrey}\claimid{meat-core/def:feasible}} & \pageref{def:feasible}\\
\addlinespace[2pt]
\mref{Proposition}{prop:trivial-alice-system} & The two source modes coincide at a trivial Alice system\newline{\footnotesize\color{metaGrey}\claimid{meat-core/prop:trivial-alice-system}} & \pageref{prop:trivial-alice-system}\\
\addlinespace[2pt]
\mref{Remark}{rem:modes-da1} & The two source modes on the program\newline{\footnotesize\color{metaGrey}\claimid{meat-core/rem:modes-da1}} & \pageref{rem:modes-da1}\\
\addlinespace[2pt]
\mref{Theorem}{thm:reduction} & Reduction of the rate to the convex program\newline{\footnotesize\color{metaGrey}\claimid{meat-core/thm:reduction}} & \pageref{thm:reduction}\\
\addlinespace[2pt]
\mref{Proposition}{prop:a-neutral} & Value neutrality along Alice's sectors\newline{\footnotesize\color{metaGrey}\claimid{meat-core/prop:a-neutral}} & \pageref{prop:a-neutral}\\
\addlinespace[2pt]
\mref{Corollary}{cor:reduction-constrained} & The reduction for arbitrary constraint lists\newline{\footnotesize\color{metaGrey}\claimid{meat-core/cor:reduction-constrained}} & \pageref{cor:reduction-constrained}\\
\addlinespace[2pt]
\mref{Remark}{rem:losses} & Where each step of the reduction loses\newline{\footnotesize\color{metaGrey}\claimid{meat-core/rem:losses}} & \pageref{rem:losses}\\
\addlinespace[2pt]
\mref{Proposition}{prop:value-range} & Range of the value and its finite entropy part\newline{\footnotesize\color{metaGrey}\claimid{meat-core/prop:value-range}} & \pageref{prop:value-range}\\
\addlinespace[2pt]
\mref{Remark}{rem:hlow-range} & Reading an infinite value\newline{\footnotesize\color{metaGrey}\claimid{meat-core/rem:hlow-range}} & \pageref{rem:hlow-range}\\
\addlinespace[2pt]
\mref{Theorem}{thm:convexity} & Joint convexity of the program\newline{\footnotesize\color{metaGrey}\claimid{meat-core/thm:convexity}} & \pageref{thm:convexity}\\
\addlinespace[2pt]
\mref{Definition}{def:classes} & Sector classes and the removal checklist\newline{\footnotesize\color{metaGrey}\claimid{meat-core/def:classes}} & \pageref{def:classes}\\
\addlinespace[2pt]
\mref{Theorem}{thm:removal} & Sound removal of a state-blind sector\newline{\footnotesize\color{metaGrey}\claimid{meat-core/thm:removal}} & \pageref{thm:removal}\\
\addlinespace[2pt]
\mref{Proposition}{prop:retention} & Unsoundness of dropping a constraint-active sector\newline{\footnotesize\color{metaGrey}\claimid{meat-core/prop:retention}} & \pageref{prop:retention}\\
\addlinespace[2pt]
\mref{Proposition}{prop:no-averaging} & Unsoundness of averaging sectors\newline{\footnotesize\color{metaGrey}\claimid{meat-core/prop:no-averaging}} & \pageref{prop:no-averaging}\\
\addlinespace[2pt]
\mref{Proposition}{prop:removal-compose} & Deleting a set of removable sectors\newline{\footnotesize\color{metaGrey}\claimid{meat-core/prop:removal-compose}} & \pageref{prop:removal-compose}\\
\addlinespace[2pt]
\mref{Remark}{rem:target-conditions} & The two index conditions of the removal checklist\newline{\footnotesize\color{metaGrey}\claimid{meat-core/rem:target-conditions}} & \pageref{rem:target-conditions}\\
\addlinespace[2pt]
\mref{Remark}{rem:witness-exactness} & Exactness of the exhibited weights and states\newline{\footnotesize\color{metaGrey}\claimid{meat-core/rem:witness-exactness}} & \pageref{rem:witness-exactness}\\
\addlinespace[2pt]
\mref{Remark}{rem:removal-general-format} & Removal in the general format\newline{\footnotesize\color{metaGrey}\claimid{meat-core/rem:removal-general-format}} & \pageref{rem:removal-general-format}\\
\addlinespace[2pt]
\mref{Remark}{rem:removal-sets} & The set hypothesis of the composition rule\newline{\footnotesize\color{metaGrey}\claimid{meat-core/rem:removal-sets}} & \pageref{rem:removal-sets}\\
\addlinespace[2pt]
\mref{Remark}{rem:classes} & Retention and deletion of sectors\newline{\footnotesize\color{metaGrey}\claimid{meat-core/rem:classes}} & \pageref{rem:classes}\\
\addlinespace[2pt]
\mref{Lemma}{lem:admissible-alice-subspace} & The admissible Alice subspace\newline{\footnotesize\color{metaGrey}\claimid{meat-core/lem:admissible-alice-subspace}} & \pageref{lem:admissible-alice-subspace}\\
\addlinespace[2pt]
\mref{Definition}{def:support-frame} & Support frame and declared isometry\newline{\footnotesize\color{metaGrey}\claimid{meat-core/def:support-frame}} & \pageref{def:support-frame}\\
\addlinespace[2pt]
\mref{Lemma}{lem:choi-marginal} & Fixed-marginal parametrisation on a support frame\newline{\footnotesize\color{metaGrey}\claimid{meat-core/lem:choi-marginal}} & \pageref{lem:choi-marginal}\\
\addlinespace[2pt]
\mref{Lemma}{lem:choi-factor} & Fixed-marginal parametrisation with any factor\newline{\footnotesize\color{metaGrey}\claimid{meat-core/lem:choi-factor}} & \pageref{lem:choi-factor}\\
\addlinespace[2pt]
\mref{Theorem}{thm:choi} & The Choi program\newline{\footnotesize\color{metaGrey}\claimid{meat-core/thm:choi}} & \pageref{thm:choi}\\
\addlinespace[2pt]
\mref{Proposition}{prop:frame-free-rows} & Frame-free constraint rows on a tight frame\newline{\footnotesize\color{metaGrey}\claimid{meat-core/prop:frame-free-rows}} & \pageref{prop:frame-free-rows}\\
\addlinespace[2pt]
\mref{Remark}{rem:tightness-by-construction} & Tightness read off the construction of the frame\newline{\footnotesize\color{metaGrey}\claimid{meat-core/rem:tightness-by-construction}} & \pageref{rem:tightness-by-construction}\\
\addlinespace[2pt]
\mref{Remark}{rem:ordering} & Ordering convention and channel reading\newline{\footnotesize\color{metaGrey}\claimid{meat-core/rem:ordering}} & \pageref{rem:ordering}\\
\addlinespace[2pt]
\mref{Remark}{rem:basis-invariance} & Invariance under the declared basis of the Choi input\newline{\footnotesize\color{metaGrey}\claimid{meat-core/rem:basis-invariance}} & \pageref{rem:basis-invariance}\\
\addlinespace[2pt]
\mref{Definition}{def:untrusted-choi} & The program in untrusted-source mode\newline{\footnotesize\color{metaGrey}\claimid{meat-core/def:untrusted-choi}} & \pageref{def:untrusted-choi}\\
\addlinespace[2pt]
\mref{Lemma}{lem:estimate} & Upper estimate of the value at any feasible point\newline{\footnotesize\color{metaGrey}\claimid{meat-core/lem:estimate}} & \pageref{lem:estimate}\\
\addlinespace[2pt]
\mref{Definition}{def:branch-support} & Maximal branch support\newline{\footnotesize\color{metaGrey}\claimid{meat-core/def:branch-support}} & \pageref{def:branch-support}\\
\addlinespace[2pt]
\mref{Lemma}{lem:branch-support} & Properties of the maximal branch support\newline{\footnotesize\color{metaGrey}\claimid{meat-core/lem:branch-support}} & \pageref{lem:branch-support}\\
\addlinespace[2pt]
\mref{Lemma}{lem:dk} & Divided-difference derivative of a matrix function\newline{\footnotesize\color{metaGrey}\claimid{meat-core/lem:dk}} & \pageref{lem:dk}\\
\addlinespace[2pt]
\mref{Theorem}{thm:gradient} & Gradient of the pinched trace power on a face\newline{\footnotesize\color{metaGrey}\claimid{meat-core/thm:gradient}} & \pageref{thm:gradient}\\
\addlinespace[2pt]
\mref{Corollary}{cor:full-gradient} & Gradient of the objective along feasible directions\newline{\footnotesize\color{metaGrey}\claimid{meat-core/cor:full-gradient}} & \pageref{cor:full-gradient}\\
\addlinespace[2pt]
\mref{Lemma}{lem:kraus-freedom} & Kraus freedom in the branch maps\newline{\footnotesize\color{metaGrey}\claimid{meat-core/lem:kraus-freedom}} & \pageref{lem:kraus-freedom}\\
\addlinespace[2pt]
\mref{Remark}{rem:directional-differentiability} & Scope of the directional differentiability of the objective\newline{\footnotesize\color{metaGrey}\claimid{meat-core/rem:directional-differentiability}} & \pageref{rem:directional-differentiability}\\
\addlinespace[2pt]
\mref{Lemma}{lem:b-derivative} & Derivative of the pinched trace power at order one\newline{\footnotesize\color{metaGrey}\claimid{meat-core/lem:b-derivative}} & \pageref{lem:b-derivative}\\
\addlinespace[2pt]
\mref{Theorem}{thm:oracle} & Limit of the value as the Rényi order tends to one\newline{\footnotesize\color{metaGrey}\claimid{meat-core/thm:oracle}} & \pageref{thm:oracle}\\
\addlinespace[2pt]
\mref{Remark}{rem:oracle} & Validation against the order-one limit\newline{\footnotesize\color{metaGrey}\claimid{meat-core/rem:oracle}} & \pageref{rem:oracle}\\
\addlinespace[2pt]
\mref{Lemma}{lem:rank-one-sector} & Linear deficit on a rank-one Alice sector\newline{\footnotesize\color{metaGrey}\claimid{meat-core/lem:rank-one-sector}} & \pageref{lem:rank-one-sector}\\
\addlinespace[2pt]
\mref{Definition}{def:dropped} & Dropped set and dropped program\newline{\footnotesize\color{metaGrey}\claimid{meat-core/def:dropped}} & \pageref{def:dropped}\\
\addlinespace[2pt]
\mref{Proposition}{prop:drop-bound} & The dropped program bounds the value from below\newline{\footnotesize\color{metaGrey}\claimid{meat-core/prop:drop-bound}} & \pageref{prop:drop-bound}\\
\addlinespace[2pt]
\mref{Corollary}{cor:drop-program} & Convexity, gradient and limit of the dropped program\newline{\footnotesize\color{metaGrey}\claimid{meat-core/cor:drop-program}} & \pageref{cor:drop-program}\\
\addlinespace[2pt]
\mref{Proposition}{prop:minimal-divergence} & The minimal divergence as the value with every pair dropped\newline{\footnotesize\color{metaGrey}\claimid{meat-core/prop:minimal-divergence}} & \pageref{prop:minimal-divergence}\\
\addlinespace[2pt]
\mref{Lemma}{lem:objective-homogeneous} & Positive homogeneity of the objective\newline{\footnotesize\color{metaGrey}\claimid{meat-core/lem:objective-homogeneous}} & \pageref{lem:objective-homogeneous}\\
\addlinespace[2pt]
\mref{Definition}{def:postprocessing-model} & Fixed-length post-processing model\newline{\footnotesize\color{metaGrey}\claimid{meat-core/def:postprocessing-model}} & \pageref{def:postprocessing-model}\\
\addlinespace[2pt]
\mref{Definition}{def:authenticated-channel} & Authenticated classical channel and authentication post-processing\newline{\footnotesize\color{metaGrey}\claimid{meat-core/def:authenticated-channel}} & \pageref{def:authenticated-channel}\\
\addlinespace[2pt]
\mref{Definition}{def:execution-end} & The end of an execution\newline{\footnotesize\color{metaGrey}\claimid{meat-core/def:execution-end}} & \pageref{def:execution-end}\\
\addlinespace[2pt]
\mref{Lemma}{lem:ev-correctness} & Correctness from error verification\newline{\footnotesize\color{metaGrey}\claimid{meat-core/lem:ev-correctness}} & \pageref{lem:ev-correctness}\\
\addlinespace[2pt]
\mref{Lemma}{lem:memory-lift} & Approximating an attack by all-at-once attacks with finite memory\newline{\footnotesize\color{metaGrey}\claimid{meat-core/lem:memory-lift}} & \pageref{lem:memory-lift}\\
\addlinespace[2pt]
\mref{Lemma}{lem:honest-attacks} & Adversaries of the honest-authentication setting as round attacks\newline{\footnotesize\color{metaGrey}\claimid{meat-core/lem:honest-attacks}} & \pageref{lem:honest-attacks}\\
\addlinespace[2pt]
\mref{Theorem}{thm:keylength} & Secrecy of the key from the accumulated entropy\newline{\footnotesize\color{metaGrey}\claimid{meat-core/thm:keylength}} & \pageref{thm:keylength}\\
\addlinespace[2pt]
\mref{Corollary}{cor:composed-security} & Security of the protocol and its net key\newline{\footnotesize\color{metaGrey}\claimid{meat-core/cor:composed-security}} & \pageref{cor:composed-security}\\
\addlinespace[2pt]
\mref{Corollary}{cor:separate-executions} & Security parameters of separate executions add\newline{\footnotesize\color{metaGrey}\claimid{meat-core/cor:separate-executions}} & \pageref{cor:separate-executions}\\
\addlinespace[2pt]
\mref{Remark}{rem:repeated-executions} & One execution's security parameter and a reused key\newline{\footnotesize\color{metaGrey}\claimid{meat-core/rem:repeated-executions}} & \pageref{rem:repeated-executions}\\
\addlinespace[2pt]
\mref{Remark}{rem:completeness-scope} & Completeness covers the acceptance test only\newline{\footnotesize\color{metaGrey}\claimid{meat-core/rem:completeness-scope}} & \pageref{rem:completeness-scope}\\
\addlinespace[2pt]
\mref{Proposition}{prop:model-pass-bound} & The pass probability of devices the declared model admits\newline{\footnotesize\color{metaGrey}\claimid{meat-core/prop:model-pass-bound}} & \pageref{prop:model-pass-bound}\\
\addlinespace[2pt]
\mref{Definition}{def:max-key-length} & Maximal key length\newline{\footnotesize\color{metaGrey}\claimid{meat-core/def:max-key-length}} & \pageref{def:max-key-length}\\
\addlinespace[2pt]
\mref{Remark}{rem:plus2} & The constant in the key length and the source's convention\newline{\footnotesize\color{metaGrey}\claimid{meat-core/rem:plus2}} & \pageref{rem:plus2}\\
\addlinespace[2pt]
\mref{Remark}{rem:hash-penalty} & One further bit for a merely two-universal family\newline{\footnotesize\color{metaGrey}\claimid{meat-core/rem:hash-penalty}} & \pageref{rem:hash-penalty}\\
\addlinespace[2pt]
\mref{Remark}{rem:alpha-domain} & The Rényi order of the key-length statements\newline{\footnotesize\color{metaGrey}\claimid{meat-core/rem:alpha-domain}} & \pageref{rem:alpha-domain}\\
\addlinespace[2pt]
\mref{Proposition}{prop:monotone} & Monotonicity of the maximal key length\newline{\footnotesize\color{metaGrey}\claimid{meat-core/prop:monotone}} & \pageref{prop:monotone}\\
\addlinespace[2pt]
\mref{Corollary}{cor:keylength-bound} & Key length from any lower bound on the value\newline{\footnotesize\color{metaGrey}\claimid{meat-core/cor:keylength-bound}} & \pageref{cor:keylength-bound}\\
\addlinespace[2pt]
\mref{Corollary}{cor:constrained-key-length} & Key length from the constrained value\newline{\footnotesize\color{metaGrey}\claimid{meat-core/cor:constrained-key-length}} & \pageref{cor:constrained-key-length}\\
\addlinespace[2pt]
\mref{Remark}{rem:cgate} & Degeneracy of the licensed-class half of the gate\newline{\footnotesize\color{metaGrey}\claimid{meat-core/rem:cgate}} & \pageref{rem:cgate}\\
\addlinespace[2pt]
\mref{Proposition}{prop:optol-nobound} & Instability of the rate under measurement perturbations\newline{\footnotesize\color{metaGrey}\claimid{meat-core/prop:optol-nobound}} & \pageref{prop:optol-nobound}\\
\addlinespace[2pt]
\mref{Remark}{rem:representability-ceiling} & The representability ceiling\newline{\footnotesize\color{metaGrey}\claimid{meat-core/rem:representability-ceiling}} & \pageref{rem:representability-ceiling}\\
\addlinespace[2pt]
\mref{Remark}{rem:optol-direction} & Scope of the perturbation gap\newline{\footnotesize\color{metaGrey}\claimid{meat-core/rem:optol-direction}} & \pageref{rem:optol-direction}\\
\addlinespace[2pt]
\mref{Remark}{rem:declared-measurement} & The declared measurement as a premise of the rate\newline{\footnotesize\color{metaGrey}\claimid{meat-core/rem:declared-measurement}} & \pageref{rem:declared-measurement}\\
\addlinespace[2pt]
\mref{Proposition}{prop:reference-companion} & The compiled reference companion is an admissible declaration\newline{\footnotesize\color{metaGrey}\claimid{meat-core/prop:reference-companion}} & \pageref{prop:reference-companion}\\
\addlinespace[2pt]
\mref{Corollary}{cor:reference-companion-factors} & Factor enclosures of the compiled reference companion\newline{\footnotesize\color{metaGrey}\claimid{meat-core/cor:reference-companion-factors}} & \pageref{cor:reference-companion-factors}\\
\addlinespace[2pt]
\mref{Assumption}{ass:keyreg} & Dedicated key register and one-way reconciliation\newline{\footnotesize\color{metaGrey}\claimid{meat-core/ass:keyreg}} & \pageref{ass:keyreg}\\
\addlinespace[2pt]
\mref{Assumption}{ass:hash} & Ideal two-universal privacy amplification\newline{\footnotesize\color{metaGrey}\claimid{meat-core/ass:hash}} & \pageref{ass:hash}\\
\addlinespace[2pt]
\mref{Assumption}{ass:seed} & Private, independent and fresh local randomness\newline{\footnotesize\color{metaGrey}\claimid{meat-core/ass:seed}} & \pageref{ass:seed}\\
\addlinespace[2pt]
\mref{Assumption}{ass:fixed} & Protocol parameters fixed before the run\newline{\footnotesize\color{metaGrey}\claimid{meat-core/ass:fixed}} & \pageref{ass:fixed}\\
\addlinespace[2pt]
\mref{Assumption}{ass:conform} & The implementation runs the analysed post-processing\newline{\footnotesize\color{metaGrey}\claimid{meat-core/ass:conform}} & \pageref{ass:conform}\\
\addlinespace[2pt]
\mref{Assumption}{ass:auth} & Authentication of the classical channel by a keyed scheme\newline{\footnotesize\color{metaGrey}\claimid{meat-core/ass:auth}} & \pageref{ass:auth}\\
\addlinespace[2pt]
\mref{Assumption}{ass:separate-executions} & Executions that share no key and no device state\newline{\footnotesize\color{metaGrey}\claimid{meat-core/ass:separate-executions}} & \pageref{ass:separate-executions}\\
\addlinespace[2pt]
\mref{Assumption}{ass:lab} & Each party's laboratory is sealed\newline{\footnotesize\color{metaGrey}\claimid{meat-core/ass:lab}} & \pageref{ass:lab}\\
\addlinespace[2pt]
\mref{Assumption}{ass:eve-dim} & The adversary's registers are finite-dimensional\newline{\footnotesize\color{metaGrey}\claimid{meat-core/ass:eve-dim}} & \pageref{ass:eve-dim}\\
\addlinespace[2pt]
\mref{Requirement}{req:exact-admission} & Exact admission of declared data\newline{\footnotesize\color{metaGrey}\claimid{meat-core/req:exact-admission}} & \pageref{req:exact-admission}\\
\addlinespace[2pt]
\mref{Requirement}{req:declaration-admission} & Admission of the declaration\newline{\footnotesize\color{metaGrey}\claimid{meat-core/req:declaration-admission}} & \pageref{req:declaration-admission}\\
\addlinespace[2pt]
\mref{Requirement}{req:companion-admission} & Admission of the reference companion beside its declaration\newline{\footnotesize\color{metaGrey}\claimid{meat-core/req:companion-admission}} & \pageref{req:companion-admission}\\
\addlinespace[2pt]
\mref{Requirement}{req:acceptance-set-admission} & Admission of the acceptance set\newline{\footnotesize\color{metaGrey}\claimid{meat-core/req:acceptance-set-admission}} & \pageref{req:acceptance-set-admission}\\
\addlinespace[2pt]
\mref{Requirement}{req:branch-register-components} & Register components of an implemented branch map\newline{\footnotesize\color{metaGrey}\claimid{meat-core/req:branch-register-components}} & \pageref{req:branch-register-components}\\
\addlinespace[2pt]
\mref{Requirement}{req:sector-report} & The sector report\newline{\footnotesize\color{metaGrey}\claimid{meat-core/req:sector-report}} & \pageref{req:sector-report}\\
\addlinespace[2pt]
\mref{Requirement}{req:state-blind-test} & Bit-level test of state-blindness\newline{\footnotesize\color{metaGrey}\claimid{meat-core/req:state-blind-test}} & \pageref{req:state-blind-test}\\
\addlinespace[2pt]
\mref{Requirement}{req:removal-checklist} & Setting the removable flag\newline{\footnotesize\color{metaGrey}\claimid{meat-core/req:removal-checklist}} & \pageref{req:removal-checklist}\\
\addlinespace[2pt]
\mref{Requirement}{req:structure-operations} & Operations on the declared sectors\newline{\footnotesize\color{metaGrey}\claimid{meat-core/req:structure-operations}} & \pageref{req:structure-operations}\\
\addlinespace[2pt]
\mref{Requirement}{req:single-sector-deletion} & At most one deleted sector without the composition check\newline{\footnotesize\color{metaGrey}\claimid{meat-core/req:single-sector-deletion}} & \pageref{req:single-sector-deletion}\\
\addlinespace[2pt]
\mref{Requirement}{req:frame-constructed} & Construction and exact check of the support frame\newline{\footnotesize\color{metaGrey}\claimid{meat-core/req:frame-constructed}} & \pageref{req:frame-constructed}\\
\addlinespace[2pt]
\mref{Requirement}{req:diagonal-frame-test} & Exact diagonality test for the diagonal frame\newline{\footnotesize\color{metaGrey}\claimid{meat-core/req:diagonal-frame-test}} & \pageref{req:diagonal-frame-test}\\
\addlinespace[2pt]
\mref{Requirement}{req:frame-recorded} & Record of the support frame\newline{\footnotesize\color{metaGrey}\claimid{meat-core/req:frame-recorded}} & \pageref{req:frame-recorded}\\
\addlinespace[2pt]
\mref{Requirement}{req:choi-convention} & Factor order and transpose basis of the Choi blocks\newline{\footnotesize\color{metaGrey}\claimid{meat-core/req:choi-convention}} & \pageref{req:choi-convention}\\
\addlinespace[2pt]
\mref{Requirement}{req:row-form-by-construction} & Frame-free rows only where the construction exhibits tightness\newline{\footnotesize\color{metaGrey}\claimid{meat-core/req:row-form-by-construction}} & \pageref{req:row-form-by-construction}\\
\addlinespace[2pt]
\mref{Requirement}{req:row-form-record} & Record of the form each constraint row was imposed in\newline{\footnotesize\color{metaGrey}\claimid{meat-core/req:row-form-record}} & \pageref{req:row-form-record}\\
\addlinespace[2pt]
\mref{Requirement}{req:split-form-gradient} & Split-form evaluation of the trace-power gradient\newline{\footnotesize\color{metaGrey}\claimid{meat-core/req:split-form-gradient}} & \pageref{req:split-form-gradient}\\
\addlinespace[2pt]
\mref{Requirement}{req:gradient-at-maximal-support} & Gradient evaluation at maximal support\newline{\footnotesize\color{metaGrey}\claimid{meat-core/req:gradient-at-maximal-support}} & \pageref{req:gradient-at-maximal-support}\\
\addlinespace[2pt]
\mref{Requirement}{req:limit-not-a-bound} & The order-one limit as a validation reference only\newline{\footnotesize\color{metaGrey}\claimid{meat-core/req:limit-not-a-bound}} & \pageref{req:limit-not-a-bound}\\
\addlinespace[2pt]
\mref{Requirement}{req:dropped-set-record} & Recording the dropped set\newline{\footnotesize\color{metaGrey}\claimid{meat-core/req:dropped-set-record}} & \pageref{req:dropped-set-record}\\
\addlinespace[2pt]
\mref{Requirement}{req:infinite-value} & Propagation of an infinite value\newline{\footnotesize\color{metaGrey}\claimid{meat-core/req:infinite-value}} & \pageref{req:infinite-value}\\
\addlinespace[2pt]
\mref{Requirement}{req:unreachable-acceptance-set} & The check of the acceptance set against the feasible statistics\newline{\footnotesize\color{metaGrey}\claimid{meat-core/req:unreachable-acceptance-set}} & \pageref{req:unreachable-acceptance-set}\\
\addlinespace[2pt]
\mref{Requirement}{req:objective-alpha-domain} & Rényi order of an evaluated objective\newline{\footnotesize\color{metaGrey}\claimid{meat-core/req:objective-alpha-domain}} & \pageref{req:objective-alpha-domain}\\
\addlinespace[2pt]
\mref{Requirement}{req:reported-key-length} & Reported key length\newline{\footnotesize\color{metaGrey}\claimid{meat-core/req:reported-key-length}} & \pageref{req:reported-key-length}\\
\addlinespace[2pt]
\mref{Requirement}{req:hash-penalty} & Key-length penalty for a merely two-universal family\newline{\footnotesize\color{metaGrey}\claimid{meat-core/req:hash-penalty}} & \pageref{req:hash-penalty}\\
\addlinespace[2pt]
\mref{Requirement}{req:key-length-alpha-domain} & Rényi order of a reported key length\newline{\footnotesize\color{metaGrey}\claimid{meat-core/req:key-length-alpha-domain}} & \pageref{req:key-length-alpha-domain}\\
\addlinespace[2pt]
\mref{Requirement}{req:compiled-declaration-only} & Key lengths only on declarations the compiler emitted\newline{\footnotesize\color{metaGrey}\claimid{meat-core/req:compiled-declaration-only}} & \pageref{req:compiled-declaration-only}\\
\addlinespace[2pt]
\mref{Requirement}{req:security-statement-not-vacuous} & Refusal of budgets with an empty security statement\newline{\footnotesize\color{metaGrey}\claimid{meat-core/req:security-statement-not-vacuous}} & \pageref{req:security-statement-not-vacuous}\\
\addlinespace[2pt]
\mref{Requirement}{req:secrecy-convention} & Declared choice of the source's secrecy convention\newline{\footnotesize\color{metaGrey}\claimid{meat-core/req:secrecy-convention}} & \pageref{req:secrecy-convention}\\
\addlinespace[2pt]
\mref{Requirement}{req:constrained-key-gate} & Two-part gate on a key length with declared constraints\newline{\footnotesize\color{metaGrey}\claimid{meat-core/req:constrained-key-gate}} & \pageref{req:constrained-key-gate}\\
\addlinespace[2pt]
\mref{Requirement}{req:degenerate-gate-report} & A degenerate licensed-class check reported as no evidence\newline{\footnotesize\color{metaGrey}\claimid{meat-core/req:degenerate-gate-report}} & \pageref{req:degenerate-gate-report}\\
\addlinespace[2pt]
\mref{Requirement}{req:marginal-licence-gate} & Source licence for a declared marginal\newline{\footnotesize\color{metaGrey}\claimid{meat-core/req:marginal-licence-gate}} & \pageref{req:marginal-licence-gate}\\
\addlinespace[2pt]
\mref{Definition}{def:relaxed-feasible-set} & The relaxed feasible set\newline{\footnotesize\color{metaGrey}\claimid{meat-core/frank-wolfe/def:relaxed-feasible-set}} & \pageref{def:relaxed-feasible-set}\\
\addlinespace[2pt]
\mref{Proposition}{prop:relaxed-set-properties} & Properties of the relaxed feasible set\newline{\footnotesize\color{metaGrey}\claimid{meat-core/frank-wolfe/prop:relaxed-set-properties}} & \pageref{prop:relaxed-set-properties}\\
\addlinespace[2pt]
\mref{Remark}{rem:feas-tol} & Scope of reading the feasible set at a tolerance\newline{\footnotesize\color{metaGrey}\claimid{meat-core/frank-wolfe/rem:feas-tol}} & \pageref{rem:feas-tol}\\
\addlinespace[2pt]
\mref{Definition}{def:first-order-setting} & Ambient program, pairings, acceptance support and linear minimisation\newline{\footnotesize\color{metaGrey}\claimid{meat-core/frank-wolfe/def:first-order-setting}} & \pageref{def:first-order-setting}\\
\addlinespace[2pt]
\mref{Lemma}{lem:fo-khyp} & Positivity needed only on the acceptance support\newline{\footnotesize\color{metaGrey}\claimid{meat-core/frank-wolfe/lem:fo-khyp}} & \pageref{lem:fo-khyp}\\
\addlinespace[2pt]
\mref{Lemma}{lem:fo-face} & Every direction is admissible at maximal support\newline{\footnotesize\color{metaGrey}\claimid{meat-core/frank-wolfe/lem:fo-face}} & \pageref{lem:fo-face}\\
\addlinespace[2pt]
\mref{Theorem}{thm:firstorder} & The estimate and the first-order lower bound\newline{\footnotesize\color{metaGrey}\claimid{meat-core/frank-wolfe/thm:firstorder}} & \pageref{thm:firstorder}\\
\addlinespace[2pt]
\mref{Remark}{rem:iterate-relaxed-oracle-exact} & Iterate in the relaxed set, minima over the exact set\newline{\footnotesize\color{metaGrey}\claimid{meat-core/frank-wolfe/rem:iterate-relaxed-oracle-exact}} & \pageref{rem:iterate-relaxed-oracle-exact}\\
\addlinespace[2pt]
\mref{Remark}{rem:fo-danger} & Failure of maximal support lies on the unsafe side\newline{\footnotesize\color{metaGrey}\claimid{meat-core/frank-wolfe/rem:fo-danger}} & \pageref{rem:fo-danger}\\
\addlinespace[2pt]
\mref{Remark}{rem:fo-nopert} & A perturbed gradient bounds another function\newline{\footnotesize\color{metaGrey}\claimid{meat-core/frank-wolfe/rem:fo-nopert}} & \pageref{rem:fo-nopert}\\
\addlinespace[2pt]
\mref{Corollary}{cor:fo-keylength} & Key length from the first-order bound\newline{\footnotesize\color{metaGrey}\claimid{meat-core/frank-wolfe/cor:fo-keylength}} & \pageref{cor:fo-keylength}\\
\addlinespace[2pt]
\mref{Corollary}{cor:fo-dropped} & First-order bound for the dropped program\newline{\footnotesize\color{metaGrey}\claimid{meat-core/frank-wolfe/cor:fo-dropped}} & \pageref{cor:fo-dropped}\\
\addlinespace[2pt]
\mref{Corollary}{cor:fo-reference} & The first-order bound on the reference companion\newline{\footnotesize\color{metaGrey}\claimid{meat-core/frank-wolfe/cor:fo-reference}} & \pageref{cor:fo-reference}\\
\addlinespace[2pt]
\mref{Remark}{rem:fo-general-format} & The first-order bound in the general format\newline{\footnotesize\color{metaGrey}\claimid{meat-core/frank-wolfe/rem:fo-general-format}} & \pageref{rem:fo-general-format}\\
\addlinespace[2pt]
\mref{Lemma}{lem:feasible-side-dual} & A lower bound on the feasible-side minimum from any multipliers\newline{\footnotesize\color{metaGrey}\claimid{meat-core/frank-wolfe/lem:feasible-side-dual}} & \pageref{lem:feasible-side-dual}\\
\addlinespace[2pt]
\mref{Lemma}{lem:certified-bracket} & The certified bracket from enclosures of the gradient\newline{\footnotesize\color{metaGrey}\claimid{meat-core/frank-wolfe/lem:certified-bracket}} & \pageref{lem:certified-bracket}\\
\addlinespace[2pt]
\mref{Lemma}{lem:certified-bracket-enclosed} & The certified bracket over a measured family known through enclosures\newline{\footnotesize\color{metaGrey}\claimid{meat-core/frank-wolfe/lem:certified-bracket-enclosed}} & \pageref{lem:certified-bracket-enclosed}\\
\addlinespace[2pt]
\mref{Corollary}{cor:certified-keylength} & Key length from the certified bracket\newline{\footnotesize\color{metaGrey}\claimid{meat-core/frank-wolfe/cor:certified-keylength}} & \pageref{cor:certified-keylength}\\
\addlinespace[2pt]
\mref{Corollary}{cor:certified-keylength-reference} & Key length on a reference companion from the certified bracket\newline{\footnotesize\color{metaGrey}\claimid{meat-core/frank-wolfe/cor:certified-keylength-reference}} & \pageref{cor:certified-keylength-reference}\\
\addlinespace[2pt]
\mref{Assumption}{ass:lmo} & Exact linear minimisation over the feasible set\newline{\footnotesize\color{metaGrey}\claimid{meat-core/frank-wolfe/ass:lmo}} & \pageref{ass:lmo}\\
\addlinespace[2pt]
\mref{Remark}{rem:acc-side} & Exact minimisation on the acceptance side\newline{\footnotesize\color{metaGrey}\claimid{meat-core/frank-wolfe/rem:acc-side}} & \pageref{rem:acc-side}\\
\addlinespace[2pt]
\mref{Remark}{rem:pairing-on-acceptance-support} & Pairing over the acceptance support\newline{\footnotesize\color{metaGrey}\claimid{meat-core/frank-wolfe/rem:pairing-on-acceptance-support}} & \pageref{rem:pairing-on-acceptance-support}\\
\addlinespace[2pt]
\mref{Remark}{rem:fo-arithmetic} & Exact real arithmetic in the first-order bound\newline{\footnotesize\color{metaGrey}\claimid{meat-core/frank-wolfe/rem:fo-arithmetic}} & \pageref{rem:fo-arithmetic}\\
\addlinespace[2pt]
\mref{Remark}{rem:certified-scope} & What a certified number rests on\newline{\footnotesize\color{metaGrey}\claimid{meat-core/frank-wolfe/rem:certified-scope}} & \pageref{rem:certified-scope}\\
\addlinespace[2pt]
\mref{Remark}{rem:fo-interface} & Two reported numbers and one first-order claim\newline{\footnotesize\color{metaGrey}\claimid{meat-core/frank-wolfe/rem:fo-interface}} & \pageref{rem:fo-interface}\\
\addlinespace[2pt]
\mref{Remark}{rem:fo-iterate-checks} & What the iterate checks establish\newline{\footnotesize\color{metaGrey}\claimid{meat-core/frank-wolfe/rem:fo-iterate-checks}} & \pageref{rem:fo-iterate-checks}\\
\addlinespace[2pt]
\mref{Remark}{rem:fo-returned-point} & What the measurement of the returned point establishes\newline{\footnotesize\color{metaGrey}\claimid{meat-core/frank-wolfe/rem:fo-returned-point}} & \pageref{rem:fo-returned-point}\\
\addlinespace[2pt]
\mref{Remark}{rem:fo-one-ruler} & Maximal support as a rank equality read with one ruler\newline{\footnotesize\color{metaGrey}\claimid{meat-core/frank-wolfe/rem:fo-one-ruler}} & \pageref{rem:fo-one-ruler}\\
\addlinespace[2pt]
\mref{Remark}{rem:fo-reachability} & Reachability of positivity and maximal support\newline{\footnotesize\color{metaGrey}\claimid{meat-core/frank-wolfe/rem:fo-reachability}} & \pageref{rem:fo-reachability}\\
\addlinespace[2pt]
\mref{Remark}{rem:fo-check-exposures} & One-sided exposures of the runtime checks\newline{\footnotesize\color{metaGrey}\claimid{meat-core/frank-wolfe/rem:fo-check-exposures}} & \pageref{rem:fo-check-exposures}\\
\addlinespace[2pt]
\mref{Remark}{rem:fo-oracle-repair} & A repaired oracle return as a step of the search\newline{\footnotesize\color{metaGrey}\claimid{meat-core/frank-wolfe/rem:fo-oracle-repair}} & \pageref{rem:fo-oracle-repair}\\
\addlinespace[2pt]
\mref{Remark}{rem:fo-iterate-repair} & A repaired iterate as a construction step\newline{\footnotesize\color{metaGrey}\claimid{meat-core/frank-wolfe/rem:fo-iterate-repair}} & \pageref{rem:fo-iterate-repair}\\
\addlinespace[2pt]
\mref{Remark}{rem:fo-negative-values} & Opposite treatment of a negative bound and a negative estimate\newline{\footnotesize\color{metaGrey}\claimid{meat-core/frank-wolfe/rem:fo-negative-values}} & \pageref{rem:fo-negative-values}\\
\addlinespace[2pt]
\mref{Requirement}{req:first-order-withdrawal} & Withdrawal of the first-order bound\newline{\footnotesize\color{metaGrey}\claimid{meat-core/frank-wolfe/req:first-order-withdrawal}} & \pageref{req:first-order-withdrawal}\\
\addlinespace[2pt]
\mref{Requirement}{req:returned-point-measured} & Measurement of the returned point\newline{\footnotesize\color{metaGrey}\claimid{meat-core/frank-wolfe/req:returned-point-measured}} & \pageref{req:returned-point-measured}\\
\addlinespace[2pt]
\mref{Requirement}{req:first-order-clamp} & Clamping of the first-order bound at zero\newline{\footnotesize\color{metaGrey}\claimid{meat-core/frank-wolfe/req:first-order-clamp}} & \pageref{req:first-order-clamp}\\
\addlinespace[2pt]
\mref{Requirement}{req:negative-estimate-refused} & Refusal of a negative estimate\newline{\footnotesize\color{metaGrey}\claimid{meat-core/frank-wolfe/req:negative-estimate-refused}} & \pageref{req:negative-estimate-refused}\\
\addlinespace[2pt]
\mref{Requirement}{req:bound-estimate-difference} & The recorded difference between the bound and the estimate\newline{\footnotesize\color{metaGrey}\claimid{meat-core/frank-wolfe/req:bound-estimate-difference}} & \pageref{req:bound-estimate-difference}\\
\addlinespace[2pt]
\mref{Requirement}{req:no-imported-penalty} & The program's own objective, with no imported penalty\newline{\footnotesize\color{metaGrey}\claimid{meat-core/frank-wolfe/req:no-imported-penalty}} & \pageref{req:no-imported-penalty}\\
\addlinespace[2pt]
\mref{Requirement}{req:bracket-on-reference} & The first-order bound evaluated on the reference companion\newline{\footnotesize\color{metaGrey}\claimid{meat-core/frank-wolfe/req:bracket-on-reference}} & \pageref{req:bracket-on-reference}\\
\addlinespace[2pt]
\mref{Requirement}{req:iterate-checks} & Feasibility checks at each iterate\newline{\footnotesize\color{metaGrey}\claimid{meat-core/frank-wolfe/req:iterate-checks}} & \pageref{req:iterate-checks}\\
\addlinespace[2pt]
\mref{Requirement}{req:support-rank-check} & Rank test for maximal support\newline{\footnotesize\color{metaGrey}\claimid{meat-core/frank-wolfe/req:support-rank-check}} & \pageref{req:support-rank-check}\\
\addlinespace[2pt]
\mref{Requirement}{req:oracle-repair} & Repair of an oracle return\newline{\footnotesize\color{metaGrey}\claimid{meat-core/frank-wolfe/req:oracle-repair}} & \pageref{req:oracle-repair}\\
\addlinespace[2pt]
\mref{Requirement}{req:iterate-repair} & Repair of a formed iterate\newline{\footnotesize\color{metaGrey}\claimid{meat-core/frank-wolfe/req:iterate-repair}} & \pageref{req:iterate-repair}\\
\addlinespace[2pt]
\mref{Requirement}{req:per-sector-coupling} & Coupling residual and record per Alice sector\newline{\footnotesize\color{metaGrey}\claimid{meat-core/frank-wolfe/req:per-sector-coupling}} & \pageref{req:per-sector-coupling}\\
\addlinespace[2pt]
\mref{Requirement}{req:acceptance-geometries} & The admitted acceptance geometries\newline{\footnotesize\color{metaGrey}\claimid{meat-core/frank-wolfe/req:acceptance-geometries}} & \pageref{req:acceptance-geometries}\\
\addlinespace[2pt]
\mref{Requirement}{req:acceptance-iterate-exact} & Exact formation of the acceptance-side iterate\newline{\footnotesize\color{metaGrey}\claimid{meat-core/frank-wolfe/req:acceptance-iterate-exact}} & \pageref{req:acceptance-iterate-exact}\\
\addlinespace[2pt]
\mref{Requirement}{req:bacc-recorded} & Record of the acceptance-side quantity in the bracket\newline{\footnotesize\color{metaGrey}\claimid{meat-core/frank-wolfe/req:bacc-recorded}} & \pageref{req:bacc-recorded}\\
\addlinespace[2pt]
\mref{Requirement}{req:verified-feasible-side} & The verified feasible side of every bracket\newline{\footnotesize\color{metaGrey}\claimid{meat-core/frank-wolfe/req:verified-feasible-side}} & \pageref{req:verified-feasible-side}\\
\addlinespace[2pt]
\mref{Requirement}{req:primal-pairing} & The primal pairing in the first-order bracket\newline{\footnotesize\color{metaGrey}\claimid{meat-core/frank-wolfe/req:primal-pairing}} & \pageref{req:primal-pairing}\\
\addlinespace[2pt]
\mref{Requirement}{req:certified-evaluation} & The certified evaluation of the first-order bound\newline{\footnotesize\color{metaGrey}\claimid{meat-core/frank-wolfe/req:certified-evaluation}} & \pageref{req:certified-evaluation}\\
\bottomrule
\end{xltabular}}

\clearpage
\bookmarksetupnext{level=part}
\printindex

\clearpage\phantomsection\bookmarksetupnext{level=part}%
\addcontentsline{toc}{section}{References}
\bibliographystyle{plain}
\bibliography{../../../refs/references}

\end{document}
